\documentclass[10pt,twoside,titlepage]{report}
\usepackage[utf8x]{inputenc}
\usepackage[fleqn]{amsmath}
\usepackage{array}
\usepackage{varwidth}
\usepackage{ltablex}
\usepackage{vmargin}
\usepackage{fancyhdr}
\usepackage[pdftex]{graphicx}
\usepackage[pdftex,pagebackref=true,unicode]{hyperref}
\pdfcompresslevel=9
\hypersetup{colorlinks=true,linkcolor=[rgb]{0.01176,0.38039,0.67058},urlcolor=[rgb]{0.01176,0.38039,0.67058}}

% Style
% -----
\parindent0mm
\parskip2ex plus1ex minus1ex
\frenchspacing
\sloppy
\setcounter{secnumdepth}{3}
\setcounter{tocdepth}{1}
\setpapersize{A4}
\setmarginsrb{2.5cm}{1cm}{2cm}{2cm}{8mm}{15mm}{5mm}{15mm}

\pagestyle{fancy}
\addtolength{\headwidth}{0.7cm}
\renewcommand{\chaptermark}[1]{\markboth{\slshape\thechapter.\ #1}{}}
\renewcommand{\sectionmark}[1]{\markright{\slshape\thesection.\ #1}}
\fancyhead[L,R]{\slshape\thepage}
\fancyhead[LO]{\slshape\nouppercase{\rightmark}}
\fancyhead[RE]{\slshape\nouppercase{\leftmark}}
\cfoot{}

\fancypagestyle{plain}{%
\pagestyle{fancy}
\fancyhead[RO]{\slshape\thepage}
\fancyhead[LE]{\slshape\thepage}
\fancyhead[LO]{}
\fancyhead[RE]{}
}

% Config table style
% ------------------
\newcolumntype{N}{>{\begin{varwidth}{6cm}}l<{\end{varwidth}}} % Name column (with max. width of 6cm)
\newcolumntype{T}{l}                                          % Type column
\newcolumntype{A}{X}                                          % Annotation column

% Commands
% --------
\newcommand{\config}[1]{\protect\includegraphics[width=0.9em]{element.pdf}\,\textbf{#1}}
\newcommand{\configClass}[2]{\protect\includegraphics[width=0.9em]{element.pdf}\,\hyperref[#2]{\textbf{#1}}}
\newcommand{\configFile}[2]{\protect\includegraphics[width=0.9em]{element.pdf}\,\hyperref[general.fileFormat:#2]{\textbf{#1}}}
\newcommand{\program}[1]{\protect\includegraphics[width=0.9em]{program.pdf}\,\hyperref[#1]{\textbf{#1}}}
\newcommand{\file}[2]{\protect\includegraphics[height=1.7ex]{file.pdf}\,\hyperref[general.fileFormat:#2]{#1}}
\newcommand{\reference}[2]{\hyperref[#2]{#1~(\ref{#2})}}
\newcommand{\M}[1]{\mathbf{#1}}

% \fig{position}{width as fraction of textwidth}{filename}{label}{caption}
% ------------------------------------------------------------------------
\newcommand{\fig}[5]{%
\begin{figure}[#1]%
\centering%
\includegraphics*[width=#2\textwidth]{#3}%
\caption{#5}\label{#4}%
\end{figure}%
}

%=========================================

\begin{document}
\hfuzz=5000pt
\vfuzz=5000pt
\hbadness=10000
\vbadness=10000

\graphicspath{{../figures/}}

%---------------------------- title page ------------------------------------
\thispagestyle{empty}
\begin{center}
\vspace*{6cm}
\begin{figure}[!h]
\centering
\includegraphics*[width=0.75\textwidth]{groops_banner}
\end{figure}
\vspace*{0.2cm}
{\huge\sf Gravity Recovery Object Oriented Programming System} \\
\vspace*{2cm}
{\huge\sf Documentation} \\
\vspace*{2cm}
{\Large\sf GitHub repository: \url{https://github.com/groops-devs/groops}}
\end{center}

\newpage
\clearpage{\pagestyle{empty}} %\cleardoublepage}%
%\pagenumbering{Roman}
\tableofcontents%
\clearpage{\pagestyle{empty}\cleardoublepage\clearpage}

\chapter{General information}
The Gravity Recovery Object Oriented Programming System (GROOPS) is a software toolkit written in C++
that enables the user to perform core geodetic tasks.
Key features of the software include gravity field recovery from satellite and terrestrial data,
the determination of satellite orbits from global navigation satellite system (GNSS) measurements,
and the processing of GNSS constellations and ground station networks. Most tasks and algorithms are
(optionally) parallelized through the Message Passing Interface (MPI), thus the software enables
a smooth transition from single-CPU desktop computers to large distributed
computing environments for resource intensive tasks.

Parts of GROOPS originate from developments in the Astronomical, Physical and Mathematical Geodesy Group
at the University of Bonn, Germany. Since 2010 it is developed and maintained at Graz University of Technology, Austria.
It is licensed under GPLv3 and hosted at GitHub, \url{https://github.com/groops-devs/groops}.

If you use data sets computed with GROOPS in a publication or publish the data itself, please cite our reference paper:\\
Mayer-Guerr, T., Behzadpour, S., Eicker, A., Ellmer, M., Koch, B., Krauss, S., Pock, C., Rieser, D., Strasser, S., Suesser-Rechberger,
B., Zehentner, N., Kvas, A. (2021). GROOPS: A software toolkit for gravity field recovery and GNSS processing.
Computers and Geosciences, 104864. \url{https://doi.org/10.1016/j.cageo.2021.104864}.

GROOPS depends on data files such as Earth rotation, Love numbers, and GNSS meta information.
An initial data set that is regularly updated is available on our FTP server,
\url{https://ftp.tugraz.at/outgoing/ITSG/groops}. You can choose between downloading the data directory
or a single \href{https://ftp.tugraz.at/outgoing/ITSG/groops/data.zip}{data.zip} with the same content.

\section{Config files}\label{general.configFiles}
GROOPS is controlled by XML configuration files. One or more configuration files
must be passed as arguments to GROOPS:
\begin{verbatim}
  groops config1.xml config2.xml [...]
\end{verbatim}
These files can be created with the graphical user interface program \verb|groopsGui|
in a convenient way (see section~\reference{GUI}{general.gui}).
A complete formal (computer readable) description of a configuration file
in the form of an XSD schema file can be created with the command
\begin{verbatim}
  groops --xsd groops.xsd
\end{verbatim}

A configuration file consists of a list of \reference{programs}{programType}
that are executed in sequential order. Each program comes with its own config options
and they work independently without any internal communication between programs.
Data flow between programs is realized via files. An \config{outputfile} of one program can serve as
an \config{inputfile} for the next program.
Most programs are deliberately kept small and focused on a specific task. This modularity combined with the
general purpose design of many programs enables the creation of complex workflows with little effort.
Including \reference{loops and conditions}{general.loopsAndConditions} in a config file provides even more flexibility.

Individual programs (and also other optional config elements) can be disabled
and are ignored during execution. Mandatory config elements are indicated by a star (\verb|*|).
Empty optional elements are ignored or a meaningful default value is assumed.

The elements of a configuration file can be one of the following basic data types:
\begin{itemize}
\item \verb|int|: integer number
\item \verb|uint|: unsigned integer number
\item \verb|double|: floating point number
\item \verb|angle|: given in degree
\item \verb|time|: given in modified Julian date (MJD)
\item \verb|boolean|: 0: false, 1: true
\item \verb|string|: text
\item \verb|filename|: absolute path to a file or path relative to the working directory
\item \verb|expression|: numerical expression evaluated during execution
\item \verb|doodson|: Doodson number or Darwin's name of a tidal frequency
\item \verb|gnssType|: GNSS observation type according to the RINEX 3 definition
\end{itemize}
The first 5 data types also allow numerical expressions as input in addition to pure numbers.
In addition to these basic types, there are a large number of complex data types called classes,
which are described in section~\reference{Classes}{classes}.

\subsection{Variables}\label{general.configFiles:variables}
In addition to programs a config file can also include elements called variables.
These elements are comparable to read-only variables in programming and can be referenced from any program and config element.
This can be done by either linking an element directly to a variable or
by using the name as a variable in an expression of an input field (see section~\reference{Parsers and variables}{general.parser}).
While elements can only be directly linked to variables of the same type, this also supports complex data types
such as \configClass{gravityfield}{gravityfieldType}. Thus it is possible to, for example, define a reference gravity field once
in the global section and use it multiple times in different programs.

Variables can be declared anywhere in the configuration file. Variables in locations other than the global section
have a local scope and hide global variables or variables from a hierarchy level above. They are valid after
declaration until the end of the hierarchy level is reached or a new variable with the same name is declared.

Variables are not evaluated directly when they are declared, but only later when they are used in a config element.
This means, for example, that a variable \verb|satelliteFile| with \verb|data/swarm_orbit_{loopTime:%D}.dat|
can be declared in the global section without the variable \verb|loopTime| having to be known at this time.

One special variable is \verb|groopsDataDir|, which is used as a variable in most default
file paths throughout many GROOPS programs. Since this variable is going to be needed in
most config files, it is recommended to define it in a template file that is used when creating
new config files in the GUI. See section~\reference{Graphical User Interface (GUI)}{general.gui}
for details on how to set up a template file.

Global variables can be manipulated when running a config file by passing the argument \verb|--global <name>=<value>|.
For example, running the command
\begin{verbatim}
  groops --global timeStart=58849 --global satellite=swarm config.xml
\end{verbatim}
runs the config file \verb|config.xml| but replaces the values of the global variable \verb|timeStart|
and \verb|satellite| with \verb|58849| and \verb|swarm|, respectively. If a global variable passed as
an argument does not already exist in the config file, it will be added with the type \verb|string|.
Only the basic data types listed above are supported. This feature can be useful when running GROOPS
from the command line or from an external script file.

\clearpage
\section{Parsers and variables}\label{general.parser}
The XML configuration file is evaluated by two parsers. In a first step a text parser is applied.
In the second step mathematical expressions are resolved to a number.
Variables (see section~\reference{variables}{general.configFiles:variables}) can be referenced via their
name directly for the expression parser or in the form \verb|{name}| for the text parser.

% auto generated by GROOPS

\subsection{Mathematical expression parser}\label{general.parser:expression}
In all input fields that accept numbers (int, uint, double, angle, time) numerical
expressions are also allowed. Values from the global section can be used as variables. The following
operations are defined:
\begin{itemize}
\item Constants:    \verb|pi|, \verb|rho=180/pi|, \verb|nan|
\item Mathematical: \verb|+|, \verb|-|, \verb|*|, \verb|/|, \verb|^|
\item Comparison:   \verb|==|, \verb|!=|, \verb|<|, \verb|<=|, \verb|>|, \verb|>=|, result is 1 or 0
\item Logical:      not \verb|!|, and \verb|&&|, \verb'||', or \verb|isnan()|, result is 1 or 0
\item Functions:    \verb|sqrt()|, \verb|exp()|,
                    \verb|sin()|,  \verb|cos()|, \verb|tan()|,
                    \verb|asin()|,  \verb|acos()|,  \verb|atan()|,
                    \verb|abs()|,  \verb|round()|,  \verb|ceil()|,  \verb|floor()|,
\item Functions with 2 arguments: \verb|atan2(,)|, \verb|min(,)|, \verb|max(,)|, \verb|mod(,)|
\item Functions with 3 arguments: \verb|date2mjd(year, month, day)|
\item Condition: \verb|if(,,)|: If the first argument is true (not 0), the second argument is evalutated, otherwise the third.
\end{itemize}


\subsection{Text parser}\label{general.parser:text}
Before the mathematical expression parser evaluates the expression, a simple text parser is applied.
The text parser is used for all input fields (also file names). It scans the text for terms like
\verb|{variable}| and replaces it by the text content of the \verb|variable| defined in the global section.

The text parser also evaluates terms in the form \verb|{expression:format}| and replaces it by a formated
output. The \verb|format| contains the text to be written as output.
It can contain embedded format specifiers that are replaced by the value of the expression
and formated as requested (also multiple times). In the following, the resulting formated output is given in the
brackets for an expression with the example value of 57493.8:
\begin{itemize}
\item \verb|%i|: Integer [57494]
\item \verb|%f|: Decimal floating point [57493.800000]
\item \verb|%e|: Scientific notation [5.749380e+04]
\item \verb|%g|: Use the shortest representation: \verb|%e| or \verb|%f| [57493.8]
\item \verb|%c|: Interpret number as ASCII character
\item \verb|%%|: Write a single literal \verb|%| character
\end{itemize}
The following specifiers interpret the value of the expression as MJD (modified Julian date):
\begin{itemize}
\item \verb|%y|: Four digit year [2016]
\item \verb|%Y|: Two digit year [16]
\item \verb|%m|: Month [04]
\item \verb|%d|: Day of month [15]
\item \verb|%H|: Hour [19]
\item \verb|%M|: Minute [12]
\item \verb|%S|: Second [00]
\item \verb|%D|: Date (same as \verb|%y-%m-%d|) [2016-04-15]
\item \verb|%T|: Time (same as \verb|%H-%M-%S|) [19-12-00]
\item \verb|%W|: GPS week [1892]
\item \verb|%w|: Day of GPS week (0..6) [5]
\item \verb|%O|: Day of year (1..366)
\end{itemize}
The format can be specified further with \verb|%[width][.precision]specifier|,
where \verb|[width]| is the minimum number of characters to be printed.
If the value to be printed is shorter than this number, the result is padded with blank spaces
(or zeros if \verb|[width]| starts with a zero).
The \verb|[.precision]| defines the number of digits after the period (for \verb|%g| the number of
significant digits instead).

Example:
Two variables \config{time}=\verb|57493+19/24+12/1440| and \config{satellite}=\verb|swarm| are
set in the global section. The \config{inputfile}=\verb|data/{time:%y}/{satellite}_{time:%D}.dat|
is expanded to \verb|"data/2016/swarm_2016-04-15.dat"|.

Example:
The variable \config{x}=\verb|3+5| is set in the global section.
The expression \config{number}=\verb|2*x| is evaluated by the expression parser to \verb|=16|.
In contrast if we use brackets like in \config{number}=\verb|2*{x}| the expression is first evaluated
by the text parser to \verb|"2*3+5"| and the expression parser now gives the result \verb|=11|.


\subsection{Variables for data}\label{general.parser:dataVariables}
Some programs (e.g. \program{FunctionsCalculate}, \program{InstrumentArcCalculate},
\program{GriddedDataCalculate}, or the plot programs)
read data (\file{matrix}{matrix}) or \file{gridded data}{griddedData}
and evaluate input/output expressions for each data row.
For these kind of expressions additional variables are automatically defined for each data column
(\verb|X| stands for the data column number: $0\ldots n$):
\begin{itemize}
\item \verb|index|: the row number, starting with zero
\item \verb|dataX|: the value itself
\item \verb|dataXcount|: number of rows
\item \verb|dataXmin|
\item \verb|dataXmax|
\item \verb|dataXrms|: root mean square
\item \verb|dataXstd|: standard deviation
\item \verb|dataXmedian|
\item \verb|dataXmad|: median absolute deviation
\item \verb|dataXstep|: the minimal difference between two neighboring data points in the column
\end{itemize}
For \file{gridded data}{griddedData} input the following variables are additionally defined for each data point:
\begin{itemize}
\item \verb|longitude| in degrees
\item \verb|latitude| in degrees
\item \verb|height| in meters
\item \verb|cartesianX| coordinate in meters
\item \verb|cartesianY| coordinate in meters
\item \verb|cartesianZ| coordinate in meters
\item \verb|area| of the unit sphere
\item \verb|dataXwmean|: area-weighted mean
\item \verb|dataXwrms|: area-weighted root mean square
\item \verb|dataXwstd|: area-weighted standard deviation
\end{itemize}



\clearpage
\section{Loops and conditions}\label{general.loopsAndConditions}
The program flow within a config file can be controlled by the classes \configClass{loop}{loopType}
and \configClass{condition}{conditionType}. The easiest way to access these classes is with the programs
\program{LoopPrograms} and \program{IfPrograms}.

The programs defined in \program{IfPrograms} are only executed if the defined
\configClass{condition}{conditionType} is met. A typical example is to check whether a file that
should have been created in previous programs actually exists. Further options are string comparisons and
checking the result of a numerical expression or the return value of an external command.

With \program{LoopPrograms} it is possible to repeat the programs defined inside within a loop.
The class \configClass{loop}{loopType} creates a sequence to loop over and defines \reference{variables}{general.parser}
that contain the index and element for the current iteration.

If \configClass{loop}{loopType} and \configClass{condition}{conditionType} are defined in the global
section (see section~\reference{Config files}{general.configFiles}) they become powerful tools as the
they can be attributed to single config elements (including programs).
Config elements with an assigned loop are repeated, with the loop variables being evaluated for each element.
If a \configClass{condition}{conditionType} is attributed to a config element in addition to a loop, each element within the loop
is only created if the condition is met. Conditions can also be attributed to optional elements without an associated loop.
If the condition is not met, the optional element will be treated as if it was not provided.

Example: A program needs all files in a download directory as input.
All the \config{inputfile}s can be selected manually of course, but it is much easier to assign
a loop variable with \config{inputfile}=\verb|{loopFile}| and attribute a
\configClass{loop:directoryListing}{loopType:directoryListing} from the global section.
The loop lists the content of the download directory and assigns each file name to the
\config{variableLoopFile}=\verb|loopFile|.

\clearpage
\section{Constants and the settings file}\label{general.constants}
GROOPS uses some built-in constants like \verb|DEFAULT_GM| or the definition
of leap seconds, which are defined in  \verb|source/base/constants.cpp|.

A complete list of the constants can be written to an XML file with:
\begin{verbatim}
  groops --write-settings <groopsDefaults.xml>
\end{verbatim}

The built-in constants can be overwritten by a \verb|groopsDefaults.xml| file
in the working directory or by explicitly passing the file as an argument at execution:
\begin{verbatim}
  groops --settings <groopsDefaults.xml> <config.xml>
\end{verbatim}

It might also be useful to adjust the default values in the schema file used by the \reference{GUI}{general.gui}:
\begin{verbatim}
  groops --settings <groopsDefaults.xml> --xsd <groops.xsd>
\end{verbatim}

Example file:
\begin{verbatim}
<?xml version="1.0" encoding="UTF-8"?>
<groops>
  <LIGHT_VELOCITY>299792458</LIGHT_VELOCITY>
  <DEFAULT_GM>3.986004415e+14</DEFAULT_GM>
  <DEFAULT_R>6378136.3</DEFAULT_R>
  <GRS80_a>6378137.0</GRS80_a>
  <GRS80_f>298.2572221010</GRS80_f>
  <GRAVITATIONALCONSTANT>6.6730e-11</GRAVITATIONALCONSTANT>
  <R_Earth>6.37813630000000e+06</R_Earth>
  <R_Moon>1.73800000000000e+06</R_Moon>
  <GM_Earth>3.98600441500000e+14</GM_Earth>
  <GM_Sun>1.32712442076000e+20</GM_Sun>
  <GM_Moon>4.90280105600000e+12</GM_Moon>
  <GM_MERCURY>2.20320808280762e+13</GM_MERCURY>
  <GM_VENUS>3.24858603864143e+14</GM_VENUS>
  <GM_MARS>4.28283149222192e+13</GM_MARS>
  <GM_JUPITER>1.26712769822770e+17</GM_JUPITER>
  <GM_SATURN>3.79406266494906e+16</GM_SATURN>
  <TIME_EPSILON>1.0e-05</TIME_EPSILON>
  <DELTA_TAI_GPS>19</DELTA_TAI_GPS>
  <DELTA_TT_GPS>51.184</DELTA_TT_GPS>
  <J2000>51544.5</J2000>
  <leapSecond>
    <MJD>57754</MJD>
    <DELTA_UTC_GPS>-18</DELTA_UTC_GPS>
  </leapSecond>
  <leapSecond>
    <MJD>57204</MJD>
    <DELTA_UTC_GPS>-17</DELTA_UTC_GPS>
  </leapSecond>
  <leapSecond>
    <MJD>56109</MJD>
    <DELTA_UTC_GPS>-16</DELTA_UTC_GPS>
  </leapSecond>

  ...

  <leapSecond>
    <MJD>41317</MJD>
    <DELTA_UTC_GPS>9</DELTA_UTC_GPS>
  </leapSecond>
  <leapSecond>
    <MJD>0</MJD>
    <DELTA_UTC_GPS>10</DELTA_UTC_GPS>
  </leapSecond>
</groops>
\end{verbatim}

\clearpage
\section{Parallelization}\label{general.parallelization}
If GROOPS is compiled with the
\href{https://en.wikipedia.org/wiki/Message_Passing_Interface}{Message Passing Interface} (MPI),
most GROOPS \reference{programs}{programType} can be run in parallel on multiple processor cores.
Processing on computer clusters with distributed memory is also supported.

Many loops are parallelized by computing each loop step at a different core.
Usually the first node distributes the work load, assigns loop steps to different cores,
and is not participating on the actual loop computation. This means running
GROOPS with only two nodes has no advantages in almost all cases.
Non-parallel parts and \reference{programs}{programType} without parallel support
are executed at the first node only.

Large systems of \file{normal equations}{normalEquation}, which are divided into blocks,
are distributed over the nodes to reduce the memory consumption on each single node.

As all nodes may read and write files (at least reading the \reference{config files}{general.configFiles})
the required part of the file system must be available on all participating computers.

\clearpage
\section{Graphical User Interface (GUI)}\label{general.gui}

The graphical user interface program \verb|groopsGui| enables the convenient creation of GROOPS config files.
It uses the \href{https://qt.io}{Qt5 framework} for cross-platform support.

\fig{!ht}{1.0}{gui_overview}{fig:gui_overview}{Overview of the GUI (components mentioned below marked in red)}

\subsection{Settings and first-time setup}

The GUI depends on an XSD schema file containing the complete formal (computer readable) description of a GROOPS config file.
This schema file can be created with the command:
\begin{verbatim}
  groops --xsd <groopsDir>/groops.xsd
\end{verbatim}
At least one schema file has to be set via
the menu \verb|Settings - Default Paths and Files|. Setting more than one schema files enables the \emph{schema selector} in the toolbar.
The selected schema will be used when (re-)opening or creating a config file.
This feature is useful when working with different versions of GROOPS at the same time.

It is possible to set a \emph{template file} via the menu \verb|Settings - Default Paths and Files|. This can be any GROOPS config file.
Whenever a new config file is created via the GUI, all global elements and programs defined in the template file are automatically created in the new config file.
It is highly recommended to create a template file containing at least the global element \verb|groopsDataDir| of type \verb|filename|.
This element is used as a \reference{variable}{general.parser} in most default file paths throughout many GROOPS programs.
Thus, setting the path to the base directory containing all GROOPS data once in the template file, for example as
\config{groopsDataDir}=\verb|/home/<user>/groops/data|, is the most convenient way to handle default paths in GROOPS.
The template file can also contain other often-used global elements, for example \verb|tmpDir| or  \verb|timeStart| and  \verb|timeEnd|.

A \emph{working directory} can be set via \verb|Settings - Default Paths and Files|.
This directory is used as the default directory in the save dialog of new config files.

The GUI offers the option to open the GROOPS documentation for a selected program. To use this feature,
the GROOPS documentation must be generated (if not already present) with the command:
\begin{verbatim}
  groops --doc <groopsDir>/docs/
\end{verbatim}
In the menu \verb|Settings - Default Paths and Files| the path to the HTML version of the documentation must be set (i.e. \verb|<groopsDir>/docs/html|).
Selecting any program and pressing \verb|F1| opens the documentation for this program in an external browser.
Pressing \verb|F1| without having any program selected opens the main page of the GROOPS documentation.

Executing a config file from the GUI requires the setup of a run command in the menu \verb|Settings - Commands|.
It is recommended for this command to open a new terminal in which GROOPS is executed with the config file given as an argument.
The placeholders \verb|%w| and \verb|%f| are replaced by the directory and file name of the selected config file, respectively.
Multiple commands can be set up, with the option to choose one of them in the run dialog.

Example commands:
\begin{itemize}
	\item Windows: \verb|cd /d %w && groops.exe %f|
	\item Linux (KDE): \verb|konsole --workdir %w -e bash -ic "groops %f; bash"|
	\item Linux (GNOME): \verb|gnome-terminal --working-directory=%w -x bash -ic "groops %f; bash"|
	\item Windows, MPI with 4 processes: \verb|cd /d %w && mpiexec -n 4 groopsMPI.exe %f|
	\item Linux (KDE), MPI with 4 processes: \verb|konsole --workdir %w -e bash -ic "mpiexec -n 4 groopsMPI %f; bash"|
	\item Linux (GNOME), MPI with 4 processes: \verb|gnome-terminal --working-directory=%w -x bash -ic "mpiexec -n 4 groopsMPI %f; bash"|
\end{itemize}

\subsection{Basic features}

Most basic features used to manipulate a config element are accessible via the context menu,
for example attributing \reference{loops and conditions}{general.loopsAndConditions} or setting an element global.
Global elements automatically appear in the dropdown value list of config elements of the same type.
Selecting a global element from the dropdown list as a value links this config element to the global element.
In case the global element is removed, all linked elements' values are replaced by the value of the deleted global element.

The sidebar features three widgets:
\begin{itemize}
	\item \verb|Open Files|: An overview of all open config files (select to change current tree)
	\item \verb|Program List|: A list of all programs defined in the schema of the active tree (filterable, supports drag and drop to tree, double click appends program)
	\item \verb|Undo Stack|: Tracks all changes in a config file (select to change state of tree)
\end{itemize}

In case the names of programs or config elements change over time, the GUI offers a rename feature to update outdated config files.
The changes must be documented in the schema using GROOPS' rename feature. Affected elements will be marked with an icon and
the context menu item \verb|Update name| will be available to change the element to the new name defined in the schema.

\subsection{Additional keyboard shortcuts}

\begin{itemize}
	\item Tree navigation:
	\begin{itemize}
		\item \verb|Enter|: Switch focus from tree to input field of selected row
		\item \verb|Escape|: Switch focus from input field back to tree
		\item \verb|Tab|: Next sibling element (or next sibling of parent if there is no next sibling, or next child otherwise)
		\item \verb|Shift+Tab|: Previous sibling element (or parent if there is no previous sibling)
		\item \verb|Ctrl+Tab|: Next tab/tree
		\item \verb|Ctrl+Shift+Tab|: Previous tab/tree
		\item \verb|Ctrl+Space|: Interact with the element (e.g. filename/program: open dialog; time: switch focus between input fields)
		\item \verb|Ctrl+Up/Down|: Next/previous sibling element
		\item \verb|Ctrl+Left/Right|: Fold/expand (complex) element
	\end{itemize}
	\item Tree manipulation:
	\begin{itemize}
		\item \verb|Ctrl+Shift+Up/Down|: Move unbounded list element (e.g. program, layer) up/down
	\end{itemize}
	\item Drag and Drop of tabs to other programs (i.e. text editors) or other GUI windows:
	\begin{itemize}
		\item \verb|Drag|: Copy tab (= keep in source window)
		\item \verb|Shift+Drag|: Move tab (= remove from source window)
	\end{itemize}
	\item Drag and Drop GROOPS config file(s) into GUI:
	\begin{itemize}
		\item \verb|Drag|: Open file(s) in new tab(s)
		\item \verb|Shift+Drag|: Open file in current tab (replaces current tab, only works with a single file)
	\end{itemize}
\end{itemize}

\clearpage
\section{File formats}\label{general.fileFormat}
All GROOPS files are written either in XML, JSON, binary, or ASCII format depending on the filename extension.
\begin{itemize}
\item \verb|.xml|: XML format
\item \verb|.json|: JSON format
\item \verb|.dat|: binary format
\item \verb|.txt| and all other extensions: ASCII format
\end{itemize}

With an additional extension of '\verb|.gz|' files are directly compressed and uncompressed. It is also possible to directly uncompress and read (but not write) \href{https://en.wikipedia.org/wiki/Compress}{Unix compress}'d files ('\verb|.Z|').

Comments are allowed in ASCII files and all the text starting from the character '\verb|#|' to the end of the line is ignored.

The program \program{FileConvert} can be used to convert between the different formats. This program is also useful to get
some general information of files in binary format.

The following special file types are used in GROOPS:

% auto generated by GROOPS
\subsection{Admittance}\label{general.fileFormat:admittance}
Interpolation matrix to create ocean minor tides from modeled major tides.
The file can be created with \program{DoodsonHarmonicsCalculateAdmittance} and used e.g. in
\configClass{doodsonHarmonicTide}{tidesType:doodsonHarmonicTide}.


%==================================
\subsection{ArcList}\label{general.fileFormat:arcList}
With the \program{InstrumentSynchronize} an \file{instrument file}{instrument} can
be divided into time intervals and within the intervals into arcs.
This file provides the information about the mapping of arcs to time intervals.

This file can be used for the variational equation approach or \program{KalmanBuildNormals}.

\begin{verbatim}
groops arclist version=20200123
         32  # number of times
# time [MJD]               first arc
# ==================================
 58909.000000000000000000          0
 58910.000000000000000000          8
 58911.000000000000000000         17
 58912.000000000000000000         25
 58913.000000000000000000         29
 58914.000000000000000000         37
 58915.000000000000000000         45
 58916.000000000000000000         53
 58917.000000000000000000         61
 58918.000000000000000000         69
 58919.000000000000000000         78
 58920.000000000000000000         86
 58921.000000000000000000         95
 58922.000000000000000000        103
 58923.000000000000000000        112
 58924.000000000000000000        120
 58925.000000000000000000        128
 58926.000000000000000000        136
 58927.000000000000000000        144
 58928.000000000000000000        153
 58929.000000000000000000        161
 58930.000000000000000000        169
 58931.000000000000000000        177
 58932.000000000000000000        185
 58933.000000000000000000        193
 58934.000000000000000000        201
 58935.000000000000000000        210
 58936.000000000000000000        218
 58937.000000000000000000        226
 58938.000000000000000000        234
 58939.000000000000000000        242
 58940.000000000000000000        250
\end{verbatim}


%==================================
\subsection{DoodsonEarthOrientationParameter}\label{general.fileFormat:doodsonEarthOrientationParameter}
Corrections for Earth orientation parameters (EOP) ($x_p, y_p, UT1, LOD$)
as cos/sin oscillations for a list of doodson tidal frequencies.

\begin{verbatim}
groops doodsonEarthOrientationParameter version=20200123
        11 # number of constituents
# dood.   xpCos [arcsec]            xpSin [arcsec]            ypCos [arcsec]            ypSin [arcsec]            ut1Cos [sec]              ut1Sin [sec]              lodCos [sec]              lodSin [sec]             name
# ===========================================================================================================================================================================================================================
 155.645  2.399999999999999786e-07  1.799999999999999971e-07  1.799999999999999971e-07 -2.399999999999999786e-07 -2.000000000000000042e-08  2.000000000000000042e-08 -1.499999999999999932e-07 -1.400000000000000095e-07 ""
 155.655 -8.220000000000000920e-06 -6.280000000000000012e-06 -6.280000000000000012e-06  8.220000000000000920e-06  7.899999999999999537e-07 -8.599999999999999187e-07  5.219999999999999150e-06  4.829999999999999496e-06 m1
 155.665 -1.649999999999999872e-06 -1.260000000000000007e-06 -1.260000000000000007e-06  1.649999999999999872e-06  1.600000000000000033e-07 -1.700000000000000135e-07  1.049999999999999900e-06  9.700000000000000302e-07 ""
 157.455 -1.539999999999999867e-06 -1.199999999999999946e-06 -1.199999999999999946e-06  1.539999999999999867e-06  1.499999999999999932e-07 -1.600000000000000033e-07  9.799999999999999345e-07  8.899999999999999491e-07 chi1
 157.465 -3.400000000000000270e-07 -2.599999999999999988e-07 -2.599999999999999988e-07  3.400000000000000270e-07  2.999999999999999732e-08 -4.000000000000000084e-08  2.099999999999999746e-07  1.999999999999999909e-07 ""
 161.557  9.999999999999999547e-08  8.000000000000000167e-08  8.000000000000000167e-08 -9.999999999999999547e-08 -1.000000000000000021e-08  1.000000000000000021e-08 -7.000000000000000477e-08 -4.999999999999999774e-08 ""
 162.556  2.549999999999999726e-06  2.019999999999999718e-06  2.019999999999999718e-06 -2.549999999999999726e-06 -2.099999999999999746e-07  2.899999999999999763e-07 -1.799999999999999919e-06 -1.289999999999999931e-06 pi1
 163.545 -4.899999999999999672e-07 -3.799999999999999616e-07 -3.799999999999999616e-07  4.899999999999999672e-07  4.000000000000000084e-08 -5.999999999999999464e-08  3.499999999999999842e-07  2.399999999999999786e-07 ""
 163.555  4.272999999999999238e-05  3.010999999999999801e-05  3.010999999999999801e-05 -4.272999999999999238e-05 -3.079999999999999734e-06  5.219999999999999150e-06 -3.270999999999999684e-05 -1.929999999999999817e-05 p1
 164.554 -3.599999999999999943e-07 -2.800000000000000191e-07 -2.800000000000000191e-07  3.599999999999999943e-07  2.999999999999999732e-08 -4.000000000000000084e-08  2.700000000000000090e-07  1.700000000000000135e-07 ""
 164.556 -1.029999999999999879e-06 -7.999999999999999638e-07 -7.999999999999999638e-07  1.029999999999999879e-06  8.000000000000000167e-08 -1.199999999999999893e-07  7.599999999999999233e-07  4.899999999999999672e-07 ""
\end{verbatim}


%==================================
\subsection{DoodsonHarmonic}\label{general.fileFormat:doodsonHarmonic}
Ocean tides are represented as time variable gravitational potential
and is given by a fourier expansion
\begin{equation}
V(\M x,t) = \sum_{f} V_f^c(\M x)\cos(\Theta_f(t)) + V_f^s(\M x)\sin(\Theta_f(t)),
\end{equation}
where $V_f^c(\M x)$ and $V_f^s(\M x)$ are spherical harmonics.
The $\Theta_f(t)$ are the arguments of the tide constituents $f$:
\begin{equation}
\Theta_f(t) = \sum_{i=1}^6 n_f^i\beta_i(t),
\end{equation}
where $\beta_i(t)$ are the Doodson's fundamental arguments ($\tau,s,h,p,N',p_s$) and $n_f^i$
are the Doodson multipliers for the term at frequency~$f$.

To extract the potential coefficients of $V_f^c$ and $V_f^s$
for each frequency $f$ use \program{DoodsonHarmonics2PotentialCoefficients}.

See also \program{PotentialCoefficients2DoodsonHarmonics}.


%==================================
\subsection{EarthOrientationParameter}\label{general.fileFormat:earthOrientationParameter}
Earth Orientation Parameter (EOP) as provided by the International Earth Rotation and Reference Systems Service (IERS) (e.g \verb|EOP 14 C04 (IAU2000A)|).

See \program{IersC04IAU2000EarthOrientationParameter}, \program{IersRapidIAU2000EarthOrientationParameter}.

\begin{verbatim}
groops earthOrientationParameter version=20200123
       9641 # number of epochs
# UTC [MJD]                 xp [arcsec]               yp [arcsec]               deltUT [sec]              LOD [sec]                 dX [arcsec]               dY [arcsec]
# ====================================================================================================================================================================================
  5.894700000000000000e+04  5.690599999999999825e-02  4.099130000000000273e-01 -2.316246000000000138e-01  1.636400000000000094e-03 -2.900000000000000017e-05  5.800000000000000034e-05
  5.894800000000000000e+04  5.771400000000000141e-02  4.110159999999999925e-01 -2.332083000000000073e-01  1.520099999999999923e-03 -6.000000000000000152e-05  2.199999999999999943e-05
  5.894900000000000000e+04  5.813000000000000111e-02  4.120099999999999874e-01 -2.346157000000000104e-01  1.293099999999999935e-03 -7.200000000000000182e-05  3.199999999999999855e-05
  5.895000000000000000e+04  5.854100000000000276e-02  4.129849999999999910e-01 -2.357567999999999886e-01  9.872999999999999832e-04 -7.600000000000000418e-05  5.899999999999999754e-05
  5.895100000000000000e+04  5.908599999999999963e-02  4.139869999999999939e-01 -2.366149999999999920e-01  7.075000000000000126e-04 -8.000000000000000654e-05  8.600000000000000331e-05
  5.895200000000000000e+04  5.976900000000000268e-02  4.154180000000000095e-01 -2.372105999999999937e-01  4.798000000000000073e-04 -8.399999999999999535e-05  1.129999999999999955e-04
  5.895300000000000000e+04  6.095400000000000124e-02  4.167310000000000181e-01 -2.375994999999999913e-01  3.118999999999999919e-04 -8.700000000000000051e-05  1.399999999999999877e-04
  5.895400000000000000e+04  6.210199999999999748e-02  4.180929999999999924e-01 -2.378588000000000091e-01  1.710000000000000094e-04 -9.100000000000000287e-05  1.669999999999999935e-04
  5.895500000000000000e+04  6.290999999999999370e-02  4.196619999999999795e-01 -2.380454999999999932e-01  1.719000000000000042e-04 -6.000000000000000152e-06  8.100000000000000375e-05
  5.895600000000000000e+04  6.385599999999999610e-02  4.214060000000000028e-01 -2.382557999999999898e-01  2.683000000000000163e-04  1.029999999999999964e-04 -3.600000000000000091e-05
  5.895700000000000000e+04  6.455500000000000127e-02  4.229890000000000039e-01 -2.385857000000000117e-01  4.040000000000000080e-04  1.019999999999999992e-04 -2.000000000000000164e-05
  5.895800000000000000e+04  6.440300000000000191e-02  4.242549999999999932e-01 -2.390210000000000112e-01  4.910999999999999567e-04  5.899999999999999754e-05  4.399999999999999886e-05
\end{verbatim}


%==================================
\subsection{EarthTide}\label{general.fileFormat:earthTide}
Containing the Love numbers together with frequency corrections to compute
the gravitational potential and the geometric displacements due to solid Earth tides.
It is used by \configClass{tides}{tidesType}.


%==================================
\subsection{Ephemerides}\label{general.fileFormat:ephemerides}
Ephemerides of sun, moon, and planets stored as coefficients of Chebyshev polynomials.
Used in \configClass{Ephemerides:jpl}{ephemeridesType:jpl}.

See also \program{JplAscii2Ephemerides}.


%==================================
\subsection{GnssAntennaDefinition}\label{general.fileFormat:gnssAntennaDefinition}
Contains a list of GNSS antennas which are identified by its
name (type), serial, and radome. Each antenna consists of
antenna center offsets (ACO) and antenna center variations (ACV)
for different signal \configClass{types}{gnssType} (code and phase).
The ACV values for each type are stored in an elevation and azimuth dependent grid.

%New \config{antenna} with \config{pattern}s for code (\config{type}=\verb|C**|) and phase
%(\config{type}=\verb|L**|).
%The standard deviation is expressed e.g. with \config{values}=\verb|0.001/cos(2*PI/180*zenith)|.

See also \program{GnssAntennaDefinitionCreate}, \program{GnssAntex2AntennaDefinition}.

\fig{!hb}{1.0}{fileFormatGnssAntennaDefinition}{fig:fileFormatGnssAntennaDefinition}{Antenna center variations of ASH701945D\_M for two frequencies of GPS and GLONASS}

\begin{verbatim}
<?xml version="1.0" encoding="UTF-8"?>
<groops type="antennaDefinition" version="20190429">
    <antennaCount>65</antennaCount>
    ...
    <antenna>
        <name>BLOCK IIIA</name>
        <serial>G074</serial>
        <radome>2018-109A</radome>
        <comment>PCO provided by the Aerospace Corporation, PV from estimations by ESA/CODE</comment>
        <pattern>
            <count>3</count>
            <cell>
                <type>*1*G**</type>
                <offset>
                    <x>-1.23333333333333e-03</x>
                    <y>4.33333333333333e-04</y>
                    <z>3.15200000000000e-01</z>
                </offset>
                <dZenit>1.00000000000000e+00</dZenit>
                <pattern>
                    <type>0</type>
                    <rows>1</rows>
                    <columns>18</columns>
                    <cell row="0" col="0">1.39000000000000e-02</cell>
                    <cell row="0" col="1">1.28000000000000e-02</cell>
                    <cell row="0" col="2">1.02000000000000e-02</cell>
                    <cell row="0" col="3">5.80000000000000e-03</cell>
                    <cell row="0" col="4">1.10000000000000e-03</cell>
                    <cell row="0" col="5">-4.50000000000000e-03</cell>
                    <cell row="0" col="6">-9.70000000000000e-03</cell>
                    <cell row="0" col="7">-1.28000000000000e-02</cell>
                    <cell row="0" col="8">-1.34000000000000e-02</cell>
                    <cell row="0" col="9">-1.18000000000000e-02</cell>
                    <cell row="0" col="10">-8.90000000000000e-03</cell>
                    <cell row="0" col="11">-4.50000000000000e-03</cell>
                    <cell row="0" col="12">1.20000000000000e-03</cell>
                    <cell row="0" col="13">7.20000000000000e-03</cell>
                    <cell row="0" col="14">1.33000000000000e-02</cell>
                    <cell row="0" col="15">1.33000000000000e-02</cell>
                    <cell row="0" col="16">1.33000000000000e-02</cell>
                    <cell row="0" col="17">1.33000000000000e-02</cell>
                </pattern>
            </cell>
            ...
        </pattern>
    </antenna>
</groops>
\end{verbatim}


%==================================
\subsection{GnssReceiverDefinition}\label{general.fileFormat:gnssReceiverDefinition}
Contains a list of GNSS receivers which are identified by its
name, serial, and version. Defines for each receiver a list of
signal \configClass{types}{gnssType} which can be observed.
Can also be used for GNSS transmitters to define a list of
transmitted signal types. For GLONASS satellites the frequency
number can be stored in the \emph{version} field.

See \program{GnssReceiverDefinitionCreate} and \program{GnssGlonassFrequencyNumberUpdate}.

\begin{verbatim}
<?xml version="1.0" encoding="UTF-8"?>
<groops type="receiverDefinition" version="20190429">
    <receiverCount>112</receiverCount>
    <receiver>
        <name>GLONASS</name>
        <serial>R779</serial>
        <version>2</version>
        <comment/>
        <types>
            <count>4</count>
            <cell>*1CR**J</cell>
            <cell>*1PR**J</cell>
            <cell>*2CR**J</cell>
            <cell>*2PR**J</cell>
        </types>
    </receiver>
    ...
    <receiver>
        <name>GLONASS-K1</name>
        <serial>R802</serial>
        <version>7</version>
        <comment/>
        <types>
            <count>10</count>
            <cell>*1CR**O</cell>
            <cell>*1PR**O</cell>
            <cell>*2CR**O</cell>
            <cell>*2PR**O</cell>
            <cell>*3IR**</cell>
            <cell>*3QR**</cell>
            <cell>*4AR**</cell>
            <cell>*4BR**</cell>
            <cell>*6AR**</cell>
            <cell>*6BR**</cell>
        </types>
    </receiver>
</groops>
\end{verbatim}


%==================================
\subsection{GnssSignalBias}\label{general.fileFormat:gnssSignalBias}
Signal biases of GNSS transmitters or receivers for different \configClass{gnssType}{gnssType}.

\begin{verbatim}
groops gnssSignalBias version=20200123
          5 # number of signals
# type   bias [m]
# ===============================
 C1CG06 -1.752461109688110974e-01
 C1WG06  4.005884595055994590e-02
 C2WG06  6.597469378913034532e-02
 L1*G06 -2.736169875580296909e-02
 L2*G06  3.422596762686257871e-02
 \end{verbatim}

See also \program{GnssProcessing}, \program{GnssSimulateReceiver}, \program{GnssSignalBias2Matrix}, \program{GnssSignalBias2SinexBias}.


%==================================
\subsection{GnssStationInfo}\label{general.fileFormat:gnssStationInfo}
Information file for a station or satellite equipped with a GNSS receiver
or for a GNSS satellite.

For stations and receiving satellites it defines the time intervals and
local coordinates of installed antennas and receivers.
For GNSS satellites it defines the PRN-SVN mapping and antenna offsets
and orientations.

GNSS station infos can be created from station log files with
\program{GnssStationLog2StationInfo}. GNSS satellite transmitter infos
can be created from an ANTEX file with \program{GnssAntex2AntennaDefinition}.

See \program{GnssStationInfoCreate}.

\begin{verbatim}
<?xml version="1.0" encoding="UTF-8"?>
<groops type="stationInfo" version="20190429">
    <stationCount>1</stationCount>
    <station>
        <markerName>GRAZ</markerName>
        <markerNumber>11001M002</markerNumber>
        <comment>Graz, AUSTRIA</comment>
        <approxPosition>
            <x>4.19442412400000e+06</x>
            <y>1.16270246000000e+06</y>
            <z>4.64724519300000e+06</z>
        </approxPosition>
        <antenna>
            <count>8</count>
            ...
            <cell>
                <name>LEIAR25.R4</name>
                <serial>726444</serial>
                <radome>LEIT</radome>
                <comment/>
                <timeStart>58080.4041666666666650</timeStart>
                <timeEnd>234166.0000000000000000</timeEnd>
                <position>
                    <x>0.00000000000000e+00</x>
                    <y>0.00000000000000e+00</y>
                    <z>1.96400000000000e+00</z>
                </position>
                <local2antennaFrame>
                    <xx>1.00000000000000e+00</xx>
                    <xy>0.00000000000000e+00</xy>
                    <xz>0.00000000000000e+00</xz>
                    <yx>0.00000000000000e+00</yx>
                    <yy>1.00000000000000e+00</yy>
                    <yz>0.00000000000000e+00</yz>
                    <zx>0.00000000000000e+00</zx>
                    <zy>0.00000000000000e+00</zy>
                    <zz>1.00000000000000e+00</zz>
                </local2antennaFrame>
            </cell>
        </antenna>
        <receiver>
            <count>25</count>
            ...
            <cell>
                <name>SEPT POLARX5</name>
                <serial>4501501</serial>
                <version>5.2.0</version>
                <comment/>
                <timeStart>58903.2958333333333343</timeStart>
                <timeEnd>234166.0000000000000000</timeEnd>
            </cell>
        </receiver>
        <referencePoint>
            <count>0</count>
        </referencePoint>
    </station>
</groops>
\end{verbatim}


%==================================
\subsection{GriddedData}\label{general.fileFormat:griddedData}
List of arbitrarily distributed points defined by geographic coordinates and ellipsoidal
height. Each point can also have an associated area
(projected on the unit sphere with a total area of $4\pi$).
This file format supports multiple values per point (called \verb|data0|, \verb|data1| and so on).

For regular gridded data and binary format (\verb|*.dat|) a more efficent storage scheme is used.

See also: \program{GriddedDataCreate}.

\begin{verbatim}
groops griddedData version=20200123
 1  2  6.378137000000000000e+06  6.356752314140356146e+06 72 # hasArea, data columns, ellipoid a, ellipoid b, data rows
# longitude [deg]           latitude [deg]            height [m]                unit areas [-]             data0                     data1
# ===========================================================================================================================================================
 -1.650000000000000000e+02  7.500000000000000000e+01  0.000000000000000000e+00  7.014893453974438420e-02  1.000000000000000000e+00  2.000000000000000000e+00
 -1.350000000000000000e+02  7.500000000000000000e+01  0.000000000000000000e+00  7.014893453974438420e-02  1.000000000000000000e+00  2.000000000000000000e+00
 -1.050000000000000142e+02  7.500000000000000000e+01  0.000000000000000000e+00  7.014893453974438420e-02  1.000000000000000000e+00  2.000000000000000000e+00
 -7.500000000000001421e+01  7.500000000000000000e+01  0.000000000000000000e+00  7.014893453974438420e-02  1.000000000000000000e+00  2.000000000000000000e+00
 -4.500000000000002132e+01  7.500000000000000000e+01  0.000000000000000000e+00  7.014893453974438420e-02  1.000000000000000000e+00  2.000000000000000000e+00
 -1.500000000000002132e+01  7.500000000000000000e+01  0.000000000000000000e+00  7.014893453974438420e-02  1.000000000000000000e+00  2.000000000000000000e+00
  1.499999999999997691e+01  7.500000000000000000e+01  0.000000000000000000e+00  7.014893453974438420e-02  1.000000000000000000e+00  2.000000000000000000e+00
  4.499999999999997868e+01  7.500000000000000000e+01  0.000000000000000000e+00  7.014893453974438420e-02  1.000000000000000000e+00  2.000000000000000000e+00
  7.499999999999997158e+01  7.500000000000000000e+01  0.000000000000000000e+00  7.014893453974438420e-02  1.000000000000000000e+00  2.000000000000000000e+00
  1.049999999999999574e+02  7.500000000000000000e+01  0.000000000000000000e+00  7.014893453974438420e-02  1.000000000000000000e+00  2.000000000000000000e+00
  1.349999999999999432e+02  7.500000000000000000e+01  0.000000000000000000e+00  7.014893453974438420e-02  1.000000000000000000e+00  2.000000000000000000e+00
  1.649999999999999432e+02  7.500000000000000000e+01  0.000000000000000000e+00  7.014893453974438420e-02  1.000000000000000000e+00  2.000000000000000000e+00
 -1.650000000000000000e+02  4.500000000000000711e+01  0.000000000000000000e+00  1.916504532594049681e-01  1.000000000000000000e+00  2.000000000000000000e+00
 -1.350000000000000000e+02  4.500000000000000711e+01  0.000000000000000000e+00  1.916504532594049681e-01  1.000000000000000000e+00  2.000000000000000000e+00
\end{verbatim}



%==================================
\subsection{GriddedDataTimeSeries}\label{general.fileFormat:griddedDataTimeSeries}
Time series of data for arbitrarily distributed points defined by geographic coordinates and ellipsoidal
height. The data can be temporal interpolated by \reference{basis splines}{fundamentals.basisSplines}.
The file format consists of a \file{griddedData}{griddedData}, a time series, and
for each spatial point and spline node pair multiple values called \verb|data0|, \verb|data1|, \ldots.


%==================================
\subsection{Instrument}\label{general.fileFormat:instrument}
This template file format can store different observations in a epoch wise manner. Each epoch consists of a time and
additional data, e.g orbits, accelerometer data, star camera quaternions (see \configClass{InstrumentType}{instrumentTypeType}).
The time series can be divided in several arcs (see \program{InstrumentSynchronize}).

Also a \file{matrix}{matrix} file is allowed as one single arc. The first column must contain times [MJD]. Without any extra column
the instrument type is INSTRUMENTTIME, with one additional column the type is MISCVALUE, and for more columns the type
MISCVALUES is used.

\begin{verbatim}
groops instrument version=20200123
# SATELLITETRACKING
         -9         60  # instrument type, number of arcs
# Time [MJD]               data0: range [m]          data1: range-rate [m/s]   data2: range-acc [m/s^2]
# =====================================================================================================
         12 # number of epochs of 1. arc
 54588.000000000000000000 -5.074649470097549492e+05  5.755440207134928654e-01  1.877605261528093308e-03
 54588.000057870370255841 -5.074620458130163024e+05  5.849357691551860805e-01  1.878948916234051596e-03
 54588.000115740740966430 -5.074590976427756250e+05  5.943331739937073310e-01  1.879937220634776869e-03
 54588.000173611111222272 -5.074561024756557308e+05  6.037340169611068452e-01  1.880370529387525701e-03
 54588.000231481481478113 -5.074530602992626373e+05  6.131368121270999172e-01  1.880680632122925426e-03
 54588.000289351851733954 -5.074499711071007187e+05  6.225398878861636565e-01  1.880495369480403561e-03
 54588.000347222222444543 -5.074468349029610981e+05  6.319414138081351773e-01  1.880073731783055927e-03
 54588.000405092592700385 -5.074436516971451929e+05  6.413404243585696385e-01  1.879464843086203459e-03
 54588.000462962962956226 -5.074404215058300761e+05  6.507353310092597320e-01  1.878578987216372124e-03
 54588.000520833333212067 -5.074371443491023383e+05  6.601267978060636477e-01  1.877878184949659246e-03
 54588.000578703703922656 -5.074338202460713219e+05  6.695136489207137442e-01  1.876962042758626532e-03
 54588.000636574074178498 -5.074304492190054734e+05  6.788964444122400632e-01  1.876091925462087043e-03
         12 # number of epochs of 2. arc
 54588.000694444444434339 -5.074270312892858055e+05  6.882748400534359767e-01  1.875376456928801432e-03
 54588.000752314814690180 -5.074235664742725785e+05  6.976508178537534910e-01  1.874929898412159559e-03
 54588.000810185185400769 -5.074200547868391732e+05  7.070236200716006891e-01  1.874312324351668077e-03
 54588.000868055555656611 -5.074164962409950094e+05  7.163943828291452487e-01  1.873924188388115340e-03
 54588.000925925925912452 -5.074128908454515622e+05  7.257639682023964145e-01  1.874025826380292404e-03
 54588.000983796296168293 -5.074092386012640782e+05  7.351333608427884636e-01  1.873680487441316657e-03
 54588.001041666666878882 -5.074055395130896359e+05  7.445020815182646912e-01  1.873849502509668122e-03
 54588.001099537037134724 -5.074017935789784533e+05  7.538716732272922050e-01  1.873971633320137753e-03
 54588.001157407407390565 -5.073980007962241652e+05  7.632414098560330595e-01  1.873984767500571974e-03
 54588.001215277777646406 -5.073941611626467784e+05  7.726123093411200182e-01  1.874295246964456478e-03
 54588.001273148148356995 -5.073902746728868224e+05  7.819835205798950639e-01  1.874226146744964808e-03
 54588.001331018518612836 -5.073863413272026228e+05  7.913547196412918927e-01  1.874173804634685515e-03
\end{verbatim}


%==================================
\subsection{Matrix}\label{general.fileFormat:matrix}
Stores matrices and vectors. Only one triangle is written for symmetric or triangular matrices.

The header (the matrix definition) is optional.
Therefore a pure text with only numbers in columns are also allowed.
This simplifies the handling of external data.

Instead of a matrix file also an \file{instrument}{instrument} file is allowed.
The first column is the time [MJD], the other columns depends on the instrument type.

\begin{verbatim}
groops matrix version=20200123
LowerSymmetricMatrix( 4 x 4 )
  1.000000000000000000e+00
  0.000000000000000000e+00  1.000000000000000000e+00
  0.000000000000000000e+00  0.000000000000000000e+00  1.000000000000000000e+00
  0.000000000000000000e+00  0.000000000000000000e+00  0.000000000000000000e+00  1.000000000000000000e+00
\end{verbatim}



%==================================
\subsection{MeanPolarMotion}\label{general.fileFormat:meanPolarMotion}
The mean pole of the Earth rotation is represented by a polynomial in a time interval.

\begin{verbatim}
<?xml version="1.0" encoding="UTF-8"?>
<groops type="meanPolarMotion" version="20200123">
  <meanPolarMotion>
    <intervalCount>2</intervalCount>
    <interval>
      <timeStart>42778.0000000000000000</timeStart>
      <degree>3</degree>
      <xp>5.59741000000000e-02</xp>
      <xp>1.82430000000000e-03</xp>
      <xp>1.84130000000000e-04</xp>
      <xp>7.02400000000000e-06</xp>
      <yp>3.46346000000000e-01</yp>
      <yp>1.78960000000000e-03</yp>
      <yp>-1.07290000000000e-04</yp>
      <yp>-9.08000000000000e-07</yp>
    </interval>
    <interval>
      <timeStart>55197.0000000000000000</timeStart>
      <degree>1</degree>
      <xp>2.35130000000000e-02</xp>
      <xp>7.61410000000000e-03</xp>
      <yp>3.58891000000000e-01</yp>
      <yp>-6.28700000000000e-04</yp>
    </interval>
  </meanPolarMotion>
</groops>
\end{verbatim}


%==================================
\subsection{NormalEquation}\label{general.fileFormat:normalEquation}
Stores a \reference{system of normal equations}{normalEquationType}
\begin{equation}
  \M N \hat{\M x} = \M n.
\end{equation}.
This file format consists of multiple files.
The file name \verb|normals.dat.gz| corresponds to the following files:
\begin{itemize}
\item \verb|normals.dat.gz| or \verb|normals.00.00.dat.gz| ... \verb|normals.0n.0n.dat.gz|:
      the normal matrix $\M N$ as \file{matrix}{matrix},
\item \verb|normals.rightHandSide.dat.gz|:
      the right hand side(s) $\M n$ as \file{matrix}{matrix},
\item \verb|normals.parameterNames.txt|: \file{parameter names}{parameterName},
\item \verb|normals.info.xml|:
\     u.a. containing the number of observations and the quadratic sum of (reduced) observations $\M l^T\M P\M l$.
\end{itemize}
A large normal matrix may be splitted into blocks and stored in multiple files.
The block row/column number is indicated in the file name.
Only the upper blocks of the sysmmetric matrix are considered.
Matrix in blocks can be distributed on muliple nodes in parallel mode to efficiently use distributed memory.


%==================================
\subsection{OceanPoleTide}\label{general.fileFormat:oceanPoleTide}
Describes the reaction of the ocean mass to the change
of the centrifugal potential (polar wobble) in terms spherical harmonics.


%==================================
\subsection{ParameterName}\label{general.fileFormat:parameterName}
Name of parameters of a system of \file{normal equations}{normalEquation} or \file{solution vector}{matrix}.

See \program{ParameterNamesCreate}.

\begin{verbatim}
groops parameterName version=20200123
# object:type:temporal:interval
# =============================
      10080 # number of parameters
 karr:position.x::2018-06-01_00-00-00_2018-06-02_00-00-00
 karr:position.y::2018-06-01_00-00-00_2018-06-02_00-00-00
 karr:position.z::2018-06-01_00-00-00_2018-06-02_00-00-00
 karr:troposphereWet:spline.n1:2018-06-01_00-00-00_2018-06-01_02-00-00
 karr:troposphereWet:spline.n1:2018-06-01_00-00-00_2018-06-01_04-00-00
 karr:troposphereWet:spline.n1:2018-06-01_02-00-00_2018-06-01_06-00-00
 karr:troposphereWet:spline.n1:2018-06-01_04-00-00_2018-06-01_08-00-00
 karr:troposphereWet:spline.n1:2018-06-01_06-00-00_2018-06-01_10-00-00
 karr:troposphereWet:spline.n1:2018-06-01_08-00-00_2018-06-01_12-00-00
 karr:troposphereWet:spline.n1:2018-06-01_10-00-00_2018-06-01_14-00-00
 karr:troposphereWet:spline.n1:2018-06-01_12-00-00_2018-06-01_16-00-00
 karr:troposphereWet:spline.n1:2018-06-01_14-00-00_2018-06-01_18-00-00
 karr:troposphereWet:spline.n1:2018-06-01_16-00-00_2018-06-01_20-00-00
 karr:troposphereWet:spline.n1:2018-06-01_18-00-00_2018-06-01_22-00-00
 karr:troposphereWet:spline.n1:2018-06-01_20-00-00_2018-06-02_00-00-00
 karr:troposphereWet:spline.n1:2018-06-01_22-00-00_2018-06-02_00-00-00
 karr:troposphereGradient.x:spline.n1:2018-06-01_00-00-00_2018-06-02_00-00-00
 karr:troposphereGradient.y:spline.n1:2018-06-01_00-00-00_2018-06-02_00-00-00
 karr:troposphereGradient.x:spline.n1:2018-06-01_00-00-00_2018-06-02_00-00-00
 karr:troposphereGradient.y:spline.n1:2018-06-01_00-00-00_2018-06-02_00-00-00
 karr:signalBias01(+1.00L1CG**)::
 karr:signalBias02(+1.00L2WG**)::
 karr:signalBias03(+1.00L2XG**)::
 G01:solarRadiationPressure.ECOM.D0::
 G01:solarRadiationPressure.ECOM.DC2::
 G01:solarRadiationPressure.ECOM.DS2::
 G01:solarRadiationPressure.ECOM.Y0::
 G01:solarRadiationPressure.ECOM.B0::
 G01:solarRadiationPressure.ECOM.BC1::
 G01:solarRadiationPressure.ECOM.BS1::
 G01:stochasticPulse.x::2018-06-01_12-00-00
 G01:stochasticPulse.y::2018-06-01_12-00-00
 G01:stochasticPulse.z::2018-06-01_12-00-00
 G01:arc0.position0.x::
 G01:arc0.position0.y::
 G01:arc0.position0.z::
 G01:arc0.velocity0.x::
 G01:arc0.velocity0.y::
 G01:arc0.velocity0.z::
 G01:signalBias01(-1.00C1CG01)::
 G01:signalBias02(+1.00L1*G01)::
 G01:signalBias03(+1.00L2*G01)::
\end{verbatim}


%==================================
\subsection{Polygon}\label{general.fileFormat:polygon}
List of longitude and latitudes to describe borders, e.g. river basins or continents.
It is used in \configClass{border:polygon}{borderType:polygon}.

\begin{verbatim}
groops polygon version=20200123
          2  # number of polygons
          6  # number of points (1. polygon)
# longitude [deg]           latitude [deg]
# ==================================================
 -1.598200000000000216e+02  2.203000000000000114e+01
 -1.596200000000000045e+02  2.189999999999999858e+01
 -1.593799999999999955e+02  2.189999999999999858e+01
 -1.593000000000000114e+02  2.221999999999999886e+01
 -1.595800000000000125e+02  2.221999999999999886e+01
 -1.598200000000000216e+02  2.203000000000000114e+01
          5  # number of points (2. polygon)
# longitude [deg]           latitude [deg]
# ==================================================
 -7.900000000000000000e+01  2.669999999999999929e+01
 -7.870000000000000284e+01  2.650000000000000000e+01
 -7.823000000000000398e+01  2.667000000000000171e+01
 -7.793000000000000682e+01  2.667000000000000171e+01
 -7.779999999999999716e+01  2.646999999999999886e+01
\end{verbatim}



%==================================
\subsection{PotentialCoefficients}\label{general.fileFormat:potentialCoefficients}
The standard \verb|.gfc| format as defined by the ICGEM is used in ASCII the format.
Only the static part is used and temporal variations (e.g. trend) are ignored.
To write additional information and temporal variations use \program{PotentialCoefficients2Icgem}.


%==================================
\subsection{SatelliteModel}\label{general.fileFormat:satelliteModel}
Properties of a satellite to model non-conservative forces (e.g. \configClass{miscAccelerations}{miscAccelerationsType}).
The file may contain surface properties, mass, drag coefficients, and antenna thrust values.

See \program{SatelliteModelCreate} and \program{SinexMetadata2SatelliteModel}.

\begin{verbatim}
<?xml version="1.0" encoding="UTF-8"?>
<groops type="satelliteModel" version="20190429">
   <satelliteCount>1</satelliteCount>
   <satellite>
       <satelliteName>GALILEO-2</satelliteName>
       <mass>7.00000000000000e+02</mass>
       <coefficientDrag>0.00000000000000e+00</coefficientDrag>
       <surfaceCount>15</surfaceCount>
       <surface>
           <type>0</type>
           <normal>
               <x>-1.00000000000000e+00</x>
               <y>0.00000000000000e+00</y>
               <z>0.00000000000000e+00</z>
           </normal>
           <area>4.40000000000000e-01</area>
           <reflexionVisible>0.00000000000000e+00</reflexionVisible>
           <diffusionVisible>7.00000000000000e-02</diffusionVisible>
           <absorptionVisible>9.30000000000000e-01</absorptionVisible>
           <reflexionInfrared>1.00000000000000e-01</reflexionInfrared>
           <diffusionInfrared>1.00000000000000e-01</diffusionInfrared>
           <absorptionInfrared>8.00000000000000e-01</absorptionInfrared>
           <hasThermalReemission>1</hasThermalReemission>
       </surface>
       ...
       <modulCount>2</modulCount>
       <modul>
           <type>1</type>
           <rotationAxis>
               <x>0.00000000000000e+00</x>
               <y>1.00000000000000e+00</y>
               <z>0.00000000000000e+00</z>
           </rotationAxis>
           <normal>
               <x>0.00000000000000e+00</x>
               <y>0.00000000000000e+00</y>
               <z>1.00000000000000e+00</z>
           </normal>
           <surface>
               <count>4</count>
               <cell>11</cell>
               <cell>12</cell>
               <cell>13</cell>
               <cell>14</cell>
           </surface>
       </modul>
       <modul>
           <type>2</type>
           <antennaThrust>
               <x>0.00000000000000e+00</x>
               <y>0.00000000000000e+00</y>
               <z>2.65000000000000e+02</z>
           </antennaThrust>
       </modul>
   </satellite>
</groops>
\end{verbatim}


%==================================
\subsection{StringList}\label{general.fileFormat:stringList}
White space separated list of strings.
Comments are allowed and all the text from the character '\verb|#|' to the end of the line is ignored.
Strings containg white spaces or the '\verb|#|' character must be set in quotes ('\verb|""|').

\begin{verbatim}
# IGSR3 stations
abmf
abpo
ade1
adis
ajac
albh
algo
alic
alrt
amc2
aoml
areq
arev
artu
asc1
\end{verbatim}


%==================================
\subsection{StringTable}\label{general.fileFormat:stringTable}
White space separated table of strings in row and columns.
Additional columns in a row may represent alternatives, e.g. for core stations in a GNSS network.
Comments are allowed and all the text from the character '\verb|#|' to the end of the line is ignored.
Strings containing white spaces or the '\verb|#|' character must be set in quotes  ('\verb|""|').

\begin{verbatim}
# core network with alternative stations
artu mdvj mdvo nril
asc1 sthl
bahr bhr1 yibl nama
chat chti auck
chpi braz ufpr savo
ckis nium
coco xmis dgar dgav
cro1 scub abmf lmmf aoml
daej taej suwn osn1
darw kat1 tow2 alic
dav1 maw1
drao albh will holb nano
fair whit
glps guat
gode godz usno usn3
goug
\end{verbatim}


%==================================
\subsection{TideGeneratingPotential}\label{general.fileFormat:tideGeneratingPotential}

\begin{verbatim}
groops tideGeneratingPotential version=20200123
       3606
# Dood.   cos                       sin                      name
# ===============================================================
 055.563 -3.122600001000726621e-06  0.000000000000000000e+00 ""
 055.565  7.719644799947265879e-02  0.000000000000000000e+00 om1
 055.573  1.975999999999999959e-07  0.000000000000000000e+00 ""
 055.575 -7.535264999889109729e-04  0.000000000000000000e+00 om2
 055.654 -4.037500000326006771e-06  0.000000000000000000e+00 ""
 055.666  6.590000000000000115e-08  0.000000000000000000e+00 ""
 055.753  6.023000000000000398e-07  0.000000000000000000e+00 ""
 056.475  7.920000000000000775e-08  0.000000000000000000e+00 ""
 056.554 -1.360322439950949897e-02  0.000000000000000000e+00 sa
 056.556  7.156881999672112154e-04  0.000000000000000000e+00 ""
 056.564  8.641769999583327993e-05  0.000000000000000000e+00 ""
\end{verbatim}



%==================================
\subsection{TimeSplinesCovariance}\label{general.fileFormat:timeSplinesCovariance}
Stores covariance information for \file{TimeSplinesGravityField}{timeSplinesGravityField}.
It can be the variances of the potential coefficients or the full covariance matrix for each
temporal nodal point.


%==================================
\subsection{TimeSplinesGravityField}\label{general.fileFormat:timeSplinesGravityField}
Temporal changing gravity field, parametrized as spherical harmonics in the spatial domain and
parametrized as basis splines in the time domain (see~\reference{basis splines}{fundamentals.basisSplines}).
It is evaluated with \configClass{gravityfield:timeSplines}{gravityfieldType:timeSplines}.

See also: \program{Gravityfield2TimeSplines}, \program{PotentialCoefficients2BlockMeanTimeSplines}.


%==================================
\subsection{VariationalEquation}\label{general.fileFormat:variationalEquation}
Arcs with reference orbit and state transition matrices.

The file contains a reference orbit (position and velocity),
the derivatives of the orbit with respect to the satellite state vector for each arc,
transformations (rotations) between the satellite, celestial, and terrestrial frame
and a satellite macro model (see \file{SatelliteModel}{satelliteModel}).

The reference orbit can be extracted with \program{Variational2Orbit}.

See also: \program{PreprocessingVariationalEquation}.



%==================================

\clearpage

\chapter{Mathematical fundamentals}
This chapter describes the methods and mathematical fundamentals used throughout GROOPS.
\section{Robust least squares adjustment}\label{fundamentals.robustLeastSquares}
The robust least squares adjustment used in GROOPS is based on a modified Huber
estimator. It downweights observations with large otuliers iteratively.

The algorithm starts with a first solution with equal weights $\M P =\M I$
\begin{equation}
  \hat{\M x} = (\M A^T\M P\M A)^{-1}\M A^T\M P\M l.
\end{equation}
The solution is used to compute the residuals
\begin{equation}
  \hat{e}_i = \left(\M l - \M A \hat{\M x}\right)_i
\end{equation}
and the redundancies of all observations
\begin{equation}
  r_i = \left(\M I - \M A(\M A^T\M P\M A)^{-1}\M A^T\M P\right)_{ii}.
\end{equation}
For observations with large residuals a new standard deviation is assigned
\begin{equation}
  \sigma_i =
  \begin{cases}
    1                                     & \text{for } \left|\frac{\hat{e}_i}{r_i}\right| \le h\cdot\hat{\sigma} \\
    \left|\frac{\hat{e}_i}{r_ih}\right|^p & \text{for } \left|\frac{\hat{e}_i}{r_i}\right| > h\cdot\hat{\sigma},
  \end{cases}
\end{equation}
where $h$ is \config{huber}, $p$ is \config{huberPower}, and $\hat{\sigma}^2$ a robust overall variance factor
computed from all residuals.
The estimation is repeated \config{huberMaxIteration} times with a new weight matrix
\begin{equation}
  \M P = \text{diag}\left(\frac{1}{\sigma_1^2}, \frac{1}{\sigma_2^2}, \ldots, \frac{1}{\sigma_n^2}\right)
\end{equation}
 or until convergence is reached.

\clearpage
\section{Basis splines}\label{fundamentals.basisSplines}
A time variable function is given by
\begin{equation}
f(x,t) =  \sum_i f_i(x)\Psi_i(t),
\end{equation}
with the (spatial) coefficients $f_i(x)$ as parameters and the temporal basis functions~$\Psi_i(t)$.
Basis splines are defined as polynomials of degree~$n$ in intervals between nodal points in time $t_i$:
\begin{itemize}
\item Block mean values ($n=0$)
\begin{equation}
  \Psi_i(t) = \begin{cases}
  1 & \text{if } t\in[t_i,t_{i+1}), \\
  0 & \text{otherwise}
\end{cases}
\end{equation}
\item Linear splines ($n=1$)
\begin{equation}
  \Phi_i(t) = \begin{cases}
  \tau_{i-1}   & \text{if } t_{i-1}\le t \le t_i, \\
  1-\tau_i     & \text{if } t_{i}\le t \le t_{i+1}, \\
  0 & \text{otherwise}.
\end{cases}
\end{equation}
\item Quadratic splines ($n=2$)
\begin{equation}
  \Phi_i(t) = \begin{cases}
  \frac{1}{2}\tau^2_{i-1}   & \text{if } t_{i-1}\le t \le t_i, \\
  -\tau^2_{i-1}+\tau_{i-1}+\frac{1}{2}     & \text{if } t_{i}\le t \le t_{i+1}, \\
  \frac{1}{2}\tau^2_{i-1}-\tau_{i-1}+\frac{1}{2}     & \text{if } t_{i}\le t \le t_{i+1}, \\
  0 & \text{otherwise}.
\end{cases}
\end{equation}
\item Cubic splines ($n=3$)
\begin{equation}
  \Phi_i(t) = \begin{cases}
   \frac{1}{6}\tau^3                                                        & \text{if } t_{i-1}\le t \le t_i, \\
  -\frac{3}{6}\tau^3 +\frac{3}{6}\tau^2 +\frac{3}{6}\tau  +\frac{1}{6} & \text{if } t_{i-1}\le t \le t_i, \\
   \frac{3}{6}\tau^3 -            \tau^2                     +\frac{4}{6} & \text{if } t_{i-1}\le t \le t_i, \\
  -\frac{1}{6}\tau^3 +\frac{3}{6}\tau^2 -\frac{3}{6}\tau  +\frac{1}{6} & \text{if } t_{i-1}\le t \le t_i, \\
  0 & \text{otherwise}.
\end{cases}
\end{equation}
\end{itemize}
where $\tau$ is the normlized time in each time interval
\begin{equation}
  \tau_i = \frac{t-t_i}{t_{i+1}-t_i}.
\end{equation}
The total number of coefficients $f_i(x)$ is $N=N_t+n-1$,
where $N_t$ is the count of nodal time points~$t_i$ and $n$ is the degree.

\fig{!hb}{0.8}{basissplines}{fig:basissplines}{Basis splines for different degrees with nodal points every 6 hours.}

\clearpage
\section{Autoregressive Models}\label{fundamentals.autoregressiveModel}
A multivariate (or vector) autoregressive model is one possible representation of a random process.
It specifies, that the output at epoch $t$ depends on the $p$ previous epochs, where $p$ is denoted process order,
plus a stochastic term.
In the following, finite order vector autoregressive - VAR($p$) in short - models as implemented in GROOPS will be described.

\subsection{Definition}

A finite order VAR($p$) model is defined as
\begin{equation}
  \mathbf{y}_e(t_i) = \sum_{k=1}^p \mathbf{\Phi}^{(p)}_k\mathbf{y}_e(t_{i-k}) + \mathbf{w}(t_i),
  \hspace{5pt} \mathbf{w}(t_i) \sim \mathcal{N}(0, \mathbf{\Sigma}^{(p)}_\mathbf{w}),
\end{equation}
where $\mathbf{y}_e(t_i)$ are realizations of a random vector process
Subtracting the right hand side and substituting the stochastic term $-\mathbf{w}(t_i)$ with the residual $\mathbf{v}(t_i)$ gives us
\begin{equation}
  \mathbf{0}  = \mathbf{y}_e(t_i) - \sum_{k=1}^p \mathbf{\Phi}^{(p)}_k\mathbf{y}_e(t_{i-k}) + \mathbf{v}(t_i)
\end{equation}
which can be used as pseudo-observation equations in the determination of the parameters $\mathbf{y}_e(t_i)$.
In matrix notation this reads
\begin{equation}
  0 =
  \begin{bmatrix}
    \mathbf{I} & -\mathbf{\Phi}^{(p)}_1 & \cdots & -\mathbf{\Phi}^{(p)}_p \\
  \end{bmatrix}
  \begin{bmatrix}
    \mathbf{y}_e(t_i) \\
    \mathbf{y}_e(t_{i-1}) \\
    \vdots \\
    \mathbf{y}_e(t_{i-p}) \\
  \end{bmatrix}
  + \mathbf{v}(t_i).
\end{equation}
After rearranging the vectors $\mathbf{x}_t$ to have ascending time stamps
\begin{equation}
  0 =
  \begin{bmatrix}
    -\mathbf{\Phi}^{(p)}_p & \cdots & -\mathbf{\Phi}^{(p)}_1 & \mathbf{I} \\
  \end{bmatrix}
  \begin{bmatrix}
    \mathbf{y}_e(t_{i-p}) \\
    \vdots \\
    \mathbf{y}_e(t_{i-1}) \\
    \mathbf{y}_e(t_i) \\
  \end{bmatrix}
  + \mathbf{v}(t_i)
\end{equation}
For practical purposes, the residuals above are further decorrelated using the
inverse square root of the white noise covariance matrix, leading to
\begin{equation}
  \bar{\mathbf{v}}(t_i) = \underbrace{\mathbf{\Sigma}^{(p)^{-\frac{1}{2}}}_\mathbf{w}}_{=\mathbf{W}}\mathbf{v}(t_i), \hspace{25pt}  \bar{\mathbf{v}}(t_i) \sim \mathcal{N}(0, \mathbf{I}).
\end{equation}
The used square root is in principle arbitrary, but should satisfy $\mathbf{W}^T\mathbf{W} = \mathbf{\Sigma}^{(p)}_\mathbf{w} $.
This means that both eigendecomposition based roots and Cholesky factors can be used.
After the applying the matrix from the left, we arrive at the observation equations
\begin{equation}
  0 =
  \begin{bmatrix}
    -\mathbf{W}\mathbf{\Phi}^{(p)}_p & \cdots & -\mathbf{W}\mathbf{\Phi}^{(p)}_1 & \mathbf{W} \\
  \end{bmatrix}
  \begin{bmatrix}
    \mathbf{y}_e(t_{i-p}) \\
    \vdots \\
    \mathbf{y}_e(t_{i-1}) \\
    \mathbf{y}_e(t_i) \\
  \end{bmatrix}
  + \bar{\mathbf{v}}(t_i)
\end{equation}
which yields fully decorrelated residuals.
Currenty, VAR($p$) models are saved to a single  file which contains this matrix.


\chapter{Cookbook}\label{general.cookbook}
This chapter presents recipes to explore the GROOPS feature set.
\section{Instrument data handling}\label{cookbook.instrument}

GROOPS provides functions and programs to read/write, preprocess, analyze and visualize uniformly and non-uniformly sampled instrument data.
This includes tools for filter design and analysis, re-sampling, smoothing, detrending, and power spectrum estimation.
This tutorial goes through exemplary steps for data handling procedures.

\subsection{Reading data}

\begin{itemize}
\item GROOPS is able to read and convert relevant data from various LEO and GNSS satellites. Instrument files need to be converted into the respective GROOPS format using conversion programs.Depending on the content of the input file, the data is stored with a specific \configClass{instrument type}{instrumentTypeType}. User also has the option to change the type later on with \program{InstrumentSetType}.

\item Multiple files can be concatenated to one file using \program{InstrumentConcatenate}. Using this program, it is also possible to sort the epochs, remove the duplicates and NaN values.
\item \emph{Example: Concatenating instrument files}
  \begin{itemize}
  \item Create three successively daily sinusoidal signals with \program{TimeSeriesCreate} and set their type to MISCVALUE with \program{InstrumentSetType}. In this example, each data set has an overlap of 1 hour with their following dataset.
  \item Merge all datasets to one single file with \program{InstrumentConcatenate}.
  \fig{!ht}{0.8}{instrumentDataHandeling_concatenating}{fig:instrumentDataHandeling_concatenating}{Example 1: Concatenating instrument files into one dataset.}
  \end{itemize}
 \item Many measurements involve data collected asynchronously by multiple sensors with different sampling. Use \program{InstrumentSynchronize} for a continuous harmonization of the data over time or segmentation of the data into arcs.
\end{itemize}

\subsection{Preprocessing}
Real-world data is often incomplete, inconsistent, and/or lacking in certain behaviors or trends, and is likely to contain many errors. Data preprocessing is a proven method of resolving such issues. Following steps are usually required to be taken:

\begin{itemize}
\item Gross outlier removal:
  \begin{itemize}
  \item Create reference values to compare the input data with. Depending on the instrument type, this can be done by simulation programs such as \program{SimulateAccelerometer} or \program{SimulateStarCamera}. If no reference data is available, the outlier detection is based on the data itself. If needed, synchronize the reference data file and the input data with \program{InstrumentSynchronize}.

  \item In case of star camera data, compute the differences between the input data and the reference data with \program{InstrumentStarCameraMultiply}.

  \item Set a threshold for outlier detection in \program{InstrumentRemoveEpochsByCriteria}. The threshold is defined empirically according to the accuracy characteristics of each data products. If the differences exceed a predefined threshold, the corresponding epochs are removed. An arbitrary margin can be defined to additionally remove epochs before and after the identified outliers. It is also possible to remove epochs at specific times using \program{InstrumentRemoveEpochsByTimes}.
  \item Missing epochs can be filled by reference data with \program{InstrumentConcatenate}.
  \item It is also possible to interpolate the missing epochs with \program{InstrumentResample}.
  \item \emph{Example: Removing outliers in a synthetic data.}
  \begin{itemize}
  \item Create a sinusoidal signal with an amplitude of 1.0 using \program{TimeSeriesCreate} and set its type to MISCVALUE with \program{InstrumentSetType}.
  \item Add zero-mean, white Gaussian noise with a standard deviation of 0.1 with \program{NoiseInstrument}. Interpret this data as a real measurement file.
  \item Set the threshold criteria to 0.2 in \program{InstrumentRemoveEpochsByCriteria} and remove the outliers and their nearest epochs in 20 second interval.
  \item Fill the data gaps with \program{InstrumentResample}.
  \fig{!ht}{0.8}{instrumentDataHandeling_outlierRemoval}{fig:instrumentDataHandeling_outlierRemoval}{Example 2: Removing gross outliers.}
  \end{itemize}
  \end{itemize}


\item Downsampling:
  \begin{itemize}
  \item If the sampling is irregular use \program{InstrumentResample} to make the sampling equidistant.
  \item Use \program{InstrumentSynchronize} to divide the data at gaps into arcs.
  \item Apply a lowpass filter (e.g. Butterworth) with the Nyquist frequency of the target sampling as cutoff with \program{InstrumentFilter}. Apply the filter in both directions to avoid phase shifts.
  \item Use \program{InstrumentReduceSampling} to down-sample the data.
  \end{itemize}

\item Calibration:
  \begin{itemize}
  \item For a general instrument file, \program{InstrumentDetrend} subtracts offsets or linear/nonlinear trends from the input data. This can be achieved also with \program{FunctionsCalculate} or \program{InstrumentArcCalculate} by applying determined calibration factors or solving a least-square adjustment.
  \item For accelerometer data, \program{InstrumentAccelerometerEstimateBiasScale} is designed to estimate and subtract complex biases or scales with respect to simulated accelerometer data. If a thruster file is given, the corresponding epochs are eliminated during estimation process.
  \item \emph{Example: GRACE-C accelerometer calibration}
  \begin{itemize}
  \item For one particular date, read and convert Level-1B GRACE-C orbit, star camera, accelerometer, and thruster data with \program{GraceL1b2Orbit}, \program{GraceL1b2StarCamera}, \program{GraceL1b2Accelerometer}, and \program{GraceL1b2Thruster} respectively. It is also required to read the macro-model data of the satellite using the related information in the official document and convert it to GROOPS format with \program{SatelliteModelCreate}.
  \item Use \program{SimulateAccelerometer} to generate simulated accelerations due to non-gravitational force models including: \configClass{miscAccelerations:atmosphericDrag}{miscAccelerationsType:atmosphericDrag}, \configClass{miscAccelerations:solarRadiationPressure}{miscAccelerationsType:solarRadiationPressure}, and \configClass{miscAccelerations:albedo}{miscAccelerationsType:albedo}.
  \item Calibrate the real measurements with a daily constant accelerometer bias by choosing a constant parameter per axis in \configClass{parametrizationAcceleration:accBias}{parametrizationAccelerationType:accBias}.
  \fig{!ht}{0.8}{instrumentDataHandeling_calibration}{fig:instrumentDataHandeling_calibration}{Example 3: Calibrating GRACE-C ACT1B data.}
  \end{itemize}
  \end{itemize}
\end{itemize}

\subsection{Statistical analysis}
\begin{itemize}
\item \program{InstrumentStatisticsTimeSeries} returns statistics for one or more instrument files. \program{InstrumentArcCalculate} is also able to generate a \config{statistics} file with one mid epoch per arc.
\end{itemize}

\subsection{Spectral analysis}
Spectral analysis studies the frequency spectrum contained in discrete, uniformly sampled data. The Fourier transform is a tool that reveals frequency components of a signal by representing it in frequency space. The Power Spectral Density (PSD) is a measurement of the energy at each frequency.

\begin{itemize}
\item If the sampling is irregular use \program{InstrumentResample} to make the sampling equidistant.
\item Use \program{Instrument2PowerSpectralDensity} to compute PSD.
\item If covariance function of a dataset is available, use \program{CovarianceFunction2PowerSpectralDensity}.
\item \emph{Example: Spectral analysis of a synthetic signal.}
  \begin{itemize}
  \item Create a sinusoidal signal with an amplitude of 1.0 using \program{TimeSeriesCreate} and set its type to MISCVALUE with \program{InstrumentSetType}. Interpret this data as a simulation data file.
  \item Add zero-mean, white Gaussian noise with a standard deviation of 0.1 with \program{NoiseInstrument}. Interpret this data as a real measurement file.
  \item Compute PSD of the simulated and measurement data and represent the results with \program{PlotGraph}.
  \fig{!ht}{0.8}{instrumentDataHandeling_psd}{fig:instrumentDataHandeling_psd}{Example 4: Spectral analysis of a synthetic signal.}
  \end{itemize}
\end{itemize}

\subsection{Data visualization}
\begin{itemize}
\item Argument of latitude plot
Plotting instrument data as a function of satellite position in orbit and time reveals features related to the orbit geometry or environmental conditions. For circular orbits, the position of  satellite can be specified by the argument of latitude.
\begin{itemize}
\item Synchronize the instrument data file with the related orbit data using \program{InstrumentSynchronize}.
\item Use \program{Orbit2ArgumentOfLatitude} to compute argument of latitude at each epoch.
\item Plot the instrument data versus argument of latitude and time with \program{PlotGraph}.
\item \emph{Example: Argument of latitude representation of GRACE-C eclipse factors}
  \begin{itemize}
  \item Compute eclipse factors at each epoch of GRACE-C orbit at an arbitrary time using \program{Orbit2EclipseFactor}.
  \item Synchronize the eclipse factor data file with the related orbit data using \program{InstrumentSynchronize}.
  \item Use \program{Orbit2ArgumentOfLatitude} to compute argument of latitude at each epoch and visualize the results with \program{PlotGraph}.
  \fig{!ht}{0.8}{instrumentDataHandeling_aol}{fig:instrumentDataHandeling_aol}{Example 5: GRACE-C eclipse factors represented in argument of latitude plot.}
  \end{itemize}
\end{itemize}

\item Ground-track plot
Plotting instrument data with respect to the satellite ground track is useful to identify any features of geophysical origin in the data.
\begin{itemize}
\item Synchronize the instrument data file with the related orbit data using \program{InstrumentSynchronize}.
\item Use \program{Orbit2Groundtracks} to map instrument data to satellite ground-track.
\item Visualize the output with \program{PlotMap}.
\item \emph{Example 6: Ground-track representation of GRACE-C eclipse factors}
  \begin{itemize}
  \item Compute eclipse factors at each epoch of GRACE-C orbit at an arbitrary time using \program{Orbit2EclipseFactor}.
  \item Synchronize the eclipse factor data file with the related orbit data using \program{InstrumentSynchronize}.
  \item Use \program{Orbit2Groundtracks} to generate the gridded data. Each grid value represents the mean value of eclipse factor over the instrument time period (1 month).  visualize the results with \program{PlotMap}.
  \end{itemize}
  \fig{!ht}{0.8}{instrumentDataHandeling_groundTrack}{fig:instrumentDataHandeling_groundTrack}{Example 6: GRACE-C eclipse factors represented in ground-track plot.}
\end{itemize}

\end{itemize}
%% =============================================================

\clearpage
\section{GNSS satellite orbit determination and station network analysis}\label{cookbook.gnssNetwork}
This cookbook chapter describes an example of global GNSS processing as done by analysis centers of the
International GNSS Service (IGS). Resulting products usually comprise:
\begin{itemize}
  \item Satellite orbits, clocks, and signal biases
  \item Station positions, clocks, signal biases, and troposphere estimates
  \item Earth orientation parameters
\end{itemize}

Scientific details about the underlying processing approach and the applied parametrizations, models, and corrections
can be found in a doctoral thesis available under DOI \href{https://doi.org/10.3217/978-3-85125-885-1}{10.3217/978-3-85125-885-1}.

An example scenario for this task is available at \url{https://ftp.tugraz.at/outgoing/ITSG/groops/scenario/scenarioGnssNetwork.zip}.
It includes GROOPS scripts and data for the example, but not the general GROOPS data and metadata found at \url{https://ftp.tugraz.at/outgoing/ITSG/groops} (data folder or zipped archive).
The scenario generally represents what is described in this cookbook, but may slightly differ in certain settings.

\emph{Note: Global GNSS processing can become very computationally intensive. Depending on the number of satellites
and stations, the observation and processing sampling, and parametrizations it can quickly exceed the capabilities
of a normal desktop computer and may require computer clusters or number crunchers (see section}
\reference{Parallelization}{general.parallelization}\emph{).}

\subsection{Data preparation}\label{cookbook.gnssNetwork:metadata}
Most of the required metadata files are provided in GROOPS file formats at \url{https://ftp.tugraz.at/outgoing/ITSG/groops}.
These files are regularly updated.

Data that has to be gathered from other sources comprises:
\begin{itemize}
  \item \textbf{Receiver observations}: GNSS measurements converted from RINEX format (see \program{RinexObservation2GnssReceiver}).
  \item \textbf{Approximate orbits}: broadcast or precise orbits in CRF for orbit integration (see \program{GnssRinexNavigation2OrbitClock} or \program{Sp3Format2Orbit}).
  \item \textbf{Approximate clocks}: broadcast or precise clocks (see \program{GnssRinexNavigation2OrbitClock} or \program{GnssClockRinex2InstrumentClock})
\end{itemize}
Receiver observations, broadcast ephemerides, and precise satellite orbits and clocks can be downloaded from the
\href{https://igs.org/data-products-overview/}{IGS Data Centers}.
GPS, GLONASS, and Galileo orbits and clocks for the period 1994-2020 are also available as part of
\href{https://doi.org/10.3217/dataset-4528-0723-0867}{Graz University of Technology's contribution to IGS repro3}.

The \href{https://ftp.tugraz.at/outgoing/ITSG/groops/scenario/scenarioGnssNetwork.zip}{example scenario} includes a small set of this data.
The script \verb|010groopsConvert.xml| can be used to convert these external formats into GROOPS formats.

Prepare a \file{station list file}{stringTable} that contains the stations to be processed.
Each line can contain more than one station. The first station in each line that has data available is used for the processing.
If your network contains more than 60-70 stations, it is recommended to start processing with a core network (see \reference{Advanced}{cookbook.gnssNetwork:advanced}).
In this case, define an additional \file{core station list file}{stringTable} that can also have multiple stations per line.

\subsection{Preprocessing: Orbit integration}\label{cookbook.gnssNetwork:orbitIntegration}
Numerical integration of the satellite orbits is the first step in global GNSS processing.
Dynamic orbits are integrated based on \configClass{force models}{forcesType} and then fitted to the approximate orbits
by estimating their initial state and additional empirical parameters for solar radiation pressure to improve the orbit fit.
The resulting \file{variational equations}{variationalEquation} file contains the integrated orbit, derivatives
with respect to the satellite state vector, attitude, Earth rotation and satellite model.

Orbit preprocessing is covered by the script \verb|020groopsGnssPreprocessing.xml| in the \href{https://ftp.tugraz.at/outgoing/ITSG/groops/scenario/scenarioGnssNetwork.zip}{example scenario}.

It is recommended to perform the steps below in a \reference{loop}{general.loopsAndConditions} over all
satellites/PRNs using \program{LoopPrograms}.

\begin{itemize}
  \item \program{GnssPrn2SvnBlockVariables}
  \item \program{InstrumentResample}: resample approximate orbits from \reference{data preparation}{cookbook.gnssPpp:metadata} to target sampling (e.g., 1 minute) by defining a \configClass{timeSeries}{timeSeriesType} based on a \config{method:polynomial} (\config{polynomialDegree}=\verb|7|, \config{maxDataPointRange}=\verb|7200|, \config{maxExtrapolationDistance}=\verb|900|).
  \item \program{OrbitAddVelocityAndAcceleration}: add velocity via running polynomial (\config{polynomialDegree}=\verb|2|) derivation (needed for attitude computation)
  \item \program{SimulateStarCameraGnss}
  \item \program{PreprocessingVariationalEquation}:
        \begin{itemize}
          \item \configFile{inputfileSatelliteModel}{satelliteModel}=\verb|{groopsDataDir}/gnss/transmitter/satelliteModel/satelliteModel_boxWing.{svn}.xml|
          \item \configFile{inputfileOrbit}{instrument}: the resampled approximate orbit from \program{InstrumentResample}
          \item \configFile{inputfileStarCamera}{instrument}: the attitude file from \program{SimulateStarCameraGnss}
          \item \config{forces}: see below
          \item \configClass{gradientfield:potentialCoefficients}{gravityfieldType:potentialCoefficients}: a static gravity field (e.g. GOCO06s) with \config{maxDegree}=\verb|4|.
        \end{itemize}
  \item \program{PreprocessingVariationalEquationOrbitFit}: fit the integrated orbit (\configFile{inputfileVariational}{variationalEquation}) to the approximate orbit (\configFile{inputfileOrbit}{instrument}) by least squares adjustment.
        Add \configClass{parametrizationAcceleration:gnssSolarRadiation}{parametrizationAccelerationType:gnssSolarRadiation}
        and select the \href{https://doi.org/10.1007/s00190-015-0814-4}{ECOM2} parameters to be estimated.
\end{itemize}

Force models usually include:
\begin{itemize}
  \item \configClass{gravityfield:potentialCoefficients}{gravityfieldType:potentialCoefficients}:
        static gravity field (e.g. GOCO06s)
  \item \configClass{gravityfield:trend}{gravityfieldType:trend}
  \begin{itemize}
    \item \configClass{gravityfield:potentialCoefficients}{gravityfieldType:potentialCoefficients}:
          trend component of time-variable gravity field (e.g. GOCO06s)
  \end{itemize}
  \item \configClass{gravityfield:oscillation}{gravityfieldType:oscillation}
  \begin{itemize}
    \item \configClass{gravityfieldCos:potentialCoefficients}{gravityfieldType:potentialCoefficients}:
          annual cosine component of time-variable gravity field (e.g. GOCO06s)
    \item \configClass{gravityfieldSin:potentialCoefficients}{gravityfieldType:potentialCoefficients}:
          annual sine component of time-variable gravity field (e.g. GOCO06s)
  \end{itemize}
  \item \configClass{tides:astronomicalTide}{tidesType:astronomicalTide}: astronomical tides (e.g. based on JPL ephemeris)
  \item \configClass{tides:earthTide}{tidesType:earthTide}: Earth tide (IERS conventions)
  \item \configClass{tides:doodsonHarmonicTide}{tidesType:doodsonHarmonicTide}: ocean tides (e.g. FES 2014b)
  \item \configClass{tides:poleTide}{tidesType:poleTide}: pole tides (IERS conventions)
  \item \configClass{tides:poleOceanTide}{tidesType:oceanPoleTide}: ocean pole tides (IERS conventions)
  \item \configClass{miscAccelerations:solarRadiationPressure}{miscAccelerationsType:solarRadiationPressure}:
        solar radiation pressure (box-wing model)
  \item \configClass{miscAccelerations:albedo}{miscAccelerationsType:albedo}: Earth radiation pressure (albedo model)
  \item \configClass{miscAccelerations:antennaThrust}{miscAccelerationsType:antennaThrust}:
        antenna thrust (e.g. from IGS metadata SINEX file)
  \item \configClass{miscAccelerations:relativisticEffect}{miscAccelerationsType:relativisticEffect}:
        relativistic effects (IERS conventions)
\end{itemize}
For the spherical harmonics expansions a \config{maxDegree}=\verb|60| is more than enough.

The result of the preprocessing should be a \file{variational equations file}{variationalEquation},
a \file{reduced dynamic orbit file}{instrument} from \program{PreprocessingVariationalEquationOrbitFit}
and an \file{attitude file}{instrument} from \program{SimulateStarCameraGnss} for each satellite.

\subsection{GNSS processing}\label{cookbook.gnssNetwork:processing}
The script \verb|030groopsGnssProcessing.xml| in the \href{https://ftp.tugraz.at/outgoing/ITSG/groops/scenario/scenarioGnssNetwork.zip}{example scenario}
implements the following steps and settings.

These are the settings for \program{GnssProcessing}. If not otherwise stated use the default values.

The first step is setting the processing sampling, in this example it is 30~seconds.
The processing interval usually is a single 24-hour day, therefore define
\configClass{timeSeries:uniformSampling}{timeSeriesType:uniformSampling} with \config{timeStart}=\verb|<mjd>|,
\config{timeEnd}=\verb|<mjd>+1|, \config{sampling}=\verb|30/86400| (processing sampling).

Add the appropriate \configClass{transmitters:gnss}{gnssTransmitterGeneratorType:gnss} (e.g. GPS, GLONASS, and Galileo)
and provide the required files:
\begin{itemize}
  \item \configFile{inputfileOrbit}{instrument} from \reference{preprocessing}{cookbook.gnssNetwork:orbitIntegration}
  \item \configFile{inputfileAttitude}{instrument} from \reference{preprocessing}{cookbook.gnssNetwork:orbitIntegration}
  \item \configFile{inputfileClock}{instrument} from \reference{data preparation}{cookbook.gnssPpp:metadata}
\end{itemize}

The following settings are needed in \configClass{receiver:stationNetwork}{gnssReceiverGeneratorType:stationNetwork}:
\begin{itemize}
   \item \configFile{inputfileStationList}{stringTable}: list of all stations to be processed
   \item \configFile{inputfileObservations}{instrument}: The converted RINEX observation files.
   \item \configClass{tidalDisplacement}{tidesType}: Use the settings described in \reference{receiver:stationNetwork}{gnssReceiverGeneratorType:stationNetwork}.
   \item \configClass{excludeType}{gnssType}: Signals you might want to exclude are \verb|C*?G| (old unknown GPS code observations), \verb|*3*R| (GLONASS G3 freq.), \verb|*6*E| (Galileo E6 freq.).
\end{itemize}

Add the following \configClass{parametrizations}{gnssParametrizationType}
and define the \config{outputfiles} you are interested in inside each of them:
\begin{itemize}
  \item \configClass{ionosphereSTEC}{gnssParametrizationType:ionosphereSTEC}: add a constraint of \config{sigmaSTEC}=\verb|40|
  \item \configClass{ionosphereMap}{gnssParametrizationType:ionosphereMap}: add \configClass{temporal:splines}{parametrizationTemporalType:splines}
        with linear (\config{degree}=\verb|1|) 2-hourly splines
  \item \configClass{clocks}{gnssParametrizationType:clocks}: optionally change \configClass{selectTransmitterZeroMeans}{gnssTransceiverSelectorType} to \configClass{wildcard}{gnssTransceiverSelectorType:wildcard} with \config{name}=\verb|G*| to align clocks to mean over GPS (instead of all) satellites
  \item \configClass{signalBiases}{gnssParametrizationType:signalBiases}
  \item \configClass{ambiguities}{gnssParametrizationType:ambiguities}
  \item \configClass{codeBiases}{gnssParametrizationType:codeBiases}
  \item \configClass{tecBiases}{gnssParametrizationType:tecBiases}
  \item \configClass{temporalBias}{gnssParametrizationType:temporalBias}: time-variable GPS L5 phase bias with \configClass{type}{gnssType}=\verb|L5*G|
        and \configClass{parametrizationTemporal:splines}{parametrizationTemporalType:splines} with degree 3 and hourly nodes.
  \item \configClass{staticPositions}{gnssParametrizationType:staticPositions}
  \item \configClass{troposphere}{gnssParametrizationType:troposphere}: select \configClass{troposphere:viennaMapping}{troposphereType:viennaMapping}
        with the appropriate vmf3grid file. Add \configClass{troposphereWetEstimation:splines}{parametrizationTemporalType:splines}
        with linear (\config{degree}=\verb|1|) 2-hourly splines
        and \configClass{troposphereGradientEstimation:splines}{parametrizationTemporalType:splines} with linear daily splines.
  \item \configClass{transmitterDynamicOrbit}{gnssParametrizationType:transmitterDynamicOrbits}:
        provide \configFile{inputfileVariational}{variationalEquation}=\verb|preprocessing/variational.{prn}.dat|
        from the preprocessing step.
        Add \configClass{parametrizationAcceleration:gnssSolarRadiation}{parametrizationAccelerationType:gnssSolarRadiation} and
        \configClass{stochasticPulse:irregular}{timeSeriesType:irregular} parameter at center of day to further improve orbit fit.
  \item \configClass{earthRotation}{gnssParametrizationType:earthRotation}
        \begin{itemize}
          \item \configClass{estimatePole:constant}{parametrizationTemporalType:constant}: polar motion
          \item \configClass{estimatePole:trend}{parametrizationTemporalType:trend} polar motion rate
                (at center of day with \config{timeStep}=\verb|1|)
          \item \configClass{estimateUT1:trend}{parametrizationTemporalType:trend}: length of day (at center of day with
                \config{timeStep}=\verb|-1| to match IGS sign convention)
        \end{itemize}
    \item \configClass{constraints}{gnssParametrizationType:constraints}: loose constraint on GPS L5 phase biases,
        \configClass{parameters:wildcard:type}{parameterSelectorType:wildcard}=\verb|signalBias.L5*|, \config{sigma}=\verb|5| meters, and \config{relativeToApriori}=\verb|yes|
    \item \configClass{constraints}{gnssParametrizationType:constraints}: loose constraint troposphere estimates,
        \configClass{parameters:wildcard:type}{parameterSelectorType:wildcard}=\verb|troposphere*|, \config{sigma}=\verb|5| meters, and \config{relativeToApriori}=\verb|yes|
    \item \configClass{constraints}{gnssParametrizationType:constraints}: constraint on stochastic pulses,
        \configClass{parameters:wildcard:type}{parameterSelectorType:wildcard}=\verb|stochasticPulse*|, \config{sigma}=\verb|0.1| micrometers/second.
\end{itemize}

Finally, define the \configClass{processingSteps}{gnssProcessingStepType}.
This can be overwhelming at first, but offers a lot of flexibility.
The example script uses a 5-minute processing sampling with subsequent clock densification to 30 seconds.

\begin{itemize}
      \item \configClass{selectEpochs}{gnssProcessingStepType:selectEpochs}: with \config{nthEpoch}=\verb|10| to reduce sampling to 5 minutes.
      \item \configClass{selectParametrizations}{gnssProcessingStepType:selectParametrizations}:
            disable \verb|constraint.STEC|, \verb|*VTEC|, \verb|*.tecBiases| as the ionosphere parameters are estimated in the final steps only.
      \item \configClass{estimate}{gnssProcessingStepType:estimate}: with \config{maxIterationCount}=\verb|6|
      \item \configClass{resolveAmbiguities}{gnssProcessingStepType:resolveAmbiguities}
      \item \configClass{selectParametrizations}{gnssProcessingStepType:selectParametrizations}:
            enable \verb|*| (all) parameters
      \item \configClass{estimate}{gnssProcessingStepType:estimate}: with \config{maxIterationCount}=\verb|4|: final iterations (with 5-minute sampling) and ionosphere parameters
      \item \configClass{group}{gnssProcessingStepType:group}: clock densification to 30-second sampling
      \begin{itemize}
            \item \configClass{selectEpochs}{gnssProcessingStepType:selectEpochs}: with \config{nthEpoch}=\verb|1| to set full 30-second sampling
            \item \configClass{selectParametrizations}{gnssProcessingStepType:selectParametrizations}:
            disable \verb|*| (all) parameters and reenable \verb|*.clock*| and \verb|*.STEC| parameters
            \item \configClass{estimate}{gnssProcessingStepType:estimate}: with \config{maxIterationCount}=\verb|6|
            \item \configClass{writeResults}{gnssProcessingStepType:writeResults}: with \config{suffix}=\verb|30s| to write 30-second clock files
      \end{itemize}
      \item \configClass{writeResults}{gnssProcessingStepType:writeResults}: write the final results
\end{itemize}

With some additional steps, the full 30-second sampling can be used to estimate all parameters (not only the clocks).
These steps are disabled in the example script, as they require at least 16~GB of system memory.
In this case, it is not necessary to separately write the 30-second clock files as listed above.

\begin{itemize}
      \item \configClass{selectEpochs}{gnssProcessingStepType:selectEpochs}: with \config{nthEpoch}=\verb|1| to set full 30-second sampling
      \item \configClass{selectNormalsBlockStructure}{gnssProcessingStepType:selectNormalsBlockStructure}: As the system of normal
      equations can be very large, the memory consumption might be reduced with \config{keepEpochNormalsinMemory}=\verb|no|.
      In this case the epoch parameters are directly eliminated during the accumulation and reconstructed in the solving step.
      This might lead to longer computation times.
      \item \configClass{estimate}{gnssProcessingStepType:estimate}: with \config{maxIterationCount}=\verb|2|: final iterations with full sampling
\end{itemize}

\subsection{Advanced: Processing large station networks}\label{cookbook.gnssNetwork:advanced}

Processing large station networks requires some additional steps to keep the computational load to a reasonable degree.
The general processing strategy is to first process a well-distributed subset of stations (i.e. a core network)
to get good estimates of all satellite parameters, which then enables integer ambiguity resolution (IAR). Once
the ambiguities of the core network are resolved and stable estimates for satellite phase biases are available, all
other (non-core) stations can be processed individually (including IAR) while keeping the satellite parameters fixed.
At last, all stations can be processed together with all satellite parameters and ionosphere parameters.

Let's start with the \configClass{processingSteps}{gnssProcessingStepType} of the core network:
\begin{itemize}
      \item \configClass{selectReceivers}{gnssProcessingStepType:selectReceivers} with \configClass{selectReceivers:file}{gnssTransceiverSelectorType:file}
      using the core network station list file from \reference{data preparation}{cookbook.gnssPpp:metadata} as \configFile{inputfileStringTable}{stringTable}.
      \item \configClass{selectEpochs}{gnssProcessingStepType:selectEpochs}: with \config{nthEpoch}=\verb|10| to reduce sampling to 5 minutes.
      \item \configClass{selectParametrizations}{gnssProcessingStepType:selectParametrizations}:
            disable \verb|constraint.STEC|, \verb|*VTEC|, \verb|*.tecBiases| as the ionosphere parameters are estimated in the final steps only.
      \item \configClass{estimate}{gnssProcessingStepType:estimate}: with \config{maxIterationCount}=\verb|6|
      \item \configClass{resolveAmbiguities}{gnssProcessingStepType:resolveAmbiguities}
      \item \configClass{estimate}{gnssProcessingStepType:estimate}: with \config{maxIterationCount}=\verb|4|: final iterations (with 5-minute sampling)
\end{itemize}

Now all other (non-core) stations can be processed separately:
\begin{itemize}
      \item \configClass{forEachReceiverSeparately}{gnssProcessingStepType:forEachReceiverSeparately}:
            with \configClass{selectReceivers:file}{gnssTransceiverSelectorType:file} inside \configClass{selectReceivers:exclude}{gnssTransceiverSelectorType:exclude}
            using the station list from the core network above to process all non-core stations individually with fixed transmitter parameters
      \begin{itemize}
      \item \configClass{processingStep:estimate}{gnssProcessingStepType:estimate}: with \config{maxIterationCount}=\verb|6|
      \item \configClass{processingStep:resolveAmbiguities}{gnssProcessingStepType:resolveAmbiguities}
      \item \configClass{processingStep:estimate}{gnssProcessingStepType:estimate}: with \config{maxIterationCount}=\verb|4|
      \end{itemize}
\end{itemize}

Next all stations are processed together with all parameters:
\begin{itemize}
      \item \configClass{selectReceivers}{gnssProcessingStepType:selectReceivers}: with \configClass{selectReceivers:all}{gnssTransceiverSelectorType:all}
      \item \configClass{selectEpochs}{gnssProcessingStepType:selectEpochs}: with \config{nthEpoch}=\verb|1| to set full 30-second sampling
      \item \configClass{group}{gnssProcessingStepType:group}: clock densification to 30-second sampling
      \begin{itemize}
            \item \configClass{selectParametrizations}{gnssProcessingStepType:selectParametrizations}:
            disable \verb|*| (all) parameters and reenable \verb|*.clock*| and \verb|*.STEC| parameters
            \item \configClass{estimate}{gnssProcessingStepType:estimate}: with \config{maxIterationCount}=\verb|6|
      \end{itemize}
      \item \configClass{selectParametrizations}{gnssProcessingStepType:selectParametrizations}: enable \verb|*| (all) parameters.
      \item \configClass{selectNormalsBlockStructure}{gnssProcessingStepType:selectNormalsBlockStructure}: with \config{keepEpochNormalsinMemory}=\verb|no|
      \item \configClass{estimate}{gnssProcessingStepType:estimate}: with \config{maxIterationCount}=\verb|4|: final iterations with full sampling and all parameters
      \item \configClass{writeResults}{gnssProcessingStepType:writeResults}: write the final results
\end{itemize}

% =============================================================

\clearpage
\section{GNSS precise point positioning (PPP)}\label{cookbook.gnssPpp}

This cookbook chapter describes an example of precise point positioning (PPP) of a ground station using GPS and Galileo.
For information on how to generate the GNSS products (orbits, clocks, signal biases, etc.) required for PPP see the
cookbook chapter on \reference{GNSS satellite orbit determination and station network analysis}{cookbook.gnssNetwork}.

\subsection{Data and metadata preparation}\label{cookbook.gnssPpp:metadata}

PPP requires some data and metadata, most of which are
provided in GROOPS file formats at \url{ftp://ftp.tugraz.at/outgoing/ITSG/groops/data}
(marked with \textbf{*} below). These files are regularly updated. Alternatively, check the linked
file types and programs for information on how to create these files manually.

Satellite metadata:
\begin{itemize}
  \item \file{Transmitter info}{gnssStationInfo}\textbf{*}: PRN-SVN mapping, antenna offsets and orientations
  \item \file{Antenna definition}{gnssAntennaDefinition}\textbf{*}: antenna center variations
  \item \file{Receiver definition}{gnssReceiverDefinition}\textbf{*}: transmitted signal types (recommended)
  \item \file{Satellite macro model}{satelliteModel}\textbf{*}: box-wing model, mass, antenna thrust
\end{itemize}

Station metadata:
\begin{itemize}
  \item \file{Station info}{gnssStationInfo}\textbf{*}: antenna and receiver information, antenna reference point
        offsets, antenna orientations
  \item \file{Antenna definition}{gnssAntennaDefinition}\textbf{*}: antenna center offsets and variations
  \item \file{Receiver definition}{gnssReceiverDefinition}: observed signal types (optional)
  \item \file{Accuracy definition}{gnssAntennaDefinition}\textbf{*}: elevation and azimuth dependent accuracy
\end{itemize}

Most data files are not provided in GROOPS formats and have to be converted manually from external file formats.

Satellite data:
\begin{itemize}
  \item \textbf{Precise orbits}: precise orbits in CRF for orbit integration (see \program{Sp3Format2Orbit})
  \item \textbf{Precise clocks}: precise clocks (see \program{GnssClockRinex2InstrumentClock})
  \item \textbf{Attitude}: rotation from body frame to CRF (see \program{SimulateStarCameraGnss} or \program{GnssOrbex2StarCamera})
  \item \textbf{Signal biases}: code (and phase) biases (see \program{GnssSinexBias2SignalBias})
\end{itemize}

Station data:
\begin{itemize}
  \item \textbf{Observations}: receiver measurements (see \program{RinexObservation2GnssReceiver})
  \item \textbf{Troposphere model}\textbf{*}: mapping functions, i.e. VMF3 (see \program{ViennaMappingFunctionGrid2File},
        \program{ViennaMappingFunctionStation2File})
\end{itemize}

Additional data:
\begin{itemize}
  \item \textbf{Earth orientation parameters}\textbf{*}: IERS Earth rotation data
        (see \configClass{EarthRotation}{earthRotationType})
  \item \textbf{Gravity field}\textbf{*}: gravity field model for orbit preparation
        (see \configClass{Gravityfield}{gravityfieldType})
  \item \textbf{Tides and tidal deformation}\textbf{*}: models for astronomical, Earth, ocean, pole, ocean pole,
        and atmospheric tides (see \configClass{Tides}{tidesType})
\end{itemize}

\subsection{Preprocessing: Orbit preparation}\label{cookbook.gnssPpp:orbitPreparation}

Simple numerical interpolation of externally sourced precise orbits is not recommended, especially for low but commonly
used sampling rates of 15~minutes. Since \program{GnssProcessing} requires \file{variational equations}{variationalEquation}
as input anyway, a simple numerical integration and orbit fit is the recommended procedure to prepare
precise orbits.

The steps, which can be done for each satellite individually using \program{LoopPrograms}, are:
\begin{itemize}
  \item Create a \reference{global element}{general.configFiles:globalSection} \config{orbitSampling}=\verb|15*60|
        (or whatever the sampling of the orbits) of type \verb|double|.
  \item Resample the orbit to at least 1-minute sampling over the full 24~hours using \program{InstrumentResample}
        (polynomial interpolation with \config{polynomialDegree}=\verb|11|, \config{maxDataPointRange}=\verb|12*{orbitSampling}|, and
        \config{maxExtrapolationDistance}=\verb|{orbitSampling}| to extrapolate to the midnight epoch in case it is not included).
  \item In case attitude was not provided, use \program{SimulateStarCameraGnss} to simulate the attitude for the resampled orbit.
  \item Use \program{GnssPrn2SvnBlockVariables} to create a \reference{variable}{general.parser} containing the space vehicle number (SVN, e.g. G056).
  \item Integrate the orbit using \program{PreprocessingVariationalEquation}.
  \begin{itemize}
    \item Provide the \configFile{inputfileSatelliteModel}{satelliteModel} for the SVN, the resampled \configFile{inputfileOrbit}{instrument},
          and the simulated or provided \configFile{inputfileStarCamera}{instrument}.
    \item Add \configClass{gravityfield:potentialCoefficients}{gravityfieldType:potentialCoefficients} (e.g. GOCO06s static up to degree 20),
          \configClass{tides:astronomicalTide}{tidesType:astronomicalTide}, \configClass{tides:earthTide}{tidesType:earthTide}, and
          \configClass{miscAccelerations:relativisticEffect}{miscAccelerationsType:relativisticEffect} as \config{forces}.
    \item For \configClass{gradientField:potentialCoefficients}{gravityfieldType:potentialCoefficients} provide the same gravity field as above.
    \item Also provide \configClass{earthRotation}{earthRotationType} and \configClass{ephemerides}{ephemeridesType}.
  \end{itemize}
  \item Use \program{PreprocessingVariationalEquationOrbitFit} to fit the integrated orbit to the resampled reference orbit by estimating
        \configClass{parametrizationAcceleration:gnssSolarRadiation}{parametrizationAccelerationType:gnssSolarRadiation}
        (selecting at least the D0, D2, Y0, B1, B1 parameters). Set \config{iterationCount}=\verb|1|.
  \item Use \program{PreprocessingVariationalEquationOrbitFit} again, but this time instead of solar radiation parameters set up
        \configClass{stochasticPulse:uniformSampling}{timeSeriesType:uniformSampling} with
        \config{timeStart}=\verb|<mjd>+{orbitSampling}/86400|, \config{timeEnd}=\verb|<mjd>+1-{orbitSampling}/86400|,
        and \config{sampling}=\verb|{orbitSampling}/86400| (replace \verb|<mjd>| with MJD of the day to process).
\end{itemize}

The resulting \file{variational equations}{variationalEquation} files contain the integrated and fitted orbit, attitude, and Earth rotation
that are reqired by \program{GnssProcessing} for PPP.

\subsection{Processing of a ground station}\label{cookbook.gnssPpp:processing}

Once all metadata, data, and the orbits are prepared, the actual processing can begin.
Use \program{GnssProcessing} and define the \configClass{transmitter}{gnssParametrizationTransmitterType}
constellations, a \configClass{receiver:stationNetwork}{gnssParametrizationReceiverType:stationNetwork}, and
the list of \configClass{processingSteps}{gnssProcessingStepType}.

Since this program is very complex, an overview of the most important settings used for the example
processing in this cookbook chapter follows. Only config elements that deviate from the default values are listed.
For most input files the default value and annotation give a good hint of what is expected, so they are not
listed here.

The first step is setting the processing sampling, in this example it is 30~seconds.
The processing interval usually is a single 24-hour day.
\begin{itemize}
  \item \configClass{intervals:uniformSampling}{timeSeriesType:uniformSampling} with \config{timeStart}=\verb|<mjd>|,
        \config{timeEnd}=\verb|<mjd>+1|, \config{sampling}=\verb|1| (replace \verb|<mjd>| with MJD of the day to process)
  \item \configClass{timeSeries:uniformSampling}{timeSeriesType:uniformSampling} with \config{timeStart}=\verb|<mjd>|,
        \config{timeEnd}=\verb|<mjd>+1|, \config{sampling}=\verb|30/86400| (processing sampling)
\end{itemize}

Next define the GNSS constellations, in this case GPS and Galileo.
\begin{itemize}
  \item \configClass{transmitter:GPS}{gnssParametrizationTransmitterType:gps}
  \begin{itemize}
    \item \configFile{inputfileSignalBias}{gnssSignalBias}: code (and phase) biases
    \item \configFile{inputfileVariational}{variationalEquation}=\verb|preprocessing/variational.{prn}.dat| (prepared orbits)
    \item \config{estimateCodeBias}=\verb|no| (code biases should be provided)
    \item \config{estimatePhaseBias}=\verb|no| (code biases must be provided in case integer ambiguities are to be resolved)
  \end{itemize}
  \item \configClass{transmitter:Galileo}{gnssParametrizationTransmitterType:galileo} (use same settings as for GPS)
\end{itemize}

Define the station(s) to be processed. In this example it is a single stations, but you can also process multiple
stations at the same time.
Prepare the respective \file{station list files}{stringList} (one station name per line) beforehand.
\begin{itemize}
  \item \configClass{receiver:stationNetwork}{gnssParametrizationReceiverType:stationNetwork}
  \begin{itemize}
    \item \configFile{inputfileStationList}{stringList}: list of all stations to be processed
    \item \configFile{inputfileSignalBias}{gnssSignalBias}: leave empty, since we estimate our own biases
    \item \config{kinematicPositionEstimation}=\verb|yes| (since we want to estimate a position each epoch)
    \item \config{estimateCodeBias}=\verb|yes|
    \item \config{estimatePhaseBias}=\verb|yes|
    \item \configClass{troposphere:viennaMapping}{troposphereType} (e.g. VMF3 operational gridded)
    \item \configClass{troposphereWetEstimation:splines}{parametrizationTemporalType:splines} (e.g. degree 1 with 2-hourly
          nodes and daily intervals)
    \item \configClass{troposphereGradientEstimation:constant}{parametrizationTemporalType:constant}
    \item \configClass{troposphereGradientEstimation:trend}{parametrizationTemporalType:trend} (at center of day with
          \config{timeStep}=\verb|1|)
    \item \configClass{tides:earthTide}{tidesType:earthTide}: Earth tidal deformations (IERS conventions)
    \item \configClass{tides:doodsonHarmonicTide}{tidesType:doodsonHarmonicTide}: ocean tidal deformations
          (e.g. FES 2014b, \config{minDegree}=\verb|1|)
    \item \configClass{tides:doodsonHarmonicTide}{tidesType:doodsonHarmonicTide}: atmospheric tidal deformation
          (e.g. AOD1B RL06, \config{minDegree}=\verb|1|)
    \item \configClass{tides:poleTide}{tidesType:poleTide}: pole tidal deformations (IERS conventions)
    \item \configClass{tides:poleOceanTide}{tidesType:oceanPoleTide}: ocean pole tidal deformations (IERS conventions)
    \item \config{preprocessing}: settings for station-wise observation preprocessing
    \begin{itemize}
      \item We recommend to explicitly specify the signals to be processed and to make sure that at least transmitter
            code biases are provided for each of them, e.g. \configClass{useType}{gnssType}=\verb|C1CG|,
            \verb|C1CW|, \verb|C2WG|, \verb|L1*G|, \verb|L2*G|, ...).
      \item Set \config{codeMaxPositionDiff}=\verb|200| or larger in case your approximate position is not accurate or
            too many observations/epochs get disabled.
      \item Set \config{minEstimableEpochsRatio}=\verb|0.25| or lower in case too many observations/epochs get disabled
            and therefore the station is not processed at all.
    \end{itemize}
  \end{itemize}
\end{itemize}

In the next step, define the Earth rotation. The selected Earth rotation includes a high-frequency EOP model.
\begin{itemize}
  \item \configClass{earthRotation}{gnssParametrizationEarthRotationType}
  \begin{itemize}
    \item \configClass{earthRotation:file}{earthRotationType:file}
          with \configFile{inputfileEOP}{earthOrientationParameter}=\verb|timeSeries_EOP_IAU2000_desai.dat|
  \end{itemize}
\end{itemize}

Some parametrizations don't offer config elements for constraints on the respective parameters, so they have to be added
as general constraints. The parameters to be constrained can be selected by name (including support for wildcards).
The sigma for general constraints is always of the same unit as the parameter.
In this example only loose constraints on tropospheric parameters are required.
\begin{itemize}
  \item \configClass{constraints}{gnssParametrizationConstraintsType}: loose constraint troposphere estimates
  \begin{itemize}
    \item \configClass{parameters:wildcard}{parameterSelectorType:wildcard} with \config{object}=\verb|*|,
          \config{type}=\verb|troposphere*|, \config{temporal}=\verb|*|, \config{interval}=\verb|*|
    \item \config{sigma}=\verb|5| meters (wet)/degrees (gradient)
    \item \config{relativeToApriori}=\verb|yes|
  \end{itemize}
\end{itemize}

Finally, define the processing steps. This can be overwhelming at first, but offers a lot of flexibility.
The general steps are that an iterative least squares adjustment shall be performed for each station
(in this example there is only one) individually. Optionally, the integer ambiguities can be resolved as well
(only possible if transmitter phase biases are provided).
With all processing done, print the residual statistics and write the resulting output files.
\begin{itemize}
  \item \configClass{processingStep:forEachReceiverSeparately}{gnssProcessingStepType:forEachReceiverSeparately}:
        process all stations individually with fixed transmitter parameters
  \begin{itemize}
    \item \configClass{processingStep:selectParameters}{gnssProcessingStepType:selectParameters}: parameters
          for least squares adjustment (TEC biases cannot be estimated together with signal biases)
    \begin{itemize}
      \item \config{estimateReceiverTecBias}=\verb|no|
      \item \config{estimateTransmitterTecBias}=\verb|no|
      \item \config{constraintIonosphereSTEC}=\verb|no|
      \item \config{estimateIonosphereVTec}=\verb|no|
      \item \config{estimateIonosphereTecMap}=\verb|no|
    \end{itemize}
    \item \configClass{processingStep:estimate}{gnssProcessingStepType:estimate}: initial iterative estimation
          of all parameters (except TEC biases)
    \begin{itemize}
      \item \config{computeResiduals}=\verb|yes|
      \item \config{adjustSigma0}=\verb|groupedPhases|
      \item \config{computeWeights}=\verb|groupedPhases| (downweight code observations individually,
            but phase observations all grouped together)
      \item \config{convergenceThreshold}=\verb|0.005| meters
      \item \config{maxIterationCount}=\verb|6|
    \end{itemize}
    \item \configClass{processingStep:group}{gnssProcessingStepType:group}: estimate TEC biases in an intermediate step
    \begin{itemize}
      \item \configClass{processingStep:selectParameters}{gnssProcessingStepType:selectParameters}: set everything
            to \verb|no| except for the elements below (settings apply only to steps in current
            \configClass{group}{gnssProcessingStepType:group})
      \begin{itemize}
        \item \config{estimateReceiverTecBias}=\verb|yes|
        \item \config{constraintTransmitterSignalBiasZeroMean}=\verb|yes|
        \item \config{estimateTransmitterTecBias}=\verb|yes|
        \item \config{estimateIonosphereVTec}=\verb|yes|
        \item \config{estimateIonosphereTecMap}=\verb|yes|
        \item \config{constraintOther}=\verb|yes|
      \end{itemize}
      \item \configClass{processingStep:estimate}{gnssProcessingStepType:estimate}
      \begin{itemize}
        \item \config{computeResiduals}=\verb|yes|
        \item \config{convergenceThreshold}=\verb|0.005| meters
        \item \config{maxIterationCount}=\verb|3|
      \end{itemize}
    \end{itemize}
    \item \configClass{processingStep:estimate}{gnssProcessingStepType:estimate}: some more least squares iterations
          (same settings as first time except for \config{maxIterationCount}=\verb|2|)
    \item \configClass{processingStep:resolveAmbiguities}{gnssProcessingStepType:resolveAmbiguities}: (optional) integer
          ambiguity resolution (only possible if transmitter phase biases have been provided)
    \begin{itemize}
      \item \config{computeResiduals}=\verb|yes|
      \item \config{adjustSigma0}=\verb|groupedPhases|
      \item \config{computeWeights}=\verb|groupedPhases|
    \end{itemize}
    \item \configClass{processingStep:estimate}{gnssProcessingStepType:estimate}: (in case of ambiguity resolution) some more least squares iterations
          (same settings as first time except for \config{maxIterationCount}=\verb|2|)
  \end{itemize}
  \item \configClass{processingStep:printResidualStatistics}{gnssProcessingStepType:printResidualStatistics}
  \item \configClass{processingStep:writeResults}{gnssProcessingStepType:writeResults}
\end{itemize}

% =============================================================

\clearpage
\section{Kinematic orbit determination of LEO satellites}\label{cookbook.kinematicOrbit}

This cookbook chapter describes exemplarily the steps for determining kinematic orbits of low-Earth orbit (LEO) satellites.

\subsection{Prepare LEO meta data}
\begin{enumerate}
  \item For creating the \file{GnssStationInfo file}{gnssStationInfo} use \program{GnssStationInfoCreate}. The rotation from the
  satellite reference frame into the antenna reference frame as well as the change of the center of mass due to fuel consumption has to be considered here.
  \item The \file{GnssReceiverDefinition file}{gnssReceiverDefinition} can be created by using \program{GnssReceiverDefinitionCreate}. In this program the different GNSS types has to be specified.
  \item For creating the \file{GnssAntennaDefinition file}{gnssAntennaDefinition}\ use \program{GnssAntennaDefinitionCreate}. This program considers the phase center offsets of L1 and L2 frequency from the used antenna.
  \item For determining the elevation dependent accuracies the program \program{GnssAntennaDefinitionCreate} is used again.
  \begin{itemize}
    \item \config{antenna}: set to new
    \item Set the \config{pattern}s for code (\config{type}=\verb|C**|) and phase
    (\config{type}=\verb|L**|). The standard deviation is expressed e.g. with \config{values}=\verb|0.001/cos(2*PI/180*zenith)|.
  \end{itemize}
\end{enumerate}

\subsection{Read LEO data}
\begin{enumerate}
  \item Conversion of the approximated orbit and star camera data into GROOPS format using a conversion program.
  \item If no attitude data is given the star camera data can be simulated by using \program{SimulateStarCamera} or \program{SimulateStarCameraSentinel1}.
  \item The GNSS observation data (given in RINEX format) can be converted with \program{RinexObservation2GnssReceiver}.
  \item Suitable programs to get daily data are \program{InstrumentConcatenate} and \program{InstrumentSynchronize}.
  \item For interpolating the orbit and star camera data to GNSS receiver epochs use \program{InstrumentResample} and specify GNSS receiver as input for
  \config{timeSeries}.
  \item For synchronizing these data use \program{InstrumentSynchronize}
\end{enumerate}
Detailed description of instrument data handling can be found in \reference{Instrument data handling}{cookbook.instrument}.

\subsection{Prepare GNSS satellite metadata, orbits, and clocks}
The cookbook chapter on \reference{GNSS precise point positioning (PPP)}{cookbook.gnssPpp} details how to prepare the GNSS satellite
\reference{metadata}{cookbook.gnssPpp:metadata} and \reference{orbits}{cookbook.gnssPpp:orbitPreparation}.
Use \program{GnssClockRinex2InstrumentClock} to convert satellite clock errors.

\subsection{Estimate kinematic orbits}

For kinematic orbit determination use \program{GnssProcessing}. The GNSS constellations have to be defined via
\configClass{transmitter}{gnssParametrizationTransmitterType} which requires for example the signal bias and clock data.
The LEO satellite can be defined via \configClass{receiver:lowEarthOrbiter}{gnssParametrizationReceiverType:lowEarthOrbiter} which requires for example the
GNSS observations, the star camera and approximated orbit data. The impact of the ionosphere can be considered via
\configClass{ionosphere}{gnssParametrizationIonosphereType} and the estimation of the earth rotation parameters is done
with \configClass{earthRotation}{gnssParametrizationEarthRotationType}. The ambiguity resolution settings can be specified via
\configClass{ambiguities}{gnssParametrizationAmbiguitiesType}.

The processing flow is controlled via a list of processing steps \configClass{processingStep}{gnssProcessingStepType}. The processing steps are
processed consecutively. Some steps allow the selection of parameters,
epochs, or the normal equation structure, which affects all subsequent steps.

\begin{enumerate}
    \item \configClass{processingStep:selectParameters}{gnssProcessingStepType}:  Parameters which should be estimated have to be enabled.
    Consider the linear dependencies between the STEC and the VTEC as well as between the STEC and particular signal biases. This means, if STEC parameter
    is enabled it is not possible to estimate the VTEC and TEC bias, so these parameters have to be be disabled.
    \item \configClass{processingStep:estimate}{gnssProcessingStepType}: Estimation of the parameters. The number of iterations can be set with
    \config{maxIterationCount}. The weights will be iteratively adjusted.
    \item Estimation of the TEC bias parameter using \configClass{processingStep:selectParameters}{gnssProcessingStepType} and
    \configClass{processingStep:estimate}{gnssProcessingStepType}. Only \config{estimateReceiverTecBias} parameter has to be selected.
    \item \configClass{processingStep:estimate}{gnssProcessingStepType}: Further estimation of the parameters.
    \item Solving the ambiguities via \configClass{processingStep:resolveAmbiguities}{gnssProcessingStepType}.
    \item Further estimations of the parameters and the TEC bias parameter by using \configClass{processingStep:estimate}{gnssProcessingStepType}.
    \item Computation of the epoch-wise covariance matrix can be done via \configClass{processingStep:computeCovarianceMatrix}{gnssProcessingStepType}.
    \item \configClass{processingStep:writeResults}{gnssProcessingStepType}: Write estimation results into output files.
\end{enumerate}




% =============================================================

\clearpage
\section{Gravity field determination from precise orbit data (POD)}\label{cookbook.gravityFieldPod}

This cookbook chapter describes exemplarily the steps for determining the monthly gravity variations from precise orbit data (POD).

\subsection{Step 1: Preperation of data}

Following data have to be prepared monthly with an adequate sampling, e.g. 10 s using
\program{InstrumentConcatenate}:
\begin{itemize}
  \item Precise (kinematic) orbit data
  \item 3x3 covariance matrices data
  \item Initial orbit data used for precise orbit determination
  \item Star camera data
  \item Accelerometer data
\end{itemize}
Reduced sampling can be achieved by \program{InstrumentReduceSampling}. If the satellite mission does not provide any required
accelerometer data, these data can be generated via \program{SimulateAccelerometer}.

For satellite missions with less knowledge about the acting forces, it make sense to consider more than one state vector within an orbit revolution.
Otherwise the accuracy of the estimated parameters will decrease. This implies that shorter arcs are necessary. The assignment of the kinematic orbit
data as well as the 3x3 covariance matrices data to the arcs can be done with \program{InstrumentSynchronize}.

\subsection{Step 2: Conversion of the background gravity field}
\program{Gravityfield2SphericalHarmonicsVector} converts the static respectively background gravity field into spherical harmonics.

\subsection{Step 3: Preprocessing POD}
For determining the accuracies and weights of the kinematic orbits it is sufficient to make a least-square estimation with only certain parameters, due
to the fact that some parameters do not influence the estimation of the accuracies and weights.
This estimation is done with \program{PreprocessingPod}. Additional this program determines the temporal correlation of the kinematic orbit positions
x,y and z. If short arcs are used the setting \configClass{observation:podIntegral}{observationType:podIntegral} shall be used. This setting
considers the frictional forces by means of a macro model as well as the conservative and non-conservative forces.

\subsection{Step 4: Solving of normal equations system}
\program{NormalsSolverVCE} sets up the observation equations and summarized them to a normal equations system. The subsequent least-square estimation delivers
the parameters surcharges.

\subsection{Step 5: Determination of the estimated gravity field parameters}
The estimated parameters result from the re-addition of the background field, which is done in \program{MatrixCalculate}.

\subsection{Step 6: Conversion of the gravity field parameters}
\program{Gravityfield2PotentialCoefficients} converts the gravity field parameters into spherical harmonics.

% =============================================================

\clearpage
\section{GRACE gravity field recovery}\label{cookbook.gravityFieldGrace}
This cookbook chapter describes an example of estimating a gravity field solution using GRACE observation
data. For the respective month a set of spherical harmonic coefficients up to a maximum degree is determined.
An example scenario for this task can be found at
\url{https://ftp.tugraz.at/outgoing/ITSG/groops/scenario/scenarioGraceGravityfieldRecovery.zip}
including the required GROOPS scripts and data sets for the gravity field recovery process.
The background models are provided at \url{https://ftp.tugraz.at/outgoing/ITSG/groops/data/}.


\subsection{Background models}\label{cookbook.gravityFieldGrace:backgroundModels}
The following background models were used during the data processing:
\begin{itemize}
  \item \textbf{Earth rotation}: \href{https://www.iers.org/IERS/EN/Publications/TechnicalNotes/tn36.html}{IERS 2010} % IERS 2010
  \item \textbf{Moon, sun and planets ephemerides}: \href{https://ipnpr.jpl.nasa.gov/progress_report/42-196/196C.pdf}{JPL DE432} %JPL DE432
  \item \textbf{Earth tide}: \href{https://www.iers.org/IERS/EN/Publications/TechnicalNotes/tn36.html}{IERS 2010} % IERS 2010
  \item \textbf{Ocean tide}: \href{https://www.aviso.altimetry.fr/es/data/products/auxiliary-products/global-tide-fes/description-fes2014.html}{FES2014b} % FES2014b
  \item \textbf{Pole tide}: \href{https://www.iers.org/IERS/EN/Publications/TechnicalNotes/tn36.html}{IERS 2010} % IERS 2010
  \item \textbf{Ocean pole tide}: \href{https://doi.org/10.1029/2001JC001224}{Desai 2004} % Desai 2004 (IERS 2010)
  \item \textbf{Atmospheric tides}: \href{https://doi.org/10.1093/gji/ggx302}{AOD1B RL06} % AOD1B RL06
  \item \textbf{Atmosphere and Ocean Dealiasing}: \href{https://doi.org/10.1093/gji/ggx302}{AOD1B RL06} % AOD1B RL06
  \item \textbf{Sub-monthly continental hydrology}: \href{https://doi.org/10.2312/GFZ.b103-08095}{LSDM (ESMGFZ)} % LSDM (ESMGFZ)
  \item \textbf{Relativistic corrections}: \href{https://www.iers.org/IERS/EN/Publications/TechnicalNotes/tn36.html}{IERS 2010} % IERS 2010
\end{itemize}
These models were reduced during the analysis process and are not present in the solution.
The \href{https://doi.org/10.5194/essd-13-99-2021}{GOCO06s}
model was used as the static gravity field as well as for the trend component and annual oscillation.
In the script \verb|000groopsBackgroundModels.xml| a monthly mean of the GOCO06s including the time-variable components is determined
in form of time splines using \program{Gravityfield2TimeSplines}. This model is later added back to the final gravity solution.

\subsection{Instrument data preparation}\label{cookbook.gravityFieldGrace:metadata}
The ITSG gravity field solutions are computed from the official GRACE L1B \href{https://doi.org/10.5067/GRJPL-L1B03}{JPL (2018)}
and GRACE-FO L1B \href{https://doi.org/10.5067/GFL1B-ASJ04}{JPL (2019)} observation data. The data sets for this example are provided
in GROOPS file format in the scenario folder.

The satellite-to-satellite-tracking (SST) data consists of:
\begin{itemize}
  \item \textbf{K-band range rates}
  \item \textbf{Light time correction}
  \item \textbf{Antenna offset corrections}
\end{itemize}

Additional observation data required for the processing comprises:
\begin{itemize}
  \item \textbf{Star camera observations}
  \item \textbf{Accelerometer data}
  \item \textbf{Approximate orbits} % \textbf{Initial dynamic orbits}
  \item \textbf{Thruster data}
\end{itemize}

The determination of
\begin{itemize}
  \item \textbf{Kinematic orbits}
  \item \textbf{3x3 epoch covariances}
\end{itemize}
is depicted in \reference{Kinematic orbit determination of LEO satellites}{cookbook.kinematicOrbit}.
These data sets are also provided in the scenario folder.

Data preparation is handled in the script \verb|010groopsInstruments.xml|. The approximate orbits (initial dynamic orbits)
of the satellites, the star camera observations, the accelerometer data and the thruster data are resampled with a 5s sampling
and small gaps in the data are filled using \program{InstrumentResample}. Gross outliers are removed using \program{InstrumentRemoveEpochsByCriteria}
and the data is synchronized using  \program{InstrumentSynchronize}.

The approximate orbits are later used as a priori information for the dynamic orbit integration.
In addition to the observed orientation of the spacecrafts (star camera observations),
the nominal orientation is computed using \program{SimulateStarCameraGrace}. The difference between
observed and simulated orientation is determined using \program{InstrumentStarCameraMultiply}
and is employed in the outlier detection.

The accelerometer data is initially calibrated by estimating a bias using
\program{InstrumentAccelerometerEstimateBiasScale} with respect to simulated data created
with \program{SimulateAccelerometer}. For simulating accelerometer data a satellite model
implying the satellite's mass and surfaces is required. Such a model can be created with
\program{SatelliteModelCreate}. Models for the GRACE and GRACE-FO satellites are also provided
at \url{https://ftp.tugraz.at/outgoing/ITSG/groops/data/satelliteModel/}. Non-gravitational forces
comprising atmospheric drag, solar radiation pressure and albedo have to modeled when simulating
the accelerometer data. The acceleration bias parameters are determined as degree 3 time splines
with 6h nodes. When determining these parameters the thruster events are excluded from the estimation.

The SST observations, the light time corrections and the antenna center corrections are synchronized
with a 5s sampling together with simulated SST data created with \program{SimulateSatelliteTracking}.
Simulated data is used for the outlier detection of the original SST observations.

The sampling of the kinematic orbits is reduced to 60s using \program{InstrumentReduceSampling} and
an outlier detection is performed using the approximate dynamic orbits.

The approximate orbits, the star camera observations and the accelerometer data are divided into 24h arcs
(variational arcs). The kinematic orbits, its 3x3 epoch covariances, KBR observations,
light time corrections, antenna center corrections and star camera observations are divided into 3h arcs
per day (short arcs). Additionally the approximate orbits and the star camera observations are also
synchronized to short arcs.

Further information on instrument data preparation can be found in
\reference{Instrument data handling}{cookbook.instrument}.

\subsection{Variational equations}\label{cookbook.gravityFieldGrace:variational}
In this processing step dynamic orbits are computed for a complete 24h orbit arc by integrating
the forces acting on the GRACE/GRACE-FO satellites. Additionally, the state transition matrix is set up.
The dynamic orbits are then fitted to kinematic orbits and SST observations in a least squares adjustment
by co-estimating additional accelerometer calibration parameters together with the initial state vector.
The newly estimated parameters are then used to re-estimate the dynamic orbits and setting up the new
state transition matrix.

The script \verb|020groopsVariational.xml| in the scenario folder implements the required processing steps.
Time splines from a time-variable gravity field are estimated using \program{Gravityfield2TimeSplines}.
In this step the static gravity field (GOCO06s) is combined with the following time-variable components:
\begin{itemize}
    \item \configClass{gravityfield:potentialCoefficients}{gravityfieldType:potentialCoefficients}:
    static gravity field
    \item \configClass{gravityfield:trend}{gravityfieldType:trend}
    \begin{itemize}
        \item \configClass{gravityfield:potentialCoefficients}{gravityfieldType:potentialCoefficients}:
      trend component of gravity field
    \end{itemize}
    \item \configClass{gravityfield:oscillation}{gravityfieldType:oscillation}
    \begin{itemize}
        \item \configClass{gravityfieldCos:potentialCoefficients}{gravityfieldType:potentialCoefficients}:
      annual cosine component of gravity field
        \item \configClass{gravityfieldSin:potentialCoefficients}{gravityfieldType:potentialCoefficients}:
      annual sine component of gravity field
    \end{itemize}
    \item \configClass{gravityfield:timeSplines}{gravityfieldType:timeSplines}: atmosphere and ocean dealiasing (AOD1B RL06)
    \item \configClass{gravityfield:timeSplines}{gravityfieldType:timeSplines}: hydrology dealiasing (LSDM)
    \item \configClass{doodsonHarmonicTide}{tidesType:doodsonHarmonicTide}: ocean tides (FES2014b)
    \item \configClass{doodsonHarmonicTide}{tidesType:doodsonHarmonicTide}: atmospheric tides (AOD1B RL06)
    \item \configClass{tides:poleTide}{tidesType:poleTide}: pole tides (IERS 2010)
    \item \configClass{tides:poleOceanTide}{tidesType:oceanPoleTide}: ocean pole tides (IERS 2010)
\end{itemize}
\config{maxDegree}=\verb|220| and \config{sampling}=\verb|10/1440| is sufficient.

In \program{PreprocessingVariationalEquation} the \file{variational equations}{variationalEquation}
comprising the integrated orbit together with the state transition matrix are stored in
\configFile{outputfileVariational}{variationalEquation}.

This program has to be executed for both GRACE or GRACE-FO satellites and it is recommended
to use \program{LoopPrograms}.
\begin{itemize}
    \item \configFile{inputfileSatelliteModel}{satelliteModel}: satellite model from \verb|020groopsInstruments.xml|
    \item \configFile{inputfileOrbit}{instrument}: the approximate orbits from \verb|020groopsInstruments.xml|
    \item \configFile{inputfileStarCamera}{instrument}: the attitude file from \verb|020groopsInstruments.xml|
    \item \configFile{inputfileAccelerometer}{instrument}: the accelerometer data from
    \verb|020groopsInstruments.xml|
    \item \config{forces}: see below
    \item \configClass{ephemerides}{ephemeridesType}: JPL DE432
    \item \configClass{gradientfield:potentialCoefficients}{gravityfieldType:potentialCoefficients}: a static gravity field (GOCO06s) with \config{maxDegree}=\verb|10| is more than sufficient.
\end{itemize}

The \configClass{force models}{forcesType} include:
\begin{itemize}
  \item \configClass{gravityfield:timeSplines}{gravityfieldType:timeSplines}:
        the previously estimated time-variable gravity field
  \item \configClass{tides:astronomicalTide}{tidesType:astronomicalTide}: astronomical tides (based on JPL DE432 ephemerides)
  \item \configClass{tides:earthTide}{tidesType:earthTide}: Earth tide (IERS conventions)
  \item \configClass{miscAccelerations:relativisticEffect}{miscAccelerationsType:relativisticEffect}:
        relativistic effects (IERS conventions)
\end{itemize}

In \program{PreprocessingVariationalEquationOrbitFit} the integrated orbit
(\configFile{inputfileVariational}{variationalEquation}) is fitted to the kinematic orbit
(\configFile{inputfileOrbit}{instrument}) by least squares adjustment. The additional accelerometer
calibration parameters can be defined by
\begin{itemize}
  \item \configClass{parametrizationAcceleration}{parametrizationAccelerationType}: accelerometer scale factor (once per day)
  \item \configClass{parametrizationAcceleration}{parametrizationAccelerationType}: accelerometer bias (time spline with 6h nodes)
\end{itemize}

The observation equations (parameter sensitivity matrix) are computed by integration of the variational
equations (\configFile{inputfileVariational}{variationalEquation}) using a polynomial with
\config{integrationDegree}=\verb|7|. \program{PreprocessingVariationalEquationOrbitFit} has to be
executed per satellite.

\program{PreprocessingVariationalEquationSstFit} fits two dynamic orbits
\configFile{inputfileVariational1/2}{variationalEquation} to the SST observations and the kinematic orbits.
\begin{itemize}
    \item \config{rightHandSide}: input for observation vectors
        \begin{itemize}
            \item \configFile{inputfileSatelliteTracking}{instrument}: K-band range rate observations
            \item \configFile{inputfileSatelliteTracking}{instrument}: light time correction
            \item \configFile{inputfileSatelliteTracking}{instrument}: antenna offset corrections
            \item \configFile{inputfileOrbit1}{instrument}: kinematic orbit of satellite 1
            \item \configFile{inputfileOrbit2}{instrument}: kinematic orbit of satellite 2
        \end{itemize}
    \item \config{sstType}: rangeRate
    \item \configFile{inputfileVariational1}{variationalEquation}: dynamic orbit and integrated state matrix of satellite 1
    \item \configFile{inputfileVariational2}{variationalEquation}: dynamic orbit and integrated state matrix of satellite 2
    \item \configClass{parametrizationAcceleration1}{parametrizationAccelerationType}: same as in In \program{PreprocessingVariationalEquationOrbitFit}
    \item \configClass{parametrizationAcceleration2}{parametrizationAccelerationType}: same as in In \program{PreprocessingVariationalEquationOrbitFit}
    \item \config{integrationDegree}: \verb|7|
    \item \config{interpolationDegree}: \verb|7|
    \item \configClass{covarianceSst}{covarianceSstType}
        \begin{itemize}
            \item \config{sigma}: \verb|1|
        \end{itemize}
    \item \configClass{covariancePod1}{covariancePodType}
        \begin{itemize}
            \item \config{sigma}: 1
            \item \configFile{inputfileCovariancePodEpoch}{instrument}: 3x3 epoch covariances
        \end{itemize}
    \item \configClass{covariancePod2}{covariancePodType}
        \begin{itemize}
            \item \config{sigma}: 1
            \item \configFile{inputfileCovariancePodEpoch}{instrument}: 3x3 epoch covariances
        \end{itemize}
\end{itemize}
The estimated accelerometer calibration parameters from \program{PreprocessingVariationalEquationOrbitFit}
and \program{PreprocessingVariationalEquationSstFit} are determined as corrections and stored
in \configFile{outputfileSolution}{matrix}. Both correction estimates have to be summed up using
\program{FunctionsCalculate}.

The dynamic orbit and the resulting accelerometer calibration parameters are now used to re-integrate
the orbit once more using \program{PreprocessingVariationalEquation} and introducing
\configClass{parametrizationAcceleration}{parametrizationAccelerationType} as \config{estimatedParameters}.
This step usually ensures convergence. If the maximum orbit difference is still not sufficient this step
can be repeated again.

\subsection{Preprocessing}\label{cookbook.gravityFieldGrace:preprocessing}
The script \verb|030groopsPreprocessing.xml| implements the following steps and settings.
The program \program{PreprocessingSst} processes SST observations and kinematic orbit data and performs
a complete least squares adjustment for gravity field determination by computing the observations equations.
Force model parameters (gravitational potential coefficients and accelerometer calibration parameters)
are computed by integrating the parameter sensitivity matrix from the variational equations.
Parameters describing effects due to the SST observation system and geometry (KBR antenna phase
center variations) are computed using the dynamic orbits as a Taylor point.
Short time gravity variations can be co-estimated together with the monthly mean gravity field.
The autoregressive model sequence constraining the short time parameters is provided in the data folder.
See \href{https://doi.org/10.1007/s00190-019-01314-1}{Kvas 2019} for more information about
this co-estimation.
\begin{itemize}
    \item \configClass{observation}{observationType}: sstVariational
        \begin{itemize}
            \item \config{rightHandSide}:
            \begin{itemize}
                \item \configFile{inputfileSatelliteTracking}{instrument}: KBR range rates
                \item \configFile{inputfileSatelliteTracking}{instrument}: light time correction
                \item \configFile{inputfileSatelliteTracking}{instrument}: antenna offset corrections
                \item \configFile{inputfileOrbit1}{instrument}: kinematic orbit of satellite 1
                \item \configFile{inputfileOrbit2}{instrument}: kinematic orbit of satellite 2
            \end{itemize}
            \item \config{sstType}: rangeRate
            \item \configFile{inputfileVariational1}{variationalEquation}: dynamic orbit and integrated state matrix of satellite 1
            \item \configFile{inputfileVariational2}{variationalEquation}: dynamic orbit and integrated state matrix of satellite 2
            \item \configClass{ephemerides}{ephemeridesType}: JPL DE432
            \item \configClass{parametrizationGravity}{parametrizationGravityType}: spherical harmonics from \config{minDegree}=\verb|2| to \config{maxDegree}=\verb|60|
            \item \configClass{parametrizationGravity}{parametrizationGravityType}: high frequency parametrization
            \item \configClass{parametrizationAcceleration1}{parametrizationAccelerationType}: same as in In \program{PreprocessingVariationalEquationOrbitFit}
            \item \configClass{parametrizationAcceleration2}{parametrizationAccelerationType}: same as in In \program{PreprocessingVariationalEquationOrbitFit}
            \item \configClass{parametrizationSst}{parametrizationSatelliteTrackingType}: antenna phase center variations (y and z for both satellites)
            \item \config{integrationDegree}: 7
            \item \config{interpolationDegree}: 7
        \end{itemize}
    \item \configClass{covarianceSst}{covarianceSstType}
        \begin{itemize}
            \item \config{sigma}: \verb|1e-7|
            \item \config{sampling}: 5 [seconds]
        \end{itemize}
    \item \configClass{covariancePod1/2}{covariancePodType}
        \begin{itemize}
            \item \config{sigma}:  \verb|2|
            \item \configFile{inputfileCovariancePodEpoch}{instrument}: 3x3 epoch covariances
            \item \config{sampling}: 60 [seconds]
        \end{itemize}
    \item \config{estimateShortTimeVariations}
        \begin{itemize}
            \item \configClass{autoregressiveModelSequence}{autoregressiveModelSequenceType}: AR models
            \item \configClass{parameterSelection}{parameterSelectorType}: names
                \begin{itemize}
                    \item \configClass{parameterName}{parameterNamesType}: parametrizationGravity
                    \begin{itemize}
                        \item \configClass{parametrization}{parametrizationGravityType}: high frequency parametrization
                    \end{itemize}
                \end{itemize}
        \end{itemize}
\end{itemize}

\program{ParameterSelection2IndexVector} and \program{MatrixCalculate} with
\configClass{matrix:reorder}{matrixGeneratorType:reorder} can be used to extract the desired
spherical harmonic coefficients from \configFile{outputfileSolution}{matrix} and the respective
standard deviations from \configFile{outputfileSigmax}{matrix} up to a certain degree.

In the program \program{Gravityfield2PotentialCoefficients} the estimated spherical harmonics
coefficients are read with
\configClass{gravityfield:fromParametrization}{gravityfieldType:fromParametrization}.
The monthly mean gravity field can be added back by additionaly selecting the time splines created
in \verb|000groopsBackgroundModels.xml| using
\configClass{gravityfield:timeSplines}{gravityfieldType:timeSplines}. The preprocessing solution
is saved as a \file{spherical harmonics file}{potentialCoefficients}.

\subsection{Setting up normal equations}\label{cookbook.gravityFieldGrace:monthlyNormals}
Normal equations are set up in the script \verb|040groopsMonthlyNormals120.xml| using
the program \program{NormalsBuildShortTimeStaticLongTime}. The time intervals which the normal
equations are divided into are defined in \configFile{inputfileArcList}{arcList}.
The normal equations are based on \configClass{observation}{observationType} including the SST data,
the kinematic orbits and the variational equations. The parametrization of the gravity field can
be set with \configClass{observation:parametrizationGravity}{parametrizationGravityType}
(e.g. spherical harmonics up to degree and order 120). Accelerometer calibration parameters
and KBR antenna phase center variations can be parameterized using
\configClass{parametrizationAcceleration}{parametrizationAccelerationType} and
\configClass{parametrizationSst}{parametrizationSatelliteTrackingType}.
With \config{estimateShortTimeVariations} short time variations of the gravity
field can be co-estimated. The parameters selected by
\configClass{parameterSelection}{parameterSelectorType} (e.g. linear splines with 6h nodes) are
constrained by an \configClass{autoregressiveModelSequence}{autoregressiveModelSequenceType}.
Additional temporal variations (e.g. trend and annual oscillation) could be estimated with
\config{estimateLongTimeVariations}.

\subsection{Solving normal equations}\label{cookbook.gravityFieldGrace:monhtlySolve}
The desired spherical harmonic coefficients are determined in the script \verb|050groopsMonthlySolve.xml|.
\program{NormalsSolverVCE} accumulates \configClass{normalEquation}{normalEquationType} and solves
the total combined system. Variance component estimation is used to determine the relative weighting
of the individual normals. The estimated parameter vector (\configFile{outputfileSolution}{matrix})
the estimated accuracies (\configFile{outputfileSigmax}{matrix}) and the full covariance matrix
(\configFile{outputfileCovariance}{matrix}) can be saved.
Using \program{Gravityfield2PotentialCoefficients} the final solution can be saved as
a \file{spherical harmonics file}{potentialCoefficients} by adding back the monthly mean gravity
field to the estimated spherical harmonic coefficients.

% =============================================================

\clearpage
\section{Regional geoid determination}\label{cookbook.regionalGeoid}
This shows exemplary the computation of a regional geoid using terrestrial gravimetric observations
in combination with a global satellite model such as GOCO06s. The geoid is estimated in a least squares adjustment with
a parametrization using radial basis functions. A detailed desciption of the method is given in
Christian Pock (2017), Consistent Combination of Satellite and Terrestrial Gravity Field Observations in Regional Geoid Modeling.
Dissertation TU Graz.


\subsection{Gravimetric data}
Here it is assumed that the measured absolute gravity data is given at points in ellipsoidal coordinates.
The observed values should be converted to SI units~$m/s^2$.
\begin{itemize}
  \item \program{Matrix2GriddedData} to convert data from text file in tabular form.
\end{itemize}

\fig{!ht}{0.5}{regionalGeoidGravimetricData}{fig:regionalGeoidGravimetricData}{Distribution of gravimetric observations}

\subsection{Topography}
A high resolution topography model is needed to reduce the observations.
As the model heights are usually given in physical heights a reference geoid is needed to compute the correct ellipsoidal height.
\begin{itemize}
  \item \program{NetCdf2GridRectangular} convert into groops format.
  \item \program{Gravityfield2GriddedData}: Compute geoid heights using the GOCO06s model with
  \begin{itemize}
    \item \configClass{grid:file}{gridType:file} The topography grid
    \item \configClass{kernel:geoidHeight}{kernelType:geoidHeight}
    \item \configClass{gravityfield:potentialCoefficients}{gravityfieldType:potentialCoefficients} the GOCO06s model
    \item \configClass{gravityfield:potentialCoefficients}{gravityfieldType:potentialCoefficients}
    Subtract (\config{factor}=-1) GRS680 normal field.
  \end{itemize}
  \item \program{GriddedDataCalculate}: Generate a new combined griddedData file
        with the orthometric height (\verb|data0|) and the geoid height (\verb|data1|).
  \item \program{GriddedTopography2PotentialCoefficients}: Compute the gravitational potential in terms
        of spherical harmonics up to a maximum degree of the global satellite model.
        This is the part of the topography, which is already included in the global satellite model.
        The integration boundaries are \config{radialUpperBound}=\verb|data0+data1|
        and \config{radialLowerBound}=\verb|data1|.
\end{itemize}

\fig{!ht}{1.0}{regionalGeoidTopography}{fig:regionalGeoidTopography}{Topography and geoid heights}

\subsection{Reduce}
Calculate approximate reference gravity to reduce it from the observations.
\begin{itemize}
  \item \program{Gravityfield2AbsoluteGravity}
  \begin{itemize}
    \item \configClass{grid:File}{gridType:file} at observation positions
    \item \configClass{gravityfield:tides}{gravityfieldType:tides} Centrifugal potential
    \item \configClass{gravityfield:potentialCoefficients}{gravityfieldType:potentialCoefficients} full GOCO06s model
    \item \configClass{gravityfield:topography}{gravityfieldType:topography}
          (\config{radialUpperBound}=\verb|data0+data1|, \config{radialLowerBound}=\verb|data1|)
    \item \configClass{gravityfield:potentialCoefficients}{gravityfieldType:potentialCoefficients}
          Subtract (\config{factor}=-1) the potential part of the topography already included in the GOCO06s model.
  \end{itemize}
  \item \program{GriddedDataCalculate} to calculate observed minus computed.
  \item Large outliers can be removed in \program{GriddedDataCalculate} with \config{removalCriteria}.
\end{itemize}

\fig{!hb}{0.5}{regionalGeoidObserved-computed}{fig:regionalGeoidObserved-computed}{Gravity disturbances: observed minus computed}


\subsection{Radial Basis Functions (RBF)}
The residual gravity is parametrized in terms of Radial Basis Functions
\configClass{parametrizationGravity:radialBasis}{parametrizationGravityType:radialBasis}.
The basis functions should be distributed on a regular \configClass{grid}{gridType} covering
a somewhat larger area than the observations, see \configClass{border}{borderType}.
The shape of the functions \configClass{kernel:coefficients}{kernelType:coefficients} should reflect
the signal content of reduced observations and are defined by the coefficients.
\begin{itemize}
  \item \program{RadialBasisSplines2KernelCoefficients}
  \begin{itemize}
    \item \configClass{gravityfield:potentialCoefficients}{gravityfieldType:potentialCoefficients} accuracies of GOCO06s model
    \item \config{maxDegree}=7000. Complemented by Kaula's rule of thumb
  \end{itemize}
\end{itemize}
The maximum degree should correspond to the spatial resolution.
Rule of thumb: the number of spherical harmonic coefficients $(\text{maxDegree}+1)^2$ should roughly agree
to the number of grid points if they would cover the complete Earth.

\fig{!hb}{0.5}{regionalGeoidDegreeAmplitudes}{fig:regionalGeoidDegreeAmplitudes}{Degree amplitudes for the shape of the radial basis functions}

\subsection{Compute: Estimate parameters in a least squares adjustment}
Setup the observation equations and accumulate the system of normal equations.
\begin{itemize}
\item \program{NormalsBuild}
  \begin{itemize}
    \item \configClass{normalEquation:design}{normalEquationType:design} with
          \configClass{observation:terrestrial}{observationType:terrestrial}
    \begin{itemize}
      \item \configClass{kernel:disturbance}{kernelType:disturbance}
      \item \configClass{parametrizationGravity:radialBasis}{parametrizationGravityType:radialBasis}
            with \configClass{kernel:coefficients}{kernelType:coefficients}
    \end{itemize}
  \end{itemize}
\item \program{NormalsSolverVCE}
  \begin{itemize}
    \item \configClass{normalEquation:file}{normalEquationType:file} from \program{NormalsBuild}
    \item \configClass{normalEquation:regularization}{normalEquationType:regularization}
           towards zero means regularization towards the GOCO06s, which is reduced from the data.
  \end{itemize}
\end{itemize}

\subsection{Restore: Calculate the geoid solution}
Evaluate the estimated parameters and add back the reduced reference models.
\begin{itemize}
  \item \program{Gravityfield2GriddedData}: Compute approximate geoid heights using the GOCO06s model with
  \begin{itemize}
    \item \configClass{grid}{gridType} select a grid with target resolution at elliposid
    \item \configClass{kernel:geoidHeight}{kernelType:geoidHeight}
    \item \configClass{gravityfield:potentialCoefficients}{gravityfieldType:potentialCoefficients} GOCO06s model
    \item \configClass{gravityfield:potentialCoefficients}{gravityfieldType:potentialCoefficients}
    Subtract (\config{factor}=-1) GRS80 normal field.
  \end{itemize}
  \item \program{GriddedDataCalculate}: Move points from ellipsoid to geoid with \config{height}=\verb|data0|
  \item \program{Gravityfield2GriddedData}
  \begin{itemize}
    \item \configClass{grid:file}{gridType:file} from above.
    \item \configClass{kernel:geoidHeight}{kernelType:geoidHeight}
    \item \configClass{gravityfield:fromParametrization}{gravityfieldType:fromParametrization}
          The solution vector~$\M x$ together with the RBF parametrization
    \item \configClass{gravityfield:potentialCoefficients}{gravityfieldType:potentialCoefficients} GOCO06s model
    \item \configClass{gravityfield:topography}{gravityfieldType:topography}
          (\config{radialUpperBound}=\verb|data0+data1|, \config{radialLowerBound}=\verb|data1|)
    \item \configClass{gravityfield:potentialCoefficients}{gravityfieldType:potentialCoefficients}
          Subtract (\config{factor}=-1) the potential part of the topography already included in the GOCO06s model.
    \item \configClass{gravityfield:potentialCoefficients}{gravityfieldType:potentialCoefficients}
          Subtract (\config{factor}=-1) GRS80 normal field.
    \item \config{convertToHarmonics}=\verb|no|, otherwise the RBF are converted to harmonics up to degree 7000.
  \end{itemize}
  \item \program{GriddedDataCalculate}: Set \config{height}=0 of the computed geoid grid.
  \item \program{GridRectangular2NetCdf}
\end{itemize}

\fig{!hb}{0.6}{regionalGeoidGeoid}{fig:regionalGeoidGeoid}{Estimated geoid}


% =============================================================


\chapter{Programs}\label{programType}
This chapter details programs included in GROOPS, describing what they are and what they do. For usage examples see the cookbook in Chapter~\ref{general.cookbook}.
% auto generated by GROOPS
\section{Programs: Covariance}
\subsection{AutoregressiveModel2CovarianceMatrix}\label{AutoregressiveModel2CovarianceMatrix}
This program computes the covariance structure of a random process represented by an AR model sequence.
The covariance matrix is determined by accumulating the normal equations of all AR models in \config{autoregressiveModelSequence}
and inverting the combined normal equation matrix.
For each output file in \configFile{outputfileCovarianceMatrix}{matrix},
the covariance matrix of appropriate time lag is saved (the first file contains the auto-covariance,
second file cross covariance and so on). The matrix for lag $h$ describes the covariance between $x_{t-h}$ and $x_{t}$, i.e. $\Sigma(t-h, t)$.


\keepXColumns
\begin{tabularx}{\textwidth}{N T A}
\hline
Name & Type & Annotation\\
\hline
\hfuzz=500pt\includegraphics[width=1em]{element-mustset-unbounded.pdf}~outputfileCovarianceMatrix & \hfuzz=500pt filename & \hfuzz=500pt covariance matrix for each lag\\
\hfuzz=500pt\includegraphics[width=1em]{element-mustset.pdf}~autoregressiveModelSequence & \hfuzz=500pt \hyperref[autoregressiveModelSequenceType]{autoregressiveModelSequence} & \hfuzz=500pt AR model sequence\\
\hline
\end{tabularx}

\clearpage
%==================================
\subsection{CovarianceFunction2DigitalFilter}\label{CovarianceFunction2DigitalFilter}
Computes digital filter coefficients for a \configClass{digital filter}{digitalFilterType} of given degree and
order. The filter coefficients are computed by fitting them to an approximated
impulse response represented by the cholesky factor of the covariance matrix.

The parameter \config{warmup} determines from which element of the cholesky matrix the
coefficients (default: half the covariance length) are fitted.

Per default, the program computes filter coefficients which generate colored noise
when applied to a white noise sequence. When \config{decorrelationFilter} is set,
a decorrelation filter is computed which yields white noise when applied to colored noise.


\keepXColumns
\begin{tabularx}{\textwidth}{N T A}
\hline
Name & Type & Annotation\\
\hline
\hfuzz=500pt\includegraphics[width=1em]{element-mustset.pdf}~outputfileFilter & \hfuzz=500pt filename & \hfuzz=500pt filter coefficients\\
\hfuzz=500pt\includegraphics[width=1em]{element-mustset.pdf}~inputfileCovariance & \hfuzz=500pt filename & \hfuzz=500pt first column: time steps, following columns: covariance functions\\
\hfuzz=500pt\includegraphics[width=1em]{element.pdf}~column & \hfuzz=500pt uint & \hfuzz=500pt Column with covariance function to be fitted\\
\hfuzz=500pt\includegraphics[width=1em]{element.pdf}~warmup & \hfuzz=500pt uint & \hfuzz=500pt number of samples until diagonal of Cholesky factor is flat (default: half covariance length)\\
\hfuzz=500pt\includegraphics[width=1em]{element.pdf}~numeratorDegree & \hfuzz=500pt uint & \hfuzz=500pt Maximum degree of numerator polynomial (MA constituent)\\
\hfuzz=500pt\includegraphics[width=1em]{element.pdf}~denominatorDegree & \hfuzz=500pt uint & \hfuzz=500pt Maximum degree of denominator polynomial (AR constitutent)\\
\hfuzz=500pt\includegraphics[width=1em]{element.pdf}~decorrelationFilter & \hfuzz=500pt boolean & \hfuzz=500pt compute a decorrelation filter\\
\hline
\end{tabularx}

\clearpage
%==================================
\subsection{CovarianceFunction2PowerSpectralDensity}\label{CovarianceFunction2PowerSpectralDensity}
One sided Power Spectral Density (PSD) from a covariance function. The first column of \configFile{inputfileCovarianceFunction}{matrix}
should contain the time lag in seconds.
Multiple covariance functions (in the following column)s are supported.
The output is a \file{matrix}{matrix} with first column contains the frequency $[Hz]$ and the other columns the PSD $[unit^2/Hz]$.

Conversion between covariance function $c_j$ and PSD $p_k$ is performed by discrete cosine transformation:
\begin{equation}
p_k = 2\Delta t\left(c_0 + c_{n-1} (-1)^k + \sum_{j=1}^{n-2} 2 c_j \cos(\pi jk/(n-1))\right).
\end{equation}

See also \program{PowerSpectralDensity2CovarianceFunction}.


\keepXColumns
\begin{tabularx}{\textwidth}{N T A}
\hline
Name & Type & Annotation\\
\hline
\hfuzz=500pt\includegraphics[width=1em]{element-mustset.pdf}~outputfilePSD & \hfuzz=500pt filename & \hfuzz=500pt first column: frequency [Hz], other columns PSD [unit\textasciicircum{}2/Hz]\\
\hfuzz=500pt\includegraphics[width=1em]{element-mustset.pdf}~inputfileCovarianceFunction & \hfuzz=500pt filename & \hfuzz=500pt first column: time steps, following columns: covariance functions\\
\hline
\end{tabularx}

\clearpage
%==================================
\subsection{CovarianceMatrix2AutoregressiveModel}\label{CovarianceMatrix2AutoregressiveModel}
This program computes a VAR(p) model from empirical covariance matrices.
The \configFile{inputfileCovarianceMatrix}{matrix} represent the covariance structure of the process:
the first file should contain the auto-covariance, the second the cross-covariance of lag one,
the next cross-covariance of lag two and so on.

Cross-covariance matrices $\Sigma_{\Delta_k}$ are defined as the cross-covariance between epoch $t-k$ and $t$.
If the process realizations $x_{t}$ are arrange by ascending time stamps
($\{\dots, x_{t-2}, x_{t-1}, x_{t}, x_{t+1}, x_{t+2},\dots\}$),
the covariance structure of the (stationary) process is therefore given by
\begin{equation}
\begin{bmatrix}
\Sigma & \Sigma_{\Delta_1} & \Sigma_{\Delta_2} & \cdots \\
\Sigma_{\Delta_1}^T & \Sigma & \Sigma_{\Delta_1} &  \cdots \\
\Sigma_{\Delta_2}^T & \Sigma_{\Delta_1}^T & \Sigma & \cdots \\
\vdots & \vdots & \vdots & \ddots \\
\end{bmatrix}.
\end{equation}

The estimate AR model is saved as single matrix \config{outputfileAutoregressiveModel} according to the GROOPS AR model conventions.


\keepXColumns
\begin{tabularx}{\textwidth}{N T A}
\hline
Name & Type & Annotation\\
\hline
\hfuzz=500pt\includegraphics[width=1em]{element-mustset.pdf}~outputfileAutoregressiveModel & \hfuzz=500pt filename & \hfuzz=500pt coefficients and white noise covariance of AR(p) model\\
\hfuzz=500pt\includegraphics[width=1em]{element-mustset-unbounded.pdf}~inputfileCovarianceMatrix & \hfuzz=500pt filename & \hfuzz=500pt file name of covariance matrix\\
\hline
\end{tabularx}

\clearpage
%==================================
\subsection{CovarianceMatrix2Correlation}\label{CovarianceMatrix2Correlation}
This program computes the pearson correlation coefficient
\begin{equation}
  \rho_{ij} = \frac{\sigma_{ij}}{\sigma_i \sigma_j}
\end{equation}
from a given covariance matrix stored in \configFile{inputfileCovarianceMatrix}{matrix}.
The result is stored in \configFile{outputfileCorrelationMatrix}{matrix}.


\keepXColumns
\begin{tabularx}{\textwidth}{N T A}
\hline
Name & Type & Annotation\\
\hline
\hfuzz=500pt\includegraphics[width=1em]{element-mustset.pdf}~outputfileCorrelationMatrix & \hfuzz=500pt filename & \hfuzz=500pt correlation matrix\\
\hfuzz=500pt\includegraphics[width=1em]{element-mustset.pdf}~inputfileCovarianceMatrix & \hfuzz=500pt filename & \hfuzz=500pt covariance matrix\\
\hline
\end{tabularx}

This program is \reference{parallelized}{general.parallelization}.
\clearpage
%==================================
\subsection{PowerSpectralDensity2CovarianceFunction}\label{PowerSpectralDensity2CovarianceFunction}
Covariance function from Power Spectral Density (PSD).
The \configFile{inputfilePSD}{matrix} contains in the first column the frequency $[Hz]$, followed by (possibly multiple) PSDs $[unit^2/Hz]$.
The output is a \file{matrix}{matrix}, the first column containing time lag $[s]$ and the other columns the covariance functions $[unit^2]$.
Conversion between PSD $p_j$ and covariance function $c_k$ is performed by discrete cosine transformation:
\begin{equation}
c_k = \frac{1}{4\Delta t (n-1)}\left(p_0 + p_{n-1} (-1)^k + \sum_{j=1}^{n-2} 2 p_j \cos(\pi jk/(n-1))\right).
\end{equation}

See also \program{CovarianceFunction2PowerSpectralDensity}.


\keepXColumns
\begin{tabularx}{\textwidth}{N T A}
\hline
Name & Type & Annotation\\
\hline
\hfuzz=500pt\includegraphics[width=1em]{element-mustset.pdf}~outputfileCovarianceFunction & \hfuzz=500pt filename & \hfuzz=500pt first column: time steps [seconds], following columns: covariance functions\\
\hfuzz=500pt\includegraphics[width=1em]{element-mustset.pdf}~inputfilePSD & \hfuzz=500pt filename & \hfuzz=500pt first column: frequency [Hz], following columns PSD [unit\textasciicircum{}2/Hz]\\
\hline
\end{tabularx}

\clearpage
%==================================
\section{Programs: DoodsonHarmonics}
\subsection{DoodsonAdmittanceInterpolation}\label{DoodsonAdmittanceInterpolation}
To visualize the interpolation of the minor tides.
The output is a \file{matrix}{matrix} with the first column containing the tidal frequency,
the second column is the tide generating amplitude (from \configFile{inputfileTideGeneratingPotential}{tideGeneratingPotential}), and the following
columns the contribution of the major tides to the this tidal frequency as defined in in \configFile{inputfileAdmittance}{admittance}.

\fig{!hb}{0.8}{doodsonAdmittanceInterpolation}{fig:doodsonAdmittanceInterpolation}{Linear interpolation of minor tides in the diurnal band.}


\keepXColumns
\begin{tabularx}{\textwidth}{N T A}
\hline
Name & Type & Annotation\\
\hline
\hfuzz=500pt\includegraphics[width=1em]{element-mustset.pdf}~outputfile & \hfuzz=500pt filename & \hfuzz=500pt \\
\hfuzz=500pt\includegraphics[width=1em]{element-mustset.pdf}~inputfileAdmittance & \hfuzz=500pt filename & \hfuzz=500pt interpolation of minor constituents\\
\hfuzz=500pt\includegraphics[width=1em]{element.pdf}~inputfileTideGeneratingPotential & \hfuzz=500pt filename & \hfuzz=500pt \\
\hline
\end{tabularx}

\clearpage
%==================================
\subsection{DoodsonAdmittanceTimeSeries}\label{DoodsonAdmittanceTimeSeries}
To visualize the interpolation of the minor tides it computes cosine multipliers of all major tides.
Without admittance this would be a simple cos oscillation.
The \config{outputfileTimeSeries} is an \file{instrument file}{instrument} (MISCVALUES) containining the cos of all the major tides.

\fig{!hb}{0.8}{doodsonAdmittanceTimeSeries}{fig:doodsonAdmittanceTimeSeries}{Cosine of the Mf tidal frequency with modulation from the interpolated minor tides.}


\keepXColumns
\begin{tabularx}{\textwidth}{N T A}
\hline
Name & Type & Annotation\\
\hline
\hfuzz=500pt\includegraphics[width=1em]{element-mustset.pdf}~outputfileTimeSeries & \hfuzz=500pt filename & \hfuzz=500pt MISCVALUES (cos of major tides, ...)\\
\hfuzz=500pt\includegraphics[width=1em]{element-mustset.pdf}~inputfileAdmittance & \hfuzz=500pt filename & \hfuzz=500pt cos/sin multipliers of the major tides\\
\hfuzz=500pt\includegraphics[width=1em]{element-mustset-unbounded.pdf}~timeSeries & \hfuzz=500pt \hyperref[timeSeriesType]{timeSeries} & \hfuzz=500pt \\
\hline
\end{tabularx}

\clearpage
%==================================
\subsection{DoodsonArguments2TimeSeries}\label{DoodsonArguments2TimeSeries}
Time series of doodson/fundamental arguments.
The \configFile{outputfileTimeSeries}{instrument} contains the six Doodson arguments,
followed by the five fundamental arguments in radians.


\keepXColumns
\begin{tabularx}{\textwidth}{N T A}
\hline
Name & Type & Annotation\\
\hline
\hfuzz=500pt\includegraphics[width=1em]{element-mustset.pdf}~outputfileTimeSeries & \hfuzz=500pt filename & \hfuzz=500pt each epoch: 6 doodson args, 5 fundamental args [rad]\\
\hfuzz=500pt\includegraphics[width=1em]{element-mustset-unbounded.pdf}~timeSeries & \hfuzz=500pt \hyperref[timeSeriesType]{timeSeries} & \hfuzz=500pt \\
\hline
\end{tabularx}

\clearpage
%==================================
\subsection{DoodsonHarmonics2GriddedAmplitudeAndPhase}\label{DoodsonHarmonics2GriddedAmplitudeAndPhase}
This program reads a \configFile{inputfileDoodsonHarmonics}{doodsonHarmonic} and evaluates a single tidal
constituent selected by \config{dooddson} (Doodson number or Darwin´s name, e.g. 255.555 or M2).
This program computes the amplitude and phase from the cos and sin coefficients on
a given \configClass{grid}{gridType}. The type of functional (e.g gravity anomalies or geoid heights)
can be choosen with \configClass{kernel}{kernelType}.
The values will be saved together with points expressed as ellipsoidal coordinates (longitude, latitude, height)
based on a reference ellipsoid with parameters \config{R} and \config{inverseFlattening}.
To visualize the results use \program{PlotMap}.

\fig{!hb}{1.}{doodsonHarmonics2GriddedAmplitudeAndPhase}{fig:doodsonHarmonics2GriddedAmplitudeAndPhase}{M2 amplitude and phase of FES2014b.}


\keepXColumns
\begin{tabularx}{\textwidth}{N T A}
\hline
Name & Type & Annotation\\
\hline
\hfuzz=500pt\includegraphics[width=1em]{element-mustset.pdf}~outputfileGrid & \hfuzz=500pt filename & \hfuzz=500pt ampl, phase [-pi,pi], cos, sin\\
\hfuzz=500pt\includegraphics[width=1em]{element-mustset.pdf}~inputfileDoodsonHarmonics & \hfuzz=500pt filename & \hfuzz=500pt \\
\hfuzz=500pt\includegraphics[width=1em]{element-mustset.pdf}~doodson & \hfuzz=500pt \hyperref[doodson]{doodson} & \hfuzz=500pt tidal constituent\\
\hfuzz=500pt\includegraphics[width=1em]{element-unbounded.pdf}~filter & \hfuzz=500pt \hyperref[sphericalHarmonicsFilterType]{sphericalHarmonicsFilter} & \hfuzz=500pt \\
\hfuzz=500pt\includegraphics[width=1em]{element-mustset-unbounded.pdf}~grid & \hfuzz=500pt \hyperref[gridType]{grid} & \hfuzz=500pt \\
\hfuzz=500pt\includegraphics[width=1em]{element-mustset.pdf}~kernel & \hfuzz=500pt \hyperref[kernelType]{kernel} & \hfuzz=500pt \\
\hfuzz=500pt\includegraphics[width=1em]{element.pdf}~minDegree & \hfuzz=500pt uint & \hfuzz=500pt \\
\hfuzz=500pt\includegraphics[width=1em]{element.pdf}~maxDegree & \hfuzz=500pt uint & \hfuzz=500pt \\
\hfuzz=500pt\includegraphics[width=1em]{element.pdf}~factor & \hfuzz=500pt double & \hfuzz=500pt the values on grid are multiplied by this factor\\
\hfuzz=500pt\includegraphics[width=1em]{element.pdf}~R & \hfuzz=500pt double & \hfuzz=500pt reference radius for ellipsoidal coordinates on output\\
\hfuzz=500pt\includegraphics[width=1em]{element.pdf}~inverseFlattening & \hfuzz=500pt double & \hfuzz=500pt reference flattening for ellipsoidal coordinates on output, 0: spherical coordinates\\
\hline
\end{tabularx}

This program is \reference{parallelized}{general.parallelization}.
\clearpage
%==================================
\subsection{DoodsonHarmonics2PotentialCoefficients}\label{DoodsonHarmonics2PotentialCoefficients}
The \configFile{inputfileDoodsonHarmonics}{doodsonHarmonic} contains a Fourier series of a time variable
gravitational potential at specific tidal frequencies (tides)
\begin{equation}
V(\M x,t) = \sum_{f} V_f^c(\M x)\cos(\theta_f(t)) + V_f^s(\M x)\sin(\theta_f(t)),
\end{equation}
where $V_f^c(\M x)$ and $V_f^s(\M x)$ are spherical harmonics expansions.
If set the expansions are limited in the range between \config{minDegree}
and \config{maxDegree} inclusivly. The coefficients are related to the reference radius~\config{R}
and the Earth gravitational constant \config{GM}.

The \configFile{outputfilePotentialCoefficients}{potentialCoefficients} is not a single file but a series of files.
For each spherical harmonics expansion $V_f^c(\M x)$ and $V_f^s(\M x)$ a separate file is created
where the variables \config{variableLoopName}, \config{variableLoopDoodson}, \config{variableLoopCosSin} are set accordingly.
The file name should contain these variables, e.g. \verb|coeff.{name}.{doodson}.{cossin}.gfc|.

If \config{applyXi} the Doodson-Warburg phase correction (see IERS conventions) is applied to the cos/sin
potentialCoefficients before.


\keepXColumns
\begin{tabularx}{\textwidth}{N T A}
\hline
Name & Type & Annotation\\
\hline
\hfuzz=500pt\includegraphics[width=1em]{element-mustset.pdf}~outputfilePotentialCoefficients & \hfuzz=500pt filename & \hfuzz=500pt \\
\hfuzz=500pt\includegraphics[width=1em]{element.pdf}~variableLoopName & \hfuzz=500pt string & \hfuzz=500pt variable with darwins's name of each constituent\\
\hfuzz=500pt\includegraphics[width=1em]{element.pdf}~variableLoopDoodson & \hfuzz=500pt string & \hfuzz=500pt variable with doodson code of each constituent\\
\hfuzz=500pt\includegraphics[width=1em]{element.pdf}~variableLoopCosSin & \hfuzz=500pt string & \hfuzz=500pt variable with 'cos' or 'sin' of each constituent\\
\hfuzz=500pt\includegraphics[width=1em]{element.pdf}~variableLoopIndex & \hfuzz=500pt string & \hfuzz=500pt variable with index of each constituent (starts with zero)\\
\hfuzz=500pt\includegraphics[width=1em]{element.pdf}~variableLoopCount & \hfuzz=500pt string & \hfuzz=500pt variable with total number of constituents\\
\hfuzz=500pt\includegraphics[width=1em]{element-mustset.pdf}~inputfileDoodsonHarmonics & \hfuzz=500pt filename & \hfuzz=500pt \\
\hfuzz=500pt\includegraphics[width=1em]{element.pdf}~inputfileTideGeneratingPotential & \hfuzz=500pt filename & \hfuzz=500pt to compute Xi phase correction\\
\hfuzz=500pt\includegraphics[width=1em]{element.pdf}~minDegree & \hfuzz=500pt uint & \hfuzz=500pt \\
\hfuzz=500pt\includegraphics[width=1em]{element.pdf}~maxDegree & \hfuzz=500pt uint & \hfuzz=500pt \\
\hfuzz=500pt\includegraphics[width=1em]{element.pdf}~GM & \hfuzz=500pt double & \hfuzz=500pt Geocentric gravitational constant\\
\hfuzz=500pt\includegraphics[width=1em]{element.pdf}~R & \hfuzz=500pt double & \hfuzz=500pt reference radius\\
\hfuzz=500pt\includegraphics[width=1em]{element.pdf}~applyXi & \hfuzz=500pt boolean & \hfuzz=500pt apply Doodson-Warburg phase correction (see IERS conventions)\\
\hline
\end{tabularx}

\clearpage
%==================================
\subsection{DoodsonHarmonicsCalculateAdmittance}\label{DoodsonHarmonicsCalculateAdmittance}
Computes the admittance function to interpolate minor tides from
tides given in \configFile{inputfileDoodsonHarmonics}{doodsonHarmonic}
using \configFile{inputfileTideGeneratingPotential}{tideGeneratingPotential}.


\keepXColumns
\begin{tabularx}{\textwidth}{N T A}
\hline
Name & Type & Annotation\\
\hline
\hfuzz=500pt\includegraphics[width=1em]{element-mustset.pdf}~outputfileAdmittance & \hfuzz=500pt filename & \hfuzz=500pt \\
\hfuzz=500pt\includegraphics[width=1em]{element-mustset.pdf}~inputfileDoodsonHarmonics & \hfuzz=500pt filename & \hfuzz=500pt \\
\hfuzz=500pt\includegraphics[width=1em]{element.pdf}~inputfileTideGeneratingPotential & \hfuzz=500pt filename & \hfuzz=500pt TGP\\
\hfuzz=500pt\includegraphics[width=1em]{element.pdf}~threshold & \hfuzz=500pt double & \hfuzz=500pt [m\textasciicircum{}2/s\textasciicircum{}2] only interpolate tides with TGP greater than threshold\\
\hfuzz=500pt\includegraphics[width=1em]{element.pdf}~degreeInterpolation & \hfuzz=500pt uint & \hfuzz=500pt polynomial degree for interpolation\\
\hfuzz=500pt\includegraphics[width=1em]{element.pdf}~degreeExtrapolation & \hfuzz=500pt uint & \hfuzz=500pt polynomial degree for extrapolation\\
\hfuzz=500pt\includegraphics[width=1em]{element-unbounded.pdf}~excludeDoodsonForInterpolation & \hfuzz=500pt \hyperref[doodson]{doodson} & \hfuzz=500pt major tides not used for interpolation\\
\hline
\end{tabularx}

\clearpage
%==================================
\subsection{DoodsonHarmonicsChangePartialTides}\label{DoodsonHarmonicsChangePartialTides}
Reads a file \configFile{inputfileDoodsonHarmonic}{doodsonHarmonic} and write it to
\configFile{outputfileDoodsonHarmonics}{doodsonHarmonic}. If set the spherical harmonics
expansion is limited in the range between \config{minDegree} and \config{maxDegree} inclusivly.
The \configClass{useDoodson}{doodson} and \configClass{ignoreDoodson}{doodson} can be used
to filter the partial types that will be exported.
Additional partial tides can be interpolated using the file \configFile{inputfileAdmittance}{admittance}.


\keepXColumns
\begin{tabularx}{\textwidth}{N T A}
\hline
Name & Type & Annotation\\
\hline
\hfuzz=500pt\includegraphics[width=1em]{element-mustset.pdf}~outputfileDoodsonHarmonics & \hfuzz=500pt filename & \hfuzz=500pt \\
\hfuzz=500pt\includegraphics[width=1em]{element-mustset.pdf}~inputfileDoodsonHarmonics & \hfuzz=500pt filename & \hfuzz=500pt \\
\hfuzz=500pt\includegraphics[width=1em]{element.pdf}~inputfileAdmittance & \hfuzz=500pt filename & \hfuzz=500pt interpolation of minor constituents\\
\hfuzz=500pt\includegraphics[width=1em]{element-unbounded.pdf}~useDoodson & \hfuzz=500pt \hyperref[doodson]{doodson} & \hfuzz=500pt use only these partial tides (additional tides will be interpolated)\\
\hfuzz=500pt\includegraphics[width=1em]{element-unbounded.pdf}~ignoreDoodson & \hfuzz=500pt \hyperref[doodson]{doodson} & \hfuzz=500pt ignore these partial tides\\
\hfuzz=500pt\includegraphics[width=1em]{element.pdf}~minDegree & \hfuzz=500pt uint & \hfuzz=500pt \\
\hfuzz=500pt\includegraphics[width=1em]{element.pdf}~maxDegree & \hfuzz=500pt uint & \hfuzz=500pt \\
\hfuzz=500pt\includegraphics[width=1em]{element.pdf}~GM & \hfuzz=500pt double & \hfuzz=500pt Geocentric gravitational constant\\
\hfuzz=500pt\includegraphics[width=1em]{element.pdf}~R & \hfuzz=500pt double & \hfuzz=500pt reference radius\\
\hline
\end{tabularx}

\clearpage
%==================================
\subsection{ModelEquilibriumTide}\label{ModelEquilibriumTide}
Computes the equilibrium ocean tide of the long periodic \config{tideGeneratingPotential}.

The ocean surface is represented by \configClass{gridOcean}{gridType} and the gravitational
effect is numerical integrated to spherical harmonics using \config{maxDegree}, \config{GM},
and \config{R}.

It takes self attraction and loading into account using the Love numbers
\configFile{inputfilePotentialLoadLoveNumber}{matrix} and
\configFile{inputfileDeformationLoadLoveNumber}{matrix}.

Additionally the effects of the solid Earth tide are considered,
both the gravitational (Love numbers \config{k20}, \config{k20plus})
and the geometrical (Love numbers \config{h20,0}, \config{h20,2}) effect.

See also \program{PotentialCoefficients2DoodsonHarmonics}.

\fig{!hb}{0.8}{modelEquilibriumTide}{fig:modelEquilibriumTide}{Equilibrium tide of SA constituent}


\keepXColumns
\begin{tabularx}{\textwidth}{N T A}
\hline
Name & Type & Annotation\\
\hline
\hfuzz=500pt\includegraphics[width=1em]{element-mustset.pdf}~outputfilePotentialCoefficients & \hfuzz=500pt filename & \hfuzz=500pt includes the loading\\
\hfuzz=500pt\includegraphics[width=1em]{element-mustset-unbounded.pdf}~gridOcean & \hfuzz=500pt \hyperref[gridType]{grid} & \hfuzz=500pt \\
\hfuzz=500pt\includegraphics[width=1em]{element.pdf}~maxDegree & \hfuzz=500pt uint & \hfuzz=500pt \\
\hfuzz=500pt\includegraphics[width=1em]{element.pdf}~GM & \hfuzz=500pt double & \hfuzz=500pt Geocentric gravitational constant\\
\hfuzz=500pt\includegraphics[width=1em]{element.pdf}~R & \hfuzz=500pt double & \hfuzz=500pt reference radius\\
\hfuzz=500pt\includegraphics[width=1em]{element.pdf}~density & \hfuzz=500pt double & \hfuzz=500pt [kg/m**3] density of sea water\\
\hfuzz=500pt\includegraphics[width=1em]{element-mustset.pdf}~tideGeneratingPotential & \hfuzz=500pt double & \hfuzz=500pt [m**2/s**2]\\
\hfuzz=500pt\includegraphics[width=1em]{element.pdf}~k20 & \hfuzz=500pt double & \hfuzz=500pt earth tide love number\\
\hfuzz=500pt\includegraphics[width=1em]{element.pdf}~k20plus & \hfuzz=500pt double & \hfuzz=500pt earth tide love number\\
\hfuzz=500pt\includegraphics[width=1em]{element.pdf}~h20\_0 & \hfuzz=500pt double & \hfuzz=500pt earth tide love number\\
\hfuzz=500pt\includegraphics[width=1em]{element.pdf}~h20\_2 & \hfuzz=500pt double & \hfuzz=500pt earth tide love number\\
\hfuzz=500pt\includegraphics[width=1em]{element-mustset.pdf}~inputfilePotentialLoadLoveNumber & \hfuzz=500pt filename & \hfuzz=500pt \\
\hfuzz=500pt\includegraphics[width=1em]{element-mustset.pdf}~inputfileDeformationLoadLoveNumber & \hfuzz=500pt filename & \hfuzz=500pt \\
\hline
\end{tabularx}

This program is \reference{parallelized}{general.parallelization}.
\clearpage
%==================================
\subsection{PotentialCoefficients2DoodsonHarmonics}\label{PotentialCoefficients2DoodsonHarmonics}
Create a \file{DoodsonHarmonic file}{doodsonHarmonic} from a list of
cos/sin \file{potentialCoefficients}{potentialCoefficients} for given \config{doodson}
(Doodson number or Darwin´s name, e.g. 255.555 or M2) tidal constituents.
If \config{applyXi} the Doodson-Warburg phase correction (see IERS conventions) is applied before.


\keepXColumns
\begin{tabularx}{\textwidth}{N T A}
\hline
Name & Type & Annotation\\
\hline
\hfuzz=500pt\includegraphics[width=1em]{element-mustset.pdf}~outputfileDoodsonHarmonics & \hfuzz=500pt filename & \hfuzz=500pt \\
\hfuzz=500pt\includegraphics[width=1em]{element.pdf}~inputfileTideGeneratingPotential & \hfuzz=500pt filename & \hfuzz=500pt to compute Xi phase correction\\
\hfuzz=500pt\includegraphics[width=1em]{element-mustset-unbounded.pdf}~constituent & \hfuzz=500pt sequence & \hfuzz=500pt \\
\hfuzz=500pt\includegraphics[width=1em]{connector.pdf}\includegraphics[width=1em]{element-mustset.pdf}~doodson & \hfuzz=500pt \hyperref[doodson]{doodson} & \hfuzz=500pt \\
\hfuzz=500pt\includegraphics[width=1em]{connector.pdf}\includegraphics[width=1em]{element-mustset.pdf}~inputfileCosPotentialCoefficients & \hfuzz=500pt filename & \hfuzz=500pt \\
\hfuzz=500pt\includegraphics[width=1em]{connector.pdf}\includegraphics[width=1em]{element-mustset.pdf}~inputfileSinPotentialCoefficients & \hfuzz=500pt filename & \hfuzz=500pt \\
\hfuzz=500pt\includegraphics[width=1em]{element.pdf}~minDegree & \hfuzz=500pt uint & \hfuzz=500pt \\
\hfuzz=500pt\includegraphics[width=1em]{element.pdf}~maxDegree & \hfuzz=500pt uint & \hfuzz=500pt \\
\hfuzz=500pt\includegraphics[width=1em]{element.pdf}~GM & \hfuzz=500pt double & \hfuzz=500pt Geocentric gravitational constant\\
\hfuzz=500pt\includegraphics[width=1em]{element.pdf}~R & \hfuzz=500pt double & \hfuzz=500pt reference radius\\
\hfuzz=500pt\includegraphics[width=1em]{element.pdf}~applyXi & \hfuzz=500pt boolean & \hfuzz=500pt apply Doodson-Warburg phase correction (see IERS conventions)\\
\hline
\end{tabularx}

\clearpage
%==================================
\section{Programs: Gnss}
\subsection{GnssAntennaDefinition2ParameterVector}\label{GnssAntennaDefinition2ParameterVector}
Estimates parameters of a parametrization of \configClass{antennaCenterVariations}{parametrizationGnssAntennaType},
which represents all antennas from \configFile{inputfileAntennaDefinition}{gnssAntennaDefinition}
matching the wildcard patterns of \config{name}, \config{serial}, \config{radome}.

The provided values at the area weighted grid points of the pattern of each gnssType are used as pseudo-observations.
A subset of patterns can be selected with \configClass{types}{gnssType}.

The \file{GnssAntennaDefinition file}{gnssAntennaDefinition} can be modified to the demands before with
\program{GnssAntennaDefinitionCreate}.

See also \program{ParameterVector2GnssAntennaDefinition}.


\keepXColumns
\begin{tabularx}{\textwidth}{N T A}
\hline
Name & Type & Annotation\\
\hline
\hfuzz=500pt\includegraphics[width=1em]{element-mustset.pdf}~outputfileSolution & \hfuzz=500pt filename & \hfuzz=500pt \\
\hfuzz=500pt\includegraphics[width=1em]{element.pdf}~outputfileParameterNames & \hfuzz=500pt filename & \hfuzz=500pt \\
\hfuzz=500pt\includegraphics[width=1em]{element-mustset-unbounded.pdf}~antennaCenterVariations & \hfuzz=500pt \hyperref[parametrizationGnssAntennaType]{parametrizationGnssAntenna} & \hfuzz=500pt \\
\hfuzz=500pt\includegraphics[width=1em]{element-mustset.pdf}~inputfileAntennaDefinition & \hfuzz=500pt filename & \hfuzz=500pt \\
\hfuzz=500pt\includegraphics[width=1em]{element.pdf}~name & \hfuzz=500pt string & \hfuzz=500pt \\
\hfuzz=500pt\includegraphics[width=1em]{element.pdf}~serial & \hfuzz=500pt string & \hfuzz=500pt \\
\hfuzz=500pt\includegraphics[width=1em]{element.pdf}~radome & \hfuzz=500pt string & \hfuzz=500pt \\
\hfuzz=500pt\includegraphics[width=1em]{element-unbounded.pdf}~types & \hfuzz=500pt \hyperref[gnssType]{gnssType} & \hfuzz=500pt if not set, all types in the file are used\\
\hfuzz=500pt\includegraphics[width=1em]{element.pdf}~zeroNaN & \hfuzz=500pt boolean & \hfuzz=500pt treat NaN values as zero, otherwise values are ignored\\
\hline
\end{tabularx}

\clearpage
%==================================
\subsection{GnssAntennaDefinition2Skyplot}\label{GnssAntennaDefinition2Skyplot}
Produce a \file{skyplot}{griddedData} of antenna center variations
which can be plotted with \program{PlotMap}.

The first antenna from \configFile{inputfileAntennaDefinition}{gnssAntennaDefinition}
matching the wildcard patterns of \config{name}, \config{serial}, \config{radome} is used.

For each antenna pattern (gnssType) a separate data column is computed.
A subset of patterns can be selected with \configClass{types}{gnssType}.

Azimuth and elevation are written as ellipsoidal longitude and latitude in a \file{griddedData file}{griddedData}.
The choosen ellipsoid parameters \config{R} and \config{inverseFlattening} are arbitrary but should be the same
as in \configClass{grid}{gridType} and \program{PlotMap}.

\fig{!hb}{1.0}{fileFormatGnssAntennaDefinition}{fig:gnssAntennaDefinition2Skyplot}{Antenna Center Variations of ASH701945D\_M for two frequencies of GPS and GLONASS}


\keepXColumns
\begin{tabularx}{\textwidth}{N T A}
\hline
Name & Type & Annotation\\
\hline
\hfuzz=500pt\includegraphics[width=1em]{element-mustset.pdf}~outputfileGriddedData & \hfuzz=500pt filename & \hfuzz=500pt data column for each gnssType\\
\hfuzz=500pt\includegraphics[width=1em]{element-mustset.pdf}~inputfileAntennaDefinition & \hfuzz=500pt filename & \hfuzz=500pt \\
\hfuzz=500pt\includegraphics[width=1em]{element-mustset-unbounded.pdf}~grid & \hfuzz=500pt \hyperref[gridType]{grid} & \hfuzz=500pt \\
\hfuzz=500pt\includegraphics[width=1em]{element.pdf}~name & \hfuzz=500pt string & \hfuzz=500pt \\
\hfuzz=500pt\includegraphics[width=1em]{element.pdf}~serial & \hfuzz=500pt string & \hfuzz=500pt \\
\hfuzz=500pt\includegraphics[width=1em]{element.pdf}~radome & \hfuzz=500pt string & \hfuzz=500pt \\
\hfuzz=500pt\includegraphics[width=1em]{element-unbounded.pdf}~types & \hfuzz=500pt \hyperref[gnssType]{gnssType} & \hfuzz=500pt if not set, all types in the file are used\\
\hfuzz=500pt\includegraphics[width=1em]{element.pdf}~R & \hfuzz=500pt double & \hfuzz=500pt reference radius for ellipsoidal coordinates\\
\hfuzz=500pt\includegraphics[width=1em]{element.pdf}~inverseFlattening & \hfuzz=500pt double & \hfuzz=500pt reference flattening for ellipsoidal coordinates\\
\hline
\end{tabularx}

\clearpage
%==================================
\subsection{GnssAntennaDefinitionCreate}\label{GnssAntennaDefinitionCreate}
Create a \file{GNSS antenna definition file}{gnssAntennaDefinition} (Antenna Center Variations, ACV) consisting of multiple antennas.
The antennas can be created from scratch or can be selected from existing files.
This program can also be used to modify existing files.

Furthermore it can be used to create accuracy definition files containing azimuth and elevation dependent accuracy values for antennas.
To create an accuracy pattern for phase observations with \verb|1 mm| accuracy at zenith and no azimuth dependency, define a
pattern with \config{type}=\verb|L|, \config{values}=\verb|0.001/cos(zenith/rho)|.

The antennas in \configFile{outputfileAntennaDefinition}{gnssAntennaDefinition}
are sorted by names and duplicates are removed (first one is kept).


\keepXColumns
\begin{tabularx}{\textwidth}{N T A}
\hline
Name & Type & Annotation\\
\hline
\hfuzz=500pt\includegraphics[width=1em]{element-mustset.pdf}~outputfileAntennaDefinition & \hfuzz=500pt filename & \hfuzz=500pt \\
\hfuzz=500pt\includegraphics[width=1em]{element-mustset-unbounded.pdf}~antenna & \hfuzz=500pt \hyperref[gnssAntennaDefintionListType]{gnssAntennaDefintionList} & \hfuzz=500pt \\
\hline
\end{tabularx}

\clearpage
%==================================
\subsection{GnssAntennaNormalsConstraint}\label{GnssAntennaNormalsConstraint}
Apply constraints to \file{normal equations}{normalEquation}
containing \configClass{antennaCenterVariations}{parametrizationGnssAntennaType}.
Usually the antenna center variations are estimated together with other parameters
like station coordinates, signal biases and slant TEC in \program{GnssProcessing}.
This results in a rank deficient matrix as not all parameters can be separated.
The deficient can be solved by adding pseudo observation equations as constraints.

To separate antenna center variations and signal biases
apply \config{constraint:mean} for each GNSS \configClass{type}{gnssType}.
The observation equation for the integral mean of antenna center variations (ACV)
in all azimuth~$A$ and elevation~$E$ dependent directions
\begin{equation}
  0 = \iint ACV(A,E)\, d\Phi \approx \sum_i ACV(A_i,E_i)\, \Delta\Phi_i
\end{equation}
is approximated by a grid defined by
\config{deltaAzimuth}, \config{deltaZenith}, and \config{maxZenith}.

To separate from station coordinates use \config{constraint:centerMean}
and from slant TEC parameters use \config{constraint:TEC}.

The constraints are applied separately to all antennas matching
the wildcard patterns of \config{name}, \config{serial}, \config{radome}.

See also \program{ParameterVector2GnssAntennaDefinition}.


\keepXColumns
\begin{tabularx}{\textwidth}{N T A}
\hline
Name & Type & Annotation\\
\hline
\hfuzz=500pt\includegraphics[width=1em]{element-mustset.pdf}~outputfileNormalEquation & \hfuzz=500pt filename & \hfuzz=500pt with applied constraints\\
\hfuzz=500pt\includegraphics[width=1em]{element-mustset.pdf}~inputfileNormalEquation & \hfuzz=500pt filename & \hfuzz=500pt \\
\hfuzz=500pt\includegraphics[width=1em]{element-mustset-unbounded.pdf}~constraint & \hfuzz=500pt choice & \hfuzz=500pt \\
\hfuzz=500pt\includegraphics[width=1em]{connector.pdf}\includegraphics[width=1em]{element-mustset.pdf}~center & \hfuzz=500pt sequence & \hfuzz=500pt zero center (x,y,z) of a single pattern\\
\hfuzz=500pt\quad\includegraphics[width=1em]{connector.pdf}\includegraphics[width=1em]{element-mustset.pdf}~type & \hfuzz=500pt \hyperref[gnssType]{gnssType} & \hfuzz=500pt applied for each matching types\\
\hfuzz=500pt\quad\includegraphics[width=1em]{connector.pdf}\includegraphics[width=1em]{element.pdf}~applyWeight & \hfuzz=500pt boolean & \hfuzz=500pt from normal equations\\
\hfuzz=500pt\quad\includegraphics[width=1em]{connector.pdf}\includegraphics[width=1em]{element.pdf}~sigma & \hfuzz=500pt double & \hfuzz=500pt [m]\\
\hfuzz=500pt\includegraphics[width=1em]{connector.pdf}\includegraphics[width=1em]{element-mustset.pdf}~centerMean & \hfuzz=500pt sequence & \hfuzz=500pt zero center (x,y,z) as (weighted) mean of all patterns\\
\hfuzz=500pt\quad\includegraphics[width=1em]{connector.pdf}\includegraphics[width=1em]{element.pdf}~applyWeight & \hfuzz=500pt boolean & \hfuzz=500pt from normal equations\\
\hfuzz=500pt\quad\includegraphics[width=1em]{connector.pdf}\includegraphics[width=1em]{element.pdf}~sigma & \hfuzz=500pt double & \hfuzz=500pt [m]\\
\hfuzz=500pt\includegraphics[width=1em]{connector.pdf}\includegraphics[width=1em]{element-mustset.pdf}~constant & \hfuzz=500pt sequence & \hfuzz=500pt zero constant (mean of all directions) of a single pattern\\
\hfuzz=500pt\quad\includegraphics[width=1em]{connector.pdf}\includegraphics[width=1em]{element-mustset.pdf}~type & \hfuzz=500pt \hyperref[gnssType]{gnssType} & \hfuzz=500pt applied for each matching types\\
\hfuzz=500pt\quad\includegraphics[width=1em]{connector.pdf}\includegraphics[width=1em]{element.pdf}~applyWeight & \hfuzz=500pt boolean & \hfuzz=500pt from normal equations\\
\hfuzz=500pt\quad\includegraphics[width=1em]{connector.pdf}\includegraphics[width=1em]{element.pdf}~sigma & \hfuzz=500pt double & \hfuzz=500pt [m]\\
\hfuzz=500pt\includegraphics[width=1em]{connector.pdf}\includegraphics[width=1em]{element-mustset.pdf}~constantMean & \hfuzz=500pt sequence & \hfuzz=500pt zero constant (mean of all directions) as (weighted) mean of all patterns\\
\hfuzz=500pt\quad\includegraphics[width=1em]{connector.pdf}\includegraphics[width=1em]{element.pdf}~applyWeight & \hfuzz=500pt boolean & \hfuzz=500pt from normal equations\\
\hfuzz=500pt\quad\includegraphics[width=1em]{connector.pdf}\includegraphics[width=1em]{element.pdf}~sigma & \hfuzz=500pt double & \hfuzz=500pt [m]\\
\hfuzz=500pt\includegraphics[width=1em]{connector.pdf}\includegraphics[width=1em]{element-mustset.pdf}~TEC & \hfuzz=500pt sequence & \hfuzz=500pt zero TEC computed as (weighetd) least squares from all types\\
\hfuzz=500pt\quad\includegraphics[width=1em]{connector.pdf}\includegraphics[width=1em]{element.pdf}~type & \hfuzz=500pt \hyperref[gnssType]{gnssType} & \hfuzz=500pt applied for combination of matching types\\
\hfuzz=500pt\quad\includegraphics[width=1em]{connector.pdf}\includegraphics[width=1em]{element.pdf}~applyWeight & \hfuzz=500pt boolean & \hfuzz=500pt from normal equations\\
\hfuzz=500pt\quad\includegraphics[width=1em]{connector.pdf}\includegraphics[width=1em]{element.pdf}~sigma & \hfuzz=500pt double & \hfuzz=500pt [TECU]\\
\hfuzz=500pt\includegraphics[width=1em]{element-mustset-unbounded.pdf}~antennaCenterVariations & \hfuzz=500pt \hyperref[parametrizationGnssAntennaType]{parametrizationGnssAntenna} & \hfuzz=500pt \\
\hfuzz=500pt\includegraphics[width=1em]{element.pdf}~antennaName & \hfuzz=500pt string & \hfuzz=500pt apply constraints to all machting antennas\\
\hfuzz=500pt\includegraphics[width=1em]{element.pdf}~antennaSerial & \hfuzz=500pt string & \hfuzz=500pt apply constraints to all machting antennas\\
\hfuzz=500pt\includegraphics[width=1em]{element.pdf}~antennaRadome & \hfuzz=500pt string & \hfuzz=500pt apply constraints to all machting antennas\\
\hfuzz=500pt\includegraphics[width=1em]{element.pdf}~deltaAzimuth & \hfuzz=500pt angle & \hfuzz=500pt [degree] sampling of pattern to estimate center/constant\\
\hfuzz=500pt\includegraphics[width=1em]{element.pdf}~deltaZenith & \hfuzz=500pt angle & \hfuzz=500pt [degree] sampling of pattern to estimate center/constant\\
\hfuzz=500pt\includegraphics[width=1em]{element.pdf}~maxZenith & \hfuzz=500pt angle & \hfuzz=500pt [degree] sampling of pattern to estimate center/constant\\
\hline
\end{tabularx}

\clearpage
%==================================
\subsection{GnssAttitudeInfoCreate}\label{GnssAttitudeInfoCreate}
Creates attitude info file (\file{Instrument(MISCVALUES)}{instrument})
used by \program{SimulateStarCameraGnss}. One or more \config{attitudeInfo}s can be specified.
They are valid from \config{timeStart} until the start of the subsequent \config{attitudeInfo}.
\config{maxManeuverTime} is used by \program{SimulateStarCameraGnss} to look
for ongoing orbit maneuvers before/after the given orbit that might affect the attitude at
the beginning or end of a given orbit.

\fig{!hb}{0.9}{gnssAttitudeModes}{fig:gnssAttitudeModes2}{Overview of attitude modes used by GNSS satellites}

Here is a list of GNSS satellite types for which the attitude behavior is known and their
respective attitude modes and required parameters:
\begin{itemize}
\item \textbf{GPS-II/IIA} [1]
\begin{itemize}
  \item \config{defaultMode}: nominalYawSteering
  \item \config{midnightMode}: shadowMaxYawSteeringAndRecovery
  \item \config{noonMode}: catchUpYawSteering
  \item \config{maxYawRate}: 0.12~deg/s
  \item \config{yawBias}: 0.5~deg
  \item \config{maxManeuverTime}: 2~h
\end{itemize}
\item \textbf{GPS-IIR/IIR-M} [1]
\begin{itemize}
  \item \config{defaultMode}: nominalYawSteering
  \item \config{midnightMode}: catchUpYawSteering
  \item \config{noonMode}: catchUpYawSteering
  \item \config{maxYawRate}: 0.2~deg/s
  \item \config{maxManeuverTime}: 30~min
\end{itemize}
\item \textbf{GPS-IIF} [2]
\begin{itemize}
  \item \config{defaultMode}: nominalYawSteering
  \item \config{midnightMode}: shadowConstantYawSteering
  \item \config{noonMode}: catchUpYawSteering
  \item \config{maxYawRate}: 0.11~deg/s
  \item \config{yawBias}: -0.7~deg
  \item \config{maxManeuverTime}: 1.5~h
\end{itemize}
\item \textbf{GLO-M} [3]
\begin{itemize}
  \item \config{defaultMode}: nominalYawSteering
  \item \config{midnightMode}: shadowMaxYawSteeringAndStop
  \item \config{noonMode}: centeredMaxYawSteering
  \item \config{maxYawRate}: 0.25~deg/s
  \item \config{noonBetaThreshold}: 2~deg
  \item \config{maxManeuverTime}: 1.5~h
\end{itemize}
\item \textbf{GAL-1} [4]
\begin{itemize}
  \item \config{defaultMode}: nominalYawSteering
  \item \config{midnightMode}: smoothedYawSteering1
  \item \config{noonMode}: smoothedYawSteering1
  \item \config{maxManeuverTime}: 1.5~h
\end{itemize}
\item \textbf{GAL-2} [4]
\begin{itemize}
  \item \config{defaultMode}: nominalYawSteering
  \item \config{midnightMode}: smoothedYawSteering2
  \item \config{noonMode}: smoothedYawSteering2
  \item \config{midnightBetaThreshold}: 4.1~deg
  \item \config{noonBetaThreshold}: 4.1~deg
  \item \config{activationThreshold}: 10~deg
  \item \config{maxManeuverTime}: 5656~s
\end{itemize}
\item \textbf{BDS-2G/3G} [5, 6]
\begin{itemize}
  \item \config{defaultMode}: orbitNormal
  \item \config{midnightMode}: orbitNormal
  \item \config{noonMode}: orbitNormal
\end{itemize}
\item \textbf{BDS-2I} [5]
\begin{itemize}
  \item \config{defaultMode}: nominalYawSteering
  \item \config{midnightMode}: betaDependentOrbitNormal
  \item \config{noonMode}: betaDependentOrbitNormal
  \item \config{maxYawRate}: 0.085~deg/s
  \item \config{midnightBetaThreshold}: 4~deg
  \item \config{noonBetaThreshold}: 4~deg
  \item \config{activationThreshold}: 5~deg
  \item \config{maxManeuverTime}: 24~h
\end{itemize}
\item \textbf{BDS-2M} [5]
\begin{itemize}
  \item \config{defaultMode}: nominalYawSteering
  \item \config{midnightMode}: betaDependentOrbitNormal
  \item \config{noonMode}: betaDependentOrbitNormal
  \item \config{maxYawRate}: 0.159~deg/s
  \item \config{midnightBetaThreshold}: 4~deg
  \item \config{noonBetaThreshold}: 4~deg
  \item \config{activationThreshold}: 5~deg
  \item \config{maxManeuverTime}: 13~h
\end{itemize}
\item \textbf{BDS-3I/3SI} [6]
\begin{itemize}
  \item \config{defaultMode}: nominalYawSteering
  \item \config{midnightMode}: smoothedYawSteering2
  \item \config{noonMode}: smoothedYawSteering2
  \item \config{midnightBetaThreshold}: 3~deg
  \item \config{noonBetaThreshold}: 3~deg
  \item \config{activationThreshold}: 6~deg
  \item \config{maxManeuverTime}: 5740~s
\end{itemize}
\item \textbf{BDS-3M/3SM} [6]
\begin{itemize}
  \item \config{defaultMode}: nominalYawSteering
  \item \config{midnightMode}: smoothedYawSteering2
  \item \config{noonMode}: smoothedYawSteering2
  \item \config{midnightBetaThreshold}: 3~deg
  \item \config{noonBetaThreshold}: 3~deg
  \item \config{activationThreshold}: 6~deg
  \item \config{maxManeuverTime}: 3090~s
\end{itemize}
\item \textbf{QZS-1} [7]
\begin{itemize}
  \item \config{defaultMode}: nominalYawSteering
  \item \config{midnightMode}: betaDependentOrbitNormal
  \item \config{noonMode}: betaDependentOrbitNormal
  \item \config{maxYawRate}: 0.01~deg/s
  \item \config{yawBias}: 180~deg
  \item \config{midnightBetaThreshold}: 20~deg
  \item \config{noonBetaThreshold}: 20~deg
  \item \config{activationThreshold}: 18.5~deg
  \item \config{maxManeuverTime}: 24~h
\end{itemize}
\item \textbf{QZS-2G} [7]
\begin{itemize}
  \item \config{defaultMode}: orbitNormal
  \item \config{midnightMode}: orbitNormal
  \item \config{noonMode}: orbitNormal
  \item \config{yawBias}: 180~deg
\end{itemize}
\item \textbf{QZS-2I} [7]
\begin{itemize}
  \item \config{defaultMode}: nominalYawSteering
  \item \config{midnightMode}: centeredMaxYawSteering
  \item \config{noonMode}: centeredMaxYawSteering
  \item \config{maxYawRate}: 0.055~deg/s
  \item \config{midnightBetaThreshold}: 5~deg
  \item \config{noonBetaThreshold}: 5~deg
  \item \config{maxManeuverTime}: 1.5~h
\end{itemize}
\end{itemize}

Some specific satellites may deviate in their attitude behavior or parameters
(e.g. G013-G040, R713, C005, C015, C017, J001).

References for the attitude behavior information:
\begin{enumerate}
\item \href{https://doi.org/10.1007/s10291-008-0092-1}{Kouba (2009)}
\item \href{https://doi.org/10.1007/s10291-016-0562-9}{Kuang et al. (2017)}
\item \href{https://doi.org/10.1016/j.asr.2010.09.007}{Dilssner et al. (2011)}
\item \url{https://www.gsc-europa.eu/support-to-developers/galileo-satellite-metadata#3}
\item \href{https://doi.org/10.1007/s10291-018-0783-1}{Wang et al. (2018)}
\item \href{https://doi.org/10.1017/S0373463318000103}{Li et al. (2018)}
\item \url{https://qzss.go.jp/en/technical/qzssinfo/index.html}
\end{enumerate}


\keepXColumns
\begin{tabularx}{\textwidth}{N T A}
\hline
Name & Type & Annotation\\
\hline
\hfuzz=500pt\includegraphics[width=1em]{element-mustset.pdf}~outputfileAttitudeInfo & \hfuzz=500pt filename & \hfuzz=500pt \\
\hfuzz=500pt\includegraphics[width=1em]{element-mustset-unbounded.pdf}~attitudeInfo & \hfuzz=500pt sequence & \hfuzz=500pt \\
\hfuzz=500pt\includegraphics[width=1em]{connector.pdf}\includegraphics[width=1em]{element-mustset.pdf}~timeStart & \hfuzz=500pt time & \hfuzz=500pt \\
\hfuzz=500pt\includegraphics[width=1em]{connector.pdf}\includegraphics[width=1em]{element-mustset.pdf}~defaultMode & \hfuzz=500pt choice & \hfuzz=500pt default attitude mode\\
\hfuzz=500pt\quad\includegraphics[width=1em]{connector.pdf}\includegraphics[width=1em]{element-mustset.pdf}~nominalYawSteering & \hfuzz=500pt  & \hfuzz=500pt yaw to keep solar panels aligned to Sun (e.g. most GNSS satellites outside eclipse)\\
\hfuzz=500pt\quad\includegraphics[width=1em]{connector.pdf}\includegraphics[width=1em]{element-mustset.pdf}~orbitNormal & \hfuzz=500pt  & \hfuzz=500pt keep fixed yaw angle, for example point X-axis in flight direction (e.g. BDS-2G, BDS-3G, QZS-2G)\\
\hfuzz=500pt\includegraphics[width=1em]{connector.pdf}\includegraphics[width=1em]{element-mustset.pdf}~midnightMode & \hfuzz=500pt choice & \hfuzz=500pt attitude mode for maneuvers around orbit midnight\\
\hfuzz=500pt\quad\includegraphics[width=1em]{connector.pdf}\includegraphics[width=1em]{element-mustset.pdf}~nominalYawSteering & \hfuzz=500pt  & \hfuzz=500pt yaw to keep solar panels aligned to Sun (e.g. most GNSS satellites outside eclipse)\\
\hfuzz=500pt\quad\includegraphics[width=1em]{connector.pdf}\includegraphics[width=1em]{element-mustset.pdf}~orbitNormal & \hfuzz=500pt  & \hfuzz=500pt keep fixed yaw angle, for example point X-axis in flight direction (e.g. BDS-2G, BDS-3G, QZS-2G)\\
\hfuzz=500pt\quad\includegraphics[width=1em]{connector.pdf}\includegraphics[width=1em]{element-mustset.pdf}~catchUpYawSteering & \hfuzz=500pt  & \hfuzz=500pt yaw at maximum yaw rate to catch up to nominal yaw angle (e.g. GPS-* (noon), GPS-IIR (midnight))\\
\hfuzz=500pt\quad\includegraphics[width=1em]{connector.pdf}\includegraphics[width=1em]{element-mustset.pdf}~shadowMaxYawSteeringAndRecovery & \hfuzz=500pt  & \hfuzz=500pt yaw at maximum yaw rate from shadow start to end, recover after shadow (e.g. GPS-IIA (midnight))\\
\hfuzz=500pt\quad\includegraphics[width=1em]{connector.pdf}\includegraphics[width=1em]{element-mustset.pdf}~shadowMaxYawSteeringAndStop & \hfuzz=500pt  & \hfuzz=500pt yaw at maximum yaw rate from shadow start until nominal yaw angle at shadow end is reached, then stop (e.g. GLO-M (midnight))\\
\hfuzz=500pt\quad\includegraphics[width=1em]{connector.pdf}\includegraphics[width=1em]{element-mustset.pdf}~shadowConstantYawSteering & \hfuzz=500pt  & \hfuzz=500pt yaw at constant yaw rate from shadow start to end (e.g. GPS-IIF (midnight))\\
\hfuzz=500pt\quad\includegraphics[width=1em]{connector.pdf}\includegraphics[width=1em]{element-mustset.pdf}~centeredMaxYawSteering & \hfuzz=500pt  & \hfuzz=500pt yaw at maximum yaw rate centered around noon/midnight (e.g. QZS-2I, GLO-M (noon))\\
\hfuzz=500pt\quad\includegraphics[width=1em]{connector.pdf}\includegraphics[width=1em]{element-mustset.pdf}~smoothedYawSteering1 & \hfuzz=500pt  & \hfuzz=500pt yaw based on an auxiliary Sun vector for a smooth yaw maneuver (e.g. GAL-1)\\
\hfuzz=500pt\quad\includegraphics[width=1em]{connector.pdf}\includegraphics[width=1em]{element-mustset.pdf}~smoothedYawSteering2 & \hfuzz=500pt  & \hfuzz=500pt yaw based on a modified yaw-steering law for a smooth yaw maneuver (e.g. GAL-2, BDS-3M, BDS-3I)\\
\hfuzz=500pt\quad\includegraphics[width=1em]{connector.pdf}\includegraphics[width=1em]{element-mustset.pdf}~betaDependentOrbitNormal & \hfuzz=500pt  & \hfuzz=500pt switch to orbit normal mode if below beta angle threshold (e.g. BDS-2M, BDS-2I, QZS-1)\\
\hfuzz=500pt\includegraphics[width=1em]{connector.pdf}\includegraphics[width=1em]{element-mustset.pdf}~noonMode & \hfuzz=500pt choice & \hfuzz=500pt attitude mode for maneuvers around orbit noon\\
\hfuzz=500pt\quad\includegraphics[width=1em]{connector.pdf}\includegraphics[width=1em]{element-mustset.pdf}~nominalYawSteering & \hfuzz=500pt  & \hfuzz=500pt yaw to keep solar panels aligned to Sun (e.g. most GNSS satellites outside eclipse)\\
\hfuzz=500pt\quad\includegraphics[width=1em]{connector.pdf}\includegraphics[width=1em]{element-mustset.pdf}~orbitNormal & \hfuzz=500pt  & \hfuzz=500pt keep fixed yaw angle, for example point X-axis in flight direction (e.g. BDS-2G, BDS-3G, QZS-2G)\\
\hfuzz=500pt\quad\includegraphics[width=1em]{connector.pdf}\includegraphics[width=1em]{element-mustset.pdf}~catchUpYawSteering & \hfuzz=500pt  & \hfuzz=500pt yaw at maximum yaw rate to catch up to nominal yaw angle (e.g. GPS-* (noon), GPS-IIR (midnight))\\
\hfuzz=500pt\quad\includegraphics[width=1em]{connector.pdf}\includegraphics[width=1em]{element-mustset.pdf}~centeredMaxYawSteering & \hfuzz=500pt  & \hfuzz=500pt yaw at maximum yaw rate centered around noon/midnight (e.g. QZS-2I, GLO-M (noon))\\
\hfuzz=500pt\quad\includegraphics[width=1em]{connector.pdf}\includegraphics[width=1em]{element-mustset.pdf}~smoothedYawSteering1 & \hfuzz=500pt  & \hfuzz=500pt yaw based on an auxiliary Sun vector for a smooth yaw maneuver (e.g. GAL-1)\\
\hfuzz=500pt\quad\includegraphics[width=1em]{connector.pdf}\includegraphics[width=1em]{element-mustset.pdf}~smoothedYawSteering2 & \hfuzz=500pt  & \hfuzz=500pt yaw based on a modified yaw-steering law for a smooth yaw maneuver (e.g. GAL-2, BDS-3M, BDS-3I)\\
\hfuzz=500pt\quad\includegraphics[width=1em]{connector.pdf}\includegraphics[width=1em]{element-mustset.pdf}~betaDependentOrbitNormal & \hfuzz=500pt  & \hfuzz=500pt switch to orbit normal mode if below beta angle threshold (e.g. BDS-2M, BDS-2I, QZS-1)\\
\hfuzz=500pt\includegraphics[width=1em]{connector.pdf}\includegraphics[width=1em]{element.pdf}~maxYawRate & \hfuzz=500pt double & \hfuzz=500pt [degree/s] maximum yaw rate of the satellite\\
\hfuzz=500pt\includegraphics[width=1em]{connector.pdf}\includegraphics[width=1em]{element.pdf}~yawBias & \hfuzz=500pt double & \hfuzz=500pt [degree] yaw bias applied in satellite attitude control system\\
\hfuzz=500pt\includegraphics[width=1em]{connector.pdf}\includegraphics[width=1em]{element.pdf}~midnightBetaThreshold & \hfuzz=500pt double & \hfuzz=500pt [degree] limit midnight maneuver to this absolute angle of the Sun above/below the satellite orbital plane\\
\hfuzz=500pt\includegraphics[width=1em]{connector.pdf}\includegraphics[width=1em]{element.pdf}~noonBetaThreshold & \hfuzz=500pt double & \hfuzz=500pt [degree] limit noon maneuver to this absolute angle of the Sun above/below the satellite orbital plane\\
\hfuzz=500pt\includegraphics[width=1em]{connector.pdf}\includegraphics[width=1em]{element.pdf}~activationThreshold & \hfuzz=500pt double & \hfuzz=500pt [degree] limit maneuver to this yaw/Earth-spacecraft-Sun angle (depending on mode)\\
\hfuzz=500pt\includegraphics[width=1em]{connector.pdf}\includegraphics[width=1em]{element.pdf}~maxManeuverTime & \hfuzz=500pt double & \hfuzz=500pt [s] maximum duration of maneuver or maximum maneuver lookup time before/after orbit start/end\\
\hline
\end{tabularx}

\clearpage
%==================================
\subsection{GnssBiasClockAlignment}\label{GnssBiasClockAlignment}
This program can be used to absolutely align GNSS transmitter clocks to reference clocks (i.e. broadcast clocks).
Each 'group' of \config{transmitter}s, usually a system like GPS or Galileo, is aligned individually by a constant shift over all transmitters.
If \config{alignClocksByFreqNo} is set, GLONASS transmitters will be divided by frequency number into groups of nominally two transmitters.
The offset between clocks and reference clocks will be shifted into receiver code biases, if \config{receiver} is provided."

By setting \config{alignFreqNoBiasesAtReceiver} and providing \config{receiver}, this program can further align GLONASS transmitter signal
biases so that the differences between frequency number-dependent receiver signal biases are minimal, which helps if PPP users don't set
up individual signal biases per frequency number at the receiver. Alignment is done by computing signal bias residuals to the mean over all
frequency numbers of a signal type at each receiver and then computing the means over all receivers for each frequency number and shifting
those from the receiver signal biases to the transmitter signal biases. Internal consistency of the biases is not affected by this.

If you only want to align GLONASS frequency numbers, provide the same clocks in
\configFile{inputfileClock}{instrument} and \configFile{inputfileReferenceClock}{instrument}.


\keepXColumns
\begin{tabularx}{\textwidth}{N T A}
\hline
Name & Type & Annotation\\
\hline
\hfuzz=500pt\includegraphics[width=1em]{element-mustset-unbounded.pdf}~transmitter & \hfuzz=500pt sequence & \hfuzz=500pt one element per satellite\\
\hfuzz=500pt\includegraphics[width=1em]{connector.pdf}\includegraphics[width=1em]{element-mustset.pdf}~outputfileClock & \hfuzz=500pt filename & \hfuzz=500pt aligned clock instrument file\\
\hfuzz=500pt\includegraphics[width=1em]{connector.pdf}\includegraphics[width=1em]{element.pdf}~outputfileSignalBias & \hfuzz=500pt filename & \hfuzz=500pt (GLONASS only) aligned signal bias file\\
\hfuzz=500pt\includegraphics[width=1em]{connector.pdf}\includegraphics[width=1em]{element-mustset.pdf}~inputfileClock & \hfuzz=500pt filename & \hfuzz=500pt clock instrument file\\
\hfuzz=500pt\includegraphics[width=1em]{connector.pdf}\includegraphics[width=1em]{element-mustset.pdf}~inputfileReferenceClock & \hfuzz=500pt filename & \hfuzz=500pt reference clock instrument file\\
\hfuzz=500pt\includegraphics[width=1em]{connector.pdf}\includegraphics[width=1em]{element.pdf}~inputfileSignalBias & \hfuzz=500pt filename & \hfuzz=500pt (GLONASS only) signal bias file\\
\hfuzz=500pt\includegraphics[width=1em]{connector.pdf}\includegraphics[width=1em]{element-mustset.pdf}~inputfileTransmitterInfo & \hfuzz=500pt filename & \hfuzz=500pt transmitter platform file\\
\hfuzz=500pt\includegraphics[width=1em]{element-unbounded.pdf}~receiver & \hfuzz=500pt sequence & \hfuzz=500pt one element per station\\
\hfuzz=500pt\includegraphics[width=1em]{connector.pdf}\includegraphics[width=1em]{element-mustset.pdf}~outputfileSignalBias & \hfuzz=500pt filename & \hfuzz=500pt aligned signal bias file\\
\hfuzz=500pt\includegraphics[width=1em]{connector.pdf}\includegraphics[width=1em]{element-mustset.pdf}~inputfileSignalBias & \hfuzz=500pt filename & \hfuzz=500pt signal bias file\\
\hfuzz=500pt\includegraphics[width=1em]{element.pdf}~alignClocksByFreqNo & \hfuzz=500pt boolean & \hfuzz=500pt align clocks for each GLONASS frequency number separately\\
\hfuzz=500pt\includegraphics[width=1em]{element.pdf}~alignFreqNoBiasesAtReceiver & \hfuzz=500pt boolean & \hfuzz=500pt align frequency number-dependent code biases for each receiver\\
\hline
\end{tabularx}

\clearpage
%==================================
\subsection{GnssEstimateClockShift}\label{GnssEstimateClockShift}
This program estimates an epoch-wise clock shift in a constellation of GNSS satellites.
Each separate \config{data} represents a satellite... (e.g. 32 GPS satellites).
The shift to reference clocks can be estimated by providing \configFile{inputfileInstrumentRef}{instrument}.
Clock shifts are estimated for each epoch given by \configClass{timeSeries}{timeSeriesType}.


\keepXColumns
\begin{tabularx}{\textwidth}{N T A}
\hline
Name & Type & Annotation\\
\hline
\hfuzz=500pt\includegraphics[width=1em]{element.pdf}~outputfileShiftTimeSeries & \hfuzz=500pt filename & \hfuzz=500pt columns: mjd, clock shift\\
\hfuzz=500pt\includegraphics[width=1em]{element-mustset-unbounded.pdf}~data & \hfuzz=500pt sequence & \hfuzz=500pt e.g. satellite\\
\hfuzz=500pt\includegraphics[width=1em]{connector.pdf}\includegraphics[width=1em]{element.pdf}~outputfileInstrument & \hfuzz=500pt filename & \hfuzz=500pt corrected clocks\\
\hfuzz=500pt\includegraphics[width=1em]{connector.pdf}\includegraphics[width=1em]{element.pdf}~outputfileInstrumentDiff & \hfuzz=500pt filename & \hfuzz=500pt clock difference after correction\\
\hfuzz=500pt\includegraphics[width=1em]{connector.pdf}\includegraphics[width=1em]{element-mustset.pdf}~inputfileInstrument & \hfuzz=500pt filename & \hfuzz=500pt input clocks\\
\hfuzz=500pt\includegraphics[width=1em]{connector.pdf}\includegraphics[width=1em]{element.pdf}~inputfileInstrumentRef & \hfuzz=500pt filename & \hfuzz=500pt reference clocks (subtracted from input clocks)\\
\hfuzz=500pt\includegraphics[width=1em]{element-mustset-unbounded.pdf}~timeSeries & \hfuzz=500pt \hyperref[timeSeriesType]{timeSeries} & \hfuzz=500pt clock epochs\\
\hfuzz=500pt\includegraphics[width=1em]{element.pdf}~margin & \hfuzz=500pt double & \hfuzz=500pt [s] margin for time comparison\\
\hline
\end{tabularx}

\clearpage
%==================================
\subsection{GnssGlonassFrequencyNumberUpdate}\label{GnssGlonassFrequencyNumberUpdate}
Update/set GLONASS frequency number in \configFile{inputfileTransmitterInfo}{platform} files.

PRN/SVN to frequency number source: \url{http://semisys.gfz-potsdam.de/semisys/api/?symname=2002&format=json&satellite=GLO}.

See also \program{GnssAntex2AntennaDefinition}.


\keepXColumns
\begin{tabularx}{\textwidth}{N T A}
\hline
Name & Type & Annotation\\
\hline
\hfuzz=500pt\includegraphics[width=1em]{element.pdf}~outputfileTransmitterInfo & \hfuzz=500pt filename & \hfuzz=500pt templated for PRN list (variableNamePrn)\\
\hfuzz=500pt\includegraphics[width=1em]{element-mustset.pdf}~inputfileTransmitterInfo & \hfuzz=500pt filename & \hfuzz=500pt templated for PRN list (variableNamePrn)\\
\hfuzz=500pt\includegraphics[width=1em]{element-mustset.pdf}~inputfilePrn2FrequencyNumber & \hfuzz=500pt filename & \hfuzz=500pt GROOPS matrix with columns: GLONASS PRN, SVN, mjdStart, mjdEnd, frequencyNumber\\
\hfuzz=500pt\includegraphics[width=1em]{element-unbounded.pdf}~prn & \hfuzz=500pt string & \hfuzz=500pt PRN (e.g. R01) for transmitter info files\\
\hfuzz=500pt\includegraphics[width=1em]{element.pdf}~variableNamePrn & \hfuzz=500pt string & \hfuzz=500pt variable name for PRN in transmitter info files\\
\hline
\end{tabularx}

\clearpage
%==================================
\subsection{GnssProcessing}\label{GnssProcessing}
This program processes GNSS observations. It calculates the linearized observation equations,
accumulates them into a system of normal equations and solves it.

The primary use cases of this program are:
\begin{itemize}
  \item \reference{GNSS satellite orbit determination and station network analysis}{cookbook.gnssNetwork}
  \item \reference{Kinematic orbit determination of LEO satellites}{cookbook.kinematicOrbit}
  \item \reference{GNSS precise point positioning (PPP)}{cookbook.gnssPpp}
\end{itemize}

The observation epochs are defined by \configClass{timeSeries}{timeSeriesType}
and only observations at these epochs (within a \config{timeMargin}) are considered.

To calculate observation equations from the tracks, the model parameters or unknown parameters need to be
defined beforehand. These unknown parameters can be chosen arbitrarily by the user with an adequate list of defined
\configClass{parametrization}{gnssParametrizationType}.
Some of the \configClass{parametrization}{gnssParametrizationType} also include a priori models.

Lastly it is required to define the process flow of the gnssProcessing. This is accomplished
with a list of \configClass{processingSteps}{gnssProcessingStepType}.
Each step is processed consecutively. Some steps allow the selection of parameters, epochs,
or the normal equation structure, which affects all subsequent steps.
A minimal example consists of following steps:
\begin{itemize}
  \item \configClass{estimate}{gnssProcessingStepType:estimate}: iterative float solution with outlier downeighting
  \item \configClass{resolveAmbiguities}{gnssProcessingStepType:resolveAmbiguities}:
        fix ambiguities to integer and remove them from the normals
  \item \configClass{estimate}{gnssProcessingStepType:estimate}: few iteration for final outlier downweighting
  \item \configClass{writeResults}{gnssProcessingStepType:writeResults}:
        write the output files defined in \configClass{parametrization}{gnssParametrizationType}
\end{itemize}

If the program is run on multiple processes the \configClass{receiver}{gnssReceiverGeneratorType}s
(stations or LEO satellites) are distributed over the processes.

See also \program{GnssSimulateReceiver}.


\keepXColumns
\begin{tabularx}{\textwidth}{N T A}
\hline
Name & Type & Annotation\\
\hline
\hfuzz=500pt\includegraphics[width=1em]{element-mustset-unbounded.pdf}~timeSeries & \hfuzz=500pt \hyperref[timeSeriesType]{timeSeries} & \hfuzz=500pt defines observation epochs\\
\hfuzz=500pt\includegraphics[width=1em]{element.pdf}~timeMargin & \hfuzz=500pt double & \hfuzz=500pt [seconds] margin to consider two times identical\\
\hfuzz=500pt\includegraphics[width=1em]{element-mustset-unbounded.pdf}~transmitter & \hfuzz=500pt \hyperref[gnssTransmitterGeneratorType]{gnssTransmitterGenerator} & \hfuzz=500pt constellation of GNSS satellites\\
\hfuzz=500pt\includegraphics[width=1em]{element-mustset-unbounded.pdf}~receiver & \hfuzz=500pt \hyperref[gnssReceiverGeneratorType]{gnssReceiverGenerator} & \hfuzz=500pt ground station network or LEO satellite\\
\hfuzz=500pt\includegraphics[width=1em]{element-mustset.pdf}~earthRotation & \hfuzz=500pt \hyperref[earthRotationType]{earthRotation} & \hfuzz=500pt apriori earth rotation\\
\hfuzz=500pt\includegraphics[width=1em]{element-mustset-unbounded.pdf}~parametrization & \hfuzz=500pt \hyperref[gnssParametrizationType]{gnssParametrization} & \hfuzz=500pt models and parameters\\
\hfuzz=500pt\includegraphics[width=1em]{element-mustset-unbounded.pdf}~processingStep & \hfuzz=500pt \hyperref[gnssProcessingStepType]{gnssProcessingStep} & \hfuzz=500pt steps are processed consecutively\\
\hline
\end{tabularx}

This program is \reference{parallelized}{general.parallelization}.
\clearpage
%==================================
\subsection{GnssReceiverDefinitionCreate}\label{GnssReceiverDefinitionCreate}
Create a \file{GNSS receiver definition file}{gnssReceiverDefinition}.


\keepXColumns
\begin{tabularx}{\textwidth}{N T A}
\hline
Name & Type & Annotation\\
\hline
\hfuzz=500pt\includegraphics[width=1em]{element-mustset.pdf}~outputfileGnssReceiverDefinition & \hfuzz=500pt filename & \hfuzz=500pt \\
\hfuzz=500pt\includegraphics[width=1em]{element-mustset-unbounded.pdf}~receiverDefinition & \hfuzz=500pt sequence & \hfuzz=500pt \\
\hfuzz=500pt\includegraphics[width=1em]{connector.pdf}\includegraphics[width=1em]{element-mustset.pdf}~name & \hfuzz=500pt string & \hfuzz=500pt \\
\hfuzz=500pt\includegraphics[width=1em]{connector.pdf}\includegraphics[width=1em]{element.pdf}~serial & \hfuzz=500pt string & \hfuzz=500pt \\
\hfuzz=500pt\includegraphics[width=1em]{connector.pdf}\includegraphics[width=1em]{element.pdf}~version & \hfuzz=500pt string & \hfuzz=500pt \\
\hfuzz=500pt\includegraphics[width=1em]{connector.pdf}\includegraphics[width=1em]{element.pdf}~comment & \hfuzz=500pt string & \hfuzz=500pt \\
\hfuzz=500pt\includegraphics[width=1em]{connector.pdf}\includegraphics[width=1em]{element-unbounded.pdf}~gnssType & \hfuzz=500pt \hyperref[gnssType]{gnssType} & \hfuzz=500pt \\
\hline
\end{tabularx}

\clearpage
%==================================
\subsection{GnssResiduals2AccuracyDefinition}\label{GnssResiduals2AccuracyDefinition}
Compute antenna accuracies from observation \configFile{inputfileResiduals}{instrument}.
The \configFile{inputfileStationInfo}{platform} is needed to assign
the residuals to the equipped antenna at observation times.

The \configFile{outputfileAccuracyDefinition}{gnssAntennaDefinition} contains
at first step the same accuracy information for all antennas as the input file.
Only the azimuth~$A$ and elevation~$E$ dependent grid points of the patterns
where enough residuals are available ($>$ \config{minRedundancy})
are replaced by estimated accuracy
\begin{equation}
 \sigma(A,E) = \sqrt{\frac{\sum_i e_i^2(A,E)}{\sum_i r_i(A,E)}},
\end{equation}
where $e_i$ are the azimuth and elevation dependent residuals and $r_i$ the
corresponding redundancies (number of observations minus the contribution to
the estimated parameters).

The \configFile{inputfileAccuracyDefinition}{gnssAntennaDefinition} can be modified
to the demands before with \program{GnssAntennaDefinitionCreate}
(e.g. with \config{antenna:resample}).

To verify the results the \configFile{outputfileAntennaMean}{gnssAntennaDefinition}
and the accumulated \configFile{outputfileAntennaRedundancy}{gnssAntennaDefinition}
of the computed pattern grid points can be written.

Example: Analysis of TerraSAR-X residuals of one month shows that low elevation
GPS satellites are not tracked by the onboard receiver. An estimation of accuracies
for these directions is not possible from the residuals and the apriori accuracies
are left untouched. The other directions show very low phase noise hardly elevation
and azimuth dependent for L2W. A nearly zero mean indicates the use of adequate antennca
center variations in the processing.

\fig{!hb}{0.8}{gnssResiduals2AccuracyDefinition}{fig:gnssResiduals2AccuracyDefinition}{L2W accuracies of TerraSAR-X determined from residuals of one month}

See also \program{GnssResiduals2TransmitterAccuracyDefinition}.


\keepXColumns
\begin{tabularx}{\textwidth}{N T A}
\hline
Name & Type & Annotation\\
\hline
\hfuzz=500pt\includegraphics[width=1em]{element.pdf}~outputfileAccuracyDefinition & \hfuzz=500pt filename & \hfuzz=500pt elevation and azimuth dependent accuracy\\
\hfuzz=500pt\includegraphics[width=1em]{element.pdf}~outputfileAntennaMean & \hfuzz=500pt filename & \hfuzz=500pt weighted mean of the residuals\\
\hfuzz=500pt\includegraphics[width=1em]{element.pdf}~outputfileAntennaRedundancy & \hfuzz=500pt filename & \hfuzz=500pt redundancy of adjustment\\
\hfuzz=500pt\includegraphics[width=1em]{element-mustset.pdf}~inputfileAccuracyDefinition & \hfuzz=500pt filename & \hfuzz=500pt apriori accuracies\\
\hfuzz=500pt\includegraphics[width=1em]{element-mustset.pdf}~inputfileStationInfo & \hfuzz=500pt filename & \hfuzz=500pt to assign residuals to antennas\\
\hfuzz=500pt\includegraphics[width=1em]{element.pdf}~isTransmitter & \hfuzz=500pt boolean & \hfuzz=500pt stationInfo is of a transmitter\\
\hfuzz=500pt\includegraphics[width=1em]{element.pdf}~thresholdOutlier & \hfuzz=500pt double & \hfuzz=500pt ignore residuals with sigma/sigma0 greater than threshold\\
\hfuzz=500pt\includegraphics[width=1em]{element.pdf}~minRedundancy & \hfuzz=500pt double & \hfuzz=500pt min number of residuals. to estimate sigma\\
\hfuzz=500pt\includegraphics[width=1em]{element-mustset-unbounded.pdf}~inputfileResiduals & \hfuzz=500pt filename & \hfuzz=500pt GNSS receiver residuals\\
\hline
\end{tabularx}

\clearpage
%==================================
\subsection{GnssResiduals2Skyplot}\label{GnssResiduals2Skyplot}
Write GNSS residuals together with azimuth and elevation to be plotted with \program{PlotMap}.
Azimuth and elevation are written as ellipsoidal longitude and latitude in a \file{griddedData file}{griddedData}.
The choosen ellipsoid parameters \config{R} and \config{inverseFlattening} are arbitrary but should be the same
as in \program{PlotMap}. If with \configClass{typeTransmitter}{gnssType} (e.g. '\verb|***G18|')
a single transmitter is selected the azimuth and elevation are computed from the transmitter point of view.

For each GNSS \configClass{type}{gnssType} an extra data column is created.

A \file{GNSS residual file}{instrument} includes additional information
besides the residuals, which can also be selected with \configClass{type}{gnssType}
\begin{itemize}
\item \verb|A1*|, \verb|E1*|: azimuth and elevation at receiver
\item \verb|A2*|, \verb|E2*|: azimuth and elevation at transmitter
\item \verb|I**|: Estimated slant total electron content (STEC)
\end{itemize}

Furthermore these files may include for each residual \configClass{type}{gnssType}
information about the redundancy and the accuracy relation $\sigma/\sigma_0$
of the estimated $\sigma$ versus the apriori $\sigma_0$ from the least squares adjustment.
The 3 values (residuals, redundancy, $\sigma/\sigma_0$) are coded with the same type.
To get access to all values the corresponding type must be repeated in \configClass{type}{gnssType}.

\fig{!hb}{0.5}{gnssResiduals2Skyplot}{fig:gnssResiduals2Skyplot}{GPS C2W residuals of GRAZ station at 2012-01-01}


\keepXColumns
\begin{tabularx}{\textwidth}{N T A}
\hline
Name & Type & Annotation\\
\hline
\hfuzz=500pt\includegraphics[width=1em]{element-mustset.pdf}~outputfileGriddedData & \hfuzz=500pt filename & \hfuzz=500pt \\
\hfuzz=500pt\includegraphics[width=1em]{element-mustset-unbounded.pdf}~type & \hfuzz=500pt \hyperref[gnssType]{gnssType} & \hfuzz=500pt \\
\hfuzz=500pt\includegraphics[width=1em]{element.pdf}~typeTransmitter & \hfuzz=500pt \hyperref[gnssType]{gnssType} & \hfuzz=500pt choose transmitter view, e.g. '***G18'\\
\hfuzz=500pt\includegraphics[width=1em]{element-mustset-unbounded.pdf}~inputfileResiduals & \hfuzz=500pt filename & \hfuzz=500pt GNSS receiver residuals\\
\hfuzz=500pt\includegraphics[width=1em]{element.pdf}~R & \hfuzz=500pt double & \hfuzz=500pt reference radius for ellipsoidal coordinates\\
\hfuzz=500pt\includegraphics[width=1em]{element.pdf}~inverseFlattening & \hfuzz=500pt double & \hfuzz=500pt reference flattening for ellipsoidal coordinates\\
\hline
\end{tabularx}

\clearpage
%==================================
\subsection{GnssResiduals2TransmitterAccuracyDefinition}\label{GnssResiduals2TransmitterAccuracyDefinition}
Compute antenna accuracies from observation \configFile{inputfileResiduals}{instrument}.
The \configFile{inputfileTransmitterInfo}{platform} is needed to assign
the residuals to the equipped antenna at observation times.

The \configFile{outputfileAccuracyDefinition}{gnssAntennaDefinition} contains
at first step the same accuracy information for all antennas as the input file.
Only the azimuth~$A$ and elevation~$E$ dependent grid points of the patterns
where enough residuals are available ($>$ \config{minRedundancy})
are replaced by estimated accuracy
\begin{equation}
 \sigma(A,E) = \sqrt{\frac{\sum_i e_i^2(A,E)}{\sum_i r_i(A,E)}},
\end{equation}
where $e_i$ are the azimuth and elevation dependent residuals and $r_i$ the
corresponding redundancies (number of observations minus the contribution to
the estimated parameters).

The \configFile{inputfileAccuracyDefinition}{gnssAntennaDefinition} can be modified
to the demands before with \program{GnssAntennaDefinitionCreate}
(e.g. with \config{antenna:resample}).

To verify the results the \configFile{outputfileAntennaMean}{gnssAntennaDefinition}
and the accumulated \configFile{outputfileAntennaRedundancy}{gnssAntennaDefinition}
of the computed pattern grid points can be written.

See also \program{GnssResiduals2AccuracyDefinition}.


\keepXColumns
\begin{tabularx}{\textwidth}{N T A}
\hline
Name & Type & Annotation\\
\hline
\hfuzz=500pt\includegraphics[width=1em]{element.pdf}~outputfileAccuracyDefinition & \hfuzz=500pt filename & \hfuzz=500pt elevation and azimuth dependent accuracy\\
\hfuzz=500pt\includegraphics[width=1em]{element.pdf}~outputfileAntennaMean & \hfuzz=500pt filename & \hfuzz=500pt weighted mean of the residuals\\
\hfuzz=500pt\includegraphics[width=1em]{element.pdf}~outputfileAntennaRedundancy & \hfuzz=500pt filename & \hfuzz=500pt redundancy of adjustment\\
\hfuzz=500pt\includegraphics[width=1em]{element-mustset.pdf}~inputfileAccuracyDefinition & \hfuzz=500pt filename & \hfuzz=500pt apriori accuracies\\
\hfuzz=500pt\includegraphics[width=1em]{element-mustset.pdf}~inputfileTransmitterInfo & \hfuzz=500pt filename & \hfuzz=500pt to assign residuals to antennas\\
\hfuzz=500pt\includegraphics[width=1em]{element.pdf}~minRedundancy & \hfuzz=500pt double & \hfuzz=500pt min number of residuals. to estimate sigma\\
\hfuzz=500pt\includegraphics[width=1em]{element-mustset-unbounded.pdf}~inputfileResiduals & \hfuzz=500pt filename & \hfuzz=500pt GNSS receiver residuals\\
\hline
\end{tabularx}

\clearpage
%==================================
\subsection{GnssSignalBias2Matrix}\label{GnssSignalBias2Matrix}
Computes signal biases for a given list of \configClass{types}{gnssType}.
If the type list is empty, all types contained in \configFile{inputfileSignalBias}{gnssSignalBias} are used.
The resulting \configFile{outputfileMatrix}{matrix} contains a vector with an entry for each type.


\keepXColumns
\begin{tabularx}{\textwidth}{N T A}
\hline
Name & Type & Annotation\\
\hline
\hfuzz=500pt\includegraphics[width=1em]{element-mustset.pdf}~outputfileMatrix & \hfuzz=500pt filename & \hfuzz=500pt \\
\hfuzz=500pt\includegraphics[width=1em]{element.pdf}~outputfileTypes & \hfuzz=500pt filename & \hfuzz=500pt ASCII list of types\\
\hfuzz=500pt\includegraphics[width=1em]{element-mustset.pdf}~inputfileSignalBias & \hfuzz=500pt filename & \hfuzz=500pt \\
\hfuzz=500pt\includegraphics[width=1em]{element-unbounded.pdf}~types & \hfuzz=500pt \hyperref[gnssType]{gnssType} & \hfuzz=500pt if not set, all types in the file are used\\
\hline
\end{tabularx}

\clearpage
%==================================
\subsection{GnssSimulateReceiver}\label{GnssSimulateReceiver}
This program simulates observations from receivers to GNSS satellites.
These simulated observations can then be used in \program{GnssProcessing}, for example to conduct closed-loop simulations.

One or more GNSS constellations must be defined via \configClass{transmitter}{gnssTransmitterGeneratorType}.
Receivers such as ground station networks or Low Earth Orbit (LEO) satellites can be defined via \configClass{receiver}{gnssReceiverGeneratorType}.

If multiple receivers defined an \configFile{outputfileGnssReceiver}{instrument} and \configFile{outputfileClock}{instrument}
are written for each single receiver with the \reference{variable}{general.parser} \verb|{station}| being replaced by the receiver name.

A list of simulated observation types can be defined via \configClass{observationType}{gnssType}. Noise can be added to both observations and clock errors
via \configClass{noiseObervation}{noiseGeneratorType} and \configClass{noiseClockReceiver}{noiseGeneratorType}, respectively. Observation noise is
interpreted as a factor that is multiplied to the accuracy derived from the accuracy pattern of the respective observation type
(see \configFile{inputfileAccuracyDefinition}{gnssAntennaDefinition} in \configClass{receiver}{gnssReceiverGeneratorType}).

The \configClass{parametrization}{gnssParametrizationType} are used to simulate a priori models (e.g. troposphere, signal biases).
Parameter settings and outputfiles are ignored.

If the program is run on multiple processes the \configClass{receiver}{gnssReceiverGeneratorType}s
(stations or LEO satellites) are distributed over the processes.


\keepXColumns
\begin{tabularx}{\textwidth}{N T A}
\hline
Name & Type & Annotation\\
\hline
\hfuzz=500pt\includegraphics[width=1em]{element-mustset.pdf}~outputfileGnssReceiver & \hfuzz=500pt filename & \hfuzz=500pt variable \{station\} available, simulated observations\\
\hfuzz=500pt\includegraphics[width=1em]{element.pdf}~outputfileClock & \hfuzz=500pt filename & \hfuzz=500pt variable \{station\} available, simulated receiver clock errors\\
\hfuzz=500pt\includegraphics[width=1em]{element-mustset-unbounded.pdf}~timeSeries & \hfuzz=500pt \hyperref[timeSeriesType]{timeSeries} & \hfuzz=500pt defines observation epochs\\
\hfuzz=500pt\includegraphics[width=1em]{element.pdf}~timeMargin & \hfuzz=500pt double & \hfuzz=500pt [seconds] margin to consider two times identical\\
\hfuzz=500pt\includegraphics[width=1em]{element-mustset-unbounded.pdf}~transmitter & \hfuzz=500pt \hyperref[gnssTransmitterGeneratorType]{gnssTransmitterGenerator} & \hfuzz=500pt constellation of GNSS satellites\\
\hfuzz=500pt\includegraphics[width=1em]{element-mustset-unbounded.pdf}~receiver & \hfuzz=500pt \hyperref[gnssReceiverGeneratorType]{gnssReceiverGenerator} & \hfuzz=500pt ground station network or LEO satellite\\
\hfuzz=500pt\includegraphics[width=1em]{element-mustset.pdf}~earthRotation & \hfuzz=500pt \hyperref[earthRotationType]{earthRotation} & \hfuzz=500pt apriori earth rotation\\
\hfuzz=500pt\includegraphics[width=1em]{element-unbounded.pdf}~parametrization & \hfuzz=500pt \hyperref[gnssParametrizationType]{gnssParametrization} & \hfuzz=500pt models and parameters\\
\hfuzz=500pt\includegraphics[width=1em]{element-mustset-unbounded.pdf}~observationType & \hfuzz=500pt \hyperref[gnssType]{gnssType} & \hfuzz=500pt simulated observation types\\
\hfuzz=500pt\includegraphics[width=1em]{element-unbounded.pdf}~noiseObservation & \hfuzz=500pt \hyperref[noiseGeneratorType]{noiseGenerator} & \hfuzz=500pt [-] noise is multiplied with type accuracy pattern of receiver\\
\hfuzz=500pt\includegraphics[width=1em]{element-unbounded.pdf}~noiseClockReceiver & \hfuzz=500pt \hyperref[noiseGeneratorType]{noiseGenerator} & \hfuzz=500pt [m] noise added to the simulated receiver clock\\
\hline
\end{tabularx}

This program is \reference{parallelized}{general.parallelization}.
\clearpage
%==================================
\subsection{GnssStationInfoCreate}\label{GnssStationInfoCreate}
Create a \file{GnssStationInfo file}{gnssStationInfo} from scratch by defining attributes such as
\config{markerName}, \config{markerNumber}, \config{comment}, \config{approxPosition},
\config{antenna} and \config{receiver}.

See also \program{GnssAntex2AntennaDefinition} and \program{GnssStationLog2Platform}.


\keepXColumns
\begin{tabularx}{\textwidth}{N T A}
\hline
Name & Type & Annotation\\
\hline
\hfuzz=500pt\includegraphics[width=1em]{element-mustset.pdf}~outputfileStationInfo & \hfuzz=500pt filename & \hfuzz=500pt \\
\hfuzz=500pt\includegraphics[width=1em]{element-mustset.pdf}~markerName & \hfuzz=500pt string & \hfuzz=500pt \\
\hfuzz=500pt\includegraphics[width=1em]{element.pdf}~markerNumber & \hfuzz=500pt string & \hfuzz=500pt \\
\hfuzz=500pt\includegraphics[width=1em]{element.pdf}~comment & \hfuzz=500pt string & \hfuzz=500pt \\
\hfuzz=500pt\includegraphics[width=1em]{element.pdf}~approxPositionX & \hfuzz=500pt double & \hfuzz=500pt [m] in TRF\\
\hfuzz=500pt\includegraphics[width=1em]{element.pdf}~approxPositionY & \hfuzz=500pt double & \hfuzz=500pt [m] in TRF\\
\hfuzz=500pt\includegraphics[width=1em]{element.pdf}~approxPositionZ & \hfuzz=500pt double & \hfuzz=500pt [m] in TRF\\
\hfuzz=500pt\includegraphics[width=1em]{element-mustset-unbounded.pdf}~antenna & \hfuzz=500pt sequence & \hfuzz=500pt \\
\hfuzz=500pt\includegraphics[width=1em]{connector.pdf}\includegraphics[width=1em]{element-mustset.pdf}~name & \hfuzz=500pt string & \hfuzz=500pt \\
\hfuzz=500pt\includegraphics[width=1em]{connector.pdf}\includegraphics[width=1em]{element.pdf}~serial & \hfuzz=500pt string & \hfuzz=500pt \\
\hfuzz=500pt\includegraphics[width=1em]{connector.pdf}\includegraphics[width=1em]{element.pdf}~radome & \hfuzz=500pt string & \hfuzz=500pt \\
\hfuzz=500pt\includegraphics[width=1em]{connector.pdf}\includegraphics[width=1em]{element.pdf}~comment & \hfuzz=500pt string & \hfuzz=500pt \\
\hfuzz=500pt\includegraphics[width=1em]{connector.pdf}\includegraphics[width=1em]{element.pdf}~timeStart & \hfuzz=500pt time & \hfuzz=500pt \\
\hfuzz=500pt\includegraphics[width=1em]{connector.pdf}\includegraphics[width=1em]{element.pdf}~timeEnd & \hfuzz=500pt time & \hfuzz=500pt \\
\hfuzz=500pt\includegraphics[width=1em]{connector.pdf}\includegraphics[width=1em]{element-mustset.pdf}~positionX & \hfuzz=500pt double & \hfuzz=500pt [m] ARP in north, east, up or vehicle system\\
\hfuzz=500pt\includegraphics[width=1em]{connector.pdf}\includegraphics[width=1em]{element-mustset.pdf}~positionY & \hfuzz=500pt double & \hfuzz=500pt [m] ARP in north, east, up or vehicle system\\
\hfuzz=500pt\includegraphics[width=1em]{connector.pdf}\includegraphics[width=1em]{element-mustset.pdf}~positionZ & \hfuzz=500pt double & \hfuzz=500pt [m] ARP in north, east, up or vehicle system\\
\hfuzz=500pt\includegraphics[width=1em]{connector.pdf}\includegraphics[width=1em]{element.pdf}~rotationX & \hfuzz=500pt angle & \hfuzz=500pt [degree] from local/vehicle to left-handed antenna system\\
\hfuzz=500pt\includegraphics[width=1em]{connector.pdf}\includegraphics[width=1em]{element.pdf}~rotationY & \hfuzz=500pt angle & \hfuzz=500pt [degree] from local/vehicle to left-handed antenna system\\
\hfuzz=500pt\includegraphics[width=1em]{connector.pdf}\includegraphics[width=1em]{element.pdf}~rotationZ & \hfuzz=500pt angle & \hfuzz=500pt [degree] from local/vehicle to left-handed antenna system\\
\hfuzz=500pt\includegraphics[width=1em]{connector.pdf}\includegraphics[width=1em]{element.pdf}~flipX & \hfuzz=500pt boolean & \hfuzz=500pt flip x-axis (after rotation)\\
\hfuzz=500pt\includegraphics[width=1em]{connector.pdf}\includegraphics[width=1em]{element.pdf}~flipY & \hfuzz=500pt boolean & \hfuzz=500pt flip y-axis (after rotation)\\
\hfuzz=500pt\includegraphics[width=1em]{connector.pdf}\includegraphics[width=1em]{element.pdf}~flipZ & \hfuzz=500pt boolean & \hfuzz=500pt flip z-axis (after rotation)\\
\hfuzz=500pt\includegraphics[width=1em]{element-unbounded.pdf}~receiver & \hfuzz=500pt sequence & \hfuzz=500pt \\
\hfuzz=500pt\includegraphics[width=1em]{connector.pdf}\includegraphics[width=1em]{element-mustset.pdf}~name & \hfuzz=500pt string & \hfuzz=500pt \\
\hfuzz=500pt\includegraphics[width=1em]{connector.pdf}\includegraphics[width=1em]{element.pdf}~serial & \hfuzz=500pt string & \hfuzz=500pt \\
\hfuzz=500pt\includegraphics[width=1em]{connector.pdf}\includegraphics[width=1em]{element.pdf}~version & \hfuzz=500pt string & \hfuzz=500pt \\
\hfuzz=500pt\includegraphics[width=1em]{connector.pdf}\includegraphics[width=1em]{element.pdf}~comment & \hfuzz=500pt string & \hfuzz=500pt \\
\hfuzz=500pt\includegraphics[width=1em]{connector.pdf}\includegraphics[width=1em]{element.pdf}~timeStart & \hfuzz=500pt time & \hfuzz=500pt \\
\hfuzz=500pt\includegraphics[width=1em]{connector.pdf}\includegraphics[width=1em]{element.pdf}~timeEnd & \hfuzz=500pt time & \hfuzz=500pt \\
\hfuzz=500pt\includegraphics[width=1em]{element-unbounded.pdf}~referencePoint & \hfuzz=500pt sequence & \hfuzz=500pt e.g. center of mass in satellite frame\\
\hfuzz=500pt\includegraphics[width=1em]{connector.pdf}\includegraphics[width=1em]{element.pdf}~comment & \hfuzz=500pt string & \hfuzz=500pt \\
\hfuzz=500pt\includegraphics[width=1em]{connector.pdf}\includegraphics[width=1em]{element-mustset.pdf}~xStart & \hfuzz=500pt double & \hfuzz=500pt [m] in north, east, up or vehicle system\\
\hfuzz=500pt\includegraphics[width=1em]{connector.pdf}\includegraphics[width=1em]{element-mustset.pdf}~yStart & \hfuzz=500pt double & \hfuzz=500pt linear motion between start and end\\
\hfuzz=500pt\includegraphics[width=1em]{connector.pdf}\includegraphics[width=1em]{element-mustset.pdf}~zStart & \hfuzz=500pt double & \hfuzz=500pt \\
\hfuzz=500pt\includegraphics[width=1em]{connector.pdf}\includegraphics[width=1em]{element-mustset.pdf}~xEnd & \hfuzz=500pt double & \hfuzz=500pt [m] in north, east, up or vehicle system\\
\hfuzz=500pt\includegraphics[width=1em]{connector.pdf}\includegraphics[width=1em]{element-mustset.pdf}~yEnd & \hfuzz=500pt double & \hfuzz=500pt linear motion between start and end\\
\hfuzz=500pt\includegraphics[width=1em]{connector.pdf}\includegraphics[width=1em]{element-mustset.pdf}~zEnd & \hfuzz=500pt double & \hfuzz=500pt \\
\hfuzz=500pt\includegraphics[width=1em]{connector.pdf}\includegraphics[width=1em]{element.pdf}~timeStart & \hfuzz=500pt time & \hfuzz=500pt \\
\hfuzz=500pt\includegraphics[width=1em]{connector.pdf}\includegraphics[width=1em]{element.pdf}~timeEnd & \hfuzz=500pt time & \hfuzz=500pt \\
\hline
\end{tabularx}

\clearpage
%==================================
\subsection{InstrumentGnssReceiver2TimeSeries}\label{InstrumentGnssReceiver2TimeSeries}
Convert selected GNSS observations or residuals into a simpler time series format.
The \config{outputfileTimeSeries} is an \file{instrument file}{instrument} (MISCVALUES).
For each epoch the first data column contains the PRN, the second the satellite system,
followed by a column for each GNSS \configClass{type}{gnssType}.
As normally more than one GNSS transmitter is tracked per epoch, the output file
has several lines per observed epoch (epochs with the same time, one for each transmitter).

The second data column of the output contains a number representating the system
\begin{itemize}
\item 71: 'G', GPS
\item 82: 'R', GLONASS
\item 69: 'E', GALILEO
\item 67: 'C', BDS
\item 83: 'S', SBAS
\item 74: 'J', QZSS
\item 73: 'I', IRNSS .
\end{itemize}

A \file{GNSS residual file}{instrument} includes additional information
besides the residuals, which can also be selected with \configClass{type}{gnssType}
\begin{itemize}
\item \verb|A1*|, \verb|E1*|: azimuth and elevation at receiver
\item \verb|A2*|, \verb|E2*|: azimuth and elevation at transmitter
\item \verb|I**|: Estimated slant total electron content (STEC)
\end{itemize}

Furthermore these files may include for each residual \configClass{type}{gnssType}
information about the redundancy and the accuracy relation $\sigma/\sigma_0$
of the estimated $\sigma$ versus the apriori $\sigma_0$ from the least squares adjustment.
The three values (residuals, redundancy, $\sigma/\sigma_0$) are coded with the same type.
To get access to all values the corresponding type must be repeated in \configClass{type}{gnssType}.

Example: Selected GPS phase residuals (\configClass{type}{gnssType}='\verb|L1*G|' and \configClass{type}{gnssType}='\verb|L2*G|').
Plotted with \program{PlotGraph} with two \configClass{layer:linesAndPoints}{plotGraphLayerType}
(\config{valueX}='\verb|data0|',  \config{valueY}='\verb|100*data3+data1|' and \config{valueY}='\verb|100*data4+data1|' respectively).
\fig{!hb}{0.8}{instrumentGnssReceiver2TimeSeries}{fig:instrumentGnssReceiver2TimeSeries}{GPS residuals in cm, shifted by PRN}


\keepXColumns
\begin{tabularx}{\textwidth}{N T A}
\hline
Name & Type & Annotation\\
\hline
\hfuzz=500pt\includegraphics[width=1em]{element-mustset.pdf}~outputfileTimeSeries & \hfuzz=500pt filename & \hfuzz=500pt Instrument (MISCVALUES): prn, system, values for each type\\
\hfuzz=500pt\includegraphics[width=1em]{element-mustset-unbounded.pdf}~inputfileGnssReceiver & \hfuzz=500pt filename & \hfuzz=500pt GNSS receiver observations or residuals\\
\hfuzz=500pt\includegraphics[width=1em]{element-mustset-unbounded.pdf}~type & \hfuzz=500pt \hyperref[gnssType]{gnssType} & \hfuzz=500pt \\
\hline
\end{tabularx}

\clearpage
%==================================
\subsection{ParameterVector2GnssAntennaDefinition}\label{ParameterVector2GnssAntennaDefinition}
Updates an \file{GnssAntennaDefinition file}{gnssAntennaDefinition} with estimated parameters which belongs
to the parametrization \configClass{antennaCenterVariations}{parametrizationGnssAntennaType}.
The \configFile{outfileAntennaDefinition}{gnssAntennaDefinition} contains all antennas
from \configFile{inputfileAntennaDefinition}{gnssAntennaDefinition}.
The antenna center variations representend by the \configFile{inputfileSolution}{matrix} are added
to the matching antennas.

The \file{GnssAntennaDefinition file}{gnssAntennaDefinition} can be modified to the demands before with
\program{GnssAntennaDefinitionCreate}.

The following steps are used to estimate antenna center variations:
\begin{itemize}
\item \program{GnssAntennaDefinitionCreate} or \program{GnssAntex2AntennaDefinition}
\item \program{GnssProcessing} with \config{inputfileAntennaDefinition} as apriori
      and writing \file{normal equations}{normalEquation} with
      parametrization of \configClass{antennaCenterVariations}{parametrizationGnssAntennaType}
\item \program{NormalsEliminate}: eliminate all other than antenna parameters
\item \program{NormalsAccumulate}: accumulate normals over a sufficient long period
\item \program{GnssAntennaNormalsConstraint}: constrain unsolvable parameter linear combinations
\item \program{NormalsSolverVCE}: estimate the parameter vector
\item \program{ParameterVector2GnssAntennaDefinition}: update \config{inputfileAntennaDefinition}
\end{itemize}

See also \program{ParameterVector2GnssAntennaDefinition}, \program{GnssAntennaNormalsConstraint}.


\keepXColumns
\begin{tabularx}{\textwidth}{N T A}
\hline
Name & Type & Annotation\\
\hline
\hfuzz=500pt\includegraphics[width=1em]{element-mustset.pdf}~outfileAntennaDefinition & \hfuzz=500pt filename & \hfuzz=500pt all apriori antennas\\
\hfuzz=500pt\includegraphics[width=1em]{element-mustset.pdf}~inputfileAntennaDefinition & \hfuzz=500pt filename & \hfuzz=500pt apriori antennas\\
\hfuzz=500pt\includegraphics[width=1em]{element-mustset-unbounded.pdf}~antennaCenterVariations & \hfuzz=500pt \hyperref[parametrizationGnssAntennaType]{parametrizationGnssAntenna} & \hfuzz=500pt \\
\hfuzz=500pt\includegraphics[width=1em]{element-mustset.pdf}~inputfileSolution & \hfuzz=500pt filename & \hfuzz=500pt \\
\hfuzz=500pt\includegraphics[width=1em]{element-mustset.pdf}~inputfileParameterNames & \hfuzz=500pt filename & \hfuzz=500pt \\
\hline
\end{tabularx}

\clearpage
%==================================
\section{Programs: Grace}
\subsection{EnsembleAveragingScaleModel}\label{EnsembleAveragingScaleModel}
This programs estimate satellite-to-satellite-tracking (SST) deterministic signals due to eclipse transits from residuals.
The ensemble averaging method is used to characterize the average properties of signal shapes across all transit events.
Each shape is assigned to one arc of 3 hours (default). This can be modefied by enabling \config{averagingInterval}.


\keepXColumns
\begin{tabularx}{\textwidth}{N T A}
\hline
Name & Type & Annotation\\
\hline
\hfuzz=500pt\includegraphics[width=1em]{element-mustset.pdf}~outputfileScaleModel & \hfuzz=500pt filename & \hfuzz=500pt \\
\hfuzz=500pt\includegraphics[width=1em]{element-mustset.pdf}~inputfileGrace1EclipseFactor & \hfuzz=500pt filename & \hfuzz=500pt GRACE-A eclipse factors computed with integrated orbit\\
\hfuzz=500pt\includegraphics[width=1em]{element-mustset.pdf}~inputfileGrace2EclipseFactor & \hfuzz=500pt filename & \hfuzz=500pt GRACE-B eclipse factors computed with integrated orbit\\
\hfuzz=500pt\includegraphics[width=1em]{element-mustset.pdf}~inputfileGraceResiduals & \hfuzz=500pt filename & \hfuzz=500pt SST Residuals\\
\hfuzz=500pt\includegraphics[width=1em]{element-mustset.pdf}~timeMargin & \hfuzz=500pt uint & \hfuzz=500pt epochs before eclipse mode\\
\hfuzz=500pt\includegraphics[width=1em]{element-mustset.pdf}~waveLength & \hfuzz=500pt uint & \hfuzz=500pt length of the sample wave\\
\hfuzz=500pt\includegraphics[width=1em]{element.pdf}~averagingInterval & \hfuzz=500pt sequence & \hfuzz=500pt \\
\hfuzz=500pt\includegraphics[width=1em]{connector.pdf}\includegraphics[width=1em]{element.pdf}~nearestNeighborNumber & \hfuzz=500pt uint & \hfuzz=500pt \\
\hline
\end{tabularx}

\clearpage
%==================================
\subsection{GraceAntennaCenterCorrectionArcCovariance}\label{GraceAntennaCenterCorrectionArcCovariance}
This program computes covariance information for the non­-stationary noise of the KBR antenna offset correction (AOC)
from the orientation covariance matrices provided in Level-1B products via variance propagation.
By using the output \configFile{outputfileSatelliteTrackingCovariance}{matrix} in \program{PreprocessingSst},
noise model distinguishes between the stationary noise of ranging observations and the non­stationary AOC noise.

The covariances are derived from the partial derivative of the AOC w.r.t. the roll/pitch/yaw rotations
and star camera covariances \configFile{inputfileScaCovariance1}{matrix} and \configFile{inputfileScaCovariance2}{matrix}.

The covariances for the range-rates and range-acceleration are computed by differentiating
an interpolation polynomial of degree \config{interpolationDegree}.


\keepXColumns
\begin{tabularx}{\textwidth}{N T A}
\hline
Name & Type & Annotation\\
\hline
\hfuzz=500pt\includegraphics[width=1em]{element.pdf}~outputfileSatelliteTrackingCovariance & \hfuzz=500pt filename & \hfuzz=500pt corrections for range, range-rate, and range-accelerations\\
\hfuzz=500pt\includegraphics[width=1em]{element-mustset.pdf}~sstType & \hfuzz=500pt choice & \hfuzz=500pt \\
\hfuzz=500pt\includegraphics[width=1em]{connector.pdf}\includegraphics[width=1em]{element-mustset.pdf}~range & \hfuzz=500pt  & \hfuzz=500pt \\
\hfuzz=500pt\includegraphics[width=1em]{connector.pdf}\includegraphics[width=1em]{element-mustset.pdf}~rangeRate & \hfuzz=500pt  & \hfuzz=500pt \\
\hfuzz=500pt\includegraphics[width=1em]{connector.pdf}\includegraphics[width=1em]{element-mustset.pdf}~rangeAcceleration & \hfuzz=500pt  & \hfuzz=500pt \\
\hfuzz=500pt\includegraphics[width=1em]{element-mustset.pdf}~inputfileOrbit1 & \hfuzz=500pt filename & \hfuzz=500pt \\
\hfuzz=500pt\includegraphics[width=1em]{element-mustset.pdf}~inputfileOrbit2 & \hfuzz=500pt filename & \hfuzz=500pt \\
\hfuzz=500pt\includegraphics[width=1em]{element-mustset.pdf}~inputfileStarCamera1 & \hfuzz=500pt filename & \hfuzz=500pt \\
\hfuzz=500pt\includegraphics[width=1em]{element-mustset.pdf}~inputfileStarCamera2 & \hfuzz=500pt filename & \hfuzz=500pt \\
\hfuzz=500pt\includegraphics[width=1em]{element-mustset.pdf}~inputfileScaCovariance1 & \hfuzz=500pt filename & \hfuzz=500pt \\
\hfuzz=500pt\includegraphics[width=1em]{element-mustset.pdf}~inputfileScaCovariance2 & \hfuzz=500pt filename & \hfuzz=500pt \\
\hfuzz=500pt\includegraphics[width=1em]{element.pdf}~sigmaAccelerometerX & \hfuzz=500pt double & \hfuzz=500pt [rad/s\textasciicircum{}2]\\
\hfuzz=500pt\includegraphics[width=1em]{element.pdf}~sigmaAccelerometerY & \hfuzz=500pt double & \hfuzz=500pt [rad/s\textasciicircum{}2]\\
\hfuzz=500pt\includegraphics[width=1em]{element.pdf}~sigmaAccelerometerZ & \hfuzz=500pt double & \hfuzz=500pt [rad/s\textasciicircum{}2]\\
\hfuzz=500pt\includegraphics[width=1em]{element-mustset.pdf}~antennaCenters & \hfuzz=500pt choice & \hfuzz=500pt KBR antenna phase center\\
\hfuzz=500pt\includegraphics[width=1em]{connector.pdf}\includegraphics[width=1em]{element-mustset.pdf}~value & \hfuzz=500pt sequence & \hfuzz=500pt \\
\hfuzz=500pt\quad\includegraphics[width=1em]{connector.pdf}\includegraphics[width=1em]{element.pdf}~center1X & \hfuzz=500pt double & \hfuzz=500pt x-coordinate of antenna position in SRF [m] for GRACEA\\
\hfuzz=500pt\quad\includegraphics[width=1em]{connector.pdf}\includegraphics[width=1em]{element.pdf}~center1Y & \hfuzz=500pt double & \hfuzz=500pt y-coordinate of antenna position in SRF [m] for GRACEA\\
\hfuzz=500pt\quad\includegraphics[width=1em]{connector.pdf}\includegraphics[width=1em]{element.pdf}~center1Z & \hfuzz=500pt double & \hfuzz=500pt z-coordinate of antenna position in SRF [m] for GRACEA\\
\hfuzz=500pt\quad\includegraphics[width=1em]{connector.pdf}\includegraphics[width=1em]{element.pdf}~center2X & \hfuzz=500pt double & \hfuzz=500pt x-coordinate of antenna position in SRF [m] for GRACEB\\
\hfuzz=500pt\quad\includegraphics[width=1em]{connector.pdf}\includegraphics[width=1em]{element.pdf}~center2Y & \hfuzz=500pt double & \hfuzz=500pt y-coordinate of antenna position in SRF [m] for GRACEB\\
\hfuzz=500pt\quad\includegraphics[width=1em]{connector.pdf}\includegraphics[width=1em]{element.pdf}~center2Z & \hfuzz=500pt double & \hfuzz=500pt z-coordinate of antenna position in SRF [m] for GRACEB\\
\hfuzz=500pt\includegraphics[width=1em]{connector.pdf}\includegraphics[width=1em]{element-mustset.pdf}~file & \hfuzz=500pt sequence & \hfuzz=500pt \\
\hfuzz=500pt\quad\includegraphics[width=1em]{connector.pdf}\includegraphics[width=1em]{element-mustset.pdf}~inputAntennaCenters & \hfuzz=500pt filename & \hfuzz=500pt \\
\hfuzz=500pt\includegraphics[width=1em]{element.pdf}~interpolationDegree & \hfuzz=500pt uint & \hfuzz=500pt differentiation by polynomial approximation of degree n\\
\hline
\end{tabularx}

This program is \reference{parallelized}{general.parallelization}.
\clearpage
%==================================
\subsection{GraceOrbit2TransplantTimeOffset}\label{GraceOrbit2TransplantTimeOffset}
This program computes the time shift between two co-orbiting satellites based on dynamic orbit data.
When applied to data of the first satellite, the computed time shift virtually shifts data of first satellite into the location of the second satellite.
Note that \config{inputfileOrbit1} and \config{inputfileOrbit2} need velocity and acceleration data, which
can be computed with \program{OrbitAddVelocityAndAcceleration}.
The program tries to find a minimum of the objective function
\begin{equation}
  f(\Delta t) = \| r_1(t) - r_2(t + \Delta t) \|^2,
\end{equation}
by applying Newton's method to the first derivative, thus iteratively computing
\begin{equation}
  \Delta t_{k+1} = \Delta t_k + \frac{f'(\Delta t_k)}{f''(\Delta t_k)}.
\end{equation}
This iteration is stopped when the difference between to consecutive time shift values falls below \config{threshold} or
\config{maximumIterations} is reached. An \config{initialGuess} of the time shift can speed up convergence.

See also \program{OrbitAddVelocityAndAcceleration} and \program{InstrumentApplyTimeOffset}.


\keepXColumns
\begin{tabularx}{\textwidth}{N T A}
\hline
Name & Type & Annotation\\
\hline
\hfuzz=500pt\includegraphics[width=1em]{element-mustset.pdf}~outputfileTimeOffset & \hfuzz=500pt filename & \hfuzz=500pt estimated time offset in seconds (MISCVALUE)\\
\hfuzz=500pt\includegraphics[width=1em]{element-mustset.pdf}~inputfileOrbit1 & \hfuzz=500pt filename & \hfuzz=500pt orbit data of satellite 1\\
\hfuzz=500pt\includegraphics[width=1em]{element-mustset.pdf}~inputfileOrbit2 & \hfuzz=500pt filename & \hfuzz=500pt orbit data of satellite 2\\
\hfuzz=500pt\includegraphics[width=1em]{element.pdf}~interpolationDegree & \hfuzz=500pt uint & \hfuzz=500pt polynomial degree for the interpolation of position, velocity and acceleration\\
\hfuzz=500pt\includegraphics[width=1em]{element.pdf}~initialGuess & \hfuzz=500pt double & \hfuzz=500pt initial guess for the time shift [seconds]\\
\hfuzz=500pt\includegraphics[width=1em]{element.pdf}~maximumIterations & \hfuzz=500pt uint & \hfuzz=500pt maximum number of iterations\\
\hfuzz=500pt\includegraphics[width=1em]{element.pdf}~threshold & \hfuzz=500pt double & \hfuzz=500pt when the maximum difference between two iterations is below this value, stop [seconds]\\
\hline
\end{tabularx}

This program is \reference{parallelized}{general.parallelization}.
\clearpage
%==================================
\subsection{GraceSstResidualAnalysis}\label{GraceSstResidualAnalysis}
This program applies the Multi-Resolution Analysis (MRA) using
Discrete Wavelet Transform (DWT) to the monthly GRACE SST post-fit residuals.
First, the residuals are transferred into wavelet domain by applying an 8 level
Daubechies wavelet transform (default).
In the next step, detail coefficients are merged into three major groups
due to their approximate frequency subbands:
\begin{itemize}
\item Low scale details, corresponding to the frequency band above 10 mHz;
\item Intermediate scale details, corresponding to the approximate frequency
      range above 3 mHz up to 10 mHz;
\item High scale details, corresponding to the approximate frequency range
above 0.5 mHz up to 10 mHz.
\end{itemize}
In the last step, each group is reconstructed back into time domain.


\keepXColumns
\begin{tabularx}{\textwidth}{N T A}
\hline
Name & Type & Annotation\\
\hline
\hfuzz=500pt\includegraphics[width=1em]{element-mustset.pdf}~outputfileInstrumentHighScale & \hfuzz=500pt filename & \hfuzz=500pt High scale details\\
\hfuzz=500pt\includegraphics[width=1em]{element-mustset.pdf}~outputfileInstrumentMidScale & \hfuzz=500pt filename & \hfuzz=500pt Intermediate scale details\\
\hfuzz=500pt\includegraphics[width=1em]{element-mustset.pdf}~outputfileInstrumentLowScale & \hfuzz=500pt filename & \hfuzz=500pt Low scale details\\
\hfuzz=500pt\includegraphics[width=1em]{element-mustset.pdf}~inputfileInstrument & \hfuzz=500pt filename & \hfuzz=500pt GRACE SST Residuals\\
\hfuzz=500pt\includegraphics[width=1em]{element-mustset.pdf}~inputfileWavelet & \hfuzz=500pt filename & \hfuzz=500pt wavelet coefficients\\
\hline
\end{tabularx}

\clearpage
%==================================
\subsection{GraceSstScaleModel}\label{GraceSstScaleModel}
This programs estimate satellite-to-satellite-tracking (SST) deterministic signals
due to eclipse transits and low-SNR values from post-fit residuals.
The low-SNR effects are estimated by directly using the residual values.
The ensemble averaging method is used to characterize the average properties of eclipse transit signal shapes across all transit events.
Each shape is assigned to one arc of 3 hours (default). This can be modefied by enabling \config{averagingInterval}.


\keepXColumns
\begin{tabularx}{\textwidth}{N T A}
\hline
Name & Type & Annotation\\
\hline
\hfuzz=500pt\includegraphics[width=1em]{element-mustset.pdf}~inputfileGraceResiduals & \hfuzz=500pt filename & \hfuzz=500pt SST Residuals\\
\hfuzz=500pt\includegraphics[width=1em]{element-mustset.pdf}~timeMargin & \hfuzz=500pt uint & \hfuzz=500pt epochs before instrumental events\\
\hfuzz=500pt\includegraphics[width=1em]{element-mustset.pdf}~waveLength & \hfuzz=500pt uint & \hfuzz=500pt length of the sample wave\\
\hfuzz=500pt\includegraphics[width=1em]{element.pdf}~estimateEclipseTransitScale & \hfuzz=500pt sequence & \hfuzz=500pt \\
\hfuzz=500pt\includegraphics[width=1em]{connector.pdf}\includegraphics[width=1em]{element-mustset.pdf}~outputfileScaleModel & \hfuzz=500pt filename & \hfuzz=500pt \\
\hfuzz=500pt\includegraphics[width=1em]{connector.pdf}\includegraphics[width=1em]{element-mustset.pdf}~inputfileGrace1EclipseFactor & \hfuzz=500pt filename & \hfuzz=500pt GRACE-A eclipse factors computed with integrated orbit\\
\hfuzz=500pt\includegraphics[width=1em]{connector.pdf}\includegraphics[width=1em]{element-mustset.pdf}~inputfileGrace2EclipseFactor & \hfuzz=500pt filename & \hfuzz=500pt GRACE-B eclipse factors computed with integrated orbit\\
\hfuzz=500pt\includegraphics[width=1em]{connector.pdf}\includegraphics[width=1em]{element.pdf}~averagingInterval & \hfuzz=500pt sequence & \hfuzz=500pt \\
\hfuzz=500pt\quad\includegraphics[width=1em]{connector.pdf}\includegraphics[width=1em]{element.pdf}~nearestNeighborNumber & \hfuzz=500pt uint & \hfuzz=500pt \\
\hfuzz=500pt\includegraphics[width=1em]{element.pdf}~estimateLowSnrScale & \hfuzz=500pt sequence & \hfuzz=500pt \\
\hfuzz=500pt\includegraphics[width=1em]{connector.pdf}\includegraphics[width=1em]{element-mustset.pdf}~outputfileScaleModel & \hfuzz=500pt filename & \hfuzz=500pt \\
\hfuzz=500pt\includegraphics[width=1em]{connector.pdf}\includegraphics[width=1em]{element-mustset.pdf}~inputfileGraceSstSNR & \hfuzz=500pt filename & \hfuzz=500pt GRACE SNR values\\
\hline
\end{tabularx}

\clearpage
%==================================
\subsection{GraceSstSpecialEvents}\label{GraceSstSpecialEvents}
Time-indexing deterministic signals in the GRACE K-Band measurements caused by Sun intrusions
into the star camera baffles of GRACE-A and eclipse transits of the satellites.
The events are determined by satellites' position (\configFile{inputfileOrbit1/2}{instrument})
and orientation (\configFile{inputfileStarCamera1/2}{instrument}). Each type of event is represented
by its mid-interval point per orbit revolution and is reported in \configFile{outputfileEvents}{instrument}.

The waveform of each event is nearly constant within one month and can be approximated by a polynomial.
For the purpose of gravity field recovery, each waveform is parameterized by a polynomial and the coefficients
of this polynomial are estimated as additional instrument calibration parameters in a common adjustment
with all other instrument, satellite, and gravity field parameters,
see \configClass{parametrizationSatelliteTracking:specialEffect}{parametrizationSatelliteTrackingType:specialEffect}.


\keepXColumns
\begin{tabularx}{\textwidth}{N T A}
\hline
Name & Type & Annotation\\
\hline
\hfuzz=500pt\includegraphics[width=1em]{element-mustset.pdf}~outputfileEvents & \hfuzz=500pt filename & \hfuzz=500pt \\
\hfuzz=500pt\includegraphics[width=1em]{element.pdf}~outputfileIntervals & \hfuzz=500pt filename & \hfuzz=500pt \\
\hfuzz=500pt\includegraphics[width=1em]{element-mustset.pdf}~inputfileOrbit1 & \hfuzz=500pt filename & \hfuzz=500pt \\
\hfuzz=500pt\includegraphics[width=1em]{element-mustset.pdf}~inputfileOrbit2 & \hfuzz=500pt filename & \hfuzz=500pt \\
\hfuzz=500pt\includegraphics[width=1em]{element-mustset.pdf}~inputfileStarCamera1 & \hfuzz=500pt filename & \hfuzz=500pt \\
\hfuzz=500pt\includegraphics[width=1em]{element-mustset.pdf}~inputfileStarCamera2 & \hfuzz=500pt filename & \hfuzz=500pt \\
\hfuzz=500pt\includegraphics[width=1em]{element-mustset.pdf}~ephemerides & \hfuzz=500pt \hyperref[ephemeridesType]{ephemerides} & \hfuzz=500pt \\
\hfuzz=500pt\includegraphics[width=1em]{element-mustset.pdf}~eclipse & \hfuzz=500pt \hyperref[eclipseType]{eclipse} & \hfuzz=500pt \\
\hfuzz=500pt\includegraphics[width=1em]{element.pdf}~marginLeft & \hfuzz=500pt double & \hfuzz=500pt margin size (on both sides) [seconds]\\
\hfuzz=500pt\includegraphics[width=1em]{element.pdf}~marginRight & \hfuzz=500pt double & \hfuzz=500pt margin size (on both sides) [seconds]\\
\hline
\end{tabularx}

\clearpage
%==================================
\subsection{GraceThrusterResponse2Accelerometer}\label{GraceThrusterResponse2Accelerometer}
Add modeled thruster responses to accelerometer data.
The epochs and durations are given in the \configFile{inputfileThruster}{instrument} (THRUSTER).

The \configFile{inputfileThrusterResponse}{matrix} is a $(6\times 3)$ matrix with
the linear accelerations in the SRF ($x, y, z$) in one line per pair:
\begin{enumerate}
\item Negative Yaw,
\item Positive Pitch,
\item Positive Yaw,
\item Negative Pitch,
\item Negative Roll,
\item Positive Roll.
\end{enumerate}


\keepXColumns
\begin{tabularx}{\textwidth}{N T A}
\hline
Name & Type & Annotation\\
\hline
\hfuzz=500pt\includegraphics[width=1em]{element-mustset.pdf}~outputfileAccelerometer & \hfuzz=500pt filename & \hfuzz=500pt ACCELEROMETER\\
\hfuzz=500pt\includegraphics[width=1em]{element-mustset.pdf}~inputfileAccelerometer & \hfuzz=500pt filename & \hfuzz=500pt ACCELEROMETER\\
\hfuzz=500pt\includegraphics[width=1em]{element-mustset.pdf}~inputfileThruster & \hfuzz=500pt filename & \hfuzz=500pt THRUSTER\\
\hfuzz=500pt\includegraphics[width=1em]{element-mustset.pdf}~inputfileThrusterResponse & \hfuzz=500pt filename & \hfuzz=500pt thruster model (matrix with one line per pair)\\
\hline
\end{tabularx}

\clearpage
%==================================
\subsection{InstrumentSatelliteTrackingAntennaCenterCorrection}\label{InstrumentSatelliteTrackingAntennaCenterCorrection}
This program computes the correction due to offset of the antenna center relative the center of mass.
The offsets $\M c_A$ and $\M c_B$ in \configFile{inputfileAntennaCenters}{matrix} are given in the satellite
reference frame. These offsets are rotated into the the inertial frame with $\M D_A$ and $\M D_B$ from
\configFile{inputfileStarCamera}{instrument} and projected onto the line of sight (LOS)
\begin{equation}
  \rho_{AOC} = \M e_{AB}\cdot(\M D_A\,\M c_A - \M D_B\,\M c_B),
\end{equation}
with the unit vector in line of sight direction
\begin{equation}
  \M e_{AB} = \frac{\M r_B - \M r_A}{\left\lVert{\M r_B - \M r_A}\right\rVert}.
\end{equation}
The corrections for the range-rates and range-acceleration are computed by differentiating
an interpolation polynomial of degree \config{interpolationDegree}.


\keepXColumns
\begin{tabularx}{\textwidth}{N T A}
\hline
Name & Type & Annotation\\
\hline
\hfuzz=500pt\includegraphics[width=1em]{element.pdf}~outputfileSatelliteTracking & \hfuzz=500pt filename & \hfuzz=500pt corrections for range, range-rate, and range-accelerations\\
\hfuzz=500pt\includegraphics[width=1em]{element-mustset.pdf}~inputfileOrbit1 & \hfuzz=500pt filename & \hfuzz=500pt \\
\hfuzz=500pt\includegraphics[width=1em]{element-mustset.pdf}~inputfileOrbit2 & \hfuzz=500pt filename & \hfuzz=500pt \\
\hfuzz=500pt\includegraphics[width=1em]{element-mustset.pdf}~inputfileStarCamera1 & \hfuzz=500pt filename & \hfuzz=500pt \\
\hfuzz=500pt\includegraphics[width=1em]{element-mustset.pdf}~inputfileStarCamera2 & \hfuzz=500pt filename & \hfuzz=500pt \\
\hfuzz=500pt\includegraphics[width=1em]{element-mustset.pdf}~antennaCenters & \hfuzz=500pt choice & \hfuzz=500pt KBR antenna phase center\\
\hfuzz=500pt\includegraphics[width=1em]{connector.pdf}\includegraphics[width=1em]{element-mustset.pdf}~value & \hfuzz=500pt sequence & \hfuzz=500pt \\
\hfuzz=500pt\quad\includegraphics[width=1em]{connector.pdf}\includegraphics[width=1em]{element.pdf}~center1X & \hfuzz=500pt double & \hfuzz=500pt x-coordinate of antenna position in SRF [m] for GRACEA\\
\hfuzz=500pt\quad\includegraphics[width=1em]{connector.pdf}\includegraphics[width=1em]{element.pdf}~center1Y & \hfuzz=500pt double & \hfuzz=500pt y-coordinate of antenna position in SRF [m] for GRACEA\\
\hfuzz=500pt\quad\includegraphics[width=1em]{connector.pdf}\includegraphics[width=1em]{element.pdf}~center1Z & \hfuzz=500pt double & \hfuzz=500pt z-coordinate of antenna position in SRF [m] for GRACEA\\
\hfuzz=500pt\quad\includegraphics[width=1em]{connector.pdf}\includegraphics[width=1em]{element.pdf}~center2X & \hfuzz=500pt double & \hfuzz=500pt x-coordinate of antenna position in SRF [m] for GRACEB\\
\hfuzz=500pt\quad\includegraphics[width=1em]{connector.pdf}\includegraphics[width=1em]{element.pdf}~center2Y & \hfuzz=500pt double & \hfuzz=500pt y-coordinate of antenna position in SRF [m] for GRACEB\\
\hfuzz=500pt\quad\includegraphics[width=1em]{connector.pdf}\includegraphics[width=1em]{element.pdf}~center2Z & \hfuzz=500pt double & \hfuzz=500pt z-coordinate of antenna position in SRF [m] for GRACEB\\
\hfuzz=500pt\includegraphics[width=1em]{connector.pdf}\includegraphics[width=1em]{element-mustset.pdf}~file & \hfuzz=500pt sequence & \hfuzz=500pt \\
\hfuzz=500pt\quad\includegraphics[width=1em]{connector.pdf}\includegraphics[width=1em]{element-mustset.pdf}~inputAntennaCenters & \hfuzz=500pt filename & \hfuzz=500pt \\
\hfuzz=500pt\includegraphics[width=1em]{element.pdf}~interpolationDegree & \hfuzz=500pt uint & \hfuzz=500pt differentiation by polynomial approximation of degree n\\
\hline
\end{tabularx}

This program is \reference{parallelized}{general.parallelization}.
\clearpage
%==================================
\subsection{InstrumentStarCameraAngularAccelerometerFusion}\label{InstrumentStarCameraAngularAccelerometerFusion}
This program estimates the satellites orientation from star camera data
\configFile{inputfileStarCamera}{instrument} and angular accelerometer data
\configFile{inputfileAngularAcc}{instrument}. The combination of both observation types
is achieved in a least square adjustment. The optimal weighting between the two different
observation groups is achieved by means of VCE in combination with a robust estimator.
The system of linearized observation equations within the sensor fusion approach can be formulated as:
\begin{equation}
  \begin{bmatrix}
  \M l_{ACC1B}\\
  \M l_{SCA1B}
  \end{bmatrix}
  =
  \begin{bmatrix}
  \M A_{ACC1B} & \M B_{ACC1B}\\
  \M A_{SCA1B} & \M 0
  \end{bmatrix}
  \begin{bmatrix}
  \M q\\
  \M b
  \end{bmatrix}
  =
  \begin{bmatrix}
  \frac{\partial \dot{\boldsymbol{\omega}}}{\partial \M q} & \frac{\partial \dot{\boldsymbol{\omega}}}{\partial \M b}\\
  \M I & \M 0
  \end{bmatrix}
  \begin{bmatrix}
  \M q\\
  \M b
  \end{bmatrix}
\end{equation}
with
\begin{equation}\begin{split}
  \M l_{ACC1B}  &= \dot{\boldsymbol{\omega}}_{ACC1B} - \dot{\boldsymbol{\omega}}_{0}, \\
  \M l_{SCA1B}  &= \M q_{SCA1B} - \M q_{0}, \\
  \M q_{Fusion} &= \M q + \M q_{0}.
\end{split}\end{equation}
The reference values $\M q_{0}$ and $\dot{\boldsymbol{\omega}}_{0}$ are derived
from \configFile{inputfileStarCameraReference}{instrument}. In the course of the estimation,
the accelerometer data is calibrated, by setting a bias factor $\M b$ with \config{accBias}.


\keepXColumns
\begin{tabularx}{\textwidth}{N T A}
\hline
Name & Type & Annotation\\
\hline
\hfuzz=500pt\includegraphics[width=1em]{element-mustset.pdf}~outputfileStarCamera & \hfuzz=500pt filename & \hfuzz=500pt combined quaternions\\
\hfuzz=500pt\includegraphics[width=1em]{element.pdf}~outputfileCovariance & \hfuzz=500pt filename & \hfuzz=500pt epoch-wise covariance matrix\\
\hfuzz=500pt\includegraphics[width=1em]{element.pdf}~outputfileCovarianceMatrix & \hfuzz=500pt filename & \hfuzz=500pt full arc-wise covariance matrix per arc. arc number is appended to filename\\
\hfuzz=500pt\includegraphics[width=1em]{element.pdf}~outputfileEpochSigmaStarCamera & \hfuzz=500pt filename & \hfuzz=500pt from vce and outlier detection\\
\hfuzz=500pt\includegraphics[width=1em]{element.pdf}~outputfileEpochSigmaAccelerometer & \hfuzz=500pt filename & \hfuzz=500pt from vce and outlier detection\\
\hfuzz=500pt\includegraphics[width=1em]{element.pdf}~outputfileAngularAcc & \hfuzz=500pt filename & \hfuzz=500pt angular acceleration observations (bias removed)\\
\hfuzz=500pt\includegraphics[width=1em]{element.pdf}~outputfileSolution & \hfuzz=500pt filename & \hfuzz=500pt estimated parameter (one column for each arc)\\
\hfuzz=500pt\includegraphics[width=1em]{element-mustset.pdf}~inputfileStarCameraReference & \hfuzz=500pt filename & \hfuzz=500pt quaternions as taylor point\\
\hfuzz=500pt\includegraphics[width=1em]{element-mustset.pdf}~inputfileStarCamera & \hfuzz=500pt filename & \hfuzz=500pt star camera observations\\
\hfuzz=500pt\includegraphics[width=1em]{element.pdf}~inputfileStarCameraCovariance & \hfuzz=500pt filename & \hfuzz=500pt star camera observations\\
\hfuzz=500pt\includegraphics[width=1em]{element-mustset.pdf}~inputfileAngularAcc & \hfuzz=500pt filename & \hfuzz=500pt angular acceleration observations\\
\hfuzz=500pt\includegraphics[width=1em]{element.pdf}~correctAccNonQuadratic & \hfuzz=500pt boolean & \hfuzz=500pt apply correction (non-square proof mass)\\
\hfuzz=500pt\includegraphics[width=1em]{element-unbounded.pdf}~accBias & \hfuzz=500pt \hyperref[parametrizationTemporalType]{parametrizationTemporal} & \hfuzz=500pt accelerometer bias per interval and axis\\
\hfuzz=500pt\includegraphics[width=1em]{element-unbounded.pdf}~accScale & \hfuzz=500pt \hyperref[parametrizationTemporalType]{parametrizationTemporal} & \hfuzz=500pt accelerometer scale per interval and axis\\
\hfuzz=500pt\includegraphics[width=1em]{element.pdf}~sigmaStarcamera & \hfuzz=500pt double & \hfuzz=500pt [rad]\\
\hfuzz=500pt\includegraphics[width=1em]{element.pdf}~sigmaAccelerometerX & \hfuzz=500pt double & \hfuzz=500pt [rad/s\textasciicircum{}2]\\
\hfuzz=500pt\includegraphics[width=1em]{element.pdf}~sigmaAccelerometerY & \hfuzz=500pt double & \hfuzz=500pt [rad/s\textasciicircum{}2]\\
\hfuzz=500pt\includegraphics[width=1em]{element.pdf}~sigmaAccelerometerZ & \hfuzz=500pt double & \hfuzz=500pt [rad/s\textasciicircum{}2]\\
\hfuzz=500pt\includegraphics[width=1em]{element.pdf}~estimateSigmaScaPerAxis & \hfuzz=500pt boolean & \hfuzz=500pt separate variance factor for roll, pitch, yaw, instead of one common factor.\\
\hfuzz=500pt\includegraphics[width=1em]{element.pdf}~estimateSigmaAccPerAxis & \hfuzz=500pt boolean & \hfuzz=500pt separate variance factor for each accelerometer axis, instead of one common factor.\\
\hfuzz=500pt\includegraphics[width=1em]{element.pdf}~huber & \hfuzz=500pt double & \hfuzz=500pt residuals \$>\$ huber*sigma0 are downweighted\\
\hfuzz=500pt\includegraphics[width=1em]{element.pdf}~huberPower & \hfuzz=500pt double & \hfuzz=500pt residuals \$>\$ huber: sigma=(e/huber)\textasciicircum{}power*sigma0\\
\hfuzz=500pt\includegraphics[width=1em]{element.pdf}~interpolationDegree & \hfuzz=500pt uint & \hfuzz=500pt \\
\hfuzz=500pt\includegraphics[width=1em]{element.pdf}~iterationCount & \hfuzz=500pt uint & \hfuzz=500pt non linear equation solved iteratively\\
\hline
\end{tabularx}

This program is \reference{parallelized}{general.parallelization}.
\clearpage
%==================================
\section{Programs: Gravityfield}
\subsection{Gravityfield2AbsoluteGravity}\label{Gravityfield2AbsoluteGravity}
This program computes the absolute value of gravity $\left\lVert{\M g}\right\rVert$
of a \configClass{gravityfield}{gravityfieldType} on a given \configClass{grid}{gridType}.
The result is multiplicated with \config{factor}.
To get the full gravity vector in a terrestrial frame add
the centrifugal part, see \configClass{gravityfield:tides:centrifugal}{tidesType:centrifugal}.

The values will be saved together with points expressed as ellipsoidal coordinates (longitude, latitude, height)
based on a reference ellipsoid with parameters \config{R} and \config{inverseFlattening}.

It is intended to compute gravity anomalies from absolute gravity observations.
To visualize the results use \program{PlotMap}.


\keepXColumns
\begin{tabularx}{\textwidth}{N T A}
\hline
Name & Type & Annotation\\
\hline
\hfuzz=500pt\includegraphics[width=1em]{element-mustset.pdf}~outputfileGriddedData & \hfuzz=500pt filename & \hfuzz=500pt \\
\hfuzz=500pt\includegraphics[width=1em]{element-mustset-unbounded.pdf}~grid & \hfuzz=500pt \hyperref[gridType]{grid} & \hfuzz=500pt \\
\hfuzz=500pt\includegraphics[width=1em]{element-mustset-unbounded.pdf}~gravityfield & \hfuzz=500pt \hyperref[gravityfieldType]{gravityfield} & \hfuzz=500pt \\
\hfuzz=500pt\includegraphics[width=1em]{element.pdf}~factor & \hfuzz=500pt double & \hfuzz=500pt the result is multiplied by this factor, set -1 to subtract the field\\
\hfuzz=500pt\includegraphics[width=1em]{element.pdf}~time & \hfuzz=500pt time & \hfuzz=500pt at this time the gravity field will be evaluated\\
\hfuzz=500pt\includegraphics[width=1em]{element.pdf}~R & \hfuzz=500pt double & \hfuzz=500pt reference radius for ellipsoidal coordinates on output\\
\hfuzz=500pt\includegraphics[width=1em]{element.pdf}~inverseFlattening & \hfuzz=500pt double & \hfuzz=500pt reference flattening for ellipsoidal coordinates on output, 0: spherical coordinates\\
\hline
\end{tabularx}

This program is \reference{parallelized}{general.parallelization}.
\clearpage
%==================================
\subsection{Gravityfield2AreaMeanTimeSeries}\label{Gravityfield2AreaMeanTimeSeries}
This program computes a time series of time variable
\configClass{gravityfield}{gravityfieldType} functionals averaged over a given area,
e.g. equivalent water heights in the amazon basin.  The type of functional
(e.g gravity anomalies or geoid heights) can be choosen with \configClass{kernel}{kernelType}.
The average is performed at each time step by a weigthed average over all \configClass{grid}{gridType}
points where the weight is the associated area at each point. If \config{removeMean} is set
the temporal mean is removed from the time series. To speed up the computation
the gravity field can be converted to spherical harmonics before the computation
with \config{convertToHarmonics}.

Additionally the root mean square of the values in the area at each time step
can is computed if \config{compueRms} is set.

Additionally the accuracy of the value at each time step can be computed if \config{compueSigma} is set.

The \configFile{outputfileTimeSeries}{instrument} is an instrument file with one, two, or three data columns.
First data column contains the computed functionals and the following columns contain the RMS and the accuracies (optionally).

To visualize the results use \program{PlotGraph}.

\fig{!hb}{0.8}{gravityfield2AreaMeanTimeSeries}{fig:gravityfield2AreaMeanTimeSeries}{Amazon basin: Area mean values of ITSG-Grace2016 monthly solutions with error bars from computed sigma.}


\keepXColumns
\begin{tabularx}{\textwidth}{N T A}
\hline
Name & Type & Annotation\\
\hline
\hfuzz=500pt\includegraphics[width=1em]{element-mustset.pdf}~outputfileTimeSeries & \hfuzz=500pt filename & \hfuzz=500pt \\
\hfuzz=500pt\includegraphics[width=1em]{element-mustset-unbounded.pdf}~grid & \hfuzz=500pt \hyperref[gridType]{grid} & \hfuzz=500pt \\
\hfuzz=500pt\includegraphics[width=1em]{element-mustset-unbounded.pdf}~timeSeries & \hfuzz=500pt \hyperref[timeSeriesType]{timeSeries} & \hfuzz=500pt \\
\hfuzz=500pt\includegraphics[width=1em]{element-mustset.pdf}~kernel & \hfuzz=500pt \hyperref[kernelType]{kernel} & \hfuzz=500pt \\
\hfuzz=500pt\includegraphics[width=1em]{element-mustset-unbounded.pdf}~gravityfield & \hfuzz=500pt \hyperref[gravityfieldType]{gravityfield} & \hfuzz=500pt \\
\hfuzz=500pt\includegraphics[width=1em]{element.pdf}~convertToHarmonics & \hfuzz=500pt boolean & \hfuzz=500pt gravityfield is converted to spherical harmonics before evaluation, may accelerate the computation\\
\hfuzz=500pt\includegraphics[width=1em]{element.pdf}~multiplyWithArea & \hfuzz=500pt boolean & \hfuzz=500pt multiply time series with total area (useful for mass estimates)\\
\hfuzz=500pt\includegraphics[width=1em]{element.pdf}~removeMean & \hfuzz=500pt boolean & \hfuzz=500pt remove the temporal mean of the series\\
\hfuzz=500pt\includegraphics[width=1em]{element.pdf}~computeRms & \hfuzz=500pt boolean & \hfuzz=500pt additional rms each time step\\
\hfuzz=500pt\includegraphics[width=1em]{element.pdf}~computeSigma & \hfuzz=500pt boolean & \hfuzz=500pt additional error bars at each time step\\
\hline
\end{tabularx}

This program is \reference{parallelized}{general.parallelization}.
\clearpage
%==================================
\subsection{Gravityfield2Deflections}\label{Gravityfield2Deflections}
This program computes the deflections of the vertical $\xi$ in north direction
and $\eta$ in east direction in radian
according to
\begin{equation}
\xi = g_x/\gamma \qquad\text{and}\qquad \eta=g_y/\gamma,
\end{equation}
where $\M g=\nabla V$ is the gravity vector from \configClass{gravityfield}{gravityfieldType} in
the local ellipsoidal system (north, east, up) and $\gamma$ is the normal gravity at that point.

The values will be saved together with points expressed as ellipsoidal coordinates (longitude, latitude, height)
based on a reference ellipsoid with parameters \config{R} and \config{inverseFlattening}.


\keepXColumns
\begin{tabularx}{\textwidth}{N T A}
\hline
Name & Type & Annotation\\
\hline
\hfuzz=500pt\includegraphics[width=1em]{element-mustset.pdf}~outputfileGriddedData & \hfuzz=500pt filename & \hfuzz=500pt xi (north), eta (east) [rad]\\
\hfuzz=500pt\includegraphics[width=1em]{element-mustset-unbounded.pdf}~grid & \hfuzz=500pt \hyperref[gridType]{grid} & \hfuzz=500pt \\
\hfuzz=500pt\includegraphics[width=1em]{element-mustset-unbounded.pdf}~gravityfield & \hfuzz=500pt \hyperref[gravityfieldType]{gravityfield} & \hfuzz=500pt \\
\hfuzz=500pt\includegraphics[width=1em]{element.pdf}~time & \hfuzz=500pt time & \hfuzz=500pt at this time the gravity field will be evaluated\\
\hfuzz=500pt\includegraphics[width=1em]{element.pdf}~R & \hfuzz=500pt double & \hfuzz=500pt reference radius for ellipsoidal coordinates on output\\
\hfuzz=500pt\includegraphics[width=1em]{element.pdf}~inverseFlattening & \hfuzz=500pt double & \hfuzz=500pt reference flattening for ellipsoidal coordinates on output, 0: spherical coordinates\\
\hline
\end{tabularx}

This program is \reference{parallelized}{general.parallelization}.
\clearpage
%==================================
\subsection{Gravityfield2DegreeAmplitudes}\label{Gravityfield2DegreeAmplitudes}
This program computes degree amplitudes from a \configClass{gravityfield}{gravityfieldType}
and saves them to a \file{matrix}{matrix} file with three columns: the degree, the degree amplitude, and the formal errors.

The coefficients can be converted to different functionals with \configClass{kernel}{kernelType}.
The gravity field can be evaluated at different altitudes by specifying \config{evaluationRadius}.
Polar regions can be excluded by setting \config{polarGap}.
If set the expansion is limited in the range between \config{minDegree}
and \config{maxDegree} inclusivly.
The coefficients are related to the reference radius~\config{R}
and the Earth gravitational constant \config{GM}.

See also \program{PotentialCoefficients2DegreeAmplitudes}.


\keepXColumns
\begin{tabularx}{\textwidth}{N T A}
\hline
Name & Type & Annotation\\
\hline
\hfuzz=500pt\includegraphics[width=1em]{element-mustset.pdf}~outputfileMatrix & \hfuzz=500pt filename & \hfuzz=500pt three column matrix with degree, signal amplitude, formal error\\
\hfuzz=500pt\includegraphics[width=1em]{element-mustset-unbounded.pdf}~gravityfield & \hfuzz=500pt \hyperref[gravityfieldType]{gravityfield} & \hfuzz=500pt \\
\hfuzz=500pt\includegraphics[width=1em]{element-mustset.pdf}~kernel & \hfuzz=500pt \hyperref[kernelType]{kernel} & \hfuzz=500pt \\
\hfuzz=500pt\includegraphics[width=1em]{element-mustset.pdf}~type & \hfuzz=500pt choice & \hfuzz=500pt type of variances\\
\hfuzz=500pt\includegraphics[width=1em]{connector.pdf}\includegraphics[width=1em]{element-mustset.pdf}~rms & \hfuzz=500pt  & \hfuzz=500pt degree amplitudes (square root of degree variances)\\
\hfuzz=500pt\includegraphics[width=1em]{connector.pdf}\includegraphics[width=1em]{element-mustset.pdf}~accumulation & \hfuzz=500pt  & \hfuzz=500pt cumulate variances over degrees\\
\hfuzz=500pt\includegraphics[width=1em]{connector.pdf}\includegraphics[width=1em]{element-mustset.pdf}~median & \hfuzz=500pt  & \hfuzz=500pt median of absolute values per degree\\
\hfuzz=500pt\includegraphics[width=1em]{element.pdf}~time & \hfuzz=500pt time & \hfuzz=500pt at this time the gravity field will be evaluated\\
\hfuzz=500pt\includegraphics[width=1em]{element.pdf}~evaluationRadius & \hfuzz=500pt double & \hfuzz=500pt evaluate the gravity field at this radius (default: evaluate at surface\\
\hfuzz=500pt\includegraphics[width=1em]{element.pdf}~polarGap & \hfuzz=500pt angle & \hfuzz=500pt exclude polar regions (aperture angle in degrees)\\
\hfuzz=500pt\includegraphics[width=1em]{element.pdf}~minDegree & \hfuzz=500pt uint & \hfuzz=500pt \\
\hfuzz=500pt\includegraphics[width=1em]{element.pdf}~maxDegree & \hfuzz=500pt uint & \hfuzz=500pt \\
\hfuzz=500pt\includegraphics[width=1em]{element.pdf}~GM & \hfuzz=500pt double & \hfuzz=500pt Geocentric gravitational constant\\
\hfuzz=500pt\includegraphics[width=1em]{element.pdf}~R & \hfuzz=500pt double & \hfuzz=500pt reference radius\\
\hline
\end{tabularx}

\clearpage
%==================================
\subsection{Gravityfield2DegreeAmplitudesPlotGrid}\label{Gravityfield2DegreeAmplitudesPlotGrid}
This program computes a \configClass{timeSeries}{timeSeriesType}
of a time variable \configClass{gravityfield}{gravityfieldType} and saves it as degree amplitudes.
The expansion is limited in the range between \config{minDegree} and \config{maxDegree} inclusivly
\begin{equation}
  \sigma_n = \frac{GM}{R}\left(\frac{R}{r}\right)^{n+1}k_n\sqrt{\sum_{m=0}^n c_{nm}^2+s_{nm}^2}.
\end{equation}

The \configFile{outputfileTimeSeries}{matrix} is a matrix with
every row containing the time, degree, degree amplitude, and the formal error.

To visualize the results use \program{PlotGraph}).

See also \program{Gravityfield2DegreeAmplitudes}.

\fig{!hb}{0.5}{gravityfield2DegreeAmplitudesPlotGrid}{fig:gravityfield2DegreeAmplitudesPlotGrid}{Degree amplitudes of monthly ITSG-Grace2016 solutions relative to GOCO05s.}


\keepXColumns
\begin{tabularx}{\textwidth}{N T A}
\hline
Name & Type & Annotation\\
\hline
\hfuzz=500pt\includegraphics[width=1em]{element-mustset.pdf}~outputfileTimeSeries & \hfuzz=500pt filename & \hfuzz=500pt each row: mjd, degree, amplitude, formal error\\
\hfuzz=500pt\includegraphics[width=1em]{element-mustset-unbounded.pdf}~gravityfield & \hfuzz=500pt \hyperref[gravityfieldType]{gravityfield} & \hfuzz=500pt \\
\hfuzz=500pt\includegraphics[width=1em]{element-mustset.pdf}~kernel & \hfuzz=500pt \hyperref[kernelType]{kernel} & \hfuzz=500pt \\
\hfuzz=500pt\includegraphics[width=1em]{element-mustset-unbounded.pdf}~timeSeries & \hfuzz=500pt \hyperref[timeSeriesType]{timeSeries} & \hfuzz=500pt \\
\hfuzz=500pt\includegraphics[width=1em]{element.pdf}~evaluationRadius & \hfuzz=500pt double & \hfuzz=500pt evaluate the gravity field at this radius (default: evaluate at surface\\
\hfuzz=500pt\includegraphics[width=1em]{element.pdf}~polarGap & \hfuzz=500pt angle & \hfuzz=500pt exclude polar regions (aperture angle in degrees)\\
\hfuzz=500pt\includegraphics[width=1em]{element-mustset.pdf}~minDegree & \hfuzz=500pt uint & \hfuzz=500pt minimal degree\\
\hfuzz=500pt\includegraphics[width=1em]{element-mustset.pdf}~maxDegree & \hfuzz=500pt uint & \hfuzz=500pt maximal degree\\
\hfuzz=500pt\includegraphics[width=1em]{element.pdf}~GM & \hfuzz=500pt double & \hfuzz=500pt Geocentric gravitational constant\\
\hfuzz=500pt\includegraphics[width=1em]{element.pdf}~R & \hfuzz=500pt double & \hfuzz=500pt reference radius\\
\hline
\end{tabularx}

\clearpage
%==================================
\subsection{Gravityfield2DisplacementTimeSeries}\label{Gravityfield2DisplacementTimeSeries}
This program computes a time series of displacements of a list of stations (\configClass{grid}{gridType})
due to the effect of time variable loading masses. The displacement~$\M u$ of a station is calculated according to
\begin{equation}\label{eq:displacement}
\M u(\M r) = \frac{1}{\gamma}\sum_{n=0}^\infty \left[\frac{h_n}{1+k_n}V_n(\M r)\,\M e_{up}
+ R\frac{l_n}{1+k_n}\left(
 \frac{\partial V_n(\M r)}{\partial \M e_{north}}\M e_{north}
+\frac{\partial V_n(\M r)}{\partial \M e_{east}} \M e_{east}\right)\right],
\end{equation}
where $\gamma$ is the normal gravity, the load Love and Shida numbers $h_n,l_n$ are given by
\configFile{inputfileDeformationLoadLoveNumber}{matrix} and the load Love numbers $k_n$ are given by
\configFile{inputfilePotentialLoadLoveNumber}{matrix}. The $V_n$ are the spherical harmonics expansion of
the full time variable gravitational potential (potential of the loading mass + deformation potential):
\begin{equation}
V(\M r) = \sum_{n=0}^\infty V_n(\M r).
\end{equation}
Deformations due to Earth tide and due to polar tides are computed using the IERS conventions.
Eq.~\eqref{eq:displacement} is not used in these cases.

The \config{outputfileTimeSeries} is an \file{instrument file}{instrument}, MISCVALUES.
The data columns contain the deformation of each station in $x,y,z$ in a global terrestrial
reference frame or alternatively in a local elliposidal frame (north, east, up)
if \config{localReferenceFrame} is set.


\keepXColumns
\begin{tabularx}{\textwidth}{N T A}
\hline
Name & Type & Annotation\\
\hline
\hfuzz=500pt\includegraphics[width=1em]{element-mustset.pdf}~outputfileTimeSeries & \hfuzz=500pt filename & \hfuzz=500pt x,y,z [m] per station\\
\hfuzz=500pt\includegraphics[width=1em]{element-mustset-unbounded.pdf}~grid & \hfuzz=500pt \hyperref[gridType]{grid} & \hfuzz=500pt station list\\
\hfuzz=500pt\includegraphics[width=1em]{element-mustset-unbounded.pdf}~timeSeries & \hfuzz=500pt \hyperref[timeSeriesType]{timeSeries} & \hfuzz=500pt \\
\hfuzz=500pt\includegraphics[width=1em]{element-unbounded.pdf}~gravityfield & \hfuzz=500pt \hyperref[gravityfieldType]{gravityfield} & \hfuzz=500pt \\
\hfuzz=500pt\includegraphics[width=1em]{element-unbounded.pdf}~tides & \hfuzz=500pt \hyperref[tidesType]{tides} & \hfuzz=500pt \\
\hfuzz=500pt\includegraphics[width=1em]{element-mustset.pdf}~earthRotation & \hfuzz=500pt \hyperref[earthRotationType]{earthRotation} & \hfuzz=500pt \\
\hfuzz=500pt\includegraphics[width=1em]{element.pdf}~ephemerides & \hfuzz=500pt \hyperref[ephemeridesType]{ephemerides} & \hfuzz=500pt \\
\hfuzz=500pt\includegraphics[width=1em]{element-mustset.pdf}~inputfileDeformationLoadLoveNumber & \hfuzz=500pt filename & \hfuzz=500pt \\
\hfuzz=500pt\includegraphics[width=1em]{element.pdf}~inputfilePotentialLoadLoveNumber & \hfuzz=500pt filename & \hfuzz=500pt if full potential is given and not only loading potential\\
\hfuzz=500pt\includegraphics[width=1em]{element.pdf}~removeMean & \hfuzz=500pt boolean & \hfuzz=500pt remove the temporal mean of each coordinate\\
\hfuzz=500pt\includegraphics[width=1em]{element.pdf}~localReferenceFrame & \hfuzz=500pt boolean & \hfuzz=500pt local left handed reference frame (north, east, up)\\
\hline
\end{tabularx}

\clearpage
%==================================
\subsection{Gravityfield2EmpiricalCovariance}\label{Gravityfield2EmpiricalCovariance}
This program estimates an spatial and temporal covariance matrix from
a time series of gravity fields.

Firstly for every time step $t_i$
a spherical harmonics vector $\M x_i$ from the time variable gravity field
is generated. The coefficients of the spherical harmonics expansion are
in the sequence given by \configClass{numbering}{sphericalHarmonicsNumberingType}.
If set the expansion is limited in the range between \config{minDegree}
and \config{maxDegree} inclusivly. The coefficients are related to the
reference radius~\config{R} and the Earth gravitational constant \config{GM}.

In the next step the empirical covariance matrix is estimated according to
\begin{equation}
\M\Sigma(\Delta i)_{full} = \frac{1}{N}\sum_{i=1}^N \M x_i \M x_{i+\Delta i}^T,
\end{equation}
where $\Delta i$ is given by \config{differenceStep}.

From the diagonal elements of $\M\Sigma(\Delta i)$ the isotropic accuracies
are computed
\begin{equation}
\sigma_n^2 = \frac{1}{2n+1}\sum_{m=0}^n \sigma_{cnm}^2+\sigma_{snm}^2,
\end{equation}
and a diagonal matrix is constructed $\Sigma_{iso} = \text{diag}(\sigma_2^2,\ldots,\sigma_N^2)$.
The result is computed:
\begin{equation}
\M\Sigma(\Delta i) = \alpha_{full}\M\Sigma(\Delta i)_{full}+\alpha_{iso}\M\Sigma(\Delta i)_{iso}.
\end{equation}


\keepXColumns
\begin{tabularx}{\textwidth}{N T A}
\hline
Name & Type & Annotation\\
\hline
\hfuzz=500pt\includegraphics[width=1em]{element-mustset.pdf}~outputfileCovarianceMatrix & \hfuzz=500pt filename & \hfuzz=500pt \\
\hfuzz=500pt\includegraphics[width=1em]{element.pdf}~outputfilePotentialCoefficients & \hfuzz=500pt filename & \hfuzz=500pt \\
\hfuzz=500pt\includegraphics[width=1em]{element-mustset-unbounded.pdf}~gravityfield & \hfuzz=500pt \hyperref[gravityfieldType]{gravityfield} & \hfuzz=500pt \\
\hfuzz=500pt\includegraphics[width=1em]{element.pdf}~minDegree & \hfuzz=500pt uint & \hfuzz=500pt \\
\hfuzz=500pt\includegraphics[width=1em]{element-mustset.pdf}~maxDegree & \hfuzz=500pt uint & \hfuzz=500pt \\
\hfuzz=500pt\includegraphics[width=1em]{element.pdf}~GM & \hfuzz=500pt double & \hfuzz=500pt Geocentric gravitational constant\\
\hfuzz=500pt\includegraphics[width=1em]{element.pdf}~R & \hfuzz=500pt double & \hfuzz=500pt reference radius\\
\hfuzz=500pt\includegraphics[width=1em]{element-mustset.pdf}~numbering & \hfuzz=500pt \hyperref[sphericalHarmonicsNumberingType]{sphericalHarmonicsNumbering} & \hfuzz=500pt numbering scheme for solution vector\\
\hfuzz=500pt\includegraphics[width=1em]{element.pdf}~removeMean & \hfuzz=500pt boolean & \hfuzz=500pt \\
\hfuzz=500pt\includegraphics[width=1em]{element-mustset-unbounded.pdf}~timeSeries & \hfuzz=500pt \hyperref[timeSeriesType]{timeSeries} & \hfuzz=500pt sampling of the gravityfield\\
\hfuzz=500pt\includegraphics[width=1em]{element.pdf}~differenceStep & \hfuzz=500pt uint & \hfuzz=500pt choose dt for: x,i(t) - x,j(t+dt)\\
\hfuzz=500pt\includegraphics[width=1em]{element.pdf}~factorFullMatrixPart & \hfuzz=500pt double & \hfuzz=500pt \\
\hfuzz=500pt\includegraphics[width=1em]{element.pdf}~factorIsotropicPart & \hfuzz=500pt double & \hfuzz=500pt \\
\hfuzz=500pt\includegraphics[width=1em]{element-unbounded.pdf}~intervals & \hfuzz=500pt \hyperref[timeSeriesType]{timeSeries} & \hfuzz=500pt \\
\hline
\end{tabularx}

This program is \reference{parallelized}{general.parallelization}.
\clearpage
%==================================
\subsection{Gravityfield2Gradients}\label{Gravityfield2Gradients}
This program computes gravity gradients from \configClass{gravityfield}{gravityfieldType}
on a \configClass{grid}{gridType} in a global terrestrial reference frame
or alternatively in a local elliposidal frame (north, east, up) if \config{localReferenceFrame} is set.
In \configFile{outputfileGriddedData}{griddedData} the values $[Vxx, Vyy, Vzz, Vxy, Vxz, Vyz]$
will be saved together with points expressed as ellipsoidal coordinates
(longitude, latitude, height) based on a reference ellipsoid with parameters \config{R} and \config{inverseFlattening}.


\keepXColumns
\begin{tabularx}{\textwidth}{N T A}
\hline
Name & Type & Annotation\\
\hline
\hfuzz=500pt\includegraphics[width=1em]{element-mustset.pdf}~outputfileGriddedData & \hfuzz=500pt filename & \hfuzz=500pt Vxx Vyy Vzz Vxy Vxz Vyz\\
\hfuzz=500pt\includegraphics[width=1em]{element-mustset-unbounded.pdf}~grid & \hfuzz=500pt \hyperref[gridType]{grid} & \hfuzz=500pt \\
\hfuzz=500pt\includegraphics[width=1em]{element-mustset-unbounded.pdf}~gravityfield & \hfuzz=500pt \hyperref[gravityfieldType]{gravityfield} & \hfuzz=500pt \\
\hfuzz=500pt\includegraphics[width=1em]{element.pdf}~localReferenceFrame & \hfuzz=500pt boolean & \hfuzz=500pt local left handed reference frame (north, east, up)\\
\hfuzz=500pt\includegraphics[width=1em]{element.pdf}~time & \hfuzz=500pt time & \hfuzz=500pt at this time the gravity field will be evaluated\\
\hfuzz=500pt\includegraphics[width=1em]{element.pdf}~R & \hfuzz=500pt double & \hfuzz=500pt reference radius for ellipsoidal coordinates on output\\
\hfuzz=500pt\includegraphics[width=1em]{element.pdf}~inverseFlattening & \hfuzz=500pt double & \hfuzz=500pt reference flattening for ellipsoidal coordinates on output, 0: spherical coordinates\\
\hline
\end{tabularx}

This program is \reference{parallelized}{general.parallelization}.
\clearpage
%==================================
\subsection{Gravityfield2GridCovarianceMatrix}\label{Gravityfield2GridCovarianceMatrix}
This program propagates the covariance matrix of a \configClass{gravityfield}{gravityfieldType}
evaluated at \config{time} to a \configClass{grid}{gridType}. The full variance-covariance matrix is computed
and written to a \file{matrix file}{matrix}:
\begin{equation}
\mathbf{\Sigma}_\mathbf{y} = \mathbf{F}\mathbf{\Sigma}_\mathbf{x}\mathbf{F}^T
\end{equation}
The \configClass{kernel}{kernelType} determines the quantity of the grid values, for example,
\configClass{kernel:waterHeight}{kernelType:waterHeight}.

See also \program{GravityfieldCovariancesPropagation2GriddedData}, \program{GravityfieldVariancesPropagation2GriddedData}.


\keepXColumns
\begin{tabularx}{\textwidth}{N T A}
\hline
Name & Type & Annotation\\
\hline
\hfuzz=500pt\includegraphics[width=1em]{element-mustset.pdf}~outputfileMatrix & \hfuzz=500pt filename & \hfuzz=500pt symmetric grid covariance matrix\\
\hfuzz=500pt\includegraphics[width=1em]{element-mustset-unbounded.pdf}~grid & \hfuzz=500pt \hyperref[gridType]{grid} & \hfuzz=500pt \\
\hfuzz=500pt\includegraphics[width=1em]{element-mustset.pdf}~kernel & \hfuzz=500pt \hyperref[kernelType]{kernel} & \hfuzz=500pt \\
\hfuzz=500pt\includegraphics[width=1em]{element-mustset-unbounded.pdf}~gravityfield & \hfuzz=500pt \hyperref[gravityfieldType]{gravityfield} & \hfuzz=500pt \\
\hfuzz=500pt\includegraphics[width=1em]{element.pdf}~time & \hfuzz=500pt time & \hfuzz=500pt at this time the gravity field will be evaluated\\
\hline
\end{tabularx}

\clearpage
%==================================
\subsection{Gravityfield2GriddedData}\label{Gravityfield2GriddedData}
This program computes values of a \configClass{gravityfield}{gravityfieldType} on a given \configClass{grid}{gridType}.
The type of value (e.g gravity anomalies or geoid heights) can be choosen with \configClass{kernel}{kernelType}.
If a time is given the gravity field will be evaluated at this point of time otherwise only the static part will be used.
The values will be saved together with points expressed as ellipsoidal coordinates (longitude, latitude, height)
based on a reference ellipsoid with parameters \config{R} and \config{inverseFlattening}.
To speed up the computation the gravity field can be converted to spherical harmonics before the computation
with \config{convertToHarmonics}.

To visualize the results use \program{PlotMap}.


\keepXColumns
\begin{tabularx}{\textwidth}{N T A}
\hline
Name & Type & Annotation\\
\hline
\hfuzz=500pt\includegraphics[width=1em]{element-mustset.pdf}~outputfileGriddedData & \hfuzz=500pt filename & \hfuzz=500pt \\
\hfuzz=500pt\includegraphics[width=1em]{element-mustset-unbounded.pdf}~grid & \hfuzz=500pt \hyperref[gridType]{grid} & \hfuzz=500pt \\
\hfuzz=500pt\includegraphics[width=1em]{element-mustset.pdf}~kernel & \hfuzz=500pt \hyperref[kernelType]{kernel} & \hfuzz=500pt \\
\hfuzz=500pt\includegraphics[width=1em]{element-mustset-unbounded.pdf}~gravityfield & \hfuzz=500pt \hyperref[gravityfieldType]{gravityfield} & \hfuzz=500pt \\
\hfuzz=500pt\includegraphics[width=1em]{element.pdf}~convertToHarmonics & \hfuzz=500pt boolean & \hfuzz=500pt gravityfield is converted to spherical harmonics before evaluation, may accelerate the computation\\
\hfuzz=500pt\includegraphics[width=1em]{element.pdf}~time & \hfuzz=500pt time & \hfuzz=500pt at this time the gravity field will be evaluated\\
\hfuzz=500pt\includegraphics[width=1em]{element.pdf}~R & \hfuzz=500pt double & \hfuzz=500pt reference radius for ellipsoidal coordinates on output\\
\hfuzz=500pt\includegraphics[width=1em]{element.pdf}~inverseFlattening & \hfuzz=500pt double & \hfuzz=500pt reference flattening for ellipsoidal coordinates on output, 0: spherical coordinates\\
\hline
\end{tabularx}

This program is \reference{parallelized}{general.parallelization}.
\clearpage
%==================================
\subsection{Gravityfield2GriddedDataTimeSeries}\label{Gravityfield2GriddedDataTimeSeries}
This program computes values of a \configClass{gravityfield}{gravityfieldType} on a given \configClass{grid}{gridType}
for each time step of \configClass{timeSeries}{timeSeriesType}.
The type of value (e.g gravity anomalies or geoid heights) can be choosen with \configClass{kernel}{kernelType}.
To speed up the computation the gravity field can be converted to spherical harmonics before the computation
with \config{convertToHarmonics}.
The \configFile{outputfileTimeSeries}{instrument} is an instrument (MISCVALUES) file with a data column
for each grid point per epoch.

This program enables the use of all instrument programs like \program{InstrumentFilter},
\program{InstrumentArcStatistics} or \program{InstrumentDetrend} to analyze and manipulate time series of gridded data.

See also \program{TimeSeries2GriddedData}, \program{Gravityfield2GriddedData}


\keepXColumns
\begin{tabularx}{\textwidth}{N T A}
\hline
Name & Type & Annotation\\
\hline
\hfuzz=500pt\includegraphics[width=1em]{element-mustset.pdf}~outputfileTimeSeries & \hfuzz=500pt filename & \hfuzz=500pt each epoch: data of grid points (MISCVALUES)\\
\hfuzz=500pt\includegraphics[width=1em]{element-mustset-unbounded.pdf}~grid & \hfuzz=500pt \hyperref[gridType]{grid} & \hfuzz=500pt \\
\hfuzz=500pt\includegraphics[width=1em]{element-mustset.pdf}~kernel & \hfuzz=500pt \hyperref[kernelType]{kernel} & \hfuzz=500pt \\
\hfuzz=500pt\includegraphics[width=1em]{element-mustset-unbounded.pdf}~gravityfield & \hfuzz=500pt \hyperref[gravityfieldType]{gravityfield} & \hfuzz=500pt \\
\hfuzz=500pt\includegraphics[width=1em]{element.pdf}~convertToHarmonics & \hfuzz=500pt boolean & \hfuzz=500pt gravityfield is converted to spherical harmonics before evaluation, may accelerate the computation\\
\hfuzz=500pt\includegraphics[width=1em]{element-mustset-unbounded.pdf}~timeSeries & \hfuzz=500pt \hyperref[timeSeriesType]{timeSeries} & \hfuzz=500pt \\
\hline
\end{tabularx}

This program is \reference{parallelized}{general.parallelization}.
\clearpage
%==================================
\subsection{Gravityfield2PotentialCoefficients}\label{Gravityfield2PotentialCoefficients}
This program evaluates a time variable \configClass{gravityfield}{gravityfieldType}
at a given \config{time} and saves it as a \file{spherical harmonics file}{potentialCoefficients}.
If set the expansion is limited in the range between \config{minDegree}
and \config{maxDegree} inclusivly.
The coefficients are related to the reference radius~\config{R}
and the Earth gravitational constant \config{GM}.


\keepXColumns
\begin{tabularx}{\textwidth}{N T A}
\hline
Name & Type & Annotation\\
\hline
\hfuzz=500pt\includegraphics[width=1em]{element-mustset.pdf}~outputfilePotentialCoefficients & \hfuzz=500pt filename & \hfuzz=500pt \\
\hfuzz=500pt\includegraphics[width=1em]{element-mustset-unbounded.pdf}~gravityfield & \hfuzz=500pt \hyperref[gravityfieldType]{gravityfield} & \hfuzz=500pt \\
\hfuzz=500pt\includegraphics[width=1em]{element.pdf}~minDegree & \hfuzz=500pt uint & \hfuzz=500pt \\
\hfuzz=500pt\includegraphics[width=1em]{element.pdf}~maxDegree & \hfuzz=500pt uint & \hfuzz=500pt \\
\hfuzz=500pt\includegraphics[width=1em]{element.pdf}~GM & \hfuzz=500pt double & \hfuzz=500pt Geocentric gravitational constant\\
\hfuzz=500pt\includegraphics[width=1em]{element.pdf}~R & \hfuzz=500pt double & \hfuzz=500pt reference radius\\
\hfuzz=500pt\includegraphics[width=1em]{element.pdf}~time & \hfuzz=500pt time & \hfuzz=500pt at this time the gravity field will be evaluated\\
\hline
\end{tabularx}

\clearpage
%==================================
\subsection{Gravityfield2PotentialCoefficientsTimeSeries}\label{Gravityfield2PotentialCoefficientsTimeSeries}
This program computes a \configClass{timeSeries}{timeSeriesType}
of a time variable \configClass{gravityfield}{gravityfieldType}
and converts to coefficients of a spherical harmonics expansion.
The expansion is limited in the range between \config{minDegree}
and \config{maxDegree} inclusivly.
The coefficients are related to the reference radius~\config{R}
and the Earth gravitational constant \config{GM}.

The \configFile{outputfileTimeSeries}{instrument} contains the potential coefficients
as data columns for each epoch in the sequence given by
\configClass{numbering}{sphericalHarmonicsNumberingType}.


\keepXColumns
\begin{tabularx}{\textwidth}{N T A}
\hline
Name & Type & Annotation\\
\hline
\hfuzz=500pt\includegraphics[width=1em]{element-mustset.pdf}~outputfileTimeSeries & \hfuzz=500pt filename & \hfuzz=500pt instrument file (MISCVALUES)\\
\hfuzz=500pt\includegraphics[width=1em]{element-mustset-unbounded.pdf}~gravityfield & \hfuzz=500pt \hyperref[gravityfieldType]{gravityfield} & \hfuzz=500pt \\
\hfuzz=500pt\includegraphics[width=1em]{element-mustset-unbounded.pdf}~timeSeries & \hfuzz=500pt \hyperref[timeSeriesType]{timeSeries} & \hfuzz=500pt \\
\hfuzz=500pt\includegraphics[width=1em]{element-mustset.pdf}~minDegree & \hfuzz=500pt uint & \hfuzz=500pt \\
\hfuzz=500pt\includegraphics[width=1em]{element-mustset.pdf}~maxDegree & \hfuzz=500pt uint & \hfuzz=500pt \\
\hfuzz=500pt\includegraphics[width=1em]{element.pdf}~GM & \hfuzz=500pt double & \hfuzz=500pt Geocentric gravitational constant\\
\hfuzz=500pt\includegraphics[width=1em]{element.pdf}~R & \hfuzz=500pt double & \hfuzz=500pt reference radius\\
\hfuzz=500pt\includegraphics[width=1em]{element-mustset.pdf}~numbering & \hfuzz=500pt \hyperref[sphericalHarmonicsNumberingType]{sphericalHarmonicsNumbering} & \hfuzz=500pt numbering scheme\\
\hline
\end{tabularx}

\clearpage
%==================================
\subsection{Gravityfield2SphericalHarmonicsVector}\label{Gravityfield2SphericalHarmonicsVector}
This program evaluates a time variable \configClass{gravityfield}{gravityfieldType} at a given \config{time}
and saves a \file{vector}{matrix} with the coefficients of a spherical harmonics expansion in the sequence given by
\configClass{numbering}{sphericalHarmonicsNumberingType}.
If set the expansion is limited in the range between \config{minDegree} and \config{maxDegree} inclusively.
The coefficients are related to the reference radius~\config{R} and the Earth gravitational constant \config{GM}.

This coefficients vector can be used as a approximate solution, see \program{NormalsMultiplyAdd},
or as pseudo oberservations for regularization,
see \configClass{normalEquation:regularization}{normalEquationType:regularization}.

For back transformation use \program{Gravityfield2PotentialCoefficients}
with \configClass{gravityfield:fromParametrization}{gravityfieldType:fromParametrization}.


\keepXColumns
\begin{tabularx}{\textwidth}{N T A}
\hline
Name & Type & Annotation\\
\hline
\hfuzz=500pt\includegraphics[width=1em]{element-mustset.pdf}~outputfileVector & \hfuzz=500pt filename & \hfuzz=500pt \\
\hfuzz=500pt\includegraphics[width=1em]{element-mustset-unbounded.pdf}~gravityfield & \hfuzz=500pt \hyperref[gravityfieldType]{gravityfield} & \hfuzz=500pt \\
\hfuzz=500pt\includegraphics[width=1em]{element.pdf}~startIndex & \hfuzz=500pt uint & \hfuzz=500pt start index to put the coefficients in the solution vector\\
\hfuzz=500pt\includegraphics[width=1em]{element-mustset.pdf}~minDegree & \hfuzz=500pt uint & \hfuzz=500pt \\
\hfuzz=500pt\includegraphics[width=1em]{element-mustset.pdf}~maxDegree & \hfuzz=500pt uint & \hfuzz=500pt \\
\hfuzz=500pt\includegraphics[width=1em]{element.pdf}~GM & \hfuzz=500pt double & \hfuzz=500pt Geocentric gravitational constant\\
\hfuzz=500pt\includegraphics[width=1em]{element.pdf}~R & \hfuzz=500pt double & \hfuzz=500pt reference radius\\
\hfuzz=500pt\includegraphics[width=1em]{element-mustset.pdf}~numbering & \hfuzz=500pt \hyperref[sphericalHarmonicsNumberingType]{sphericalHarmonicsNumbering} & \hfuzz=500pt numbering scheme for solution vector\\
\hfuzz=500pt\includegraphics[width=1em]{element.pdf}~time & \hfuzz=500pt time & \hfuzz=500pt at this time the gravity field will be evaluated\\
\hfuzz=500pt\includegraphics[width=1em]{element.pdf}~useSigma & \hfuzz=500pt boolean & \hfuzz=500pt use formal errors instead of coefficients\\
\hline
\end{tabularx}

\clearpage
%==================================
\subsection{Gravityfield2TimeSplines}\label{Gravityfield2TimeSplines}
This program estimates splines in time domain from a time variable gravity field
and writes \configFile{outputfileTimeSplines}{timeSplinesGravityField}.
The \configClass{gravityfield}{gravityfieldType} is sampled at \configClass{sampling}{timeSeriesType}, converted to potential coefficients
in the range between \config{minDegree} and \config{maxDegree} inclusively.
The time series of spherical harmonics can be temporal filtered with \configClass{temporalFilter}{digitalFilterType}.

In the next step temporal splines with \config{splineDegree} and nodal points given
at \configClass{splineTimeSeries}{timeSeriesType} are adjusted to the time series in a least squares sense.
This is very fast for block means (splineDegree = 0) but for other splines a large systems of equations
must be solved. In the adjustment process the time series of gravity fields can be interpreted as samples
at the given times or as continuous function with linear behaviour between sampled points (\config{linearInterpolation}).

To combine a series of potential coefficients to a spline file with block means (splineDegree = 0)
use the fast \program{PotentialCoefficients2BlockMeanTimeSplines} instead.


\keepXColumns
\begin{tabularx}{\textwidth}{N T A}
\hline
Name & Type & Annotation\\
\hline
\hfuzz=500pt\includegraphics[width=1em]{element-mustset.pdf}~outputfileTimeSplines & \hfuzz=500pt filename & \hfuzz=500pt \\
\hfuzz=500pt\includegraphics[width=1em]{element-mustset-unbounded.pdf}~gravityfield & \hfuzz=500pt \hyperref[gravityfieldType]{gravityfield} & \hfuzz=500pt \\
\hfuzz=500pt\includegraphics[width=1em]{element-unbounded.pdf}~temporalFilter & \hfuzz=500pt \hyperref[digitalFilterType]{digitalFilter} & \hfuzz=500pt filter sampled gravity field in time\\
\hfuzz=500pt\includegraphics[width=1em]{element.pdf}~minDegree & \hfuzz=500pt uint & \hfuzz=500pt \\
\hfuzz=500pt\includegraphics[width=1em]{element.pdf}~maxDegree & \hfuzz=500pt uint & \hfuzz=500pt \\
\hfuzz=500pt\includegraphics[width=1em]{element.pdf}~GM & \hfuzz=500pt double & \hfuzz=500pt Geocentric gravitational constant\\
\hfuzz=500pt\includegraphics[width=1em]{element.pdf}~R & \hfuzz=500pt double & \hfuzz=500pt reference radius\\
\hfuzz=500pt\includegraphics[width=1em]{element-mustset-unbounded.pdf}~sampling & \hfuzz=500pt \hyperref[timeSeriesType]{timeSeries} & \hfuzz=500pt gravity field is sampled at these times\\
\hfuzz=500pt\includegraphics[width=1em]{element.pdf}~removeMean & \hfuzz=500pt boolean & \hfuzz=500pt remove the temporal mean of the series before estimating the splines\\
\hfuzz=500pt\includegraphics[width=1em]{element.pdf}~linearInterpolation & \hfuzz=500pt boolean & \hfuzz=500pt assume linear behavior between sampled points\\
\hfuzz=500pt\includegraphics[width=1em]{element-mustset.pdf}~splineDegree & \hfuzz=500pt uint & \hfuzz=500pt degree of splines\\
\hfuzz=500pt\includegraphics[width=1em]{element-mustset-unbounded.pdf}~splineTimeSeries & \hfuzz=500pt \hyperref[timeSeriesType]{timeSeries} & \hfuzz=500pt nodal points of splines in time domain\\
\hline
\end{tabularx}

\clearpage
%==================================
\subsection{Gravityfield2TrendPotentialCoefficients}\label{Gravityfield2TrendPotentialCoefficients}
This program estimates \configClass{parametrizationTemporal}{parametrizationTemporalType}
(e.g. mean, trend, annual) from a time variable gravity field.

In a first step a time variable \configClass{gravityfield}{gravityfieldType}
is sampled at \configClass{timeSeries}{timeSeriesType}
and converted to coefficients of a spherical harmonics expansion.
The expansion is limited in the range between \config{minDegree}
and \config{maxDegree} inclusively.
The coefficients are related to the reference radius~\config{R}
and the Earth gravitational constant \config{GM}.

These coefficients serves as observations of
a \reference{robust least squares adjustment}{fundamentals.robustLeastSquares} to estimate
\configClass{parametrizationTemporal}{parametrizationTemporalType} parameters.
For each temporal parameter an
\configFile{outputfilePotentialCoefficients}{potentialCoefficients}
is generated.


\keepXColumns
\begin{tabularx}{\textwidth}{N T A}
\hline
Name & Type & Annotation\\
\hline
\hfuzz=500pt\includegraphics[width=1em]{element-mustset-unbounded.pdf}~outputfilePotentialCoefficients & \hfuzz=500pt filename & \hfuzz=500pt for each temporal parameter\\
\hfuzz=500pt\includegraphics[width=1em]{element-mustset-unbounded.pdf}~gravityfield & \hfuzz=500pt \hyperref[gravityfieldType]{gravityfield} & \hfuzz=500pt \\
\hfuzz=500pt\includegraphics[width=1em]{element-mustset-unbounded.pdf}~timeSeries & \hfuzz=500pt \hyperref[timeSeriesType]{timeSeries} & \hfuzz=500pt \\
\hfuzz=500pt\includegraphics[width=1em]{element-mustset-unbounded.pdf}~parametrizationTemporal & \hfuzz=500pt \hyperref[parametrizationTemporalType]{parametrizationTemporal} & \hfuzz=500pt \\
\hfuzz=500pt\includegraphics[width=1em]{element-mustset.pdf}~minDegree & \hfuzz=500pt uint & \hfuzz=500pt \\
\hfuzz=500pt\includegraphics[width=1em]{element.pdf}~maxDegree & \hfuzz=500pt uint & \hfuzz=500pt \\
\hfuzz=500pt\includegraphics[width=1em]{element.pdf}~GM & \hfuzz=500pt double & \hfuzz=500pt Geocentric gravitational constant\\
\hfuzz=500pt\includegraphics[width=1em]{element.pdf}~R & \hfuzz=500pt double & \hfuzz=500pt reference radius\\
\hfuzz=500pt\includegraphics[width=1em]{element.pdf}~huber & \hfuzz=500pt double & \hfuzz=500pt for robust least squares\\
\hfuzz=500pt\includegraphics[width=1em]{element.pdf}~huberPower & \hfuzz=500pt double & \hfuzz=500pt for robust least squares\\
\hfuzz=500pt\includegraphics[width=1em]{element.pdf}~huberMaxIteration & \hfuzz=500pt uint & \hfuzz=500pt (maximum) number of iterations for robust estimation\\
\hline
\end{tabularx}

\clearpage
%==================================
\subsection{GravityfieldCovariancesPropagation2GriddedData}\label{GravityfieldCovariancesPropagation2GriddedData}
This program computes the covariance between a source point given
by longitude/latitude (\config{L}, \config{B}) and the points of a \configClass{grid}{gridType}
in terms of the functional given by \configClass{kernel}{kernelType} from the variance-covariance
matrix of a \configClass{gravityfield}{gravityfieldType} evaluated at \config{time}.

If \config{computeCorrelation} is set, the program returns the correlation according to
\begin{equation}
r_{ij} = \frac{\sigma_{ij}}{\sigma_i \sigma_j}
\end{equation}
in the range of [-1, 1] instead of the covariance.

See also \program{Gravityfield2GridCovarianceMatrix}, \program{GravityfieldVariancesPropagation2GriddedData}.


\keepXColumns
\begin{tabularx}{\textwidth}{N T A}
\hline
Name & Type & Annotation\\
\hline
\hfuzz=500pt\includegraphics[width=1em]{element-mustset.pdf}~outputfileGriddedData & \hfuzz=500pt filename & \hfuzz=500pt gridded data file containing the covariance betwenn source point and grid points\\
\hfuzz=500pt\includegraphics[width=1em]{element-mustset-unbounded.pdf}~grid & \hfuzz=500pt \hyperref[gridType]{grid} & \hfuzz=500pt \\
\hfuzz=500pt\includegraphics[width=1em]{element-mustset.pdf}~kernel & \hfuzz=500pt \hyperref[kernelType]{kernel} & \hfuzz=500pt functional\\
\hfuzz=500pt\includegraphics[width=1em]{element-mustset-unbounded.pdf}~gravityfield & \hfuzz=500pt \hyperref[gravityfieldType]{gravityfield} & \hfuzz=500pt \\
\hfuzz=500pt\includegraphics[width=1em]{element.pdf}~time & \hfuzz=500pt time & \hfuzz=500pt at this time the gravity field will be evaluated\\
\hfuzz=500pt\includegraphics[width=1em]{element.pdf}~L & \hfuzz=500pt angle & \hfuzz=500pt longitude of variance point\\
\hfuzz=500pt\includegraphics[width=1em]{element.pdf}~B & \hfuzz=500pt angle & \hfuzz=500pt latitude of variance point\\
\hfuzz=500pt\includegraphics[width=1em]{element.pdf}~height & \hfuzz=500pt double & \hfuzz=500pt ellipsoidal height of source point\\
\hfuzz=500pt\includegraphics[width=1em]{element.pdf}~computeCorrelation & \hfuzz=500pt boolean & \hfuzz=500pt compute correlations instead of covariances\\
\hfuzz=500pt\includegraphics[width=1em]{element.pdf}~R & \hfuzz=500pt double & \hfuzz=500pt reference radius for ellipsoidal coordinates on output\\
\hfuzz=500pt\includegraphics[width=1em]{element.pdf}~inverseFlattening & \hfuzz=500pt double & \hfuzz=500pt reference flattening for ellipsoidal coordinates on output, 0: spherical coordinates\\
\hline
\end{tabularx}

This program is \reference{parallelized}{general.parallelization}.
\clearpage
%==================================
\subsection{GravityfieldReplacePotentialCoefficients}\label{GravityfieldReplacePotentialCoefficients}
Replaces single potential coefficients in a gravity field.
Both \configClass{gravityfield}{gravityfieldType}
and \configClass{gravityfieldReplacement}{gravityfieldType} are evaluated
at \config{time} and converted to spherical harmonic coefficients.
Single \config{coefficients} are then replaced in \configClass{gravityfield}{gravityfieldType}
by the values from \configClass{gravityfieldReplacement}{gravityfieldType}
and the result is written to \configFile{outputfilePotentialCoefficients}{potentialCoefficients}
from \config{minDegree} to \config{maxDegree},


\keepXColumns
\begin{tabularx}{\textwidth}{N T A}
\hline
Name & Type & Annotation\\
\hline
\hfuzz=500pt\includegraphics[width=1em]{element-mustset.pdf}~outputfilePotentialCoefficients & \hfuzz=500pt filename & \hfuzz=500pt \\
\hfuzz=500pt\includegraphics[width=1em]{element-mustset-unbounded.pdf}~gravityfield & \hfuzz=500pt \hyperref[gravityfieldType]{gravityfield} & \hfuzz=500pt single coefficients are replaced by the other gravityfield\\
\hfuzz=500pt\includegraphics[width=1em]{element-mustset-unbounded.pdf}~gravityfieldReplacement & \hfuzz=500pt \hyperref[gravityfieldType]{gravityfield} & \hfuzz=500pt contains the coefficients for replacement\\
\hfuzz=500pt\includegraphics[width=1em]{element-mustset-unbounded.pdf}~coefficients & \hfuzz=500pt choice & \hfuzz=500pt \\
\hfuzz=500pt\includegraphics[width=1em]{connector.pdf}\includegraphics[width=1em]{element-mustset.pdf}~cnm & \hfuzz=500pt sequence & \hfuzz=500pt \\
\hfuzz=500pt\quad\includegraphics[width=1em]{connector.pdf}\includegraphics[width=1em]{element-mustset.pdf}~degree & \hfuzz=500pt uint & \hfuzz=500pt \\
\hfuzz=500pt\quad\includegraphics[width=1em]{connector.pdf}\includegraphics[width=1em]{element-mustset.pdf}~order & \hfuzz=500pt uint & \hfuzz=500pt \\
\hfuzz=500pt\includegraphics[width=1em]{connector.pdf}\includegraphics[width=1em]{element-mustset.pdf}~snm & \hfuzz=500pt sequence & \hfuzz=500pt \\
\hfuzz=500pt\quad\includegraphics[width=1em]{connector.pdf}\includegraphics[width=1em]{element-mustset.pdf}~degree & \hfuzz=500pt uint & \hfuzz=500pt \\
\hfuzz=500pt\quad\includegraphics[width=1em]{connector.pdf}\includegraphics[width=1em]{element-mustset.pdf}~order & \hfuzz=500pt uint & \hfuzz=500pt \\
\hfuzz=500pt\includegraphics[width=1em]{element.pdf}~minDegree & \hfuzz=500pt uint & \hfuzz=500pt \\
\hfuzz=500pt\includegraphics[width=1em]{element.pdf}~maxDegree & \hfuzz=500pt uint & \hfuzz=500pt \\
\hfuzz=500pt\includegraphics[width=1em]{element.pdf}~GM & \hfuzz=500pt double & \hfuzz=500pt Geocentric gravitational constant\\
\hfuzz=500pt\includegraphics[width=1em]{element.pdf}~R & \hfuzz=500pt double & \hfuzz=500pt reference radius\\
\hfuzz=500pt\includegraphics[width=1em]{element.pdf}~time & \hfuzz=500pt time & \hfuzz=500pt at this time the gravity field will be evaluated\\
\hline
\end{tabularx}

\clearpage
%==================================
\subsection{GravityfieldVariancesPropagation2GriddedData}\label{GravityfieldVariancesPropagation2GriddedData}
This program propagates variance-covariance matrix of a \configClass{gravityfield}{gravityfieldType}
evaluated at \config{time} to the points of a \configClass{grid}{gridType} in terms of the functional
given by \configClass{kernel}{kernelType}.
The resulting \file{outputfileGriddedData}{griddedData} contains the standard deviations of the grid
points.

See also \program{Gravityfield2GridCovarianceMatrix}, \program{GravityfieldCovariancesPropagation2GriddedData}.


\keepXColumns
\begin{tabularx}{\textwidth}{N T A}
\hline
Name & Type & Annotation\\
\hline
\hfuzz=500pt\includegraphics[width=1em]{element-mustset.pdf}~outputfileGriddedData & \hfuzz=500pt filename & \hfuzz=500pt standard deviation at each grid point\\
\hfuzz=500pt\includegraphics[width=1em]{element-mustset-unbounded.pdf}~grid & \hfuzz=500pt \hyperref[gridType]{grid} & \hfuzz=500pt \\
\hfuzz=500pt\includegraphics[width=1em]{element-mustset.pdf}~kernel & \hfuzz=500pt \hyperref[kernelType]{kernel} & \hfuzz=500pt functional\\
\hfuzz=500pt\includegraphics[width=1em]{element-mustset-unbounded.pdf}~gravityfield & \hfuzz=500pt \hyperref[gravityfieldType]{gravityfield} & \hfuzz=500pt \\
\hfuzz=500pt\includegraphics[width=1em]{element.pdf}~time & \hfuzz=500pt time & \hfuzz=500pt at this time the gravity field will be evaluated\\
\hfuzz=500pt\includegraphics[width=1em]{element.pdf}~R & \hfuzz=500pt double & \hfuzz=500pt reference radius for ellipsoidal coordinates on output\\
\hfuzz=500pt\includegraphics[width=1em]{element.pdf}~inverseFlattening & \hfuzz=500pt double & \hfuzz=500pt reference flattening for ellipsoidal coordinates on output, 0: spherical coordinates\\
\hline
\end{tabularx}

This program is \reference{parallelized}{general.parallelization}.
\clearpage
%==================================
\section{Programs: Grid}
\subsection{GriddedData2AreaMeanTimeSeries}\label{GriddedData2AreaMeanTimeSeries}
This program computes a time series of area mean values
in a basin represented by \configClass{border}{borderType} from a sequence of grid files.
If a file is not found, the epoch is skipped. Per default
the weighted average of all points in the given border is computed where the points are weighted by their area element.

If \config{computeMean} is set, the time average of each grid points is subtracted before the computation.
If \config{multiplyWithArea} is set, the weighted average is multiplied with the total basin area.
This is useful for computing the total mass in the basin.

The \configFile{outputfileTimeSeries}{instrument} is an instrument file, where the first columns are the
mean value each data column in the grid files, followed by the the weighted RMS
for each data column in the grid files if \config{computeRms} is set.
If the number of data columns differs between the grid files, the remaining columns are padded with zeros.

See also \program{Gravityfield2AreaMeanTimeSeries}.


\keepXColumns
\begin{tabularx}{\textwidth}{N T A}
\hline
Name & Type & Annotation\\
\hline
\hfuzz=500pt\includegraphics[width=1em]{element-mustset.pdf}~outputfileTimeSeries & \hfuzz=500pt filename & \hfuzz=500pt \\
\hfuzz=500pt\includegraphics[width=1em]{element-mustset-unbounded.pdf}~inputfileGriddedData & \hfuzz=500pt filename & \hfuzz=500pt \\
\hfuzz=500pt\includegraphics[width=1em]{element-mustset-unbounded.pdf}~border & \hfuzz=500pt \hyperref[borderType]{border} & \hfuzz=500pt \\
\hfuzz=500pt\includegraphics[width=1em]{element-mustset-unbounded.pdf}~timeSeries & \hfuzz=500pt \hyperref[timeSeriesType]{timeSeries} & \hfuzz=500pt \\
\hfuzz=500pt\includegraphics[width=1em]{element.pdf}~multiplyWithArea & \hfuzz=500pt boolean & \hfuzz=500pt multiply time series with total area (useful for mass estimates)\\
\hfuzz=500pt\includegraphics[width=1em]{element.pdf}~removeMean & \hfuzz=500pt boolean & \hfuzz=500pt remove the temporal mean of the series\\
\hfuzz=500pt\includegraphics[width=1em]{element.pdf}~computeRms & \hfuzz=500pt boolean & \hfuzz=500pt additional rms each time step\\
\hline
\end{tabularx}

\clearpage
%==================================
\subsection{GriddedData2GriddedDataStatistics}\label{GriddedData2GriddedDataStatistics}
This program assigns values \configFile{inputfileGriddedData}{griddedData} to the nearest points
of a new \configClass{grid}{gridType}. If some of the new points are not filled in with data
\config{emptyValue} is used instead. If multiple points of the input fall on the same node
the result can be selected with \config{statistics} (e.g. mean, root mean square, min, max, \ldots).
It also is possible to simply count the number of data points that were assigned to each point.

Be aware in case borders are given within \configClass{grid}{gridType}, the \configFile{outputfileGriddedData}{griddedData} will have points excluded before the assignement of old points to the new points.
The data from \configFile{inputfileGriddedData}{griddedData} will not be limited by the given borders! See \reference{GriddedDataConcatenate}{GriddedDataConcatenate} to limit the
\configFile{inputfileGriddedData}{griddedData} to given borders.

\fig{!hb}{0.8}{griddedData2GriddedDataStatistics}{fig:griddedData2GriddedDataStatistics}{Assignement of irregular distributed data to grid.}


\keepXColumns
\begin{tabularx}{\textwidth}{N T A}
\hline
Name & Type & Annotation\\
\hline
\hfuzz=500pt\includegraphics[width=1em]{element-mustset.pdf}~outputfileGriddedData & \hfuzz=500pt filename & \hfuzz=500pt \\
\hfuzz=500pt\includegraphics[width=1em]{element-mustset.pdf}~inputfileGriddedData & \hfuzz=500pt filename & \hfuzz=500pt \\
\hfuzz=500pt\includegraphics[width=1em]{element-mustset-unbounded.pdf}~grid & \hfuzz=500pt \hyperref[gridType]{grid} & \hfuzz=500pt \\
\hfuzz=500pt\includegraphics[width=1em]{element-mustset.pdf}~statistic & \hfuzz=500pt choice & \hfuzz=500pt statistic used if multiple values fall on the same cell\\
\hfuzz=500pt\includegraphics[width=1em]{connector.pdf}\includegraphics[width=1em]{element-mustset.pdf}~mean & \hfuzz=500pt  & \hfuzz=500pt mean\\
\hfuzz=500pt\includegraphics[width=1em]{connector.pdf}\includegraphics[width=1em]{element-mustset.pdf}~wmean & \hfuzz=500pt  & \hfuzz=500pt area weighted mean\\
\hfuzz=500pt\includegraphics[width=1em]{connector.pdf}\includegraphics[width=1em]{element-mustset.pdf}~rms & \hfuzz=500pt  & \hfuzz=500pt root mean square\\
\hfuzz=500pt\includegraphics[width=1em]{connector.pdf}\includegraphics[width=1em]{element-mustset.pdf}~wrms & \hfuzz=500pt  & \hfuzz=500pt area weighted root mean square\\
\hfuzz=500pt\includegraphics[width=1em]{connector.pdf}\includegraphics[width=1em]{element-mustset.pdf}~std & \hfuzz=500pt  & \hfuzz=500pt standard deviation\\
\hfuzz=500pt\includegraphics[width=1em]{connector.pdf}\includegraphics[width=1em]{element-mustset.pdf}~wstd & \hfuzz=500pt  & \hfuzz=500pt area weighted standard deviation\\
\hfuzz=500pt\includegraphics[width=1em]{connector.pdf}\includegraphics[width=1em]{element-mustset.pdf}~sum & \hfuzz=500pt  & \hfuzz=500pt sum\\
\hfuzz=500pt\includegraphics[width=1em]{connector.pdf}\includegraphics[width=1em]{element-mustset.pdf}~min & \hfuzz=500pt  & \hfuzz=500pt minimum value\\
\hfuzz=500pt\includegraphics[width=1em]{connector.pdf}\includegraphics[width=1em]{element-mustset.pdf}~max & \hfuzz=500pt  & \hfuzz=500pt maximum value\\
\hfuzz=500pt\includegraphics[width=1em]{connector.pdf}\includegraphics[width=1em]{element-mustset.pdf}~count & \hfuzz=500pt  & \hfuzz=500pt number of values\\
\hfuzz=500pt\includegraphics[width=1em]{connector.pdf}\includegraphics[width=1em]{element-mustset.pdf}~first & \hfuzz=500pt  & \hfuzz=500pt first value\\
\hfuzz=500pt\includegraphics[width=1em]{connector.pdf}\includegraphics[width=1em]{element-mustset.pdf}~last & \hfuzz=500pt  & \hfuzz=500pt last value\\
\hfuzz=500pt\includegraphics[width=1em]{element.pdf}~emptyValue & \hfuzz=500pt double & \hfuzz=500pt value for nodes without data\\
\hfuzz=500pt\includegraphics[width=1em]{element.pdf}~R & \hfuzz=500pt double & \hfuzz=500pt reference radius for ellipsoidal coordinates\\
\hfuzz=500pt\includegraphics[width=1em]{element.pdf}~inverseFlattening & \hfuzz=500pt double & \hfuzz=500pt reference flattening for ellipsoidal coordinates\\
\hline
\end{tabularx}

\clearpage
%==================================
\subsection{GriddedData2GriddedDataTimeSeries}\label{GriddedData2GriddedDataTimeSeries}
Write a series of \configFile{inputfileGriddedData}{griddedData}
with the corresponding \configClass{timeSeries}{timeSeriesType}
as a single \file{gridded data time series file}{griddedDataTimeSeries}.
The \config{splineDegree} defines the possible temporal interpolation of data in the output file.
For a file with spline degree 0 (temporal block means) the time intervals
in which the grids are valid are defined between adjacent points in time.
Therefore one more point in time is needed than the number of input grid files for degree 0.

See also \program{GriddedDataTimeSeries2GriddedData}.


\keepXColumns
\begin{tabularx}{\textwidth}{N T A}
\hline
Name & Type & Annotation\\
\hline
\hfuzz=500pt\includegraphics[width=1em]{element-mustset.pdf}~outputfileGriddedDataTimeSeries & \hfuzz=500pt filename & \hfuzz=500pt \\
\hfuzz=500pt\includegraphics[width=1em]{element-mustset-unbounded.pdf}~inputfileGriddedData & \hfuzz=500pt filename & \hfuzz=500pt file count must agree with number of times+splineDegre-1\\
\hfuzz=500pt\includegraphics[width=1em]{element-mustset-unbounded.pdf}~timeSeries & \hfuzz=500pt \hyperref[timeSeriesType]{timeSeries} & \hfuzz=500pt \\
\hfuzz=500pt\includegraphics[width=1em]{element.pdf}~splineDegree & \hfuzz=500pt uint & \hfuzz=500pt degree of splines\\
\hline
\end{tabularx}

\clearpage
%==================================
\subsection{GriddedData2Matrix}\label{GriddedData2Matrix}
This program converts \configFile{inputfileGriddedData}{griddedData}
to \configFile{outputfileMatrix}{matrix} with data columns.
The grid is expressed as ellipsoidal coordinates (longitude, latitude, height)
based on a reference ellipsoid with parameters \config{R} and \config{inverseFlattening}.
The content of the output matrix can be controlled by \config{outColumn} expressions
applied to every grid point. The common data variables for grids are available,
see \reference{dataVariables}{general.parser:dataVariables}.


\keepXColumns
\begin{tabularx}{\textwidth}{N T A}
\hline
Name & Type & Annotation\\
\hline
\hfuzz=500pt\includegraphics[width=1em]{element-mustset.pdf}~outputfileMatrix & \hfuzz=500pt filename & \hfuzz=500pt point list as matrix with longitude and latitude values in columns and possible additional columns\\
\hfuzz=500pt\includegraphics[width=1em]{element-mustset.pdf}~inputfileGriddedData & \hfuzz=500pt filename & \hfuzz=500pt \\
\hfuzz=500pt\includegraphics[width=1em]{element.pdf}~R & \hfuzz=500pt double & \hfuzz=500pt reference radius for ellipsoidal coordinates on output\\
\hfuzz=500pt\includegraphics[width=1em]{element.pdf}~inverseFlattening & \hfuzz=500pt double & \hfuzz=500pt reference flattening for ellipsoidal coordinates on output, 0: spherical coordinates\\
\hfuzz=500pt\includegraphics[width=1em]{element-unbounded.pdf}~outColumn & \hfuzz=500pt expression & \hfuzz=500pt expression (variables: longitude, latitude, height, area, data0, data1, ...)\\
\hline
\end{tabularx}

\clearpage
%==================================
\subsection{GriddedData2PotentialCoefficients}\label{GriddedData2PotentialCoefficients}
This program estimate potential coefficients from \configFile{inputfileGriddedData}{griddedData}
gravity field functionals. It used a simple quadrature formular
\begin{equation}
  c_{nm} = \frac{1}{4\pi}\frac{R}{GM} \sum_i f_i \left(\frac{r_i}{R}\right)^{n+1} k_n C_{nm}(\lambda_i,\vartheta_i)\,\Delta\Phi_i
\end{equation}
or a \config{leastSquares} adjustment with block diagonal normal matrix (order by order).
For the latter one the data must be regular distributed.

The \config{value}s $f_i$ and the \config{weight}s $\Delta\Phi_i$ are expressions
using the common data variables for grids, see \reference{dataVariables}{general.parser:dataVariables}.
Multiple \configFile{outputfilePotentialCoefficients}{potentialCoefficients} can be estimated in one step.
For each an indivdual \config{value} must be specified.
The type of the gridded data (e.g gravity anomalies or geoid heights)
must be set with \configClass{kernel}{kernelType} $k_n$.

The expansion is limited in the range between \config{minDegree}
and \config{maxDegree} inclusively. The coefficients are related
to the reference radius~\config{R} and the Earth gravitational constant \config{GM}.

For irregular distributed data and using the full variance covariance matrix use
\program{NormalsSolverVCE} together with \configClass{oberservation:terrestrial}{observationType:terrestrial}
and \configClass{parametrizationGravity:sphericalHarmonics}{parametrizationGravityType:sphericalHarmonics}.

See also \program{GriddedDataTimeSeries2PotentialCoefficients}.


\keepXColumns
\begin{tabularx}{\textwidth}{N T A}
\hline
Name & Type & Annotation\\
\hline
\hfuzz=500pt\includegraphics[width=1em]{element-mustset-unbounded.pdf}~outputfilePotentialCoefficients & \hfuzz=500pt filename & \hfuzz=500pt one file for each value expression\\
\hfuzz=500pt\includegraphics[width=1em]{element-mustset.pdf}~inputfileGriddedData & \hfuzz=500pt filename & \hfuzz=500pt \\
\hfuzz=500pt\includegraphics[width=1em]{element-mustset-unbounded.pdf}~value & \hfuzz=500pt expression & \hfuzz=500pt expression to compute values (input columns are named data0, data1, ...)\\
\hfuzz=500pt\includegraphics[width=1em]{element-mustset.pdf}~weight & \hfuzz=500pt expression & \hfuzz=500pt expression to compute values (input columns are named data0, data1, ...)\\
\hfuzz=500pt\includegraphics[width=1em]{element-mustset.pdf}~kernel & \hfuzz=500pt \hyperref[kernelType]{kernel} & \hfuzz=500pt data type of input values\\
\hfuzz=500pt\includegraphics[width=1em]{element.pdf}~minDegree & \hfuzz=500pt uint & \hfuzz=500pt \\
\hfuzz=500pt\includegraphics[width=1em]{element-mustset.pdf}~maxDegree & \hfuzz=500pt uint & \hfuzz=500pt \\
\hfuzz=500pt\includegraphics[width=1em]{element.pdf}~GM & \hfuzz=500pt double & \hfuzz=500pt Geocentric gravitational constant\\
\hfuzz=500pt\includegraphics[width=1em]{element.pdf}~R & \hfuzz=500pt double & \hfuzz=500pt reference radius\\
\hfuzz=500pt\includegraphics[width=1em]{element.pdf}~leastSquares & \hfuzz=500pt boolean & \hfuzz=500pt false: quadrature formular, true: least squares adjustment order by order\\
\hline
\end{tabularx}

This program is \reference{parallelized}{general.parallelization}.
\clearpage
%==================================
\subsection{GriddedData2SphericalDistance}\label{GriddedData2SphericalDistance}
Compute the spherical distance on the unit sphere in radians between all point pairs of two grids.
The spherical distance is computed by
\begin{equation}
  \psi_{12} = \arccos(\M n_1 \cdot \M n_2),
\end{equation}
where $\M n_i$ is the (normalized) position. This implies that all points are projected onto the unit sphere.


\keepXColumns
\begin{tabularx}{\textwidth}{N T A}
\hline
Name & Type & Annotation\\
\hline
\hfuzz=500pt\includegraphics[width=1em]{element-mustset.pdf}~outputfileMatrix & \hfuzz=500pt filename & \hfuzz=500pt matrix containing the spherical distance between all point pairs [rad]\\
\hfuzz=500pt\includegraphics[width=1em]{element-mustset-unbounded.pdf}~grid1 & \hfuzz=500pt \hyperref[gridType]{grid} & \hfuzz=500pt \\
\hfuzz=500pt\includegraphics[width=1em]{element-mustset-unbounded.pdf}~grid2 & \hfuzz=500pt \hyperref[gridType]{grid} & \hfuzz=500pt \\
\hline
\end{tabularx}

\clearpage
%==================================
\subsection{GriddedData2TimeSeries}\label{GriddedData2TimeSeries}
Write a series of \configFile{inputfileGriddedData}{griddedData} with the corresponding
\configClass{timeSeries}{timeSeriesType} as a single time series file
(\file{instrument}{instrument}, MISCVALUES).

If \config{groupDataByPoints} is true the \config{outputfileTimeSeries} starts
for each epoch with all data (\verb|data0|, \verb|data1|\ldots) for the first point,
followed by all data of the second point and so on.
If \config{groupDataByPoints} is false, the file starts with \verb|data0|
for all points, followed by all \verb|data1| and so on.

This enables the use of all instrument programs like \program{InstrumentFilter} or
\program{InstrumentDetrend} to analyze and manipulate time series of gridded data.

See also \program{TimeSeries2GriddedData}.


\keepXColumns
\begin{tabularx}{\textwidth}{N T A}
\hline
Name & Type & Annotation\\
\hline
\hfuzz=500pt\includegraphics[width=1em]{element-mustset.pdf}~outputfileTimeSeries & \hfuzz=500pt filename & \hfuzz=500pt each epoch: multiple data for points (MISCVALUES)\\
\hfuzz=500pt\includegraphics[width=1em]{element-mustset-unbounded.pdf}~inputfileGriddedData & \hfuzz=500pt filename & \hfuzz=500pt file count must agree with number of times\\
\hfuzz=500pt\includegraphics[width=1em]{element-mustset-unbounded.pdf}~timeSeries & \hfuzz=500pt \hyperref[timeSeriesType]{timeSeries} & \hfuzz=500pt \\
\hfuzz=500pt\includegraphics[width=1em]{element.pdf}~groupDataByPoints & \hfuzz=500pt boolean & \hfuzz=500pt multiple data are given point by point, otherwise: data0 for all points, followed by all data1\\
\hline
\end{tabularx}

\clearpage
%==================================
\subsection{GriddedDataCalculate}\label{GriddedDataCalculate}
This program manipulates \file{grid files}{griddedData} with data in columns similar to
\program{FunctionsCalculate}, see there for more details.
If several \config{inputfile}s are given the data columns are copied side by side.
All \config{inputfile}s must contain the same grid points.
The columns are enumerated by \verb|data0|,~\verb|data1|,~\ldots.

The content of \configFile{outputfileGriddedData}{griddedData} is controlled by \config{outColumn}.
The algorithm to compute the output is as follows:
The expressions in \config{outColumn} are evaluated once for each grid point of the input.
The variables \verb|data0|,~\verb|data1|,~\ldots are replaced by the according values
from the input columns before.
Additional variables are available, e.g. \verb|index|, \verb|data0rms|,
see~\reference{dataVariables}{general.parser:dataVariables}.

For a simplified handling \config{constant}s can be defined by \verb|name=value|.
It is also possible to estimate \config{parameter}s in a least squares adjustment.
The \config{leastSquares} serves as template for observation equations for every point.
The expression \config{leastSquares} is evaluated for each grid point.
The variables \verb|data0|,~\verb|data1|,~\ldots are replaced by the according values from the input columns before.
In the next step the parameters are estimated in order to minimize the expressions in \config{leastSquares}
in the sense of least squares.

Afterwards grid points are removed if one of the \config{removalCriteria} expressions
for this grid point evaluates true (not zero).

An extra \configFile{statistics:outputfile}{matrix} can be generated with one row of data.
For the computation of the \config{outColumn} values
all~\reference{dataVariables}{general.parser:dataVariables} are available
(e.g. \verb|data3mean|, \verb|data4std|) inclusively the \config{constant}s and
estimated \config{parameter}s but without the \verb|data0|,~\verb|data1|,~\ldots itself.
The variables and the numbering of the columns refers to the \configFile{outputfileGriddedData}{griddedData}.

See also \program{FunctionsCalculate}, \program{InstrumentArcCalculate}, \program{MatrixCalculate}.


\keepXColumns
\begin{tabularx}{\textwidth}{N T A}
\hline
Name & Type & Annotation\\
\hline
\hfuzz=500pt\includegraphics[width=1em]{element.pdf}~outputfileGriddedData & \hfuzz=500pt filename & \hfuzz=500pt \\
\hfuzz=500pt\includegraphics[width=1em]{element-mustset-unbounded.pdf}~inputfileGriddedData & \hfuzz=500pt filename & \hfuzz=500pt \\
\hfuzz=500pt\includegraphics[width=1em]{element-unbounded.pdf}~constant & \hfuzz=500pt expression & \hfuzz=500pt define a constant by name=value\\
\hfuzz=500pt\includegraphics[width=1em]{element-unbounded.pdf}~parameter & \hfuzz=500pt expression & \hfuzz=500pt define a parameter by name[=value]\\
\hfuzz=500pt\includegraphics[width=1em]{element-unbounded.pdf}~leastSquares & \hfuzz=500pt expression & \hfuzz=500pt try to minimize the expression by adjustment of the parameters\\
\hfuzz=500pt\includegraphics[width=1em]{element-unbounded.pdf}~removalCriteria & \hfuzz=500pt expression & \hfuzz=500pt points are removed if one criterion evaluates true. data0 is the first data field.\\
\hfuzz=500pt\includegraphics[width=1em]{element-mustset.pdf}~longitude & \hfuzz=500pt expression & \hfuzz=500pt expression\\
\hfuzz=500pt\includegraphics[width=1em]{element-mustset.pdf}~latitude & \hfuzz=500pt expression & \hfuzz=500pt expression\\
\hfuzz=500pt\includegraphics[width=1em]{element-mustset.pdf}~height & \hfuzz=500pt expression & \hfuzz=500pt expression\\
\hfuzz=500pt\includegraphics[width=1em]{element.pdf}~area & \hfuzz=500pt expression & \hfuzz=500pt expression: e.g. deltaL * 2.0 * sin(deltaB/2.0) * cos(latitude/rho)\\
\hfuzz=500pt\includegraphics[width=1em]{element-unbounded.pdf}~value & \hfuzz=500pt expression & \hfuzz=500pt expression to compute values (input columns are named data0, data1, ...)\\
\hfuzz=500pt\includegraphics[width=1em]{element.pdf}~computeArea & \hfuzz=500pt boolean & \hfuzz=500pt automatically area computation of rectangular grids (overwrite area)\\
\hfuzz=500pt\includegraphics[width=1em]{element.pdf}~R & \hfuzz=500pt double & \hfuzz=500pt reference radius for ellipsoidal coordinates\\
\hfuzz=500pt\includegraphics[width=1em]{element.pdf}~inverseFlattening & \hfuzz=500pt double & \hfuzz=500pt reference flattening for ellipsoidal coordinates\\
\hfuzz=500pt\includegraphics[width=1em]{element.pdf}~statistics & \hfuzz=500pt sequence & \hfuzz=500pt \\
\hfuzz=500pt\includegraphics[width=1em]{connector.pdf}\includegraphics[width=1em]{element-mustset.pdf}~outputfile & \hfuzz=500pt filename & \hfuzz=500pt matrix with one row, columns are user defined\\
\hfuzz=500pt\includegraphics[width=1em]{connector.pdf}\includegraphics[width=1em]{element-mustset-unbounded.pdf}~outColumn & \hfuzz=500pt expression & \hfuzz=500pt expression to compute statistics columns, data* are the outputColumns\\
\hline
\end{tabularx}

\clearpage
%==================================
\subsection{GriddedDataConcatenate}\label{GriddedDataConcatenate}
This program concatenate grid from several \configFile{inputfileGriddedData}{griddedData}
and write it to a new \configFile{outputfileGriddedData}{griddedData}.
Input files must have the same number of data columns.
If \config{sort} is enabled, the points are sorted by latitudes starting from north/west to south east.
Identical points (within a \config{margin}) can be removed with \config{removeDuplicates}.


\keepXColumns
\begin{tabularx}{\textwidth}{N T A}
\hline
Name & Type & Annotation\\
\hline
\hfuzz=500pt\includegraphics[width=1em]{element-mustset.pdf}~outputfileGriddedData & \hfuzz=500pt filename & \hfuzz=500pt \\
\hfuzz=500pt\includegraphics[width=1em]{element-mustset-unbounded.pdf}~inputfileGriddedData & \hfuzz=500pt filename & \hfuzz=500pt \\
\hfuzz=500pt\includegraphics[width=1em]{element-unbounded.pdf}~border & \hfuzz=500pt \hyperref[borderType]{border} & \hfuzz=500pt \\
\hfuzz=500pt\includegraphics[width=1em]{element.pdf}~sortPoints & \hfuzz=500pt boolean & \hfuzz=500pt sort from north/west to south east\\
\hfuzz=500pt\includegraphics[width=1em]{element.pdf}~removeDuplicates & \hfuzz=500pt choice & \hfuzz=500pt remove duplicate points\\
\hfuzz=500pt\includegraphics[width=1em]{connector.pdf}\includegraphics[width=1em]{element-mustset.pdf}~keepFirst & \hfuzz=500pt sequence & \hfuzz=500pt keep first point, remove all other identicals\\
\hfuzz=500pt\quad\includegraphics[width=1em]{connector.pdf}\includegraphics[width=1em]{element.pdf}~margin & \hfuzz=500pt double & \hfuzz=500pt margin distance for identical points [m]\\
\hfuzz=500pt\includegraphics[width=1em]{connector.pdf}\includegraphics[width=1em]{element-mustset.pdf}~keepLast & \hfuzz=500pt sequence & \hfuzz=500pt keep last point, remove all other identicals\\
\hfuzz=500pt\quad\includegraphics[width=1em]{connector.pdf}\includegraphics[width=1em]{element.pdf}~margin & \hfuzz=500pt double & \hfuzz=500pt margin distance for identical points [m]\\
\hfuzz=500pt\includegraphics[width=1em]{element.pdf}~R & \hfuzz=500pt double & \hfuzz=500pt reference radius for ellipsoidal coordinates\\
\hfuzz=500pt\includegraphics[width=1em]{element.pdf}~inverseFlattening & \hfuzz=500pt double & \hfuzz=500pt reference flattening for ellipsoidal coordinates\\
\hline
\end{tabularx}

\clearpage
%==================================
\subsection{GriddedDataCreate}\label{GriddedDataCreate}
This program creates a \configClass{grid}{gridType} and writes it to \configFile{outputfileGrid}{griddedData}.
The grid is expressed as ellipsoidal coordinates (longitude, latitude, height)
based on a reference ellipsoid with parameters \config{R} and \config{inverseFlattening}.
Extra \config{value} columns can be appended using expressions
with the common \reference{data variables}{general.parser:dataVariables} for gridded data.


\keepXColumns
\begin{tabularx}{\textwidth}{N T A}
\hline
Name & Type & Annotation\\
\hline
\hfuzz=500pt\includegraphics[width=1em]{element-mustset.pdf}~outputfileGrid & \hfuzz=500pt filename & \hfuzz=500pt \\
\hfuzz=500pt\includegraphics[width=1em]{element-mustset-unbounded.pdf}~grid & \hfuzz=500pt \hyperref[gridType]{grid} & \hfuzz=500pt \\
\hfuzz=500pt\includegraphics[width=1em]{element.pdf}~R & \hfuzz=500pt double & \hfuzz=500pt reference radius for ellipsoidal coordinates on output\\
\hfuzz=500pt\includegraphics[width=1em]{element.pdf}~inverseFlattening & \hfuzz=500pt double & \hfuzz=500pt reference flattening for ellipsoidal coordinates on output, 0: spherical coordinates\\
\hfuzz=500pt\includegraphics[width=1em]{element-unbounded.pdf}~value & \hfuzz=500pt expression & \hfuzz=500pt expression (variables as 'longitude', 'height', 'area' are taken from the gridded data)\\
\hline
\end{tabularx}

\clearpage
%==================================
\subsection{GriddedDataInterpolate}\label{GriddedDataInterpolate}
Interpolate values of a regular rectangular \configFile{inputfileGriddedData}{griddedData}
to new points given by \configClass{grid}{gridType} and write as \configFile{outputfileGriddedData}{griddedData}.
Only longitude and latitude of points are considered; the height is ignored for interpolation.

(Only nearest neighbor method is implemented at the moment.)

\fig{!hb}{0.8}{griddedDataInterpolate}{fig:griddedDataInterpolate}{Interpolation of point data from rectangular gridded data.}


\keepXColumns
\begin{tabularx}{\textwidth}{N T A}
\hline
Name & Type & Annotation\\
\hline
\hfuzz=500pt\includegraphics[width=1em]{element-mustset.pdf}~outputfileGriddedData & \hfuzz=500pt filename & \hfuzz=500pt \\
\hfuzz=500pt\includegraphics[width=1em]{element-mustset.pdf}~inputfileGriddedData & \hfuzz=500pt filename & \hfuzz=500pt must be rectangular\\
\hfuzz=500pt\includegraphics[width=1em]{element-mustset-unbounded.pdf}~grid & \hfuzz=500pt \hyperref[gridType]{grid} & \hfuzz=500pt \\
\hfuzz=500pt\includegraphics[width=1em]{element-mustset.pdf}~method & \hfuzz=500pt choice & \hfuzz=500pt \\
\hfuzz=500pt\includegraphics[width=1em]{connector.pdf}\includegraphics[width=1em]{element-mustset.pdf}~nearestNeighbor & \hfuzz=500pt  & \hfuzz=500pt \\
\hline
\end{tabularx}

\clearpage
%==================================
\subsection{GriddedDataReduceSampling}\label{GriddedDataReduceSampling}
Generate coarse grid by computing area weighted mean values.
The number of points is decimated by averaging integer multiplies of grid points
(\config{multiplierLongitude}, \config{multiplierLatitude}).

if \config{volumeConserving} is set, data are interpreted as heights above ellipsoid
and the tesseroid volume
\begin{equation}
  V=\int_r^{r+H}\int_{\varphi_1}^{\varphi_2}\int_{\lambda_1}^{\lambda_2} r^2\cos\varphi\,d\varphi\,d\lambda\,dr
\end{equation}
is conserved, where $r$ is the radius of the ellipsoid at grid center and
$(\varphi_1-\varphi_2)\times(\lambda_1-\lambda_2)$ are the grid cell boundaries.
This is meaninful for Digital Elevation Models (DEM).

The fine grid can be written, where the first coarse grid values (data0) are additionally appended.


\keepXColumns
\begin{tabularx}{\textwidth}{N T A}
\hline
Name & Type & Annotation\\
\hline
\hfuzz=500pt\includegraphics[width=1em]{element.pdf}~outputfileCoarseGridRectangular & \hfuzz=500pt filename & \hfuzz=500pt coarse grid\\
\hfuzz=500pt\includegraphics[width=1em]{element.pdf}~outputfileFineGridRectangular & \hfuzz=500pt filename & \hfuzz=500pt fine grid with additional coarse grid values\\
\hfuzz=500pt\includegraphics[width=1em]{element-mustset.pdf}~inputfileFineGridRectangular & \hfuzz=500pt filename & \hfuzz=500pt Digital Terrain Model\\
\hfuzz=500pt\includegraphics[width=1em]{element-mustset.pdf}~multiplierLongitude & \hfuzz=500pt uint & \hfuzz=500pt Generalizing factor\\
\hfuzz=500pt\includegraphics[width=1em]{element-mustset.pdf}~multiplierLatitude & \hfuzz=500pt uint & \hfuzz=500pt Generalizing factor\\
\hfuzz=500pt\includegraphics[width=1em]{element.pdf}~volumeConserving & \hfuzz=500pt boolean & \hfuzz=500pt data are interpreted as heights above ellipsoid\\
\hline
\end{tabularx}

\clearpage
%==================================
\subsection{GriddedDataTimeSeries2GriddedData}\label{GriddedDataTimeSeries2GriddedData}
Read a \configFile{inputfileGriddedDataTimeSeries}{griddedDataTimeSeries}
and write for each epoch a \file{gridded data file}{griddedData} where
the \config{variableLoopTime} and \config{variableLoopIndex} are expanded for
each point of the given \configClass{timeSeries}{timeSeriesType}
to create the file name for this epoch (see \reference{text parser}{general.parser:text}).

If \configClass{timeSeries}{timeSeriesType} is not set
the temporal nodal points from the inputfile are used.

See also \program{GriddedData2GriddedDataTimeSeries}.


\keepXColumns
\begin{tabularx}{\textwidth}{N T A}
\hline
Name & Type & Annotation\\
\hline
\hfuzz=500pt\includegraphics[width=1em]{element-mustset.pdf}~outputfilesGriddedData & \hfuzz=500pt filename & \hfuzz=500pt for each epoch\\
\hfuzz=500pt\includegraphics[width=1em]{element.pdf}~variableLoopTime & \hfuzz=500pt string & \hfuzz=500pt variable with time of each epoch\\
\hfuzz=500pt\includegraphics[width=1em]{element.pdf}~variableLoopIndex & \hfuzz=500pt string & \hfuzz=500pt variable with index of current epoch (starts with zero)\\
\hfuzz=500pt\includegraphics[width=1em]{element.pdf}~variableLoopCount & \hfuzz=500pt string & \hfuzz=500pt variable with total number of epochs\\
\hfuzz=500pt\includegraphics[width=1em]{element-mustset.pdf}~inputfileGriddedDataTimeSeries & \hfuzz=500pt filename & \hfuzz=500pt \\
\hfuzz=500pt\includegraphics[width=1em]{element-unbounded.pdf}~timeSeries & \hfuzz=500pt \hyperref[timeSeriesType]{timeSeries} & \hfuzz=500pt otherwise times from inputfile are used\\
\hline
\end{tabularx}

\clearpage
%==================================
\subsection{GriddedDataTimeSeries2PotentialCoefficients}\label{GriddedDataTimeSeries2PotentialCoefficients}
This program estimate potential coefficients from
\configFile{inputfileGriddedDataTimeSeries}{griddedDataTimeSeries}
in the same way as \program{GriddedData2PotentialCoefficients}
but not only for one grid but for each epoch of
\configClass{timeSeries}{timeSeriesType} of if not set
for the temporal nodal points from the inputfile.
The \configFile{outputfilePotentialCoefficients}{potentialCoefficients}
(one for each \config{value}) are written for each epoch with the expansion
of \config{variableLoopTime} and \config{variableLoopIndex}
(see \reference{text parser}{general.parser:text}).

See also \program{GriddedData2PotentialCoefficients}.


\keepXColumns
\begin{tabularx}{\textwidth}{N T A}
\hline
Name & Type & Annotation\\
\hline
\hfuzz=500pt\includegraphics[width=1em]{element-mustset-unbounded.pdf}~outputfilesPotentialCoefficients & \hfuzz=500pt filename & \hfuzz=500pt for each epoch\\
\hfuzz=500pt\includegraphics[width=1em]{element.pdf}~variableLoopTime & \hfuzz=500pt string & \hfuzz=500pt variable with time of each epoch\\
\hfuzz=500pt\includegraphics[width=1em]{element.pdf}~variableLoopIndex & \hfuzz=500pt string & \hfuzz=500pt variable with index of current epoch (starts with zero)\\
\hfuzz=500pt\includegraphics[width=1em]{element.pdf}~variableLoopCount & \hfuzz=500pt string & \hfuzz=500pt variable with total number of epochs\\
\hfuzz=500pt\includegraphics[width=1em]{element-mustset.pdf}~inputfileGriddedDataTimeSeries & \hfuzz=500pt filename & \hfuzz=500pt \\
\hfuzz=500pt\includegraphics[width=1em]{element-unbounded.pdf}~timeSeries & \hfuzz=500pt \hyperref[timeSeriesType]{timeSeries} & \hfuzz=500pt otherwise times from inputfile are used\\
\hfuzz=500pt\includegraphics[width=1em]{element-mustset-unbounded.pdf}~value & \hfuzz=500pt expression & \hfuzz=500pt expression (variables: longitude, latitude, height, area, data0, data1, ...)\\
\hfuzz=500pt\includegraphics[width=1em]{element-mustset.pdf}~weight & \hfuzz=500pt expression & \hfuzz=500pt expression to compute values (input columns are named data0, data1, ...)\\
\hfuzz=500pt\includegraphics[width=1em]{element-mustset.pdf}~kernel & \hfuzz=500pt \hyperref[kernelType]{kernel} & \hfuzz=500pt kernel in which the grid values are given\\
\hfuzz=500pt\includegraphics[width=1em]{element.pdf}~minDegree & \hfuzz=500pt uint & \hfuzz=500pt \\
\hfuzz=500pt\includegraphics[width=1em]{element-mustset.pdf}~maxDegree & \hfuzz=500pt uint & \hfuzz=500pt \\
\hfuzz=500pt\includegraphics[width=1em]{element.pdf}~GM & \hfuzz=500pt double & \hfuzz=500pt Geocentric gravitational constant\\
\hfuzz=500pt\includegraphics[width=1em]{element.pdf}~R & \hfuzz=500pt double & \hfuzz=500pt reference radius for potential coefficients\\
\hfuzz=500pt\includegraphics[width=1em]{element.pdf}~leastSquares & \hfuzz=500pt boolean & \hfuzz=500pt false: quadrature formular, true: least squares adjustment order by order\\
\hline
\end{tabularx}

This program is \reference{parallelized}{general.parallelization}.
\clearpage
%==================================
\subsection{GriddedTopography2AtmospherePotentialCoefficients}\label{GriddedTopography2AtmospherePotentialCoefficients}
Estimate interior and exterior potential coefficients for atmosphere above digital terrain models.
Coefficients for interior $(1/r)^{n+1}$ and exterior ($r^n$) are computed.
The density of the atmosphere is assumed to be (Sjöberg, 1998)
\begin{equation}
\rho_0\left(\frac{R}{R+h}\right)^\nu,
\end{equation}
where $R$ is the radial distance of the ellipsoid at each point, $h$ the radial height above the ellipsoid,
$\rho_0$ is \config{densitySeaLevel} and \config{nu} $\nu$ is a constant factor. The density is integrated
from \config{radialLowerBound} and \config{upperAtmosphericBoundary} above the ellipsoid.
The \config{radialLowerBound} is typically the topography and can be computed as expression at every point
from \configFile{inputfileGriddedData}{griddedData}.


\keepXColumns
\begin{tabularx}{\textwidth}{N T A}
\hline
Name & Type & Annotation\\
\hline
\hfuzz=500pt\includegraphics[width=1em]{element.pdf}~outputfilePotentialCoefficientsExterior & \hfuzz=500pt filename & \hfuzz=500pt \\
\hfuzz=500pt\includegraphics[width=1em]{element.pdf}~outputfilePotentialCoefficientsInterior & \hfuzz=500pt filename & \hfuzz=500pt \\
\hfuzz=500pt\includegraphics[width=1em]{element-mustset.pdf}~inputfileGriddedData & \hfuzz=500pt filename & \hfuzz=500pt Digital Terrain Model\\
\hfuzz=500pt\includegraphics[width=1em]{element.pdf}~densitySeaLevel & \hfuzz=500pt double & \hfuzz=500pt [kg/m**3]\\
\hfuzz=500pt\includegraphics[width=1em]{element.pdf}~ny & \hfuzz=500pt double & \hfuzz=500pt Constant for Atmosphere\\
\hfuzz=500pt\includegraphics[width=1em]{element.pdf}~radialLowerBound & \hfuzz=500pt expression & \hfuzz=500pt expression (variables 'L', 'B', 'height', 'data', and 'area' are taken from the gridded data\\
\hfuzz=500pt\includegraphics[width=1em]{element.pdf}~upperAtmosphericBoundary & \hfuzz=500pt double & \hfuzz=500pt constant upper bound [m]\\
\hfuzz=500pt\includegraphics[width=1em]{element.pdf}~factor & \hfuzz=500pt double & \hfuzz=500pt the result is multiplied by this factor, set -1 to subtract the field\\
\hfuzz=500pt\includegraphics[width=1em]{element.pdf}~minDegree & \hfuzz=500pt uint & \hfuzz=500pt \\
\hfuzz=500pt\includegraphics[width=1em]{element-mustset.pdf}~maxDegree & \hfuzz=500pt uint & \hfuzz=500pt \\
\hfuzz=500pt\includegraphics[width=1em]{element.pdf}~GM & \hfuzz=500pt double & \hfuzz=500pt Geocentric gravitational constant\\
\hfuzz=500pt\includegraphics[width=1em]{element.pdf}~R & \hfuzz=500pt double & \hfuzz=500pt reference radius\\
\hline
\end{tabularx}

This program is \reference{parallelized}{general.parallelization}.
\clearpage
%==================================
\subsection{GriddedTopography2PotentialCoefficients}\label{GriddedTopography2PotentialCoefficients}
Estimate potential coefficients from digital terrain models.
Coefficients for interior $(1/r)^{n+1}$ and exterior ($r^n$) are computed.


\keepXColumns
\begin{tabularx}{\textwidth}{N T A}
\hline
Name & Type & Annotation\\
\hline
\hfuzz=500pt\includegraphics[width=1em]{element.pdf}~outputfilePotentialCoefficients & \hfuzz=500pt filename & \hfuzz=500pt \\
\hfuzz=500pt\includegraphics[width=1em]{element.pdf}~outputfilePotentialCoefficientsInterior & \hfuzz=500pt filename & \hfuzz=500pt \\
\hfuzz=500pt\includegraphics[width=1em]{element-mustset.pdf}~inputfileGriddedData & \hfuzz=500pt filename & \hfuzz=500pt Digital Terrain Model\\
\hfuzz=500pt\includegraphics[width=1em]{element.pdf}~density & \hfuzz=500pt expression & \hfuzz=500pt expression [kg/m\textasciicircum{}3]\\
\hfuzz=500pt\includegraphics[width=1em]{element.pdf}~radialUpperBound & \hfuzz=500pt expression & \hfuzz=500pt expression (variables 'L', 'B', 'height', 'data', and 'area' are taken from the gridded data\\
\hfuzz=500pt\includegraphics[width=1em]{element.pdf}~radialLowerBound & \hfuzz=500pt expression & \hfuzz=500pt expression (variables 'L', 'B', 'height', 'data', and 'area' are taken from the gridded data\\
\hfuzz=500pt\includegraphics[width=1em]{element.pdf}~factor & \hfuzz=500pt double & \hfuzz=500pt the result is multiplied by this factor\\
\hfuzz=500pt\includegraphics[width=1em]{element.pdf}~minDegree & \hfuzz=500pt uint & \hfuzz=500pt \\
\hfuzz=500pt\includegraphics[width=1em]{element-mustset.pdf}~maxDegree & \hfuzz=500pt uint & \hfuzz=500pt \\
\hfuzz=500pt\includegraphics[width=1em]{element.pdf}~GM & \hfuzz=500pt double & \hfuzz=500pt Geocentric gravitational constant\\
\hfuzz=500pt\includegraphics[width=1em]{element.pdf}~R & \hfuzz=500pt double & \hfuzz=500pt reference radius\\
\hline
\end{tabularx}

This program is \reference{parallelized}{general.parallelization}.
\clearpage
%==================================
\subsection{GriddedTopographyEllipsoidal2Radial}\label{GriddedTopographyEllipsoidal2Radial}
Interpolate digital terrain models from ellipoidal heights to radial heights.


\keepXColumns
\begin{tabularx}{\textwidth}{N T A}
\hline
Name & Type & Annotation\\
\hline
\hfuzz=500pt\includegraphics[width=1em]{element.pdf}~outputfileGriddedData & \hfuzz=500pt filename & \hfuzz=500pt \\
\hfuzz=500pt\includegraphics[width=1em]{element-mustset.pdf}~inputfileGriddedData & \hfuzz=500pt filename & \hfuzz=500pt Digital Terrain Model\\
\hline
\end{tabularx}

\clearpage
%==================================
\subsection{Matrix2GriddedData}\label{Matrix2GriddedData}
This program reads a \file{matrix file}{matrix} with data in columns
and convert into \file{gridded data}{griddedData}.
The input columns are enumerated by \verb|data0|,~\verb|data1|,~\ldots,
see~\reference{dataVariables}{general.parser:dataVariables}.


\keepXColumns
\begin{tabularx}{\textwidth}{N T A}
\hline
Name & Type & Annotation\\
\hline
\hfuzz=500pt\includegraphics[width=1em]{element-mustset.pdf}~outputfileGriddedData & \hfuzz=500pt filename & \hfuzz=500pt \\
\hfuzz=500pt\includegraphics[width=1em]{element-mustset.pdf}~inputfileMatrix & \hfuzz=500pt filename & \hfuzz=500pt \\
\hfuzz=500pt\includegraphics[width=1em]{element-mustset.pdf}~points & \hfuzz=500pt choice & \hfuzz=500pt \\
\hfuzz=500pt\includegraphics[width=1em]{connector.pdf}\includegraphics[width=1em]{element-mustset.pdf}~ellipsoidal & \hfuzz=500pt sequence & \hfuzz=500pt \\
\hfuzz=500pt\quad\includegraphics[width=1em]{connector.pdf}\includegraphics[width=1em]{element-mustset.pdf}~longitude & \hfuzz=500pt expression & \hfuzz=500pt expression\\
\hfuzz=500pt\quad\includegraphics[width=1em]{connector.pdf}\includegraphics[width=1em]{element-mustset.pdf}~latitude & \hfuzz=500pt expression & \hfuzz=500pt expression\\
\hfuzz=500pt\quad\includegraphics[width=1em]{connector.pdf}\includegraphics[width=1em]{element-mustset.pdf}~height & \hfuzz=500pt expression & \hfuzz=500pt expression\\
\hfuzz=500pt\includegraphics[width=1em]{connector.pdf}\includegraphics[width=1em]{element-mustset.pdf}~cartesian & \hfuzz=500pt sequence & \hfuzz=500pt \\
\hfuzz=500pt\quad\includegraphics[width=1em]{connector.pdf}\includegraphics[width=1em]{element-mustset.pdf}~x & \hfuzz=500pt expression & \hfuzz=500pt expression\\
\hfuzz=500pt\quad\includegraphics[width=1em]{connector.pdf}\includegraphics[width=1em]{element-mustset.pdf}~y & \hfuzz=500pt expression & \hfuzz=500pt expression\\
\hfuzz=500pt\quad\includegraphics[width=1em]{connector.pdf}\includegraphics[width=1em]{element-mustset.pdf}~z & \hfuzz=500pt expression & \hfuzz=500pt expression\\
\hfuzz=500pt\includegraphics[width=1em]{element.pdf}~area & \hfuzz=500pt expression & \hfuzz=500pt expression (e.g. deltaL*2*sin(deltaB/2)*cos(data1/RHO))\\
\hfuzz=500pt\includegraphics[width=1em]{element-unbounded.pdf}~value & \hfuzz=500pt expression & \hfuzz=500pt expression\\
\hfuzz=500pt\includegraphics[width=1em]{element.pdf}~sortPoints & \hfuzz=500pt boolean & \hfuzz=500pt sort from north/west to south east\\
\hfuzz=500pt\includegraphics[width=1em]{element.pdf}~computeArea & \hfuzz=500pt boolean & \hfuzz=500pt the area can be computed automatically for rectangular grids\\
\hfuzz=500pt\includegraphics[width=1em]{element.pdf}~R & \hfuzz=500pt double & \hfuzz=500pt reference radius for ellipsoidal coordinates\\
\hfuzz=500pt\includegraphics[width=1em]{element.pdf}~inverseFlattening & \hfuzz=500pt double & \hfuzz=500pt reference flattening for ellipsoidal coordinates\\
\hline
\end{tabularx}

\clearpage
%==================================
\subsection{MatrixRectangular2GriddedData}\label{MatrixRectangular2GriddedData}
Read gridded data (matrix).


\keepXColumns
\begin{tabularx}{\textwidth}{N T A}
\hline
Name & Type & Annotation\\
\hline
\hfuzz=500pt\includegraphics[width=1em]{element-mustset.pdf}~outputfileGriddedData & \hfuzz=500pt filename & \hfuzz=500pt \\
\hfuzz=500pt\includegraphics[width=1em]{element-mustset.pdf}~inputfileMatrix & \hfuzz=500pt filename & \hfuzz=500pt \\
\hfuzz=500pt\includegraphics[width=1em]{element.pdf}~rowMajor & \hfuzz=500pt boolean & \hfuzz=500pt true: data is ordered row by row, false: columnwise\\
\hfuzz=500pt\includegraphics[width=1em]{element-mustset.pdf}~startLongitude & \hfuzz=500pt angle & \hfuzz=500pt longitude of upper left corner of the grid\\
\hfuzz=500pt\includegraphics[width=1em]{element-mustset.pdf}~startLatitude & \hfuzz=500pt angle & \hfuzz=500pt latitude of upper left corner of the grid\\
\hfuzz=500pt\includegraphics[width=1em]{element-mustset.pdf}~deltaLongitude & \hfuzz=500pt angle & \hfuzz=500pt sampling, negative for east to west data\\
\hfuzz=500pt\includegraphics[width=1em]{element-mustset.pdf}~deltaLatitude & \hfuzz=500pt angle & \hfuzz=500pt sampling, negative for south to north data\\
\hfuzz=500pt\includegraphics[width=1em]{element.pdf}~R & \hfuzz=500pt double & \hfuzz=500pt reference radius for ellipsoidal coordinates\\
\hfuzz=500pt\includegraphics[width=1em]{element.pdf}~inverseFlattening & \hfuzz=500pt double & \hfuzz=500pt reference flattening for ellipsoidal coordinates\\
\hline
\end{tabularx}

\clearpage
%==================================
\subsection{TimeSeries2GriddedData}\label{TimeSeries2GriddedData}
Interpret the data columns of \configFile{inputfileTimeSeries}{instrument}
as data points of a corresponding \configClass{grid}{gridType}.

For each epoch a \file{gridded data file}{griddedData} is written where
the \config{variableLoopTime} and \config{variableLoopIndex} are expanded for
each point of the given time series to create the file name for this epoch
(see \reference{text parser}{general.parser:text}).

The number of input data columns must be a multiple of the number $n$ of grid points.
If \config{isGroupedDataByPoint} is true the \configFile{inputfileTimeSeries}{instrument} starts
with all data (\verb|data0|, \verb|data1|\ldots) for the first point, followed by all data of the second point and so on.
If \config{isGroupedDataByPoint} is false, the file starts with \verb|data0| for all points, followed by all \verb|data1| and so on.

See also \program{GriddedData2TimeSeries}.


\keepXColumns
\begin{tabularx}{\textwidth}{N T A}
\hline
Name & Type & Annotation\\
\hline
\hfuzz=500pt\includegraphics[width=1em]{element-mustset.pdf}~outputfilesGriddedData & \hfuzz=500pt filename & \hfuzz=500pt for each epoch\\
\hfuzz=500pt\includegraphics[width=1em]{element.pdf}~variableLoopTime & \hfuzz=500pt string & \hfuzz=500pt variable with time of each epoch\\
\hfuzz=500pt\includegraphics[width=1em]{element.pdf}~variableLoopIndex & \hfuzz=500pt string & \hfuzz=500pt variable with index of current epoch (starts with zero)\\
\hfuzz=500pt\includegraphics[width=1em]{element.pdf}~variableLoopCount & \hfuzz=500pt string & \hfuzz=500pt variable with total number of epochs\\
\hfuzz=500pt\includegraphics[width=1em]{element-mustset.pdf}~inputfileTimeSeries & \hfuzz=500pt filename & \hfuzz=500pt each epoch: multiple data for points (MISCVALUES)\\
\hfuzz=500pt\includegraphics[width=1em]{element-mustset-unbounded.pdf}~grid & \hfuzz=500pt \hyperref[gridType]{grid} & \hfuzz=500pt corresponding grid points\\
\hfuzz=500pt\includegraphics[width=1em]{element.pdf}~isDataGroupedByPoint & \hfuzz=500pt boolean & \hfuzz=500pt multiple data are given point by point, otherwise: first data0 for all points, followed by all data1\\
\hfuzz=500pt\includegraphics[width=1em]{element.pdf}~R & \hfuzz=500pt double & \hfuzz=500pt reference radius for ellipsoidal coordinates on output\\
\hfuzz=500pt\includegraphics[width=1em]{element.pdf}~inverseFlattening & \hfuzz=500pt double & \hfuzz=500pt reference flattening for ellipsoidal coordinates on output, 0: spherical coordinates\\
\hline
\end{tabularx}

\clearpage
%==================================
\section{Programs: Instrument}
\subsection{Instrument2AllanVariance}\label{Instrument2AllanVariance}
This program computes the overlapping Allan variance from an
\configFile{inputfileInstrument}{instrument}.
The estimate is averaged over all arcs (arcs are assumed to contain no data gaps).

The overlapping Allan variance is defined as
\begin{equation}
  \sigma^2(m\tau_0) = \frac{1}{2(m\tau_0)^2(N-2m)} \sum_{n=1}^{N-2m}(x_{n+2m}-2x_{n+m}+x_n)^2,
\end{equation}
where $m\tau_0$ is the averaging interval defined by the median sampling $\tau_0$.


\keepXColumns
\begin{tabularx}{\textwidth}{N T A}
\hline
Name & Type & Annotation\\
\hline
\hfuzz=500pt\includegraphics[width=1em]{element-mustset.pdf}~outputfileAllanVariance & \hfuzz=500pt filename & \hfuzz=500pt column 0: averaging interval [seconds], column 1-(n-1): Allan variance for each data column\\
\hfuzz=500pt\includegraphics[width=1em]{element-mustset.pdf}~inputfileInstrument & \hfuzz=500pt filename & \hfuzz=500pt \\
\hline
\end{tabularx}

This program is \reference{parallelized}{general.parallelization}.
\clearpage
%==================================
\subsection{Instrument2CovarianceFunctionVCE}\label{Instrument2CovarianceFunctionVCE}
This estimates a covariance function of \configFile{inputfileInstrument}{instrument}
for all selected columns with \config{startDataFields} and \config{countDataFields}.
The estimation is performed robustly via variance component estimation.
Bad arcs are downweigthed and the accuracies can be written with \configFile{outputfileSigmasPerArc}{matrix}.
The length of the covariance functions are determined by the longest arc.
Additionaly the data can be detrended with \configClass{parameter}{parametrizationTemporalType}
and \configClass{parameterPerArc}{parametrizationTemporalType}.


\keepXColumns
\begin{tabularx}{\textwidth}{N T A}
\hline
Name & Type & Annotation\\
\hline
\hfuzz=500pt\includegraphics[width=1em]{element-mustset.pdf}~outputfileCovarianceFunction & \hfuzz=500pt filename & \hfuzz=500pt covariance functions\\
\hfuzz=500pt\includegraphics[width=1em]{element.pdf}~outputfileSigmasPerArc & \hfuzz=500pt filename & \hfuzz=500pt accuracies of each arc\\
\hfuzz=500pt\includegraphics[width=1em]{element.pdf}~outputfileResiduals & \hfuzz=500pt filename & \hfuzz=500pt \\
\hfuzz=500pt\includegraphics[width=1em]{element.pdf}~outputfileSolution & \hfuzz=500pt filename & \hfuzz=500pt estimated parameter vector (global part only)\\
\hfuzz=500pt\includegraphics[width=1em]{element-mustset.pdf}~inputfileInstrument & \hfuzz=500pt filename & \hfuzz=500pt \\
\hfuzz=500pt\includegraphics[width=1em]{element.pdf}~startDataFields & \hfuzz=500pt uint & \hfuzz=500pt start\\
\hfuzz=500pt\includegraphics[width=1em]{element.pdf}~countDataFields & \hfuzz=500pt uint & \hfuzz=500pt number of data fields (default: all after start)\\
\hfuzz=500pt\includegraphics[width=1em]{element-unbounded.pdf}~parameter & \hfuzz=500pt \hyperref[parametrizationTemporalType]{parametrizationTemporal} & \hfuzz=500pt data is reduced by temporal representation\\
\hfuzz=500pt\includegraphics[width=1em]{element-unbounded.pdf}~parameterPerArc & \hfuzz=500pt \hyperref[parametrizationTemporalType]{parametrizationTemporal} & \hfuzz=500pt data is reduced by temporal representation\\
\hfuzz=500pt\includegraphics[width=1em]{element.pdf}~iterationCount & \hfuzz=500pt uint & \hfuzz=500pt number of iterations for the estimation\\
\hline
\end{tabularx}

This program is \reference{parallelized}{general.parallelization}.
\clearpage
%==================================
\subsection{Instrument2CrossCorrelationFunction}\label{Instrument2CrossCorrelationFunction}
This program computes the cross correlation between all corresponding data columns
in two \file{instrument files}{instrument}. The instrument files must be synchronized (\program{InstrumentSynchronize}).
The \configFile{outputfileCorrelation}{matrix} is a matrix with the first column containing the time lag followed by
cross-correlation function for each data column. The maximum lag is defined by the maximum arc length.

The correlation is based on the unbiased estimate of the cross-covariance between data columns $x$ and $y$,
\begin{equation}
  \sigma_{xy}(h) = \frac{1}{N}\sum_{k=1} x_{k+h} y_k,
\end{equation}
which is averaged over all arcs. From this estimate, the correlation for each lag is then computed via
\begin{equation}
  r_{xy}(h) = \frac{\sigma_{xy}(h)}{\sigma_x(0)\sigma_y(0)},
\end{equation}
which is the ratio between the biased estimates of the cross-covariance at lag $h$ and the auto-covariance of the individual data columns.

For instrument with data gaps, lag bins without any data are set to NAN.


\keepXColumns
\begin{tabularx}{\textwidth}{N T A}
\hline
Name & Type & Annotation\\
\hline
\hfuzz=500pt\includegraphics[width=1em]{element-mustset.pdf}~outputfileCorrelation & \hfuzz=500pt filename & \hfuzz=500pt column 1: time lag, column 2..n cross-correlation\\
\hfuzz=500pt\includegraphics[width=1em]{element-mustset.pdf}~inputfileInstrument & \hfuzz=500pt filename & \hfuzz=500pt \\
\hfuzz=500pt\includegraphics[width=1em]{element-mustset.pdf}~inputfileInstrumentReference & \hfuzz=500pt filename & \hfuzz=500pt \\
\hline
\end{tabularx}

This program is \reference{parallelized}{general.parallelization}.
\clearpage
%==================================
\subsection{Instrument2Histogram}\label{Instrument2Histogram}
This program computes the arc-wise histogram from an \file{instrument file}{instrument}.
The output is a \file{matrix}{matrix} with the first column containing the lower bound of each bin.
The other columns contain the histograms for each arc.

\fig{!hb}{0.8}{instrument2Histogram}{fig:instrument2Histogram}{GRACE range-rate residuals of one month (one arc) divided into 50 bins.}


\keepXColumns
\begin{tabularx}{\textwidth}{N T A}
\hline
Name & Type & Annotation\\
\hline
\hfuzz=500pt\includegraphics[width=1em]{element-mustset.pdf}~outputfileMatrix & \hfuzz=500pt filename & \hfuzz=500pt column 1: lower bin bound; columns 2 to N: histogram of each arc\\
\hfuzz=500pt\includegraphics[width=1em]{element-mustset.pdf}~inputfileInstrument & \hfuzz=500pt filename & \hfuzz=500pt \\
\hfuzz=500pt\includegraphics[width=1em]{element.pdf}~selectDataField & \hfuzz=500pt uint & \hfuzz=500pt select channel for histogram computation\\
\hfuzz=500pt\includegraphics[width=1em]{element.pdf}~binCount & \hfuzz=500pt uint & \hfuzz=500pt (default: Freedman-Diaconis' choice, maximum of all channels)\\
\hfuzz=500pt\includegraphics[width=1em]{element.pdf}~lowerBound & \hfuzz=500pt expression & \hfuzz=500pt lower bound for bins (default: global minimum, data values outside are ignored)\\
\hfuzz=500pt\includegraphics[width=1em]{element.pdf}~upperBound & \hfuzz=500pt expression & \hfuzz=500pt upper bound for bins (default: global maximum, data values outside are ignored)\\
\hfuzz=500pt\includegraphics[width=1em]{element.pdf}~relative & \hfuzz=500pt boolean & \hfuzz=500pt output relative frequencies\\
\hfuzz=500pt\includegraphics[width=1em]{element.pdf}~cumulative & \hfuzz=500pt boolean & \hfuzz=500pt accumulate frequencies\\
\hline
\end{tabularx}

This program is \reference{parallelized}{general.parallelization}.
\clearpage
%==================================
\subsection{Instrument2PowerSpectralDensity}\label{Instrument2PowerSpectralDensity}
This program computes the power spectral density (PSD) for all data fields in an \file{instrument file}{instrument}.
The PSD is computed using Lomb's method. For each arc and each frequency $f$, a sinusoid is fit to the data
\begin{equation}
  l_i = a \cos(2\pi f t_i) + b \sin(2\pi f t_i) + e_i
\end{equation}

The PSD for this frequency is then computed by forming the square sum of adjusted observations:
\begin{equation}
  P(f) = \sum_i \hat{l}^2_i.
\end{equation}

The resulting PSD is the average over all arcs. For regularly sampled time series,
this method yields the same results as FFT based PSD estimates.

A regular frequency grid based on the longest arc and the median sampling is computed.
The maximum number of epochs per arc is determined by
\begin{equation}
  N = \frac{t_{\text{end}} - t_{\text{start}}}{\Delta t_{\text{median}} } + 1,
\end{equation}
the Nyquist frequency is given by
\begin{equation}
  f_{\text{nyq}} = \frac{1}{2\Delta t_{\text{median}}}.
\end{equation}

If it is suspected that \configFile{inputfileInstrument}{instrument} contains secular variations,
the input should be detrended using \program{InstrumentDetrend}.

See also \program{Instrument2CovarianceFunctionVCE},
\program{CovarianceFunction2PowerSpectralDensity}, \program{PowerSpectralDensity2CovarianceFunction}.


\keepXColumns
\begin{tabularx}{\textwidth}{N T A}
\hline
Name & Type & Annotation\\
\hline
\hfuzz=500pt\includegraphics[width=1em]{element-mustset.pdf}~outputfilePSD & \hfuzz=500pt filename & \hfuzz=500pt estimated PSD: column 0: frequency vector, column 1-(n-1): PSD estimate for each channel\\
\hfuzz=500pt\includegraphics[width=1em]{element-mustset.pdf}~inputfileInstrument & \hfuzz=500pt filename & \hfuzz=500pt \\
\hline
\end{tabularx}

This program is \reference{parallelized}{general.parallelization}.
\clearpage
%==================================
\subsection{Instrument2RmsPlotGrid}\label{Instrument2RmsPlotGrid}
This program computes an RMS plot grid from one or more \configFile{inputfileInstrument}{instrument}
containing 3D data (e.g. orbits or station positions), which can then be plotted as gridded data in \program{PlotGraph}.
The RMS is computed from the difference between \configFile{inputfileInstrument}{instrument} and
\configFile{inputfileInstrumentReference}{instrument}.
All instrument files must be synchronized (see \program{InstrumentSynchronize}).

Each separate \configFile{inputfileInstrument}{instrument} represents an entry (e.g. a satellite or station)
in the resulting grid. Therefore, providing, for example, 32 orbit files of GPS satellites
results in a grid with columns: mjd, id (0-31), rms.

The first three data columns of the instrument data are considered for computation of the RMS values.
The \config{factor} can be set to, for example, sqrt(3) to get 3D instead of 1D RMS values.

If \configClass{timeIntervals}{timeSeriesType} are provided, each \configFile{inputfileInstrument}{instrument}
and \configFile{inputfileInstrumentReference}{instrument} serves as a template with variable \verb|loopTime|.
This allows concatenation of instrument files, for example to create a month-long RMS plot grid from daily GPS
orbit files (see below).

Helmert parameters between the two frames can be estimated each epoch optionally if
\config{estimateShift}, \config{estimateScale}, or \config{estimateRotation} are set.
It uses a \reference{robust least squares adjustment}{fundamentals.robustLeastSquares}.

\fig{!hb}{0.8}{instrument2RmsPlotGrid}{fig:instrument2RmsPlotGrid}{Comparison of estimated GPS orbits with IGS final solution.}


\keepXColumns
\begin{tabularx}{\textwidth}{N T A}
\hline
Name & Type & Annotation\\
\hline
\hfuzz=500pt\includegraphics[width=1em]{element.pdf}~outputfileRmsPlotGrid & \hfuzz=500pt filename & \hfuzz=500pt columns: mjd, id, rms\\
\hfuzz=500pt\includegraphics[width=1em]{element.pdf}~outputfileHelmertTimeSeries & \hfuzz=500pt filename & \hfuzz=500pt columns: mjd, tx, ty, tz, scale, rx, ry, rz\\
\hfuzz=500pt\includegraphics[width=1em]{element-mustset-unbounded.pdf}~inputfileInstrument & \hfuzz=500pt filename & \hfuzz=500pt one file per satellite/station\\
\hfuzz=500pt\includegraphics[width=1em]{element-mustset-unbounded.pdf}~inputfileInstrumentReference & \hfuzz=500pt filename & \hfuzz=500pt one file per satellite/station, same order as above\\
\hfuzz=500pt\includegraphics[width=1em]{element-unbounded.pdf}~timeIntervals & \hfuzz=500pt \hyperref[timeSeriesType]{timeSeries} & \hfuzz=500pt for \{loopTime\} variable in inputfile\\
\hfuzz=500pt\includegraphics[width=1em]{element.pdf}~factor & \hfuzz=500pt double & \hfuzz=500pt e.g. sqrt(3) for 3D RMS\\
\hfuzz=500pt\includegraphics[width=1em]{element.pdf}~estimateShift & \hfuzz=500pt boolean & \hfuzz=500pt coordinate center every epoch\\
\hfuzz=500pt\includegraphics[width=1em]{element.pdf}~estimateScale & \hfuzz=500pt boolean & \hfuzz=500pt scale factor of position every epoch\\
\hfuzz=500pt\includegraphics[width=1em]{element.pdf}~estimateRotation & \hfuzz=500pt boolean & \hfuzz=500pt rotation every epoch\\
\hfuzz=500pt\includegraphics[width=1em]{element.pdf}~huber & \hfuzz=500pt double & \hfuzz=500pt for robust least squares\\
\hfuzz=500pt\includegraphics[width=1em]{element.pdf}~huberPower & \hfuzz=500pt double & \hfuzz=500pt for robust least squares\\
\hfuzz=500pt\includegraphics[width=1em]{element.pdf}~huberMaxIteration & \hfuzz=500pt uint & \hfuzz=500pt (maximum) number of iterations for robust estimation\\
\hline
\end{tabularx}

\clearpage
%==================================
\subsection{Instrument2Scaleogram}\label{Instrument2Scaleogram}
This program computes the wavelet transform of a time series up to a \config{maxLevel}.
The scalogram is written to a matrix which can be plotted by using a gridded layer in \program{PlotGraph}.
Individual detail levels can be written to matrix files by setting \configFile{outputfileLevels}{matrix}.
The data column to be decomposed must be set by \config{selectDataField}.

The wavelet transform is implemented as a filter bank, so care should be taken when the input contains data gaps.
Low/highpass wavelet filters are applied in forward and backward direction, input is padded symmetric.
See \configClass{digitalFilter}{digitalFilterType} for details.

\fig{!hb}{0.8}{instrument2Scaleogram}{fig:Instrument2Scaleogram}{GRACE range-rate residuals of one month.}


\keepXColumns
\begin{tabularx}{\textwidth}{N T A}
\hline
Name & Type & Annotation\\
\hline
\hfuzz=500pt\includegraphics[width=1em]{element-mustset.pdf}~outputfileScaleogram & \hfuzz=500pt filename & \hfuzz=500pt matrix columns: mjd, level, value\\
\hfuzz=500pt\includegraphics[width=1em]{element.pdf}~outputfileLevels & \hfuzz=500pt filename & \hfuzz=500pt use loopLevel as variable\\
\hfuzz=500pt\includegraphics[width=1em]{element-mustset.pdf}~inputfileInstrument & \hfuzz=500pt filename & \hfuzz=500pt \\
\hfuzz=500pt\includegraphics[width=1em]{element-mustset.pdf}~inputfileWavelet & \hfuzz=500pt filename & \hfuzz=500pt wavelet coefficients\\
\hfuzz=500pt\includegraphics[width=1em]{element.pdf}~selectDataField & \hfuzz=500pt uint & \hfuzz=500pt data column to transform\\
\hfuzz=500pt\includegraphics[width=1em]{element.pdf}~maxLevel & \hfuzz=500pt uint & \hfuzz=500pt maximum level of decomposition (default: full)\\
\hline
\end{tabularx}

\clearpage
%==================================
\subsection{Instrument2SpectralCoherence}\label{Instrument2SpectralCoherence}
This program computes the spectral coherence between two \file{instrument files}{instrument}.

The (magnitude-squared) coherence is defined as
\begin{equation}
  C_{xy}(f) = \frac{|P_{xy}(f)|^2}{P_{xx}(f)P_{yy}(f)}
\end{equation}
and is a measure in the range [0, 1] for the similarity of the signals $x$ and $y$ in frequency domain.
$P_{xy}$ is the cross-spectral density between $x$ and $y$ and $P_{xx}$, $P_{yy}$ are auto-spectral densities.
Auto- and cross-spectral densities are computed using Lomb's method (see \program{Instrument2PowerSpectralDensity} for details).

The resulting PSD is the average over all arcs. For regularly sampled time series,
this method yields the same results as FFT based PSD estimates.

A regular frequency grid based on the longest arc and the median sampling is computed.
The maximum number of epochs per arc is determined by
\begin{equation}
  N = \frac{t_{\text{end}} - t_{\text{start}}}{\Delta t_{\text{median}} } + 1,
\end{equation}
the Nyquist frequency is given by
\begin{equation}
  f_{\text{nyq}} = \frac{1}{2\Delta t_{\text{median}}}.
\end{equation}

If it is suspected that \configFile{inputfileInstrument}{instrument} contains secular variations,
the input should be detrended using \program{InstrumentDetrend}.

The \configFile{outputfileCoherence}{matrix} contains a matrix with the frequency vector as first column,
the coherence for each instrument channel is saved in the following columns.


\keepXColumns
\begin{tabularx}{\textwidth}{N T A}
\hline
Name & Type & Annotation\\
\hline
\hfuzz=500pt\includegraphics[width=1em]{element-mustset.pdf}~outputfileCoherence & \hfuzz=500pt filename & \hfuzz=500pt column 1: frequency, column 2-n coherence\\
\hfuzz=500pt\includegraphics[width=1em]{element-mustset.pdf}~inputfileInstrument & \hfuzz=500pt filename & \hfuzz=500pt \\
\hfuzz=500pt\includegraphics[width=1em]{element-mustset.pdf}~inputfileInstrumentReference & \hfuzz=500pt filename & \hfuzz=500pt \\
\hline
\end{tabularx}

This program is \reference{parallelized}{general.parallelization}.
\clearpage
%==================================
\subsection{Instrument2Spectrogram}\label{Instrument2Spectrogram}
This program applies the Short Time Fourier Transform (STFT) to selected data columns
of \configFile{inputfileInstrument}{instrument} and computes the spectrogram.
The STFT is computed at centered \configClass{timeSeries}{timeSeriesType} with
an (possible overlapping) rectangular window with \config{windowLength} seconds.
Data gaps are zero padded within the window.

The \configFile{outputfileSpectrogram}{matrix} is a matrix with each row the time (MJD),
the frequency $[Hz]$, and the amplitudes $[unit/\sqrt{Hz}]$ for the selected data columns.
It can be plotted with \program{PlotGraph}.

\fig{!hb}{0.8}{instrument2Spectrogram}{fig:instrument2Spectrogram}{GRACE range-rate residuals of one month (window of 6 hours).}


\keepXColumns
\begin{tabularx}{\textwidth}{N T A}
\hline
Name & Type & Annotation\\
\hline
\hfuzz=500pt\includegraphics[width=1em]{element-mustset.pdf}~outputfileSpectrogram & \hfuzz=500pt filename & \hfuzz=500pt mjd, freq, ampl0, ampl1, ...\\
\hfuzz=500pt\includegraphics[width=1em]{element-mustset.pdf}~inputfileInstrument & \hfuzz=500pt filename & \hfuzz=500pt \\
\hfuzz=500pt\includegraphics[width=1em]{element-mustset-unbounded.pdf}~timeSeries & \hfuzz=500pt \hyperref[timeSeriesType]{timeSeries} & \hfuzz=500pt center of SFFT window\\
\hfuzz=500pt\includegraphics[width=1em]{element-mustset.pdf}~windowLength & \hfuzz=500pt double & \hfuzz=500pt [seconds]\\
\hfuzz=500pt\includegraphics[width=1em]{element.pdf}~startDataFields & \hfuzz=500pt uint & \hfuzz=500pt start\\
\hfuzz=500pt\includegraphics[width=1em]{element.pdf}~countDataFields & \hfuzz=500pt uint & \hfuzz=500pt number of data fields (default: all)\\
\hline
\end{tabularx}

This program is \reference{parallelized}{general.parallelization}.
\clearpage
%==================================
\subsection{InstrumentAccelerometer2ThermosphericDensity}\label{InstrumentAccelerometer2ThermosphericDensity}
This program estimates neutral mass densities along the satellite trajectory based on \file{accelerometer data}{instrument}.
In order to determine the neutral mass density the accelerometer input should only reflect the accelerations due to drag
(e.g. \configClass{miscAccelerations:atmosphericDrag}{miscAccelerationsType:atmosphericDrag}).
Thus, influences from solar and Earth radiation pressure must be reduced beforehand.


\keepXColumns
\begin{tabularx}{\textwidth}{N T A}
\hline
Name & Type & Annotation\\
\hline
\hfuzz=500pt\includegraphics[width=1em]{element-mustset.pdf}~outputfileDensity & \hfuzz=500pt filename & \hfuzz=500pt MISCVALUE (kg/m\textasciicircum{}3)\\
\hfuzz=500pt\includegraphics[width=1em]{element.pdf}~satelliteModel & \hfuzz=500pt filename & \hfuzz=500pt satellite macro model\\
\hfuzz=500pt\includegraphics[width=1em]{element-mustset.pdf}~inputfileOrbit & \hfuzz=500pt filename & \hfuzz=500pt \\
\hfuzz=500pt\includegraphics[width=1em]{element-mustset.pdf}~inputfileStarCamera & \hfuzz=500pt filename & \hfuzz=500pt \\
\hfuzz=500pt\includegraphics[width=1em]{element-mustset.pdf}~inputfileAccelerometer & \hfuzz=500pt filename & \hfuzz=500pt add non-gravitational forces in satellite reference frame\\
\hfuzz=500pt\includegraphics[width=1em]{element-mustset.pdf}~thermosphere & \hfuzz=500pt \hyperref[thermosphereType]{thermosphere} & \hfuzz=500pt used to compute temperature and wind\\
\hfuzz=500pt\includegraphics[width=1em]{element.pdf}~considerTemperature & \hfuzz=500pt boolean & \hfuzz=500pt compute drag and lift, otherwise simple drag coefficient is used\\
\hfuzz=500pt\includegraphics[width=1em]{element.pdf}~considerWind & \hfuzz=500pt boolean & \hfuzz=500pt \\
\hfuzz=500pt\includegraphics[width=1em]{element-mustset.pdf}~earthRotation & \hfuzz=500pt \hyperref[earthRotationType]{earthRotation} & \hfuzz=500pt \\
\hfuzz=500pt\includegraphics[width=1em]{element.pdf}~ephemerides & \hfuzz=500pt \hyperref[ephemeridesType]{ephemerides} & \hfuzz=500pt \\
\hline
\end{tabularx}

This program is \reference{parallelized}{general.parallelization}.
\clearpage
%==================================
\subsection{InstrumentAccelerometerApplyEstimatedParameters}\label{InstrumentAccelerometerApplyEstimatedParameters}
This program evaluates estimated satellite parameters and writes the result to an accelerometer file.


\keepXColumns
\begin{tabularx}{\textwidth}{N T A}
\hline
Name & Type & Annotation\\
\hline
\hfuzz=500pt\includegraphics[width=1em]{element-mustset.pdf}~outputfileAccelerometer & \hfuzz=500pt filename & \hfuzz=500pt \\
\hfuzz=500pt\includegraphics[width=1em]{element.pdf}~inputfileSatelliteModel & \hfuzz=500pt filename & \hfuzz=500pt satellite macro model\\
\hfuzz=500pt\includegraphics[width=1em]{element.pdf}~inputfileOrbit & \hfuzz=500pt filename & \hfuzz=500pt \\
\hfuzz=500pt\includegraphics[width=1em]{element.pdf}~inputfileStarCamera & \hfuzz=500pt filename & \hfuzz=500pt \\
\hfuzz=500pt\includegraphics[width=1em]{element.pdf}~inputfileAccelerometer & \hfuzz=500pt filename & \hfuzz=500pt add non-gravitational forces in satellite reference frame\\
\hfuzz=500pt\includegraphics[width=1em]{element.pdf}~earthRotation & \hfuzz=500pt \hyperref[earthRotationType]{earthRotation} & \hfuzz=500pt \\
\hfuzz=500pt\includegraphics[width=1em]{element.pdf}~ephemerides & \hfuzz=500pt \hyperref[ephemeridesType]{ephemerides} & \hfuzz=500pt may be needed by parametrizationAcceleration\\
\hfuzz=500pt\includegraphics[width=1em]{element-mustset-unbounded.pdf}~parametrizationAcceleration & \hfuzz=500pt \hyperref[parametrizationAccelerationType]{parametrizationAcceleration} & \hfuzz=500pt orbit force parameters\\
\hfuzz=500pt\includegraphics[width=1em]{element-mustset.pdf}~inputfileParameter & \hfuzz=500pt filename & \hfuzz=500pt estimated orbit force parameters\\
\hfuzz=500pt\includegraphics[width=1em]{element.pdf}~indexStart & \hfuzz=500pt int & \hfuzz=500pt position in the solution vector\\
\hfuzz=500pt\includegraphics[width=1em]{element.pdf}~rightSide & \hfuzz=500pt int & \hfuzz=500pt if solution contains several right hand sides, select one\\
\hfuzz=500pt\includegraphics[width=1em]{element.pdf}~factor & \hfuzz=500pt double & \hfuzz=500pt the result is multiplied by this factor\\
\hline
\end{tabularx}

This program is \reference{parallelized}{general.parallelization}.
\clearpage
%==================================
\subsection{InstrumentAccelerometerEstimateBiasScale}\label{InstrumentAccelerometerEstimateBiasScale}
This program calibrates \configFile{inputfileAccelerometer}{instrument} with respect to
simulated accelerometer data, see \program{SimulateAccelerometer}.
The parameters \configFile{outputfileSolution}{matrix}
of \configClass{parametrizationAcceleration}{parametrizationAccelerationType}
are estimated and the effect is reduced to calibrate the \file{accelerometer data}{instrument}.

If \configFile{inputfileThruster}{instrument} is given, the corresponding epochs
(within \config{marginThruster}) are not used for the parameter estimation,
but the accelerometer epochs are still calibrated afterwards.
An arbitrary instrument file is allowed here.

The \configFile{inputfileOrbit}{instrument}, \configFile{inputfileStarCamera}{instrument},
\configClass{earthRotation}{earthRotationType}, \configClass{ephemerides}{ephemeridesType},
and \configFile{satelliteModel}{satelliteModel} are only needed for some special parametrizations.


\keepXColumns
\begin{tabularx}{\textwidth}{N T A}
\hline
Name & Type & Annotation\\
\hline
\hfuzz=500pt\includegraphics[width=1em]{element-mustset.pdf}~outputfileAccelerometer & \hfuzz=500pt filename & \hfuzz=500pt \\
\hfuzz=500pt\includegraphics[width=1em]{element.pdf}~outputfileSolution & \hfuzz=500pt filename & \hfuzz=500pt \\
\hfuzz=500pt\includegraphics[width=1em]{element-mustset.pdf}~inputfileAccelerometer & \hfuzz=500pt filename & \hfuzz=500pt \\
\hfuzz=500pt\includegraphics[width=1em]{element-mustset.pdf}~inputfileAccelerometerSim & \hfuzz=500pt filename & \hfuzz=500pt \\
\hfuzz=500pt\includegraphics[width=1em]{element.pdf}~inputfileThruster & \hfuzz=500pt filename & \hfuzz=500pt remove thruster events\\
\hfuzz=500pt\includegraphics[width=1em]{element.pdf}~marginThruster & \hfuzz=500pt double & \hfuzz=500pt margin size (on both sides) [seconds]\\
\hfuzz=500pt\includegraphics[width=1em]{element.pdf}~inputfileOrbit & \hfuzz=500pt filename & \hfuzz=500pt \\
\hfuzz=500pt\includegraphics[width=1em]{element.pdf}~inputfileStarCamera & \hfuzz=500pt filename & \hfuzz=500pt \\
\hfuzz=500pt\includegraphics[width=1em]{element.pdf}~earthRotation & \hfuzz=500pt \hyperref[earthRotationType]{earthRotation} & \hfuzz=500pt \\
\hfuzz=500pt\includegraphics[width=1em]{element.pdf}~ephemerides & \hfuzz=500pt \hyperref[ephemeridesType]{ephemerides} & \hfuzz=500pt may be needed by parametrizationAcceleration\\
\hfuzz=500pt\includegraphics[width=1em]{element.pdf}~inputfileSatelliteModel & \hfuzz=500pt filename & \hfuzz=500pt satellite macro model\\
\hfuzz=500pt\includegraphics[width=1em]{element-mustset-unbounded.pdf}~parametrizationAcceleration & \hfuzz=500pt \hyperref[parametrizationAccelerationType]{parametrizationAcceleration} & \hfuzz=500pt \\
\hline
\end{tabularx}

This program is \reference{parallelized}{general.parallelization}.
\clearpage
%==================================
\subsection{InstrumentAccelerometerEstimateParameters}\label{InstrumentAccelerometerEstimateParameters}
This program estimates calibration parameters for acceleration data given given an optional reference acceleration.
Specifically, the program solves the equation
\begin{equation}
  \mathbf{a} - \mathbf{a}_\text{ref} = \mathbf{f}(\mathbf{x}) + \mathbf{e}
\end{equation}
for the unknown parameters $\mathbf{x}$, where $\mathbf{a}$ is given in \configFile{inputfileAccelerometer}{instrument} and
$\mathbf{a}_\text{ref}$ is given in  \configFile{inputfileAccelerometerReference}{instrument}.
The parametrization of $\mathbf{x}$ can be set via \configClass{parametrizationAcceleration}{parametrizationAccelerationType}.
Optionally, the empirical covariance functions for the accelerations $\mathbf{a}$ can be estimated by enabling \config{estimateCovarianceFunctions}.

The estimated parameters are written to the file \configFile{outputfileSolution}{matrix} and can be used by
\program{InstrumentAccelerometerApplyEstimatedParameters} to calibrate accelerometer measurements.


\keepXColumns
\begin{tabularx}{\textwidth}{N T A}
\hline
Name & Type & Annotation\\
\hline
\hfuzz=500pt\includegraphics[width=1em]{element-mustset.pdf}~outputfileSolution & \hfuzz=500pt filename & \hfuzz=500pt values for estimated parameters\\
\hfuzz=500pt\includegraphics[width=1em]{element.pdf}~outputfileParameterNames & \hfuzz=500pt filename & \hfuzz=500pt names of the estimated parameters\\
\hfuzz=500pt\includegraphics[width=1em]{element.pdf}~estimateArcSigmas & \hfuzz=500pt sequence & \hfuzz=500pt \\
\hfuzz=500pt\includegraphics[width=1em]{connector.pdf}\includegraphics[width=1em]{element.pdf}~outputfileArcSigmas & \hfuzz=500pt filename & \hfuzz=500pt accuracies of each arc\\
\hfuzz=500pt\includegraphics[width=1em]{element.pdf}~estimateEpochSigmas & \hfuzz=500pt sequence & \hfuzz=500pt \\
\hfuzz=500pt\includegraphics[width=1em]{connector.pdf}\includegraphics[width=1em]{element.pdf}~outputfileEpochSigmas & \hfuzz=500pt filename & \hfuzz=500pt estimated epoch-wise sigmas\\
\hfuzz=500pt\includegraphics[width=1em]{element.pdf}~estimateCovarianceFunctions & \hfuzz=500pt sequence & \hfuzz=500pt \\
\hfuzz=500pt\includegraphics[width=1em]{connector.pdf}\includegraphics[width=1em]{element.pdf}~outputfileCovarianceFunction & \hfuzz=500pt filename & \hfuzz=500pt covariance functions for x, y, z direction\\
\hfuzz=500pt\includegraphics[width=1em]{element-mustset.pdf}~inputfileAccelerometer & \hfuzz=500pt filename & \hfuzz=500pt \\
\hfuzz=500pt\includegraphics[width=1em]{element.pdf}~inputfileAccelerometerReference & \hfuzz=500pt filename & \hfuzz=500pt if not given, reference acceleration is assumed zero\\
\hfuzz=500pt\includegraphics[width=1em]{element.pdf}~inputfileOrbit & \hfuzz=500pt filename & \hfuzz=500pt may be needed by parametrizationAcceleration\\
\hfuzz=500pt\includegraphics[width=1em]{element.pdf}~inputfileStarCamera & \hfuzz=500pt filename & \hfuzz=500pt may be needed by parametrizationAcceleration\\
\hfuzz=500pt\includegraphics[width=1em]{element.pdf}~inputfileSatelliteModel & \hfuzz=500pt filename & \hfuzz=500pt satellite macro model (may be needed by parametrizationAcceleration)\\
\hfuzz=500pt\includegraphics[width=1em]{element.pdf}~earthRotation & \hfuzz=500pt \hyperref[earthRotationType]{earthRotation} & \hfuzz=500pt may be needed by parametrizationAcceleration\\
\hfuzz=500pt\includegraphics[width=1em]{element.pdf}~ephemerides & \hfuzz=500pt \hyperref[ephemeridesType]{ephemerides} & \hfuzz=500pt may be needed by parametrizationAcceleration\\
\hfuzz=500pt\includegraphics[width=1em]{element-mustset-unbounded.pdf}~parametrizationAcceleration & \hfuzz=500pt \hyperref[parametrizationAccelerationType]{parametrizationAcceleration} & \hfuzz=500pt \\
\hfuzz=500pt\includegraphics[width=1em]{element.pdf}~sigmaX & \hfuzz=500pt double & \hfuzz=500pt apriori accuracy in x-axis\\
\hfuzz=500pt\includegraphics[width=1em]{element.pdf}~sigmaY & \hfuzz=500pt double & \hfuzz=500pt apriori accuracy in y-axis\\
\hfuzz=500pt\includegraphics[width=1em]{element.pdf}~sigmaZ & \hfuzz=500pt double & \hfuzz=500pt apriori accuracy in z-axis\\
\hfuzz=500pt\includegraphics[width=1em]{element.pdf}~iterationCount & \hfuzz=500pt uint & \hfuzz=500pt iteration count for determining the covariance function\\
\hline
\end{tabularx}

This program is \reference{parallelized}{general.parallelization}.
\clearpage
%==================================
\subsection{InstrumentApplyTimeOffset}\label{InstrumentApplyTimeOffset}
This program applies a \configFile{inputfileTimeOffset}{instrument} (MISCVALUE)
to an \configFile{inputfileInstrument}{instrument}.
The time offsets in seconds are multiplicated with a \config{factor}.
The instrument files must be synchronized (see \program{InstrumentSynchronize}).


\keepXColumns
\begin{tabularx}{\textwidth}{N T A}
\hline
Name & Type & Annotation\\
\hline
\hfuzz=500pt\includegraphics[width=1em]{element-mustset.pdf}~outputfileInstrument & \hfuzz=500pt filename & \hfuzz=500pt \\
\hfuzz=500pt\includegraphics[width=1em]{element-mustset.pdf}~inputfileInstrument & \hfuzz=500pt filename & \hfuzz=500pt \\
\hfuzz=500pt\includegraphics[width=1em]{element-mustset.pdf}~inputfileTimeOffset & \hfuzz=500pt filename & \hfuzz=500pt MISCVALUE with time offset in seconds\\
\hfuzz=500pt\includegraphics[width=1em]{element.pdf}~factor & \hfuzz=500pt double & \hfuzz=500pt applied to time offset\\
\hline
\end{tabularx}

\clearpage
%==================================
\subsection{InstrumentArcCalculate}\label{InstrumentArcCalculate}
This program manipulates the data columns every arc of an \file{instrument file}{instrument} similar to
\program{FunctionsCalculate}, see there for more details.
If several \configFile{inputfileInstrument}{instrument}s are given the data columns are copied side by side.
For this the instrument files must be synchronized (see \program{InstrumentSynchronize}). For the data
columns the standard data variables are available, see~\reference{dataVariables}{general.parser:dataVariables}.
For the time column (MJD) a variable \verb|epoch| (together with \verb|epochmean|, \verb|epochmin|, \ldots)
is defined additionally.

The content of \configFile{outputfileInstrument}{instrument} is controlled by \config{outColumn}.
The number of \config{outColumn} must agree with the selected \configClass{outType}{instrumentTypeType}.
The algorithm to compute the output is as follows:
The expressions in \config{outColumn} are evaluated once for each epoch of the input.
The variables \verb|data0|,~\verb|data1|,~\ldots are replaced by the according values from the input columns before.
If no \config{outColumn} are specified all input columns are used instead directly.
The \configClass{instrument type}{instrumentTypeType} can be specified with \config{outType} and must be agree with the number of columns.

An extra \config{statistics} file can be generated with one mid epoch per arc. For the computation of the \config{outColumn} values
all~\reference{dataVariables}{general.parser:dataVariables} are available (e.g. \verb|epochmin|, \verb|data0mean|, \verb|data1std|, \ldots)
inclusively the \config{constant}s and estimated \config{parameter}s but without the \verb|data0|,~\verb|data1|,~\ldots itself.
The variables and the numbering of the columns refers to the \configFile{outputfileInstrument}{instrument}.

See also \program{FunctionsCalculate}, \program{MatrixCalculate}.


\keepXColumns
\begin{tabularx}{\textwidth}{N T A}
\hline
Name & Type & Annotation\\
\hline
\hfuzz=500pt\includegraphics[width=1em]{element.pdf}~outputfileInstrument & \hfuzz=500pt filename & \hfuzz=500pt \\
\hfuzz=500pt\includegraphics[width=1em]{element-mustset-unbounded.pdf}~inputfileInstrument & \hfuzz=500pt filename & \hfuzz=500pt data columns are appended to the right\\
\hfuzz=500pt\includegraphics[width=1em]{element-unbounded.pdf}~constant & \hfuzz=500pt expression & \hfuzz=500pt define a constant by name=value\\
\hfuzz=500pt\includegraphics[width=1em]{element-unbounded.pdf}~parameter & \hfuzz=500pt expression & \hfuzz=500pt define a parameter by name[=value]\\
\hfuzz=500pt\includegraphics[width=1em]{element-unbounded.pdf}~leastSquares & \hfuzz=500pt expression & \hfuzz=500pt try to minimize the expression by adjustment of the parameters\\
\hfuzz=500pt\includegraphics[width=1em]{element-unbounded.pdf}~removalCriteria & \hfuzz=500pt expression & \hfuzz=500pt row is removed if one criterion evaluates true.\\
\hfuzz=500pt\includegraphics[width=1em]{element.pdf}~outType & \hfuzz=500pt \hyperref[instrumentTypeType]{instrumentType} & \hfuzz=500pt \\
\hfuzz=500pt\includegraphics[width=1em]{element-unbounded.pdf}~outColumn & \hfuzz=500pt expression & \hfuzz=500pt expression of output columns, extra 'epoch' variable\\
\hfuzz=500pt\includegraphics[width=1em]{element.pdf}~statistics & \hfuzz=500pt sequence & \hfuzz=500pt \\
\hfuzz=500pt\includegraphics[width=1em]{connector.pdf}\includegraphics[width=1em]{element-mustset.pdf}~outputfileInstrument & \hfuzz=500pt filename & \hfuzz=500pt instrument file with mid epoch per arc, data columns are user defined\\
\hfuzz=500pt\includegraphics[width=1em]{connector.pdf}\includegraphics[width=1em]{element-mustset-unbounded.pdf}~outColumn & \hfuzz=500pt expression & \hfuzz=500pt expression to compute statistics columns, data* are from outColumn\\
\hline
\end{tabularx}

This program is \reference{parallelized}{general.parallelization}.
\clearpage
%==================================
\subsection{InstrumentArcCrossStatistics}\label{InstrumentArcCrossStatistics}
Computes statistics of selected data columns between two \file{instrument files}{instrument} arc wise.
The \configFile{outputfileStatisticsTimeSeries}{instrument} contains for every arc one (mid) epoch
with statistics column(s). Possible statistics are
\begin{itemize}
  \item Correlation
  \begin{equation}
    \rho = \frac{\sum_i x_i y_i}{\sqrt{(\sum_i x_i^2) (\sum_i y_i^2})},
  \end{equation}
  \item Error RMS
  \begin{equation}
    rms = \sqrt{\frac{1}{N}\sum_i (x_i-y_i)^2},
  \end{equation}
  \item Nash-Sutcliffe coefficient (NSC)
  \begin{equation}
    nsc = 1- \frac{\sum_i (x_i-y_i)^2}{\sum_i (y_i-\bar{y})^2}.
  \end{equation}
\end{itemize}
With \config{removeArcMean} the mean of each data column of each arc is reduced before.

With \config{perColumn} separate statistics for each selected data column are computed,
otherwise an overall value is computed.

See also \program{InstrumentArcStatistics}, \program{InstrumentStatisticsTimeSeries}.


\keepXColumns
\begin{tabularx}{\textwidth}{N T A}
\hline
Name & Type & Annotation\\
\hline
\hfuzz=500pt\includegraphics[width=1em]{element-mustset.pdf}~outputfileStatisticsTimeSeries & \hfuzz=500pt filename & \hfuzz=500pt statistics column(s) per arc, MISCVALUES\\
\hfuzz=500pt\includegraphics[width=1em]{element-mustset.pdf}~inputfileInstrument & \hfuzz=500pt filename & \hfuzz=500pt \\
\hfuzz=500pt\includegraphics[width=1em]{element-mustset.pdf}~inputfileInstrumentReference & \hfuzz=500pt filename & \hfuzz=500pt \\
\hfuzz=500pt\includegraphics[width=1em]{element-mustset.pdf}~statistics & \hfuzz=500pt choice & \hfuzz=500pt \\
\hfuzz=500pt\includegraphics[width=1em]{connector.pdf}\includegraphics[width=1em]{element-mustset.pdf}~correlation & \hfuzz=500pt  & \hfuzz=500pt \\
\hfuzz=500pt\includegraphics[width=1em]{connector.pdf}\includegraphics[width=1em]{element-mustset.pdf}~errorRMS & \hfuzz=500pt  & \hfuzz=500pt rms of differences\\
\hfuzz=500pt\includegraphics[width=1em]{connector.pdf}\includegraphics[width=1em]{element-mustset.pdf}~nashSutcliffe & \hfuzz=500pt  & \hfuzz=500pt with respect to reference field\\
\hfuzz=500pt\includegraphics[width=1em]{element.pdf}~removeArcMean & \hfuzz=500pt boolean & \hfuzz=500pt \\
\hfuzz=500pt\includegraphics[width=1em]{element.pdf}~startDataFields & \hfuzz=500pt uint & \hfuzz=500pt start\\
\hfuzz=500pt\includegraphics[width=1em]{element.pdf}~countDataFields & \hfuzz=500pt uint & \hfuzz=500pt number of data fields (default: all)\\
\hfuzz=500pt\includegraphics[width=1em]{element.pdf}~perColumn & \hfuzz=500pt boolean & \hfuzz=500pt compute statistic per column\\
\hline
\end{tabularx}

This program is \reference{parallelized}{general.parallelization}.
\clearpage
%==================================
\subsection{InstrumentArcStatistics}\label{InstrumentArcStatistics}
Computes statistics of selected data columns of \configFile{inputfileInstrument}{instrument} arc wise.
The \configFile{outputfileStatisticsTimeSeries}{instrument} contains for every arc one (mid) epoch
with statistics column(s). Possible statistics are root mean square, standard deviation,
mean, median, min, and max.

With \config{perColumn} separate statistics for each selected data column are computed,
otherwise an overall value is computed.

See also \program{InstrumentArcCrossStatistics}, \program{InstrumentStatisticsTimeSeries}.


\keepXColumns
\begin{tabularx}{\textwidth}{N T A}
\hline
Name & Type & Annotation\\
\hline
\hfuzz=500pt\includegraphics[width=1em]{element-mustset.pdf}~outputfileStatisticsTimeSeries & \hfuzz=500pt filename & \hfuzz=500pt columns: mjd, statistics column(s) per instrument file\\
\hfuzz=500pt\includegraphics[width=1em]{element-mustset.pdf}~inputfileInstrument & \hfuzz=500pt filename & \hfuzz=500pt \\
\hfuzz=500pt\includegraphics[width=1em]{element-mustset.pdf}~statistics & \hfuzz=500pt choice & \hfuzz=500pt \\
\hfuzz=500pt\includegraphics[width=1em]{connector.pdf}\includegraphics[width=1em]{element-mustset.pdf}~rootMeanSquare & \hfuzz=500pt  & \hfuzz=500pt \\
\hfuzz=500pt\includegraphics[width=1em]{connector.pdf}\includegraphics[width=1em]{element-mustset.pdf}~standardDeviation & \hfuzz=500pt  & \hfuzz=500pt \\
\hfuzz=500pt\includegraphics[width=1em]{connector.pdf}\includegraphics[width=1em]{element-mustset.pdf}~mean & \hfuzz=500pt  & \hfuzz=500pt \\
\hfuzz=500pt\includegraphics[width=1em]{connector.pdf}\includegraphics[width=1em]{element-mustset.pdf}~median & \hfuzz=500pt  & \hfuzz=500pt \\
\hfuzz=500pt\includegraphics[width=1em]{connector.pdf}\includegraphics[width=1em]{element-mustset.pdf}~min & \hfuzz=500pt  & \hfuzz=500pt \\
\hfuzz=500pt\includegraphics[width=1em]{connector.pdf}\includegraphics[width=1em]{element-mustset.pdf}~max & \hfuzz=500pt  & \hfuzz=500pt \\
\hfuzz=500pt\includegraphics[width=1em]{connector.pdf}\includegraphics[width=1em]{element-mustset.pdf}~epochCount & \hfuzz=500pt  & \hfuzz=500pt \\
\hfuzz=500pt\includegraphics[width=1em]{element.pdf}~startDataFields & \hfuzz=500pt uint & \hfuzz=500pt start\\
\hfuzz=500pt\includegraphics[width=1em]{element.pdf}~countDataFields & \hfuzz=500pt uint & \hfuzz=500pt number of data fields (default: all)\\
\hfuzz=500pt\includegraphics[width=1em]{element.pdf}~perColumn & \hfuzz=500pt boolean & \hfuzz=500pt compute statistic per column\\
\hfuzz=500pt\includegraphics[width=1em]{element.pdf}~ignoreNan & \hfuzz=500pt boolean & \hfuzz=500pt ignore NaN values in input\\
\hline
\end{tabularx}

This program is \reference{parallelized}{general.parallelization}.
\clearpage
%==================================
\subsection{InstrumentConcatenate}\label{InstrumentConcatenate}
This program concatenate the arcs from several \file{instrument files}{instrument}
and write it to a new \file{file}{instrument}. Input files must be of the same type.
The arcs are merged to one arc even though there is a gap inbetween.
To split the data into arcs use \program{InstrumentSynchronize}.
Three options are available: \config{sort}, \config{removeDuplicates} and \config{checkForNaNs}.
If \config{sort} is enabled, the program reads all files, no matter if they are sorted correctly in time, and
then sorts the epochs. If \config{removeDuplicates} is enabled, the program checks the whole data set
for epochs that are contained twice. And if \config{checkForNaNs} is enabled the data set is checked for
invalid epochs containing NaNs.


\keepXColumns
\begin{tabularx}{\textwidth}{N T A}
\hline
Name & Type & Annotation\\
\hline
\hfuzz=500pt\includegraphics[width=1em]{element-mustset.pdf}~outputfile & \hfuzz=500pt filename & \hfuzz=500pt \\
\hfuzz=500pt\includegraphics[width=1em]{element-mustset-unbounded.pdf}~inputfile & \hfuzz=500pt filename & \hfuzz=500pt \\
\hfuzz=500pt\includegraphics[width=1em]{element.pdf}~sort & \hfuzz=500pt boolean & \hfuzz=500pt sort epochs with increasing time\\
\hfuzz=500pt\includegraphics[width=1em]{element.pdf}~removeDuplicates & \hfuzz=500pt choice & \hfuzz=500pt remove duplicate epochs\\
\hfuzz=500pt\includegraphics[width=1em]{connector.pdf}\includegraphics[width=1em]{element-mustset.pdf}~keepFirst & \hfuzz=500pt sequence & \hfuzz=500pt keep first epoch with the same time stamp, remove all others\\
\hfuzz=500pt\quad\includegraphics[width=1em]{connector.pdf}\includegraphics[width=1em]{element.pdf}~margin & \hfuzz=500pt double & \hfuzz=500pt margin for identical times [seconds]\\
\hfuzz=500pt\includegraphics[width=1em]{connector.pdf}\includegraphics[width=1em]{element-mustset.pdf}~keepLast & \hfuzz=500pt sequence & \hfuzz=500pt keep last epoch with the same time stamp, remove all others\\
\hfuzz=500pt\quad\includegraphics[width=1em]{connector.pdf}\includegraphics[width=1em]{element.pdf}~margin & \hfuzz=500pt double & \hfuzz=500pt margin for identical times [seconds]\\
\hfuzz=500pt\includegraphics[width=1em]{element.pdf}~checkForNaNs & \hfuzz=500pt boolean & \hfuzz=500pt remove epochs with NaN values in one of the data fields\\
\hline
\end{tabularx}

\clearpage
%==================================
\subsection{InstrumentCovarianceCheck}\label{InstrumentCovarianceCheck}
This program checks \configFile{inputfileCovariance3d}{instrument}
3x3 covariance matrices if they are invertible or not and removes the invalid epochs.


\keepXColumns
\begin{tabularx}{\textwidth}{N T A}
\hline
Name & Type & Annotation\\
\hline
\hfuzz=500pt\includegraphics[width=1em]{element-mustset.pdf}~outputfileCovariance3d & \hfuzz=500pt filename & \hfuzz=500pt \\
\hfuzz=500pt\includegraphics[width=1em]{element-mustset.pdf}~inputfileCovariance3d & \hfuzz=500pt filename & \hfuzz=500pt \\
\hline
\end{tabularx}

This program is \reference{parallelized}{general.parallelization}.
\clearpage
%==================================
\subsection{InstrumentDetrend}\label{InstrumentDetrend}
Reduces \configClass{parametrizationTemporal}{parametrizationTemporalType} (e.g. const, trend, polynomial)
per arc from selected data columns of \configFile{inputfileInstrument}{instrument}
using a robust \reference{robust least squares adjustment}{fundamentals.robustLeastSquares}.

The \configFile{outputfileTimeSeriesArcParameters}{instrument} contains for every arc one (mid) epoch
with the estimated parameters. The order is: first all data (\verb|data0|, \verb|data1|, \ldots)
of first temporal parameter, followed by all data of the second temporal parameter and so on.


\keepXColumns
\begin{tabularx}{\textwidth}{N T A}
\hline
Name & Type & Annotation\\
\hline
\hfuzz=500pt\includegraphics[width=1em]{element-mustset.pdf}~outputfileInstrument & \hfuzz=500pt filename & \hfuzz=500pt detrended instrument time series\\
\hfuzz=500pt\includegraphics[width=1em]{element.pdf}~outputfileTimeSeriesArcParameters & \hfuzz=500pt filename & \hfuzz=500pt time series of estimated parameters per arc\\
\hfuzz=500pt\includegraphics[width=1em]{element-mustset.pdf}~inputfileInstrument & \hfuzz=500pt filename & \hfuzz=500pt \\
\hfuzz=500pt\includegraphics[width=1em]{element-mustset-unbounded.pdf}~parametrizationTemporal & \hfuzz=500pt \hyperref[parametrizationTemporalType]{parametrizationTemporal} & \hfuzz=500pt per arc, data is reduced by temporal representation\\
\hfuzz=500pt\includegraphics[width=1em]{element.pdf}~startDataFields & \hfuzz=500pt uint & \hfuzz=500pt start\\
\hfuzz=500pt\includegraphics[width=1em]{element.pdf}~countDataFields & \hfuzz=500pt uint & \hfuzz=500pt number of data fields (default: all after start)\\
\hfuzz=500pt\includegraphics[width=1em]{element.pdf}~huber & \hfuzz=500pt double & \hfuzz=500pt for robust least squares\\
\hfuzz=500pt\includegraphics[width=1em]{element.pdf}~huberPower & \hfuzz=500pt double & \hfuzz=500pt for robust least squares\\
\hfuzz=500pt\includegraphics[width=1em]{element.pdf}~huberMaxIteration & \hfuzz=500pt uint & \hfuzz=500pt (maximum) number of iterations for robust estimation\\
\hline
\end{tabularx}

This program is \reference{parallelized}{general.parallelization}.
\clearpage
%==================================
\subsection{InstrumentEarthRotation}\label{InstrumentEarthRotation}
Precompute Earth rotation matrix from celestial to terrestrial frame
and save as \file{StarCamera file}{instrument}.


\keepXColumns
\begin{tabularx}{\textwidth}{N T A}
\hline
Name & Type & Annotation\\
\hline
\hfuzz=500pt\includegraphics[width=1em]{element-mustset.pdf}~outputfileStarCamera & \hfuzz=500pt filename & \hfuzz=500pt rotation from CRF to TRF\\
\hfuzz=500pt\includegraphics[width=1em]{element-mustset.pdf}~earthRotation & \hfuzz=500pt \hyperref[earthRotationType]{earthRotation} & \hfuzz=500pt \\
\hfuzz=500pt\includegraphics[width=1em]{element-mustset-unbounded.pdf}~timeSeries & \hfuzz=500pt \hyperref[timeSeriesType]{timeSeries} & \hfuzz=500pt \\
\hline
\end{tabularx}

\clearpage
%==================================
\subsection{InstrumentEstimateEmpiricalCovariance}\label{InstrumentEstimateEmpiricalCovariance}
This program estimates the empirical auto- and cross-covariance of selected data columns per arc
of \configFile{inputfileInstrument}{instrument}.
The maximum computed lag is determined by the number of \configFile{outputfileCovarianceMatrix}{matrix} specified
(for a single output file only the auto-covariance is determined, for two output files auto- and cross-covariance is computed and so on).

Stationarity is assumed for the input time series, which means the temporal covariance matrix has Toeplitz structure.
\begin{equation}
\begin{bmatrix}
\Sigma & \Sigma_{\Delta_1} & \Sigma_{\Delta_2} & \Sigma_{\Delta_3} & \Sigma_{\Delta_4} \\
       & \Sigma            & \Sigma_{\Delta_1} & \Sigma_{\Delta_2} & \Sigma_{\Delta_3} \\
       &                   & \Sigma            & \Sigma_{\Delta_1} & \Sigma_{\Delta_2} \\
       &                   &                   & \Sigma            & \Sigma_{\Delta_1} \\
       &                   &                   &                   & \Sigma            \\
\end{bmatrix}
\end{equation}

The matrix for lag $h$ describes the covariance between $x_{t-h}$ and $x_{t}$, i.e. $\Sigma(t-h, t)$.

To get a reliable estimate, \program{InstrumentDetrend} should be called first.


\keepXColumns
\begin{tabularx}{\textwidth}{N T A}
\hline
Name & Type & Annotation\\
\hline
\hfuzz=500pt\includegraphics[width=1em]{element-mustset-unbounded.pdf}~outputfileCovarianceMatrix & \hfuzz=500pt filename & \hfuzz=500pt \\
\hfuzz=500pt\includegraphics[width=1em]{element-mustset.pdf}~inputfileInstrument & \hfuzz=500pt filename & \hfuzz=500pt \\
\hfuzz=500pt\includegraphics[width=1em]{element.pdf}~startDataFields & \hfuzz=500pt uint & \hfuzz=500pt start\\
\hfuzz=500pt\includegraphics[width=1em]{element.pdf}~countDataFields & \hfuzz=500pt uint & \hfuzz=500pt number of data fields (default: all after start)\\
\hline
\end{tabularx}

This program is \reference{parallelized}{general.parallelization}.
\clearpage
%==================================
\subsection{InstrumentEstimateHelmertTransformation}\label{InstrumentEstimateHelmertTransformation}
This program estimates a 3D Helmert transformation between two networks
(frame realizations, e.g. GNSS satellite or station network).
Each separate \config{data} represents a satellite/station/\ldots (e.g. 32 GPS satellites).
The instrument data (x,y,z position) considered can be set with \config{startData}.
The Helmert parameters are set up according to \configClass{parametrizationTemporal}{parametrizationTemporalType}
for each \configClass{timeIntervals}{timeSeriesType} and are estimated using a
\reference{robust least squares adjustment}{fundamentals.robustLeastSquares}.


\keepXColumns
\begin{tabularx}{\textwidth}{N T A}
\hline
Name & Type & Annotation\\
\hline
\hfuzz=500pt\includegraphics[width=1em]{element.pdf}~outputfileHelmertTimeSeries & \hfuzz=500pt filename & \hfuzz=500pt columns: mjd, Tx,Ty,Tz,s,Rx,Ry,Rz according to temporal parametrization\\
\hfuzz=500pt\includegraphics[width=1em]{element-mustset-unbounded.pdf}~data & \hfuzz=500pt sequence & \hfuzz=500pt e.g. satellite, station\\
\hfuzz=500pt\includegraphics[width=1em]{connector.pdf}\includegraphics[width=1em]{element.pdf}~outputfileInstrument & \hfuzz=500pt filename & \hfuzz=500pt transformed positions as instrument type Vector3d\\
\hfuzz=500pt\includegraphics[width=1em]{connector.pdf}\includegraphics[width=1em]{element.pdf}~outputfileInstrumentDiff & \hfuzz=500pt filename & \hfuzz=500pt position difference as instrument type Vector3d\\
\hfuzz=500pt\includegraphics[width=1em]{connector.pdf}\includegraphics[width=1em]{element-mustset.pdf}~inputfileInstrument & \hfuzz=500pt filename & \hfuzz=500pt \\
\hfuzz=500pt\includegraphics[width=1em]{connector.pdf}\includegraphics[width=1em]{element-mustset.pdf}~inputfileInstrumentReference & \hfuzz=500pt filename & \hfuzz=500pt \\
\hfuzz=500pt\includegraphics[width=1em]{connector.pdf}\includegraphics[width=1em]{element.pdf}~startDataFields & \hfuzz=500pt uint & \hfuzz=500pt start index of position (x,y,z) columns\\
\hfuzz=500pt\includegraphics[width=1em]{element-mustset-unbounded.pdf}~timeIntervals & \hfuzz=500pt \hyperref[timeSeriesType]{timeSeries} & \hfuzz=500pt parameters are estimated per interval\\
\hfuzz=500pt\includegraphics[width=1em]{element-mustset-unbounded.pdf}~parametrizationTemporal & \hfuzz=500pt \hyperref[parametrizationTemporalType]{parametrizationTemporal} & \hfuzz=500pt temporal parametrization\\
\hfuzz=500pt\includegraphics[width=1em]{element.pdf}~estimateShift & \hfuzz=500pt boolean & \hfuzz=500pt coordinate center\\
\hfuzz=500pt\includegraphics[width=1em]{element.pdf}~estimateScale & \hfuzz=500pt boolean & \hfuzz=500pt scale factor of position\\
\hfuzz=500pt\includegraphics[width=1em]{element.pdf}~estimateRotation & \hfuzz=500pt boolean & \hfuzz=500pt rotation\\
\hfuzz=500pt\includegraphics[width=1em]{element.pdf}~huber & \hfuzz=500pt double & \hfuzz=500pt for robust least squares\\
\hfuzz=500pt\includegraphics[width=1em]{element.pdf}~huberPower & \hfuzz=500pt double & \hfuzz=500pt for robust least squares\\
\hfuzz=500pt\includegraphics[width=1em]{element.pdf}~huberMaxIteration & \hfuzz=500pt uint & \hfuzz=500pt (maximum) number of iterations for robust estimation\\
\hline
\end{tabularx}

\clearpage
%==================================
\subsection{InstrumentFilter}\label{InstrumentFilter}
This program filter selected data columns of \configFile{inputfileInstrument}{instrument}
with \configClass{digitalFilter}{digitalFilterType} arc wise.


\keepXColumns
\begin{tabularx}{\textwidth}{N T A}
\hline
Name & Type & Annotation\\
\hline
\hfuzz=500pt\includegraphics[width=1em]{element-mustset.pdf}~outputfileInstrument & \hfuzz=500pt filename & \hfuzz=500pt \\
\hfuzz=500pt\includegraphics[width=1em]{element-mustset.pdf}~inputfileInstrument & \hfuzz=500pt filename & \hfuzz=500pt \\
\hfuzz=500pt\includegraphics[width=1em]{element-mustset-unbounded.pdf}~digitalFilter & \hfuzz=500pt \hyperref[digitalFilterType]{digitalFilter} & \hfuzz=500pt \\
\hfuzz=500pt\includegraphics[width=1em]{element.pdf}~startDataFields & \hfuzz=500pt uint & \hfuzz=500pt start\\
\hfuzz=500pt\includegraphics[width=1em]{element.pdf}~countDataFields & \hfuzz=500pt uint & \hfuzz=500pt number of data fields (default: all after start)\\
\hline
\end{tabularx}

This program is \reference{parallelized}{general.parallelization}.
\clearpage
%==================================
\subsection{InstrumentInsertNAN}\label{InstrumentInsertNAN}
This program inserts NAN epochs into \configFile{inputfileInstrument}{instrument} files,
either at specific \configClass{times}{timeSeriesType} or where gaps in the instrument are detected.


\keepXColumns
\begin{tabularx}{\textwidth}{N T A}
\hline
Name & Type & Annotation\\
\hline
\hfuzz=500pt\includegraphics[width=1em]{element-mustset.pdf}~outputfileInstrument & \hfuzz=500pt filename & \hfuzz=500pt \\
\hfuzz=500pt\includegraphics[width=1em]{element-mustset.pdf}~inputfileInstrument & \hfuzz=500pt filename & \hfuzz=500pt \\
\hfuzz=500pt\includegraphics[width=1em]{element-unbounded.pdf}~times & \hfuzz=500pt \hyperref[timeSeriesType]{timeSeries} & \hfuzz=500pt Insert NAN at specific times.\\
\hfuzz=500pt\includegraphics[width=1em]{element.pdf}~atGaps & \hfuzz=500pt boolean & \hfuzz=500pt Insert NAN where epochs are more than 1.5 times the median sampling apart.\\
\hfuzz=500pt\includegraphics[width=1em]{element.pdf}~atArcEnds & \hfuzz=500pt boolean & \hfuzz=500pt Insert one epoch with data NAN at arc ends\\
\hline
\end{tabularx}

This program is \reference{parallelized}{general.parallelization}.
\clearpage
%==================================
\subsection{InstrumentMultiplyAdd}\label{InstrumentMultiplyAdd}
This program multiply \file{instrument data}{instrument} with a factor and add them together.
Afterwards the mean of each arc and data column can be removed with \config{removeArcMean}.
The instrument files must be synchronized (\program{InstrumentSynchronize}).

See also \program{InstrumentArcCalculate}.


\keepXColumns
\begin{tabularx}{\textwidth}{N T A}
\hline
Name & Type & Annotation\\
\hline
\hfuzz=500pt\includegraphics[width=1em]{element-mustset.pdf}~outputfileInstrument & \hfuzz=500pt filename & \hfuzz=500pt \\
\hfuzz=500pt\includegraphics[width=1em]{element-mustset-unbounded.pdf}~instrument & \hfuzz=500pt sequence & \hfuzz=500pt \\
\hfuzz=500pt\includegraphics[width=1em]{connector.pdf}\includegraphics[width=1em]{element-mustset.pdf}~inputfileInstrument & \hfuzz=500pt filename & \hfuzz=500pt \\
\hfuzz=500pt\includegraphics[width=1em]{connector.pdf}\includegraphics[width=1em]{element.pdf}~factor & \hfuzz=500pt double & \hfuzz=500pt \\
\hfuzz=500pt\includegraphics[width=1em]{element.pdf}~removeArcMean & \hfuzz=500pt boolean & \hfuzz=500pt remove mean value of each arc\\
\hline
\end{tabularx}

This program is \reference{parallelized}{general.parallelization}.
\clearpage
%==================================
\subsection{InstrumentReduceSampling}\label{InstrumentReduceSampling}
This program reduce the sampling of a instrument file. Only epochs with a time stamp
with a division by \config{sampling} without remainder are kept (inside \config{margin}).


\keepXColumns
\begin{tabularx}{\textwidth}{N T A}
\hline
Name & Type & Annotation\\
\hline
\hfuzz=500pt\includegraphics[width=1em]{element-mustset.pdf}~outputfileInstrument & \hfuzz=500pt filename & \hfuzz=500pt \\
\hfuzz=500pt\includegraphics[width=1em]{element-mustset.pdf}~inputfileInstrument & \hfuzz=500pt filename & \hfuzz=500pt \\
\hfuzz=500pt\includegraphics[width=1em]{element-mustset.pdf}~sampling & \hfuzz=500pt double & \hfuzz=500pt new sampling in seconds\\
\hfuzz=500pt\includegraphics[width=1em]{element.pdf}~margin & \hfuzz=500pt double & \hfuzz=500pt margin around the new sampling in seconds\\
\hfuzz=500pt\includegraphics[width=1em]{element.pdf}~relative2FirstEpoch & \hfuzz=500pt boolean & \hfuzz=500pt compute sampling relative to time of first epoch\\
\hline
\end{tabularx}

\clearpage
%==================================
\subsection{InstrumentRemoveEpochsByCriteria}\label{InstrumentRemoveEpochsByCriteria}
This program removes epochs from \configFile{inputfileInstrument}{instrument}
by evaluating a set of \config{removalCriteria} expressions. For the data
columns the standard data variables are available,
see~\reference{dataVariables}{general.parser:dataVariables}.

The instrument data can be reduced by data from \configFile{inputfileInstrumentReference}{instrument}
prior to evaluation of the expressions.

To reduce the data by its median, use an expression like \verb|data1-data1mean|.
To remove epochs that deviate by more than 3 sigma use \verb|abs(data1)>3*data1std|
or \verb|abs(data0-data0median)>3*1.4826*data0mad|.

All arcs in the input instrument file are concatenated, meaning expressions
like \verb|data1mean| refer to the complete dataset. The removed epochs can be saved
in a separate \configFile{outputfileInstrumentRemovedEpochs}{instrument}.


\keepXColumns
\begin{tabularx}{\textwidth}{N T A}
\hline
Name & Type & Annotation\\
\hline
\hfuzz=500pt\includegraphics[width=1em]{element.pdf}~outputfileInstrument & \hfuzz=500pt filename & \hfuzz=500pt all data is stored in one arc\\
\hfuzz=500pt\includegraphics[width=1em]{element.pdf}~outputfileInstrumentRemovedEpochs & \hfuzz=500pt filename & \hfuzz=500pt all data is stored in one arc\\
\hfuzz=500pt\includegraphics[width=1em]{element-mustset.pdf}~inputfileInstrument & \hfuzz=500pt filename & \hfuzz=500pt arcs are concatenated for processing\\
\hfuzz=500pt\includegraphics[width=1em]{element.pdf}~inputfileInstrumentReference & \hfuzz=500pt filename & \hfuzz=500pt if given, the reference data is reduced prior to the expressions being evaluated\\
\hfuzz=500pt\includegraphics[width=1em]{element-mustset-unbounded.pdf}~removalCriteria & \hfuzz=500pt expression & \hfuzz=500pt epochs are removed if one criterion evaluates true. data0 is the first data field.\\
\hfuzz=500pt\includegraphics[width=1em]{element.pdf}~margin & \hfuzz=500pt double & \hfuzz=500pt remove data around identified epochs (on both sides) [seconds]\\
\hline
\end{tabularx}

\clearpage
%==================================
\subsection{InstrumentRemoveEpochsByTimes}\label{InstrumentRemoveEpochsByTimes}
This program compares an \file{instrument file}{instrument} with a
\configClass{time series}{timeSeriesType}.
Epochs contained within the time series (including a defined margin)
are removed from the instrument file. The margin is added on
both sides of the epochs. The arcs of the instrument file are
concatenated to one arc. The removed epochs can be saved
in a separate instrument file.


\keepXColumns
\begin{tabularx}{\textwidth}{N T A}
\hline
Name & Type & Annotation\\
\hline
\hfuzz=500pt\includegraphics[width=1em]{element.pdf}~outputfileInstrument & \hfuzz=500pt filename & \hfuzz=500pt all epochs are concatenated in one arc\\
\hfuzz=500pt\includegraphics[width=1em]{element.pdf}~outputfileInstrumentRemovedEpochs & \hfuzz=500pt filename & \hfuzz=500pt all epochs are concatenated in one arc\\
\hfuzz=500pt\includegraphics[width=1em]{element-mustset.pdf}~inputfileInstrument & \hfuzz=500pt filename & \hfuzz=500pt \\
\hfuzz=500pt\includegraphics[width=1em]{element-mustset-unbounded.pdf}~timePoints & \hfuzz=500pt \hyperref[timeSeriesType]{timeSeries} & \hfuzz=500pt \\
\hfuzz=500pt\includegraphics[width=1em]{element.pdf}~margin & \hfuzz=500pt double & \hfuzz=500pt margin size (on both sides) [seconds]\\
\hline
\end{tabularx}

\clearpage
%==================================
\subsection{InstrumentRemoveEpochsThruster}\label{InstrumentRemoveEpochsThruster}
This program remove epochs from an \file{instrument file}{instrument}.
The epochs are defined by a \file{thruster file}{instrument}
plus a defined margin before and after the thruster firings.
The arcs of the instrument file are concatenated to one arc.
The removed epochs can be saved in a separate instrument file.


\keepXColumns
\begin{tabularx}{\textwidth}{N T A}
\hline
Name & Type & Annotation\\
\hline
\hfuzz=500pt\includegraphics[width=1em]{element.pdf}~outputfileInstrument & \hfuzz=500pt filename & \hfuzz=500pt all epochs are concatenated in one arc\\
\hfuzz=500pt\includegraphics[width=1em]{element.pdf}~outputfileInstrumentRemovedEpochs & \hfuzz=500pt filename & \hfuzz=500pt all epochs are concatenated in one arc\\
\hfuzz=500pt\includegraphics[width=1em]{element-mustset.pdf}~inputfileInstrument & \hfuzz=500pt filename & \hfuzz=500pt \\
\hfuzz=500pt\includegraphics[width=1em]{element-mustset.pdf}~inputfileThruster & \hfuzz=500pt filename & \hfuzz=500pt THRUSTER\\
\hfuzz=500pt\includegraphics[width=1em]{element.pdf}~marginBefore & \hfuzz=500pt double & \hfuzz=500pt margin before start of firing [seconds]\\
\hfuzz=500pt\includegraphics[width=1em]{element.pdf}~marginAfter & \hfuzz=500pt double & \hfuzz=500pt margin after end of firing [seconds]\\
\hline
\end{tabularx}

\clearpage
%==================================
\subsection{InstrumentResample}\label{InstrumentResample}
This program resamples \file{instrument data}{instrument} to a given
\configClass{timeSeries}{timeSeriesType} using a resampling
\configClass{method}{interpolatorTimeSeriesType}.

This program can also be used to reduce the sampling of an instrument file,
but a better way to reduce the sampling of noisy data with regular sampling
is to use a low pass filter first with \program{InstrumentFilter} and then thin
out the data with \program{InstrumentReduceSampling}.


\keepXColumns
\begin{tabularx}{\textwidth}{N T A}
\hline
Name & Type & Annotation\\
\hline
\hfuzz=500pt\includegraphics[width=1em]{element-mustset.pdf}~outputfileInstrument & \hfuzz=500pt filename & \hfuzz=500pt \\
\hfuzz=500pt\includegraphics[width=1em]{element-mustset.pdf}~inputfileInstrument & \hfuzz=500pt filename & \hfuzz=500pt \\
\hfuzz=500pt\includegraphics[width=1em]{element-mustset.pdf}~method & \hfuzz=500pt \hyperref[interpolatorTimeSeriesType]{interpolatorTimeSeries} & \hfuzz=500pt resampling method\\
\hfuzz=500pt\includegraphics[width=1em]{element-mustset-unbounded.pdf}~timeSeries & \hfuzz=500pt \hyperref[timeSeriesType]{timeSeries} & \hfuzz=500pt resampled points in time\\
\hline
\end{tabularx}

\clearpage
%==================================
\subsection{InstrumentRotate}\label{InstrumentRotate}
This program rotates \file{instrument data}{instrument} into a new reference frame
(using \configFile{inputfileStarCamera}{instrument}).
The rotation is usually done from satellite frame into inertial frame.

To apply Earth rotation to orbits use \program{Orbit2EarthFixedOrbit}
For other instrument data use \program{InstrumentEarthRotation} before.


\keepXColumns
\begin{tabularx}{\textwidth}{N T A}
\hline
Name & Type & Annotation\\
\hline
\hfuzz=500pt\includegraphics[width=1em]{element-mustset.pdf}~outputfileInstrument & \hfuzz=500pt filename & \hfuzz=500pt \\
\hfuzz=500pt\includegraphics[width=1em]{element-mustset.pdf}~inputfileInstrument & \hfuzz=500pt filename & \hfuzz=500pt \\
\hfuzz=500pt\includegraphics[width=1em]{element-mustset.pdf}~inputfileStarCamera & \hfuzz=500pt filename & \hfuzz=500pt \\
\hfuzz=500pt\includegraphics[width=1em]{element.pdf}~inverseRotate & \hfuzz=500pt boolean & \hfuzz=500pt \\
\hline
\end{tabularx}

This program is \reference{parallelized}{general.parallelization}.
\clearpage
%==================================
\subsection{InstrumentSetType}\label{InstrumentSetType}
Convert \file{instrument data}{instrument} into instrument data with new \configClass{type}{instrumentTypeType}.
The selected number of data columns must agree with the \configClass{type}{instrumentTypeType}.


\keepXColumns
\begin{tabularx}{\textwidth}{N T A}
\hline
Name & Type & Annotation\\
\hline
\hfuzz=500pt\includegraphics[width=1em]{element-mustset.pdf}~outputfileInstrument & \hfuzz=500pt filename & \hfuzz=500pt \\
\hfuzz=500pt\includegraphics[width=1em]{element-mustset.pdf}~inputfileInstrument & \hfuzz=500pt filename & \hfuzz=500pt \\
\hfuzz=500pt\includegraphics[width=1em]{element.pdf}~type & \hfuzz=500pt \hyperref[instrumentTypeType]{instrumentType} & \hfuzz=500pt \\
\hfuzz=500pt\includegraphics[width=1em]{element.pdf}~startDataFields & \hfuzz=500pt uint & \hfuzz=500pt start\\
\hfuzz=500pt\includegraphics[width=1em]{element.pdf}~countDataFields & \hfuzz=500pt uint & \hfuzz=500pt number of data fields (default: all after start)\\
\hline
\end{tabularx}

\clearpage
%==================================
\subsection{InstrumentStarCamera2AccAngularRate}\label{InstrumentStarCamera2AccAngularRate}
This program derivate from a time series of quaternions
a series of angular rates and angular accelerations.
The derivatives are computed by a polynomial interpolation
with \config{interpolationDegree} of the quaternions.


\keepXColumns
\begin{tabularx}{\textwidth}{N T A}
\hline
Name & Type & Annotation\\
\hline
\hfuzz=500pt\includegraphics[width=1em]{element.pdf}~outputfileAngularRate & \hfuzz=500pt filename & \hfuzz=500pt [rad/s], VECTOR3D\\
\hfuzz=500pt\includegraphics[width=1em]{element.pdf}~outputfileAngularAcc & \hfuzz=500pt filename & \hfuzz=500pt [rad/s**2], VECTOR3D\\
\hfuzz=500pt\includegraphics[width=1em]{element-mustset.pdf}~inputfileStarCamera & \hfuzz=500pt filename & \hfuzz=500pt \\
\hfuzz=500pt\includegraphics[width=1em]{element.pdf}~interpolationDegree & \hfuzz=500pt uint & \hfuzz=500pt derivation by polynomial interpolation of degree n\\
\hline
\end{tabularx}

\clearpage
%==================================
\subsection{InstrumentStarCamera2RollPitchYaw}\label{InstrumentStarCamera2RollPitchYaw}
Compute roll, pitch, yaw angles from \configFile{inputfileStarCamera}{instrument} data.
Optional the angles are computed relative
to a \configFile{inputfileStarCameraReference}{instrument}.

See also \program{SimulateStarCamera}.


\keepXColumns
\begin{tabularx}{\textwidth}{N T A}
\hline
Name & Type & Annotation\\
\hline
\hfuzz=500pt\includegraphics[width=1em]{element-mustset.pdf}~outputfileInstrument & \hfuzz=500pt filename & \hfuzz=500pt roll, pitch, yaw [rad], VECTOR3D\\
\hfuzz=500pt\includegraphics[width=1em]{element-mustset.pdf}~inputfileStarCamera & \hfuzz=500pt filename & \hfuzz=500pt \\
\hfuzz=500pt\includegraphics[width=1em]{element.pdf}~inputfileStarCameraReference & \hfuzz=500pt filename & \hfuzz=500pt nominal orientation\\
\hline
\end{tabularx}

This program is \reference{parallelized}{general.parallelization}.
\clearpage
%==================================
\subsection{InstrumentStarCamera2RotaryMatrix}\label{InstrumentStarCamera2RotaryMatrix}
Write \configFile{inputfileStarCamera}{instrument} rotations
as \configFile{outputfileInstrument}{instrument} rotary matrices
(for each epoch $xx, xy, xz, yx, yy, yz, zx, zy, zz$).


\keepXColumns
\begin{tabularx}{\textwidth}{N T A}
\hline
Name & Type & Annotation\\
\hline
\hfuzz=500pt\includegraphics[width=1em]{element-mustset.pdf}~outputfileInstrument & \hfuzz=500pt filename & \hfuzz=500pt xx, xy, xz, yx, yy, yz, zx, zy, zz (MISCVALUES)\\
\hfuzz=500pt\includegraphics[width=1em]{element-mustset.pdf}~inputfileStarCamera & \hfuzz=500pt filename & \hfuzz=500pt \\
\hline
\end{tabularx}

This program is \reference{parallelized}{general.parallelization}.
\clearpage
%==================================
\subsection{InstrumentStarCameraMultiply}\label{InstrumentStarCameraMultiply}
This program applies several rotations given by \configFile{inputfileStarCamera}{instrument}.
The resulting rotation is written as \configFile{outputfileStarCamera}{instrument}.
All instrument files must be synchronized (\program{InstrumentSynchronize}).


\keepXColumns
\begin{tabularx}{\textwidth}{N T A}
\hline
Name & Type & Annotation\\
\hline
\hfuzz=500pt\includegraphics[width=1em]{element-mustset.pdf}~outputfileStarCamera & \hfuzz=500pt filename & \hfuzz=500pt \\
\hfuzz=500pt\includegraphics[width=1em]{element-mustset-unbounded.pdf}~instrument & \hfuzz=500pt sequence & \hfuzz=500pt \\
\hfuzz=500pt\includegraphics[width=1em]{connector.pdf}\includegraphics[width=1em]{element-mustset.pdf}~inputfileStarCamera & \hfuzz=500pt filename & \hfuzz=500pt \\
\hfuzz=500pt\includegraphics[width=1em]{connector.pdf}\includegraphics[width=1em]{element.pdf}~inverse & \hfuzz=500pt boolean & \hfuzz=500pt \\
\hline
\end{tabularx}

\clearpage
%==================================
\subsection{InstrumentStatisticsTimeSeries}\label{InstrumentStatisticsTimeSeries}
This program computes a time series of statistics for one or more instrument files.
Possible statistics are root mean square, standard deviation, mean, median, min, and max.
The columns of the output time series are defined either as one per \configFile{inputfileInstrument}{instrument}
or, if \config{perColumn} is true, statistics are computed per column for each file.
Providing e.g. 32 orbit files of GPS satellites results in a time series matrix
with columns: mjd, statisticsG01, statisticsG02, ..., statisticsG32.
If \config{intervals} are provided, the input data is split into these intervals
and one statistic is computed per interval. Otherwise, overall statistics are computed.
The instrument data considered for computation of the component-wise statistics
can be set with \config{startDataFields} and \config{countDataFields}.
The \config{factor} can be set to e.g. sqrt(3) to get 3D instead of 1D RMS values.

See also \program{InstrumentArcStatistics}, \program{InstrumentArcCrossStatistics}.


\keepXColumns
\begin{tabularx}{\textwidth}{N T A}
\hline
Name & Type & Annotation\\
\hline
\hfuzz=500pt\includegraphics[width=1em]{element-mustset.pdf}~outputfileStatisticsTimeSeries & \hfuzz=500pt filename & \hfuzz=500pt columns: mjd, statistics column(s) per instrument file\\
\hfuzz=500pt\includegraphics[width=1em]{element-mustset-unbounded.pdf}~inputfileInstrument & \hfuzz=500pt filename & \hfuzz=500pt \\
\hfuzz=500pt\includegraphics[width=1em]{element-mustset.pdf}~statistics & \hfuzz=500pt choice & \hfuzz=500pt \\
\hfuzz=500pt\includegraphics[width=1em]{connector.pdf}\includegraphics[width=1em]{element-mustset.pdf}~rootMeanSquare & \hfuzz=500pt  & \hfuzz=500pt \\
\hfuzz=500pt\includegraphics[width=1em]{connector.pdf}\includegraphics[width=1em]{element-mustset.pdf}~standardDeviation & \hfuzz=500pt  & \hfuzz=500pt \\
\hfuzz=500pt\includegraphics[width=1em]{connector.pdf}\includegraphics[width=1em]{element-mustset.pdf}~mean & \hfuzz=500pt  & \hfuzz=500pt \\
\hfuzz=500pt\includegraphics[width=1em]{connector.pdf}\includegraphics[width=1em]{element-mustset.pdf}~median & \hfuzz=500pt  & \hfuzz=500pt \\
\hfuzz=500pt\includegraphics[width=1em]{connector.pdf}\includegraphics[width=1em]{element-mustset.pdf}~sum & \hfuzz=500pt  & \hfuzz=500pt \\
\hfuzz=500pt\includegraphics[width=1em]{connector.pdf}\includegraphics[width=1em]{element-mustset.pdf}~min & \hfuzz=500pt  & \hfuzz=500pt \\
\hfuzz=500pt\includegraphics[width=1em]{connector.pdf}\includegraphics[width=1em]{element-mustset.pdf}~max & \hfuzz=500pt  & \hfuzz=500pt \\
\hfuzz=500pt\includegraphics[width=1em]{connector.pdf}\includegraphics[width=1em]{element-mustset.pdf}~epochCount & \hfuzz=500pt  & \hfuzz=500pt \\
\hfuzz=500pt\includegraphics[width=1em]{element.pdf}~startDataFields & \hfuzz=500pt uint & \hfuzz=500pt start\\
\hfuzz=500pt\includegraphics[width=1em]{element.pdf}~countDataFields & \hfuzz=500pt uint & \hfuzz=500pt number of data fields (default: all)\\
\hfuzz=500pt\includegraphics[width=1em]{element.pdf}~perColumn & \hfuzz=500pt boolean & \hfuzz=500pt compute statistic per column\\
\hfuzz=500pt\includegraphics[width=1em]{element.pdf}~ignoreNan & \hfuzz=500pt boolean & \hfuzz=500pt ignore NaN values in statistic computation\\
\hfuzz=500pt\includegraphics[width=1em]{element-unbounded.pdf}~intervals & \hfuzz=500pt \hyperref[timeSeriesType]{timeSeries} & \hfuzz=500pt intervals for statistics computation (one statistic per interval)\\
\hfuzz=500pt\includegraphics[width=1em]{element.pdf}~factor & \hfuzz=500pt double & \hfuzz=500pt e.g. sqrt(3) for 3D RMS\\
\hline
\end{tabularx}

\clearpage
%==================================
\subsection{InstrumentSynchronize}\label{InstrumentSynchronize}
This program reads several \file{instrument files}{instrument} and synchronize the data.
Every epoch with some missing data will be deleted so the remaining epochs
have data from every instrument.

In a second step the epochs are divided into arcs with maximal epochs
(or \config{maxArcLen}) without having a gap inside an arc.
A Gap is defined by a time step with at least \config{minGap} seconds
between consecutive epochs or if not set the 1.5 of the median sampling.
Arc with an epoch count less than \config{minArcLen} will be rejected.

A specific region can be selected with \configClass{border}{borderType}.
In this case one of the instrument data must an orbit.

If \configClass{timeIntervals}{timeSeriesType} is given the data are also divided into time bins.
The assignment of arcs to the bins can be saved in \configFile{outputfileArcList}{arcList}.
This file can be used for the variational equation approach or \program{KalmanBuildNormals}.

Instrument files from \config{irregularData} are not synchronized but
divided into the same number of arcs within the same time intervals.
Data outside the defined arcs will be deleted.


\keepXColumns
\begin{tabularx}{\textwidth}{N T A}
\hline
Name & Type & Annotation\\
\hline
\hfuzz=500pt\includegraphics[width=1em]{element-mustset-unbounded.pdf}~data & \hfuzz=500pt sequence & \hfuzz=500pt \\
\hfuzz=500pt\includegraphics[width=1em]{connector.pdf}\includegraphics[width=1em]{element.pdf}~outputfileInstrument & \hfuzz=500pt filename & \hfuzz=500pt \\
\hfuzz=500pt\includegraphics[width=1em]{connector.pdf}\includegraphics[width=1em]{element-mustset.pdf}~inputfileInstrument & \hfuzz=500pt filename & \hfuzz=500pt \\
\hfuzz=500pt\includegraphics[width=1em]{element.pdf}~margin & \hfuzz=500pt double & \hfuzz=500pt margin for identical times [seconds]\\
\hfuzz=500pt\includegraphics[width=1em]{element.pdf}~minGap & \hfuzz=500pt double & \hfuzz=500pt minimal time to define a gap and to begin a new arc, 0: no dividing [seconds], if not set 1.5*median sampling is used\\
\hfuzz=500pt\includegraphics[width=1em]{element.pdf}~minArcLength & \hfuzz=500pt uint & \hfuzz=500pt minimal number of epochs of an arc\\
\hfuzz=500pt\includegraphics[width=1em]{element.pdf}~maxArcLength & \hfuzz=500pt uint & \hfuzz=500pt maximal number of epochs of an arc\\
\hfuzz=500pt\includegraphics[width=1em]{element.pdf}~arcType & \hfuzz=500pt choice & \hfuzz=500pt all arcs or only ascending or descending arcs are selected\\
\hfuzz=500pt\includegraphics[width=1em]{connector.pdf}\includegraphics[width=1em]{element-mustset.pdf}~ascending & \hfuzz=500pt  & \hfuzz=500pt \\
\hfuzz=500pt\includegraphics[width=1em]{connector.pdf}\includegraphics[width=1em]{element-mustset.pdf}~descending & \hfuzz=500pt  & \hfuzz=500pt \\
\hfuzz=500pt\includegraphics[width=1em]{element-unbounded.pdf}~border & \hfuzz=500pt \hyperref[borderType]{border} & \hfuzz=500pt only data in a specific region is selected\\
\hfuzz=500pt\includegraphics[width=1em]{element-unbounded.pdf}~timeIntervals & \hfuzz=500pt \hyperref[timeSeriesType]{timeSeries} & \hfuzz=500pt divide data into time bins\\
\hfuzz=500pt\includegraphics[width=1em]{element.pdf}~outputfileArcList & \hfuzz=500pt filename & \hfuzz=500pt arc and time bin mapping\\
\hfuzz=500pt\includegraphics[width=1em]{element-unbounded.pdf}~irregularData & \hfuzz=500pt sequence & \hfuzz=500pt instrument files with irregular sampling\\
\hfuzz=500pt\includegraphics[width=1em]{connector.pdf}\includegraphics[width=1em]{element.pdf}~outputfileInstrument & \hfuzz=500pt filename & \hfuzz=500pt \\
\hfuzz=500pt\includegraphics[width=1em]{connector.pdf}\includegraphics[width=1em]{element-mustset.pdf}~inputfileInstrument & \hfuzz=500pt filename & \hfuzz=500pt \\
\hfuzz=500pt\includegraphics[width=1em]{connector.pdf}\includegraphics[width=1em]{element.pdf}~minArcLength & \hfuzz=500pt uint & \hfuzz=500pt minimal number of epochs in an arc\\
\hline
\end{tabularx}

\clearpage
%==================================
\subsection{InstrumentWaveletDecomposition}\label{InstrumentWaveletDecomposition}
This program performs a multilevel one-dimensional wavelet analysis on one \config{selectDataField}
data column of \configFile{inputfileInstrument}{instrument}.
The \configFile{outputfileInstrument}{instrument} contains the decomposed levels in time domain ${a_J,d_J,...,d_1}$


\keepXColumns
\begin{tabularx}{\textwidth}{N T A}
\hline
Name & Type & Annotation\\
\hline
\hfuzz=500pt\includegraphics[width=1em]{element-mustset.pdf}~outputfileInstrument & \hfuzz=500pt filename & \hfuzz=500pt MISCVALUES, decomposed levels in time domain a\_J,d\_J,...,d\_1\\
\hfuzz=500pt\includegraphics[width=1em]{element-mustset.pdf}~inputfileInstrument & \hfuzz=500pt filename & \hfuzz=500pt \\
\hfuzz=500pt\includegraphics[width=1em]{element.pdf}~selectDataField & \hfuzz=500pt uint & \hfuzz=500pt select a data column for decomposition\\
\hfuzz=500pt\includegraphics[width=1em]{element-mustset.pdf}~inputfileWavelet & \hfuzz=500pt filename & \hfuzz=500pt wavelet coefficients\\
\hfuzz=500pt\includegraphics[width=1em]{element-mustset.pdf}~level & \hfuzz=500pt uint & \hfuzz=500pt level of decomposition\\
\hline
\end{tabularx}

\clearpage
%==================================
\subsection{LocalLevelFrame2StarCamera}\label{LocalLevelFrame2StarCamera}
Compute rotation (\file{StarCamera file}{instrument}) from local level frame (ellipsoidal north, east, down)
to TRF for positions given in \configFile{inputfileInstrument}{instrument} (first 3 data columns).


\keepXColumns
\begin{tabularx}{\textwidth}{N T A}
\hline
Name & Type & Annotation\\
\hline
\hfuzz=500pt\includegraphics[width=1em]{element-mustset.pdf}~outputfileStarCamera & \hfuzz=500pt filename & \hfuzz=500pt rotation matrix from local level frame (ellipsoidal north, east, down) to TRF\\
\hfuzz=500pt\includegraphics[width=1em]{element-mustset.pdf}~inputfileInstrument & \hfuzz=500pt filename & \hfuzz=500pt origin of local level frame\\
\hfuzz=500pt\includegraphics[width=1em]{element.pdf}~constantOriginPerArc & \hfuzz=500pt boolean & \hfuzz=500pt use constant origin for all epochs of an arc (median position)\\
\hfuzz=500pt\includegraphics[width=1em]{element.pdf}~R & \hfuzz=500pt double & \hfuzz=500pt reference radius for ellipsoidal coordinates\\
\hfuzz=500pt\includegraphics[width=1em]{element.pdf}~inverseFlattening & \hfuzz=500pt double & \hfuzz=500pt reference flattening for ellipsoidal coordinates, 0: spherical coordinates\\
\hline
\end{tabularx}

This program is \reference{parallelized}{general.parallelization}.
\clearpage
%==================================
\section{Programs: KalmanFilter}
\subsection{KalmanBuildNormals}\label{KalmanBuildNormals}
This program sets up normal equations based on \configClass{observation}{observationType}
for short-term gravity field variations.
It computes the normal equations based on the intervals $i \in \{1, ..., N\}$ given in the \configFile{arcList}{arcList}.
It sets up the least squares adjustment
\begin{equation}
    \begin{bmatrix}
    \mathbf{l}_1 \\
    \mathbf{l}_2 \\
    \vdots \\
    \mathbf{l}_N \\
  \end{bmatrix}
  =
  \begin{bmatrix}
    \mathbf{A}_1  &  & & \\
    & \mathbf{A}_2  & &\\
    &  & \ddots & \\
    & & & \mathbf{A}_N \\
  \end{bmatrix}
  \begin{bmatrix}
    \mathbf{x}^{(1)} \\
    \mathbf{x}^{(2)} \\
    \vdots \\
    \mathbf{x}^{(N)} \\
  \end{bmatrix}
  +
  \begin{bmatrix}
    \mathbf{e}_1 \\
    \mathbf{e}_2 \\
    \vdots \\
    \mathbf{e}_N \\
  \end{bmatrix},
\end{equation}
and subsequently computes the normal equations $\mathbf{N}_i, \mathbf{n}_i$ for each interval.
If \config{eliminateNonGravityParameters} is true, all non-gravity parameters are eliminated before the normals
are written to \configFile{outputfileNormalEquation}{normalEquation}.
For each time interval in \config{arcList} a single \file{normal equation file}{normalEquation} is written.

This program computes the input normals for \program{KalmanFilter} and \program{KalmanSmootherLeastSquares}.


\keepXColumns
\begin{tabularx}{\textwidth}{N T A}
\hline
Name & Type & Annotation\\
\hline
\hfuzz=500pt\includegraphics[width=1em]{element-mustset.pdf}~outputfileNormalEquation & \hfuzz=500pt filename & \hfuzz=500pt outputfile for normal equations\\
\hfuzz=500pt\includegraphics[width=1em]{element-mustset.pdf}~observation & \hfuzz=500pt \hyperref[observationType]{observation} & \hfuzz=500pt \\
\hfuzz=500pt\includegraphics[width=1em]{element-mustset.pdf}~inputfileArcList & \hfuzz=500pt filename & \hfuzz=500pt list to correspond points of time to arc numbers\\
\hfuzz=500pt\includegraphics[width=1em]{element.pdf}~eliminateNonGravityParameters & \hfuzz=500pt boolean & \hfuzz=500pt eliminate additional parameters from normals, 0: all parameter are saved\\
\hline
\end{tabularx}

This program is \reference{parallelized}{general.parallelization}.
\clearpage
%==================================
\subsection{KalmanFilter}\label{KalmanFilter}
The program computes time variable gravity fields using the Kalman filter approach of

Kurtenbach, E., Eicker, A., Mayer-Gürr, T., Holschneider, M., Hayn, M., Fuhrmann, M., and Kusche, J. (2012).
Improved daily GRACE gravity field solutions using a Kalman smoother. Journal of Geodynamics, 59–60, 39–48.
\url{https://doi.org/10.1016/j.jog.2012.02.006}.

The updated state $\mathbf{x}_t^+$ is determined by solving the least squares adjustment
\begin{equation}
\mathbf{l}_t = \mathbf{A}_t \mathbf{x}_t + \mathbf{e}_t \hspace{25pt} \mathbf{e}_t \sim \mathcal{N}(0, \mathbf{R}_t)\\
\mathbf{B} \mathbf{x}^+_{t-1} = \mathbf{I} \mathbf{x}_t + \mathbf{v}_t\hspace{25pt} \mathbf{v} \sim \mathcal{N}(0,\mathbf{Q} + \mathbf{B} \mathbf{P}^+_{t-1}\mathbf{B}^T).
\end{equation}
In normal equation form this can be written as
\begin{equation}
\hat{\mathbf{x}}_t = \mathbf{x}^+_t = (\mathbf{N}_t + \mathbf{P}^{-^{-1}}_t)^{-1}(\mathbf{n}_t + \mathbf{P}^{-^{-1}}_t \mathbf{x}^-_t),
\end{equation}
where $\mathbf{x}_t^- = \mathbf{B} \mathbf{x}^+_{t-1}$ and $\mathbf{P}_t^{-} = \mathbf{Q} + \mathbf{B} \mathbf{P}^+_{t-1}\mathbf{B}^T$
are the predicted state and its covariance matrix.

The process dynamic $\mathbf{B}, \mathbf{Q}$ is represented as an \reference{autoregressive model}{fundamentals.autoregressiveModel},
and passed to the program through \configFile{inputfileAutoregressiveModel}{matrix}.
The sequence of normal equations $\mathbf{N}_t, \mathbf{n}_t$ are given as list of \configFile{inputfileNormalEquations}{normalEquation},
which can be generated using \configClass{loops}{loopType}.
In the same way, the \file{matrix files}{matrix} for \config{outputfileUpdatedState} and \config{inputfileUpdatedStateCovariance}
can also be specified using \configClass{loops}{loopType}.

If no \configFile{inputfileInitialState}{matrix} is set, a zero vector with appropriate dimensions is used.
The \configFile{inputfileInitialStateCovarianceMatrix}{matrix} however must be given.

See also \program{KalmanBuildNormals}, \program{KalmanSmoother}.


\keepXColumns
\begin{tabularx}{\textwidth}{N T A}
\hline
Name & Type & Annotation\\
\hline
\hfuzz=500pt\includegraphics[width=1em]{element-mustset-unbounded.pdf}~outputfileUpdatedState & \hfuzz=500pt filename & \hfuzz=500pt estimated state x+ (nx1-matrix)\\
\hfuzz=500pt\includegraphics[width=1em]{element-unbounded.pdf}~outputfileUpdatedStateCovarianceMatrix & \hfuzz=500pt filename & \hfuzz=500pt estimated state' s covariance matrix Cov(x+)\\
\hfuzz=500pt\includegraphics[width=1em]{element-mustset-unbounded.pdf}~inputfileNormalEquations & \hfuzz=500pt filename & \hfuzz=500pt normal equations input file\\
\hfuzz=500pt\includegraphics[width=1em]{element.pdf}~inputfileInitialState & \hfuzz=500pt filename & \hfuzz=500pt initial state x0\\
\hfuzz=500pt\includegraphics[width=1em]{element-mustset.pdf}~inputfileInitialStateCovarianceMatrix & \hfuzz=500pt filename & \hfuzz=500pt initial state's covariance matrix Cov(x0)\\
\hfuzz=500pt\includegraphics[width=1em]{element-mustset.pdf}~inputfileAutoregressiveModel & \hfuzz=500pt filename & \hfuzz=500pt file name of autoregressive model\\
\hline
\end{tabularx}

\clearpage
%==================================
\subsection{KalmanSmoother}\label{KalmanSmoother}
Apply the Rauch-Tung-Striebel smoother to a gravity field time series computed by \program{KalmanFilter}.
This is the implementation of the approach presented in

Kurtenbach, E., Eicker, A., Mayer-Gürr, T., Holschneider, M., Hayn, M., Fuhrmann, M., and Kusche, J. (2012).
Improved daily GRACE gravity field solutions using a Kalman smoother. Journal of Geodynamics, 59–60, 39–48.
\url{https://doi.org/10.1016/j.jog.2012.02.006}.

The result has zero phase and the squared magnitude response of \configFile{inputfileAutoregressiveModel}{matrix}
(see \reference{autoregressiveModel}{fundamentals.autoregressiveModel} for details).
\configFile{inputfileUpdatedState}{matrix} and \configFile{inputfileUpdatedStateCovariance}{matrix}
are the output of a \program{KalmanFilter} forward sweep.
The matrix files for\configFile{outputfileUpdatedState}{matrix}, \configFile{inputfileUpdatedState}{matrix}
and \configFile{inputfileUpdatedStateCovariance}{matrix} can also be specified using \configClass{loops}{loopType}.

See also \program{KalmanBuildNormals}, \program{KalmanFilter} and \program{KalmanSmootherLeastSquares}.


\keepXColumns
\begin{tabularx}{\textwidth}{N T A}
\hline
Name & Type & Annotation\\
\hline
\hfuzz=500pt\includegraphics[width=1em]{element-mustset-unbounded.pdf}~outputfileState & \hfuzz=500pt filename & \hfuzz=500pt estimated parameters (nx1-matrix)\\
\hfuzz=500pt\includegraphics[width=1em]{element-unbounded.pdf}~outputfileStateCovarianceMatrix & \hfuzz=500pt filename & \hfuzz=500pt estimated parameters' covariance matrix\\
\hfuzz=500pt\includegraphics[width=1em]{element-mustset-unbounded.pdf}~inputfileUpdatedState & \hfuzz=500pt filename & \hfuzz=500pt \\
\hfuzz=500pt\includegraphics[width=1em]{element-mustset-unbounded.pdf}~inputfileUpdatedStateCovarianceMatrix & \hfuzz=500pt filename & \hfuzz=500pt \\
\hfuzz=500pt\includegraphics[width=1em]{element-mustset.pdf}~inputfileAutoregressiveModel & \hfuzz=500pt filename & \hfuzz=500pt file name of autoregressive model\\
\hline
\end{tabularx}

\clearpage
%==================================
\subsection{KalmanSmootherLeastSquares}\label{KalmanSmootherLeastSquares}
This program estimates temporal gravity field variations with a constraint least squares adjustment.
Prior information is introduced by means of a \configClass{autoregressiveModelSequence}{autoregressiveModelSequenceType}
which represent a stationary random process (see the \reference{autoregressive model description}{fundamentals.autoregressiveModel}) for details.

The output files for the estimated gravity field (\configFile{outputfileSolution}{matrix}), the
corresponding standard deviations (\configFile{outputfileSigmax}{matrix}) and the full covariance matrix
(\configFile{outputfileCovariance}{matrix}) can be specified using \configClass{loops}{loopType}.
Similarly, the \configFile{inputfileNormalEquations}{normalEquation}
can also be specified using \configClass{loops}{loopType}.

See also \program{KalmanBuildNormals}, \program{KalmanFilter} and\program{KalmanSmoother}


\keepXColumns
\begin{tabularx}{\textwidth}{N T A}
\hline
Name & Type & Annotation\\
\hline
\hfuzz=500pt\includegraphics[width=1em]{element-mustset-unbounded.pdf}~outputfileSolution & \hfuzz=500pt filename & \hfuzz=500pt file name of solution vector (use time tags)\\
\hfuzz=500pt\includegraphics[width=1em]{element-unbounded.pdf}~outputfileSigmax & \hfuzz=500pt filename & \hfuzz=500pt file name of sigma vector (use time tags)\\
\hfuzz=500pt\includegraphics[width=1em]{element-unbounded.pdf}~outputfileCovariance & \hfuzz=500pt filename & \hfuzz=500pt file name of full covariance matrix (use time tags)\\
\hfuzz=500pt\includegraphics[width=1em]{element-mustset-unbounded.pdf}~inputfileNormalEquations & \hfuzz=500pt filename & \hfuzz=500pt input normal equations (loopTime will be expanded)\\
\hfuzz=500pt\includegraphics[width=1em]{element-mustset.pdf}~autoregressiveModelSequence & \hfuzz=500pt \hyperref[autoregressiveModelSequenceType]{autoregressiveModelSequence} & \hfuzz=500pt file containing AR model for spatiotemporal constraint\\
\hline
\end{tabularx}

This program is \reference{parallelized}{general.parallelization}.
\clearpage
%==================================
\section{Programs: Misc}
\subsection{DigitalFilter2FrequencyResponse}\label{DigitalFilter2FrequencyResponse}
Compute amplitude-, phase-, group delay and frequency response of a \configClass{digitalFilter}{digitalFilterType} cascade.
The \configFile{outputfileResponse}{matrix} is a matrix with following columns:
freq $[Hz]$, ampl, phase $[rad]$, group delay $[-]$, real, imag.

When \config{unwrapPhase} is set to true, $2\pi$ jumps of the phase response are removed before writing the output to file.

The response of the filter cascade is given by the product of each individual frequency response:
\begin{equation}
  H(f) = \prod_f H_j(f).
\end{equation}
Amplitude and phase response are computed from the frequency response via
\begin{equation}
  A(f) = |H(f)| \hspace{5pt}\text{and}\hspace{5pt} \Phi(f) = \arctan \frac{\mathcal{I}(H(f))}{\mathcal{R}(H(f))}.
\end{equation}
The group delay is computed by numerically differentiating the phase response
\begin{equation}
  \tau_g(f_k) = \frac{1}{2} \left[\frac{\Phi(f_k) - \Phi(f_{k-1})}{2\pi(f_k-f_{k-1})} + \frac{\Phi(f_{k+1}) - \Phi(f_{k})}{2\pi(f_{k+1}-f_{k})}\right] \approx \frac{d\Phi}{df}\frac{df}{d\omega}.
\end{equation}
The frequency vector for a \config{length} $N$ and a \config{sampling} $\Delta t$ is given by
\begin{equation}
  f_k = \frac{k}{N \Delta t}, \hspace{15pt} k \in \{0, \dots, \left\lfloor\frac{N+2}{2}\right\rfloor-1\}.
\end{equation}

See also \program{DigitalFilter2ImpulseResponse}.


\keepXColumns
\begin{tabularx}{\textwidth}{N T A}
\hline
Name & Type & Annotation\\
\hline
\hfuzz=500pt\includegraphics[width=1em]{element-mustset.pdf}~outputfileResponse & \hfuzz=500pt filename & \hfuzz=500pt columns: freq [Hz], ampl, phase [rad], group delay [-], real, imag\\
\hfuzz=500pt\includegraphics[width=1em]{element-mustset-unbounded.pdf}~digitalFilter & \hfuzz=500pt \hyperref[digitalFilterType]{digitalFilter} & \hfuzz=500pt \\
\hfuzz=500pt\includegraphics[width=1em]{element.pdf}~length & \hfuzz=500pt uint & \hfuzz=500pt length of the data series in time domain\\
\hfuzz=500pt\includegraphics[width=1em]{element.pdf}~sampling & \hfuzz=500pt double & \hfuzz=500pt sampling to determine frequency [seconds]\\
\hfuzz=500pt\includegraphics[width=1em]{element.pdf}~skipZeroFrequency & \hfuzz=500pt boolean & \hfuzz=500pt omit zero frequency when writing to file\\
\hfuzz=500pt\includegraphics[width=1em]{element.pdf}~unwrapPhase & \hfuzz=500pt boolean & \hfuzz=500pt unwrap phase response\\
\hline
\end{tabularx}

\clearpage
%==================================
\subsection{DigitalFilter2ImpulseResponse}\label{DigitalFilter2ImpulseResponse}
Impulse response of a \configClass{digitalFilter}{digitalFilterType} cascade.
The impulse response is computed by filtering a sequence with \config{length} samples and a unit impulse at index \config{pulseLag}.

The \configFile{outputfileResponse}{matrix} is a matrix with the time stamp (zero at \config{pulseLag})
in the first column and the impulse response $h_k$ in the second column.

See also \program{DigitalFilter2FrequencyResponse}.


\keepXColumns
\begin{tabularx}{\textwidth}{N T A}
\hline
Name & Type & Annotation\\
\hline
\hfuzz=500pt\includegraphics[width=1em]{element-mustset.pdf}~outputfileResponse & \hfuzz=500pt filename & \hfuzz=500pt columns: time [seconds], response\\
\hfuzz=500pt\includegraphics[width=1em]{element-mustset-unbounded.pdf}~digitalFilter & \hfuzz=500pt \hyperref[digitalFilterType]{digitalFilter} & \hfuzz=500pt \\
\hfuzz=500pt\includegraphics[width=1em]{element.pdf}~length & \hfuzz=500pt uint & \hfuzz=500pt length of the impulse response\\
\hfuzz=500pt\includegraphics[width=1em]{element.pdf}~pulseLag & \hfuzz=500pt uint & \hfuzz=500pt start of the pulse in the data series\\
\hfuzz=500pt\includegraphics[width=1em]{element.pdf}~sampling & \hfuzz=500pt double & \hfuzz=500pt [seconds]\\
\hline
\end{tabularx}

\clearpage
%==================================
\subsection{EarthOrientationParameterTimeSeries}\label{EarthOrientationParameterTimeSeries}
Computes a \configClass{timeSeries}{timeSeriesType} (GPS time) of Earth Orientation Parameter (EOP).
The \file{instrument file}{instrument} (MISCVALUES) contains the elements at each epoch in the following order:
\begin{itemize}
\item $x_p$ [rad]
\item $y_p$ [rad]
\item $s_p$ [rad]
\item $UT1-UTC$ [seconds]
\item length of day (LOD) [seconds]
\item $X$ [rad]
\item $Y$ [rad]
\item $S$ [rad]
\end{itemize}
The values are in situ values with all corrections and models applied. The time series can be used to
precompute Earth rotation with a low temporal resolution (e.g. 10 min) and reuse the file in
\configClass{earthRotation:file}{earthRotationType:file} to interpolate the data to the needed epochs
(e.g. to rotate orbit data). As some Earth rotation models are quite slow this can accelerate the computation.


\keepXColumns
\begin{tabularx}{\textwidth}{N T A}
\hline
Name & Type & Annotation\\
\hline
\hfuzz=500pt\includegraphics[width=1em]{element-mustset.pdf}~outputfileEOP & \hfuzz=500pt filename & \hfuzz=500pt each row: mjd(GPS), xp, yp, sp, dUT1, LOD, X, Y, S\\
\hfuzz=500pt\includegraphics[width=1em]{element-mustset.pdf}~earthRotation & \hfuzz=500pt \hyperref[earthRotationType]{earthRotation} & \hfuzz=500pt \\
\hfuzz=500pt\includegraphics[width=1em]{element-mustset-unbounded.pdf}~timeSeries & \hfuzz=500pt \hyperref[timeSeriesType]{timeSeries} & \hfuzz=500pt \\
\hline
\end{tabularx}

This program is \reference{parallelized}{general.parallelization}.
\clearpage
%==================================
\subsection{EarthRotaryVectorTimeSeries}\label{EarthRotaryVectorTimeSeries}
Computes a \configFile{outputfileTimeSeries}{instrument} of Earth's rotary axis
and its temporal derivative at \configClass{timeSeries}{timeSeriesType} (GPS time).
The \file{instrument file}{instrument} (MISCVALUES) contains the elements at each epoch in the following order:
\begin{itemize}
\item $\omega_x [rad/s]$
\item $\omega_y [rad/s]$
\item $\omega_z [rad/s]$
\item $\dot{\omega}_x [rad/s^2]$
\item $\dot{\omega}_y [rad/s^2]$
\item $\dot{\omega}_z [rad/s^2]$.
\end{itemize}


\keepXColumns
\begin{tabularx}{\textwidth}{N T A}
\hline
Name & Type & Annotation\\
\hline
\hfuzz=500pt\includegraphics[width=1em]{element-mustset.pdf}~outputfileTimeSeries & \hfuzz=500pt filename & \hfuzz=500pt wx, wy, wz [rad], dwx, dwy, dwz [rad/s\textasciicircum{}2]\\
\hfuzz=500pt\includegraphics[width=1em]{element-mustset.pdf}~earthRotation & \hfuzz=500pt \hyperref[earthRotationType]{earthRotation} & \hfuzz=500pt \\
\hfuzz=500pt\includegraphics[width=1em]{element-mustset-unbounded.pdf}~timeSeries & \hfuzz=500pt \hyperref[timeSeriesType]{timeSeries} & \hfuzz=500pt \\
\hfuzz=500pt\includegraphics[width=1em]{element.pdf}~inTRF & \hfuzz=500pt boolean & \hfuzz=500pt terrestrial reference frame, otherwise celestial\\
\hline
\end{tabularx}

This program is \reference{parallelized}{general.parallelization}.
\clearpage
%==================================
\subsection{EclipseFactor2GriddedData}\label{EclipseFactor2GriddedData}
This program converts the output of a \configClass{eclipse}{eclipseType} model on a given
\configClass{grid}{gridType}. The time for the evaluation can be specified in \config{time}.
The values will be saved together with points expressed as ellipsoidal coordinates
(longitude, latitude, height) based on a reference ellipsoid with parameters \config{R}
and \config{inverseFlattening}.


\keepXColumns
\begin{tabularx}{\textwidth}{N T A}
\hline
Name & Type & Annotation\\
\hline
\hfuzz=500pt\includegraphics[width=1em]{element-mustset.pdf}~outputfileGriddedData & \hfuzz=500pt filename & \hfuzz=500pt eclipse factor\\
\hfuzz=500pt\includegraphics[width=1em]{element-mustset-unbounded.pdf}~grid & \hfuzz=500pt \hyperref[gridType]{grid} & \hfuzz=500pt \\
\hfuzz=500pt\includegraphics[width=1em]{element-mustset.pdf}~eclipse & \hfuzz=500pt \hyperref[eclipseType]{eclipse} & \hfuzz=500pt \\
\hfuzz=500pt\includegraphics[width=1em]{element-mustset.pdf}~ephemerides & \hfuzz=500pt \hyperref[ephemeridesType]{ephemerides} & \hfuzz=500pt \\
\hfuzz=500pt\includegraphics[width=1em]{element.pdf}~earthRotation & \hfuzz=500pt \hyperref[earthRotationType]{earthRotation} & \hfuzz=500pt \\
\hfuzz=500pt\includegraphics[width=1em]{element-mustset.pdf}~time & \hfuzz=500pt time & \hfuzz=500pt \\
\hfuzz=500pt\includegraphics[width=1em]{element.pdf}~R & \hfuzz=500pt double & \hfuzz=500pt reference radius for ellipsoidal coordinates on output\\
\hfuzz=500pt\includegraphics[width=1em]{element.pdf}~inverseFlattening & \hfuzz=500pt double & \hfuzz=500pt reference flattening for ellipsoidal coordinates on output, 0: spherical coordinates\\
\hline
\end{tabularx}

This program is \reference{parallelized}{general.parallelization}.
\clearpage
%==================================
\subsection{FilterMatrixWindowedPotentialCoefficients}\label{FilterMatrixWindowedPotentialCoefficients}
Create a spherical harmonic window matrix. The window matrix $\mathbf{W}$ is generated in space domain through
spherical harmonic synthesis and analysis matrices.
The resulting linear operator can be written as
\begin{equation}
\mathbf{W} = \mathbf{K} \mathbf{A} \mathbf{\Omega} \mathbf{S} \mathbf{K}^{-1}.
\end{equation}
Here, $\mathbf{K}$ is a diagonal matrix with the \configClass{kernel}{kernelType} coefficients on the main diagonal,
$\mathbf{S}$ is the spherical harmonic synthesis matrix, $\mathbf{\Omega}$ is defined by the values in
\file{inputfileGriddedData}{griddedData} and the
expression \config{value}, $\mathbf{A}$ is the spherical harmonic analysis matrix.
The resulting window matrix is written to a \file{matrix}{matrix} file.

The spherical harmonic degree range, and coefficient numbering are defined by
\config{minDegree}, \config{maxDegree}, and \configClass{numbering}{sphericalHarmonicsNumberingType}.

Note that a proper window function $\mathbf{\Omega}$ should contain values in the range [0, 1].
The window function $\mathbf{\Omega}$ can feature a smooth transition between 0 and 1 to avoid ringing effects.


\keepXColumns
\begin{tabularx}{\textwidth}{N T A}
\hline
Name & Type & Annotation\\
\hline
\hfuzz=500pt\includegraphics[width=1em]{element.pdf}~outputfileWindowMatrix & \hfuzz=500pt filename & \hfuzz=500pt \\
\hfuzz=500pt\includegraphics[width=1em]{element-mustset.pdf}~inputfileGriddedData & \hfuzz=500pt filename & \hfuzz=500pt gridded data which defines the window function in space domain\\
\hfuzz=500pt\includegraphics[width=1em]{element-mustset.pdf}~value & \hfuzz=500pt expression & \hfuzz=500pt expression to compute the window function (input columns are named data0, data1, ...)\\
\hfuzz=500pt\includegraphics[width=1em]{element-mustset.pdf}~kernel & \hfuzz=500pt \hyperref[kernelType]{kernel} & \hfuzz=500pt kernel for windowing\\
\hfuzz=500pt\includegraphics[width=1em]{element.pdf}~minDegree & \hfuzz=500pt uint & \hfuzz=500pt \\
\hfuzz=500pt\includegraphics[width=1em]{element-mustset.pdf}~maxDegree & \hfuzz=500pt uint & \hfuzz=500pt \\
\hfuzz=500pt\includegraphics[width=1em]{element.pdf}~GM & \hfuzz=500pt double & \hfuzz=500pt Geocentric gravitational constant\\
\hfuzz=500pt\includegraphics[width=1em]{element.pdf}~R & \hfuzz=500pt double & \hfuzz=500pt reference radius\\
\hfuzz=500pt\includegraphics[width=1em]{element-mustset.pdf}~numbering & \hfuzz=500pt \hyperref[sphericalHarmonicsNumberingType]{sphericalHarmonicsNumbering} & \hfuzz=500pt numbering scheme for solution vector\\
\hline
\end{tabularx}

This program is \reference{parallelized}{general.parallelization}.
\clearpage
%==================================
\subsection{FunctionsCalculate}\label{FunctionsCalculate}
This program manipulates \file{matrix files}{matrix} with data in columns.
If several \config{inputfile}s are given the data columns are copied side by side.
All \config{inputfile}s must contain the same number of rows.
The columns are enumerated by \verb|data0|,~\verb|data1|,~\ldots.

The content of \configFile{outputfile}{matrix} is controlled by \config{outColumn}.
The algorithm to compute the output is as follows:
The expressions in \config{outColumn} are evaluated once for each row of the input.
The variables \verb|data0|,~\verb|data1|,~\ldots are replaced by the according values from the input columns before.
Additional variables are available, e.g. \verb|index|, \verb|data0rms|, see~\reference{dataVariables}{general.parser:dataVariables}.
If no \config{outColumn} are specified all input columns are used instead directly.

For a simplified handling \config{constant}s can be defined by \verb|name=value|, e.g. \verb|annual=365.25|.
It is also possible to estimate \config{parameter}s in a least squares adjustment.
The \config{leastSquares} serves as template for observation equations for every row.
The expression \config{leastSquares} is evaluated for each row in the \config{inputfile}.
The variables \verb|data0|,~\verb|data1|,~\ldots are replaced by the according values from the input columns before.
In the next step the parameters are estimated in order to minimize the expressions in \config{leastSquares}
in the sense of least squares.

Afterwards complete rows are removed if one of the \config{removalCriteria} expressions for this row evaluates true (not zero).

An extra \config{statistics} file can be generated with one row of data. For the computation of the \config{outColumn} values
all~\reference{dataVariables}{general.parser:dataVariables} are available (e.g. \verb|data3mean|, \verb|data4std|)
inclusively the \config{constant}s and estimated \config{parameter}s but without the \verb|data0|,~\verb|data1|,~\ldots itself.
The variables and the numbering of the columns refers to the \configFile{outputfile}{matrix}.

First example: To calculate the mean of two values at each row set \config{outColumn} to \verb|0.5*(data1+data0)|.

Second example: An input file contain a column with times and a column with values.
To remove a trend from the values define the \config{parameter}s \verb|trend| and \verb|bias|.
The observation equation in \config{leastSquares} is \verb|data1 - (trend*data0+bias)|.
For output you can define the following columns for example:
\begin{itemize}
\item \config{outColumn}=\verb|data0|: points in time.
\item \config{outColumn}=\verb|data1|: the values itself.
\item \config{outColumn}=\verb|trend*data0+bias|: the linear fit.
\item \config{outColumn}=\verb|data1-trend*data0-bias|: the residuals.
\end{itemize}
The extra statistics file could contain in this case:
\begin{itemize}
\item \config{outColumn}=\verb|data0max-data0min|: time span.
\item \config{outColumn}=\verb|bias|: estimated parameter.
\item \config{outColumn}=\verb|trend|: estimated parameter.
\item \config{outColumn}=\verb|data3rms|: root mean square of the residuals.
\end{itemize}

See also \program{InstrumentArcCalculate}, \program{GriddedDataCalculate}, \program{MatrixCalculate}.


\keepXColumns
\begin{tabularx}{\textwidth}{N T A}
\hline
Name & Type & Annotation\\
\hline
\hfuzz=500pt\includegraphics[width=1em]{element.pdf}~outputfile & \hfuzz=500pt filename & \hfuzz=500pt \\
\hfuzz=500pt\includegraphics[width=1em]{element-mustset-unbounded.pdf}~inputfile & \hfuzz=500pt filename & \hfuzz=500pt \\
\hfuzz=500pt\includegraphics[width=1em]{element-unbounded.pdf}~constant & \hfuzz=500pt expression & \hfuzz=500pt define a constant by name=value\\
\hfuzz=500pt\includegraphics[width=1em]{element-unbounded.pdf}~parameter & \hfuzz=500pt expression & \hfuzz=500pt define a parameter by name[=value]\\
\hfuzz=500pt\includegraphics[width=1em]{element-unbounded.pdf}~leastSquares & \hfuzz=500pt expression & \hfuzz=500pt try to minimize the expression by adjustment of the parameters\\
\hfuzz=500pt\includegraphics[width=1em]{element-unbounded.pdf}~removalCriteria & \hfuzz=500pt expression & \hfuzz=500pt row is removed if one criterion evaluates true.\\
\hfuzz=500pt\includegraphics[width=1em]{element-unbounded.pdf}~outColumn & \hfuzz=500pt expression & \hfuzz=500pt expression to compute output columns (input columns are named data0, data1, ...)\\
\hfuzz=500pt\includegraphics[width=1em]{element.pdf}~statistics & \hfuzz=500pt sequence & \hfuzz=500pt \\
\hfuzz=500pt\includegraphics[width=1em]{connector.pdf}\includegraphics[width=1em]{element-mustset.pdf}~outputfile & \hfuzz=500pt filename & \hfuzz=500pt matrix with one row, columns are user defined\\
\hfuzz=500pt\includegraphics[width=1em]{connector.pdf}\includegraphics[width=1em]{element-mustset-unbounded.pdf}~outColumn & \hfuzz=500pt expression & \hfuzz=500pt expression to compute statistics columns, data* are the outputColumns\\
\hline
\end{tabularx}

\clearpage
%==================================
\subsection{Grs2PotentialCoefficients}\label{Grs2PotentialCoefficients}
This program creates potential coefficients from the defining constants
of a Geodetic Reference System (GRS). The potential coeffiencts excludes the centrifugal part.
The form of the reference ellipsoid is either determined by the dynamical form factor \config{J2},
or the geometric \config{inverseFlattening}. One of those form parameters must be specified.

The default values create the GRS80.


\keepXColumns
\begin{tabularx}{\textwidth}{N T A}
\hline
Name & Type & Annotation\\
\hline
\hfuzz=500pt\includegraphics[width=1em]{element-mustset.pdf}~outputfilePotentialCoefficients & \hfuzz=500pt filename & \hfuzz=500pt \\
\hfuzz=500pt\includegraphics[width=1em]{element-mustset.pdf}~maxDegree & \hfuzz=500pt uint & \hfuzz=500pt \\
\hfuzz=500pt\includegraphics[width=1em]{element.pdf}~GM & \hfuzz=500pt double & \hfuzz=500pt Geocentric gravitational constant\\
\hfuzz=500pt\includegraphics[width=1em]{element.pdf}~R & \hfuzz=500pt double & \hfuzz=500pt reference radius\\
\hfuzz=500pt\includegraphics[width=1em]{element.pdf}~omega & \hfuzz=500pt double & \hfuzz=500pt Angular velocity of rotation\\
\hfuzz=500pt\includegraphics[width=1em]{element.pdf}~J2 & \hfuzz=500pt double & \hfuzz=500pt Dynamical form factor\\
\hfuzz=500pt\includegraphics[width=1em]{element.pdf}~inverseFlattening & \hfuzz=500pt double & \hfuzz=500pt Geometric inverse flattening of reference ellipsoid (0: sphere, ignored when J2 is set)\\
\hline
\end{tabularx}

\clearpage
%==================================
\subsection{Kaula2SigmaPotentialCoefficients}\label{Kaula2SigmaPotentialCoefficients}
Create signal standard deviations of potential coefficients according Kaula's rule of thumb
\begin{equation}
  \sigma_n = \frac{f}{n^p},
\end{equation}
with the degree $n$, the \config{factor} $f$, and the \config{power} $p$.

The standard deviations are written as formal errors of
 \configFile{outputfilePotentialCoefficients}{potentialCoefficients}.


\keepXColumns
\begin{tabularx}{\textwidth}{N T A}
\hline
Name & Type & Annotation\\
\hline
\hfuzz=500pt\includegraphics[width=1em]{element-mustset.pdf}~outputfilePotentialCoefficients & \hfuzz=500pt filename & \hfuzz=500pt \\
\hfuzz=500pt\includegraphics[width=1em]{element.pdf}~minDegree & \hfuzz=500pt uint & \hfuzz=500pt \\
\hfuzz=500pt\includegraphics[width=1em]{element-mustset.pdf}~maxDegree & \hfuzz=500pt uint & \hfuzz=500pt \\
\hfuzz=500pt\includegraphics[width=1em]{element.pdf}~GM & \hfuzz=500pt double & \hfuzz=500pt Geocentric gravitational constant\\
\hfuzz=500pt\includegraphics[width=1em]{element.pdf}~R & \hfuzz=500pt double & \hfuzz=500pt reference radius\\
\hfuzz=500pt\includegraphics[width=1em]{element.pdf}~power & \hfuzz=500pt double & \hfuzz=500pt sigma = factor/degree\textasciicircum{}power\\
\hfuzz=500pt\includegraphics[width=1em]{element.pdf}~factor & \hfuzz=500pt double & \hfuzz=500pt sigma = factor/degree\textasciicircum{}power\\
\hline
\end{tabularx}

\clearpage
%==================================
\subsection{Kernel2Coefficients}\label{Kernel2Coefficients}
This program computes and returns the coefficients and inverse coefficients of a \configClass{kernel}{kernelType}
from from \config{minDegree} to \config{maxDegree} at a given \config{height}.

The main purpose is for visualization with \program{PlotGraph}.


\keepXColumns
\begin{tabularx}{\textwidth}{N T A}
\hline
Name & Type & Annotation\\
\hline
\hfuzz=500pt\includegraphics[width=1em]{element-mustset.pdf}~outputfileMatrix & \hfuzz=500pt filename & \hfuzz=500pt matrix with columns degree, coefficients and inverse coefficients\\
\hfuzz=500pt\includegraphics[width=1em]{element-mustset.pdf}~kernel & \hfuzz=500pt \hyperref[kernelType]{kernel} & \hfuzz=500pt \\
\hfuzz=500pt\includegraphics[width=1em]{element.pdf}~minDegree & \hfuzz=500pt uint & \hfuzz=500pt minimum degree of returned coefficients\\
\hfuzz=500pt\includegraphics[width=1em]{element-mustset.pdf}~maxDegre & \hfuzz=500pt uint & \hfuzz=500pt compute coefficients up to maxDegree\\
\hfuzz=500pt\includegraphics[width=1em]{element.pdf}~height & \hfuzz=500pt double & \hfuzz=500pt evaluate kernel at R+height [m]\\
\hfuzz=500pt\includegraphics[width=1em]{element.pdf}~R & \hfuzz=500pt double & \hfuzz=500pt reference radius\\
\hline
\end{tabularx}

\clearpage
%==================================
\subsection{Kernel2SigmaPotentialCoefficients}\label{Kernel2SigmaPotentialCoefficients}
Create variances of spherical harmonics by convolution a kernel with white noise,
e.g. to display filter coefficients of a Gaussian filter.
The coefficients are written as formal errors of \configFile{outputfilePotentialCoefficients}{potentialCoefficients}.


\keepXColumns
\begin{tabularx}{\textwidth}{N T A}
\hline
Name & Type & Annotation\\
\hline
\hfuzz=500pt\includegraphics[width=1em]{element-mustset.pdf}~outputfilePotentialCoefficients & \hfuzz=500pt filename & \hfuzz=500pt \\
\hfuzz=500pt\includegraphics[width=1em]{element-mustset.pdf}~kernel & \hfuzz=500pt \hyperref[kernelType]{kernel} & \hfuzz=500pt \\
\hfuzz=500pt\includegraphics[width=1em]{element.pdf}~minDegree & \hfuzz=500pt uint & \hfuzz=500pt \\
\hfuzz=500pt\includegraphics[width=1em]{element.pdf}~maxDegree & \hfuzz=500pt uint & \hfuzz=500pt \\
\hfuzz=500pt\includegraphics[width=1em]{element.pdf}~GM & \hfuzz=500pt double & \hfuzz=500pt Geocentric gravitational constant\\
\hfuzz=500pt\includegraphics[width=1em]{element.pdf}~R & \hfuzz=500pt double & \hfuzz=500pt reference radius\\
\hfuzz=500pt\includegraphics[width=1em]{element.pdf}~factor & \hfuzz=500pt double & \hfuzz=500pt \\
\hline
\end{tabularx}

\clearpage
%==================================
\subsection{KernelEvaluate}\label{KernelEvaluate}
Compute \configClass{kernel}{kernelType} values for distant angles.
The main purpose is for visualization with \program{PlotGraph}.


\keepXColumns
\begin{tabularx}{\textwidth}{N T A}
\hline
Name & Type & Annotation\\
\hline
\hfuzz=500pt\includegraphics[width=1em]{element-mustset.pdf}~outputfileMatrix & \hfuzz=500pt filename & \hfuzz=500pt matrix with first column the angle [degree], second the kernel value\\
\hfuzz=500pt\includegraphics[width=1em]{element-mustset.pdf}~kernel & \hfuzz=500pt \hyperref[kernelType]{kernel} & \hfuzz=500pt \\
\hfuzz=500pt\includegraphics[width=1em]{element.pdf}~minAngle & \hfuzz=500pt angle & \hfuzz=500pt [degree]\\
\hfuzz=500pt\includegraphics[width=1em]{element.pdf}~maxAngle & \hfuzz=500pt angle & \hfuzz=500pt [degree]\\
\hfuzz=500pt\includegraphics[width=1em]{element.pdf}~sampling & \hfuzz=500pt angle & \hfuzz=500pt [degree]\\
\hfuzz=500pt\includegraphics[width=1em]{element.pdf}~height & \hfuzz=500pt double & \hfuzz=500pt evaluate at R+height [m]\\
\hfuzz=500pt\includegraphics[width=1em]{element.pdf}~R & \hfuzz=500pt double & \hfuzz=500pt reference radius\\
\hline
\end{tabularx}

\clearpage
%==================================
\subsection{MagneticField2GriddedData}\label{MagneticField2GriddedData}
Computes x, y, z of the magentic field vector.


\keepXColumns
\begin{tabularx}{\textwidth}{N T A}
\hline
Name & Type & Annotation\\
\hline
\hfuzz=500pt\includegraphics[width=1em]{element-mustset.pdf}~outputfileGriddedData & \hfuzz=500pt filename & \hfuzz=500pt x, y, z [Tesla = kg/A/s**2]\\
\hfuzz=500pt\includegraphics[width=1em]{element-mustset.pdf}~magnetosphere & \hfuzz=500pt \hyperref[magnetosphereType]{magnetosphere} & \hfuzz=500pt \\
\hfuzz=500pt\includegraphics[width=1em]{element-mustset-unbounded.pdf}~grid & \hfuzz=500pt \hyperref[gridType]{grid} & \hfuzz=500pt \\
\hfuzz=500pt\includegraphics[width=1em]{element.pdf}~time & \hfuzz=500pt time & \hfuzz=500pt \\
\hfuzz=500pt\includegraphics[width=1em]{element.pdf}~localReferenceFrame & \hfuzz=500pt boolean & \hfuzz=500pt local left handed reference frame (north, east, up)\\
\hfuzz=500pt\includegraphics[width=1em]{element.pdf}~R & \hfuzz=500pt double & \hfuzz=500pt reference radius for ellipsoidal coordinates on output\\
\hfuzz=500pt\includegraphics[width=1em]{element.pdf}~inverseFlattening & \hfuzz=500pt double & \hfuzz=500pt reference flattening for ellipsoidal coordinates on output, 0: spherical coordinates\\
\hline
\end{tabularx}

This program is \reference{parallelized}{general.parallelization}.
\clearpage
%==================================
\subsection{MatrixCalculate}\label{MatrixCalculate}
This program creates a \file{matrix}{matrix} from multiple matrices.
All \configClass{matrices}{matrixGeneratorType} are summed up. The size of the resulting matrix is exandeded to fit all matrices.
The class \configClass{matrixGenerator}{matrixGeneratorType} allows complex matrix operations before.


\keepXColumns
\begin{tabularx}{\textwidth}{N T A}
\hline
Name & Type & Annotation\\
\hline
\hfuzz=500pt\includegraphics[width=1em]{element-mustset.pdf}~outputfileMatrix & \hfuzz=500pt filename & \hfuzz=500pt \\
\hfuzz=500pt\includegraphics[width=1em]{element-mustset-unbounded.pdf}~matrix & \hfuzz=500pt \hyperref[matrixGeneratorType]{matrixGenerator} & \hfuzz=500pt \\
\hline
\end{tabularx}

\clearpage
%==================================
\subsection{ObservationEquations2Files}\label{ObservationEquations2Files}
This program computes the linearized and decorrelated equation system for each arc $i$:
\begin{equation}
\M l_i  = \M A_i \M x + \M B_i \M y_i + \M e_i
\end{equation}
using class \configClass{observation}{observationType} and writes $\M A_i$, $\M B_i$ and $\M l_i$ as \file{matrix}{matrix} files.


\keepXColumns
\begin{tabularx}{\textwidth}{N T A}
\hline
Name & Type & Annotation\\
\hline
\hfuzz=500pt\includegraphics[width=1em]{element.pdf}~outputfileObservationVector & \hfuzz=500pt filename & \hfuzz=500pt one file for each arc\\
\hfuzz=500pt\includegraphics[width=1em]{element.pdf}~outputfileDesignMatrix & \hfuzz=500pt filename & \hfuzz=500pt one file for each arc, without arc related parameters\\
\hfuzz=500pt\includegraphics[width=1em]{element.pdf}~outputfileDesignMatrixArc & \hfuzz=500pt filename & \hfuzz=500pt one file for each arc, arc related parameters\\
\hfuzz=500pt\includegraphics[width=1em]{element.pdf}~variableArc & \hfuzz=500pt string & \hfuzz=500pt variable with arc number\\
\hfuzz=500pt\includegraphics[width=1em]{element.pdf}~outputfileParameterNames & \hfuzz=500pt filename & \hfuzz=500pt without arc related parameters\\
\hfuzz=500pt\includegraphics[width=1em]{element-mustset.pdf}~observation & \hfuzz=500pt \hyperref[observationType]{observation} & \hfuzz=500pt \\
\hline
\end{tabularx}

This program is \reference{parallelized}{general.parallelization}.
\clearpage
%==================================
\subsection{PlatformCreate}\label{PlatformCreate}
Create a \file{Platform file}{platform} from scratch by defining attributes such as
\config{markerName}, \config{markerNumber}, \config{comment}, \config{approxPosition},
\config{equipment}.

See also \program{GnssAntex2AntennaDefinition} and \program{GnssStationLog2Platform}.

\fig{!hb}{0.8}{fileFormatPlatform}{fig:platformCreate}{Platform for stations, LEOs, and GNSS satellites.}


\keepXColumns
\begin{tabularx}{\textwidth}{N T A}
\hline
Name & Type & Annotation\\
\hline
\hfuzz=500pt\includegraphics[width=1em]{element-mustset.pdf}~outputfilePlatform & \hfuzz=500pt filename & \hfuzz=500pt \\
\hfuzz=500pt\includegraphics[width=1em]{element-mustset.pdf}~markerName & \hfuzz=500pt string & \hfuzz=500pt \\
\hfuzz=500pt\includegraphics[width=1em]{element.pdf}~markerNumber & \hfuzz=500pt string & \hfuzz=500pt \\
\hfuzz=500pt\includegraphics[width=1em]{element.pdf}~comment & \hfuzz=500pt string & \hfuzz=500pt \\
\hfuzz=500pt\includegraphics[width=1em]{element.pdf}~approxPositionX & \hfuzz=500pt double & \hfuzz=500pt [m] in TRF\\
\hfuzz=500pt\includegraphics[width=1em]{element.pdf}~approxPositionY & \hfuzz=500pt double & \hfuzz=500pt [m] in TRF\\
\hfuzz=500pt\includegraphics[width=1em]{element.pdf}~approxPositionZ & \hfuzz=500pt double & \hfuzz=500pt [m] in TRF\\
\hfuzz=500pt\includegraphics[width=1em]{element-unbounded.pdf}~equipment & \hfuzz=500pt choice & \hfuzz=500pt \\
\hfuzz=500pt\includegraphics[width=1em]{connector.pdf}\includegraphics[width=1em]{element-mustset.pdf}~gnssAntenna & \hfuzz=500pt sequence & \hfuzz=500pt \\
\hfuzz=500pt\quad\includegraphics[width=1em]{connector.pdf}\includegraphics[width=1em]{element-mustset.pdf}~name & \hfuzz=500pt string & \hfuzz=500pt \\
\hfuzz=500pt\quad\includegraphics[width=1em]{connector.pdf}\includegraphics[width=1em]{element.pdf}~serial & \hfuzz=500pt string & \hfuzz=500pt \\
\hfuzz=500pt\quad\includegraphics[width=1em]{connector.pdf}\includegraphics[width=1em]{element.pdf}~radome & \hfuzz=500pt string & \hfuzz=500pt \\
\hfuzz=500pt\quad\includegraphics[width=1em]{connector.pdf}\includegraphics[width=1em]{element.pdf}~comment & \hfuzz=500pt string & \hfuzz=500pt \\
\hfuzz=500pt\quad\includegraphics[width=1em]{connector.pdf}\includegraphics[width=1em]{element.pdf}~timeStart & \hfuzz=500pt time & \hfuzz=500pt \\
\hfuzz=500pt\quad\includegraphics[width=1em]{connector.pdf}\includegraphics[width=1em]{element.pdf}~timeEnd & \hfuzz=500pt time & \hfuzz=500pt \\
\hfuzz=500pt\quad\includegraphics[width=1em]{connector.pdf}\includegraphics[width=1em]{element-mustset.pdf}~positionX & \hfuzz=500pt double & \hfuzz=500pt [m] ARP in north, east, up or vehicle system\\
\hfuzz=500pt\quad\includegraphics[width=1em]{connector.pdf}\includegraphics[width=1em]{element-mustset.pdf}~positionY & \hfuzz=500pt double & \hfuzz=500pt [m] ARP in north, east, up or vehicle system\\
\hfuzz=500pt\quad\includegraphics[width=1em]{connector.pdf}\includegraphics[width=1em]{element-mustset.pdf}~positionZ & \hfuzz=500pt double & \hfuzz=500pt [m] ARP in north, east, up or vehicle system\\
\hfuzz=500pt\quad\includegraphics[width=1em]{connector.pdf}\includegraphics[width=1em]{element.pdf}~rotationX & \hfuzz=500pt angle & \hfuzz=500pt [degree] from local/vehicle to left-handed antenna system\\
\hfuzz=500pt\quad\includegraphics[width=1em]{connector.pdf}\includegraphics[width=1em]{element.pdf}~rotationY & \hfuzz=500pt angle & \hfuzz=500pt [degree] from local/vehicle to left-handed antenna system\\
\hfuzz=500pt\quad\includegraphics[width=1em]{connector.pdf}\includegraphics[width=1em]{element.pdf}~rotationZ & \hfuzz=500pt angle & \hfuzz=500pt [degree] from local/vehicle to left-handed antenna system\\
\hfuzz=500pt\quad\includegraphics[width=1em]{connector.pdf}\includegraphics[width=1em]{element.pdf}~flipX & \hfuzz=500pt boolean & \hfuzz=500pt flip x-axis (after rotation)\\
\hfuzz=500pt\quad\includegraphics[width=1em]{connector.pdf}\includegraphics[width=1em]{element.pdf}~flipY & \hfuzz=500pt boolean & \hfuzz=500pt flip y-axis (after rotation)\\
\hfuzz=500pt\quad\includegraphics[width=1em]{connector.pdf}\includegraphics[width=1em]{element.pdf}~flipZ & \hfuzz=500pt boolean & \hfuzz=500pt flip z-axis (after rotation)\\
\hfuzz=500pt\includegraphics[width=1em]{connector.pdf}\includegraphics[width=1em]{element-mustset.pdf}~gnssReceiver & \hfuzz=500pt sequence & \hfuzz=500pt \\
\hfuzz=500pt\quad\includegraphics[width=1em]{connector.pdf}\includegraphics[width=1em]{element-mustset.pdf}~name & \hfuzz=500pt string & \hfuzz=500pt \\
\hfuzz=500pt\quad\includegraphics[width=1em]{connector.pdf}\includegraphics[width=1em]{element.pdf}~serial & \hfuzz=500pt string & \hfuzz=500pt \\
\hfuzz=500pt\quad\includegraphics[width=1em]{connector.pdf}\includegraphics[width=1em]{element.pdf}~version & \hfuzz=500pt string & \hfuzz=500pt \\
\hfuzz=500pt\quad\includegraphics[width=1em]{connector.pdf}\includegraphics[width=1em]{element.pdf}~comment & \hfuzz=500pt string & \hfuzz=500pt \\
\hfuzz=500pt\quad\includegraphics[width=1em]{connector.pdf}\includegraphics[width=1em]{element.pdf}~timeStart & \hfuzz=500pt time & \hfuzz=500pt \\
\hfuzz=500pt\quad\includegraphics[width=1em]{connector.pdf}\includegraphics[width=1em]{element.pdf}~timeEnd & \hfuzz=500pt time & \hfuzz=500pt \\
\hfuzz=500pt\includegraphics[width=1em]{connector.pdf}\includegraphics[width=1em]{element-mustset.pdf}~laserRetroReflector & \hfuzz=500pt sequence & \hfuzz=500pt \\
\hfuzz=500pt\quad\includegraphics[width=1em]{connector.pdf}\includegraphics[width=1em]{element-mustset.pdf}~name & \hfuzz=500pt string & \hfuzz=500pt e.g. GFZ, ITE, IPIE\\
\hfuzz=500pt\quad\includegraphics[width=1em]{connector.pdf}\includegraphics[width=1em]{element.pdf}~serial & \hfuzz=500pt string & \hfuzz=500pt \\
\hfuzz=500pt\quad\includegraphics[width=1em]{connector.pdf}\includegraphics[width=1em]{element.pdf}~comment & \hfuzz=500pt string & \hfuzz=500pt \\
\hfuzz=500pt\quad\includegraphics[width=1em]{connector.pdf}\includegraphics[width=1em]{element.pdf}~timeStart & \hfuzz=500pt time & \hfuzz=500pt \\
\hfuzz=500pt\quad\includegraphics[width=1em]{connector.pdf}\includegraphics[width=1em]{element.pdf}~timeEnd & \hfuzz=500pt time & \hfuzz=500pt \\
\hfuzz=500pt\quad\includegraphics[width=1em]{connector.pdf}\includegraphics[width=1em]{element-mustset.pdf}~positionX & \hfuzz=500pt double & \hfuzz=500pt [m] optial reference point RP in satellite system\\
\hfuzz=500pt\quad\includegraphics[width=1em]{connector.pdf}\includegraphics[width=1em]{element-mustset.pdf}~positionY & \hfuzz=500pt double & \hfuzz=500pt [m] optial reference point RP in satellite system\\
\hfuzz=500pt\quad\includegraphics[width=1em]{connector.pdf}\includegraphics[width=1em]{element-mustset.pdf}~positionZ & \hfuzz=500pt double & \hfuzz=500pt [m] optial reference point RP in satellite system\\
\hfuzz=500pt\quad\includegraphics[width=1em]{connector.pdf}\includegraphics[width=1em]{element.pdf}~rotationX & \hfuzz=500pt angle & \hfuzz=500pt [degree] from local/vehicle to LRR system\\
\hfuzz=500pt\quad\includegraphics[width=1em]{connector.pdf}\includegraphics[width=1em]{element.pdf}~rotationY & \hfuzz=500pt angle & \hfuzz=500pt [degree] from local/vehicle to LRR system\\
\hfuzz=500pt\quad\includegraphics[width=1em]{connector.pdf}\includegraphics[width=1em]{element.pdf}~rotationZ & \hfuzz=500pt angle & \hfuzz=500pt [degree] from local/vehicle to LRR system\\
\hfuzz=500pt\quad\includegraphics[width=1em]{connector.pdf}\includegraphics[width=1em]{element.pdf}~flipX & \hfuzz=500pt boolean & \hfuzz=500pt flip x-axis (after rotation)\\
\hfuzz=500pt\quad\includegraphics[width=1em]{connector.pdf}\includegraphics[width=1em]{element.pdf}~flipY & \hfuzz=500pt boolean & \hfuzz=500pt flip y-axis (after rotation)\\
\hfuzz=500pt\quad\includegraphics[width=1em]{connector.pdf}\includegraphics[width=1em]{element.pdf}~flipZ & \hfuzz=500pt boolean & \hfuzz=500pt flip z-axis (after rotation)\\
\hfuzz=500pt\quad\includegraphics[width=1em]{connector.pdf}\includegraphics[width=1em]{element.pdf}~range & \hfuzz=500pt double & \hfuzz=500pt [m] range bias (only without range matrix)\\
\hfuzz=500pt\quad\includegraphics[width=1em]{connector.pdf}\includegraphics[width=1em]{element.pdf}~inputfileRangeMatrix & \hfuzz=500pt filename & \hfuzz=500pt [m] (azimuth(0..360) x zenith(0..dZenit*rows)\\
\hfuzz=500pt\quad\includegraphics[width=1em]{connector.pdf}\includegraphics[width=1em]{element.pdf}~dZenit & \hfuzz=500pt angle & \hfuzz=500pt [degree] increment of range matrix\\
\hfuzz=500pt\includegraphics[width=1em]{connector.pdf}\includegraphics[width=1em]{element-mustset.pdf}~geodeticSatellite & \hfuzz=500pt sequence & \hfuzz=500pt e.g. LAGEOS\\
\hfuzz=500pt\quad\includegraphics[width=1em]{connector.pdf}\includegraphics[width=1em]{element-mustset.pdf}~name & \hfuzz=500pt string & \hfuzz=500pt \\
\hfuzz=500pt\quad\includegraphics[width=1em]{connector.pdf}\includegraphics[width=1em]{element.pdf}~serial & \hfuzz=500pt string & \hfuzz=500pt \\
\hfuzz=500pt\quad\includegraphics[width=1em]{connector.pdf}\includegraphics[width=1em]{element.pdf}~comment & \hfuzz=500pt string & \hfuzz=500pt \\
\hfuzz=500pt\quad\includegraphics[width=1em]{connector.pdf}\includegraphics[width=1em]{element.pdf}~timeStart & \hfuzz=500pt time & \hfuzz=500pt \\
\hfuzz=500pt\quad\includegraphics[width=1em]{connector.pdf}\includegraphics[width=1em]{element.pdf}~timeEnd & \hfuzz=500pt time & \hfuzz=500pt \\
\hfuzz=500pt\quad\includegraphics[width=1em]{connector.pdf}\includegraphics[width=1em]{element-mustset.pdf}~range & \hfuzz=500pt double & \hfuzz=500pt [m] standard center-of-mass correction\\
\hfuzz=500pt\includegraphics[width=1em]{connector.pdf}\includegraphics[width=1em]{element-mustset.pdf}~slrStation & \hfuzz=500pt sequence & \hfuzz=500pt \\
\hfuzz=500pt\quad\includegraphics[width=1em]{connector.pdf}\includegraphics[width=1em]{element-mustset.pdf}~name & \hfuzz=500pt string & \hfuzz=500pt CDP SOD 8-digit No.\\
\hfuzz=500pt\quad\includegraphics[width=1em]{connector.pdf}\includegraphics[width=1em]{element.pdf}~serial & \hfuzz=500pt string & \hfuzz=500pt IERS DOMES\\
\hfuzz=500pt\quad\includegraphics[width=1em]{connector.pdf}\includegraphics[width=1em]{element.pdf}~comment & \hfuzz=500pt string & \hfuzz=500pt \\
\hfuzz=500pt\quad\includegraphics[width=1em]{connector.pdf}\includegraphics[width=1em]{element.pdf}~timeStart & \hfuzz=500pt time & \hfuzz=500pt \\
\hfuzz=500pt\quad\includegraphics[width=1em]{connector.pdf}\includegraphics[width=1em]{element.pdf}~timeEnd & \hfuzz=500pt time & \hfuzz=500pt \\
\hfuzz=500pt\quad\includegraphics[width=1em]{connector.pdf}\includegraphics[width=1em]{element.pdf}~positionX & \hfuzz=500pt double & \hfuzz=500pt [m] eccentricity in north\\
\hfuzz=500pt\quad\includegraphics[width=1em]{connector.pdf}\includegraphics[width=1em]{element.pdf}~positionY & \hfuzz=500pt double & \hfuzz=500pt [m] eccentricity in east\\
\hfuzz=500pt\quad\includegraphics[width=1em]{connector.pdf}\includegraphics[width=1em]{element.pdf}~positionZ & \hfuzz=500pt double & \hfuzz=500pt [m] eccentricity in up\\
\hfuzz=500pt\includegraphics[width=1em]{connector.pdf}\includegraphics[width=1em]{element-mustset.pdf}~satelliteIdentifier & \hfuzz=500pt sequence & \hfuzz=500pt \\
\hfuzz=500pt\quad\includegraphics[width=1em]{connector.pdf}\includegraphics[width=1em]{element-mustset.pdf}~name & \hfuzz=500pt string & \hfuzz=500pt \\
\hfuzz=500pt\quad\includegraphics[width=1em]{connector.pdf}\includegraphics[width=1em]{element.pdf}~serial & \hfuzz=500pt string & \hfuzz=500pt \\
\hfuzz=500pt\quad\includegraphics[width=1em]{connector.pdf}\includegraphics[width=1em]{element.pdf}~cospar & \hfuzz=500pt string & \hfuzz=500pt Satellite COSPAR ID\\
\hfuzz=500pt\quad\includegraphics[width=1em]{connector.pdf}\includegraphics[width=1em]{element.pdf}~norad & \hfuzz=500pt string & \hfuzz=500pt Satellite Catalog (NORAD) Number\\
\hfuzz=500pt\quad\includegraphics[width=1em]{connector.pdf}\includegraphics[width=1em]{element.pdf}~sic & \hfuzz=500pt string & \hfuzz=500pt SIC Code\\
\hfuzz=500pt\quad\includegraphics[width=1em]{connector.pdf}\includegraphics[width=1em]{element.pdf}~sp3 & \hfuzz=500pt string & \hfuzz=500pt SP3\\
\hfuzz=500pt\quad\includegraphics[width=1em]{connector.pdf}\includegraphics[width=1em]{element.pdf}~comment & \hfuzz=500pt string & \hfuzz=500pt \\
\hfuzz=500pt\quad\includegraphics[width=1em]{connector.pdf}\includegraphics[width=1em]{element.pdf}~timeStart & \hfuzz=500pt time & \hfuzz=500pt \\
\hfuzz=500pt\quad\includegraphics[width=1em]{connector.pdf}\includegraphics[width=1em]{element.pdf}~timeEnd & \hfuzz=500pt time & \hfuzz=500pt \\
\hfuzz=500pt\includegraphics[width=1em]{connector.pdf}\includegraphics[width=1em]{element-mustset.pdf}~other & \hfuzz=500pt sequence & \hfuzz=500pt \\
\hfuzz=500pt\quad\includegraphics[width=1em]{connector.pdf}\includegraphics[width=1em]{element-mustset.pdf}~name & \hfuzz=500pt string & \hfuzz=500pt \\
\hfuzz=500pt\quad\includegraphics[width=1em]{connector.pdf}\includegraphics[width=1em]{element.pdf}~serial & \hfuzz=500pt string & \hfuzz=500pt \\
\hfuzz=500pt\quad\includegraphics[width=1em]{connector.pdf}\includegraphics[width=1em]{element.pdf}~comment & \hfuzz=500pt string & \hfuzz=500pt \\
\hfuzz=500pt\quad\includegraphics[width=1em]{connector.pdf}\includegraphics[width=1em]{element.pdf}~timeStart & \hfuzz=500pt time & \hfuzz=500pt \\
\hfuzz=500pt\quad\includegraphics[width=1em]{connector.pdf}\includegraphics[width=1em]{element.pdf}~timeEnd & \hfuzz=500pt time & \hfuzz=500pt \\
\hfuzz=500pt\quad\includegraphics[width=1em]{connector.pdf}\includegraphics[width=1em]{element.pdf}~positionX & \hfuzz=500pt double & \hfuzz=500pt [m] in north, east, up or vehicle system\\
\hfuzz=500pt\quad\includegraphics[width=1em]{connector.pdf}\includegraphics[width=1em]{element.pdf}~positionY & \hfuzz=500pt double & \hfuzz=500pt [m] in north, east, up or vehicle system\\
\hfuzz=500pt\quad\includegraphics[width=1em]{connector.pdf}\includegraphics[width=1em]{element.pdf}~positionZ & \hfuzz=500pt double & \hfuzz=500pt [m] in north, east, up or vehicle system\\
\hfuzz=500pt\includegraphics[width=1em]{element-unbounded.pdf}~referencePoint & \hfuzz=500pt sequence & \hfuzz=500pt e.g. center of mass in satellite frame\\
\hfuzz=500pt\includegraphics[width=1em]{connector.pdf}\includegraphics[width=1em]{element.pdf}~comment & \hfuzz=500pt string & \hfuzz=500pt \\
\hfuzz=500pt\includegraphics[width=1em]{connector.pdf}\includegraphics[width=1em]{element-mustset.pdf}~xStart & \hfuzz=500pt double & \hfuzz=500pt [m] in north, east, up or vehicle system\\
\hfuzz=500pt\includegraphics[width=1em]{connector.pdf}\includegraphics[width=1em]{element-mustset.pdf}~yStart & \hfuzz=500pt double & \hfuzz=500pt linear motion between start and end\\
\hfuzz=500pt\includegraphics[width=1em]{connector.pdf}\includegraphics[width=1em]{element-mustset.pdf}~zStart & \hfuzz=500pt double & \hfuzz=500pt \\
\hfuzz=500pt\includegraphics[width=1em]{connector.pdf}\includegraphics[width=1em]{element-mustset.pdf}~xEnd & \hfuzz=500pt double & \hfuzz=500pt [m] in north, east, up or vehicle system\\
\hfuzz=500pt\includegraphics[width=1em]{connector.pdf}\includegraphics[width=1em]{element-mustset.pdf}~yEnd & \hfuzz=500pt double & \hfuzz=500pt linear motion between start and end\\
\hfuzz=500pt\includegraphics[width=1em]{connector.pdf}\includegraphics[width=1em]{element-mustset.pdf}~zEnd & \hfuzz=500pt double & \hfuzz=500pt \\
\hfuzz=500pt\includegraphics[width=1em]{connector.pdf}\includegraphics[width=1em]{element.pdf}~timeStart & \hfuzz=500pt time & \hfuzz=500pt \\
\hfuzz=500pt\includegraphics[width=1em]{connector.pdf}\includegraphics[width=1em]{element.pdf}~timeEnd & \hfuzz=500pt time & \hfuzz=500pt \\
\hline
\end{tabularx}

\clearpage
%==================================
\subsection{PotentialCoefficients2BlockMeanTimeSplines}\label{PotentialCoefficients2BlockMeanTimeSplines}
This program is a simplified version of \program{Gravityfield2TimeSplines}.
It reads a series of potential coefficient files (\configFile{inputfilePotentialCoefficients}{potentialCoefficients})
and creates a time splines file with spline degree 0 (temporal block means) or degree 1 (linear splines).
The time intervals in which the potential coefficients are valid are defined between adjacent
points in time given by \config{splineTimeSeries}. Therefore one more point in time is needed
than the number of potential coefficient files for degree 0.

The coefficients can be filtered with \configClass{filter}{sphericalHarmonicsFilterType}.
If set the expansion is limited in the range between \config{minDegree} and \config{maxDegree} inclusivly.
The coefficients are related to the reference radius~\config{R} and the Earth gravitational constant \config{GM}.

This program is useful e.g. to combine monthly GRACE solutions to one file.


\keepXColumns
\begin{tabularx}{\textwidth}{N T A}
\hline
Name & Type & Annotation\\
\hline
\hfuzz=500pt\includegraphics[width=1em]{element-mustset.pdf}~outputfileTimeSplines & \hfuzz=500pt filename & \hfuzz=500pt \\
\hfuzz=500pt\includegraphics[width=1em]{element.pdf}~outputfileTimeSplinesCovariance & \hfuzz=500pt filename & \hfuzz=500pt only the variances are saved\\
\hfuzz=500pt\includegraphics[width=1em]{element-mustset-unbounded.pdf}~inputfilePotentialCoefficients & \hfuzz=500pt filename & \hfuzz=500pt \\
\hfuzz=500pt\includegraphics[width=1em]{element-unbounded.pdf}~filter & \hfuzz=500pt \hyperref[sphericalHarmonicsFilterType]{sphericalHarmonicsFilter} & \hfuzz=500pt \\
\hfuzz=500pt\includegraphics[width=1em]{element.pdf}~minDegree & \hfuzz=500pt uint & \hfuzz=500pt \\
\hfuzz=500pt\includegraphics[width=1em]{element.pdf}~maxDegree & \hfuzz=500pt uint & \hfuzz=500pt \\
\hfuzz=500pt\includegraphics[width=1em]{element.pdf}~GM & \hfuzz=500pt double & \hfuzz=500pt Geocentric gravitational constant\\
\hfuzz=500pt\includegraphics[width=1em]{element.pdf}~R & \hfuzz=500pt double & \hfuzz=500pt reference radius\\
\hfuzz=500pt\includegraphics[width=1em]{element.pdf}~removeMean & \hfuzz=500pt boolean & \hfuzz=500pt remove the temporal mean of the series before estimating the splines\\
\hfuzz=500pt\includegraphics[width=1em]{element.pdf}~interpolate & \hfuzz=500pt boolean & \hfuzz=500pt interpolate missing files\\
\hfuzz=500pt\includegraphics[width=1em]{element-mustset-unbounded.pdf}~splineTimeSeries & \hfuzz=500pt \hyperref[timeSeriesType]{timeSeries} & \hfuzz=500pt input files must be between points in time\\
\hfuzz=500pt\includegraphics[width=1em]{element.pdf}~splineDegree & \hfuzz=500pt uint & \hfuzz=500pt degree of splines\\
\hline
\end{tabularx}

\clearpage
%==================================
\subsection{PotentialCoefficients2DegreeAmplitudes}\label{PotentialCoefficients2DegreeAmplitudes}
This program computes degree amplitudes from
\file{potentialCoefficients files}{potentialCoefficients}
and saves them to a \file{matrix}{matrix} file.

The coefficients can be filtered with \configClass{filter}{sphericalHarmonicsFilterType} and converted
to different functionals with \configClass{kernel}{kernelType}. The gravity field can be evaluated at
different altitudes by specifying \config{evaluationRadius}. Polar regions can be excluded
by setting \config{polarGap}. If set the expansion is limited in the range between \config{minDegree}
and \config{maxDegree} inclusivly. The coefficients are related to the reference radius~\config{R}
and the Earth gravitational constant \config{GM}.

The \configFile{outputfileMatrix}{matrix} contains in the first 3 columns the degree, the degree amplitude, and
the formal errors. For each additional \configFile{inputfilePotentialCoefficients}{potentialCoefficients} three columns
are appended: the degree amplitude, the formal errors, and the difference to the first file.

For example the data columns for 4 \configFile{inputfilePotentialCoefficients}{potentialCoefficients} are
\begin{itemize}
\item degree=\verb|data0|
\item PotentialCoefficients0: signal=\verb|data1|, error=\verb|data2|,
\item PotentialCoefficients1: signal=\verb|data3|, error=\verb|data4|,  difference=\verb|data5|,
\item PotentialCoefficients2: signal=\verb|data6|, error=\verb|data7|,  difference=\verb|data8|,
\item PotentialCoefficients3: signal=\verb|data9|, error=\verb|data10|, difference=\verb|data11|.
\end{itemize}

See also \program{Gravityfield2DegreeAmplitudes}.


\keepXColumns
\begin{tabularx}{\textwidth}{N T A}
\hline
Name & Type & Annotation\\
\hline
\hfuzz=500pt\includegraphics[width=1em]{element-mustset.pdf}~outputfileMatrix & \hfuzz=500pt filename & \hfuzz=500pt matrix with degree, signal amplitude, formal error, differences\\
\hfuzz=500pt\includegraphics[width=1em]{element-mustset-unbounded.pdf}~inputfilePotentialCoefficients & \hfuzz=500pt filename & \hfuzz=500pt \\
\hfuzz=500pt\includegraphics[width=1em]{element-mustset.pdf}~kernel & \hfuzz=500pt \hyperref[kernelType]{kernel} & \hfuzz=500pt \\
\hfuzz=500pt\includegraphics[width=1em]{element-unbounded.pdf}~filter & \hfuzz=500pt \hyperref[sphericalHarmonicsFilterType]{sphericalHarmonicsFilter} & \hfuzz=500pt filter the coefficients\\
\hfuzz=500pt\includegraphics[width=1em]{element-mustset.pdf}~type & \hfuzz=500pt choice & \hfuzz=500pt type of variances\\
\hfuzz=500pt\includegraphics[width=1em]{connector.pdf}\includegraphics[width=1em]{element-mustset.pdf}~rms & \hfuzz=500pt  & \hfuzz=500pt degree amplitudes (square root of degree variances)\\
\hfuzz=500pt\includegraphics[width=1em]{connector.pdf}\includegraphics[width=1em]{element-mustset.pdf}~accumulation & \hfuzz=500pt  & \hfuzz=500pt cumulate variances over degrees\\
\hfuzz=500pt\includegraphics[width=1em]{element.pdf}~evaluationRadius & \hfuzz=500pt double & \hfuzz=500pt evaluate the gravity field at this radius (default: evaluate at surface\\
\hfuzz=500pt\includegraphics[width=1em]{element.pdf}~polarGap & \hfuzz=500pt angle & \hfuzz=500pt exclude polar regions (aperture angle in degrees)\\
\hfuzz=500pt\includegraphics[width=1em]{element.pdf}~minDegree & \hfuzz=500pt uint & \hfuzz=500pt \\
\hfuzz=500pt\includegraphics[width=1em]{element.pdf}~maxDegree & \hfuzz=500pt uint & \hfuzz=500pt \\
\hfuzz=500pt\includegraphics[width=1em]{element.pdf}~GM & \hfuzz=500pt double & \hfuzz=500pt Geocentric gravitational constant\\
\hfuzz=500pt\includegraphics[width=1em]{element.pdf}~R & \hfuzz=500pt double & \hfuzz=500pt reference radius\\
\hline
\end{tabularx}

\clearpage
%==================================
\subsection{RadialBasisSplines2KernelCoefficients}\label{RadialBasisSplines2KernelCoefficients}
This program calculates the coefficients $k_n$ of a \configClass{kernel:coefficients}{kernelType:coefficients} according to
\begin{equation}
  k_n = \frac{GM}{4\pi R}\frac{\sigma_n}{\sqrt{2n+1}}.
\end{equation}
from a given \configClass{gravityfield}{gravityfieldType},
with \config{R} and \config{GM} describing the reference radius and the geocentric constant, respectively.
The $\sigma_n$
stand for the gravity field accuracies (from degree \config{minDegree} to \config{maxDegree}), if they are given.
If no accuracies are provided, the $\sigma_n$
represent the square root of the degree variances of the gravity field.
If \config{maxDegree} exceeds the maximum degree given by \configClass{gravityfield}{gravityfieldType},
the higher degrees are complemented by Kaula's rule
The output of the coefficients is given in the file  \configFile{outputfileCoefficients}{matrix}.


\keepXColumns
\begin{tabularx}{\textwidth}{N T A}
\hline
Name & Type & Annotation\\
\hline
\hfuzz=500pt\includegraphics[width=1em]{element-mustset.pdf}~outputfileCoefficients & \hfuzz=500pt filename & \hfuzz=500pt \\
\hfuzz=500pt\includegraphics[width=1em]{element-unbounded.pdf}~gravityfield & \hfuzz=500pt \hyperref[gravityfieldType]{gravityfield} & \hfuzz=500pt use sigmas, if not given use signal (cnm,snm), if not given use kaulas rule\\
\hfuzz=500pt\includegraphics[width=1em]{element-mustset.pdf}~minDegree & \hfuzz=500pt uint & \hfuzz=500pt \\
\hfuzz=500pt\includegraphics[width=1em]{element.pdf}~maxDegree & \hfuzz=500pt uint & \hfuzz=500pt \\
\hfuzz=500pt\includegraphics[width=1em]{element.pdf}~GM & \hfuzz=500pt double & \hfuzz=500pt Geocentric gravitational constant\\
\hfuzz=500pt\includegraphics[width=1em]{element.pdf}~R & \hfuzz=500pt double & \hfuzz=500pt reference radius\\
\hfuzz=500pt\includegraphics[width=1em]{element.pdf}~kaulaPower & \hfuzz=500pt double & \hfuzz=500pt sigma = kaulaFactor/degree\textasciicircum{}kaulaPower\\
\hfuzz=500pt\includegraphics[width=1em]{element.pdf}~kaulaFactor & \hfuzz=500pt double & \hfuzz=500pt sigma = kaulaFactor/degree\textasciicircum{}kaulaPower\\
\hline
\end{tabularx}

\clearpage
%==================================
\subsection{SatelliteModelCreate}\label{SatelliteModelCreate}
This program creates a satellite macro model for the estimation of non-gravitational accelerations acting on a satellite.
Mandatory input values are the \config{satelliteName}, \config{mass}, \config{coefficientDrag} and information
about the satellite \config{surfaces}. For low Earth orbiting satellites, like GRACE for instance, a good guess
for the drag coefficient could be 2.3. Apart from that, it is latter on possible to estimate a more precise variable drag coefficient
(e.g. \configClass{miscAccelerations:atmosphericDrag}{miscAccelerationsType:atmosphericDrag}), which will override this initial guess.
Concerning the satellite surfaces an external file must be imported which must contain information about each single
 satellite plate in terms of plate \config{area}, the associated plate normal and re-radiation properties
(reflexion, diffusion and absorption) properties in the visible and IR part. Examplarily, a description of
the macro model for GRACE can be found under:
\url{https://podaac-tools.jpl.nasa.gov/drive/files/allData/grace/docs/ProdSpecDoc_v4.6.pdf}
Additionally, it is possible to add further information like antennaThrust, solar panel, temporal mass changes and
massInstrument using the modules option.


\keepXColumns
\begin{tabularx}{\textwidth}{N T A}
\hline
Name & Type & Annotation\\
\hline
\hfuzz=500pt\includegraphics[width=1em]{element-mustset.pdf}~outputfileSatelliteModel & \hfuzz=500pt filename & \hfuzz=500pt \\
\hfuzz=500pt\includegraphics[width=1em]{element-mustset-unbounded.pdf}~satellite & \hfuzz=500pt sequence & \hfuzz=500pt \\
\hfuzz=500pt\includegraphics[width=1em]{connector.pdf}\includegraphics[width=1em]{element-mustset.pdf}~satelliteName & \hfuzz=500pt string & \hfuzz=500pt \\
\hfuzz=500pt\includegraphics[width=1em]{connector.pdf}\includegraphics[width=1em]{element-mustset.pdf}~mass & \hfuzz=500pt double & \hfuzz=500pt \\
\hfuzz=500pt\includegraphics[width=1em]{connector.pdf}\includegraphics[width=1em]{element-mustset.pdf}~coefficientDrag & \hfuzz=500pt double & \hfuzz=500pt \\
\hfuzz=500pt\includegraphics[width=1em]{connector.pdf}\includegraphics[width=1em]{element-mustset.pdf}~surfaces & \hfuzz=500pt sequence & \hfuzz=500pt \\
\hfuzz=500pt\quad\includegraphics[width=1em]{connector.pdf}\includegraphics[width=1em]{element-mustset.pdf}~inputfile & \hfuzz=500pt filename & \hfuzz=500pt each line must contain one surface element\\
\hfuzz=500pt\quad\includegraphics[width=1em]{connector.pdf}\includegraphics[width=1em]{element-mustset.pdf}~type & \hfuzz=500pt expression & \hfuzz=500pt 0: plate, 1: sphere, 2: cylinder\\
\hfuzz=500pt\quad\includegraphics[width=1em]{connector.pdf}\includegraphics[width=1em]{element-mustset.pdf}~area & \hfuzz=500pt expression & \hfuzz=500pt [m\textasciicircum{}2]\\
\hfuzz=500pt\quad\includegraphics[width=1em]{connector.pdf}\includegraphics[width=1em]{element-mustset.pdf}~normalX & \hfuzz=500pt expression & \hfuzz=500pt \\
\hfuzz=500pt\quad\includegraphics[width=1em]{connector.pdf}\includegraphics[width=1em]{element-mustset.pdf}~normalY & \hfuzz=500pt expression & \hfuzz=500pt \\
\hfuzz=500pt\quad\includegraphics[width=1em]{connector.pdf}\includegraphics[width=1em]{element-mustset.pdf}~normalZ & \hfuzz=500pt expression & \hfuzz=500pt \\
\hfuzz=500pt\quad\includegraphics[width=1em]{connector.pdf}\includegraphics[width=1em]{element-mustset.pdf}~reflexionVisible & \hfuzz=500pt expression & \hfuzz=500pt \\
\hfuzz=500pt\quad\includegraphics[width=1em]{connector.pdf}\includegraphics[width=1em]{element-mustset.pdf}~diffusionVisible & \hfuzz=500pt expression & \hfuzz=500pt \\
\hfuzz=500pt\quad\includegraphics[width=1em]{connector.pdf}\includegraphics[width=1em]{element-mustset.pdf}~absorptionVisible & \hfuzz=500pt expression & \hfuzz=500pt \\
\hfuzz=500pt\quad\includegraphics[width=1em]{connector.pdf}\includegraphics[width=1em]{element-mustset.pdf}~reflexionInfrared & \hfuzz=500pt expression & \hfuzz=500pt \\
\hfuzz=500pt\quad\includegraphics[width=1em]{connector.pdf}\includegraphics[width=1em]{element-mustset.pdf}~diffusionInfrared & \hfuzz=500pt expression & \hfuzz=500pt \\
\hfuzz=500pt\quad\includegraphics[width=1em]{connector.pdf}\includegraphics[width=1em]{element-mustset.pdf}~absorptionInfrared & \hfuzz=500pt expression & \hfuzz=500pt \\
\hfuzz=500pt\quad\includegraphics[width=1em]{connector.pdf}\includegraphics[width=1em]{element-mustset.pdf}~specificHeatCapacity & \hfuzz=500pt expression & \hfuzz=500pt 0: no thermal radiation, -1: direct reemission [Ws/K/m\textasciicircum{}2]\\
\hfuzz=500pt\includegraphics[width=1em]{connector.pdf}\includegraphics[width=1em]{element-unbounded.pdf}~module & \hfuzz=500pt choice & \hfuzz=500pt \\
\hfuzz=500pt\quad\includegraphics[width=1em]{connector.pdf}\includegraphics[width=1em]{element-mustset.pdf}~antennaThrust & \hfuzz=500pt sequence & \hfuzz=500pt \\
\hfuzz=500pt\quad\quad\includegraphics[width=1em]{connector.pdf}\includegraphics[width=1em]{element.pdf}~thrustX & \hfuzz=500pt double & \hfuzz=500pt \\
\hfuzz=500pt\quad\quad\includegraphics[width=1em]{connector.pdf}\includegraphics[width=1em]{element.pdf}~thrustY & \hfuzz=500pt double & \hfuzz=500pt \\
\hfuzz=500pt\quad\quad\includegraphics[width=1em]{connector.pdf}\includegraphics[width=1em]{element.pdf}~thrustZ & \hfuzz=500pt double & \hfuzz=500pt \\
\hfuzz=500pt\quad\includegraphics[width=1em]{connector.pdf}\includegraphics[width=1em]{element-mustset.pdf}~solarPanel & \hfuzz=500pt sequence & \hfuzz=500pt \\
\hfuzz=500pt\quad\quad\includegraphics[width=1em]{connector.pdf}\includegraphics[width=1em]{element.pdf}~rotationAxisX & \hfuzz=500pt double & \hfuzz=500pt \\
\hfuzz=500pt\quad\quad\includegraphics[width=1em]{connector.pdf}\includegraphics[width=1em]{element.pdf}~rotationAxisY & \hfuzz=500pt double & \hfuzz=500pt \\
\hfuzz=500pt\quad\quad\includegraphics[width=1em]{connector.pdf}\includegraphics[width=1em]{element.pdf}~rotationAxisZ & \hfuzz=500pt double & \hfuzz=500pt \\
\hfuzz=500pt\quad\quad\includegraphics[width=1em]{connector.pdf}\includegraphics[width=1em]{element.pdf}~normalX & \hfuzz=500pt double & \hfuzz=500pt Direction to sun\\
\hfuzz=500pt\quad\quad\includegraphics[width=1em]{connector.pdf}\includegraphics[width=1em]{element.pdf}~normalY & \hfuzz=500pt double & \hfuzz=500pt Direction to sun\\
\hfuzz=500pt\quad\quad\includegraphics[width=1em]{connector.pdf}\includegraphics[width=1em]{element.pdf}~normalZ & \hfuzz=500pt double & \hfuzz=500pt Direction to sun\\
\hfuzz=500pt\quad\quad\includegraphics[width=1em]{connector.pdf}\includegraphics[width=1em]{element-mustset-unbounded.pdf}~indexSurface & \hfuzz=500pt uint & \hfuzz=500pt index of solar panel surfaces\\
\hfuzz=500pt\quad\includegraphics[width=1em]{connector.pdf}\includegraphics[width=1em]{element-mustset.pdf}~massChange & \hfuzz=500pt sequence & \hfuzz=500pt \\
\hfuzz=500pt\quad\quad\includegraphics[width=1em]{connector.pdf}\includegraphics[width=1em]{element-mustset-unbounded.pdf}~time & \hfuzz=500pt time & \hfuzz=500pt \\
\hfuzz=500pt\quad\quad\includegraphics[width=1em]{connector.pdf}\includegraphics[width=1em]{element-mustset-unbounded.pdf}~mass & \hfuzz=500pt double & \hfuzz=500pt \\
\hfuzz=500pt\quad\includegraphics[width=1em]{connector.pdf}\includegraphics[width=1em]{element-mustset.pdf}~massInstrument & \hfuzz=500pt sequence & \hfuzz=500pt \\
\hfuzz=500pt\quad\quad\includegraphics[width=1em]{connector.pdf}\includegraphics[width=1em]{element-mustset-unbounded.pdf}~inputfileInstrument & \hfuzz=500pt filename & \hfuzz=500pt \\
\hline
\end{tabularx}

\clearpage
%==================================
\subsection{SynthesisSphericalHarmonicsMatrix}\label{SynthesisSphericalHarmonicsMatrix}
This program builds a linear operator matrix for spherical harmonic analysis or synthesis based on
the points defined in \configClass{grid}{gridType}. Depending on the chosen
\config{type} (synthesis, quadrature, or leastSquares), the resulting matrix can be used to:
\begin{itemize}
  \item \textbf{synthesis}: Map spherical harmonic coefficients to values on a grid,
  \item \textbf{quadrature}: Integrate grid-based functionals into spherical harmonic coefficients by
        a simple quadrature formula,
  \item \textbf{leastSquares}: Estimate coefficients from grid data via a least squares approach.
\end{itemize}

he spherical harmonic degree range is constrained by
\config{minDegree} and \config{maxDegree}, and the ordering of the coefficients is given by
\configClass{numbering}{sphericalHarmonicsNumberingType}. The reference gravitational
constant is \config{GM}, and the reference radius is \config{R}.

The computed matrix is written to \configFile{outputfileMatrix}{matrix} with dimensions
(number of grid points)~$\times$~(number of spherical harmonic coefficients). For
\config{type} = \emph{leastSquares}, the program applies a QR-based pseudo-inverse so that the
output matrix can directly form the normal-equation building blocks for a blockwise
least-squares solution in spherical harmonic space.

See also \program{Gravityfield2GriddedData}, \program{GriddedData2PotentialCoefficients},
\program{Gravityfield2SphericalHarmonicsVector}, and \program{MatrixCalculate} for additional
tools to convert between grids and spherical harmonics.


\keepXColumns
\begin{tabularx}{\textwidth}{N T A}
\hline
Name & Type & Annotation\\
\hline
\hfuzz=500pt\includegraphics[width=1em]{element.pdf}~outputfileMatrix & \hfuzz=500pt filename & \hfuzz=500pt \\
\hfuzz=500pt\includegraphics[width=1em]{element-mustset-unbounded.pdf}~grid & \hfuzz=500pt \hyperref[gridType]{grid} & \hfuzz=500pt \\
\hfuzz=500pt\includegraphics[width=1em]{element-mustset.pdf}~kernel & \hfuzz=500pt \hyperref[kernelType]{kernel} & \hfuzz=500pt \\
\hfuzz=500pt\includegraphics[width=1em]{element.pdf}~minDegree & \hfuzz=500pt uint & \hfuzz=500pt \\
\hfuzz=500pt\includegraphics[width=1em]{element-mustset.pdf}~maxDegree & \hfuzz=500pt uint & \hfuzz=500pt \\
\hfuzz=500pt\includegraphics[width=1em]{element.pdf}~GM & \hfuzz=500pt double & \hfuzz=500pt Geocentric gravitational constant\\
\hfuzz=500pt\includegraphics[width=1em]{element.pdf}~R & \hfuzz=500pt double & \hfuzz=500pt reference radius\\
\hfuzz=500pt\includegraphics[width=1em]{element-mustset.pdf}~numbering & \hfuzz=500pt \hyperref[sphericalHarmonicsNumberingType]{sphericalHarmonicsNumbering} & \hfuzz=500pt numbering scheme of sh coefficients\\
\hfuzz=500pt\includegraphics[width=1em]{element-mustset.pdf}~type & \hfuzz=500pt choice & \hfuzz=500pt \\
\hfuzz=500pt\includegraphics[width=1em]{connector.pdf}\includegraphics[width=1em]{element-mustset.pdf}~synthesis & \hfuzz=500pt  & \hfuzz=500pt synthesize spherical harmonics on a grid\\
\hfuzz=500pt\includegraphics[width=1em]{connector.pdf}\includegraphics[width=1em]{element-mustset.pdf}~quadrature & \hfuzz=500pt  & \hfuzz=500pt calculate spherical harmonics from grid\\
\hfuzz=500pt\includegraphics[width=1em]{connector.pdf}\includegraphics[width=1em]{element-mustset.pdf}~leastSquares & \hfuzz=500pt  & \hfuzz=500pt estimated spherical harmonics from grid\\
\hline
\end{tabularx}

This program is \reference{parallelized}{general.parallelization}.
\clearpage
%==================================
\subsection{TemporalRepresentation2TimeSeries}\label{TemporalRepresentation2TimeSeries}
This program computes the design matrix of temporal representation at a given time series.
The output matrix contains the time steps in MJD in the first column, the other columns contain the design matrix.
The intention of this program is to visualize the parametrization together with \program{PlotGraph}.


\keepXColumns
\begin{tabularx}{\textwidth}{N T A}
\hline
Name & Type & Annotation\\
\hline
\hfuzz=500pt\includegraphics[width=1em]{element-mustset.pdf}~outputfileMatrix & \hfuzz=500pt filename & \hfuzz=500pt Time (MJD) in first column, design matrix follows\\
\hfuzz=500pt\includegraphics[width=1em]{element-mustset-unbounded.pdf}~timeSeries & \hfuzz=500pt \hyperref[timeSeriesType]{timeSeries} & \hfuzz=500pt \\
\hfuzz=500pt\includegraphics[width=1em]{element-mustset-unbounded.pdf}~temporal & \hfuzz=500pt \hyperref[parametrizationTemporalType]{parametrizationTemporal} & \hfuzz=500pt \\
\hline
\end{tabularx}

\clearpage
%==================================
\subsection{ThermosphericState2GriddedData}\label{ThermosphericState2GriddedData}
This program converts the output (neutral mass density,temperature) of an empirical thermosphere model (e.g. JB2008) on a given \configClass{grid}{gridType}.
Additionally, also the thermospheric winds estimated by using the horizontal wind model HWM 2014 can be assessed.
The time for the evaluation can be specified in \config{time}. The values will be saved together with points expressed as ellipsoidal coordinates
(longitude, latitude, height) based on a reference ellipsoid with parameters \config{R} and \config{inverseFlattening}.
\fig{!hb}{1.0}{thermosphericState2GriddedData}{fig:thermosphericState2GriddedData}{JB2008 model in 300 km height at 2003-07-01 12:00.}


\keepXColumns
\begin{tabularx}{\textwidth}{N T A}
\hline
Name & Type & Annotation\\
\hline
\hfuzz=500pt\includegraphics[width=1em]{element-mustset.pdf}~outputfileGriddedData & \hfuzz=500pt filename & \hfuzz=500pt density [kg/m**3], temperature [K], wind (x, y, z) [m/s**2]\\
\hfuzz=500pt\includegraphics[width=1em]{element-mustset.pdf}~thermosphere & \hfuzz=500pt \hyperref[thermosphereType]{thermosphere} & \hfuzz=500pt \\
\hfuzz=500pt\includegraphics[width=1em]{element-mustset-unbounded.pdf}~grid & \hfuzz=500pt \hyperref[gridType]{grid} & \hfuzz=500pt \\
\hfuzz=500pt\includegraphics[width=1em]{element-mustset.pdf}~time & \hfuzz=500pt time & \hfuzz=500pt \\
\hfuzz=500pt\includegraphics[width=1em]{element.pdf}~localReferenceFrame & \hfuzz=500pt boolean & \hfuzz=500pt wind in local north, east, up, otherwise global terrestrial\\
\hfuzz=500pt\includegraphics[width=1em]{element.pdf}~R & \hfuzz=500pt double & \hfuzz=500pt reference radius for ellipsoidal coordinates on output\\
\hfuzz=500pt\includegraphics[width=1em]{element.pdf}~inverseFlattening & \hfuzz=500pt double & \hfuzz=500pt reference flattening for ellipsoidal coordinates on output, 0: spherical coordinates\\
\hline
\end{tabularx}

This program is \reference{parallelized}{general.parallelization}.
\clearpage
%==================================
\subsection{TimeSeries2PotentialCoefficients}\label{TimeSeries2PotentialCoefficients}
Interpret the data columns of \configFile{inputfileTimeSeries}{instrument}
as potential coefficients. The sequence of coefficients is given by
\configClass{numbering}{sphericalHarmonicsNumberingType} starting from data column \config{startDataFields}.

For each epoch a \configFile{outputfilesPotentialCoefficients}{potentialCoefficients}
is written where the \config{variableLoopTime} and \config{variableLoopIndex} are expanded for
each point of the given time series to create the file name for this epoch,
see \reference{text parser}{general.parser:text}.

See also \program{Gravityfield2PotentialCoefficientsTimeSeries}.


\keepXColumns
\begin{tabularx}{\textwidth}{N T A}
\hline
Name & Type & Annotation\\
\hline
\hfuzz=500pt\includegraphics[width=1em]{element-mustset.pdf}~outputfilesPotentialCoefficients & \hfuzz=500pt filename & \hfuzz=500pt for each epoch\\
\hfuzz=500pt\includegraphics[width=1em]{element.pdf}~variableLoopTime & \hfuzz=500pt string & \hfuzz=500pt variable with time of each epoch\\
\hfuzz=500pt\includegraphics[width=1em]{element.pdf}~variableLoopIndex & \hfuzz=500pt string & \hfuzz=500pt variable with index of current epoch (starts with zero)\\
\hfuzz=500pt\includegraphics[width=1em]{element.pdf}~variableLoopCount & \hfuzz=500pt string & \hfuzz=500pt variable with total number of epochs\\
\hfuzz=500pt\includegraphics[width=1em]{element-mustset.pdf}~inputfileTimeSeries & \hfuzz=500pt filename & \hfuzz=500pt each epoch: multiple data for points (MISCVALUES)\\
\hfuzz=500pt\includegraphics[width=1em]{element.pdf}~startDataFields & \hfuzz=500pt uint & \hfuzz=500pt first data column\\
\hfuzz=500pt\includegraphics[width=1em]{element-mustset.pdf}~minDegree & \hfuzz=500pt uint & \hfuzz=500pt minimal degree\\
\hfuzz=500pt\includegraphics[width=1em]{element-mustset.pdf}~maxDegree & \hfuzz=500pt uint & \hfuzz=500pt maximal degree\\
\hfuzz=500pt\includegraphics[width=1em]{element.pdf}~GM & \hfuzz=500pt double & \hfuzz=500pt Geocentric gravitational constant\\
\hfuzz=500pt\includegraphics[width=1em]{element.pdf}~R & \hfuzz=500pt double & \hfuzz=500pt reference radius\\
\hfuzz=500pt\includegraphics[width=1em]{element-mustset.pdf}~numbering & \hfuzz=500pt \hyperref[sphericalHarmonicsNumberingType]{sphericalHarmonicsNumbering} & \hfuzz=500pt numbering scheme\\
\hline
\end{tabularx}

\clearpage
%==================================
\subsection{TimeSeriesCreate}\label{TimeSeriesCreate}
This program generates an \file{instrument file}{instrument},
containing a time series.


\keepXColumns
\begin{tabularx}{\textwidth}{N T A}
\hline
Name & Type & Annotation\\
\hline
\hfuzz=500pt\includegraphics[width=1em]{element-mustset.pdf}~outputfileTimeSeries & \hfuzz=500pt filename & \hfuzz=500pt instrument file\\
\hfuzz=500pt\includegraphics[width=1em]{element-mustset-unbounded.pdf}~timeSeries & \hfuzz=500pt \hyperref[timeSeriesType]{timeSeries} & \hfuzz=500pt time series to be created\\
\hfuzz=500pt\includegraphics[width=1em]{element-unbounded.pdf}~data & \hfuzz=500pt expression & \hfuzz=500pt expression of output columns, extra 'epoch' variable\\
\hline
\end{tabularx}

\clearpage
%==================================
\subsection{Variational2OrbitAndStarCamera}\label{Variational2OrbitAndStarCamera}
Extracts the reference \configFile{outputfileOrbit}{instrument}, \configFile{outputfileStarCamera}{instrument},
and \configFile{outputfileEarthRotation}{instrument} from \configFile{inputfileVariational}{variationalEquation}.


\keepXColumns
\begin{tabularx}{\textwidth}{N T A}
\hline
Name & Type & Annotation\\
\hline
\hfuzz=500pt\includegraphics[width=1em]{element.pdf}~outputfileOrbit & \hfuzz=500pt filename & \hfuzz=500pt output orbit (instrument) file\\
\hfuzz=500pt\includegraphics[width=1em]{element.pdf}~outputfileStarCamera & \hfuzz=500pt filename & \hfuzz=500pt output satellite attidude as star camera (instrument) file\\
\hfuzz=500pt\includegraphics[width=1em]{element.pdf}~outputfileEarthRotation & \hfuzz=500pt filename & \hfuzz=500pt output Earth rotation as star camera (instrument) file\\
\hfuzz=500pt\includegraphics[width=1em]{element-mustset.pdf}~inputfileVariational & \hfuzz=500pt filename & \hfuzz=500pt input variational file\\
\hline
\end{tabularx}

\clearpage
%==================================
\section{Programs: NormalEquation}
\subsection{NormalsAccumulate}\label{NormalsAccumulate}
This program accumulates normal equations and writes the total combined system to
\configFile{outputfileNormalequation}{normalEquation}.
The \configFile{inputfileNormalEquation}{normalEquation}s must have all the same size and the same block structure.
This program is the simplified and fast version of the more general program \program{NormalsBuild}.


\keepXColumns
\begin{tabularx}{\textwidth}{N T A}
\hline
Name & Type & Annotation\\
\hline
\hfuzz=500pt\includegraphics[width=1em]{element-mustset.pdf}~outputfileNormalEquation & \hfuzz=500pt filename & \hfuzz=500pt \\
\hfuzz=500pt\includegraphics[width=1em]{element-mustset-unbounded.pdf}~inputfileNormalEquation & \hfuzz=500pt filename & \hfuzz=500pt \\
\hline
\end{tabularx}

\clearpage
%==================================
\subsection{NormalsBuild}\label{NormalsBuild}
This program accumulates \configClass{normalEquation}{normalEquationType}s and
writes the total combined system to \configFile{outputfileNormalequation}{normalEquation}.
For a detailed description of the used algorithm see \configClass{normalEquation}{normalEquationType}.
Large normal equation systems can be divided into blocks with \config{normalsBlockSize}.

A simplifed and fast version of this program is \program{NormalsAccumulate}.
To solve the system of normal equations use \program{NormalsSolverVCE}.


\keepXColumns
\begin{tabularx}{\textwidth}{N T A}
\hline
Name & Type & Annotation\\
\hline
\hfuzz=500pt\includegraphics[width=1em]{element-mustset.pdf}~outputfileNormalEquation & \hfuzz=500pt filename & \hfuzz=500pt \\
\hfuzz=500pt\includegraphics[width=1em]{element-mustset-unbounded.pdf}~normalEquation & \hfuzz=500pt \hyperref[normalEquationType]{normalEquation} & \hfuzz=500pt \\
\hfuzz=500pt\includegraphics[width=1em]{element.pdf}~normalsBlockSize & \hfuzz=500pt uint & \hfuzz=500pt block size for distributing the normal equations, 0: one block\\
\hline
\end{tabularx}

This program is \reference{parallelized}{general.parallelization}.
\clearpage
%==================================
\subsection{NormalsBuildShortTimeStaticLongTime}\label{NormalsBuildShortTimeStaticLongTime}
This program sets up normal equations based on \configClass{observation}{observationType}.
Additionally short time and long time variations can be parametrized based on the static parameters
in \configClass{observation}{observationType} in an efficient way. The observation equations
are divided into time intervals $i \in \{1, ..., N\}$ (e.g. daily) as defined in
\configFile{inputfileArcList}{arcList}.

With \config{estimateLongTimeVariations} additional temporal variations can be co-estimated
for a subset of the parameters selected by \configClass{parameterSelection}{parameterSelectorType}.
These parameters might be spherical harmonic coefficients with a limited maximum degree.
The temporal variations are represented by base functions $\Phi_k(t_i)$ (e.g. trend and annual oscillation)
given by \configClass{parametrizationTemporal}{parametrizationTemporalType}.
The temporal base functions are evaluated at the mid time~$t_i$ of each interval~$i$, multiplicated
with the design matrix $\M A_i$ of the selected parameters, and the design matrix is extended
accordingly.

\fig{!hb}{0.8}{normalsBuildShortTimeStaticLongTime}{fig:normalsBuildShortTimeStaticLongTime}{Schema of the extended design matrix.}

With \config{estimateShortTimeVariations} short time variations of the gravity field can be co-estimated.
Their purpose is to mitigate temporal aliasing.
The short time parameters selected by \configClass{parameterSelection}{parameterSelectorType}
(e.g. daily constant or linear splines every 6 hour) are constrained by an
\configClass{autoregressiveModelSequence}{autoregressiveModelSequenceType}. If only a static parameter
set is selected the coressponding part of the design matrix is copied and modeled as a constant value
per interval in \configFile{inputfileArcList}{arcList} additionally so the corresponding temporal factor can be expressed as
\begin{equation}
  \Phi_i(t)  =
  \begin{cases}
    1 &\text{if} \hspace{5pt} t \in [t_i, t_{i+1}) \\
    0 & \text{otherwise}
  \end{cases}.
\end{equation}

Before writing the normal equations to \configFile{outputfileNormalEquation}{normalEquation}
short time gravity and satellite specific parameters can be eliminated with \config{eliminateParameter}.

Example: For the computation of the mean gravity field ITSG-Grace2018s with additional trend and annual signal
the normal equations are computed month by month and accumulated afterwards (see \program{NormalsAccumulate}).
The observations were divided into daily intervals with \configFile{inputfileArcList}{arcList}.
The static gravity field has been parametrized as spherical harmonics
up to degree $n=200$ in \configClass{observation:parametrizationGravity}{parametrizationGravityType}.
The trend and annual signals defined by
\configClass{estimateLongTimeVariations:parametrizationTemporal}{parametrizationTemporalType}
were estimated for selected parameters up to degree $n=120$.
To mitigate temporal aliasing daily gravity fields up to degree $n=40$ were setup and constrained
with an \configClass{autoregressiveModelSequence}{autoregressiveModelSequenceType} up to order three.

A detailed description of the approach is given in:
Kvas, A., Mayer-Gürr, T. GRACE gravity field recovery with background model uncertainties.
J Geod 93, 2543–2552 (2019). \url{https://doi.org/10.1007/s00190-019-01314-1}.


\keepXColumns
\begin{tabularx}{\textwidth}{N T A}
\hline
Name & Type & Annotation\\
\hline
\hfuzz=500pt\includegraphics[width=1em]{element-mustset.pdf}~outputfileNormalEquation & \hfuzz=500pt filename & \hfuzz=500pt outputfile for normal equations\\
\hfuzz=500pt\includegraphics[width=1em]{element-mustset.pdf}~observation & \hfuzz=500pt \hyperref[observationType]{observation} & \hfuzz=500pt \\
\hfuzz=500pt\includegraphics[width=1em]{element.pdf}~estimateShortTimeVariations & \hfuzz=500pt sequence & \hfuzz=500pt co-estimate short time gravity field variations\\
\hfuzz=500pt\includegraphics[width=1em]{connector.pdf}\includegraphics[width=1em]{element-mustset.pdf}~autoregressiveModelSequence & \hfuzz=500pt \hyperref[autoregressiveModelSequenceType]{autoregressiveModelSequence} & \hfuzz=500pt AR model sequence for constraining short time gravity variations\\
\hfuzz=500pt\includegraphics[width=1em]{connector.pdf}\includegraphics[width=1em]{element-mustset-unbounded.pdf}~parameterSelection & \hfuzz=500pt \hyperref[parameterSelectorType]{parameterSelector} & \hfuzz=500pt parameters describing the short time gravity field\\
\hfuzz=500pt\includegraphics[width=1em]{element.pdf}~estimateLongTimeVariations & \hfuzz=500pt sequence & \hfuzz=500pt co-estimate long time gravity field variations\\
\hfuzz=500pt\includegraphics[width=1em]{connector.pdf}\includegraphics[width=1em]{element-mustset-unbounded.pdf}~parametrizationTemporal & \hfuzz=500pt \hyperref[parametrizationTemporalType]{parametrizationTemporal} & \hfuzz=500pt parametrization of time variations (trend, annual, ...)\\
\hfuzz=500pt\includegraphics[width=1em]{connector.pdf}\includegraphics[width=1em]{element-mustset-unbounded.pdf}~parameterSelection & \hfuzz=500pt \hyperref[parameterSelectorType]{parameterSelector} & \hfuzz=500pt parameters describing the long time gravity field\\
\hfuzz=500pt\includegraphics[width=1em]{element-mustset.pdf}~inputfileArcList & \hfuzz=500pt filename & \hfuzz=500pt list to correspond points of time to arc numbers\\
\hfuzz=500pt\includegraphics[width=1em]{element.pdf}~defaultBlockSize & \hfuzz=500pt uint & \hfuzz=500pt block size for distributing the normal equations, 0: one block\\
\hfuzz=500pt\includegraphics[width=1em]{element.pdf}~eliminateParameter & \hfuzz=500pt boolean & \hfuzz=500pt eliminate short time and state parameter\\
\hline
\end{tabularx}

This program is \reference{parallelized}{general.parallelization}.
\clearpage
%==================================
\subsection{NormalsCreate}\label{NormalsCreate}
Create \file{normal equations}{normalEquation}
from calculated matrices (\configClass{matrixGenerator}{matrixGeneratorType}).

The \configFile{inputfileParameterNames}{parameterName} can be created with \program{ParameterNamesCreate}.

The \configClass{normalMatrix}{matrixGeneratorType} must be symmetric.
The \configClass{rightHandSide}{matrixGeneratorType} must have the same number of rows
and can contain multiple columns for multiple solutions.

The Vector $\M l^T\M P\M l$ is the quadratic sum of observations for each column of the right hand side.
It is used to determine the aposteriori accuracy
\begin{equation}
\hat{\sigma}^2 = \frac{\hat{\M e}^T\M P\hat{\M e}}{n-m} = \frac{\M l^T\M P\M l - \M n^T\hat{\M x}}{n-m}.
\end{equation}
If the vector is not given, it is automatically determined by assuming $\hat{\sigma}^2=1$.

The number of observations~$n$ is given by the expression \config{observationCount}.
The variable \verb|observationCount| can be used, if it is set by a normal equation file
\configFile{inputfileNormalEquationObsCount}{normalEquation}.


\keepXColumns
\begin{tabularx}{\textwidth}{N T A}
\hline
Name & Type & Annotation\\
\hline
\hfuzz=500pt\includegraphics[width=1em]{element-mustset.pdf}~outputfileNormalEquation & \hfuzz=500pt filename & \hfuzz=500pt \\
\hfuzz=500pt\includegraphics[width=1em]{element.pdf}~inputfileParameterNames & \hfuzz=500pt filename & \hfuzz=500pt \\
\hfuzz=500pt\includegraphics[width=1em]{element-mustset-unbounded.pdf}~normalMatrix & \hfuzz=500pt \hyperref[matrixGeneratorType]{matrixGenerator} & \hfuzz=500pt \\
\hfuzz=500pt\includegraphics[width=1em]{element-mustset-unbounded.pdf}~rightHandSide & \hfuzz=500pt \hyperref[matrixGeneratorType]{matrixGenerator} & \hfuzz=500pt \\
\hfuzz=500pt\includegraphics[width=1em]{element-unbounded.pdf}~lPl & \hfuzz=500pt \hyperref[matrixGeneratorType]{matrixGenerator} & \hfuzz=500pt vector with size of rhs columns\\
\hfuzz=500pt\includegraphics[width=1em]{element.pdf}~inputfileNormalEquationObsCount & \hfuzz=500pt filename & \hfuzz=500pt sets the variable observationCount\\
\hfuzz=500pt\includegraphics[width=1em]{element-mustset.pdf}~observationCount & \hfuzz=500pt expression & \hfuzz=500pt (variables: rows, columns (rhs), observationCount)\\
\hline
\end{tabularx}

\clearpage
%==================================
\subsection{NormalsEliminate}\label{NormalsEliminate}
This program eliminates parameters from a system of \configFile{inputfileNormalEquation}{normalEquation}s.
To just remove (cutting out) parameters use \program{NormalsReorder}.

The \configClass{remainingParameters}{parameterSelectorType} allows the selection
of parameters that will remain, all others will be eliminated. The order of remaining parameters
can be modified via the parameter selection. Block size of the output normal matrix can be adjusted with
\config{outBlockSize}. If it is set to zero, the \configFile{outputfileNormalEquation}{normalEquation}
is written to a single block file.

For example the normal equations are divided into two groups of
parameters $\hat{\M x}_1$ and $\hat{\M x}_2$ according to
\begin{equation}
\begin{pmatrix}
  \M N_{11} & \M N_{12} \\
  \M N_{21} & \M N_{22}
\end{pmatrix}
\begin{pmatrix} \hat{\M x}_1 \\ \hat{\M x}_2 \end{pmatrix}
=
\begin{pmatrix}
  \M n_1 \\
  \M n_2
\end{pmatrix}.
\end{equation}
and $\hat{\M x}_2$ shall be eliminated, the reduced system of normal equations is given by
\begin{equation}
\bar{\M N}\hat{\M x} = \bar{\M n}
\qquad\text{with}\qquad
\bar{\M N}=\M N_{11}-\M N_{12}\M N_{22}^{-1}\M N_{12}^T
\qquad\text{and}\qquad\bar{\M n} =  \M n_1 - \M N_{12}\M N_{22}^{-1}\M n_2.
\end{equation}

See also \program{NormalsReorder}.


\keepXColumns
\begin{tabularx}{\textwidth}{N T A}
\hline
Name & Type & Annotation\\
\hline
\hfuzz=500pt\includegraphics[width=1em]{element-mustset.pdf}~outputfileNormalEquation & \hfuzz=500pt filename & \hfuzz=500pt \\
\hfuzz=500pt\includegraphics[width=1em]{element-mustset.pdf}~inputfileNormalEquation & \hfuzz=500pt filename & \hfuzz=500pt \\
\hfuzz=500pt\includegraphics[width=1em]{element-mustset-unbounded.pdf}~remainingParameters & \hfuzz=500pt \hyperref[parameterSelectorType]{parameterSelector} & \hfuzz=500pt parameter order/selection of output normal equations\\
\hfuzz=500pt\includegraphics[width=1em]{element.pdf}~outBlockSize & \hfuzz=500pt uint & \hfuzz=500pt block size for distributing the normal equations, 0: one block\\
\hline
\end{tabularx}

This program is \reference{parallelized}{general.parallelization}.
\clearpage
%==================================
\subsection{NormalsMultiplyAdd}\label{NormalsMultiplyAdd}
This program modifies \configFile{inputfileNormalEquation}{normalEquation} in a way
that $\bar{\M x}$ is estimated instead of $\M x$.
\begin{equation}
 \bar{\M x} := \M x + \alpha\, \M x_0,
\end{equation}
where $\M x_0$ is \configFile{inputfileParameter}{matrix} and $\alpha$ is \config{factor}.
This can be used to re-add reduced reference fields before a combined estimation
at normal equation level.
Therefore the right hand side of the normal equations is modified by
\begin{equation}
 \bar{\M n} := \M n + \alpha\,\M N\M x_0,
\end{equation}
and the quadratic sum of observations by
\begin{equation}
 \bar{\M l^T\M P\M l} := \M l^T\M P\M l + \alpha^2\,\M x_0^T\M N\M x_0 + 2\alpha\,\M x_0^T\M n
\end{equation}

As the normal matrix itself is not modified, rewriting of the matrix can be disabled by setting
\config{writeNormalMatrix} to false.


\keepXColumns
\begin{tabularx}{\textwidth}{N T A}
\hline
Name & Type & Annotation\\
\hline
\hfuzz=500pt\includegraphics[width=1em]{element-mustset.pdf}~outputfileNormalEquation & \hfuzz=500pt filename & \hfuzz=500pt \\
\hfuzz=500pt\includegraphics[width=1em]{element-mustset.pdf}~inputfileNormalEquation & \hfuzz=500pt filename & \hfuzz=500pt \\
\hfuzz=500pt\includegraphics[width=1em]{element-mustset.pdf}~inputfileParameter & \hfuzz=500pt filename & \hfuzz=500pt x\\
\hfuzz=500pt\includegraphics[width=1em]{element.pdf}~factor & \hfuzz=500pt double & \hfuzz=500pt alpha\\
\hfuzz=500pt\includegraphics[width=1em]{element.pdf}~writeNormalMatrix & \hfuzz=500pt boolean & \hfuzz=500pt write full coefficient matrix, right hand sides and info files\\
\hline
\end{tabularx}

This program is \reference{parallelized}{general.parallelization}.
\clearpage
%==================================
\subsection{NormalsRegularizationBorders}\label{NormalsRegularizationBorders}
This program sets up two regularization matrices for two different regional areas.
For a given set of points defined by \configClass{grid}{gridType} it is evaluated, whether each point
(corresponding to an unknown parameter of a respective parameterization by space localizing basis functions)
is inside or outside a certain area given by \configClass{border}{borderType}.
Each regularization matrix is a diagonal matrix, one of them features a one if the
point is inside, and a zero if the point lies outside the area. The other matrix features
a zero if the point is inside, and a one if the point lies outside the area
This results in two regularization matrices with
\begin{equation}
\M R_1+\M R_2=\M I.
\end{equation}
The two matrices are provided as vectors of the diagonal
in the output files \configFile{outputfileOutside}{matrix} and \configFile{outputfileInside}{matrix}.
The regularization matrices are then used by \configClass{normalEquation:regularization}{normalEquationType:regularization}.
As an example, the two different areas could be oceanic regions on the one hand and continental areas on the other hand.


\keepXColumns
\begin{tabularx}{\textwidth}{N T A}
\hline
Name & Type & Annotation\\
\hline
\hfuzz=500pt\includegraphics[width=1em]{element-mustset.pdf}~outputfileInside & \hfuzz=500pt filename & \hfuzz=500pt \\
\hfuzz=500pt\includegraphics[width=1em]{element-mustset.pdf}~outputfileOutside & \hfuzz=500pt filename & \hfuzz=500pt \\
\hfuzz=500pt\includegraphics[width=1em]{element-mustset-unbounded.pdf}~grid & \hfuzz=500pt \hyperref[gridType]{grid} & \hfuzz=500pt nodal point distribution of parameters, e.g harmonics splines\\
\hfuzz=500pt\includegraphics[width=1em]{element-mustset-unbounded.pdf}~border & \hfuzz=500pt \hyperref[borderType]{border} & \hfuzz=500pt regularization areas, e.g land and ocean\\
\hfuzz=500pt\includegraphics[width=1em]{element.pdf}~R & \hfuzz=500pt double & \hfuzz=500pt reference radius for ellipsoidal coordinates for border\\
\hfuzz=500pt\includegraphics[width=1em]{element.pdf}~inverseFlattening & \hfuzz=500pt double & \hfuzz=500pt reference flattening for ellipsoidal coordinates for border, 0: spherical coordinates\\
\hline
\end{tabularx}

\clearpage
%==================================
\subsection{NormalsRegularizationSphericalHarmonics}\label{NormalsRegularizationSphericalHarmonics}
Diagonal regularization matrix from gravity field accuracies,
if not given from signal (cnm,snm), if not given from kaulas rule.
The inverse accuracies $1/\sigma_n^2$ are used as weights in the regularization matrix.
The diagonal is saved as Vector.

The corresponding pseudo observations can be computed with \program{Gravityfield2SphericalHarmonicsVector}.


\keepXColumns
\begin{tabularx}{\textwidth}{N T A}
\hline
Name & Type & Annotation\\
\hline
\hfuzz=500pt\includegraphics[width=1em]{element-mustset.pdf}~outputfileDiagonalmatrix & \hfuzz=500pt filename & \hfuzz=500pt \\
\hfuzz=500pt\includegraphics[width=1em]{element-unbounded.pdf}~gravityfield & \hfuzz=500pt \hyperref[gravityfieldType]{gravityfield} & \hfuzz=500pt use sigmas, if not given use signal (cnm,snm), if not given use kaulas rule\\
\hfuzz=500pt\includegraphics[width=1em]{element.pdf}~minRegularizationDegree & \hfuzz=500pt uint & \hfuzz=500pt \\
\hfuzz=500pt\includegraphics[width=1em]{element.pdf}~maxRegularizationDegree & \hfuzz=500pt uint & \hfuzz=500pt \\
\hfuzz=500pt\includegraphics[width=1em]{element-mustset.pdf}~minDegree & \hfuzz=500pt uint & \hfuzz=500pt \\
\hfuzz=500pt\includegraphics[width=1em]{element-mustset.pdf}~maxDegree & \hfuzz=500pt uint & \hfuzz=500pt \\
\hfuzz=500pt\includegraphics[width=1em]{element-mustset.pdf}~numbering & \hfuzz=500pt \hyperref[sphericalHarmonicsNumberingType]{sphericalHarmonicsNumbering} & \hfuzz=500pt numbering scheme for regul matrix\\
\hfuzz=500pt\includegraphics[width=1em]{element.pdf}~GM & \hfuzz=500pt double & \hfuzz=500pt Geocentric gravitational constant\\
\hfuzz=500pt\includegraphics[width=1em]{element.pdf}~R & \hfuzz=500pt double & \hfuzz=500pt reference radius\\
\hfuzz=500pt\includegraphics[width=1em]{element.pdf}~makeIsotropic & \hfuzz=500pt boolean & \hfuzz=500pt \\
\hfuzz=500pt\includegraphics[width=1em]{element.pdf}~kaulaPower & \hfuzz=500pt double & \hfuzz=500pt sigma = kaulaFactor*degree**kaulaPower\\
\hfuzz=500pt\includegraphics[width=1em]{element.pdf}~kaulaFactor & \hfuzz=500pt double & \hfuzz=500pt sigma = kaulaFactor*degree**kaulaPower\\
\hline
\end{tabularx}

\clearpage
%==================================
\subsection{NormalsReorder}\label{NormalsReorder}
Reorder \configFile{inputfileNormalEquation}{normalEquation} by selecting parameters in a specific order.
The \configClass{parameterSelection}{parameterSelectorType} also allows one to change dimension of the normal equations,
either by cutting parameters or by inserting zero rows/columns for additional parameters.
Without \configClass{parameterSelection}{parameterSelectorType} the order of parameters remains the same.
Additionally the block sizes of the files can be adjusted. If \config{outBlockSize} is set to zero,
the normal matrix is written to a single block file, which is needed by some programs.

To eliminate parameters without changing the result of the other parameters use \program{NormalsEliminate}.


\keepXColumns
\begin{tabularx}{\textwidth}{N T A}
\hline
Name & Type & Annotation\\
\hline
\hfuzz=500pt\includegraphics[width=1em]{element-mustset.pdf}~outputfileNormalEquation & \hfuzz=500pt filename & \hfuzz=500pt \\
\hfuzz=500pt\includegraphics[width=1em]{element-mustset.pdf}~inputfileNormalEquation & \hfuzz=500pt filename & \hfuzz=500pt \\
\hfuzz=500pt\includegraphics[width=1em]{element-unbounded.pdf}~parameterSelection & \hfuzz=500pt \hyperref[parameterSelectorType]{parameterSelector} & \hfuzz=500pt parameter order/selection of output normal equations\\
\hfuzz=500pt\includegraphics[width=1em]{element.pdf}~outBlockSize & \hfuzz=500pt uint & \hfuzz=500pt block size for distributing the normal equations, 0: one block\\
\hline
\end{tabularx}

This program is \reference{parallelized}{general.parallelization}.
\clearpage
%==================================
\subsection{NormalsScale}\label{NormalsScale}
Scales rows and columns of a system of \configFile{inputfileNormalEquation}{normalEquation}
given by a diagonal matrix \configFile{inputfileFactorVector}{matrix} $\M S$
\begin{equation}
  \bar{\M N} := \M S \M N \M S \qquad\text{and}\qquad \bar{\M n} := \M S \M n.
\end{equation}
The estimated solution is now
\begin{equation}
  \bar{\M x} := \M S^{-1} \M x.
\end{equation}
This is effectively the same as rescaling columns of the design matrix.
This program is useful when combining normal equations from different sources,
for example in case the units of certain parameters don't match.


\keepXColumns
\begin{tabularx}{\textwidth}{N T A}
\hline
Name & Type & Annotation\\
\hline
\hfuzz=500pt\includegraphics[width=1em]{element-mustset.pdf}~outputfileNormalEquation & \hfuzz=500pt filename & \hfuzz=500pt \\
\hfuzz=500pt\includegraphics[width=1em]{element-mustset.pdf}~inputfileNormalEquation & \hfuzz=500pt filename & \hfuzz=500pt \\
\hfuzz=500pt\includegraphics[width=1em]{element-mustset.pdf}~inputfileFactorVector & \hfuzz=500pt filename & \hfuzz=500pt Vector containing the factors\\
\hline
\end{tabularx}

This program is \reference{parallelized}{general.parallelization}.
\clearpage
%==================================
\subsection{NormalsSolverVCE}\label{NormalsSolverVCE}
This program accumulates \configClass{normalEquation}{normalEquationType}
and solves the total combined system.
The relative weigthing between the individual normals is determined iteratively
by means of variance component estimation (VCE). For a detailed description
of the used algorithm see \configClass{normalEquation}{normalEquationType}.

Besides the estimated parameter vector (\configFile{outputfileSolution}{matrix}) the
estimated accuracies (\configFile{outputfileSigmax}{matrix}) and the full covariance matrix
(\configFile{outputfileCovariance}{matrix}) can be saved. Also the combined normal system
can be written to \configFile{outputfileNormalEquation}{normalEquation}.

The \configFile{outputfileContribution}{matrix} is a matrix with rows for each estimated
parameter and columns for each \configClass{normalEquation}{normalEquationType}
and indicates the contribution of the individual normals to the estimated parameters.
Each row sum up to one.

See also \program{NormalsBuild}.


\keepXColumns
\begin{tabularx}{\textwidth}{N T A}
\hline
Name & Type & Annotation\\
\hline
\hfuzz=500pt\includegraphics[width=1em]{element.pdf}~outputfileSolution & \hfuzz=500pt filename & \hfuzz=500pt parameter vector\\
\hfuzz=500pt\includegraphics[width=1em]{element.pdf}~outputfileSigmax & \hfuzz=500pt filename & \hfuzz=500pt standard deviations of the parameters (sqrt of the diagonal of the inverse normal equation)\\
\hfuzz=500pt\includegraphics[width=1em]{element.pdf}~outputfileCovariance & \hfuzz=500pt filename & \hfuzz=500pt full covariance matrix\\
\hfuzz=500pt\includegraphics[width=1em]{element.pdf}~outputfileContribution & \hfuzz=500pt filename & \hfuzz=500pt contribution of normal system components to the solution vector\\
\hfuzz=500pt\includegraphics[width=1em]{element.pdf}~outputfileVarianceFactors & \hfuzz=500pt filename & \hfuzz=500pt estimated variance factors as vector\\
\hfuzz=500pt\includegraphics[width=1em]{element.pdf}~outputfileNormalEquation & \hfuzz=500pt filename & \hfuzz=500pt the combined normal equation system\\
\hfuzz=500pt\includegraphics[width=1em]{element-mustset-unbounded.pdf}~normalEquation & \hfuzz=500pt \hyperref[normalEquationType]{normalEquation} & \hfuzz=500pt \\
\hfuzz=500pt\includegraphics[width=1em]{element.pdf}~inputfileApproxSolution & \hfuzz=500pt filename & \hfuzz=500pt to accelerate convergence\\
\hfuzz=500pt\includegraphics[width=1em]{element.pdf}~rightHandSideNumberVCE & \hfuzz=500pt uint & \hfuzz=500pt the right hand side number for estimation of variance factors\\
\hfuzz=500pt\includegraphics[width=1em]{element.pdf}~normalsBlockSize & \hfuzz=500pt uint & \hfuzz=500pt block size for distributing the normal equations, 0: one block\\
\hfuzz=500pt\includegraphics[width=1em]{element.pdf}~maxIterationCount & \hfuzz=500pt uint & \hfuzz=500pt maximum number of iterations for variance component estimation\\
\hline
\end{tabularx}

This program is \reference{parallelized}{general.parallelization}.
\clearpage
%==================================
\subsection{NormalsTemporalCombination}\label{NormalsTemporalCombination}
This program reads a times series of \configFile{inputfileNormalequation}{normalEquation}
with asscociated \configClass{timeSeries}{timeSeriesType} and setup a new combined normal equation system.
For each parameter a \configClass{parametrizationTemporal}{parametrizationTemporalType} is used.

It can be used to estimate trend and annual spherical harmonic coefficients from monthly GRACE normal equations.


\keepXColumns
\begin{tabularx}{\textwidth}{N T A}
\hline
Name & Type & Annotation\\
\hline
\hfuzz=500pt\includegraphics[width=1em]{element-mustset.pdf}~outputfileNormalEquation & \hfuzz=500pt filename & \hfuzz=500pt \\
\hfuzz=500pt\includegraphics[width=1em]{element-mustset-unbounded.pdf}~inputfileNormalEquation & \hfuzz=500pt filename & \hfuzz=500pt normal equations for each point in time\\
\hfuzz=500pt\includegraphics[width=1em]{element-mustset-unbounded.pdf}~timeSeries & \hfuzz=500pt \hyperref[timeSeriesType]{timeSeries} & \hfuzz=500pt times of each normal equations\\
\hfuzz=500pt\includegraphics[width=1em]{element-mustset-unbounded.pdf}~parametrizationTemporal & \hfuzz=500pt \hyperref[parametrizationTemporalType]{parametrizationTemporal} & \hfuzz=500pt \\
\hline
\end{tabularx}

This program is \reference{parallelized}{general.parallelization}.
\clearpage
%==================================
\subsection{ParameterNamesCreate}\label{ParameterNamesCreate}
Generate a \configFile{outputfileParameterNames}{parameterName} by \configClass{parameterName}{parameterNamesType}.
This file can be used in \program{NormalsCreate} or in the class \configClass{parameterSelector}{parameterSelectorType}.


\keepXColumns
\begin{tabularx}{\textwidth}{N T A}
\hline
Name & Type & Annotation\\
\hline
\hfuzz=500pt\includegraphics[width=1em]{element-mustset.pdf}~outputfileParameterNames & \hfuzz=500pt filename & \hfuzz=500pt output parameter names file\\
\hfuzz=500pt\includegraphics[width=1em]{element-mustset-unbounded.pdf}~parameterName & \hfuzz=500pt \hyperref[parameterNamesType]{parameterNames} & \hfuzz=500pt \\
\hline
\end{tabularx}

\clearpage
%==================================
\subsection{ParameterSelection2IndexVector}\label{ParameterSelection2IndexVector}
Generate index vector from parameter selection in \file{matrix format}{matrix}.
This vector can be used in \program{MatrixCalculate}
with \configClass{matrix:reorder}{matrixGeneratorType:reorder}
to reorder arbitrary vectors and matrices similar to \program{NormalsReorder}.

The \configClass{parameterSelection}{parameterSelectorType} allows reordering and dimension changes,
either by cutting parameters or by inserting additional parameters.
\configFile{outputfileIndexVector}{matrix} contains indices of parameters in
\configFile{inputfileParameterNames}{parameterName} or -1 for newly added parameters.
\configFile{outputfileParameterNames}{parameterName} contains the selected parameter names.


\keepXColumns
\begin{tabularx}{\textwidth}{N T A}
\hline
Name & Type & Annotation\\
\hline
\hfuzz=500pt\includegraphics[width=1em]{element.pdf}~outputfileIndexVector & \hfuzz=500pt filename & \hfuzz=500pt indices of source parameters in target normal equations\\
\hfuzz=500pt\includegraphics[width=1em]{element.pdf}~outputfileParameterNames & \hfuzz=500pt filename & \hfuzz=500pt output parameter names file\\
\hfuzz=500pt\includegraphics[width=1em]{element-mustset.pdf}~inputfileParameterNames & \hfuzz=500pt filename & \hfuzz=500pt parameter names file of source normal equations\\
\hfuzz=500pt\includegraphics[width=1em]{element-mustset-unbounded.pdf}~parameterSelection & \hfuzz=500pt \hyperref[parameterSelectorType]{parameterSelector} & \hfuzz=500pt parameter order/selection of target normal equations\\
\hline
\end{tabularx}

\clearpage
%==================================
\section{Programs: Orbit}
\subsection{Orbit2ArgumentOfLatitude}\label{Orbit2ArgumentOfLatitude}
This program computes the argument of latitude of an \file{orbit}{instrument}
and writes it as \file{instrument file}{instrument} (MISCVALUE(S)).
The data of \configFile{inputfileInstrument}{instrument} are appended as values to each epoch.

\fig{!hb}{0.8}{instrumentOrbit2ArgumentOfLatitude}{fig:instrumentOrbit2ArgumentOfLatitude}{Derivation filtered GRACE range-rate residuals.}


\keepXColumns
\begin{tabularx}{\textwidth}{N T A}
\hline
Name & Type & Annotation\\
\hline
\hfuzz=500pt\includegraphics[width=1em]{element-mustset.pdf}~outputfileArgOfLatitude & \hfuzz=500pt filename & \hfuzz=500pt instrument file (MISCVALUE(S): argLat, ...)\\
\hfuzz=500pt\includegraphics[width=1em]{element-mustset.pdf}~inputfileOrbit & \hfuzz=500pt filename & \hfuzz=500pt \\
\hfuzz=500pt\includegraphics[width=1em]{element-unbounded.pdf}~inputfileInstrument & \hfuzz=500pt filename & \hfuzz=500pt data are appended\\
\hline
\end{tabularx}

This program is \reference{parallelized}{general.parallelization}.
\clearpage
%==================================
\subsection{Orbit2BetaPrimeAngle}\label{Orbit2BetaPrimeAngle}
This program computes the beta prime angle (between the orbital plane and earth-sun direction)
and writes it as MISCVALUE(S) \file{instrument file}{instrument}. The angle is calculated w.r.t the sun (per default),
but can be changed.
The data of \configFile{inputfileInstrument}{instrument} are appended as values to each epoch.


\keepXColumns
\begin{tabularx}{\textwidth}{N T A}
\hline
Name & Type & Annotation\\
\hline
\hfuzz=500pt\includegraphics[width=1em]{element-mustset.pdf}~outputfileBetaAngle & \hfuzz=500pt filename & \hfuzz=500pt instrument file (MISCVALUE(S): beta', ...)\\
\hfuzz=500pt\includegraphics[width=1em]{element-mustset.pdf}~inputfileOrbit & \hfuzz=500pt filename & \hfuzz=500pt \\
\hfuzz=500pt\includegraphics[width=1em]{element-unbounded.pdf}~inputfileInstrument & \hfuzz=500pt filename & \hfuzz=500pt data are appended\\
\hfuzz=500pt\includegraphics[width=1em]{element-mustset.pdf}~ephemerides & \hfuzz=500pt \hyperref[ephemeridesType]{ephemerides} & \hfuzz=500pt \\
\hfuzz=500pt\includegraphics[width=1em]{element.pdf}~planet & \hfuzz=500pt \hyperref[planetType]{planet} & \hfuzz=500pt \\
\hline
\end{tabularx}

This program is \reference{parallelized}{general.parallelization}.
\clearpage
%==================================
\subsection{Orbit2EarthFixedOrbit}\label{Orbit2EarthFixedOrbit}
Normally the orbits in GROOPS are given in the celestial reference frame (CRF) with the
origin in the center of mass (CoM). This program rotates the orbit with
\configClass{earthRotation}{earthRotationType} from CRF to the TRF.

To additionally tranform into the center of solid Earth (CE) frame (or center of Figure (CF)),
a correction can be applied by providing degree one coefficients of a
\configClass{gravityfield}{gravityfieldType} (e.g. ocean tides).

If \config{celestial2terrestrial} is set to no, the inverse transformation is applied.

See also \program{InstrumentRotate}.


\keepXColumns
\begin{tabularx}{\textwidth}{N T A}
\hline
Name & Type & Annotation\\
\hline
\hfuzz=500pt\includegraphics[width=1em]{element-mustset.pdf}~outputfileOrbit & \hfuzz=500pt filename & \hfuzz=500pt \\
\hfuzz=500pt\includegraphics[width=1em]{element-mustset.pdf}~inputfileOrbit & \hfuzz=500pt filename & \hfuzz=500pt \\
\hfuzz=500pt\includegraphics[width=1em]{element-mustset.pdf}~earthRotation & \hfuzz=500pt \hyperref[earthRotationType]{earthRotation} & \hfuzz=500pt transformation from CRF to TRF\\
\hfuzz=500pt\includegraphics[width=1em]{element-unbounded.pdf}~gravityfield & \hfuzz=500pt \hyperref[gravityfieldType]{gravityfield} & \hfuzz=500pt degree 1 fluid mantle for CM2CE correction\\
\hfuzz=500pt\includegraphics[width=1em]{element.pdf}~celestial2terrestrial & \hfuzz=500pt boolean & \hfuzz=500pt yes: crf-\$>\$trf, no: trf-\$>\$crf\\
\hline
\end{tabularx}

This program is \reference{parallelized}{general.parallelization}.
\clearpage
%==================================
\subsection{Orbit2EclipseFactor}\label{Orbit2EclipseFactor}
This program generates an \file{instrument file}{instrument} (MISCVALUE(S)) containing the eclipse factor for a given set of orbit.
The data of \configFile{inputfileInstrument}{instrument} are appended as values to each epoch.


\keepXColumns
\begin{tabularx}{\textwidth}{N T A}
\hline
Name & Type & Annotation\\
\hline
\hfuzz=500pt\includegraphics[width=1em]{element-mustset.pdf}~outputfileEclipseFactor & \hfuzz=500pt filename & \hfuzz=500pt instrument file (MISCVALUE(S): eclipse, ...)\\
\hfuzz=500pt\includegraphics[width=1em]{element-mustset.pdf}~inputfileOrbit & \hfuzz=500pt filename & \hfuzz=500pt \\
\hfuzz=500pt\includegraphics[width=1em]{element-unbounded.pdf}~inputfileInstrument & \hfuzz=500pt filename & \hfuzz=500pt data are appended\\
\hfuzz=500pt\includegraphics[width=1em]{element-mustset.pdf}~ephemerides & \hfuzz=500pt \hyperref[ephemeridesType]{ephemerides} & \hfuzz=500pt \\
\hfuzz=500pt\includegraphics[width=1em]{element-mustset.pdf}~eclipse & \hfuzz=500pt \hyperref[eclipseType]{eclipse} & \hfuzz=500pt \\
\hline
\end{tabularx}

This program is \reference{parallelized}{general.parallelization}.
\clearpage
%==================================
\subsection{Orbit2Groundtracks}\label{Orbit2Groundtracks}
This program write \file{satellites positions}{instrument} as \file{gridded data}{griddedData}
(\config{outputfileTrackGriddedData}) in a terrestrial reference frame. The points are expressed as ellipsoidal coordinates
(longitude, latitude, height) based on a reference ellipsoid with parameters \config{R} and
\config{inverseFlattening}. The orbit data are given in the celestial frame so \configClass{earthRotation}{earthRotationType}
is needed to transform the data into the terrestrial frame.
The data of \configFile{inputfileInstrument}{instrument} are appended as values to each point.


\keepXColumns
\begin{tabularx}{\textwidth}{N T A}
\hline
Name & Type & Annotation\\
\hline
\hfuzz=500pt\includegraphics[width=1em]{element-mustset.pdf}~outputfileGriddedData & \hfuzz=500pt filename & \hfuzz=500pt positions as gridded data\\
\hfuzz=500pt\includegraphics[width=1em]{element-mustset.pdf}~inputfileOrbit & \hfuzz=500pt filename & \hfuzz=500pt \\
\hfuzz=500pt\includegraphics[width=1em]{element-unbounded.pdf}~inputfileInstrument & \hfuzz=500pt filename & \hfuzz=500pt values at grid points\\
\hfuzz=500pt\includegraphics[width=1em]{element-mustset.pdf}~earthRotation & \hfuzz=500pt \hyperref[earthRotationType]{earthRotation} & \hfuzz=500pt transformation from CRF to TRF\\
\hfuzz=500pt\includegraphics[width=1em]{element.pdf}~R & \hfuzz=500pt double & \hfuzz=500pt reference radius for ellipsoidal coordinates\\
\hfuzz=500pt\includegraphics[width=1em]{element.pdf}~inverseFlattening & \hfuzz=500pt double & \hfuzz=500pt reference flattening for ellipsoidal coordinates\\
\hline
\end{tabularx}

\clearpage
%==================================
\subsection{Orbit2Kepler}\label{Orbit2Kepler}
This program computes the osculating Keplerian elements from position and velocity
of a given \configFile{inputfileOrbit}{instrument}.
The \configFile{inputfileOrbit}{instrument} must contain positions and velocities (see \program{OrbitAddVelocityAndAcceleration}).

The \config{outputfileKepler} is an \file{instrument file}{instrument} (MISCVALUES)
with the Keplerian elements at each epoch in the following order
\begin{itemize}
\item Ascending Node $\Omega$ [degree]
\item Inclination $i$ [degree]
\item Argument of perigee $\omega$ [degree]
\item major axis $a$ [m]
\item eccentricity $e$
\item mean anomaly $M$ [degree]
\item transit time of perigee $\tau$ [mjd]
\end{itemize}
The data of \configFile{inputfileInstrument}{instrument} are appended as values to each epoch.


\keepXColumns
\begin{tabularx}{\textwidth}{N T A}
\hline
Name & Type & Annotation\\
\hline
\hfuzz=500pt\includegraphics[width=1em]{element-mustset.pdf}~outputfileKepler & \hfuzz=500pt filename & \hfuzz=500pt instrument file (MISCVALUES: Omega, i, omega [degree], a [m], e, M [degree], tau [mjd], ...)\\
\hfuzz=500pt\includegraphics[width=1em]{element-mustset.pdf}~inputfileOrbit & \hfuzz=500pt filename & \hfuzz=500pt position and velocity at each epoch define the kepler orbit\\
\hfuzz=500pt\includegraphics[width=1em]{element-unbounded.pdf}~inputfileInstrument & \hfuzz=500pt filename & \hfuzz=500pt data is appended\\
\hfuzz=500pt\includegraphics[width=1em]{element.pdf}~GM & \hfuzz=500pt double & \hfuzz=500pt Geocentric gravitational constant\\
\hline
\end{tabularx}

This program is \reference{parallelized}{general.parallelization}.
\clearpage
%==================================
\subsection{Orbit2MagneticField}\label{Orbit2MagneticField}
This program computes the magentic field vector($x, y, z$ $[Tesla = kg/A/s^2]$ in CRF))
along an \file{orbit}{instrument} and writes it as \file{instrument file}{instrument} (MISCVALUES).
The data of \configFile{inputfileInstrument}{instrument} are appended as data columns to each epoch.


\keepXColumns
\begin{tabularx}{\textwidth}{N T A}
\hline
Name & Type & Annotation\\
\hline
\hfuzz=500pt\includegraphics[width=1em]{element-mustset.pdf}~outputfileMagneticField & \hfuzz=500pt filename & \hfuzz=500pt instrument file (x,y,z in CRF [Tesla = kg/A/s\textasciicircum{}2]), ...)\\
\hfuzz=500pt\includegraphics[width=1em]{element-mustset.pdf}~inputfileOrbit & \hfuzz=500pt filename & \hfuzz=500pt \\
\hfuzz=500pt\includegraphics[width=1em]{element-unbounded.pdf}~inputfileInstrument & \hfuzz=500pt filename & \hfuzz=500pt data are appended to output file\\
\hfuzz=500pt\includegraphics[width=1em]{element-mustset.pdf}~magnetosphere & \hfuzz=500pt \hyperref[magnetosphereType]{magnetosphere} & \hfuzz=500pt \\
\hfuzz=500pt\includegraphics[width=1em]{element-mustset.pdf}~earthRotation & \hfuzz=500pt \hyperref[earthRotationType]{earthRotation} & \hfuzz=500pt \\
\hline
\end{tabularx}

This program is \reference{parallelized}{general.parallelization}.
\clearpage
%==================================
\subsection{Orbit2ThermosphericState}\label{Orbit2ThermosphericState}
This program computes the thermosperic state (density, temperature, wind (x,y,z in CRF))
based on emprical models along an \file{orbit}{instrument}
and writes it as \file{instrument file}{instrument} (MISCVALUES).
The wind is given in an celestial reference frame (CRF).
The data of \configFile{inputfileInstrument}{instrument} are appended as values to each epoch.


\keepXColumns
\begin{tabularx}{\textwidth}{N T A}
\hline
Name & Type & Annotation\\
\hline
\hfuzz=500pt\includegraphics[width=1em]{element-mustset.pdf}~outputfileThermosphericState & \hfuzz=500pt filename & \hfuzz=500pt instrument file (MISCVALUES: density, temperature, wind (x,y,z in CRF), ...)\\
\hfuzz=500pt\includegraphics[width=1em]{element-mustset.pdf}~inputfileOrbit & \hfuzz=500pt filename & \hfuzz=500pt \\
\hfuzz=500pt\includegraphics[width=1em]{element-unbounded.pdf}~inputfileInstrument & \hfuzz=500pt filename & \hfuzz=500pt data are appended to output file\\
\hfuzz=500pt\includegraphics[width=1em]{element-mustset.pdf}~thermosphere & \hfuzz=500pt \hyperref[thermosphereType]{thermosphere} & \hfuzz=500pt \\
\hfuzz=500pt\includegraphics[width=1em]{element-mustset.pdf}~earthRotation & \hfuzz=500pt \hyperref[earthRotationType]{earthRotation} & \hfuzz=500pt \\
\hline
\end{tabularx}

This program is \reference{parallelized}{general.parallelization}.
\clearpage
%==================================
\subsection{OrbitAddVelocityAndAcceleration}\label{OrbitAddVelocityAndAcceleration}
This program computes velocities and accelerations from a given \file{orbit}{instrument}
by differentiating a moving polynomial.
The values are saved in one output file which then contains orbit, velocity and acceleration.


\keepXColumns
\begin{tabularx}{\textwidth}{N T A}
\hline
Name & Type & Annotation\\
\hline
\hfuzz=500pt\includegraphics[width=1em]{element-mustset.pdf}~outputfileOrbit & \hfuzz=500pt filename & \hfuzz=500pt \\
\hfuzz=500pt\includegraphics[width=1em]{element-mustset.pdf}~inputfileOrbit & \hfuzz=500pt filename & \hfuzz=500pt \\
\hfuzz=500pt\includegraphics[width=1em]{element.pdf}~polynomialDegree & \hfuzz=500pt uint & \hfuzz=500pt Polynomial degree, must be even!\\
\hline
\end{tabularx}

This program is \reference{parallelized}{general.parallelization}.
\clearpage
%==================================
\subsection{PlanetOrbit}\label{PlanetOrbit}
Creates an \file{orbit file}{instrument} of sun, moon, or planets.
The orbit is given in the celestial reference frame (CRF)
or alternatively in the terrestrial reference frame (TRF)
if \configClass{earthRotation}{earthRotationType} is provided.


\keepXColumns
\begin{tabularx}{\textwidth}{N T A}
\hline
Name & Type & Annotation\\
\hline
\hfuzz=500pt\includegraphics[width=1em]{element-mustset.pdf}~outputfileOrbit & \hfuzz=500pt filename & \hfuzz=500pt \\
\hfuzz=500pt\includegraphics[width=1em]{element-mustset.pdf}~planet & \hfuzz=500pt \hyperref[planetType]{planet} & \hfuzz=500pt \\
\hfuzz=500pt\includegraphics[width=1em]{element-mustset-unbounded.pdf}~timeSeries & \hfuzz=500pt \hyperref[timeSeriesType]{timeSeries} & \hfuzz=500pt \\
\hfuzz=500pt\includegraphics[width=1em]{element-mustset.pdf}~ephemerides & \hfuzz=500pt \hyperref[ephemeridesType]{ephemerides} & \hfuzz=500pt \\
\hfuzz=500pt\includegraphics[width=1em]{element.pdf}~earthRotation & \hfuzz=500pt \hyperref[earthRotationType]{earthRotation} & \hfuzz=500pt transform orbits into TRF\\
\hline
\end{tabularx}

\clearpage
%==================================
\section{Programs: Plot}
\subsection{PlotDegreeAmplitudes}\label{PlotDegreeAmplitudes}
Plot degree amplitudes of potential coefficients computed by \program{Gravityfield2DegreeAmplitudes}
or \program{PotentialCoefficients2DegreeAmplitudes} using the GMT Generic Mapping Tools
(\url{https://www.generic-mapping-tools.org}).
A variety of image file formats are supported (e.g. png, jpg, eps) determined by the extension of \config{outputfile}.
This is a convenience program with meaningful default values. The same plots can be generated with the more general \program{PlotGraph}.

\fig{!hb}{0.8}{plotDegreeAmplitudes}{fig:plotDegreeAmplitudes}{Comparison of GRACE solutions (2008-06) with GOCO06s.}


\keepXColumns
\begin{tabularx}{\textwidth}{N T A}
\hline
Name & Type & Annotation\\
\hline
\hfuzz=500pt\includegraphics[width=1em]{element-mustset.pdf}~outputfile & \hfuzz=500pt filename & \hfuzz=500pt *.png, *.jpg, *.eps, ...\\
\hfuzz=500pt\includegraphics[width=1em]{element.pdf}~title & \hfuzz=500pt string & \hfuzz=500pt \\
\hfuzz=500pt\includegraphics[width=1em]{element-mustset-unbounded.pdf}~layer & \hfuzz=500pt \hyperref[plotGraphLayerType]{plotGraphLayer} & \hfuzz=500pt \\
\hfuzz=500pt\includegraphics[width=1em]{element.pdf}~minDegree & \hfuzz=500pt uint & \hfuzz=500pt \\
\hfuzz=500pt\includegraphics[width=1em]{element.pdf}~maxDegree & \hfuzz=500pt uint & \hfuzz=500pt \\
\hfuzz=500pt\includegraphics[width=1em]{element.pdf}~majorTickSpacingDegree & \hfuzz=500pt double & \hfuzz=500pt boundary annotation\\
\hfuzz=500pt\includegraphics[width=1em]{element.pdf}~minorTickSpacingDegree & \hfuzz=500pt double & \hfuzz=500pt frame tick spacing\\
\hfuzz=500pt\includegraphics[width=1em]{element.pdf}~gridLineSpacingDegree & \hfuzz=500pt double & \hfuzz=500pt gridline spacing\\
\hfuzz=500pt\includegraphics[width=1em]{element.pdf}~labelDegree & \hfuzz=500pt string & \hfuzz=500pt description of the x-axis\\
\hfuzz=500pt\includegraphics[width=1em]{element.pdf}~logarithmicDegree & \hfuzz=500pt boolean & \hfuzz=500pt use logarithmic scale for the x-axis\\
\hfuzz=500pt\includegraphics[width=1em]{element.pdf}~minY & \hfuzz=500pt double & \hfuzz=500pt \\
\hfuzz=500pt\includegraphics[width=1em]{element.pdf}~maxY & \hfuzz=500pt double & \hfuzz=500pt \\
\hfuzz=500pt\includegraphics[width=1em]{element.pdf}~majorTickSpacingY & \hfuzz=500pt double & \hfuzz=500pt boundary annotation\\
\hfuzz=500pt\includegraphics[width=1em]{element.pdf}~minorTickSpacingY & \hfuzz=500pt double & \hfuzz=500pt frame tick spacing\\
\hfuzz=500pt\includegraphics[width=1em]{element.pdf}~gridLineSpacingY & \hfuzz=500pt double & \hfuzz=500pt gridline spacing\\
\hfuzz=500pt\includegraphics[width=1em]{element.pdf}~unitY & \hfuzz=500pt string & \hfuzz=500pt appended to axis values\\
\hfuzz=500pt\includegraphics[width=1em]{element.pdf}~labelY & \hfuzz=500pt string & \hfuzz=500pt description of the y-axis\\
\hfuzz=500pt\includegraphics[width=1em]{element.pdf}~logarithmicY & \hfuzz=500pt boolean & \hfuzz=500pt use logarithmic scale for the y-axis\\
\hfuzz=500pt\includegraphics[width=1em]{element.pdf}~gridLine & \hfuzz=500pt \hyperref[plotLineType]{plotLine} & \hfuzz=500pt The style of the grid lines.\\
\hfuzz=500pt\includegraphics[width=1em]{element.pdf}~legend & \hfuzz=500pt \hyperref[plotLegendType]{plotLegend} & \hfuzz=500pt \\
\hfuzz=500pt\includegraphics[width=1em]{element-mustset.pdf}~options & \hfuzz=500pt sequence & \hfuzz=500pt further options...\\
\hfuzz=500pt\includegraphics[width=1em]{connector.pdf}\includegraphics[width=1em]{element.pdf}~width & \hfuzz=500pt double & \hfuzz=500pt in cm\\
\hfuzz=500pt\includegraphics[width=1em]{connector.pdf}\includegraphics[width=1em]{element.pdf}~height & \hfuzz=500pt double & \hfuzz=500pt in cm\\
\hfuzz=500pt\includegraphics[width=1em]{connector.pdf}\includegraphics[width=1em]{element.pdf}~titleFontSize & \hfuzz=500pt uint & \hfuzz=500pt in pt\\
\hfuzz=500pt\includegraphics[width=1em]{connector.pdf}\includegraphics[width=1em]{element.pdf}~marginTitle & \hfuzz=500pt double & \hfuzz=500pt between title and figure [cm]\\
\hfuzz=500pt\includegraphics[width=1em]{connector.pdf}\includegraphics[width=1em]{element.pdf}~drawGridOnTop & \hfuzz=500pt boolean & \hfuzz=500pt grid lines above all other lines/points\\
\hfuzz=500pt\includegraphics[width=1em]{connector.pdf}\includegraphics[width=1em]{element-unbounded.pdf}~options & \hfuzz=500pt string & \hfuzz=500pt \\
\hfuzz=500pt\includegraphics[width=1em]{connector.pdf}\includegraphics[width=1em]{element.pdf}~transparent & \hfuzz=500pt boolean & \hfuzz=500pt make background transparent\\
\hfuzz=500pt\includegraphics[width=1em]{connector.pdf}\includegraphics[width=1em]{element.pdf}~dpi & \hfuzz=500pt uint & \hfuzz=500pt use this resolution when rasterizing postscript file\\
\hfuzz=500pt\includegraphics[width=1em]{connector.pdf}\includegraphics[width=1em]{element.pdf}~removeFiles & \hfuzz=500pt boolean & \hfuzz=500pt remove .gmt and script files\\
\hline
\end{tabularx}

\clearpage
%==================================
\subsection{PlotGraph}\label{PlotGraph}
Generates a two dimensional xy plot using the GMT Generic Mapping Tools (\url{https://www.generic-mapping-tools.org}).
A variety of image file formats are supported (e.g. png, jpg, eps) determined by the extension of \config{outputfile}.

The plotting area is defined by the two axes \configClass{axisX/Y}{plotAxisType}. An alternative \configClass{axisY2}{plotAxisType}
on the right hand side can be added. The content of the graph itself is defined
by one or more \configClass{layer}{plotGraphLayerType}s.

The plot programs create a temporary directory in the path of \config{outputfile}, writes all needed data into it,
generates a batch/shell script with the GMT commands, execute it, and remove the temporary directory.
With setting \config{options:removeFiles}=false the last step is skipped and it is possible to adjust the plot manually
to specific publication needs. Individual GMT settings are adjusted with \config{options:options}="\verb|FORMAT=value|",
see \url{https://docs.generic-mapping-tools.org/latest/gmt.conf.html}.

See also: \program{PlotDegreeAmplitudes}, \program{PlotMap}, \program{PlotMatrix}, \program{PlotSphericalHarmonicsTriangle}.


\keepXColumns
\begin{tabularx}{\textwidth}{N T A}
\hline
Name & Type & Annotation\\
\hline
\hfuzz=500pt\includegraphics[width=1em]{element-mustset.pdf}~outputfile & \hfuzz=500pt filename & \hfuzz=500pt *.png, *.jpg, *.eps, ...\\
\hfuzz=500pt\includegraphics[width=1em]{element.pdf}~title & \hfuzz=500pt string & \hfuzz=500pt \\
\hfuzz=500pt\includegraphics[width=1em]{element-mustset-unbounded.pdf}~layer & \hfuzz=500pt \hyperref[plotGraphLayerType]{plotGraphLayer} & \hfuzz=500pt \\
\hfuzz=500pt\includegraphics[width=1em]{element-mustset.pdf}~axisX & \hfuzz=500pt \hyperref[plotAxisType]{plotAxis} & \hfuzz=500pt \\
\hfuzz=500pt\includegraphics[width=1em]{element-mustset.pdf}~axisY & \hfuzz=500pt \hyperref[plotAxisType]{plotAxis} & \hfuzz=500pt \\
\hfuzz=500pt\includegraphics[width=1em]{element.pdf}~axisY2 & \hfuzz=500pt \hyperref[plotAxisType]{plotAxis} & \hfuzz=500pt Second y-axis on right hand side\\
\hfuzz=500pt\includegraphics[width=1em]{element.pdf}~colorbar & \hfuzz=500pt \hyperref[plotColorbarType]{plotColorbar} & \hfuzz=500pt \\
\hfuzz=500pt\includegraphics[width=1em]{element.pdf}~legend & \hfuzz=500pt \hyperref[plotLegendType]{plotLegend} & \hfuzz=500pt \\
\hfuzz=500pt\includegraphics[width=1em]{element-mustset.pdf}~options & \hfuzz=500pt sequence & \hfuzz=500pt further options...\\
\hfuzz=500pt\includegraphics[width=1em]{connector.pdf}\includegraphics[width=1em]{element.pdf}~width & \hfuzz=500pt double & \hfuzz=500pt in cm\\
\hfuzz=500pt\includegraphics[width=1em]{connector.pdf}\includegraphics[width=1em]{element.pdf}~height & \hfuzz=500pt double & \hfuzz=500pt in cm\\
\hfuzz=500pt\includegraphics[width=1em]{connector.pdf}\includegraphics[width=1em]{element.pdf}~titleFontSize & \hfuzz=500pt uint & \hfuzz=500pt in pt\\
\hfuzz=500pt\includegraphics[width=1em]{connector.pdf}\includegraphics[width=1em]{element.pdf}~marginTitle & \hfuzz=500pt double & \hfuzz=500pt between title and figure [cm]\\
\hfuzz=500pt\includegraphics[width=1em]{connector.pdf}\includegraphics[width=1em]{element.pdf}~drawGridOnTop & \hfuzz=500pt boolean & \hfuzz=500pt grid lines above all other lines/points\\
\hfuzz=500pt\includegraphics[width=1em]{connector.pdf}\includegraphics[width=1em]{element-unbounded.pdf}~options & \hfuzz=500pt string & \hfuzz=500pt \\
\hfuzz=500pt\includegraphics[width=1em]{connector.pdf}\includegraphics[width=1em]{element.pdf}~transparent & \hfuzz=500pt boolean & \hfuzz=500pt make background transparent\\
\hfuzz=500pt\includegraphics[width=1em]{connector.pdf}\includegraphics[width=1em]{element.pdf}~dpi & \hfuzz=500pt uint & \hfuzz=500pt use this resolution when rasterizing postscript file\\
\hfuzz=500pt\includegraphics[width=1em]{connector.pdf}\includegraphics[width=1em]{element.pdf}~removeFiles & \hfuzz=500pt boolean & \hfuzz=500pt remove .gmt and script files\\
\hline
\end{tabularx}

\clearpage
%==================================
\subsection{PlotMap}\label{PlotMap}
Generates a map using the GMT Generic Mapping Tools (\url{https://www.generic-mapping-tools.org}).
A variety of image file formats are supported (e.g. png, jpg, eps) determined by the extension of \config{outputfile}.

The base map is defined by a \configClass{projection}{plotMapProjectionType} of an ellipsoid (\config{R}, \config{inverseFlattening}).
The content of the map itself is defined by one or more \configClass{layer}{plotMapLayerType}s.

The plot programs create a temporary directory in the path of \config{outputfile}, writes all needed data into it,
generates a batch/shell script with the GMT commands, execute it, and remove the temporary directory.
With setting \config{options:removeFiles}=false the last step is skipped and it is possible to adjust the plot manually
to specific publication needs. Individual GMT settings are adjusted with \config{options:options}="\verb|FORMAT=value|",
see \url{https://docs.generic-mapping-tools.org/latest/gmt.conf.html}.

See also: \program{PlotDegreeAmplitudes}, \program{PlotGraph}, \program{PlotMatrix}, \program{PlotSphericalHarmonicsTriangle}.

\fig{!hb}{0.8}{plotMap}{fig:plotMap}{A Robinson projection with griddedData (geoid), coast, polygon (amazon), and points (IGS stations) layer.}


\keepXColumns
\begin{tabularx}{\textwidth}{N T A}
\hline
Name & Type & Annotation\\
\hline
\hfuzz=500pt\includegraphics[width=1em]{element-mustset.pdf}~outputfile & \hfuzz=500pt filename & \hfuzz=500pt *.png, *.jpg, *.eps, ...\\
\hfuzz=500pt\includegraphics[width=1em]{element.pdf}~title & \hfuzz=500pt string & \hfuzz=500pt \\
\hfuzz=500pt\includegraphics[width=1em]{element.pdf}~statisticInfos & \hfuzz=500pt boolean & \hfuzz=500pt \\
\hfuzz=500pt\includegraphics[width=1em]{element-mustset-unbounded.pdf}~layer & \hfuzz=500pt \hyperref[plotMapLayerType]{plotMapLayer} & \hfuzz=500pt \\
\hfuzz=500pt\includegraphics[width=1em]{element.pdf}~R & \hfuzz=500pt double & \hfuzz=500pt reference radius for ellipsoidal coordinates on output\\
\hfuzz=500pt\includegraphics[width=1em]{element.pdf}~inverseFlattening & \hfuzz=500pt double & \hfuzz=500pt reference flattening for ellipsoidal coordinates on output, 0: spherical coordinates\\
\hfuzz=500pt\includegraphics[width=1em]{element.pdf}~minLambda & \hfuzz=500pt angle & \hfuzz=500pt min. longitude (default: compute from input data)\\
\hfuzz=500pt\includegraphics[width=1em]{element.pdf}~maxLambda & \hfuzz=500pt angle & \hfuzz=500pt max. longitude (default: compute from input data)\\
\hfuzz=500pt\includegraphics[width=1em]{element.pdf}~minPhi & \hfuzz=500pt angle & \hfuzz=500pt min. latitude (default: compute from input data)\\
\hfuzz=500pt\includegraphics[width=1em]{element.pdf}~maxPhi & \hfuzz=500pt angle & \hfuzz=500pt max. latitude (default: compute from input data)\\
\hfuzz=500pt\includegraphics[width=1em]{element.pdf}~majorTickSpacing & \hfuzz=500pt angle & \hfuzz=500pt boundary annotation\\
\hfuzz=500pt\includegraphics[width=1em]{element.pdf}~minorTickSpacing & \hfuzz=500pt angle & \hfuzz=500pt frame tick spacing\\
\hfuzz=500pt\includegraphics[width=1em]{element.pdf}~gridLineSpacing & \hfuzz=500pt angle & \hfuzz=500pt gridline spacing\\
\hfuzz=500pt\includegraphics[width=1em]{element.pdf}~colorbar & \hfuzz=500pt \hyperref[plotColorbarType]{plotColorbar} & \hfuzz=500pt \\
\hfuzz=500pt\includegraphics[width=1em]{element-mustset.pdf}~projection & \hfuzz=500pt \hyperref[plotMapProjectionType]{plotMapProjection} & \hfuzz=500pt map projection\\
\hfuzz=500pt\includegraphics[width=1em]{element-mustset.pdf}~options & \hfuzz=500pt sequence & \hfuzz=500pt further options...\\
\hfuzz=500pt\includegraphics[width=1em]{connector.pdf}\includegraphics[width=1em]{element.pdf}~width & \hfuzz=500pt double & \hfuzz=500pt in cm\\
\hfuzz=500pt\includegraphics[width=1em]{connector.pdf}\includegraphics[width=1em]{element.pdf}~height & \hfuzz=500pt double & \hfuzz=500pt in cm\\
\hfuzz=500pt\includegraphics[width=1em]{connector.pdf}\includegraphics[width=1em]{element.pdf}~titleFontSize & \hfuzz=500pt uint & \hfuzz=500pt in pt\\
\hfuzz=500pt\includegraphics[width=1em]{connector.pdf}\includegraphics[width=1em]{element.pdf}~marginTitle & \hfuzz=500pt double & \hfuzz=500pt between title and figure [cm]\\
\hfuzz=500pt\includegraphics[width=1em]{connector.pdf}\includegraphics[width=1em]{element.pdf}~drawGridOnTop & \hfuzz=500pt boolean & \hfuzz=500pt grid lines above all other lines/points\\
\hfuzz=500pt\includegraphics[width=1em]{connector.pdf}\includegraphics[width=1em]{element-unbounded.pdf}~options & \hfuzz=500pt string & \hfuzz=500pt \\
\hfuzz=500pt\includegraphics[width=1em]{connector.pdf}\includegraphics[width=1em]{element.pdf}~transparent & \hfuzz=500pt boolean & \hfuzz=500pt make background transparent\\
\hfuzz=500pt\includegraphics[width=1em]{connector.pdf}\includegraphics[width=1em]{element.pdf}~dpi & \hfuzz=500pt uint & \hfuzz=500pt use this resolution when rasterizing postscript file\\
\hfuzz=500pt\includegraphics[width=1em]{connector.pdf}\includegraphics[width=1em]{element.pdf}~removeFiles & \hfuzz=500pt boolean & \hfuzz=500pt remove .gmt and script files\\
\hline
\end{tabularx}

\clearpage
%==================================
\subsection{PlotMatrix}\label{PlotMatrix}
Plot the coefficients of a \configFile{inputfileMatrix}{matrix}
using the GMT Generic Mapping Tools (\url{https://www.generic-mapping-tools.org}).
A variety of image file formats are supported (e.g. png, jpg, eps) determined by the extension of \config{outputfile}.

The plot programs create a temporary directory in the path of \config{outputfile}, writes all needed data into it,
generates a batch/shell script with the GMT commands, execute it, and remove the temporary directory.
With setting \config{options:removeFiles}=false the last step is skipped and it is possible to adjust the plot manually
to specific publication needs. Individual GMT settings are adjusted with \config{options:options}="\verb|FORMAT=value|",
see \url{https://docs.generic-mapping-tools.org/latest/gmt.conf.html}.

\fig{!hb}{0.6}{plotMatrix}{fig:plotMatrix}{Upper left part of the DDK filter matrix.}


\keepXColumns
\begin{tabularx}{\textwidth}{N T A}
\hline
Name & Type & Annotation\\
\hline
\hfuzz=500pt\includegraphics[width=1em]{element-mustset.pdf}~outputfile & \hfuzz=500pt filename & \hfuzz=500pt *.png, *.jpg, *.eps, ...\\
\hfuzz=500pt\includegraphics[width=1em]{element.pdf}~title & \hfuzz=500pt string & \hfuzz=500pt \\
\hfuzz=500pt\includegraphics[width=1em]{element-mustset.pdf}~inputfileMatrix & \hfuzz=500pt filename & \hfuzz=500pt \\
\hfuzz=500pt\includegraphics[width=1em]{element.pdf}~minColumn & \hfuzz=500pt uint & \hfuzz=500pt minimum column index to plot\\
\hfuzz=500pt\includegraphics[width=1em]{element.pdf}~maxColumn & \hfuzz=500pt uint & \hfuzz=500pt maximum column index to plot\\
\hfuzz=500pt\includegraphics[width=1em]{element.pdf}~majorTickSpacingX & \hfuzz=500pt double & \hfuzz=500pt boundary annotation\\
\hfuzz=500pt\includegraphics[width=1em]{element.pdf}~minorTickSpacingX & \hfuzz=500pt double & \hfuzz=500pt frame tick spacing\\
\hfuzz=500pt\includegraphics[width=1em]{element.pdf}~gridLineSpacingX & \hfuzz=500pt double & \hfuzz=500pt gridline spacing\\
\hfuzz=500pt\includegraphics[width=1em]{element.pdf}~minRow & \hfuzz=500pt uint & \hfuzz=500pt minimum row index to plot\\
\hfuzz=500pt\includegraphics[width=1em]{element.pdf}~maxRow & \hfuzz=500pt uint & \hfuzz=500pt maximum row index to plot\\
\hfuzz=500pt\includegraphics[width=1em]{element.pdf}~majorTickSpacingY & \hfuzz=500pt double & \hfuzz=500pt boundary annotation\\
\hfuzz=500pt\includegraphics[width=1em]{element.pdf}~minorTickSpacingY & \hfuzz=500pt double & \hfuzz=500pt frame tick spacing\\
\hfuzz=500pt\includegraphics[width=1em]{element.pdf}~gridLineSpacingY & \hfuzz=500pt double & \hfuzz=500pt gridline spacing\\
\hfuzz=500pt\includegraphics[width=1em]{element.pdf}~gridLine & \hfuzz=500pt \hyperref[plotLineType]{plotLine} & \hfuzz=500pt The style of the grid lines.\\
\hfuzz=500pt\includegraphics[width=1em]{element-mustset.pdf}~colorbar & \hfuzz=500pt \hyperref[plotColorbarType]{plotColorbar} & \hfuzz=500pt \\
\hfuzz=500pt\includegraphics[width=1em]{element-mustset.pdf}~options & \hfuzz=500pt sequence & \hfuzz=500pt further options...\\
\hfuzz=500pt\includegraphics[width=1em]{connector.pdf}\includegraphics[width=1em]{element.pdf}~width & \hfuzz=500pt double & \hfuzz=500pt in cm\\
\hfuzz=500pt\includegraphics[width=1em]{connector.pdf}\includegraphics[width=1em]{element.pdf}~height & \hfuzz=500pt double & \hfuzz=500pt in cm\\
\hfuzz=500pt\includegraphics[width=1em]{connector.pdf}\includegraphics[width=1em]{element.pdf}~titleFontSize & \hfuzz=500pt uint & \hfuzz=500pt in pt\\
\hfuzz=500pt\includegraphics[width=1em]{connector.pdf}\includegraphics[width=1em]{element.pdf}~marginTitle & \hfuzz=500pt double & \hfuzz=500pt between title and figure [cm]\\
\hfuzz=500pt\includegraphics[width=1em]{connector.pdf}\includegraphics[width=1em]{element.pdf}~drawGridOnTop & \hfuzz=500pt boolean & \hfuzz=500pt grid lines above all other lines/points\\
\hfuzz=500pt\includegraphics[width=1em]{connector.pdf}\includegraphics[width=1em]{element-unbounded.pdf}~options & \hfuzz=500pt string & \hfuzz=500pt \\
\hfuzz=500pt\includegraphics[width=1em]{connector.pdf}\includegraphics[width=1em]{element.pdf}~transparent & \hfuzz=500pt boolean & \hfuzz=500pt make background transparent\\
\hfuzz=500pt\includegraphics[width=1em]{connector.pdf}\includegraphics[width=1em]{element.pdf}~dpi & \hfuzz=500pt uint & \hfuzz=500pt use this resolution when rasterizing postscript file\\
\hfuzz=500pt\includegraphics[width=1em]{connector.pdf}\includegraphics[width=1em]{element.pdf}~removeFiles & \hfuzz=500pt boolean & \hfuzz=500pt remove .gmt and script files\\
\hline
\end{tabularx}

\clearpage
%==================================
\subsection{PlotSphericalHarmonicsTriangle}\label{PlotSphericalHarmonicsTriangle}
Plot the potential coefficients of a spherical harmonic expansion
using the GMT Generic Mapping Tools (\url{https://www.generic-mapping-tools.org}).
A variety of image file formats are supported (e.g. png, jpg, eps) determined by the extension of \config{outputfile}.

This program plots the formal errors (sigmas).
If \configClass{gravityfield}{gravityfieldType} provides no sigmas
e.g. with \config{setSigmasToZero} in \configClass{gravityfield:potentialCoefficients}{gravityfieldType:potentialCoefficients}
the coefficients itself are plotted instead.

The plot programs create a temporary directory in the path of \config{outputfile}, writes all needed data into it,
generates a batch/shell script with the GMT commands, execute it, and remove the temporary directory.
With setting \config{options:removeFiles}=false the last step is skipped and it is possible to adjust the plot manually
to specific publication needs. Individual GMT settings are adjusted with \config{options:options}="\verb|FORMAT=value|",
see \url{https://docs.generic-mapping-tools.org/latest/gmt.conf.html}.

\fig{!hb}{0.8}{plotSphericalHarmonicsTriangle}{fig:plotSphericalHarmonicsTriangle}{Formal errors of GOCO06s.}


\keepXColumns
\begin{tabularx}{\textwidth}{N T A}
\hline
Name & Type & Annotation\\
\hline
\hfuzz=500pt\includegraphics[width=1em]{element-mustset.pdf}~outputfile & \hfuzz=500pt filename & \hfuzz=500pt *.png, *.jpg, *.eps, ...\\
\hfuzz=500pt\includegraphics[width=1em]{element.pdf}~title & \hfuzz=500pt string & \hfuzz=500pt \\
\hfuzz=500pt\includegraphics[width=1em]{element-mustset-unbounded.pdf}~gravityfield & \hfuzz=500pt \hyperref[gravityfieldType]{gravityfield} & \hfuzz=500pt use sigmas, if not given use signal (cnm,snm)\\
\hfuzz=500pt\includegraphics[width=1em]{element.pdf}~time & \hfuzz=500pt time & \hfuzz=500pt at this time the gravity field will be evaluated\\
\hfuzz=500pt\includegraphics[width=1em]{element.pdf}~minDegree & \hfuzz=500pt uint & \hfuzz=500pt \\
\hfuzz=500pt\includegraphics[width=1em]{element.pdf}~maxDegree & \hfuzz=500pt uint & \hfuzz=500pt \\
\hfuzz=500pt\includegraphics[width=1em]{element.pdf}~majorTickSpacing & \hfuzz=500pt double & \hfuzz=500pt boundary annotation\\
\hfuzz=500pt\includegraphics[width=1em]{element.pdf}~minorTickSpacing & \hfuzz=500pt double & \hfuzz=500pt frame tick spacing\\
\hfuzz=500pt\includegraphics[width=1em]{element.pdf}~gridLineSpacing & \hfuzz=500pt double & \hfuzz=500pt gridline spacing\\
\hfuzz=500pt\includegraphics[width=1em]{element.pdf}~gridLine & \hfuzz=500pt \hyperref[plotLineType]{plotLine} & \hfuzz=500pt The style of the grid lines.\\
\hfuzz=500pt\includegraphics[width=1em]{element-mustset.pdf}~colorbar & \hfuzz=500pt \hyperref[plotColorbarType]{plotColorbar} & \hfuzz=500pt \\
\hfuzz=500pt\includegraphics[width=1em]{element-mustset.pdf}~options & \hfuzz=500pt sequence & \hfuzz=500pt further options...\\
\hfuzz=500pt\includegraphics[width=1em]{connector.pdf}\includegraphics[width=1em]{element.pdf}~width & \hfuzz=500pt double & \hfuzz=500pt in cm\\
\hfuzz=500pt\includegraphics[width=1em]{connector.pdf}\includegraphics[width=1em]{element.pdf}~height & \hfuzz=500pt double & \hfuzz=500pt in cm\\
\hfuzz=500pt\includegraphics[width=1em]{connector.pdf}\includegraphics[width=1em]{element.pdf}~titleFontSize & \hfuzz=500pt uint & \hfuzz=500pt in pt\\
\hfuzz=500pt\includegraphics[width=1em]{connector.pdf}\includegraphics[width=1em]{element.pdf}~marginTitle & \hfuzz=500pt double & \hfuzz=500pt between title and figure [cm]\\
\hfuzz=500pt\includegraphics[width=1em]{connector.pdf}\includegraphics[width=1em]{element.pdf}~drawGridOnTop & \hfuzz=500pt boolean & \hfuzz=500pt grid lines above all other lines/points\\
\hfuzz=500pt\includegraphics[width=1em]{connector.pdf}\includegraphics[width=1em]{element-unbounded.pdf}~options & \hfuzz=500pt string & \hfuzz=500pt \\
\hfuzz=500pt\includegraphics[width=1em]{connector.pdf}\includegraphics[width=1em]{element.pdf}~transparent & \hfuzz=500pt boolean & \hfuzz=500pt make background transparent\\
\hfuzz=500pt\includegraphics[width=1em]{connector.pdf}\includegraphics[width=1em]{element.pdf}~dpi & \hfuzz=500pt uint & \hfuzz=500pt use this resolution when rasterizing postscript file\\
\hfuzz=500pt\includegraphics[width=1em]{connector.pdf}\includegraphics[width=1em]{element.pdf}~removeFiles & \hfuzz=500pt boolean & \hfuzz=500pt remove .gmt and script files\\
\hline
\end{tabularx}

\clearpage
%==================================
\section{Programs: Preprocessing}
\subsection{PreprocessingDualSst}\label{PreprocessingDualSst}
This programs processes satellite-to-satellite-tracking (SST) and orbit observations in a GRACE like configuration.
Four different observation groups are considered separately: two types of SST and POD1/POD2 for the two satellites.
This program works similar to \program{PreprocessingSst}, see there for details. Here only the settings explained,
which are different.

Both SST observation types are reduced by the same background models and the same impact
of accelerometer measurements. The covariance matrix of the reduced observations should not consider
the the instrument noise only (\configClass{covarianceSst1/2}{covarianceSstType}) but must
take the cross correlations \configClass{covarianceAcc}{covarianceSstType} into account.
The covariance matrix of the reduced observations is given by
\begin{equation}
  \M\Sigma(\begin{bmatrix} \Delta l_{SST1} \\ \Delta l_{SST2} \end{bmatrix})
  = \begin{bmatrix} \M\Sigma_{SST1} + \M\Sigma_{ACC} & \M\Sigma_{ACC} \\
                   \M\Sigma_{ACC} & \M\Sigma_{SST2} + \M\Sigma_{ACC}
    \end{bmatrix}.
\end{equation}


\keepXColumns
\begin{tabularx}{\textwidth}{N T A}
\hline
Name & Type & Annotation\\
\hline
\hfuzz=500pt\includegraphics[width=1em]{element.pdf}~outputfileSolution & \hfuzz=500pt filename & \hfuzz=500pt estimated parameter vector (static part only)\\
\hfuzz=500pt\includegraphics[width=1em]{element.pdf}~outputfileSigmax & \hfuzz=500pt filename & \hfuzz=500pt standard deviations of the parameters (sqrt of the diagonal of the inverse normal equation)\\
\hfuzz=500pt\includegraphics[width=1em]{element.pdf}~outputfileParameterName & \hfuzz=500pt filename & \hfuzz=500pt estimated signal parameters (index is appended)\\
\hfuzz=500pt\includegraphics[width=1em]{element.pdf}~estimateArcSigmas & \hfuzz=500pt sequence & \hfuzz=500pt \\
\hfuzz=500pt\includegraphics[width=1em]{connector.pdf}\includegraphics[width=1em]{element.pdf}~outputfileSigmasPerArcSst1 & \hfuzz=500pt filename & \hfuzz=500pt accuracies of each arc (SST1)\\
\hfuzz=500pt\includegraphics[width=1em]{connector.pdf}\includegraphics[width=1em]{element.pdf}~outputfileSigmasPerArcSst2 & \hfuzz=500pt filename & \hfuzz=500pt accuracies of each arc (SST2)\\
\hfuzz=500pt\includegraphics[width=1em]{connector.pdf}\includegraphics[width=1em]{element.pdf}~outputfileSigmasPerArcAcc & \hfuzz=500pt filename & \hfuzz=500pt accuracies of each arc (ACC)\\
\hfuzz=500pt\includegraphics[width=1em]{connector.pdf}\includegraphics[width=1em]{element.pdf}~outputfileSigmasPerArcPod1 & \hfuzz=500pt filename & \hfuzz=500pt accuracies of each arc (POD1)\\
\hfuzz=500pt\includegraphics[width=1em]{connector.pdf}\includegraphics[width=1em]{element.pdf}~outputfileSigmasPerArcPod2 & \hfuzz=500pt filename & \hfuzz=500pt accuracies of each arc (POD2)\\
\hfuzz=500pt\includegraphics[width=1em]{element.pdf}~estimateEpochSigmas & \hfuzz=500pt sequence & \hfuzz=500pt \\
\hfuzz=500pt\includegraphics[width=1em]{connector.pdf}\includegraphics[width=1em]{element.pdf}~outputfileSigmasPerEpochSst1 & \hfuzz=500pt filename & \hfuzz=500pt accuracies of each epoch (SST1)\\
\hfuzz=500pt\includegraphics[width=1em]{connector.pdf}\includegraphics[width=1em]{element.pdf}~outputfileSigmasPerEpochSst2 & \hfuzz=500pt filename & \hfuzz=500pt accuracies of each epoch (SST2)\\
\hfuzz=500pt\includegraphics[width=1em]{connector.pdf}\includegraphics[width=1em]{element.pdf}~outputfileSigmasPerEpochAcc & \hfuzz=500pt filename & \hfuzz=500pt accuracies of each epoch (ACC)\\
\hfuzz=500pt\includegraphics[width=1em]{connector.pdf}\includegraphics[width=1em]{element.pdf}~outputfileSigmasPerEpochPod1 & \hfuzz=500pt filename & \hfuzz=500pt accuracies of each epoch (POD1)\\
\hfuzz=500pt\includegraphics[width=1em]{connector.pdf}\includegraphics[width=1em]{element.pdf}~outputfileSigmasPerEpochPod2 & \hfuzz=500pt filename & \hfuzz=500pt accuracies of each epoch (POD2)\\
\hfuzz=500pt\includegraphics[width=1em]{element.pdf}~estimateCovarianceFunctions & \hfuzz=500pt sequence & \hfuzz=500pt \\
\hfuzz=500pt\includegraphics[width=1em]{connector.pdf}\includegraphics[width=1em]{element.pdf}~outputfileCovarianceFunctionSst1 & \hfuzz=500pt filename & \hfuzz=500pt covariance function\\
\hfuzz=500pt\includegraphics[width=1em]{connector.pdf}\includegraphics[width=1em]{element.pdf}~outputfileCovarianceFunctionSst2 & \hfuzz=500pt filename & \hfuzz=500pt covariance function\\
\hfuzz=500pt\includegraphics[width=1em]{connector.pdf}\includegraphics[width=1em]{element.pdf}~outputfileCovarianceFunctionAcc & \hfuzz=500pt filename & \hfuzz=500pt covariance function\\
\hfuzz=500pt\includegraphics[width=1em]{connector.pdf}\includegraphics[width=1em]{element.pdf}~outputfileCovarianceFunctionPod1 & \hfuzz=500pt filename & \hfuzz=500pt covariance functions for along, cross, radial direction\\
\hfuzz=500pt\includegraphics[width=1em]{connector.pdf}\includegraphics[width=1em]{element.pdf}~outputfileCovarianceFunctionPod2 & \hfuzz=500pt filename & \hfuzz=500pt covariance functions for along, cross, radial direction\\
\hfuzz=500pt\includegraphics[width=1em]{element.pdf}~computeResiduals & \hfuzz=500pt sequence & \hfuzz=500pt \\
\hfuzz=500pt\includegraphics[width=1em]{connector.pdf}\includegraphics[width=1em]{element.pdf}~outputfileSst1Residuals & \hfuzz=500pt filename & \hfuzz=500pt \\
\hfuzz=500pt\includegraphics[width=1em]{connector.pdf}\includegraphics[width=1em]{element.pdf}~outputfileSst2Residuals & \hfuzz=500pt filename & \hfuzz=500pt \\
\hfuzz=500pt\includegraphics[width=1em]{connector.pdf}\includegraphics[width=1em]{element.pdf}~outputfileAccResiduals & \hfuzz=500pt filename & \hfuzz=500pt \\
\hfuzz=500pt\includegraphics[width=1em]{connector.pdf}\includegraphics[width=1em]{element.pdf}~outputfilePod1Residuals & \hfuzz=500pt filename & \hfuzz=500pt \\
\hfuzz=500pt\includegraphics[width=1em]{connector.pdf}\includegraphics[width=1em]{element.pdf}~outputfilePod2Residuals & \hfuzz=500pt filename & \hfuzz=500pt \\
\hfuzz=500pt\includegraphics[width=1em]{element-mustset.pdf}~observation & \hfuzz=500pt choice & \hfuzz=500pt obervation equations (Sst)\\
\hfuzz=500pt\includegraphics[width=1em]{connector.pdf}\includegraphics[width=1em]{element-mustset.pdf}~dualSstVariational & \hfuzz=500pt sequence & \hfuzz=500pt two SST observations\\
\hfuzz=500pt\quad\includegraphics[width=1em]{connector.pdf}\includegraphics[width=1em]{element-mustset.pdf}~rightHandSide & \hfuzz=500pt sequence & \hfuzz=500pt input for observation vectors\\
\hfuzz=500pt\quad\quad\includegraphics[width=1em]{connector.pdf}\includegraphics[width=1em]{element-mustset-unbounded.pdf}~inputfileSatelliteTracking1 & \hfuzz=500pt filename & \hfuzz=500pt ranging observations and corrections\\
\hfuzz=500pt\quad\quad\includegraphics[width=1em]{connector.pdf}\includegraphics[width=1em]{element-mustset-unbounded.pdf}~inputfileSatelliteTracking2 & \hfuzz=500pt filename & \hfuzz=500pt ranging observations and corrections\\
\hfuzz=500pt\quad\quad\includegraphics[width=1em]{connector.pdf}\includegraphics[width=1em]{element.pdf}~inputfileOrbit1 & \hfuzz=500pt filename & \hfuzz=500pt kinematic positions of satellite A as observations\\
\hfuzz=500pt\quad\quad\includegraphics[width=1em]{connector.pdf}\includegraphics[width=1em]{element.pdf}~inputfileOrbit2 & \hfuzz=500pt filename & \hfuzz=500pt kinematic positions of satellite B as observations\\
\hfuzz=500pt\quad\includegraphics[width=1em]{connector.pdf}\includegraphics[width=1em]{element-mustset.pdf}~sstType & \hfuzz=500pt choice & \hfuzz=500pt \\
\hfuzz=500pt\quad\quad\includegraphics[width=1em]{connector.pdf}\includegraphics[width=1em]{element-mustset.pdf}~range & \hfuzz=500pt  & \hfuzz=500pt \\
\hfuzz=500pt\quad\quad\includegraphics[width=1em]{connector.pdf}\includegraphics[width=1em]{element-mustset.pdf}~rangeRate & \hfuzz=500pt  & \hfuzz=500pt \\
\hfuzz=500pt\quad\quad\includegraphics[width=1em]{connector.pdf}\includegraphics[width=1em]{element-mustset.pdf}~none & \hfuzz=500pt  & \hfuzz=500pt \\
\hfuzz=500pt\quad\includegraphics[width=1em]{connector.pdf}\includegraphics[width=1em]{element-mustset.pdf}~inputfileVariational1 & \hfuzz=500pt filename & \hfuzz=500pt approximate position and integrated state matrix\\
\hfuzz=500pt\quad\includegraphics[width=1em]{connector.pdf}\includegraphics[width=1em]{element-mustset.pdf}~inputfileVariational2 & \hfuzz=500pt filename & \hfuzz=500pt approximate position and integrated state matrix\\
\hfuzz=500pt\quad\includegraphics[width=1em]{connector.pdf}\includegraphics[width=1em]{element.pdf}~ephemerides & \hfuzz=500pt \hyperref[ephemeridesType]{ephemerides} & \hfuzz=500pt \\
\hfuzz=500pt\quad\includegraphics[width=1em]{connector.pdf}\includegraphics[width=1em]{element-unbounded.pdf}~parametrizationGravity & \hfuzz=500pt \hyperref[parametrizationGravityType]{parametrizationGravity} & \hfuzz=500pt gravity field parametrization\\
\hfuzz=500pt\quad\includegraphics[width=1em]{connector.pdf}\includegraphics[width=1em]{element-unbounded.pdf}~parametrizationAcceleration1 & \hfuzz=500pt \hyperref[parametrizationAccelerationType]{parametrizationAcceleration} & \hfuzz=500pt orbit1 force parameters\\
\hfuzz=500pt\quad\includegraphics[width=1em]{connector.pdf}\includegraphics[width=1em]{element-unbounded.pdf}~parametrizationAcceleration2 & \hfuzz=500pt \hyperref[parametrizationAccelerationType]{parametrizationAcceleration} & \hfuzz=500pt orbit2 force parameters\\
\hfuzz=500pt\quad\includegraphics[width=1em]{connector.pdf}\includegraphics[width=1em]{element-unbounded.pdf}~parametrizationSst1 & \hfuzz=500pt \hyperref[parametrizationSatelliteTrackingType]{parametrizationSatelliteTracking} & \hfuzz=500pt satellite tracking parameter for first ranging observations\\
\hfuzz=500pt\quad\includegraphics[width=1em]{connector.pdf}\includegraphics[width=1em]{element-unbounded.pdf}~parametrizationSst2 & \hfuzz=500pt \hyperref[parametrizationSatelliteTrackingType]{parametrizationSatelliteTracking} & \hfuzz=500pt satellite tracking parameter for second ranging observations\\
\hfuzz=500pt\quad\includegraphics[width=1em]{connector.pdf}\includegraphics[width=1em]{element.pdf}~integrationDegree & \hfuzz=500pt uint & \hfuzz=500pt integration of forces by polynomial approximation of degree n\\
\hfuzz=500pt\quad\includegraphics[width=1em]{connector.pdf}\includegraphics[width=1em]{element.pdf}~interpolationDegree & \hfuzz=500pt uint & \hfuzz=500pt orbit interpolation by polynomial approximation of degree n\\
\hfuzz=500pt\includegraphics[width=1em]{element-mustset.pdf}~covarianceSst1 & \hfuzz=500pt sequence & \hfuzz=500pt \\
\hfuzz=500pt\includegraphics[width=1em]{connector.pdf}\includegraphics[width=1em]{element.pdf}~sigma & \hfuzz=500pt double & \hfuzz=500pt apriori factor of covariance function\\
\hfuzz=500pt\includegraphics[width=1em]{connector.pdf}\includegraphics[width=1em]{element.pdf}~inputfileSigmasPerArc & \hfuzz=500pt filename & \hfuzz=500pt apriori different accuaries for each arc (multiplicated with sigma)\\
\hfuzz=500pt\includegraphics[width=1em]{connector.pdf}\includegraphics[width=1em]{element.pdf}~inputfileSigmasPerEpoch & \hfuzz=500pt filename & \hfuzz=500pt apriori different accuaries for each epoch\\
\hfuzz=500pt\includegraphics[width=1em]{connector.pdf}\includegraphics[width=1em]{element.pdf}~inputfileCovarianceFunction & \hfuzz=500pt filename & \hfuzz=500pt approximate covariances in time\\
\hfuzz=500pt\includegraphics[width=1em]{connector.pdf}\includegraphics[width=1em]{element-unbounded.pdf}~inputfileCovarianceMatrixArc & \hfuzz=500pt filename & \hfuzz=500pt Must be given per sst arc with correct dimensions.\\
\hfuzz=500pt\includegraphics[width=1em]{connector.pdf}\includegraphics[width=1em]{element.pdf}~inputfileSigmasCovarianceMatrixArc & \hfuzz=500pt filename & \hfuzz=500pt Vector with one sigma for each \$<\$inputfileCovarianceMatrixArc\$>\$\\
\hfuzz=500pt\includegraphics[width=1em]{connector.pdf}\includegraphics[width=1em]{element.pdf}~sampling & \hfuzz=500pt double & \hfuzz=500pt [seconds] sampling of the covariance function\\
\hfuzz=500pt\includegraphics[width=1em]{element-mustset.pdf}~covarianceSst2 & \hfuzz=500pt sequence & \hfuzz=500pt \\
\hfuzz=500pt\includegraphics[width=1em]{connector.pdf}\includegraphics[width=1em]{element.pdf}~sigma & \hfuzz=500pt double & \hfuzz=500pt apriori factor of covariance function\\
\hfuzz=500pt\includegraphics[width=1em]{connector.pdf}\includegraphics[width=1em]{element.pdf}~inputfileSigmasPerArc & \hfuzz=500pt filename & \hfuzz=500pt apriori different accuaries for each arc (multiplicated with sigma)\\
\hfuzz=500pt\includegraphics[width=1em]{connector.pdf}\includegraphics[width=1em]{element.pdf}~inputfileSigmasPerEpoch & \hfuzz=500pt filename & \hfuzz=500pt apriori different accuaries for each epoch\\
\hfuzz=500pt\includegraphics[width=1em]{connector.pdf}\includegraphics[width=1em]{element.pdf}~inputfileCovarianceFunction & \hfuzz=500pt filename & \hfuzz=500pt approximate covariances in time\\
\hfuzz=500pt\includegraphics[width=1em]{connector.pdf}\includegraphics[width=1em]{element-unbounded.pdf}~inputfileCovarianceMatrixArc & \hfuzz=500pt filename & \hfuzz=500pt Must be given per sst arc with correct dimensions.\\
\hfuzz=500pt\includegraphics[width=1em]{connector.pdf}\includegraphics[width=1em]{element.pdf}~inputfileSigmasCovarianceMatrixArc & \hfuzz=500pt filename & \hfuzz=500pt Vector with one sigma for each \$<\$inputfileCovarianceMatrixArc\$>\$\\
\hfuzz=500pt\includegraphics[width=1em]{connector.pdf}\includegraphics[width=1em]{element.pdf}~sampling & \hfuzz=500pt double & \hfuzz=500pt [seconds] sampling of the covariance function\\
\hfuzz=500pt\includegraphics[width=1em]{element-mustset.pdf}~covarianceAcc & \hfuzz=500pt sequence & \hfuzz=500pt \\
\hfuzz=500pt\includegraphics[width=1em]{connector.pdf}\includegraphics[width=1em]{element.pdf}~sigma & \hfuzz=500pt double & \hfuzz=500pt apriori factor of covariance function\\
\hfuzz=500pt\includegraphics[width=1em]{connector.pdf}\includegraphics[width=1em]{element.pdf}~inputfileSigmasPerArc & \hfuzz=500pt filename & \hfuzz=500pt apriori different accuaries for each arc (multiplicated with sigma)\\
\hfuzz=500pt\includegraphics[width=1em]{connector.pdf}\includegraphics[width=1em]{element.pdf}~inputfileSigmasPerEpoch & \hfuzz=500pt filename & \hfuzz=500pt apriori different accuaries for each epoch\\
\hfuzz=500pt\includegraphics[width=1em]{connector.pdf}\includegraphics[width=1em]{element.pdf}~inputfileCovarianceFunction & \hfuzz=500pt filename & \hfuzz=500pt approximate covariances in time\\
\hfuzz=500pt\includegraphics[width=1em]{connector.pdf}\includegraphics[width=1em]{element-unbounded.pdf}~inputfileCovarianceMatrixArc & \hfuzz=500pt filename & \hfuzz=500pt Must be given per sst arc with correct dimensions.\\
\hfuzz=500pt\includegraphics[width=1em]{connector.pdf}\includegraphics[width=1em]{element.pdf}~inputfileSigmasCovarianceMatrixArc & \hfuzz=500pt filename & \hfuzz=500pt Vector with one sigma for each \$<\$inputfileCovarianceMatrixArc\$>\$\\
\hfuzz=500pt\includegraphics[width=1em]{connector.pdf}\includegraphics[width=1em]{element.pdf}~sampling & \hfuzz=500pt double & \hfuzz=500pt [seconds] sampling of the covariance function\\
\hfuzz=500pt\includegraphics[width=1em]{element-mustset.pdf}~covariancePod1 & \hfuzz=500pt sequence & \hfuzz=500pt \\
\hfuzz=500pt\includegraphics[width=1em]{connector.pdf}\includegraphics[width=1em]{element.pdf}~sigma & \hfuzz=500pt double & \hfuzz=500pt apriori factor of covariance function\\
\hfuzz=500pt\includegraphics[width=1em]{connector.pdf}\includegraphics[width=1em]{element.pdf}~inputfileSigmasPerArc & \hfuzz=500pt filename & \hfuzz=500pt apriori different accuaries for each arc (multiplicated with sigma)\\
\hfuzz=500pt\includegraphics[width=1em]{connector.pdf}\includegraphics[width=1em]{element.pdf}~inputfileSigmasPerEpoch & \hfuzz=500pt filename & \hfuzz=500pt apriori different accuaries for each epoch\\
\hfuzz=500pt\includegraphics[width=1em]{connector.pdf}\includegraphics[width=1em]{element.pdf}~inputfileCovarianceFunction & \hfuzz=500pt filename & \hfuzz=500pt approximate covariances in time\\
\hfuzz=500pt\includegraphics[width=1em]{connector.pdf}\includegraphics[width=1em]{element.pdf}~inputfileCovariancePodEpoch & \hfuzz=500pt filename & \hfuzz=500pt 3x3 epoch covariances\\
\hfuzz=500pt\includegraphics[width=1em]{connector.pdf}\includegraphics[width=1em]{element.pdf}~sampling & \hfuzz=500pt double & \hfuzz=500pt [seconds] sampling of the covariance function\\
\hfuzz=500pt\includegraphics[width=1em]{element-mustset.pdf}~covariancePod2 & \hfuzz=500pt sequence & \hfuzz=500pt \\
\hfuzz=500pt\includegraphics[width=1em]{connector.pdf}\includegraphics[width=1em]{element.pdf}~sigma & \hfuzz=500pt double & \hfuzz=500pt apriori factor of covariance function\\
\hfuzz=500pt\includegraphics[width=1em]{connector.pdf}\includegraphics[width=1em]{element.pdf}~inputfileSigmasPerArc & \hfuzz=500pt filename & \hfuzz=500pt apriori different accuaries for each arc (multiplicated with sigma)\\
\hfuzz=500pt\includegraphics[width=1em]{connector.pdf}\includegraphics[width=1em]{element.pdf}~inputfileSigmasPerEpoch & \hfuzz=500pt filename & \hfuzz=500pt apriori different accuaries for each epoch\\
\hfuzz=500pt\includegraphics[width=1em]{connector.pdf}\includegraphics[width=1em]{element.pdf}~inputfileCovarianceFunction & \hfuzz=500pt filename & \hfuzz=500pt approximate covariances in time\\
\hfuzz=500pt\includegraphics[width=1em]{connector.pdf}\includegraphics[width=1em]{element.pdf}~inputfileCovariancePodEpoch & \hfuzz=500pt filename & \hfuzz=500pt 3x3 epoch covariances\\
\hfuzz=500pt\includegraphics[width=1em]{connector.pdf}\includegraphics[width=1em]{element.pdf}~sampling & \hfuzz=500pt double & \hfuzz=500pt [seconds] sampling of the covariance function\\
\hfuzz=500pt\includegraphics[width=1em]{element.pdf}~estimateShortTimeVariations & \hfuzz=500pt sequence & \hfuzz=500pt co-estimate short time gravity field variations\\
\hfuzz=500pt\includegraphics[width=1em]{connector.pdf}\includegraphics[width=1em]{element.pdf}~estimateSigma & \hfuzz=500pt boolean & \hfuzz=500pt estimate standard deviation via VCE\\
\hfuzz=500pt\includegraphics[width=1em]{connector.pdf}\includegraphics[width=1em]{element-mustset.pdf}~autoregressiveModelSequence & \hfuzz=500pt \hyperref[autoregressiveModelSequenceType]{autoregressiveModelSequence} & \hfuzz=500pt AR model sequence for constraining short time gravity variations\\
\hfuzz=500pt\includegraphics[width=1em]{connector.pdf}\includegraphics[width=1em]{element-mustset-unbounded.pdf}~parameterSelection & \hfuzz=500pt \hyperref[parameterSelectorType]{parameterSelector} & \hfuzz=500pt parameters describing the short time gravity field\\
\hfuzz=500pt\includegraphics[width=1em]{element.pdf}~downweightPod & \hfuzz=500pt double & \hfuzz=500pt downweight factor for POD\\
\hfuzz=500pt\includegraphics[width=1em]{element.pdf}~inputfileArcList & \hfuzz=500pt filename & \hfuzz=500pt list to correspond points of time to arc numbers\\
\hfuzz=500pt\includegraphics[width=1em]{element.pdf}~iterationCount & \hfuzz=500pt uint & \hfuzz=500pt (maximum) number of iterations for the estimation of calibration parameter and error PSD\\
\hfuzz=500pt\includegraphics[width=1em]{element.pdf}~variableNameIterations & \hfuzz=500pt string & \hfuzz=500pt All output fileNames in preprocessing iteration are expanded with this variable prior to writing to disk\\
\hfuzz=500pt\includegraphics[width=1em]{element.pdf}~defaultBlockSize & \hfuzz=500pt uint & \hfuzz=500pt block size of static normal equation blocks\\
\hline
\end{tabularx}

This program is \reference{parallelized}{general.parallelization}.
\clearpage
%==================================
\subsection{PreprocessingGradiometer}\label{PreprocessingGradiometer}
This program estimates empirical covariance functions of the gradiometer instrument noise and determine arc wise variances to
downweight arcs with outliers. This program works similar to \program{PreprocessingPod}, see there for details.
Here only the settings explained, which are different.

...


\keepXColumns
\begin{tabularx}{\textwidth}{N T A}
\hline
Name & Type & Annotation\\
\hline
\hfuzz=500pt\includegraphics[width=1em]{element.pdf}~outputfileCovarianceFunction & \hfuzz=500pt filename & \hfuzz=500pt \\
\hfuzz=500pt\includegraphics[width=1em]{element.pdf}~outputfileSigmasPerArc & \hfuzz=500pt filename & \hfuzz=500pt accuracies of each arc\\
\hfuzz=500pt\includegraphics[width=1em]{element.pdf}~outputfileSggResiduals & \hfuzz=500pt filename & \hfuzz=500pt \\
\hfuzz=500pt\includegraphics[width=1em]{element-mustset.pdf}~rightHandSide & \hfuzz=500pt \hyperref[sggRightSideType]{sggRightSide} & \hfuzz=500pt input for the observation vector\\
\hfuzz=500pt\includegraphics[width=1em]{element-mustset.pdf}~inputfileOrbit & \hfuzz=500pt filename & \hfuzz=500pt \\
\hfuzz=500pt\includegraphics[width=1em]{element-mustset.pdf}~inputfileStarCamera & \hfuzz=500pt filename & \hfuzz=500pt \\
\hfuzz=500pt\includegraphics[width=1em]{element-mustset.pdf}~earthRotation & \hfuzz=500pt \hyperref[earthRotationType]{earthRotation} & \hfuzz=500pt \\
\hfuzz=500pt\includegraphics[width=1em]{element.pdf}~ephemerides & \hfuzz=500pt \hyperref[ephemeridesType]{ephemerides} & \hfuzz=500pt \\
\hfuzz=500pt\includegraphics[width=1em]{element-unbounded.pdf}~parametrizationBias & \hfuzz=500pt \hyperref[parametrizationTemporalType]{parametrizationTemporal} & \hfuzz=500pt per arc\\
\hfuzz=500pt\includegraphics[width=1em]{element-mustset.pdf}~covarianceSgg & \hfuzz=500pt sequence & \hfuzz=500pt \\
\hfuzz=500pt\includegraphics[width=1em]{connector.pdf}\includegraphics[width=1em]{element.pdf}~inputfileCovarianceFunction & \hfuzz=500pt filename & \hfuzz=500pt approximate covariances in time\\
\hfuzz=500pt\includegraphics[width=1em]{connector.pdf}\includegraphics[width=1em]{element-mustset.pdf}~covarianceLength & \hfuzz=500pt uint & \hfuzz=500pt counts observation epochs\\
\hfuzz=500pt\includegraphics[width=1em]{connector.pdf}\includegraphics[width=1em]{element-mustset.pdf}~sampling & \hfuzz=500pt double & \hfuzz=500pt [seconds] sampling of the covariance function\\
\hfuzz=500pt\includegraphics[width=1em]{element.pdf}~iterationCount & \hfuzz=500pt uint & \hfuzz=500pt for the estimation of calibration parameter and error PSD\\
\hline
\end{tabularx}

This program is \reference{parallelized}{general.parallelization}.
\clearpage
%==================================
\subsection{PreprocessingPod}\label{PreprocessingPod}
This program estimates empirical covariance functions of the instrument noise and determine arc wise variances to
downweight arc with outliers.

A complete least squares adjustment for gravity field determination is performed by computing the \config{observation}
equations, see \configClass{observation:podIntegral}{observationType:podIntegral} or
\configClass{observation:podVariational}{observationType:podVariational} for details. The normal equations
are accumulated and solved to \configFile{outputfileSolution}{matrix} together with the estimated accuracies
\configFile{outputfileSigmax}{matrix}. The estimated residuals~$\hat{\M e}=\M l-\M A\hat{\M x}$ can be computed with
\config{computeResiduals}.

For each component (along, cross, radial) of the kinematic orbit positions a noise covariance function is estimated
\begin{equation}
  \text{cov}(\Delta t_i) = \sum_{n=0}^{N-1} a_n^2 \cos\left(\frac{\pi}{T} n\Delta t_i\right).
\end{equation}
The covariance matrix is composed of the sum of matrices $F_n$ and unknown variance factors
\begin{equation}
  \M\Sigma = a_1^2\M F_1 + a_2^2 \M F_2 + \cdots + a_N^2\M F_N,
\end{equation}
with the cosine transformation matrices
\begin{equation}
  \M F_n = \left(\cos\left(\frac{\pi}{T} n(t_i-t_k)\right)\right)_{ik}.
\end{equation}

An additional variance factor can be computed (\config{estimateArcSigmas}) for each arc~$k$  according to
\begin{equation}
  \hat{\sigma}_k^2 = \frac{\hat{\M e}_k^T\M\Sigma^{-1}\hat{\M e}_k}{r_k},
\end{equation}
where $r_k$ is the redundancy. This variance factor should be around one for normal behaving arcs
as the noise characteristics is already considered by the covariance matrix but bad arcs get a much larger variance.
By appling this factor bad arcs or arcs with large outliers are downweighted.


\keepXColumns
\begin{tabularx}{\textwidth}{N T A}
\hline
Name & Type & Annotation\\
\hline
\hfuzz=500pt\includegraphics[width=1em]{element.pdf}~outputfileSolution & \hfuzz=500pt filename & \hfuzz=500pt estimated parameter vector (static part only)\\
\hfuzz=500pt\includegraphics[width=1em]{element.pdf}~outputfileSigmax & \hfuzz=500pt filename & \hfuzz=500pt standard deviations of the parameters (sqrt of the diagonal of the inverse normal equation)\\
\hfuzz=500pt\includegraphics[width=1em]{element.pdf}~outputfileParameterName & \hfuzz=500pt filename & \hfuzz=500pt names of estimated parameters (static part only)\\
\hfuzz=500pt\includegraphics[width=1em]{element.pdf}~estimateArcSigmas & \hfuzz=500pt sequence & \hfuzz=500pt \\
\hfuzz=500pt\includegraphics[width=1em]{connector.pdf}\includegraphics[width=1em]{element.pdf}~outputfileSigmasPerArcPod & \hfuzz=500pt filename & \hfuzz=500pt accuracies of each arc (POD2)\\
\hfuzz=500pt\includegraphics[width=1em]{element.pdf}~estimateCovarianceFunctions & \hfuzz=500pt sequence & \hfuzz=500pt \\
\hfuzz=500pt\includegraphics[width=1em]{connector.pdf}\includegraphics[width=1em]{element.pdf}~outputfileCovarianceFunctionPod & \hfuzz=500pt filename & \hfuzz=500pt covariance functions for along, cross, radial direction\\
\hfuzz=500pt\includegraphics[width=1em]{element.pdf}~computeResiduals & \hfuzz=500pt sequence & \hfuzz=500pt \\
\hfuzz=500pt\includegraphics[width=1em]{connector.pdf}\includegraphics[width=1em]{element.pdf}~outputfilePodResiduals & \hfuzz=500pt filename & \hfuzz=500pt \\
\hfuzz=500pt\includegraphics[width=1em]{element-mustset.pdf}~observation & \hfuzz=500pt choice & \hfuzz=500pt obervation equations (POD)\\
\hfuzz=500pt\includegraphics[width=1em]{connector.pdf}\includegraphics[width=1em]{element-mustset.pdf}~podIntegral & \hfuzz=500pt sequence & \hfuzz=500pt precise orbit data (integral approach)\\
\hfuzz=500pt\quad\includegraphics[width=1em]{connector.pdf}\includegraphics[width=1em]{element.pdf}~inputfileSatelliteModel & \hfuzz=500pt filename & \hfuzz=500pt satellite macro model\\
\hfuzz=500pt\quad\includegraphics[width=1em]{connector.pdf}\includegraphics[width=1em]{element-mustset-unbounded.pdf}~rightHandSide & \hfuzz=500pt \hyperref[podRightSideType]{podRightSide} & \hfuzz=500pt input for the reduced observation vector\\
\hfuzz=500pt\quad\includegraphics[width=1em]{connector.pdf}\includegraphics[width=1em]{element-mustset.pdf}~inputfileOrbit & \hfuzz=500pt filename & \hfuzz=500pt used to evaluate the observation equations, not used as observations\\
\hfuzz=500pt\quad\includegraphics[width=1em]{connector.pdf}\includegraphics[width=1em]{element-mustset.pdf}~inputfileStarCamera & \hfuzz=500pt filename & \hfuzz=500pt \\
\hfuzz=500pt\quad\includegraphics[width=1em]{connector.pdf}\includegraphics[width=1em]{element-mustset.pdf}~earthRotation & \hfuzz=500pt \hyperref[earthRotationType]{earthRotation} & \hfuzz=500pt \\
\hfuzz=500pt\quad\includegraphics[width=1em]{connector.pdf}\includegraphics[width=1em]{element.pdf}~ephemerides & \hfuzz=500pt \hyperref[ephemeridesType]{ephemerides} & \hfuzz=500pt \\
\hfuzz=500pt\quad\includegraphics[width=1em]{connector.pdf}\includegraphics[width=1em]{element-unbounded.pdf}~gradientfield & \hfuzz=500pt \hyperref[gravityfieldType]{gravityfield} & \hfuzz=500pt low order field to estimate the change of the gravity by position adjustement\\
\hfuzz=500pt\quad\includegraphics[width=1em]{connector.pdf}\includegraphics[width=1em]{element-unbounded.pdf}~parametrizationGravity & \hfuzz=500pt \hyperref[parametrizationGravityType]{parametrizationGravity} & \hfuzz=500pt gravity field parametrization\\
\hfuzz=500pt\quad\includegraphics[width=1em]{connector.pdf}\includegraphics[width=1em]{element-unbounded.pdf}~parametrizationAcceleration & \hfuzz=500pt \hyperref[parametrizationAccelerationType]{parametrizationAcceleration} & \hfuzz=500pt orbit force parameters\\
\hfuzz=500pt\quad\includegraphics[width=1em]{connector.pdf}\includegraphics[width=1em]{element.pdf}~keepSatelliteStates & \hfuzz=500pt boolean & \hfuzz=500pt set boundary values of each arc global\\
\hfuzz=500pt\quad\includegraphics[width=1em]{connector.pdf}\includegraphics[width=1em]{element.pdf}~integrationDegree & \hfuzz=500pt uint & \hfuzz=500pt integration of forces by polynomial approximation of degree n\\
\hfuzz=500pt\quad\includegraphics[width=1em]{connector.pdf}\includegraphics[width=1em]{element.pdf}~interpolationDegree & \hfuzz=500pt uint & \hfuzz=500pt orbit interpolation by polynomial approximation of degree n\\
\hfuzz=500pt\quad\includegraphics[width=1em]{connector.pdf}\includegraphics[width=1em]{element.pdf}~accelerateComputation & \hfuzz=500pt boolean & \hfuzz=500pt acceleration of computation by transforming the observations\\
\hfuzz=500pt\includegraphics[width=1em]{connector.pdf}\includegraphics[width=1em]{element-mustset.pdf}~podVariational & \hfuzz=500pt sequence & \hfuzz=500pt precise orbit data (variational equations)\\
\hfuzz=500pt\quad\includegraphics[width=1em]{connector.pdf}\includegraphics[width=1em]{element-mustset.pdf}~rightHandSide & \hfuzz=500pt sequence & \hfuzz=500pt input for observation vectors\\
\hfuzz=500pt\quad\quad\includegraphics[width=1em]{connector.pdf}\includegraphics[width=1em]{element-mustset.pdf}~inputfileOrbit & \hfuzz=500pt filename & \hfuzz=500pt kinematic positions as observations\\
\hfuzz=500pt\quad\includegraphics[width=1em]{connector.pdf}\includegraphics[width=1em]{element-mustset.pdf}~inputfileVariational & \hfuzz=500pt filename & \hfuzz=500pt approximate position and integrated state matrix\\
\hfuzz=500pt\quad\includegraphics[width=1em]{connector.pdf}\includegraphics[width=1em]{element.pdf}~ephemerides & \hfuzz=500pt \hyperref[ephemeridesType]{ephemerides} & \hfuzz=500pt \\
\hfuzz=500pt\quad\includegraphics[width=1em]{connector.pdf}\includegraphics[width=1em]{element-unbounded.pdf}~parametrizationGravity & \hfuzz=500pt \hyperref[parametrizationGravityType]{parametrizationGravity} & \hfuzz=500pt gravity field parametrization\\
\hfuzz=500pt\quad\includegraphics[width=1em]{connector.pdf}\includegraphics[width=1em]{element-unbounded.pdf}~parametrizationAcceleration & \hfuzz=500pt \hyperref[parametrizationAccelerationType]{parametrizationAcceleration} & \hfuzz=500pt orbit force parameters\\
\hfuzz=500pt\quad\includegraphics[width=1em]{connector.pdf}\includegraphics[width=1em]{element.pdf}~integrationDegree & \hfuzz=500pt uint & \hfuzz=500pt integration of forces by polynomial approximation of degree n\\
\hfuzz=500pt\quad\includegraphics[width=1em]{connector.pdf}\includegraphics[width=1em]{element.pdf}~interpolationDegree & \hfuzz=500pt uint & \hfuzz=500pt orbit interpolation by polynomial approximation of degree n\\
\hfuzz=500pt\quad\includegraphics[width=1em]{connector.pdf}\includegraphics[width=1em]{element.pdf}~accelerateComputation & \hfuzz=500pt boolean & \hfuzz=500pt acceleration of computation by transforming the observations\\
\hfuzz=500pt\includegraphics[width=1em]{element-mustset.pdf}~covariancePod & \hfuzz=500pt sequence & \hfuzz=500pt \\
\hfuzz=500pt\includegraphics[width=1em]{connector.pdf}\includegraphics[width=1em]{element.pdf}~sigma & \hfuzz=500pt double & \hfuzz=500pt apriori factor of covariance function\\
\hfuzz=500pt\includegraphics[width=1em]{connector.pdf}\includegraphics[width=1em]{element.pdf}~inputfileSigmasPerArc & \hfuzz=500pt filename & \hfuzz=500pt apriori different accuaries for each arc (multiplicated with sigma)\\
\hfuzz=500pt\includegraphics[width=1em]{connector.pdf}\includegraphics[width=1em]{element.pdf}~inputfileCovarianceFunction & \hfuzz=500pt filename & \hfuzz=500pt approximate covariances in time\\
\hfuzz=500pt\includegraphics[width=1em]{connector.pdf}\includegraphics[width=1em]{element.pdf}~inputfileCovariancePodEpoch & \hfuzz=500pt filename & \hfuzz=500pt 3x3 epoch covariances\\
\hfuzz=500pt\includegraphics[width=1em]{connector.pdf}\includegraphics[width=1em]{element.pdf}~sampling & \hfuzz=500pt double & \hfuzz=500pt [seconds] sampling of the covariance function\\
\hfuzz=500pt\includegraphics[width=1em]{element.pdf}~inputfileArcList & \hfuzz=500pt filename & \hfuzz=500pt list to correspond points of time to arc numbers\\
\hfuzz=500pt\includegraphics[width=1em]{element.pdf}~adjustmentThreshold & \hfuzz=500pt double & \hfuzz=500pt Adjustment factor threshold: Iteration will be stopped once both SST and POD adjustment factors are under this threshold\\
\hfuzz=500pt\includegraphics[width=1em]{element.pdf}~iterationCount & \hfuzz=500pt uint & \hfuzz=500pt (maximum) number of iterations for the estimation of calibration parameter and error PSD\\
\hline
\end{tabularx}

This program is \reference{parallelized}{general.parallelization}.
\clearpage
%==================================
\subsection{PreprocessingSst}\label{PreprocessingSst}
This program processes satellite-to-satellite-tracking (SST) and kinematic orbit observations in a GRACE like configuration.
Three different observation groups are considered separately: SST and POD1/POD2 for the two satellites.
This program works similar to \program{PreprocessingPod}, see there for details. Here only deviations
in the settings are explained.

Precise orbit data (POD) often contains systematic errors in addition to stochastic noise. In this case the
variance component estimation fails and assigns too much weight to the POD data. Therefore an additional
\config{downweightPod} factor can be applied to the standard deviation of POD for the next least squares adjustment
in the iteration. This factor should also applied as \config{sigma} in \configClass{observation}{observationType}
for computation of the final solution e.g. with \program{NormalsSolverVCE}.

Short time variations of the gravity field can be co-estimated together with the static/monthly
mean gravity field. The short time parameters must also be set in \configClass{observation:parametrizationGravity}{parametrizationGravityType} and
can then be selected by \configClass{estimateShortTimeVariations:parameterSelection}{parameterSelectorType}.
If these parameters are not time variable, for example when a range of static parameters is selected,
they are set up as constant for each time interval defined in \config{inputfileArcList}. The parameters are constrained by an
\configClass{estimateShortTimeVariations:autoregressiveModelSequence}{autoregressiveModelSequenceType}. The weight of
the constrain equations in terms of the standard deviation can be estimated by means of
Variance Component Estimation (VCE) if \config{estimateShortTimeVariations:estimateSigma} is set.
The mathematical background of this co-estimation can be found in:

Kvas, A., Mayer-Gürr, T. GRACE gravity field recovery with background model uncertainties.
J Geod 93, 2543–2552 (2019). \url{https://doi.org/10.1007/s00190-019-01314-1}.


\keepXColumns
\begin{tabularx}{\textwidth}{N T A}
\hline
Name & Type & Annotation\\
\hline
\hfuzz=500pt\includegraphics[width=1em]{element.pdf}~outputfileSolution & \hfuzz=500pt filename & \hfuzz=500pt estimated parameter vector (static part only)\\
\hfuzz=500pt\includegraphics[width=1em]{element.pdf}~outputfileSigmax & \hfuzz=500pt filename & \hfuzz=500pt standard deviations of the parameters (sqrt of the diagonal of the inverse normal equation)\\
\hfuzz=500pt\includegraphics[width=1em]{element.pdf}~outputfileParameterName & \hfuzz=500pt filename & \hfuzz=500pt estimated signal parameters (index is appended)\\
\hfuzz=500pt\includegraphics[width=1em]{element.pdf}~estimateArcSigmas & \hfuzz=500pt sequence & \hfuzz=500pt \\
\hfuzz=500pt\includegraphics[width=1em]{connector.pdf}\includegraphics[width=1em]{element.pdf}~outputfileSigmasPerArcSst & \hfuzz=500pt filename & \hfuzz=500pt accuracies of each arc (SST)\\
\hfuzz=500pt\includegraphics[width=1em]{connector.pdf}\includegraphics[width=1em]{element.pdf}~outputfileSigmasPerArcPod1 & \hfuzz=500pt filename & \hfuzz=500pt accuracies of each arc (POD1)\\
\hfuzz=500pt\includegraphics[width=1em]{connector.pdf}\includegraphics[width=1em]{element.pdf}~outputfileSigmasPerArcPod2 & \hfuzz=500pt filename & \hfuzz=500pt accuracies of each arc (POD2)\\
\hfuzz=500pt\includegraphics[width=1em]{element.pdf}~estimateEpochSigmas & \hfuzz=500pt sequence & \hfuzz=500pt \\
\hfuzz=500pt\includegraphics[width=1em]{connector.pdf}\includegraphics[width=1em]{element.pdf}~outputfileSigmasPerEpochSst & \hfuzz=500pt filename & \hfuzz=500pt accuracies of each epoch (SST)\\
\hfuzz=500pt\includegraphics[width=1em]{connector.pdf}\includegraphics[width=1em]{element.pdf}~outputfileSigmasPerEpochPod1 & \hfuzz=500pt filename & \hfuzz=500pt accuracies of each epoch (POD1)\\
\hfuzz=500pt\includegraphics[width=1em]{connector.pdf}\includegraphics[width=1em]{element.pdf}~outputfileSigmasPerEpochPod2 & \hfuzz=500pt filename & \hfuzz=500pt accuracies of each epoch (POD2)\\
\hfuzz=500pt\includegraphics[width=1em]{element.pdf}~estimateCovarianceFunctions & \hfuzz=500pt sequence & \hfuzz=500pt \\
\hfuzz=500pt\includegraphics[width=1em]{connector.pdf}\includegraphics[width=1em]{element.pdf}~outputfileCovarianceFunctionSst & \hfuzz=500pt filename & \hfuzz=500pt covariance function\\
\hfuzz=500pt\includegraphics[width=1em]{connector.pdf}\includegraphics[width=1em]{element.pdf}~outputfileCovarianceFunctionPod1 & \hfuzz=500pt filename & \hfuzz=500pt covariance functions for along, cross, radial direction\\
\hfuzz=500pt\includegraphics[width=1em]{connector.pdf}\includegraphics[width=1em]{element.pdf}~outputfileCovarianceFunctionPod2 & \hfuzz=500pt filename & \hfuzz=500pt covariance functions for along, cross, radial direction\\
\hfuzz=500pt\includegraphics[width=1em]{element.pdf}~estimateSstArcCovarianceSigmas & \hfuzz=500pt sequence & \hfuzz=500pt \\
\hfuzz=500pt\includegraphics[width=1em]{connector.pdf}\includegraphics[width=1em]{element.pdf}~outputfileSigmasCovarianceMatrixArc & \hfuzz=500pt filename & \hfuzz=500pt one variance factor per matrix\\
\hfuzz=500pt\includegraphics[width=1em]{element.pdf}~computeResiduals & \hfuzz=500pt sequence & \hfuzz=500pt \\
\hfuzz=500pt\includegraphics[width=1em]{connector.pdf}\includegraphics[width=1em]{element.pdf}~outputfileSstResiduals & \hfuzz=500pt filename & \hfuzz=500pt \\
\hfuzz=500pt\includegraphics[width=1em]{connector.pdf}\includegraphics[width=1em]{element.pdf}~outputfilePod1Residuals & \hfuzz=500pt filename & \hfuzz=500pt \\
\hfuzz=500pt\includegraphics[width=1em]{connector.pdf}\includegraphics[width=1em]{element.pdf}~outputfilePod2Residuals & \hfuzz=500pt filename & \hfuzz=500pt \\
\hfuzz=500pt\includegraphics[width=1em]{element-mustset.pdf}~observation & \hfuzz=500pt choice & \hfuzz=500pt obervation equations (Sst)\\
\hfuzz=500pt\includegraphics[width=1em]{connector.pdf}\includegraphics[width=1em]{element-mustset.pdf}~sstIntegral & \hfuzz=500pt sequence & \hfuzz=500pt integral approach\\
\hfuzz=500pt\quad\includegraphics[width=1em]{connector.pdf}\includegraphics[width=1em]{element.pdf}~inputfileSatelliteModel1 & \hfuzz=500pt filename & \hfuzz=500pt satellite macro model\\
\hfuzz=500pt\quad\includegraphics[width=1em]{connector.pdf}\includegraphics[width=1em]{element.pdf}~inputfileSatelliteModel2 & \hfuzz=500pt filename & \hfuzz=500pt satellite macro model\\
\hfuzz=500pt\quad\includegraphics[width=1em]{connector.pdf}\includegraphics[width=1em]{element-mustset-unbounded.pdf}~rightHandSide & \hfuzz=500pt \hyperref[sstRightSideType]{sstRightSide} & \hfuzz=500pt input for the reduced observation vector\\
\hfuzz=500pt\quad\includegraphics[width=1em]{connector.pdf}\includegraphics[width=1em]{element-mustset.pdf}~sstType & \hfuzz=500pt choice & \hfuzz=500pt \\
\hfuzz=500pt\quad\quad\includegraphics[width=1em]{connector.pdf}\includegraphics[width=1em]{element-mustset.pdf}~range & \hfuzz=500pt  & \hfuzz=500pt \\
\hfuzz=500pt\quad\quad\includegraphics[width=1em]{connector.pdf}\includegraphics[width=1em]{element-mustset.pdf}~rangeRate & \hfuzz=500pt  & \hfuzz=500pt \\
\hfuzz=500pt\quad\quad\includegraphics[width=1em]{connector.pdf}\includegraphics[width=1em]{element-mustset.pdf}~rangeAcceleration & \hfuzz=500pt  & \hfuzz=500pt \\
\hfuzz=500pt\quad\quad\includegraphics[width=1em]{connector.pdf}\includegraphics[width=1em]{element-mustset.pdf}~none & \hfuzz=500pt  & \hfuzz=500pt \\
\hfuzz=500pt\quad\includegraphics[width=1em]{connector.pdf}\includegraphics[width=1em]{element-mustset.pdf}~inputfileOrbit1 & \hfuzz=500pt filename & \hfuzz=500pt used to evaluate the observation equations, not used as observations\\
\hfuzz=500pt\quad\includegraphics[width=1em]{connector.pdf}\includegraphics[width=1em]{element-mustset.pdf}~inputfileOrbit2 & \hfuzz=500pt filename & \hfuzz=500pt used to evaluate the observation equations, not used as observations\\
\hfuzz=500pt\quad\includegraphics[width=1em]{connector.pdf}\includegraphics[width=1em]{element-mustset.pdf}~inputfileStarCamera1 & \hfuzz=500pt filename & \hfuzz=500pt \\
\hfuzz=500pt\quad\includegraphics[width=1em]{connector.pdf}\includegraphics[width=1em]{element-mustset.pdf}~inputfileStarCamera2 & \hfuzz=500pt filename & \hfuzz=500pt \\
\hfuzz=500pt\quad\includegraphics[width=1em]{connector.pdf}\includegraphics[width=1em]{element-mustset.pdf}~earthRotation & \hfuzz=500pt \hyperref[earthRotationType]{earthRotation} & \hfuzz=500pt \\
\hfuzz=500pt\quad\includegraphics[width=1em]{connector.pdf}\includegraphics[width=1em]{element.pdf}~ephemerides & \hfuzz=500pt \hyperref[ephemeridesType]{ephemerides} & \hfuzz=500pt \\
\hfuzz=500pt\quad\includegraphics[width=1em]{connector.pdf}\includegraphics[width=1em]{element-unbounded.pdf}~gradientfield & \hfuzz=500pt \hyperref[gravityfieldType]{gravityfield} & \hfuzz=500pt low order field to estimate the change of the gravity by position adjustement\\
\hfuzz=500pt\quad\includegraphics[width=1em]{connector.pdf}\includegraphics[width=1em]{element-unbounded.pdf}~parametrizationGravity & \hfuzz=500pt \hyperref[parametrizationGravityType]{parametrizationGravity} & \hfuzz=500pt gravity field parametrization\\
\hfuzz=500pt\quad\includegraphics[width=1em]{connector.pdf}\includegraphics[width=1em]{element-unbounded.pdf}~parametrizationAcceleration1 & \hfuzz=500pt \hyperref[parametrizationAccelerationType]{parametrizationAcceleration} & \hfuzz=500pt orbit1 force parameters\\
\hfuzz=500pt\quad\includegraphics[width=1em]{connector.pdf}\includegraphics[width=1em]{element-unbounded.pdf}~parametrizationAcceleration2 & \hfuzz=500pt \hyperref[parametrizationAccelerationType]{parametrizationAcceleration} & \hfuzz=500pt orbit2 force parameters\\
\hfuzz=500pt\quad\includegraphics[width=1em]{connector.pdf}\includegraphics[width=1em]{element-unbounded.pdf}~parametrizationSst & \hfuzz=500pt \hyperref[parametrizationSatelliteTrackingType]{parametrizationSatelliteTracking} & \hfuzz=500pt satellite tracking parameter\\
\hfuzz=500pt\quad\includegraphics[width=1em]{connector.pdf}\includegraphics[width=1em]{element.pdf}~keepSatelliteStates & \hfuzz=500pt boolean & \hfuzz=500pt set boundary values of each arc global\\
\hfuzz=500pt\quad\includegraphics[width=1em]{connector.pdf}\includegraphics[width=1em]{element.pdf}~integrationDegree & \hfuzz=500pt uint & \hfuzz=500pt integration of forces by polynomial approximation of degree n\\
\hfuzz=500pt\quad\includegraphics[width=1em]{connector.pdf}\includegraphics[width=1em]{element.pdf}~interpolationDegree & \hfuzz=500pt uint & \hfuzz=500pt orbit interpolation by polynomial approximation of degree n\\
\hfuzz=500pt\includegraphics[width=1em]{connector.pdf}\includegraphics[width=1em]{element-mustset.pdf}~sstVariational & \hfuzz=500pt sequence & \hfuzz=500pt variational equations\\
\hfuzz=500pt\quad\includegraphics[width=1em]{connector.pdf}\includegraphics[width=1em]{element-mustset.pdf}~rightHandSide & \hfuzz=500pt sequence & \hfuzz=500pt input for observation vectors\\
\hfuzz=500pt\quad\quad\includegraphics[width=1em]{connector.pdf}\includegraphics[width=1em]{element-mustset-unbounded.pdf}~inputfileSatelliteTracking & \hfuzz=500pt filename & \hfuzz=500pt ranging observations and corrections\\
\hfuzz=500pt\quad\quad\includegraphics[width=1em]{connector.pdf}\includegraphics[width=1em]{element.pdf}~inputfileOrbit1 & \hfuzz=500pt filename & \hfuzz=500pt kinematic positions of satellite A as observations\\
\hfuzz=500pt\quad\quad\includegraphics[width=1em]{connector.pdf}\includegraphics[width=1em]{element.pdf}~inputfileOrbit2 & \hfuzz=500pt filename & \hfuzz=500pt kinematic positions of satellite B as observations\\
\hfuzz=500pt\quad\includegraphics[width=1em]{connector.pdf}\includegraphics[width=1em]{element-mustset.pdf}~sstType & \hfuzz=500pt choice & \hfuzz=500pt \\
\hfuzz=500pt\quad\quad\includegraphics[width=1em]{connector.pdf}\includegraphics[width=1em]{element-mustset.pdf}~range & \hfuzz=500pt  & \hfuzz=500pt \\
\hfuzz=500pt\quad\quad\includegraphics[width=1em]{connector.pdf}\includegraphics[width=1em]{element-mustset.pdf}~rangeRate & \hfuzz=500pt  & \hfuzz=500pt \\
\hfuzz=500pt\quad\quad\includegraphics[width=1em]{connector.pdf}\includegraphics[width=1em]{element-mustset.pdf}~none & \hfuzz=500pt  & \hfuzz=500pt \\
\hfuzz=500pt\quad\includegraphics[width=1em]{connector.pdf}\includegraphics[width=1em]{element-mustset.pdf}~inputfileVariational1 & \hfuzz=500pt filename & \hfuzz=500pt approximate position and integrated state matrix\\
\hfuzz=500pt\quad\includegraphics[width=1em]{connector.pdf}\includegraphics[width=1em]{element-mustset.pdf}~inputfileVariational2 & \hfuzz=500pt filename & \hfuzz=500pt approximate position and integrated state matrix\\
\hfuzz=500pt\quad\includegraphics[width=1em]{connector.pdf}\includegraphics[width=1em]{element.pdf}~ephemerides & \hfuzz=500pt \hyperref[ephemeridesType]{ephemerides} & \hfuzz=500pt \\
\hfuzz=500pt\quad\includegraphics[width=1em]{connector.pdf}\includegraphics[width=1em]{element-unbounded.pdf}~parametrizationGravity & \hfuzz=500pt \hyperref[parametrizationGravityType]{parametrizationGravity} & \hfuzz=500pt gravity field parametrization\\
\hfuzz=500pt\quad\includegraphics[width=1em]{connector.pdf}\includegraphics[width=1em]{element-unbounded.pdf}~parametrizationAcceleration1 & \hfuzz=500pt \hyperref[parametrizationAccelerationType]{parametrizationAcceleration} & \hfuzz=500pt orbit1 force parameters\\
\hfuzz=500pt\quad\includegraphics[width=1em]{connector.pdf}\includegraphics[width=1em]{element-unbounded.pdf}~parametrizationAcceleration2 & \hfuzz=500pt \hyperref[parametrizationAccelerationType]{parametrizationAcceleration} & \hfuzz=500pt orbit2 force parameters\\
\hfuzz=500pt\quad\includegraphics[width=1em]{connector.pdf}\includegraphics[width=1em]{element-unbounded.pdf}~parametrizationSst & \hfuzz=500pt \hyperref[parametrizationSatelliteTrackingType]{parametrizationSatelliteTracking} & \hfuzz=500pt satellite tracking parameter\\
\hfuzz=500pt\quad\includegraphics[width=1em]{connector.pdf}\includegraphics[width=1em]{element.pdf}~integrationDegree & \hfuzz=500pt uint & \hfuzz=500pt integration of forces by polynomial approximation of degree n\\
\hfuzz=500pt\quad\includegraphics[width=1em]{connector.pdf}\includegraphics[width=1em]{element.pdf}~interpolationDegree & \hfuzz=500pt uint & \hfuzz=500pt orbit interpolation by polynomial approximation of degree n\\
\hfuzz=500pt\includegraphics[width=1em]{element-mustset.pdf}~covarianceSst & \hfuzz=500pt sequence & \hfuzz=500pt \\
\hfuzz=500pt\includegraphics[width=1em]{connector.pdf}\includegraphics[width=1em]{element.pdf}~sigma & \hfuzz=500pt double & \hfuzz=500pt apriori factor of covariance function\\
\hfuzz=500pt\includegraphics[width=1em]{connector.pdf}\includegraphics[width=1em]{element.pdf}~inputfileSigmasPerArc & \hfuzz=500pt filename & \hfuzz=500pt apriori different accuaries for each arc (multiplicated with sigma)\\
\hfuzz=500pt\includegraphics[width=1em]{connector.pdf}\includegraphics[width=1em]{element.pdf}~inputfileSigmasPerEpoch & \hfuzz=500pt filename & \hfuzz=500pt apriori different accuaries for each epoch\\
\hfuzz=500pt\includegraphics[width=1em]{connector.pdf}\includegraphics[width=1em]{element.pdf}~inputfileCovarianceFunction & \hfuzz=500pt filename & \hfuzz=500pt approximate covariances in time\\
\hfuzz=500pt\includegraphics[width=1em]{connector.pdf}\includegraphics[width=1em]{element-unbounded.pdf}~inputfileCovarianceMatrixArc & \hfuzz=500pt filename & \hfuzz=500pt Must be given per sst arc with correct dimensions.\\
\hfuzz=500pt\includegraphics[width=1em]{connector.pdf}\includegraphics[width=1em]{element.pdf}~inputfileSigmasCovarianceMatrixArc & \hfuzz=500pt filename & \hfuzz=500pt Vector with one sigma for each \$<\$inputfileCovarianceMatrixArc\$>\$\\
\hfuzz=500pt\includegraphics[width=1em]{connector.pdf}\includegraphics[width=1em]{element.pdf}~sampling & \hfuzz=500pt double & \hfuzz=500pt [seconds] sampling of the covariance function\\
\hfuzz=500pt\includegraphics[width=1em]{element-mustset.pdf}~covariancePod1 & \hfuzz=500pt sequence & \hfuzz=500pt \\
\hfuzz=500pt\includegraphics[width=1em]{connector.pdf}\includegraphics[width=1em]{element.pdf}~sigma & \hfuzz=500pt double & \hfuzz=500pt apriori factor of covariance function\\
\hfuzz=500pt\includegraphics[width=1em]{connector.pdf}\includegraphics[width=1em]{element.pdf}~inputfileSigmasPerArc & \hfuzz=500pt filename & \hfuzz=500pt apriori different accuaries for each arc (multiplicated with sigma)\\
\hfuzz=500pt\includegraphics[width=1em]{connector.pdf}\includegraphics[width=1em]{element.pdf}~inputfileSigmasPerEpoch & \hfuzz=500pt filename & \hfuzz=500pt apriori different accuaries for each epoch\\
\hfuzz=500pt\includegraphics[width=1em]{connector.pdf}\includegraphics[width=1em]{element.pdf}~inputfileCovarianceFunction & \hfuzz=500pt filename & \hfuzz=500pt approximate covariances in time\\
\hfuzz=500pt\includegraphics[width=1em]{connector.pdf}\includegraphics[width=1em]{element.pdf}~inputfileCovariancePodEpoch & \hfuzz=500pt filename & \hfuzz=500pt 3x3 epoch covariances\\
\hfuzz=500pt\includegraphics[width=1em]{connector.pdf}\includegraphics[width=1em]{element.pdf}~sampling & \hfuzz=500pt double & \hfuzz=500pt [seconds] sampling of the covariance function\\
\hfuzz=500pt\includegraphics[width=1em]{element-mustset.pdf}~covariancePod2 & \hfuzz=500pt sequence & \hfuzz=500pt \\
\hfuzz=500pt\includegraphics[width=1em]{connector.pdf}\includegraphics[width=1em]{element.pdf}~sigma & \hfuzz=500pt double & \hfuzz=500pt apriori factor of covariance function\\
\hfuzz=500pt\includegraphics[width=1em]{connector.pdf}\includegraphics[width=1em]{element.pdf}~inputfileSigmasPerArc & \hfuzz=500pt filename & \hfuzz=500pt apriori different accuaries for each arc (multiplicated with sigma)\\
\hfuzz=500pt\includegraphics[width=1em]{connector.pdf}\includegraphics[width=1em]{element.pdf}~inputfileSigmasPerEpoch & \hfuzz=500pt filename & \hfuzz=500pt apriori different accuaries for each epoch\\
\hfuzz=500pt\includegraphics[width=1em]{connector.pdf}\includegraphics[width=1em]{element.pdf}~inputfileCovarianceFunction & \hfuzz=500pt filename & \hfuzz=500pt approximate covariances in time\\
\hfuzz=500pt\includegraphics[width=1em]{connector.pdf}\includegraphics[width=1em]{element.pdf}~inputfileCovariancePodEpoch & \hfuzz=500pt filename & \hfuzz=500pt 3x3 epoch covariances\\
\hfuzz=500pt\includegraphics[width=1em]{connector.pdf}\includegraphics[width=1em]{element.pdf}~sampling & \hfuzz=500pt double & \hfuzz=500pt [seconds] sampling of the covariance function\\
\hfuzz=500pt\includegraphics[width=1em]{element.pdf}~estimateShortTimeVariations & \hfuzz=500pt sequence & \hfuzz=500pt co-estimate short time gravity field variations\\
\hfuzz=500pt\includegraphics[width=1em]{connector.pdf}\includegraphics[width=1em]{element.pdf}~estimateSigma & \hfuzz=500pt boolean & \hfuzz=500pt estimate standard deviation via VCE\\
\hfuzz=500pt\includegraphics[width=1em]{connector.pdf}\includegraphics[width=1em]{element-mustset.pdf}~autoregressiveModelSequence & \hfuzz=500pt \hyperref[autoregressiveModelSequenceType]{autoregressiveModelSequence} & \hfuzz=500pt AR model sequence for constraining short time gravity variations\\
\hfuzz=500pt\includegraphics[width=1em]{connector.pdf}\includegraphics[width=1em]{element-mustset-unbounded.pdf}~parameterSelection & \hfuzz=500pt \hyperref[parameterSelectorType]{parameterSelector} & \hfuzz=500pt parameters describing the short time gravity field\\
\hfuzz=500pt\includegraphics[width=1em]{element.pdf}~downweightPod & \hfuzz=500pt double & \hfuzz=500pt downweight factor for POD\\
\hfuzz=500pt\includegraphics[width=1em]{element.pdf}~inputfileArcList & \hfuzz=500pt filename & \hfuzz=500pt list to correspond points of time to arc numbers\\
\hfuzz=500pt\includegraphics[width=1em]{element.pdf}~iterationCount & \hfuzz=500pt uint & \hfuzz=500pt (maximum) number of iterations for the estimation of calibration parameter and error PSD\\
\hfuzz=500pt\includegraphics[width=1em]{element.pdf}~variableNameIterations & \hfuzz=500pt string & \hfuzz=500pt All output fileNames in preprocessing iteration are expanded with this variable prior to writing to disk\\
\hfuzz=500pt\includegraphics[width=1em]{element.pdf}~defaultBlockSize & \hfuzz=500pt uint & \hfuzz=500pt block size of static normal equation blocks\\
\hline
\end{tabularx}

This program is \reference{parallelized}{general.parallelization}.
\clearpage
%==================================
\subsection{PreprocessingVariationalEquation}\label{PreprocessingVariationalEquation}
This program integrates an orbit dynamically using the given forces and setup the state transition matrix
for each time step. These are the prerequisite for a least squares adjustement (e.g. gravity field determination) using
the variational equation approach. The variational equations are computed arc wise as defined by \configFile{inputfileOrbit}{instrument}.
This means for each arc new initial state parameters are setup.

In a first step the \configClass{forces}{forcesType} acting on the satellite are evaluated at the apriori positions given
by \configFile{inputfileOrbit}{instrument}. Non-conservative forces like solar radiation pressure needs the orientation of the
satellite (\configFile{inputfileStarCamera}{instrument}) and additional a satellite macro model (\config{satelliteModel})
with the surface properties. Furthermore \configFile{inputfileAccelerometer}{instrument} observations are also considered.

In a second step the accelerations are integrated twice to an dynamic orbit using a moving polynomial with the degree
\config{integrationDegree}. The orbit is corrected to be self-consistent. This means the forces should be evaluated
at the new integrated positions instead of the apriori ones. This correction is computed in a linear approximation
using the gradient of the forces with respect to the positions (\config{gradientfield}). As this term is small generally
only the largest force components has to be considered. A low degree spherical harmonic expansion of the static gravity
field (about up to degree 5) is sufficient in almost all cases. In this step also the state transition matrix (the partial
derivatices of the current state, position and velocity) with respect to the initial state is computed.
The integrated orbit together with the state transitions are stored in \configFile{outputfileVariational}{variationalEquation},
the integrated orbit only in \configFile{outputfileOrbit}{instrument}.

To improve the numerical stability a reference ellipse can be reduced beforehand using Enke's method (\config{useEnke}).
Mathematically the result is the same, but as the large central term is removed before and restored
afterwards more digits are available for the computation.

The integrated orbit should be fitted to observations afterwards by the programs
\program{PreprocessingVariationalEquationOrbitFit} and/or \program{PreprocessingVariationalEquationSstFit}.
They apply a least squares adjustment by estimating some satellite parameters (e.g. an accelerometer bias).
If the fitted orbit is too far away from the original \configFile{inputfileOrbit}{instrument} the linearization may not be
accurate enough. In this case \program{PreprocessingVariationalEquation} should be run again with the fitted orbit
as \configFile{inputfileOrbit}{instrument} and introducing the \config{estimatedParameters} as additional forces.


\keepXColumns
\begin{tabularx}{\textwidth}{N T A}
\hline
Name & Type & Annotation\\
\hline
\hfuzz=500pt\includegraphics[width=1em]{element-mustset.pdf}~outputfileVariational & \hfuzz=500pt filename & \hfuzz=500pt approximate position and integrated state matrix\\
\hfuzz=500pt\includegraphics[width=1em]{element.pdf}~outputfileOrbit & \hfuzz=500pt filename & \hfuzz=500pt integrated orbit\\
\hfuzz=500pt\includegraphics[width=1em]{element.pdf}~inputfileSatelliteModel & \hfuzz=500pt filename & \hfuzz=500pt satellite macro model\\
\hfuzz=500pt\includegraphics[width=1em]{element-mustset.pdf}~inputfileOrbit & \hfuzz=500pt filename & \hfuzz=500pt approximate position, used to evaluate the force\\
\hfuzz=500pt\includegraphics[width=1em]{element-mustset.pdf}~inputfileStarCamera & \hfuzz=500pt filename & \hfuzz=500pt rotation from body frame to CRF\\
\hfuzz=500pt\includegraphics[width=1em]{element.pdf}~inputfileAccelerometer & \hfuzz=500pt filename & \hfuzz=500pt non-gravitational forces in satellite reference frame\\
\hfuzz=500pt\includegraphics[width=1em]{element-mustset.pdf}~forces & \hfuzz=500pt \hyperref[forcesType]{forces} & \hfuzz=500pt \\
\hfuzz=500pt\includegraphics[width=1em]{element.pdf}~estimatedParameters & \hfuzz=500pt sequence & \hfuzz=500pt satellite parameters e.g. from orbit fit\\
\hfuzz=500pt\includegraphics[width=1em]{connector.pdf}\includegraphics[width=1em]{element-mustset-unbounded.pdf}~parametrizationAcceleration & \hfuzz=500pt \hyperref[parametrizationAccelerationType]{parametrizationAcceleration} & \hfuzz=500pt orbit force parameters\\
\hfuzz=500pt\includegraphics[width=1em]{connector.pdf}\includegraphics[width=1em]{element-mustset.pdf}~inputfileParameter & \hfuzz=500pt filename & \hfuzz=500pt estimated orbit force parameters\\
\hfuzz=500pt\includegraphics[width=1em]{element-mustset.pdf}~earthRotation & \hfuzz=500pt \hyperref[earthRotationType]{earthRotation} & \hfuzz=500pt \\
\hfuzz=500pt\includegraphics[width=1em]{element.pdf}~ephemerides & \hfuzz=500pt \hyperref[ephemeridesType]{ephemerides} & \hfuzz=500pt \\
\hfuzz=500pt\includegraphics[width=1em]{element-mustset-unbounded.pdf}~gradientfield & \hfuzz=500pt \hyperref[gravityfieldType]{gravityfield} & \hfuzz=500pt low order field to estimate the change of the gravity by position adjustement\\
\hfuzz=500pt\includegraphics[width=1em]{element.pdf}~integrationDegree & \hfuzz=500pt uint & \hfuzz=500pt integration of forces by polynomial approximation of degree n\\
\hfuzz=500pt\includegraphics[width=1em]{element.pdf}~useEnke & \hfuzz=500pt sequence & \hfuzz=500pt integrate differential forces to an elliptical reference trajectory\\
\hfuzz=500pt\includegraphics[width=1em]{connector.pdf}\includegraphics[width=1em]{element.pdf}~GM & \hfuzz=500pt double & \hfuzz=500pt geocentric gravitational constant used for elliptical reference orbit\\
\hline
\end{tabularx}

This program is \reference{parallelized}{general.parallelization}.
\clearpage
%==================================
\subsection{PreprocessingVariationalEquationOrbitFit}\label{PreprocessingVariationalEquationOrbitFit}
This program fits an \configFile{inputfileVariational}{variationalEquation} to an observed \configFile{inputfileOrbit}{instrument} by estimating parameters
in a least squares adjustment. Additional to the initial satellite state for each arc, these parameters can be
\configClass{parametrizationGravity}{parametrizationGravityType}, satellite \configClass{parametrizationAcceleration}{parametrizationAccelerationType}
and stochastic pulses (velocity jumps) at given times, \configClass{stochasticPulse}{timeSeriesType}. The estimated parameters can be stored with
\configFile{outputfileSolution}{matrix} and an extra file with the parameter names is created. The fitted orbit is written
as new reference in \configFile{outputfileVariational}{variationalEquation} and additionally in \configFile{outputfileOrbit}{instrument}.

The observed orbit positions (\configFile{inputfileOrbit}{instrument}) together with the epoch wise covariance matrix
(\configFile{inputfileCovariancePodEpoch}{instrument}) must be splitted in the same arcs as the variational equations but not
necessarily uniform distributed (use irregularData in \program{InstrumentSynchronize}). An iterative downweighting of
outliers is performed by M-Huber method.

The observation equations (parameter sensitity matrix) are computed by integration of the variational equations
(\configFile{inputfileVariational}{variationalEquation}) using a polynomial with \config{integrationDegree} and interpolated to the
observation epochs using a polynomial with \config{interpolationDegree}.

All parameters used here must be reestimated in the full least squares adjustment
for the gravity field determination to get a solution which is not biased towards the reference field.
The solution of additional estimations are relative (deltas) as the parameters are already used as Taylor point
in the reference orbit.

See also \program{PreprocessingVariationalEquation}.


\keepXColumns
\begin{tabularx}{\textwidth}{N T A}
\hline
Name & Type & Annotation\\
\hline
\hfuzz=500pt\includegraphics[width=1em]{element-mustset.pdf}~outputfileVariational & \hfuzz=500pt filename & \hfuzz=500pt approximate position and integrated state matrix\\
\hfuzz=500pt\includegraphics[width=1em]{element.pdf}~outputfileOrbit & \hfuzz=500pt filename & \hfuzz=500pt integrated orbit\\
\hfuzz=500pt\includegraphics[width=1em]{element.pdf}~outputfileSolution & \hfuzz=500pt filename & \hfuzz=500pt estimated calibration and state parameters\\
\hfuzz=500pt\includegraphics[width=1em]{element-mustset.pdf}~inputfileVariational & \hfuzz=500pt filename & \hfuzz=500pt approximate position and integrated state matrix\\
\hfuzz=500pt\includegraphics[width=1em]{element-mustset.pdf}~inputfileOrbit & \hfuzz=500pt filename & \hfuzz=500pt kinematic positions of satellite as observations\\
\hfuzz=500pt\includegraphics[width=1em]{element.pdf}~inputfileCovariancePodEpoch & \hfuzz=500pt filename & \hfuzz=500pt 3x3 epoch wise covariances\\
\hfuzz=500pt\includegraphics[width=1em]{element.pdf}~ephemerides & \hfuzz=500pt \hyperref[ephemeridesType]{ephemerides} & \hfuzz=500pt may be needed by parametrizationAcceleration\\
\hfuzz=500pt\includegraphics[width=1em]{element-unbounded.pdf}~parametrizationGravity & \hfuzz=500pt \hyperref[parametrizationGravityType]{parametrizationGravity} & \hfuzz=500pt gravity field parametrization\\
\hfuzz=500pt\includegraphics[width=1em]{element-unbounded.pdf}~parametrizationAcceleration & \hfuzz=500pt \hyperref[parametrizationAccelerationType]{parametrizationAcceleration} & \hfuzz=500pt orbit force parameters\\
\hfuzz=500pt\includegraphics[width=1em]{element-unbounded.pdf}~stochasticPulse & \hfuzz=500pt \hyperref[timeSeriesType]{timeSeries} & \hfuzz=500pt \\
\hfuzz=500pt\includegraphics[width=1em]{element.pdf}~integrationDegree & \hfuzz=500pt uint & \hfuzz=500pt integration of forces by polynomial approximation of degree n\\
\hfuzz=500pt\includegraphics[width=1em]{element.pdf}~interpolationDegree & \hfuzz=500pt uint & \hfuzz=500pt orbit interpolation by polynomial approximation of degree n\\
\hfuzz=500pt\includegraphics[width=1em]{element.pdf}~iterationCount & \hfuzz=500pt uint & \hfuzz=500pt for the estimation of calibration parameter and error PSD\\
\hline
\end{tabularx}

This program is \reference{parallelized}{general.parallelization}.
\clearpage
%==================================
\subsection{PreprocessingVariationalEquationSstFit}\label{PreprocessingVariationalEquationSstFit}
This program fits two \configFile{inputfileVariational1/2}{variationalEquation} to satellite-to-satellite-tracking (SST) and orbit
observations in a GRACE like configuration. It works similar to \program{PreprocessingVariationalEquationOrbitFit},
see there for details.

As the relative weighting of the observation types is important complex description of the covariances can be set with
\configClass{covarianceSst}{covarianceSstType}, \configClass{covariancePod1}{covariancePodType}, \configClass{covariancePod2}{covariancePodType}.


\keepXColumns
\begin{tabularx}{\textwidth}{N T A}
\hline
Name & Type & Annotation\\
\hline
\hfuzz=500pt\includegraphics[width=1em]{element-mustset.pdf}~outputfileVariational1 & \hfuzz=500pt filename & \hfuzz=500pt approximate position and integrated state matrix\\
\hfuzz=500pt\includegraphics[width=1em]{element-mustset.pdf}~outputfileVariational2 & \hfuzz=500pt filename & \hfuzz=500pt approximate position and integrated state matrix\\
\hfuzz=500pt\includegraphics[width=1em]{element.pdf}~outputfileOrbit1 & \hfuzz=500pt filename & \hfuzz=500pt integrated orbit\\
\hfuzz=500pt\includegraphics[width=1em]{element.pdf}~outputfileOrbit2 & \hfuzz=500pt filename & \hfuzz=500pt integrated orbit\\
\hfuzz=500pt\includegraphics[width=1em]{element.pdf}~outputfileSolution1 & \hfuzz=500pt filename & \hfuzz=500pt estimated calibration and state parameters\\
\hfuzz=500pt\includegraphics[width=1em]{element.pdf}~outputfileSolution2 & \hfuzz=500pt filename & \hfuzz=500pt estimated calibration and state parameters\\
\hfuzz=500pt\includegraphics[width=1em]{element-mustset.pdf}~rightHandSide & \hfuzz=500pt sequence & \hfuzz=500pt input for observation vectors\\
\hfuzz=500pt\includegraphics[width=1em]{connector.pdf}\includegraphics[width=1em]{element-mustset-unbounded.pdf}~inputfileSatelliteTracking & \hfuzz=500pt filename & \hfuzz=500pt ranging observations and corrections\\
\hfuzz=500pt\includegraphics[width=1em]{connector.pdf}\includegraphics[width=1em]{element.pdf}~inputfileOrbit1 & \hfuzz=500pt filename & \hfuzz=500pt kinematic positions of satellite A as observations\\
\hfuzz=500pt\includegraphics[width=1em]{connector.pdf}\includegraphics[width=1em]{element.pdf}~inputfileOrbit2 & \hfuzz=500pt filename & \hfuzz=500pt kinematic positions of satellite B as observations\\
\hfuzz=500pt\includegraphics[width=1em]{element-mustset.pdf}~sstType & \hfuzz=500pt choice & \hfuzz=500pt \\
\hfuzz=500pt\includegraphics[width=1em]{connector.pdf}\includegraphics[width=1em]{element-mustset.pdf}~range & \hfuzz=500pt  & \hfuzz=500pt \\
\hfuzz=500pt\includegraphics[width=1em]{connector.pdf}\includegraphics[width=1em]{element-mustset.pdf}~rangeRate & \hfuzz=500pt  & \hfuzz=500pt \\
\hfuzz=500pt\includegraphics[width=1em]{connector.pdf}\includegraphics[width=1em]{element-mustset.pdf}~none & \hfuzz=500pt  & \hfuzz=500pt \\
\hfuzz=500pt\includegraphics[width=1em]{element-mustset.pdf}~inputfileVariational1 & \hfuzz=500pt filename & \hfuzz=500pt approximate position and integrated state matrix\\
\hfuzz=500pt\includegraphics[width=1em]{element-mustset.pdf}~inputfileVariational2 & \hfuzz=500pt filename & \hfuzz=500pt approximate position and integrated state matrix\\
\hfuzz=500pt\includegraphics[width=1em]{element.pdf}~ephemerides & \hfuzz=500pt \hyperref[ephemeridesType]{ephemerides} & \hfuzz=500pt \\
\hfuzz=500pt\includegraphics[width=1em]{element-unbounded.pdf}~parametrizationGravity & \hfuzz=500pt \hyperref[parametrizationGravityType]{parametrizationGravity} & \hfuzz=500pt gravity field parametrization\\
\hfuzz=500pt\includegraphics[width=1em]{element-unbounded.pdf}~parametrizationAcceleration1 & \hfuzz=500pt \hyperref[parametrizationAccelerationType]{parametrizationAcceleration} & \hfuzz=500pt orbit1 force parameters\\
\hfuzz=500pt\includegraphics[width=1em]{element-unbounded.pdf}~parametrizationAcceleration2 & \hfuzz=500pt \hyperref[parametrizationAccelerationType]{parametrizationAcceleration} & \hfuzz=500pt orbit2 force parameters\\
\hfuzz=500pt\includegraphics[width=1em]{element-unbounded.pdf}~parametrizationSst & \hfuzz=500pt \hyperref[parametrizationSatelliteTrackingType]{parametrizationSatelliteTracking} & \hfuzz=500pt satellite tracking parameter\\
\hfuzz=500pt\includegraphics[width=1em]{element.pdf}~integrationDegree & \hfuzz=500pt uint & \hfuzz=500pt integration of forces by polynomial approximation of degree n\\
\hfuzz=500pt\includegraphics[width=1em]{element.pdf}~interpolationDegree & \hfuzz=500pt uint & \hfuzz=500pt orbit interpolation by polynomial approximation of degree n\\
\hfuzz=500pt\includegraphics[width=1em]{element-mustset.pdf}~covarianceSst & \hfuzz=500pt \hyperref[covarianceSstType]{covarianceSst} & \hfuzz=500pt covariance matrix of satellite to satellite tracking observations\\
\hfuzz=500pt\includegraphics[width=1em]{element-mustset.pdf}~covariancePod1 & \hfuzz=500pt \hyperref[covariancePodType]{covariancePod} & \hfuzz=500pt covariance matrix of kinematic orbits (satellite 1)\\
\hfuzz=500pt\includegraphics[width=1em]{element-mustset.pdf}~covariancePod2 & \hfuzz=500pt \hyperref[covariancePodType]{covariancePod} & \hfuzz=500pt covariance matrix of kinematic orbits (satellite 2)\\
\hfuzz=500pt\includegraphics[width=1em]{element.pdf}~iterationCount & \hfuzz=500pt uint & \hfuzz=500pt for the estimation of calibration parameter and error PSD\\
\hline
\end{tabularx}

This program is \reference{parallelized}{general.parallelization}.
\clearpage
%==================================
\section{Programs: Simulation}
\subsection{NoiseAccelerometer}\label{NoiseAccelerometer}
This program adds noise and biases to simulated \file{accelerometer data}{instrument}
generated by \program{SimulateAccelerometer}.
See \configClass{noiseGenerator}{noiseGeneratorType} for details on noise generation.


\keepXColumns
\begin{tabularx}{\textwidth}{N T A}
\hline
Name & Type & Annotation\\
\hline
\hfuzz=500pt\includegraphics[width=1em]{element-mustset.pdf}~outputfileAccelerometer & \hfuzz=500pt filename & \hfuzz=500pt \\
\hfuzz=500pt\includegraphics[width=1em]{element-mustset.pdf}~inputfileAccelerometer & \hfuzz=500pt filename & \hfuzz=500pt \\
\hfuzz=500pt\includegraphics[width=1em]{element.pdf}~biasAlong & \hfuzz=500pt double & \hfuzz=500pt [m/s**2]\\
\hfuzz=500pt\includegraphics[width=1em]{element.pdf}~biasCross & \hfuzz=500pt double & \hfuzz=500pt [m/s**2]\\
\hfuzz=500pt\includegraphics[width=1em]{element.pdf}~biasRadial & \hfuzz=500pt double & \hfuzz=500pt [m/s**2]\\
\hfuzz=500pt\includegraphics[width=1em]{element-unbounded.pdf}~noiseAlong & \hfuzz=500pt \hyperref[noiseGeneratorType]{noiseGenerator} & \hfuzz=500pt [m/s**2]\\
\hfuzz=500pt\includegraphics[width=1em]{element-unbounded.pdf}~noiseCross & \hfuzz=500pt \hyperref[noiseGeneratorType]{noiseGenerator} & \hfuzz=500pt [m/s**2]\\
\hfuzz=500pt\includegraphics[width=1em]{element-unbounded.pdf}~noiseRadial & \hfuzz=500pt \hyperref[noiseGeneratorType]{noiseGenerator} & \hfuzz=500pt [m/s**2]\\
\hline
\end{tabularx}

This program is \reference{parallelized}{general.parallelization}.
\clearpage
%==================================
\subsection{NoiseGriddedData}\label{NoiseGriddedData}
This program adds noise to \file{gridded data data}{griddedData}.
See \configClass{noiseGenerator}{noiseGeneratorType} for details on noise generation.


\keepXColumns
\begin{tabularx}{\textwidth}{N T A}
\hline
Name & Type & Annotation\\
\hline
\hfuzz=500pt\includegraphics[width=1em]{element-mustset.pdf}~outputfileGriddedData & \hfuzz=500pt filename & \hfuzz=500pt \\
\hfuzz=500pt\includegraphics[width=1em]{element-mustset.pdf}~inputfileGriddedData & \hfuzz=500pt filename & \hfuzz=500pt \\
\hfuzz=500pt\includegraphics[width=1em]{element-mustset-unbounded.pdf}~noise & \hfuzz=500pt \hyperref[noiseGeneratorType]{noiseGenerator} & \hfuzz=500pt \\
\hfuzz=500pt\includegraphics[width=1em]{element.pdf}~startDataFields & \hfuzz=500pt uint & \hfuzz=500pt start\\
\hfuzz=500pt\includegraphics[width=1em]{element.pdf}~countDataFields & \hfuzz=500pt uint & \hfuzz=500pt number of data fields (default: all after start)\\
\hline
\end{tabularx}

\clearpage
%==================================
\subsection{NoiseInstrument}\label{NoiseInstrument}
This program adds noise to \file{instrument data}{instrument}.
See \configClass{noiseGenerator}{noiseGeneratorType} for details on noise generation.


\keepXColumns
\begin{tabularx}{\textwidth}{N T A}
\hline
Name & Type & Annotation\\
\hline
\hfuzz=500pt\includegraphics[width=1em]{element-mustset.pdf}~outputfileInstrument & \hfuzz=500pt filename & \hfuzz=500pt \\
\hfuzz=500pt\includegraphics[width=1em]{element-mustset.pdf}~inputfileInstrument & \hfuzz=500pt filename & \hfuzz=500pt \\
\hfuzz=500pt\includegraphics[width=1em]{element-mustset-unbounded.pdf}~noise & \hfuzz=500pt \hyperref[noiseGeneratorType]{noiseGenerator} & \hfuzz=500pt \\
\hfuzz=500pt\includegraphics[width=1em]{element.pdf}~startDataFields & \hfuzz=500pt uint & \hfuzz=500pt start\\
\hfuzz=500pt\includegraphics[width=1em]{element.pdf}~countDataFields & \hfuzz=500pt uint & \hfuzz=500pt number of data fields (default: all after start)\\
\hline
\end{tabularx}

This program is \reference{parallelized}{general.parallelization}.
\clearpage
%==================================
\subsection{NoiseNormalsSolution}\label{NoiseNormalsSolution}
The inverse of the normal matrix of \configFile{inputfileNormalEquation}{normalEquation}
represents the covariance matrix of the estimated parameters. This program generates a noise vector with
\begin{equation}
\M\Sigma(\M e) = \M N^{-1},
\end{equation}
if generated input noise is standard white noise.

The noise vector is computed with
\begin{equation}
\M e = \M W^{-T} \M z,
\end{equation}
where $\M z$ is the generated \configClass{noise}{noiseGeneratorType} and
$\M W$ is the cholesky upper triangle matrix of the normal matrix $\M N=\M W^T\M W$.


\keepXColumns
\begin{tabularx}{\textwidth}{N T A}
\hline
Name & Type & Annotation\\
\hline
\hfuzz=500pt\includegraphics[width=1em]{element-mustset.pdf}~outputfileNoise & \hfuzz=500pt filename & \hfuzz=500pt generated noise as matrix: parameterCount x sampleCount\\
\hfuzz=500pt\includegraphics[width=1em]{element-mustset.pdf}~inputfileNormalEquation & \hfuzz=500pt filename & \hfuzz=500pt \\
\hfuzz=500pt\includegraphics[width=1em]{element-mustset-unbounded.pdf}~noise & \hfuzz=500pt \hyperref[noiseGeneratorType]{noiseGenerator} & \hfuzz=500pt \\
\hfuzz=500pt\includegraphics[width=1em]{element.pdf}~sampleCount & \hfuzz=500pt uint & \hfuzz=500pt number of samples to be generated\\
\hfuzz=500pt\includegraphics[width=1em]{element.pdf}~useEigenvalueDecomposition & \hfuzz=500pt boolean & \hfuzz=500pt use eigenvalue decomposition\\
\hline
\end{tabularx}

This program is \reference{parallelized}{general.parallelization}.
\clearpage
%==================================
\subsection{NoiseOrbit}\label{NoiseOrbit}
This program adds noise to simulated \file{satellite}{instrument}'s positions
and velocities generated by \program{SimulateOrbit} (along, cross, radial).
See \configClass{noiseGenerator}{noiseGeneratorType} for details on noise options.


\keepXColumns
\begin{tabularx}{\textwidth}{N T A}
\hline
Name & Type & Annotation\\
\hline
\hfuzz=500pt\includegraphics[width=1em]{element-mustset.pdf}~outputfileOrbit & \hfuzz=500pt filename & \hfuzz=500pt \\
\hfuzz=500pt\includegraphics[width=1em]{element-mustset.pdf}~inputfileOrbit & \hfuzz=500pt filename & \hfuzz=500pt \\
\hfuzz=500pt\includegraphics[width=1em]{element-unbounded.pdf}~noisePosition & \hfuzz=500pt \hyperref[noiseGeneratorType]{noiseGenerator} & \hfuzz=500pt along, cross, radial [m]\\
\hfuzz=500pt\includegraphics[width=1em]{element-unbounded.pdf}~noiseVelocity & \hfuzz=500pt \hyperref[noiseGeneratorType]{noiseGenerator} & \hfuzz=500pt along, cross, radial [m/s]\\
\hline
\end{tabularx}

This program is \reference{parallelized}{general.parallelization}.
\clearpage
%==================================
\subsection{NoiseSatelliteTracking}\label{NoiseSatelliteTracking}
This program adds noise to simulated satellite tracking data generated by \program{SimulateSatelliteTracking}.
See \configClass{noiseGenerator}{noiseGeneratorType} for details on noise generation.


\keepXColumns
\begin{tabularx}{\textwidth}{N T A}
\hline
Name & Type & Annotation\\
\hline
\hfuzz=500pt\includegraphics[width=1em]{element-mustset.pdf}~outputfileSatelliteTracking & \hfuzz=500pt filename & \hfuzz=500pt \\
\hfuzz=500pt\includegraphics[width=1em]{element-mustset.pdf}~inputfileSatelliteTracking & \hfuzz=500pt filename & \hfuzz=500pt \\
\hfuzz=500pt\includegraphics[width=1em]{element-unbounded.pdf}~noiseRange & \hfuzz=500pt \hyperref[noiseGeneratorType]{noiseGenerator} & \hfuzz=500pt [m]\\
\hfuzz=500pt\includegraphics[width=1em]{element-unbounded.pdf}~noiseRangeRate & \hfuzz=500pt \hyperref[noiseGeneratorType]{noiseGenerator} & \hfuzz=500pt [m/s]\\
\hfuzz=500pt\includegraphics[width=1em]{element-unbounded.pdf}~noiseRangeAcceleration & \hfuzz=500pt \hyperref[noiseGeneratorType]{noiseGenerator} & \hfuzz=500pt [m/s\textasciicircum{}2]\\
\hline
\end{tabularx}

This program is \reference{parallelized}{general.parallelization}.
\clearpage
%==================================
\subsection{NoiseStarCamera}\label{NoiseStarCamera}
This program adds noise to rotation observations. The noise is computed via a pseudo random sequence.
See \configClass{noiseGenerator}{noiseGeneratorType} for details on noise options.


\keepXColumns
\begin{tabularx}{\textwidth}{N T A}
\hline
Name & Type & Annotation\\
\hline
\hfuzz=500pt\includegraphics[width=1em]{element-mustset.pdf}~outputfileStarCamera & \hfuzz=500pt filename & \hfuzz=500pt \\
\hfuzz=500pt\includegraphics[width=1em]{element-mustset.pdf}~inputfileStarCamera & \hfuzz=500pt filename & \hfuzz=500pt \\
\hfuzz=500pt\includegraphics[width=1em]{element-mustset-unbounded.pdf}~noiseRoll & \hfuzz=500pt \hyperref[noiseGeneratorType]{noiseGenerator} & \hfuzz=500pt [rad]\\
\hfuzz=500pt\includegraphics[width=1em]{element-mustset-unbounded.pdf}~noisePitch & \hfuzz=500pt \hyperref[noiseGeneratorType]{noiseGenerator} & \hfuzz=500pt [rad]\\
\hfuzz=500pt\includegraphics[width=1em]{element-mustset-unbounded.pdf}~noiseYaw & \hfuzz=500pt \hyperref[noiseGeneratorType]{noiseGenerator} & \hfuzz=500pt [rad]\\
\hline
\end{tabularx}

This program is \reference{parallelized}{general.parallelization}.
\clearpage
%==================================
\subsection{NoiseTimeSeries}\label{NoiseTimeSeries}
This program generates \configFile{outputfileNoise}{instrument} with the requested characteristics.
See \configClass{noiseGenerator}{noiseGeneratorType} for details on noise options.


\keepXColumns
\begin{tabularx}{\textwidth}{N T A}
\hline
Name & Type & Annotation\\
\hline
\hfuzz=500pt\includegraphics[width=1em]{element-mustset.pdf}~outputfileNoise & \hfuzz=500pt filename & \hfuzz=500pt \\
\hfuzz=500pt\includegraphics[width=1em]{element.pdf}~outputfileCovarianceFunction & \hfuzz=500pt filename & \hfuzz=500pt \\
\hfuzz=500pt\includegraphics[width=1em]{element-mustset-unbounded.pdf}~noise & \hfuzz=500pt \hyperref[noiseGeneratorType]{noiseGenerator} & \hfuzz=500pt \\
\hfuzz=500pt\includegraphics[width=1em]{element-mustset-unbounded.pdf}~timeSeries & \hfuzz=500pt \hyperref[timeSeriesType]{timeSeries} & \hfuzz=500pt \\
\hfuzz=500pt\includegraphics[width=1em]{element-mustset.pdf}~columns & \hfuzz=500pt uint & \hfuzz=500pt number of noise series (columns)\\
\hline
\end{tabularx}

\clearpage
%==================================
\subsection{SimulateAccelerometer}\label{SimulateAccelerometer}
This program simulate \file{accelerometer data}{instrument}. The orientation of the accelerometer
is given by \configFile{inputfileStarCamera}{instrument} otherwise the celestial reference frame (CRF) is used.
For computation of non-conservative forces a \configFile{satelliteModel}{satelliteModel} is needed.


\keepXColumns
\begin{tabularx}{\textwidth}{N T A}
\hline
Name & Type & Annotation\\
\hline
\hfuzz=500pt\includegraphics[width=1em]{element-mustset.pdf}~outputfileAccelerometer & \hfuzz=500pt filename & \hfuzz=500pt \\
\hfuzz=500pt\includegraphics[width=1em]{element.pdf}~inputfileSatelliteModel & \hfuzz=500pt filename & \hfuzz=500pt satellite macro model\\
\hfuzz=500pt\includegraphics[width=1em]{element-mustset.pdf}~inputfileOrbit & \hfuzz=500pt filename & \hfuzz=500pt \\
\hfuzz=500pt\includegraphics[width=1em]{element.pdf}~inputfileStarCamera & \hfuzz=500pt filename & \hfuzz=500pt \\
\hfuzz=500pt\includegraphics[width=1em]{element-mustset.pdf}~earthRotation & \hfuzz=500pt \hyperref[earthRotationType]{earthRotation} & \hfuzz=500pt \\
\hfuzz=500pt\includegraphics[width=1em]{element.pdf}~ephemerides & \hfuzz=500pt \hyperref[ephemeridesType]{ephemerides} & \hfuzz=500pt \\
\hfuzz=500pt\includegraphics[width=1em]{element-mustset.pdf}~forces & \hfuzz=500pt \hyperref[forcesType]{forces} & \hfuzz=500pt \\
\hline
\end{tabularx}

This program is \reference{parallelized}{general.parallelization}.
\clearpage
%==================================
\subsection{SimulateAccelerometerCoMOffset}\label{SimulateAccelerometerCoMOffset}
This program generates an \file{accelerometer file}{instrument} containing perturbing accelerations
due to a given center of mass (CoM) offset. This includes centrifugal effects,
Euler forces and the effect of gravity gradients.


\keepXColumns
\begin{tabularx}{\textwidth}{N T A}
\hline
Name & Type & Annotation\\
\hline
\hfuzz=500pt\includegraphics[width=1em]{element-mustset.pdf}~outputfileAccelerometer & \hfuzz=500pt filename & \hfuzz=500pt effect of offset\\
\hfuzz=500pt\includegraphics[width=1em]{element-mustset.pdf}~inputfileOrbit & \hfuzz=500pt filename & \hfuzz=500pt \\
\hfuzz=500pt\includegraphics[width=1em]{element-mustset.pdf}~inputfileStarCamera & \hfuzz=500pt filename & \hfuzz=500pt \\
\hfuzz=500pt\includegraphics[width=1em]{element.pdf}~applyAngularRate & \hfuzz=500pt boolean & \hfuzz=500pt compute effect of centrifugal forces\\
\hfuzz=500pt\includegraphics[width=1em]{element.pdf}~applyAngularAccelerations & \hfuzz=500pt boolean & \hfuzz=500pt compute effect of Euler forces\\
\hfuzz=500pt\includegraphics[width=1em]{element-unbounded.pdf}~gradientfield & \hfuzz=500pt \hyperref[gravityfieldType]{gravityfield} & \hfuzz=500pt low order field to estimate the change of the gravity by position adjustement\\
\hfuzz=500pt\includegraphics[width=1em]{element-mustset.pdf}~earthRotation & \hfuzz=500pt \hyperref[earthRotationType]{earthRotation} & \hfuzz=500pt \\
\hfuzz=500pt\includegraphics[width=1em]{element.pdf}~interpolationDegree & \hfuzz=500pt uint & \hfuzz=500pt derivation of quaternions by polynomial interpolation of degree n\\
\hfuzz=500pt\includegraphics[width=1em]{element-mustset.pdf}~CoMOffsetX & \hfuzz=500pt double & \hfuzz=500pt offset [m]\\
\hfuzz=500pt\includegraphics[width=1em]{element-mustset.pdf}~CoMOffsetY & \hfuzz=500pt double & \hfuzz=500pt offset [m]\\
\hfuzz=500pt\includegraphics[width=1em]{element-mustset.pdf}~CoMOffsetZ & \hfuzz=500pt double & \hfuzz=500pt offset [m]\\
\hline
\end{tabularx}

This program is \reference{parallelized}{general.parallelization}.
\clearpage
%==================================
\subsection{SimulateGradiometer}\label{SimulateGradiometer}
This program simulates error free \file{gradiometer data}{instrument} along a satellite's orbit.
The orientation of the full tensor gradiometer is given by \configFile{inputfileStarCamera}{instrument}
otherwise the celestial reference frame (CRF) is used.
The gravity gradients are given by \configClass{gravityfield}{gravityfieldType} and
\configClass{tides}{tidesType}.


\keepXColumns
\begin{tabularx}{\textwidth}{N T A}
\hline
Name & Type & Annotation\\
\hline
\hfuzz=500pt\includegraphics[width=1em]{element-mustset.pdf}~outputfileGradiometer & \hfuzz=500pt filename & \hfuzz=500pt \\
\hfuzz=500pt\includegraphics[width=1em]{element-mustset.pdf}~inputfileOrbit & \hfuzz=500pt filename & \hfuzz=500pt \\
\hfuzz=500pt\includegraphics[width=1em]{element.pdf}~inputfileStarCamera & \hfuzz=500pt filename & \hfuzz=500pt \\
\hfuzz=500pt\includegraphics[width=1em]{element-mustset.pdf}~earthRotation & \hfuzz=500pt \hyperref[earthRotationType]{earthRotation} & \hfuzz=500pt \\
\hfuzz=500pt\includegraphics[width=1em]{element.pdf}~ephemerides & \hfuzz=500pt \hyperref[ephemeridesType]{ephemerides} & \hfuzz=500pt \\
\hfuzz=500pt\includegraphics[width=1em]{element-unbounded.pdf}~gravityfield & \hfuzz=500pt \hyperref[gravityfieldType]{gravityfield} & \hfuzz=500pt \\
\hfuzz=500pt\includegraphics[width=1em]{element-unbounded.pdf}~tides & \hfuzz=500pt \hyperref[tidesType]{tides} & \hfuzz=500pt \\
\hline
\end{tabularx}

This program is \reference{parallelized}{general.parallelization}.
\clearpage
%==================================
\subsection{SimulateKeplerOrbit}\label{SimulateKeplerOrbit}
This program simulates a Keplerian \file{orbit}{instrument} at a given \config{timeSeries}
starting from the given \config{integrationConstants}.


\keepXColumns
\begin{tabularx}{\textwidth}{N T A}
\hline
Name & Type & Annotation\\
\hline
\hfuzz=500pt\includegraphics[width=1em]{element-mustset.pdf}~outputfileOrbit & \hfuzz=500pt filename & \hfuzz=500pt \\
\hfuzz=500pt\includegraphics[width=1em]{element-mustset-unbounded.pdf}~timeSeries & \hfuzz=500pt \hyperref[timeSeriesType]{timeSeries} & \hfuzz=500pt \\
\hfuzz=500pt\includegraphics[width=1em]{element.pdf}~GM & \hfuzz=500pt double & \hfuzz=500pt Geocentric gravitational constant\\
\hfuzz=500pt\includegraphics[width=1em]{element-mustset.pdf}~integrationConstants & \hfuzz=500pt choice & \hfuzz=500pt \\
\hfuzz=500pt\includegraphics[width=1em]{connector.pdf}\includegraphics[width=1em]{element-mustset.pdf}~kepler & \hfuzz=500pt sequence & \hfuzz=500pt \\
\hfuzz=500pt\quad\includegraphics[width=1em]{connector.pdf}\includegraphics[width=1em]{element-mustset.pdf}~majorAxis & \hfuzz=500pt double & \hfuzz=500pt [m]\\
\hfuzz=500pt\quad\includegraphics[width=1em]{connector.pdf}\includegraphics[width=1em]{element-mustset.pdf}~eccentricity & \hfuzz=500pt double & \hfuzz=500pt [-]\\
\hfuzz=500pt\quad\includegraphics[width=1em]{connector.pdf}\includegraphics[width=1em]{element-mustset.pdf}~inclination & \hfuzz=500pt angle & \hfuzz=500pt [degree]\\
\hfuzz=500pt\quad\includegraphics[width=1em]{connector.pdf}\includegraphics[width=1em]{element-mustset.pdf}~ascendingNode & \hfuzz=500pt angle & \hfuzz=500pt [degree]\\
\hfuzz=500pt\quad\includegraphics[width=1em]{connector.pdf}\includegraphics[width=1em]{element-mustset.pdf}~argumentOfPerigee & \hfuzz=500pt angle & \hfuzz=500pt [degree]\\
\hfuzz=500pt\quad\includegraphics[width=1em]{connector.pdf}\includegraphics[width=1em]{element-mustset.pdf}~meanAnomaly & \hfuzz=500pt angle & \hfuzz=500pt [degree]\\
\hfuzz=500pt\quad\includegraphics[width=1em]{connector.pdf}\includegraphics[width=1em]{element-mustset.pdf}~time & \hfuzz=500pt time & \hfuzz=500pt integration constants are valid at this epoch\\
\hfuzz=500pt\includegraphics[width=1em]{connector.pdf}\includegraphics[width=1em]{element-mustset.pdf}~positionAndVelocity & \hfuzz=500pt sequence & \hfuzz=500pt \\
\hfuzz=500pt\quad\includegraphics[width=1em]{connector.pdf}\includegraphics[width=1em]{element-mustset.pdf}~position0x & \hfuzz=500pt double & \hfuzz=500pt [m] in CRF\\
\hfuzz=500pt\quad\includegraphics[width=1em]{connector.pdf}\includegraphics[width=1em]{element-mustset.pdf}~position0y & \hfuzz=500pt double & \hfuzz=500pt [m] in CRF\\
\hfuzz=500pt\quad\includegraphics[width=1em]{connector.pdf}\includegraphics[width=1em]{element-mustset.pdf}~position0z & \hfuzz=500pt double & \hfuzz=500pt [m] in CRF\\
\hfuzz=500pt\quad\includegraphics[width=1em]{connector.pdf}\includegraphics[width=1em]{element-mustset.pdf}~velocity0x & \hfuzz=500pt double & \hfuzz=500pt [m/s]\\
\hfuzz=500pt\quad\includegraphics[width=1em]{connector.pdf}\includegraphics[width=1em]{element-mustset.pdf}~velocity0y & \hfuzz=500pt double & \hfuzz=500pt [m/s]\\
\hfuzz=500pt\quad\includegraphics[width=1em]{connector.pdf}\includegraphics[width=1em]{element-mustset.pdf}~velocity0z & \hfuzz=500pt double & \hfuzz=500pt [m/s]\\
\hfuzz=500pt\quad\includegraphics[width=1em]{connector.pdf}\includegraphics[width=1em]{element-mustset.pdf}~time & \hfuzz=500pt time & \hfuzz=500pt integration constants are valid at this epoch\\
\hline
\end{tabularx}

\clearpage
%==================================
\subsection{SimulateOrbit}\label{SimulateOrbit}
This program integrates an \file{orbit}{instrument} from a given force function (dynamic orbit).
The force functions are given by \configClass{forces}{forcesType}.
For computation of non-conservative forces a \file{satelliteModel}{satelliteModel} is needed.
The integration method must be selected with \configClass{propagator}{orbitPropagatorType}.
Because the orbit data are calculated in the celestial reference frame (CRF) you need
\configClass{earthRotation}{earthRotationType} to transform the force function
from the terrestrial reference frame (TRF).
The integration start and end time, as well as the sampling, are derived from
the \config{timeSeries} option. It is possible to integrate the arc in \config{reverse},
where the initial conditions are assumed to be met at the end time of the \config{timeSeries}.


\keepXColumns
\begin{tabularx}{\textwidth}{N T A}
\hline
Name & Type & Annotation\\
\hline
\hfuzz=500pt\includegraphics[width=1em]{element-mustset.pdf}~outputfileOrbit & \hfuzz=500pt filename & \hfuzz=500pt orbit file to be written.\\
\hfuzz=500pt\includegraphics[width=1em]{element.pdf}~inputfileSatelliteModel & \hfuzz=500pt filename & \hfuzz=500pt satellite macro model\\
\hfuzz=500pt\includegraphics[width=1em]{element-mustset-unbounded.pdf}~timeSeries & \hfuzz=500pt \hyperref[timeSeriesType]{timeSeries} & \hfuzz=500pt time points for simulated orbit epochs.\\
\hfuzz=500pt\includegraphics[width=1em]{element-mustset.pdf}~integrationConstants & \hfuzz=500pt choice & \hfuzz=500pt \\
\hfuzz=500pt\includegraphics[width=1em]{connector.pdf}\includegraphics[width=1em]{element-mustset.pdf}~kepler & \hfuzz=500pt sequence & \hfuzz=500pt \\
\hfuzz=500pt\quad\includegraphics[width=1em]{connector.pdf}\includegraphics[width=1em]{element-mustset.pdf}~majorAxis & \hfuzz=500pt double & \hfuzz=500pt [m]\\
\hfuzz=500pt\quad\includegraphics[width=1em]{connector.pdf}\includegraphics[width=1em]{element-mustset.pdf}~eccentricity & \hfuzz=500pt double & \hfuzz=500pt [-]\\
\hfuzz=500pt\quad\includegraphics[width=1em]{connector.pdf}\includegraphics[width=1em]{element-mustset.pdf}~inclination & \hfuzz=500pt angle & \hfuzz=500pt [degree]\\
\hfuzz=500pt\quad\includegraphics[width=1em]{connector.pdf}\includegraphics[width=1em]{element-mustset.pdf}~ascendingNode & \hfuzz=500pt angle & \hfuzz=500pt [degree]\\
\hfuzz=500pt\quad\includegraphics[width=1em]{connector.pdf}\includegraphics[width=1em]{element-mustset.pdf}~argumentOfPerigee & \hfuzz=500pt angle & \hfuzz=500pt [degree]\\
\hfuzz=500pt\quad\includegraphics[width=1em]{connector.pdf}\includegraphics[width=1em]{element-mustset.pdf}~meanAnomaly & \hfuzz=500pt angle & \hfuzz=500pt [degree]\\
\hfuzz=500pt\quad\includegraphics[width=1em]{connector.pdf}\includegraphics[width=1em]{element.pdf}~GM & \hfuzz=500pt double & \hfuzz=500pt Geocentric gravitational constant\\
\hfuzz=500pt\includegraphics[width=1em]{connector.pdf}\includegraphics[width=1em]{element-mustset.pdf}~positionAndVelocity & \hfuzz=500pt sequence & \hfuzz=500pt \\
\hfuzz=500pt\quad\includegraphics[width=1em]{connector.pdf}\includegraphics[width=1em]{element-mustset.pdf}~position0x & \hfuzz=500pt double & \hfuzz=500pt [m] in CRF\\
\hfuzz=500pt\quad\includegraphics[width=1em]{connector.pdf}\includegraphics[width=1em]{element-mustset.pdf}~position0y & \hfuzz=500pt double & \hfuzz=500pt [m] in CRF\\
\hfuzz=500pt\quad\includegraphics[width=1em]{connector.pdf}\includegraphics[width=1em]{element-mustset.pdf}~position0z & \hfuzz=500pt double & \hfuzz=500pt [m] in CRF\\
\hfuzz=500pt\quad\includegraphics[width=1em]{connector.pdf}\includegraphics[width=1em]{element-mustset.pdf}~velocity0x & \hfuzz=500pt double & \hfuzz=500pt [m/s]\\
\hfuzz=500pt\quad\includegraphics[width=1em]{connector.pdf}\includegraphics[width=1em]{element-mustset.pdf}~velocity0y & \hfuzz=500pt double & \hfuzz=500pt [m/s]\\
\hfuzz=500pt\quad\includegraphics[width=1em]{connector.pdf}\includegraphics[width=1em]{element-mustset.pdf}~velocity0z & \hfuzz=500pt double & \hfuzz=500pt [m/s]\\
\hfuzz=500pt\includegraphics[width=1em]{connector.pdf}\includegraphics[width=1em]{element-mustset.pdf}~file & \hfuzz=500pt sequence & \hfuzz=500pt \\
\hfuzz=500pt\quad\includegraphics[width=1em]{connector.pdf}\includegraphics[width=1em]{element-mustset.pdf}~inputfileOrbit & \hfuzz=500pt filename & \hfuzz=500pt only epoch at timeStart is used\\
\hfuzz=500pt\quad\includegraphics[width=1em]{connector.pdf}\includegraphics[width=1em]{element.pdf}~margin & \hfuzz=500pt double & \hfuzz=500pt [seconds] used when finding initial epoch in orbitFile\\
\hfuzz=500pt\includegraphics[width=1em]{element-mustset.pdf}~propagator & \hfuzz=500pt \hyperref[orbitPropagatorType]{orbitPropagator} & \hfuzz=500pt orbit propagation method.\\
\hfuzz=500pt\includegraphics[width=1em]{element-mustset.pdf}~earthRotation & \hfuzz=500pt \hyperref[earthRotationType]{earthRotation} & \hfuzz=500pt \\
\hfuzz=500pt\includegraphics[width=1em]{element.pdf}~ephemerides & \hfuzz=500pt \hyperref[ephemeridesType]{ephemerides} & \hfuzz=500pt \\
\hfuzz=500pt\includegraphics[width=1em]{element-mustset.pdf}~forces & \hfuzz=500pt \hyperref[forcesType]{forces} & \hfuzz=500pt considered in orbit propagation.\\
\hfuzz=500pt\includegraphics[width=1em]{element.pdf}~reverse & \hfuzz=500pt boolean & \hfuzz=500pt start integration at last epoch in timeSeries, going backward in time.\\
\hline
\end{tabularx}

\clearpage
%==================================
\subsection{SimulateSatelliteTracking}\label{SimulateSatelliteTracking}
This program simulates \file{tracking data}{instrument} (range, range-rate, range-accelerations)
between 2 satellites. The range is given by
\begin{equation}
\rho(t) = \left\lVert{\M r_B(t) - \M r_A(t)}\right\rVert = \M e_{AB}(t)\cdot\M r_{AB}(t),
\end{equation}
with $\M r_{AB} = \M r_B - \M r_A$ and the unit vector in line of sight (LOS) direction
\begin{equation}\label{sst.los}
\M e_{AB} = \frac{\M r_{AB}}{\left\lVert{\M r_{AB}}\right\rVert}=\frac{\M r_{AB}}{\rho}.
\end{equation}
Range-rates~$\dot{\rho}$ and range accelrations~$\ddot{\rho}$ are obtained by differentation
\begin{equation}\label{obsRangeRate}
\dot{\rho}  = \M e_{AB}\cdot\dot{\M r}_{AB} + \dot{\M e}_{AB}\cdot\M r_{AB}
            = \M e_{AB}\cdot\dot{\M r}_{AB},
\end{equation}
\begin{equation}\label{obsRangeAccl}
\begin{split}
\ddot{\rho} &= \M e_{AB}\cdot\ddot{\M r}_{AB} +\dot{\M e}_{AB}\cdot\dot{\M r}_{AB}
            = \M e_{AB}\cdot\ddot{\M r}_{AB} +
   \frac{1}{\rho}\left(\dot{\M r}_{AB}^2-\dot{\rho}^2\right). \\
\end{split}
\end{equation}
with the derivative of the unit vector
\begin{equation}
\dot{\M e}_{AB}=\frac{d}{dt}\left(\frac{\M r_{AB}}{\rho}\right)
=\frac{\dot{\M r}_{AB}}{\rho}-\frac{\dot{\rho}\cdot\M r_{AB}}{\rho^2}
=\frac{1}{\rho}\left({\dot{\M r}_{AB}-\dot{\rho}\cdot\M e_{AB}}\right).
\end{equation}
The \configFile{inputfileOrbit}{instrument}s must contain positions, velocities, and acceleration
(see \program{OrbitAddVelocityAndAcceleration}).


\keepXColumns
\begin{tabularx}{\textwidth}{N T A}
\hline
Name & Type & Annotation\\
\hline
\hfuzz=500pt\includegraphics[width=1em]{element-mustset.pdf}~outputfileSatelliteTracking & \hfuzz=500pt filename & \hfuzz=500pt \\
\hfuzz=500pt\includegraphics[width=1em]{element-mustset.pdf}~inputfileOrbit1 & \hfuzz=500pt filename & \hfuzz=500pt \\
\hfuzz=500pt\includegraphics[width=1em]{element-mustset.pdf}~inputfileOrbit2 & \hfuzz=500pt filename & \hfuzz=500pt \\
\hline
\end{tabularx}

This program is \reference{parallelized}{general.parallelization}.
\clearpage
%==================================
\subsection{SimulateStarCamera}\label{SimulateStarCamera}
This program simulates \file{star camera}{instrument} measurements at each satellite's position.
The satellite's orientation follows a local orbit frame with the x-axis in along track (along velocity),
y-axis is cross track (normal to position and velocity vector) and z-axis pointing nadir (negative position vector).
As for non circular orbit the position and velocity are not exact normal, the default is the x-axis to be exact
along velocity and the z-axis forms a right hand system (not exact nadir) or with \config{nadirPointing} the z-axis
is exact nadir and x-axis approximates along.
The resulting rotation matrices rotate from satellite frame to inertial frame.


\keepXColumns
\begin{tabularx}{\textwidth}{N T A}
\hline
Name & Type & Annotation\\
\hline
\hfuzz=500pt\includegraphics[width=1em]{element-mustset.pdf}~outputfileStarCamera & \hfuzz=500pt filename & \hfuzz=500pt rotation from satellite to inertial frame (x: along, y: cross, z: nadir)\\
\hfuzz=500pt\includegraphics[width=1em]{element-mustset.pdf}~inputfileOrbit & \hfuzz=500pt filename & \hfuzz=500pt position and velocity defines the orientation of the satellite at each epoch\\
\hfuzz=500pt\includegraphics[width=1em]{element.pdf}~nadirPointing & \hfuzz=500pt boolean & \hfuzz=500pt false: exact along and nearly nadir, true: nearly along and exact nadir\\
\hline
\end{tabularx}

This program is \reference{parallelized}{general.parallelization}.
\clearpage
%==================================
\subsection{SimulateStarCameraGnss}\label{SimulateStarCameraGnss}
This program simulates \file{star camera}{instrument} measurements at each satellite position
of \configFile{inputfileOrbit}{instrument}.
The resulting rotation matrices rotate from body frame to inertial frame. The body frame refers
to the IGS-specific (not the manufacturer-specific) body frame, as described by
\href{https://doi.org/10.1016/j.asr.2015.06.019}{Montenbruck et al. (2015)}.
The \configFile{inputfileOrbit}{instrument} must contain velocities
(use \program{OrbitAddVelocityAndAcceleration} if needed).

Information about the attitude mode(s) used by the GNSS satellite may be provided via
\configFile{inputfileAttitudeInfo}{instrument}. This file can be created with
\program{GnssAttitudeInfoCreate}. It contains one or more time-dependent entries,
each defining the default attitude mode, the attitude modes used around orbit noon and
midnight, and some parameters required by the various modes.
If no \configFile{inputfileAttitudeInfo}{instrument} is selected, the program defaults
to a nominal yaw-steering attitude model.
A sufficiently high \config{modelingResolution} ensures that the attitude behavior is modeled properly
at all times.

The attitude behavior is defined by the respective mode. Here is a list of the supported
modes with a brief explanation and references:
\begin{itemize}
\item \textbf{nominalYawSteering}:
      Yaw to keep solar panels aligned to Sun (e.g. most GNSS satellites outside eclipse) [1]
\item \textbf{orbitNormal}:
      Keep fixed yaw angle, for example point X-axis in flight direction (e.g. BDS-2G, BDS-3G, QZS-2G) [1]
\item \textbf{catchUpYawSteering}:
      Yaw at maximum yaw rate to catch up to nominal yaw angle (e.g. GPS-* (noon), GPS-IIR (midnight)) [2, 3]
\item \textbf{shadowMaxYawSteeringAndRecovery}:
      Yaw at maximum yaw rate from shadow start to end, recover after shadow (e.g. GPS-IIA (midnight)) [2]
\item \textbf{shadowMaxYawSteeringAndStop}:
      Yaw at maximum yaw rate from shadow start until nominal yaw angle at shadow end is reached,
      then stop (e.g. GLO-M (midnight)) [4]
\item \textbf{shadowConstantYawSteering}:
      Yaw at constant yaw rate from shadow start to end (e.g. GPS-IIF (midnight)) [3]
\item \textbf{centeredMaxYawSteering}:
      Yaw at maximum yaw rate centered around noon/midnight (e.g. QZS-2I, GLO-M (noon)) [4, 8]
\item \textbf{smoothedYawSteering1}:
      Yaw based on an auxiliary Sun vector for a smooth yaw maneuver (e.g. GAL-1) [5]
\item \textbf{smoothedYawSteering2}:
      Yaw based on a modified yaw-steering law for a smooth yaw maneuver (e.g. GAL-2, BDS-3M, BDS-3I) [5, 6]
\item \textbf{betaDependentOrbitNormal}:
      Switch to orbit normal mode if below beta angle threshold (e.g. BDS-2M, BDS-2I, QZS-1) [7, 8]
\end{itemize}

\fig{!hb}{0.9}{gnssAttitudeModes}{fig:gnssAttitudeModes1}{Overview of attitude modes used by GNSS satellites}

See \program{GnssAttitudeInfoCreate} for more details on which satellite uses which attitude modes
and the required parameters for each mode.

References for the attitude modes:
\begin{enumerate}
\item \href{https://doi.org/10.1016/j.asr.2015.06.019}{Montenbruck et al. (2015)}
\item \href{https://doi.org/10.1007/s10291-008-0092-1}{Kouba (2009)}
\item \href{https://doi.org/10.1007/s10291-016-0562-9}{Kuang et al. (2017)}
\item \href{https://doi.org/10.1016/j.asr.2010.09.007}{Dilssner et al. (2011)}
\item \url{https://www.gsc-europa.eu/support-to-developers/galileo-satellite-metadata#3}
\item \href{https://doi.org/10.1007/s10291-018-0783-1}{Wang et al. (2018)}
\item \href{https://doi.org/10.1017/S0373463318000103}{Li et al. (2018)}
\item \url{https://qzss.go.jp/en/technical/qzssinfo/index.html}
\end{enumerate}


\keepXColumns
\begin{tabularx}{\textwidth}{N T A}
\hline
Name & Type & Annotation\\
\hline
\hfuzz=500pt\includegraphics[width=1em]{element-mustset.pdf}~outputfileStarCamera & \hfuzz=500pt filename & \hfuzz=500pt rotation from body frame to CRF\\
\hfuzz=500pt\includegraphics[width=1em]{element-mustset.pdf}~inputfileOrbit & \hfuzz=500pt filename & \hfuzz=500pt attitude is modeled based on this orbit\\
\hfuzz=500pt\includegraphics[width=1em]{element.pdf}~inputfileAttitudeInfo & \hfuzz=500pt filename & \hfuzz=500pt attitude modes used by the satellite and respective parameters\\
\hfuzz=500pt\includegraphics[width=1em]{element.pdf}~interpolationDegree & \hfuzz=500pt uint & \hfuzz=500pt polynomial degree for orbit interpolation\\
\hfuzz=500pt\includegraphics[width=1em]{element.pdf}~modelingResolution & \hfuzz=500pt double & \hfuzz=500pt [s] resolution for attitude model evaluation\\
\hfuzz=500pt\includegraphics[width=1em]{element-mustset.pdf}~ephemerides & \hfuzz=500pt \hyperref[ephemeridesType]{ephemerides} & \hfuzz=500pt \\
\hfuzz=500pt\includegraphics[width=1em]{element-mustset.pdf}~eclipse & \hfuzz=500pt \hyperref[eclipseType]{eclipse} & \hfuzz=500pt model to determine if satellite is in Earth's shadow\\
\hline
\end{tabularx}

\clearpage
%==================================
\subsection{SimulateStarCameraGrace}\label{SimulateStarCameraGrace}
Simulates \file{star camera data}{instrument} of the two GRACE satellites.
\begin{itemize}
\item x: the antenna center pointing to the other satellite.
\item y: normal to line of sight and the radial direction.
\item z: forms a right handed system.
\end{itemize}


\keepXColumns
\begin{tabularx}{\textwidth}{N T A}
\hline
Name & Type & Annotation\\
\hline
\hfuzz=500pt\includegraphics[width=1em]{element-mustset.pdf}~outputfileStarCamera1 & \hfuzz=500pt filename & \hfuzz=500pt \\
\hfuzz=500pt\includegraphics[width=1em]{element-mustset.pdf}~outputfileStarCamera2 & \hfuzz=500pt filename & \hfuzz=500pt \\
\hfuzz=500pt\includegraphics[width=1em]{element-mustset.pdf}~inputfileOrbit1 & \hfuzz=500pt filename & \hfuzz=500pt position define the orientation of the satellite at each epoch\\
\hfuzz=500pt\includegraphics[width=1em]{element-mustset.pdf}~inputfileOrbit2 & \hfuzz=500pt filename & \hfuzz=500pt position define the orientation of the satellite at each epoch\\
\hfuzz=500pt\includegraphics[width=1em]{element.pdf}~antennaCenters & \hfuzz=500pt choice & \hfuzz=500pt KBR antenna phase center\\
\hfuzz=500pt\includegraphics[width=1em]{connector.pdf}\includegraphics[width=1em]{element-mustset.pdf}~value & \hfuzz=500pt sequence & \hfuzz=500pt \\
\hfuzz=500pt\quad\includegraphics[width=1em]{connector.pdf}\includegraphics[width=1em]{element.pdf}~center1X & \hfuzz=500pt double & \hfuzz=500pt x-coordinate of antenna position in SRF [m] for GRACEA\\
\hfuzz=500pt\quad\includegraphics[width=1em]{connector.pdf}\includegraphics[width=1em]{element.pdf}~center1Y & \hfuzz=500pt double & \hfuzz=500pt y-coordinate of antenna position in SRF [m] for GRACEA\\
\hfuzz=500pt\quad\includegraphics[width=1em]{connector.pdf}\includegraphics[width=1em]{element.pdf}~center1Z & \hfuzz=500pt double & \hfuzz=500pt z-coordinate of antenna position in SRF [m] for GRACEA\\
\hfuzz=500pt\quad\includegraphics[width=1em]{connector.pdf}\includegraphics[width=1em]{element.pdf}~center2X & \hfuzz=500pt double & \hfuzz=500pt x-coordinate of antenna position in SRF [m] for GRACEB\\
\hfuzz=500pt\quad\includegraphics[width=1em]{connector.pdf}\includegraphics[width=1em]{element.pdf}~center2Y & \hfuzz=500pt double & \hfuzz=500pt y-coordinate of antenna position in SRF [m] for GRACEB\\
\hfuzz=500pt\quad\includegraphics[width=1em]{connector.pdf}\includegraphics[width=1em]{element.pdf}~center2Z & \hfuzz=500pt double & \hfuzz=500pt z-coordinate of antenna position in SRF [m] for GRACEB\\
\hfuzz=500pt\includegraphics[width=1em]{connector.pdf}\includegraphics[width=1em]{element-mustset.pdf}~file & \hfuzz=500pt sequence & \hfuzz=500pt \\
\hfuzz=500pt\quad\includegraphics[width=1em]{connector.pdf}\includegraphics[width=1em]{element-mustset.pdf}~inputAntennaCenters & \hfuzz=500pt filename & \hfuzz=500pt \\
\hline
\end{tabularx}

\clearpage
%==================================
\subsection{SimulateStarCameraSentinel1}\label{SimulateStarCameraSentinel1}
This program simulates \file{star camera}{instrument} measurements at each satellite's position for the Sentinel 1A satellite.
The \configFile{inputfileOrbit}{instrument} must contain positions and velocities (see \program{OrbitAddVelocityAndAcceleration}).
The resulting rotation matrices rotate from satellite frame to inertial frame.


\keepXColumns
\begin{tabularx}{\textwidth}{N T A}
\hline
Name & Type & Annotation\\
\hline
\hfuzz=500pt\includegraphics[width=1em]{element-mustset.pdf}~outputfileStarCamera & \hfuzz=500pt filename & \hfuzz=500pt \\
\hfuzz=500pt\includegraphics[width=1em]{element-mustset.pdf}~inputfileOrbit & \hfuzz=500pt filename & \hfuzz=500pt position and velocity defines the orientation of the satellite at each epoch\\
\hline
\end{tabularx}

This program is \reference{parallelized}{general.parallelization}.
\clearpage
%==================================
\subsection{SimulateStarCameraTerrasar}\label{SimulateStarCameraTerrasar}
This program simulates \configFile{outputfileStarCamera}{instrument} measurements at each satellite's position for the Terrasar satellite.
The \configFile{inputfileOrbit}{instrument} must contain positions and velocities (see \program{OrbitAddVelocityAndAcceleration}).
The resulting rotation matrices rotate from satellite frame to inertial frame.

H. Fiedler, E. Boerner, J. Mittermayer and G. Krieger,
Total zero Doppler Steering-a new method for minimizing the Doppler centroid,
in IEEE Geoscience and Remote Sensing Letters, vol. 2, no. 2, pp. 141-145, April 2005, \url{https://www.doi.org/10.1109/LGRS.2005.844591}.


\keepXColumns
\begin{tabularx}{\textwidth}{N T A}
\hline
Name & Type & Annotation\\
\hline
\hfuzz=500pt\includegraphics[width=1em]{element-mustset.pdf}~outputfileStarCamera & \hfuzz=500pt filename & \hfuzz=500pt rotation from satellite to inertial frame (x: along, y: cross, z: nadir)\\
\hfuzz=500pt\includegraphics[width=1em]{element-mustset.pdf}~inputfileOrbit & \hfuzz=500pt filename & \hfuzz=500pt position and velocity defines the orientation of the satellite at each epoch\\
\hline
\end{tabularx}

This program is \reference{parallelized}{general.parallelization}.
\clearpage
%==================================
\section{Programs: Slr}
\subsection{SlrProcessing}\label{SlrProcessing}
This program processes SLR normal point or full rate observations. It calculates the linearized observation equations,
accumulates them into a system of normal equations and solves it.

To calculate observation equations from the passes, the model parameters or unknown parameters need to be
defined beforehand. These unknown parameters can be chosen arbitrarily by the user with an adequate list of defined
\configClass{parametrization}{slrParametrizationType}.
Some of the \configClass{parametrization}{slrParametrizationType} also include a priori models.

Lastly it is required to define the process flow of the SLR processing. This is accomplished
with a list of \configClass{processingSteps}{slrProcessingStepType}.
Each step is processed consecutively. Some steps allow the selection of parameters, station, or satellites,
which affects all subsequent steps.

The \configClass{timeSeries}{timeSeriesType} is used to precompute Earth rotation and station displacements
with a uniform sampling. In a second step these values are interpolated to the observation epochs.
A sampling of about 10 minutes should be adequate.

It should be noted that GROOPS uses GPS time format, but normal point/full rate data files and CPF files, provided by ILRS data centers
are given in UTC time format.
%See also \program{SlrSimulation}.


\keepXColumns
\begin{tabularx}{\textwidth}{N T A}
\hline
Name & Type & Annotation\\
\hline
\hfuzz=500pt\includegraphics[width=1em]{element-mustset-unbounded.pdf}~timeSeries & \hfuzz=500pt \hyperref[timeSeriesType]{timeSeries} & \hfuzz=500pt defines station movements and earth rotation epochs\\
\hfuzz=500pt\includegraphics[width=1em]{element-mustset-unbounded.pdf}~satellite & \hfuzz=500pt \hyperref[slrSatelliteGeneratorType]{slrSatelliteGenerator} & \hfuzz=500pt satellites\\
\hfuzz=500pt\includegraphics[width=1em]{element-mustset-unbounded.pdf}~station & \hfuzz=500pt \hyperref[slrStationGeneratorType]{slrStationGenerator} & \hfuzz=500pt ground station network\\
\hfuzz=500pt\includegraphics[width=1em]{element-mustset.pdf}~earthRotation & \hfuzz=500pt \hyperref[earthRotationType]{earthRotation} & \hfuzz=500pt apriori earth rotation\\
\hfuzz=500pt\includegraphics[width=1em]{element-mustset-unbounded.pdf}~parametrization & \hfuzz=500pt \hyperref[slrParametrizationType]{slrParametrization} & \hfuzz=500pt models and parameters\\
\hfuzz=500pt\includegraphics[width=1em]{element-mustset-unbounded.pdf}~processingStep & \hfuzz=500pt \hyperref[slrProcessingStepType]{slrProcessingStep} & \hfuzz=500pt steps are processed consecutively\\
\hline
\end{tabularx}

\clearpage
%==================================
\section{Programs: System}
\subsection{FileConvert}\label{FileConvert}
Converts GROOPS file between different file formats (ASCII, XML, binary),
see \reference{file formats}{general.fileFormat} for details.
It prints also some information about the content.
Therefore it can be used to get an idea about the content of binary files.


\keepXColumns
\begin{tabularx}{\textwidth}{N T A}
\hline
Name & Type & Annotation\\
\hline
\hfuzz=500pt\includegraphics[width=1em]{element.pdf}~outputfile & \hfuzz=500pt filename & \hfuzz=500pt GROOPS formats: .xml, .txt, .dat\\
\hfuzz=500pt\includegraphics[width=1em]{element-mustset.pdf}~inputfile & \hfuzz=500pt filename & \hfuzz=500pt GROOPS formats: .xml, .txt, .dat\\
\hline
\end{tabularx}

\clearpage
%==================================
\subsection{FileCreateDirectories}\label{FileCreateDirectories}
Creates the directory and parent directories as needed.


\keepXColumns
\begin{tabularx}{\textwidth}{N T A}
\hline
Name & Type & Annotation\\
\hline
\hfuzz=500pt\includegraphics[width=1em]{element-mustset-unbounded.pdf}~directory & \hfuzz=500pt filename & \hfuzz=500pt \\
\hline
\end{tabularx}

\clearpage
%==================================
\subsection{FileMove}\label{FileMove}
Move/rename file or directory. If the \config{outputfile} is an existing directory
the \config{inputfile} is moved into it.


\keepXColumns
\begin{tabularx}{\textwidth}{N T A}
\hline
Name & Type & Annotation\\
\hline
\hfuzz=500pt\includegraphics[width=1em]{element-mustset.pdf}~outputfile & \hfuzz=500pt filename & \hfuzz=500pt target name or directory for the move/rename\\
\hfuzz=500pt\includegraphics[width=1em]{element-mustset.pdf}~inputfile & \hfuzz=500pt filename & \hfuzz=500pt \\
\hline
\end{tabularx}

\clearpage
%==================================
\subsection{FileRemove}\label{FileRemove}
Remove files or directories.
Deletes also the content recursivley if one of \config{files} is a directory.


\keepXColumns
\begin{tabularx}{\textwidth}{N T A}
\hline
Name & Type & Annotation\\
\hline
\hfuzz=500pt\includegraphics[width=1em]{element-mustset-unbounded.pdf}~files & \hfuzz=500pt filename & \hfuzz=500pt \\
\hline
\end{tabularx}

\clearpage
%==================================
\subsection{FileTextCreate}\label{FileTextCreate}
Create text \config{outputfile} containing \config{line}s.
This program can be a powerful tool,
if the \config{line} is repeated with a \configClass{loop}{loopType}
together with the \reference{text parser}{general.parser:text}.


\keepXColumns
\begin{tabularx}{\textwidth}{N T A}
\hline
Name & Type & Annotation\\
\hline
\hfuzz=500pt\includegraphics[width=1em]{element-mustset.pdf}~outputfile & \hfuzz=500pt filename & \hfuzz=500pt \\
\hfuzz=500pt\includegraphics[width=1em]{element-unbounded.pdf}~line & \hfuzz=500pt string & \hfuzz=500pt \\
\hline
\end{tabularx}

\clearpage
%==================================
\subsection{GroupPrograms}\label{GroupPrograms}
Runs \config{program}s in a group, which can be used to structure a config file.
If \config{catchErrors} is enabled and an error occurs, the remaining \config{program}s
are skipped and execution continues with \config{errorProgram}s, in case any are defined.
Otherwise an exception is thrown.

The \config{silently} option disables the screen ouput of the \config{program}s.
With \config{outputfileLog} a log file is written for this group additional to a global log file.
This might be helpful within \program{LoopPrograms} with parallel iterations.


\keepXColumns
\begin{tabularx}{\textwidth}{N T A}
\hline
Name & Type & Annotation\\
\hline
\hfuzz=500pt\includegraphics[width=1em]{element.pdf}~outputfileLog & \hfuzz=500pt filename & \hfuzz=500pt additional log file\\
\hfuzz=500pt\includegraphics[width=1em]{element.pdf}~silently & \hfuzz=500pt boolean & \hfuzz=500pt without showing the output.\\
\hfuzz=500pt\includegraphics[width=1em]{element-unbounded.pdf}~program & \hfuzz=500pt programType & \hfuzz=500pt \\
\hfuzz=500pt\includegraphics[width=1em]{element.pdf}~catchErrors & \hfuzz=500pt sequence & \hfuzz=500pt \\
\hfuzz=500pt\includegraphics[width=1em]{connector.pdf}\includegraphics[width=1em]{element-unbounded.pdf}~errorProgram & \hfuzz=500pt programType & \hfuzz=500pt executed if an error occured\\
\hline
\end{tabularx}

This program is \reference{parallelized}{general.parallelization}.
\clearpage
%==================================
\subsection{IfPrograms}\label{IfPrograms}
Runs a list of \config{program}s if a \configClass{condition}{conditionType} is met.
Otherwise \config{elseProgram}s are executed.


\keepXColumns
\begin{tabularx}{\textwidth}{N T A}
\hline
Name & Type & Annotation\\
\hline
\hfuzz=500pt\includegraphics[width=1em]{element-mustset.pdf}~condition & \hfuzz=500pt \hyperref[conditionType]{condition} & \hfuzz=500pt \\
\hfuzz=500pt\includegraphics[width=1em]{element-unbounded.pdf}~program & \hfuzz=500pt programType & \hfuzz=500pt executed if condition evaluates to true\\
\hfuzz=500pt\includegraphics[width=1em]{element-unbounded.pdf}~elseProgram & \hfuzz=500pt programType & \hfuzz=500pt executed if condition evaluates to false\\
\hline
\end{tabularx}

This program is \reference{parallelized}{general.parallelization}.
\clearpage
%==================================
\subsection{LoopPrograms}\label{LoopPrograms}
This program runs a list of programs in a \configClass{loop}{loopType}.

If \config{continueAfterError}=\verb|yes| and an error occurs, the remaining programs in the current iteration
are skipped and the loop continues with the next iteration. Otherwise an exception is thrown.

If this program is executed on multiple processing nodes, the iterations can be computed in parallel,
see \reference{parallelization}{general.parallelization}. The first process serves as load balancer
and the other processes are assigned to iterations according to \config{processCountPerIteration}.
For example, running a loop containing three iterations on 13 processes with \config{processCountPerIteration}=\verb|4|,
runs the three iterations in parallel, with each iteration being assigned four processes.
With \config{parallelLog}=\verb|yes| all processes write output to screen and the log file.
As the output can be quite confusing in this case, running \program{GroupPrograms} with an extra \config{outputfileLog}
for each iteration (use the loop variables for the name of the log files) might be helpful.


\keepXColumns
\begin{tabularx}{\textwidth}{N T A}
\hline
Name & Type & Annotation\\
\hline
\hfuzz=500pt\includegraphics[width=1em]{element-mustset.pdf}~loop & \hfuzz=500pt \hyperref[loopType]{loop} & \hfuzz=500pt subprograms are called for every iteration\\
\hfuzz=500pt\includegraphics[width=1em]{element.pdf}~continueAfterError & \hfuzz=500pt boolean & \hfuzz=500pt continue with next iteration after error, otherwise throw exception\\
\hfuzz=500pt\includegraphics[width=1em]{element.pdf}~processCountPerIteration & \hfuzz=500pt uint & \hfuzz=500pt 0: use all processes for each iteration\\
\hfuzz=500pt\includegraphics[width=1em]{element.pdf}~parallelLog & \hfuzz=500pt boolean & \hfuzz=500pt write to screen/log file from all processing nodes in parallelized loops\\
\hfuzz=500pt\includegraphics[width=1em]{element-unbounded.pdf}~program & \hfuzz=500pt programType & \hfuzz=500pt \\
\hline
\end{tabularx}

This program is \reference{parallelized}{general.parallelization}.
\clearpage
%==================================
\subsection{RunCommand}\label{RunCommand}
Execute system \config{command}s. If \config{executeParallel} is set and
multiple \config{command}s are given they are executed in parallel at
distributed nodes, otherwise they are executed consecutively at master node only.


\keepXColumns
\begin{tabularx}{\textwidth}{N T A}
\hline
Name & Type & Annotation\\
\hline
\hfuzz=500pt\includegraphics[width=1em]{element-mustset-unbounded.pdf}~command & \hfuzz=500pt filename & \hfuzz=500pt \\
\hfuzz=500pt\includegraphics[width=1em]{element.pdf}~silently & \hfuzz=500pt boolean & \hfuzz=500pt without showing the output.\\
\hfuzz=500pt\includegraphics[width=1em]{element.pdf}~continueAfterError & \hfuzz=500pt boolean & \hfuzz=500pt continue with next command after error, otherwise throw exception\\
\hfuzz=500pt\includegraphics[width=1em]{element.pdf}~executeParallel & \hfuzz=500pt boolean & \hfuzz=500pt execute several commands in parallel\\
\hline
\end{tabularx}

This program is \reference{parallelized}{general.parallelization}.
\clearpage
%==================================
\section{Programs: Conversion}
\subsection{BerneseKinematic2Orbit}\label{BerneseKinematic2Orbit}
Read kinematic orbits in Bernese format.


\keepXColumns
\begin{tabularx}{\textwidth}{N T A}
\hline
Name & Type & Annotation\\
\hline
\hfuzz=500pt\includegraphics[width=1em]{element-mustset.pdf}~outputfileOrbit & \hfuzz=500pt filename & \hfuzz=500pt \\
\hfuzz=500pt\includegraphics[width=1em]{element-mustset.pdf}~outputfileCovariance & \hfuzz=500pt filename & \hfuzz=500pt \\
\hfuzz=500pt\includegraphics[width=1em]{element.pdf}~earthRotation & \hfuzz=500pt \hyperref[earthRotationType]{earthRotation} & \hfuzz=500pt from TRF to CRF\\
\hfuzz=500pt\includegraphics[width=1em]{element-mustset-unbounded.pdf}~inputfile & \hfuzz=500pt filename & \hfuzz=500pt \\
\hline
\end{tabularx}

\clearpage
%==================================
\subsection{Champ2AccStar}\label{Champ2AccStar}
This program reads in CHAMP accelerometer and star camera data given in the special CHAMP format.
In case of CHAMP accelerometer and star camera data is both stored in one file.
A description of the format can be found under: \url{http://op.gfz-potsdam.de/champ/docs_CHAMP/CH-GFZ-FD-001.pdf}.


\keepXColumns
\begin{tabularx}{\textwidth}{N T A}
\hline
Name & Type & Annotation\\
\hline
\hfuzz=500pt\includegraphics[width=1em]{element.pdf}~outputfileAccelerometer & \hfuzz=500pt filename & \hfuzz=500pt \\
\hfuzz=500pt\includegraphics[width=1em]{element.pdf}~outputfileAngularAcceleration & \hfuzz=500pt filename & \hfuzz=500pt \\
\hfuzz=500pt\includegraphics[width=1em]{element.pdf}~outputfileStarCamera & \hfuzz=500pt filename & \hfuzz=500pt \\
\hfuzz=500pt\includegraphics[width=1em]{element-mustset-unbounded.pdf}~inputfile & \hfuzz=500pt filename & \hfuzz=500pt \\
\hline
\end{tabularx}

\clearpage
%==================================
\subsection{Champ2Orbit}\label{Champ2Orbit}
This program reads in CHAMP precise science orbits in the special CHORB format.
A description of the format can be found under: \url{http://op.gfz-potsdam.de/champ/docs_CHAMP/CH-GFZ-FD-002.pdf}


\keepXColumns
\begin{tabularx}{\textwidth}{N T A}
\hline
Name & Type & Annotation\\
\hline
\hfuzz=500pt\includegraphics[width=1em]{element-mustset.pdf}~outputfileOrbit & \hfuzz=500pt filename & \hfuzz=500pt \\
\hfuzz=500pt\includegraphics[width=1em]{element.pdf}~earthRotation & \hfuzz=500pt \hyperref[earthRotationType]{earthRotation} & \hfuzz=500pt \\
\hfuzz=500pt\includegraphics[width=1em]{element-mustset-unbounded.pdf}~inputfile & \hfuzz=500pt filename & \hfuzz=500pt \\
\hline
\end{tabularx}

\clearpage
%==================================
\subsection{Cosmic2OrbitStar}\label{Cosmic2OrbitStar}
This program reads in cosmic orbit and star camera data given in the CHAMP format.
In case of cosmic orbit and star camera data is stored in one file.
A description of the format can be found under: \url{http://op.gfz-potsdam.de/champ/docs_CHAMP/CH-GFZ-FD-001.pdf}


\keepXColumns
\begin{tabularx}{\textwidth}{N T A}
\hline
Name & Type & Annotation\\
\hline
\hfuzz=500pt\includegraphics[width=1em]{element-mustset.pdf}~outputfileOrbit & \hfuzz=500pt filename & \hfuzz=500pt \\
\hfuzz=500pt\includegraphics[width=1em]{element-mustset.pdf}~outputfileStarCamera & \hfuzz=500pt filename & \hfuzz=500pt \\
\hfuzz=500pt\includegraphics[width=1em]{element-mustset-unbounded.pdf}~inputfile & \hfuzz=500pt filename & \hfuzz=500pt \\
\hline
\end{tabularx}

\clearpage
%==================================
\subsection{Cpf2Orbit}\label{Cpf2Orbit}
Converts \href{https://ilrs.gsfc.nasa.gov/data_and_products/formats/cpf.html}{CPF file}
and writes an \file{instrument file (ORBIT)}{instrument}.

The time format of the CPF file is UTC.
The coordinate system used in the CPF format is usually represented in TRF.
If \configClass{earthRotation}{earthRotationType} is provided the data are transformed
from terrestrial (TRF) to celestial reference frame (CRF).

See also \program{Orbit2Cpf}


\keepXColumns
\begin{tabularx}{\textwidth}{N T A}
\hline
Name & Type & Annotation\\
\hline
\hfuzz=500pt\includegraphics[width=1em]{element-mustset.pdf}~outputfileOrbit & \hfuzz=500pt filename & \hfuzz=500pt \\
\hfuzz=500pt\includegraphics[width=1em]{element-mustset.pdf}~earthRotation & \hfuzz=500pt \hyperref[earthRotationType]{earthRotation} & \hfuzz=500pt \\
\hfuzz=500pt\includegraphics[width=1em]{element-mustset.pdf}~inputfile & \hfuzz=500pt filename & \hfuzz=500pt SLR CPF file\\
\hline
\end{tabularx}

\clearpage
%==================================
\subsection{Crd2NormalPoints}\label{Crd2NormalPoints}
Converts \href{https://ilrs.gsfc.nasa.gov/data_and_products/formats/crd.html}{CRD file}
and writes an \file{instrument file (METEOROLOGICAL)}{instrument} including meteorological data like
temperature, air pressure and humidity as well as an \file{instrument file (SATELLITELASERRANGING)}{instrument}
including normal point data like range, accuracy, redundancy, wavelength and window size.


\keepXColumns
\begin{tabularx}{\textwidth}{N T A}
\hline
Name & Type & Annotation\\
\hline
\hfuzz=500pt\includegraphics[width=1em]{element-mustset.pdf}~outputfileNormalPoints & \hfuzz=500pt filename & \hfuzz=500pt variable \{station\} available\\
\hfuzz=500pt\includegraphics[width=1em]{element-mustset.pdf}~outputfileMeteorological & \hfuzz=500pt filename & \hfuzz=500pt variable \{station\} available\\
\hfuzz=500pt\includegraphics[width=1em]{element-mustset.pdf}~inputfileSlrData & \hfuzz=500pt filename & \hfuzz=500pt SLR CRD files\\
\hline
\end{tabularx}

\clearpage
%==================================
\subsection{Cstg2NormalPoints}\label{Cstg2NormalPoints}
Converts \href{https://ilrs.gsfc.nasa.gov/data_and_products/data/npt/npt_format.html}{CSTG file} provided by the \href{https://ilrs.gsfc.nasa.gov/}{ILRS}
and writes an \file{instrument file (METEOROLOGICAL)}{instrument} including meteorological data like
temperature, air pressure and humidity as well as an \file{instrument file (SATELLITELASERRANGING)}{instrument}
including normal point data like range, accuracy, redundancy, wavelength and window size.


\keepXColumns
\begin{tabularx}{\textwidth}{N T A}
\hline
Name & Type & Annotation\\
\hline
\hfuzz=500pt\includegraphics[width=1em]{element-mustset.pdf}~outputfileNormalPoints & \hfuzz=500pt filename & \hfuzz=500pt variable \{station\} available\\
\hfuzz=500pt\includegraphics[width=1em]{element-mustset.pdf}~outputfileMeteorological & \hfuzz=500pt filename & \hfuzz=500pt variable \{station\} available\\
\hfuzz=500pt\includegraphics[width=1em]{element-mustset.pdf}~inputfileSlrData & \hfuzz=500pt filename & \hfuzz=500pt SLR CSTG file\\
\hline
\end{tabularx}

\clearpage
%==================================
\subsection{DoodsonAdmittance2SupplementaryFiles}\label{DoodsonAdmittance2SupplementaryFiles}
The publication of an ocean tide model includes not only the atlas
in the form of spherical harmonics coefficients,
but also the matrix of Doodson multipliers (\config{outputfileDoodsonMatrix})
and the \config{outputfileAdmittanceMatrix}.

The \config{outputfileMajorTideList} contains the \config{fileNames}
for each contituent.
The required information is taken from the
\configFile{inputfileAdmittance}{admittance}.

See also \program{DoodsonHarmonics2PotentialCoefficients}.


\keepXColumns
\begin{tabularx}{\textwidth}{N T A}
\hline
Name & Type & Annotation\\
\hline
\hfuzz=500pt\includegraphics[width=1em]{element.pdf}~outputfileMajorTideList & \hfuzz=500pt filename & \hfuzz=500pt \\
\hfuzz=500pt\includegraphics[width=1em]{element.pdf}~fileNames & \hfuzz=500pt filename & \hfuzz=500pt template for fileList, variables: doodson, name, cossin\\
\hfuzz=500pt\includegraphics[width=1em]{element.pdf}~outputfileDoodsonMatrix & \hfuzz=500pt filename & \hfuzz=500pt \\
\hfuzz=500pt\includegraphics[width=1em]{element.pdf}~outputfileAdmittanceMatrix & \hfuzz=500pt filename & \hfuzz=500pt \\
\hfuzz=500pt\includegraphics[width=1em]{element-mustset.pdf}~inputfileAdmittance & \hfuzz=500pt filename & \hfuzz=500pt interpolation of minor constituents\\
\hline
\end{tabularx}

\clearpage
%==================================
\subsection{DoodsonHarmonics2IersPotential}\label{DoodsonHarmonics2IersPotential}
Convert doodson harmonics to IERS conventions according to FES2004.
cf. \url{ftp://tai.bipm.org/iers/conv2010/chapter6/tidemodels/fes2004.dat}.


\keepXColumns
\begin{tabularx}{\textwidth}{N T A}
\hline
Name & Type & Annotation\\
\hline
\hfuzz=500pt\includegraphics[width=1em]{element-mustset.pdf}~outputfile & \hfuzz=500pt filename & \hfuzz=500pt according to IERS2010, chapter 6.3.2, footnote 7\\
\hfuzz=500pt\includegraphics[width=1em]{element-mustset.pdf}~inputfileDoodsonHarmoncis & \hfuzz=500pt filename & \hfuzz=500pt \\
\hfuzz=500pt\includegraphics[width=1em]{element-mustset-unbounded.pdf}~header & \hfuzz=500pt string & \hfuzz=500pt info for output header\\
\hfuzz=500pt\includegraphics[width=1em]{element.pdf}~factor & \hfuzz=500pt double & \hfuzz=500pt \\
\hfuzz=500pt\includegraphics[width=1em]{element.pdf}~minDegree & \hfuzz=500pt uint & \hfuzz=500pt \\
\hfuzz=500pt\includegraphics[width=1em]{element.pdf}~maxDegree & \hfuzz=500pt uint & \hfuzz=500pt \\
\hline
\end{tabularx}

\clearpage
%==================================
\subsection{DoodsonHarmonics2IersWaterHeight}\label{DoodsonHarmonics2IersWaterHeight}
Convert doodson harmonics to IERS conventions according to FES2004.
cf. \url{ftp://tai.bipm.org/iers/conv2010/chapter6/tidemodels/fes2004.dat}.


\keepXColumns
\begin{tabularx}{\textwidth}{N T A}
\hline
Name & Type & Annotation\\
\hline
\hfuzz=500pt\includegraphics[width=1em]{element-mustset.pdf}~outputfile & \hfuzz=500pt filename & \hfuzz=500pt according to IERS2010, chapter 6.3.2, footnote 7\\
\hfuzz=500pt\includegraphics[width=1em]{element-mustset.pdf}~inputfileDoodsonHarmoncis & \hfuzz=500pt filename & \hfuzz=500pt \\
\hfuzz=500pt\includegraphics[width=1em]{element-mustset.pdf}~inputfileTideGeneratingPotential & \hfuzz=500pt filename & \hfuzz=500pt to compute Xi phase correction\\
\hfuzz=500pt\includegraphics[width=1em]{element-mustset-unbounded.pdf}~header & \hfuzz=500pt string & \hfuzz=500pt info for output header\\
\hfuzz=500pt\includegraphics[width=1em]{element-mustset.pdf}~kernel & \hfuzz=500pt \hyperref[kernelType]{kernel} & \hfuzz=500pt data type of output values\\
\hfuzz=500pt\includegraphics[width=1em]{element.pdf}~factor & \hfuzz=500pt double & \hfuzz=500pt e.g. from [m] to [cm]\\
\hfuzz=500pt\includegraphics[width=1em]{element.pdf}~minDegree & \hfuzz=500pt uint & \hfuzz=500pt \\
\hfuzz=500pt\includegraphics[width=1em]{element.pdf}~maxDegree & \hfuzz=500pt uint & \hfuzz=500pt \\
\hline
\end{tabularx}

\clearpage
%==================================
\subsection{GnssAntex2AntennaDefinition}\label{GnssAntex2AntennaDefinition}
Converts metadata and antenna definitions from the \href{https://files.igs.org/pub/data/format/antex14.txt}{IGS ANTEX format}.
to \configFile{transmitterInfo}{platform}, \configFile{antennaDefinition}{gnssAntennaDefinition}, \configFile{svnBlockTable}{stringTable},
and \configFile{transmitterList}{stringList} files for the respective GNSS and for the list of ground station antennas.


\keepXColumns
\begin{tabularx}{\textwidth}{N T A}
\hline
Name & Type & Annotation\\
\hline
\hfuzz=500pt\includegraphics[width=1em]{element.pdf}~outputfileAntennaDefinitionStation & \hfuzz=500pt filename & \hfuzz=500pt antenna center variations\\
\hfuzz=500pt\includegraphics[width=1em]{element.pdf}~outputfileAntennaDefinitionTransmitter & \hfuzz=500pt filename & \hfuzz=500pt antenna center variations\\
\hfuzz=500pt\includegraphics[width=1em]{element.pdf}~outputfileTransmitterInfo & \hfuzz=500pt filename & \hfuzz=500pt PRN is appended to file name\\
\hfuzz=500pt\includegraphics[width=1em]{element.pdf}~outputfileSvnBlockTableGps & \hfuzz=500pt filename & \hfuzz=500pt SVN to satellite block mapping\\
\hfuzz=500pt\includegraphics[width=1em]{element.pdf}~outputfileSvnBlockTableGlonass & \hfuzz=500pt filename & \hfuzz=500pt SVN to satellite block mapping\\
\hfuzz=500pt\includegraphics[width=1em]{element.pdf}~outputfileSvnBlockTableGalileo & \hfuzz=500pt filename & \hfuzz=500pt SVN to satellite block mapping\\
\hfuzz=500pt\includegraphics[width=1em]{element.pdf}~outputfileSvnBlockTableBeiDou & \hfuzz=500pt filename & \hfuzz=500pt SVN to satellite block mapping\\
\hfuzz=500pt\includegraphics[width=1em]{element.pdf}~outputfileSvnBlockTableQzss & \hfuzz=500pt filename & \hfuzz=500pt SVN to satellite block mapping\\
\hfuzz=500pt\includegraphics[width=1em]{element.pdf}~outputfileTransmitterListGps & \hfuzz=500pt filename & \hfuzz=500pt list of PRNs\\
\hfuzz=500pt\includegraphics[width=1em]{element.pdf}~outputfileTransmitterListGlonass & \hfuzz=500pt filename & \hfuzz=500pt list of PRNs\\
\hfuzz=500pt\includegraphics[width=1em]{element.pdf}~outputfileTransmitterListGalileo & \hfuzz=500pt filename & \hfuzz=500pt list of PRNs\\
\hfuzz=500pt\includegraphics[width=1em]{element.pdf}~outputfileTransmitterListBeiDou & \hfuzz=500pt filename & \hfuzz=500pt list of PRNs\\
\hfuzz=500pt\includegraphics[width=1em]{element.pdf}~outputfileTransmitterListQzss & \hfuzz=500pt filename & \hfuzz=500pt list of PRNs\\
\hfuzz=500pt\includegraphics[width=1em]{element-mustset.pdf}~inputfileAntex & \hfuzz=500pt filename & \hfuzz=500pt \\
\hfuzz=500pt\includegraphics[width=1em]{element.pdf}~timeStart & \hfuzz=500pt time & \hfuzz=500pt ignore older antenna definitions\\
\hfuzz=500pt\includegraphics[width=1em]{element.pdf}~createZeroModel & \hfuzz=500pt boolean & \hfuzz=500pt create empty antenna patterns\\
\hline
\end{tabularx}

\clearpage
%==================================
\subsection{GnssClock2ClockRinex}\label{GnssClock2ClockRinex}
Converts GNSS clocks from GROOPS format to \href{https://files.igs.org/pub/data/format/rinex_clock304.txt}{IGS clock RINEX format}.
Clocks can be provided via \config{satelliteData} and/or \config{stationData}.
Observed signal types are inferred from \configFile{inputfileSignalBias}{gnssSignalBias}.
Satellites/stations used as clock references can be provided via \config{referenceClock}.

See IGS clock RINEX format description for further details on header information.


\keepXColumns
\begin{tabularx}{\textwidth}{N T A}
\hline
Name & Type & Annotation\\
\hline
\hfuzz=500pt\includegraphics[width=1em]{element-mustset.pdf}~outputfileClockRinex & \hfuzz=500pt filename & \hfuzz=500pt \\
\hfuzz=500pt\includegraphics[width=1em]{element-unbounded.pdf}~satelliteData & \hfuzz=500pt sequence & \hfuzz=500pt one element per satellite\\
\hfuzz=500pt\includegraphics[width=1em]{connector.pdf}\includegraphics[width=1em]{element-mustset.pdf}~inputfileClock & \hfuzz=500pt filename & \hfuzz=500pt clock instrument file\\
\hfuzz=500pt\includegraphics[width=1em]{connector.pdf}\includegraphics[width=1em]{element-mustset.pdf}~inputfileSignalBias & \hfuzz=500pt filename & \hfuzz=500pt signal bias file\\
\hfuzz=500pt\includegraphics[width=1em]{connector.pdf}\includegraphics[width=1em]{element-mustset.pdf}~identifier & \hfuzz=500pt string & \hfuzz=500pt PRN (e.g. G23)\\
\hfuzz=500pt\includegraphics[width=1em]{element-unbounded.pdf}~stationData & \hfuzz=500pt sequence & \hfuzz=500pt one element per station\\
\hfuzz=500pt\includegraphics[width=1em]{connector.pdf}\includegraphics[width=1em]{element-mustset.pdf}~inputfileClock & \hfuzz=500pt filename & \hfuzz=500pt clock instrument file\\
\hfuzz=500pt\includegraphics[width=1em]{connector.pdf}\includegraphics[width=1em]{element-mustset.pdf}~inputfilePosition & \hfuzz=500pt filename & \hfuzz=500pt station position file\\
\hfuzz=500pt\includegraphics[width=1em]{connector.pdf}\includegraphics[width=1em]{element-mustset.pdf}~inputfileStationInfo & \hfuzz=500pt filename & \hfuzz=500pt station info file\\
\hfuzz=500pt\includegraphics[width=1em]{connector.pdf}\includegraphics[width=1em]{element-mustset.pdf}~identifier & \hfuzz=500pt string & \hfuzz=500pt station name (e.g. wtzz)\\
\hfuzz=500pt\includegraphics[width=1em]{element-unbounded.pdf}~comment & \hfuzz=500pt string & \hfuzz=500pt comment in header\\
\hfuzz=500pt\includegraphics[width=1em]{element-mustset.pdf}~program & \hfuzz=500pt string & \hfuzz=500pt name of program (for first line)\\
\hfuzz=500pt\includegraphics[width=1em]{element-mustset.pdf}~institution & \hfuzz=500pt string & \hfuzz=500pt name of agency (for first line)\\
\hfuzz=500pt\includegraphics[width=1em]{element-mustset.pdf}~analysisCenter & \hfuzz=500pt string & \hfuzz=500pt name of analysis center\\
\hfuzz=500pt\includegraphics[width=1em]{element.pdf}~differentialCodeBias & \hfuzz=500pt string & \hfuzz=500pt program and source for applied differential code bias\\
\hfuzz=500pt\includegraphics[width=1em]{element-mustset.pdf}~phaseCenterVariations & \hfuzz=500pt string & \hfuzz=500pt program and source for applied phase center variations\\
\hfuzz=500pt\includegraphics[width=1em]{element-mustset-unbounded.pdf}~referenceClock & \hfuzz=500pt string & \hfuzz=500pt identifier of reference satellite/station\\
\hfuzz=500pt\includegraphics[width=1em]{element-mustset.pdf}~referenceFrame & \hfuzz=500pt string & \hfuzz=500pt terrestrial reference frame for the stations\\
\hline
\end{tabularx}

\clearpage
%==================================
\subsection{GnssClockRinex2InstrumentClock}\label{GnssClockRinex2InstrumentClock}
This program converts clocks from the \href{https://files.igs.org/pub/data/format/rinex_clock304.txt}{IGS clock RINEX format},
which contains the clocks of all satellites and stations in a single file,
into an \file{instrument file (MISCVALUE)}{instrument} for each \config{identifier}
(satellite and/or station).


\keepXColumns
\begin{tabularx}{\textwidth}{N T A}
\hline
Name & Type & Annotation\\
\hline
\hfuzz=500pt\includegraphics[width=1em]{element-mustset.pdf}~outputfileInstrument & \hfuzz=500pt filename & \hfuzz=500pt identifier is appended to each file\\
\hfuzz=500pt\includegraphics[width=1em]{element-mustset-unbounded.pdf}~inputfileClockRinex & \hfuzz=500pt filename & \hfuzz=500pt \\
\hfuzz=500pt\includegraphics[width=1em]{element-mustset-unbounded.pdf}~identifier & \hfuzz=500pt string & \hfuzz=500pt satellite or station identifier, e.g. G23 or alic\\
\hfuzz=500pt\includegraphics[width=1em]{element-mustset-unbounded.pdf}~intervals & \hfuzz=500pt \hyperref[timeSeriesType]{timeSeries} & \hfuzz=500pt \\
\hfuzz=500pt\includegraphics[width=1em]{element.pdf}~minEpochsPerInterval & \hfuzz=500pt uint & \hfuzz=500pt minimum number of epochs in an interval\\
\hline
\end{tabularx}

\clearpage
%==================================
\subsection{GnssEop2IgsErp}\label{GnssEop2IgsErp}
Write GNSS Earth orientation parameters to \href{https://files.igs.org/pub/data/format/erp.txt}{IGS ERP file format}.

Requires polar motion, polar motion rate, dUT1 and LOD parameters in the solution
vector \configFile{inputfileSolution}{matrix} and their sigmas in \configFile{inputfileSigmax}{matrix}.
Solution usually comes out of \program{GnssProcessing}.


\keepXColumns
\begin{tabularx}{\textwidth}{N T A}
\hline
Name & Type & Annotation\\
\hline
\hfuzz=500pt\includegraphics[width=1em]{element-mustset.pdf}~outputfileIgsErp & \hfuzz=500pt filename & \hfuzz=500pt IGS ERP file\\
\hfuzz=500pt\includegraphics[width=1em]{element-mustset-unbounded.pdf}~epoch & \hfuzz=500pt sequence & \hfuzz=500pt e.g. daily solution\\
\hfuzz=500pt\includegraphics[width=1em]{connector.pdf}\includegraphics[width=1em]{element-mustset.pdf}~inputfileSolution & \hfuzz=500pt filename & \hfuzz=500pt parameter vector\\
\hfuzz=500pt\includegraphics[width=1em]{connector.pdf}\includegraphics[width=1em]{element-mustset.pdf}~inputfileSigmax & \hfuzz=500pt filename & \hfuzz=500pt standard deviations of the parameters (sqrt of the diagonal of the inverse normal equation)\\
\hfuzz=500pt\includegraphics[width=1em]{connector.pdf}\includegraphics[width=1em]{element-mustset.pdf}~inputfileParameterNames & \hfuzz=500pt filename & \hfuzz=500pt parameter names\\
\hfuzz=500pt\includegraphics[width=1em]{connector.pdf}\includegraphics[width=1em]{element-unbounded.pdf}~inputfileTransmitterList & \hfuzz=500pt filename & \hfuzz=500pt transmitter PRNs used in solution (used for transmitter count)\\
\hfuzz=500pt\includegraphics[width=1em]{connector.pdf}\includegraphics[width=1em]{element.pdf}~inputfileStationList & \hfuzz=500pt filename & \hfuzz=500pt stations used in solution (used for station count)\\
\hfuzz=500pt\includegraphics[width=1em]{connector.pdf}\includegraphics[width=1em]{element-mustset.pdf}~time & \hfuzz=500pt time & \hfuzz=500pt reference time for epoch\\
\hfuzz=500pt\includegraphics[width=1em]{element-unbounded.pdf}~comment & \hfuzz=500pt string & \hfuzz=500pt \\
\hline
\end{tabularx}

\clearpage
%==================================
\subsection{GnssGriddedDataTimeSeries2Ionex}\label{GnssGriddedDataTimeSeries2Ionex}
Converts TEC maps from GROOPS \file{gridded data time series}{griddedDataTimeSeries} format
to IGS \href{https://files.igs.org/pub/data/format/ionex1.pdf}{IONEX file format}.

Currently only supports 2D TEC maps.

See also \program{GnssIonex2GriddedDataTimeSeries}, \configClass{IonosphereMap}{gnssParametrizationType:ionosphereMap}.


\keepXColumns
\begin{tabularx}{\textwidth}{N T A}
\hline
Name & Type & Annotation\\
\hline
\hfuzz=500pt\includegraphics[width=1em]{element-mustset.pdf}~outputfileIonex & \hfuzz=500pt filename & \hfuzz=500pt \\
\hfuzz=500pt\includegraphics[width=1em]{element-mustset.pdf}~inputfileGriddedDataTimeSeries & \hfuzz=500pt filename & \hfuzz=500pt must contain regular grid\\
\hfuzz=500pt\includegraphics[width=1em]{element-mustset.pdf}~value & \hfuzz=500pt expression & \hfuzz=500pt expression (e.g. data column)\\
\hfuzz=500pt\includegraphics[width=1em]{element-unbounded.pdf}~timeSeries & \hfuzz=500pt \hyperref[timeSeriesType]{timeSeries} & \hfuzz=500pt (empty = use input file time series)\\
\hfuzz=500pt\includegraphics[width=1em]{element-mustset.pdf}~program & \hfuzz=500pt string & \hfuzz=500pt name of program (for first line)\\
\hfuzz=500pt\includegraphics[width=1em]{element-mustset.pdf}~institution & \hfuzz=500pt string & \hfuzz=500pt name of agency (for first line)\\
\hfuzz=500pt\includegraphics[width=1em]{element-unbounded.pdf}~description & \hfuzz=500pt string & \hfuzz=500pt description in header\\
\hfuzz=500pt\includegraphics[width=1em]{element-unbounded.pdf}~comment & \hfuzz=500pt string & \hfuzz=500pt comment in header\\
\hfuzz=500pt\includegraphics[width=1em]{element.pdf}~mappingFunction & \hfuzz=500pt string & \hfuzz=500pt see IONEX documentation\\
\hfuzz=500pt\includegraphics[width=1em]{element.pdf}~elevationCutoff & \hfuzz=500pt double & \hfuzz=500pt see IONEX documentation (0 if unknown)\\
\hfuzz=500pt\includegraphics[width=1em]{element.pdf}~observablesUsed & \hfuzz=500pt string & \hfuzz=500pt see IONEX documentation\\
\hfuzz=500pt\includegraphics[width=1em]{element.pdf}~exponent & \hfuzz=500pt int & \hfuzz=500pt factor 10\textasciicircum{}exponent is applied to all values\\
\hline
\end{tabularx}

\clearpage
%==================================
\subsection{GnssIonex2GriddedDataTimeSeries}\label{GnssIonex2GriddedDataTimeSeries}
Converts TEC maps from IGS \href{https://files.igs.org/pub/data/format/ionex1.pdf}{IONEX file format}
to GROOPS \file{gridded data time series}{griddedDataTimeSeries} format.

Currently only supports 2D TEC maps.

See also \program{GnssGriddedDataTimeSeries2Ionex}, \configClass{IonosphereMap}{gnssParametrizationType:ionosphereMap}.


\keepXColumns
\begin{tabularx}{\textwidth}{N T A}
\hline
Name & Type & Annotation\\
\hline
\hfuzz=500pt\includegraphics[width=1em]{element-mustset.pdf}~outputfileGriddedDataTimeSeries & \hfuzz=500pt filename & \hfuzz=500pt \\
\hfuzz=500pt\includegraphics[width=1em]{element-mustset-unbounded.pdf}~inputfileIonex & \hfuzz=500pt filename & \hfuzz=500pt \\
\hline
\end{tabularx}

\clearpage
%==================================
\subsection{GnssNormals2Sinex}\label{GnssNormals2Sinex}
Write GNSS data/metadata and \file{normal equations}{normalEquation} to
\href{http://www.iers.org/IERS/EN/Organization/AnalysisCoordinator/SinexFormat/sinex.html}{SINEX format}.

Normal equations usually come from \program{GnssProcessing}
(e.g. from \reference{GNSS satellite orbit determination and station network analysis}{cookbook.gnssNetwork}).
Metadata input files include \configFile{stationInfo/transmitterInfo}{platform}, \configFile{antennaDefinition}{gnssAntennaDefinition},
and \configFile{stationList/transmitterList}{stringList}, see \program{GnssAntex2AntennaDefinition}.

See also \program{Sinex2Normals} and \program{NormalsSphericalHarmonics2Sinex}.


\keepXColumns
\begin{tabularx}{\textwidth}{N T A}
\hline
Name & Type & Annotation\\
\hline
\hfuzz=500pt\includegraphics[width=1em]{element.pdf}~outputfileSinexNormals & \hfuzz=500pt filename & \hfuzz=500pt full SINEX file including normal equations\\
\hfuzz=500pt\includegraphics[width=1em]{element.pdf}~outputfileSinexCoordinates & \hfuzz=500pt filename & \hfuzz=500pt SINEX file without normal equations (station coordinates file)\\
\hfuzz=500pt\includegraphics[width=1em]{element-mustset.pdf}~inputfileNormals & \hfuzz=500pt filename & \hfuzz=500pt normal equation matrix\\
\hfuzz=500pt\includegraphics[width=1em]{element.pdf}~inputfileSolution & \hfuzz=500pt filename & \hfuzz=500pt parameter vector\\
\hfuzz=500pt\includegraphics[width=1em]{element.pdf}~inputfileSigmax & \hfuzz=500pt filename & \hfuzz=500pt standard deviations of the parameters (sqrt of the diagonal of the inverse normal equation)\\
\hfuzz=500pt\includegraphics[width=1em]{element-mustset.pdf}~inputfileApriori & \hfuzz=500pt filename & \hfuzz=500pt apriori parameter vector\\
\hfuzz=500pt\includegraphics[width=1em]{element.pdf}~inputfileAprioriSigma & \hfuzz=500pt filename & \hfuzz=500pt constraint sigmas for apriori parameter vector\\
\hfuzz=500pt\includegraphics[width=1em]{element.pdf}~inputfileAprioriMatrix & \hfuzz=500pt filename & \hfuzz=500pt normal equation matrix of applied constraints\\
\hfuzz=500pt\includegraphics[width=1em]{element-mustset-unbounded.pdf}~transmitterConstellation & \hfuzz=500pt sequence & \hfuzz=500pt transmitter constellation metadata\\
\hfuzz=500pt\includegraphics[width=1em]{connector.pdf}\includegraphics[width=1em]{element-mustset.pdf}~inputfileTransmitterList & \hfuzz=500pt filename & \hfuzz=500pt transmitter PRNs used in solution\\
\hfuzz=500pt\includegraphics[width=1em]{connector.pdf}\includegraphics[width=1em]{element-mustset.pdf}~inputfileTransmitterInfo & \hfuzz=500pt filename & \hfuzz=500pt transmitter info file template\\
\hfuzz=500pt\includegraphics[width=1em]{connector.pdf}\includegraphics[width=1em]{element-mustset.pdf}~inputfileAntennaDefinition & \hfuzz=500pt filename & \hfuzz=500pt transmitter phase centers and variations (ANTEX)\\
\hfuzz=500pt\includegraphics[width=1em]{connector.pdf}\includegraphics[width=1em]{element.pdf}~variablePrn & \hfuzz=500pt string & \hfuzz=500pt loop variable for PRNs from transmitter list\\
\hfuzz=500pt\includegraphics[width=1em]{element-mustset.pdf}~stations & \hfuzz=500pt sequence & \hfuzz=500pt \\
\hfuzz=500pt\includegraphics[width=1em]{connector.pdf}\includegraphics[width=1em]{element-mustset.pdf}~inputfileStationList & \hfuzz=500pt filename & \hfuzz=500pt stations contained in normal equations\\
\hfuzz=500pt\includegraphics[width=1em]{connector.pdf}\includegraphics[width=1em]{element-mustset.pdf}~inputfileStationInfo & \hfuzz=500pt filename & \hfuzz=500pt station info file template\\
\hfuzz=500pt\includegraphics[width=1em]{connector.pdf}\includegraphics[width=1em]{element-mustset.pdf}~inputfileAntennaDefinition & \hfuzz=500pt filename & \hfuzz=500pt station phase centers and variations (ANTEX)\\
\hfuzz=500pt\includegraphics[width=1em]{connector.pdf}\includegraphics[width=1em]{element.pdf}~variableStationName & \hfuzz=500pt string & \hfuzz=500pt loop variable for station names from station list\\
\hfuzz=500pt\includegraphics[width=1em]{connector.pdf}\includegraphics[width=1em]{element-mustset.pdf}~observationTimeStart & \hfuzz=500pt time & \hfuzz=500pt start time for which solution has observations\\
\hfuzz=500pt\includegraphics[width=1em]{connector.pdf}\includegraphics[width=1em]{element-mustset.pdf}~observationTimeEnd & \hfuzz=500pt time & \hfuzz=500pt end time for which solution has observations\\
\hfuzz=500pt\includegraphics[width=1em]{element-mustset.pdf}~time & \hfuzz=500pt time & \hfuzz=500pt reference time for parameters\\
\hfuzz=500pt\includegraphics[width=1em]{element.pdf}~sampling & \hfuzz=500pt double & \hfuzz=500pt [seconds] observation sampling\\
\hfuzz=500pt\includegraphics[width=1em]{element-mustset.pdf}~antennaCalibrationModel & \hfuzz=500pt string & \hfuzz=500pt e.g. IGS14\_WWWW (WWWW = ANTEX release GPS week)\\
\hfuzz=500pt\includegraphics[width=1em]{element-mustset.pdf}~sinexHeader & \hfuzz=500pt sequence & \hfuzz=500pt \\
\hfuzz=500pt\includegraphics[width=1em]{connector.pdf}\includegraphics[width=1em]{element.pdf}~agencyCode & \hfuzz=500pt string & \hfuzz=500pt identify the agency providing the data\\
\hfuzz=500pt\includegraphics[width=1em]{connector.pdf}\includegraphics[width=1em]{element.pdf}~timeStart & \hfuzz=500pt time & \hfuzz=500pt start time of the data\\
\hfuzz=500pt\includegraphics[width=1em]{connector.pdf}\includegraphics[width=1em]{element.pdf}~timeEnd & \hfuzz=500pt time & \hfuzz=500pt end time of the data \\
\hfuzz=500pt\includegraphics[width=1em]{connector.pdf}\includegraphics[width=1em]{element.pdf}~observationCode & \hfuzz=500pt string & \hfuzz=500pt technique used to generate the SINEX solution\\
\hfuzz=500pt\includegraphics[width=1em]{connector.pdf}\includegraphics[width=1em]{element.pdf}~constraintCode & \hfuzz=500pt string & \hfuzz=500pt 0: tight constraint, 1: siginficant constraint, 2: unconstrained\\
\hfuzz=500pt\includegraphics[width=1em]{connector.pdf}\includegraphics[width=1em]{element.pdf}~solutionContent & \hfuzz=500pt string & \hfuzz=500pt solution types contained in the SINEX solution (S O E T C A)\\
\hfuzz=500pt\includegraphics[width=1em]{connector.pdf}\includegraphics[width=1em]{element.pdf}~description & \hfuzz=500pt string & \hfuzz=500pt organizitions gathering/alerting the file contents\\
\hfuzz=500pt\includegraphics[width=1em]{connector.pdf}\includegraphics[width=1em]{element.pdf}~contact & \hfuzz=500pt string & \hfuzz=500pt Address of the relevant contact. e-mail\\
\hfuzz=500pt\includegraphics[width=1em]{connector.pdf}\includegraphics[width=1em]{element.pdf}~output & \hfuzz=500pt string & \hfuzz=500pt Description of the file contents\\
\hfuzz=500pt\includegraphics[width=1em]{connector.pdf}\includegraphics[width=1em]{element.pdf}~input & \hfuzz=500pt string & \hfuzz=500pt Brief description of the input used to generate this solution\\
\hfuzz=500pt\includegraphics[width=1em]{connector.pdf}\includegraphics[width=1em]{element.pdf}~software & \hfuzz=500pt string & \hfuzz=500pt Software used to generate the file\\
\hfuzz=500pt\includegraphics[width=1em]{connector.pdf}\includegraphics[width=1em]{element.pdf}~hardware & \hfuzz=500pt string & \hfuzz=500pt Computer hardware on which above software was run\\
\hfuzz=500pt\includegraphics[width=1em]{connector.pdf}\includegraphics[width=1em]{element.pdf}~inputfileComment & \hfuzz=500pt filename & \hfuzz=500pt comments in the comment block from a file (truncated at 80 characters)\\
\hfuzz=500pt\includegraphics[width=1em]{connector.pdf}\includegraphics[width=1em]{element-unbounded.pdf}~comment & \hfuzz=500pt string & \hfuzz=500pt comments in the comment block\\
\hline
\end{tabularx}

\clearpage
%==================================
\subsection{GnssOrbex2StarCamera}\label{GnssOrbex2StarCamera}
Converts GNSS satellite attitude from \href{http://acc.igs.org/misc/proposal_orbex_april2019.pdf}{ORBEX file format}
(quaternions) to \file{instrument file (STARCAMERA)}{instrument}.
The resulting star camera files contain the rotation from satellite body frame to TRF, or to CRF in case
\configClass{earthRotation}{earthRotationType} is provided.

See also \program{GnssAttitude2Orbex}.


\keepXColumns
\begin{tabularx}{\textwidth}{N T A}
\hline
Name & Type & Annotation\\
\hline
\hfuzz=500pt\includegraphics[width=1em]{element-mustset.pdf}~outputfileStarCamera & \hfuzz=500pt filename & \hfuzz=500pt rotation from body frame to TRF/CRF, identifier is appended to each file\\
\hfuzz=500pt\includegraphics[width=1em]{element-mustset.pdf}~inputfileOrbex & \hfuzz=500pt filename & \hfuzz=500pt \\
\hfuzz=500pt\includegraphics[width=1em]{element-unbounded.pdf}~identifier & \hfuzz=500pt string & \hfuzz=500pt (empty = all) satellite identifier, e.g. G23 or E05\\
\hfuzz=500pt\includegraphics[width=1em]{element.pdf}~earthRotation & \hfuzz=500pt \hyperref[earthRotationType]{earthRotation} & \hfuzz=500pt rotation from TRF to CRF\\
\hline
\end{tabularx}

\clearpage
%==================================
\subsection{GnssReceiver2RinexObservation}\label{GnssReceiver2RinexObservation}
Converts a \configFile{inputfileGnssReceiver}{instrument} into a
\href{https://files.igs.org/pub/data/format/rinex_4.00.pdf}{RINEX} observation file.
The \configFile{inputfileStationInfo}{platform} contains the antenna
and receiver information for the RINEX header.
The \configClass{useType}{gnssType} and \configClass{ignoreType}{gnssType} can be used to filter
the observation types that will be exported.


\keepXColumns
\begin{tabularx}{\textwidth}{N T A}
\hline
Name & Type & Annotation\\
\hline
\hfuzz=500pt\includegraphics[width=1em]{element-mustset.pdf}~outputfileRinexObservation & \hfuzz=500pt filename & \hfuzz=500pt RINEX observation file\\
\hfuzz=500pt\includegraphics[width=1em]{element-mustset.pdf}~inputfileGnssReceiver & \hfuzz=500pt filename & \hfuzz=500pt GNSS instrument file\\
\hfuzz=500pt\includegraphics[width=1em]{element-mustset.pdf}~inputfileStationInfo & \hfuzz=500pt filename & \hfuzz=500pt antenna and receiver info\\
\hfuzz=500pt\includegraphics[width=1em]{element-unbounded.pdf}~comment & \hfuzz=500pt string & \hfuzz=500pt write comments at begin of header\\
\hfuzz=500pt\includegraphics[width=1em]{element.pdf}~observer & \hfuzz=500pt string & \hfuzz=500pt header information\\
\hfuzz=500pt\includegraphics[width=1em]{element.pdf}~angency & \hfuzz=500pt string & \hfuzz=500pt header information\\
\hfuzz=500pt\includegraphics[width=1em]{element-unbounded.pdf}~useType & \hfuzz=500pt \hyperref[gnssType]{gnssType} & \hfuzz=500pt only use observations that match any of these patterns\\
\hfuzz=500pt\includegraphics[width=1em]{element-unbounded.pdf}~ignoreType & \hfuzz=500pt \hyperref[gnssType]{gnssType} & \hfuzz=500pt ignore observations that match any of these patterns\\
\hline
\end{tabularx}

\clearpage
%==================================
\subsection{GnssRinexNavigation2OrbitClock}\label{GnssRinexNavigation2OrbitClock}
Evaluates orbit and clock parameters from \href{https://files.igs.org/pub/data/format/rinex_4.00.pdf}{RINEX} (version 2, 3, and 4)
navigation file \config{inputfileRinex} at epochs given by \configClass{timeSeries}{timeSeriesType} and writes them to
\configFile{outputfileOrbit}{instrument} and \configFile{outputfileClock}{instrument}, respectively.

Orbits are rotated from TRF (as broadcasted) to CRF via \configClass{earthRotation}{earthRotationType},
but system-specific TRFs (WGS84, PZ-90, etc.) are not aligned to a common TRF.

Furthermore, option is available to remove any satellite ephemeris data that has their satellite flag set to unhealthy.

See also \program{OrbitAddVelocityAndAcceleration}.


\keepXColumns
\begin{tabularx}{\textwidth}{N T A}
\hline
Name & Type & Annotation\\
\hline
\hfuzz=500pt\includegraphics[width=1em]{element.pdf}~outputfileOrbit & \hfuzz=500pt filename & \hfuzz=500pt PRN is appended to file name\\
\hfuzz=500pt\includegraphics[width=1em]{element.pdf}~outputfileClock & \hfuzz=500pt filename & \hfuzz=500pt PRN is appended to file name\\
\hfuzz=500pt\includegraphics[width=1em]{element-mustset.pdf}~inputfileRinex & \hfuzz=500pt filename & \hfuzz=500pt RINEX navigation file\\
\hfuzz=500pt\includegraphics[width=1em]{element-mustset-unbounded.pdf}~timeSeries & \hfuzz=500pt \hyperref[timeSeriesType]{timeSeries} & \hfuzz=500pt orbit and clock evaluation epochs\\
\hfuzz=500pt\includegraphics[width=1em]{element-mustset.pdf}~earthRotation & \hfuzz=500pt \hyperref[earthRotationType]{earthRotation} & \hfuzz=500pt for rotation from TRF to CRF\\
\hfuzz=500pt\includegraphics[width=1em]{element-unbounded.pdf}~useType & \hfuzz=500pt \hyperref[gnssType]{gnssType} & \hfuzz=500pt (e.g. ***G12) only use satellites with PRN that match any of these patterns\\
\hfuzz=500pt\includegraphics[width=1em]{element-unbounded.pdf}~ignoreType & \hfuzz=500pt \hyperref[gnssType]{gnssType} & \hfuzz=500pt (e.g. ***R**) ignore satellites PRN that match any of these patterns\\
\hfuzz=500pt\includegraphics[width=1em]{element.pdf}~removeUnhealthySatellites & \hfuzz=500pt boolean & \hfuzz=500pt Remove satellite ephemeris that have their sat flags set to unhealthy\\
\hline
\end{tabularx}

\clearpage
%==================================
\subsection{GnssSignalBias2SinexBias}\label{GnssSignalBias2SinexBias}
Convert \file{GNSS signal biases}{gnssSignalBias} from GROOPS format to \href{https://files.igs.org/pub/data/format/sinex_bias_100.pdf}{IGS SINEX Bias format}.
Biases can be provided via \config{transmitterBiases} and/or \config{receiverBiases}.
Phase biases without attribute (e.g. \verb|L1*|) are automatically expanded so each code
bias has a corresponding phase bias
(Example: \verb|C1C|, \verb|C1W|, \verb|L1*| are converted to \verb|C1C|, \verb|C1W|, \verb|L1C|, \verb|L1W|).

Time-variable biases (e.g. GPS L5 satellite phase bias) can be provided via \config{timeVariableBias}.
Their time span will be based on the provided epochs ($t \pm \Delta t / 2$).
The slope of the bias can be optionally provided in the second data column.

If GLONASS receiver biases depend on frequency number, those must be defined in \configFile{inputfileTransmitterInfo}{platform}
to get the correct PRN/SVN assignment to the biases.

See IGS SINEX Bias format description for further details on header information.

See also \program{GnssSinexBias2SignalBias} and \program{GnssBiasClockAlignment}.


\keepXColumns
\begin{tabularx}{\textwidth}{N T A}
\hline
Name & Type & Annotation\\
\hline
\hfuzz=500pt\includegraphics[width=1em]{element-mustset.pdf}~outputfileSinexBias & \hfuzz=500pt filename & \hfuzz=500pt \\
\hfuzz=500pt\includegraphics[width=1em]{element-mustset-unbounded.pdf}~inputfileTransmitterInfo & \hfuzz=500pt filename & \hfuzz=500pt one file per satellite\\
\hfuzz=500pt\includegraphics[width=1em]{element-unbounded.pdf}~transmitterBiases & \hfuzz=500pt sequence & \hfuzz=500pt one element per satellite\\
\hfuzz=500pt\includegraphics[width=1em]{connector.pdf}\includegraphics[width=1em]{element-mustset.pdf}~inputfileSignalBias & \hfuzz=500pt filename & \hfuzz=500pt signal bias file\\
\hfuzz=500pt\includegraphics[width=1em]{connector.pdf}\includegraphics[width=1em]{element-unbounded.pdf}~timeVariableBias & \hfuzz=500pt sequence & \hfuzz=500pt one entry per time variable bias type\\
\hfuzz=500pt\quad\includegraphics[width=1em]{connector.pdf}\includegraphics[width=1em]{element-mustset.pdf}~inputfileSignalBias & \hfuzz=500pt filename & \hfuzz=500pt columns: mjd, bias [m], (biasSlope [m/s])\\
\hfuzz=500pt\quad\includegraphics[width=1em]{connector.pdf}\includegraphics[width=1em]{element-mustset.pdf}~type & \hfuzz=500pt \hyperref[gnssType]{gnssType} & \hfuzz=500pt bias type\\
\hfuzz=500pt\includegraphics[width=1em]{connector.pdf}\includegraphics[width=1em]{element-mustset.pdf}~identifier & \hfuzz=500pt string & \hfuzz=500pt PRN or station name (e.g. G23 or wtzz)\\
\hfuzz=500pt\includegraphics[width=1em]{element-unbounded.pdf}~receiverBiases & \hfuzz=500pt sequence & \hfuzz=500pt one element per station\\
\hfuzz=500pt\includegraphics[width=1em]{connector.pdf}\includegraphics[width=1em]{element-mustset.pdf}~inputfileSignalBias & \hfuzz=500pt filename & \hfuzz=500pt signal bias file\\
\hfuzz=500pt\includegraphics[width=1em]{connector.pdf}\includegraphics[width=1em]{element-unbounded.pdf}~timeVariableBias & \hfuzz=500pt sequence & \hfuzz=500pt one entry per time variable bias type\\
\hfuzz=500pt\quad\includegraphics[width=1em]{connector.pdf}\includegraphics[width=1em]{element-mustset.pdf}~inputfileSignalBias & \hfuzz=500pt filename & \hfuzz=500pt columns: mjd, bias [m], (biasSlope [m/s])\\
\hfuzz=500pt\quad\includegraphics[width=1em]{connector.pdf}\includegraphics[width=1em]{element-mustset.pdf}~type & \hfuzz=500pt \hyperref[gnssType]{gnssType} & \hfuzz=500pt bias type\\
\hfuzz=500pt\includegraphics[width=1em]{connector.pdf}\includegraphics[width=1em]{element-mustset.pdf}~identifier & \hfuzz=500pt string & \hfuzz=500pt PRN or station name (e.g. G23 or wtzz)\\
\hfuzz=500pt\includegraphics[width=1em]{element-mustset.pdf}~agencyCode & \hfuzz=500pt string & \hfuzz=500pt identify the agency providing the data\\
\hfuzz=500pt\includegraphics[width=1em]{element-mustset.pdf}~fileAgencyCode & \hfuzz=500pt string & \hfuzz=500pt identify the agency creating the file\\
\hfuzz=500pt\includegraphics[width=1em]{element-mustset.pdf}~timeStart & \hfuzz=500pt time & \hfuzz=500pt start time of the data\\
\hfuzz=500pt\includegraphics[width=1em]{element-mustset.pdf}~timeEnd & \hfuzz=500pt time & \hfuzz=500pt end time of the data \\
\hfuzz=500pt\includegraphics[width=1em]{element-mustset.pdf}~biasMode & \hfuzz=500pt choice & \hfuzz=500pt absolute or relative bias estimates\\
\hfuzz=500pt\includegraphics[width=1em]{connector.pdf}\includegraphics[width=1em]{element-mustset.pdf}~absolute & \hfuzz=500pt  & \hfuzz=500pt \\
\hfuzz=500pt\includegraphics[width=1em]{connector.pdf}\includegraphics[width=1em]{element-mustset.pdf}~relative & \hfuzz=500pt  & \hfuzz=500pt \\
\hfuzz=500pt\includegraphics[width=1em]{element.pdf}~observationSampling & \hfuzz=500pt uint & \hfuzz=500pt [seconds]\\
\hfuzz=500pt\includegraphics[width=1em]{element.pdf}~intervalLength & \hfuzz=500pt uint & \hfuzz=500pt [seconds] interval for bias parameter representation\\
\hfuzz=500pt\includegraphics[width=1em]{element-mustset.pdf}~determinationMethod & \hfuzz=500pt string & \hfuzz=500pt determination method used to generate the bias results (see SINEX Bias format description)\\
\hfuzz=500pt\includegraphics[width=1em]{element.pdf}~receiverClockReferenceGnss & \hfuzz=500pt string & \hfuzz=500pt (G, R, E, C) reference GNSS used for receiver clock estimation\\
\hfuzz=500pt\includegraphics[width=1em]{element-unbounded.pdf}~satelliteClockReferenceObservables & \hfuzz=500pt string & \hfuzz=500pt one per system, reference code observable on first and second frequency (RINEX3 format)\\
\hfuzz=500pt\includegraphics[width=1em]{element-mustset.pdf}~description & \hfuzz=500pt string & \hfuzz=500pt organizition gathering/altering the file contents\\
\hfuzz=500pt\includegraphics[width=1em]{element-mustset.pdf}~contact & \hfuzz=500pt string & \hfuzz=500pt contact name and/or email address\\
\hfuzz=500pt\includegraphics[width=1em]{element-mustset.pdf}~input & \hfuzz=500pt string & \hfuzz=500pt brief description of the input used to generate this solution\\
\hfuzz=500pt\includegraphics[width=1em]{element-mustset.pdf}~output & \hfuzz=500pt string & \hfuzz=500pt description of the file contents\\
\hfuzz=500pt\includegraphics[width=1em]{element-mustset.pdf}~software & \hfuzz=500pt string & \hfuzz=500pt software used to generate the file\\
\hfuzz=500pt\includegraphics[width=1em]{element-mustset.pdf}~hardware & \hfuzz=500pt string & \hfuzz=500pt computer hardware on which above software was run\\
\hfuzz=500pt\includegraphics[width=1em]{element-unbounded.pdf}~comment & \hfuzz=500pt string & \hfuzz=500pt comments in the comment block\\
\hline
\end{tabularx}

\clearpage
%==================================
\subsection{GnssSinexBias2SignalBias}\label{GnssSinexBias2SignalBias}
Converts GNSS signal biases from \href{https://files.igs.org/pub/data/format/sinex_bias_100.pdf}{IGS SINEX Bias format}
to \file{GnssSignalBias format}{gnssSignalBias}.

Only satellite observable-specific signal biases (OSB) are supported at the moment.
If multiple entries exist for the same bias, the weighted average (based on time span) of all entries is used.
Time-variable biases are not supported at the moment.

See also \program{GnssSignalBias2SinexBias}.


\keepXColumns
\begin{tabularx}{\textwidth}{N T A}
\hline
Name & Type & Annotation\\
\hline
\hfuzz=500pt\includegraphics[width=1em]{element-mustset.pdf}~outputfileSignalBias & \hfuzz=500pt filename & \hfuzz=500pt identifier is appended to file name\\
\hfuzz=500pt\includegraphics[width=1em]{element-mustset.pdf}~inputfileSinexBias & \hfuzz=500pt filename & \hfuzz=500pt \\
\hfuzz=500pt\includegraphics[width=1em]{element.pdf}~inputfileGlonassSignalDefinition & \hfuzz=500pt filename & \hfuzz=500pt GLONASS frequency number mapping\\
\hfuzz=500pt\includegraphics[width=1em]{element-unbounded.pdf}~identifier & \hfuzz=500pt string & \hfuzz=500pt (empty = all) satellite PRN, e.g. G23 or E05\\
\hline
\end{tabularx}

\clearpage
%==================================
\subsection{GnssStationLog2Platform}\label{GnssStationLog2Platform}
Converts \href{https://files.igs.org/pub/station/general/blank.log}{IGS station log format} to \configFile{outputfileStationPlatform}{platform}.

If \configFile{inputfileAntennaDefinition}{gnssAntennaDefinition} is provided, station log data is cross-checked with the given antenna definitions.
Cross-checking station log data with a \href{https://www.iers.org/IERS/EN/Organization/AnalysisCoordinator/SinexFormat/sinex.html}{SINEX file} is
possible with \program{CheckStationsPlatformsWithSinex}.


\keepXColumns
\begin{tabularx}{\textwidth}{N T A}
\hline
Name & Type & Annotation\\
\hline
\hfuzz=500pt\includegraphics[width=1em]{element-mustset.pdf}~outputfileStationPlatform & \hfuzz=500pt filename & \hfuzz=500pt \\
\hfuzz=500pt\includegraphics[width=1em]{element-mustset.pdf}~inputfileStationLog & \hfuzz=500pt filename & \hfuzz=500pt \\
\hfuzz=500pt\includegraphics[width=1em]{element.pdf}~inputfileAntennaDefinition & \hfuzz=500pt filename & \hfuzz=500pt used to check antennas\\
\hline
\end{tabularx}

\clearpage
%==================================
\subsection{GnssStationLog2StationInfo}\label{GnssStationLog2StationInfo}
Converts \href{https://files.igs.org/pub/station/general/blank.log}{IGS station log format} to \configFile{outputfileStationInfo}{gnssStationInfo}.

If \configFile{inputfileAntennaDefinition}{gnssAntennaDefinition} is provided, station log data is cross-checked with the given antenna definitions.
Cross-checking station log data with a \href{https://www.iers.org/IERS/EN/Organization/AnalysisCoordinator/SinexFormat/sinex.html}{SINEX file} is also
possible by providing \config{inputfileSinex}. Any failed checks result in warnings in the output log.


\keepXColumns
\begin{tabularx}{\textwidth}{N T A}
\hline
Name & Type & Annotation\\
\hline
\hfuzz=500pt\includegraphics[width=1em]{element-mustset.pdf}~outputfileStationInfo & \hfuzz=500pt filename & \hfuzz=500pt \\
\hfuzz=500pt\includegraphics[width=1em]{element-mustset.pdf}~inputfileStationLog & \hfuzz=500pt filename & \hfuzz=500pt \\
\hfuzz=500pt\includegraphics[width=1em]{element.pdf}~inputfileAntennaDefinition & \hfuzz=500pt filename & \hfuzz=500pt used to check antennas\\
\hfuzz=500pt\includegraphics[width=1em]{element.pdf}~inputfileSinex & \hfuzz=500pt filename & \hfuzz=500pt used to cross-check station log with SINEX file\\
\hline
\end{tabularx}

\clearpage
%==================================
\subsection{GnssTroposphere2TropoSinex}\label{GnssTroposphere2TropoSinex}
Convert GNSS troposphere data from GROOPS format to \href{https://files.igs.org/pub/data/format/sinex_tro_v2.00.pdf}{IGS SINEX TRO} format.

Specification of the station list is done via \config{inputfileStationList}.
\config{inputfileTroposphereData} needs the troposphere data provided from \program{GnssProcessing}.
Additional following station metadata are required: \config{inputfileStationInfo} using file type \file{Station info}{platform},
\config{inputfileAntennaDefinition} using file type \file{Antenna definition}{gnssAntennaDefinition} and \config{inputfileGridPos} which
uses the stations positions provided from \program{GnssProcessing}.
For considering the geoid height use \config{inputfileGeoidHeight}. The geoid height is provided by \program{Gravityfield2GriddedData}.


\keepXColumns
\begin{tabularx}{\textwidth}{N T A}
\hline
Name & Type & Annotation\\
\hline
\hfuzz=500pt\includegraphics[width=1em]{element-mustset.pdf}~outputfileTropoSinex & \hfuzz=500pt filename & \hfuzz=500pt \\
\hfuzz=500pt\includegraphics[width=1em]{element-mustset.pdf}~stations & \hfuzz=500pt sequence & \hfuzz=500pt \\
\hfuzz=500pt\includegraphics[width=1em]{connector.pdf}\includegraphics[width=1em]{element-mustset.pdf}~inputfileStationList & \hfuzz=500pt filename & \hfuzz=500pt ASCII file with station names\\
\hfuzz=500pt\includegraphics[width=1em]{connector.pdf}\includegraphics[width=1em]{element-mustset.pdf}~inputfileTroposphereData & \hfuzz=500pt filename & \hfuzz=500pt Troposphere data estimates template (columns: mjd, trodry, trowet, tgndry, tgnwet, tgedry, tgewet)\\
\hfuzz=500pt\includegraphics[width=1em]{connector.pdf}\includegraphics[width=1em]{element.pdf}~inputfileTroposphereSigmas & \hfuzz=500pt filename & \hfuzz=500pt Troposphere data sigmas template (columns: mjd, sigma\_trowet, sigma\_tgnwet, sigma\_tgewet)\\
\hfuzz=500pt\includegraphics[width=1em]{connector.pdf}\includegraphics[width=1em]{element-mustset.pdf}~inputfileStationInfo & \hfuzz=500pt filename & \hfuzz=500pt station info file template\\
\hfuzz=500pt\includegraphics[width=1em]{connector.pdf}\includegraphics[width=1em]{element.pdf}~inputfileGeoidHeight & \hfuzz=500pt filename & \hfuzz=500pt File including geoid height\\
\hfuzz=500pt\includegraphics[width=1em]{connector.pdf}\includegraphics[width=1em]{element-mustset.pdf}~inputfileGridPos & \hfuzz=500pt filename & \hfuzz=500pt File including stations positions\\
\hfuzz=500pt\includegraphics[width=1em]{connector.pdf}\includegraphics[width=1em]{element-mustset.pdf}~inputfileAntennaDefinition & \hfuzz=500pt filename & \hfuzz=500pt station phase centers and variations (ANTEX)\\
\hfuzz=500pt\includegraphics[width=1em]{connector.pdf}\includegraphics[width=1em]{element.pdf}~variableStationName & \hfuzz=500pt filename & \hfuzz=500pt Loop variable for station names from station list\\
\hfuzz=500pt\includegraphics[width=1em]{connector.pdf}\includegraphics[width=1em]{element-mustset.pdf}~observationTimeStart & \hfuzz=500pt time & \hfuzz=500pt Start time for which solution has observations\\
\hfuzz=500pt\includegraphics[width=1em]{connector.pdf}\includegraphics[width=1em]{element-mustset.pdf}~observationTimeEnd & \hfuzz=500pt time & \hfuzz=500pt End time for which solution has observations\\
\hfuzz=500pt\includegraphics[width=1em]{element-mustset.pdf}~dataSamplingInterval & \hfuzz=500pt double & \hfuzz=500pt [sec] GNSS data sampling rate\\
\hfuzz=500pt\includegraphics[width=1em]{element-mustset.pdf}~tropoSamplingInterval & \hfuzz=500pt double & \hfuzz=500pt [sec] Tropospheric parameter sampling interval\\
\hfuzz=500pt\includegraphics[width=1em]{element-mustset.pdf}~tropoModelingMethod & \hfuzz=500pt string & \hfuzz=500pt Tropospheric estimation method: Filter, Smoother, Least Squares, Piece-Wise Linear Interpolation\\
\hfuzz=500pt\includegraphics[width=1em]{element-mustset.pdf}~aPrioriTropoModel & \hfuzz=500pt string & \hfuzz=500pt A priori tropospheric model used\\
\hfuzz=500pt\includegraphics[width=1em]{element-mustset.pdf}~tropoMappingFunction & \hfuzz=500pt string & \hfuzz=500pt Name of mapping function used for hydrostatic and wet delay\\
\hfuzz=500pt\includegraphics[width=1em]{element-mustset.pdf}~gradientMappingFunction & \hfuzz=500pt string & \hfuzz=500pt Name of mapping function used for gradients\\
\hfuzz=500pt\includegraphics[width=1em]{element.pdf}~metDataSource & \hfuzz=500pt string & \hfuzz=500pt source of surface meteorological observations used (see format desc.)\\
\hfuzz=500pt\includegraphics[width=1em]{element-mustset.pdf}~observationWeighting & \hfuzz=500pt string & \hfuzz=500pt observation weighting model applied\\
\hfuzz=500pt\includegraphics[width=1em]{element-mustset.pdf}~elevationCutoff & \hfuzz=500pt double & \hfuzz=500pt [deg]\\
\hfuzz=500pt\includegraphics[width=1em]{element-mustset.pdf}~gnssSystems & \hfuzz=500pt string & \hfuzz=500pt G=GPS, R=GLONASS, E=Galileo, C=BeiDou\\
\hfuzz=500pt\includegraphics[width=1em]{element-mustset.pdf}~timeSystem & \hfuzz=500pt string & \hfuzz=500pt G (GPS) or UTC\\
\hfuzz=500pt\includegraphics[width=1em]{element.pdf}~oceanTideModel & \hfuzz=500pt string & \hfuzz=500pt Ocean tide loading model applied\\
\hfuzz=500pt\includegraphics[width=1em]{element.pdf}~atmosphericTideModel & \hfuzz=500pt string & \hfuzz=500pt Atmospheric tide loading model applied\\
\hfuzz=500pt\includegraphics[width=1em]{element.pdf}~geoidModel & \hfuzz=500pt string & \hfuzz=500pt Geoid model name for undulation values\\
\hfuzz=500pt\includegraphics[width=1em]{element-mustset.pdf}~systemCode & \hfuzz=500pt string & \hfuzz=500pt Terrestrial reference system code\\
\hfuzz=500pt\includegraphics[width=1em]{element-mustset.pdf}~remark & \hfuzz=500pt string & \hfuzz=500pt Remark used to identify the origin of the coordinates (AC acronym)\\
\hfuzz=500pt\includegraphics[width=1em]{element-mustset.pdf}~antennaCalibrationModel & \hfuzz=500pt string & \hfuzz=500pt e.g. IGS14\_WWWW (WWWW = ANTEX release GPS week)\\
\hfuzz=500pt\includegraphics[width=1em]{element-mustset.pdf}~sinexTroHeader & \hfuzz=500pt sequence & \hfuzz=500pt \\
\hfuzz=500pt\includegraphics[width=1em]{connector.pdf}\includegraphics[width=1em]{element-mustset.pdf}~agencyCode & \hfuzz=500pt string & \hfuzz=500pt Identify the agency providing the data\\
\hfuzz=500pt\includegraphics[width=1em]{connector.pdf}\includegraphics[width=1em]{element-mustset.pdf}~timeStart & \hfuzz=500pt time & \hfuzz=500pt Start time of the data\\
\hfuzz=500pt\includegraphics[width=1em]{connector.pdf}\includegraphics[width=1em]{element-mustset.pdf}~timeEnd & \hfuzz=500pt time & \hfuzz=500pt End time of the data \\
\hfuzz=500pt\includegraphics[width=1em]{connector.pdf}\includegraphics[width=1em]{element-mustset.pdf}~observationCode & \hfuzz=500pt string & \hfuzz=500pt Technique used to generate the SINEX solution\\
\hfuzz=500pt\includegraphics[width=1em]{connector.pdf}\includegraphics[width=1em]{element-mustset.pdf}~solutionContents & \hfuzz=500pt string & \hfuzz=500pt Marker name if single station, MIX if multiple stations\\
\hfuzz=500pt\includegraphics[width=1em]{connector.pdf}\includegraphics[width=1em]{element-mustset.pdf}~description & \hfuzz=500pt string & \hfuzz=500pt Organizitions gathering/alerting the file contents\\
\hfuzz=500pt\includegraphics[width=1em]{connector.pdf}\includegraphics[width=1em]{element-mustset.pdf}~output & \hfuzz=500pt string & \hfuzz=500pt Description of the file contents\\
\hfuzz=500pt\includegraphics[width=1em]{connector.pdf}\includegraphics[width=1em]{element-mustset.pdf}~contact & \hfuzz=500pt string & \hfuzz=500pt Address of the relevant contact e-mail\\
\hfuzz=500pt\includegraphics[width=1em]{connector.pdf}\includegraphics[width=1em]{element-mustset.pdf}~software & \hfuzz=500pt string & \hfuzz=500pt Software used to generate the file\\
\hfuzz=500pt\includegraphics[width=1em]{connector.pdf}\includegraphics[width=1em]{element-mustset.pdf}~hardware & \hfuzz=500pt string & \hfuzz=500pt Computer hardware on which above software was run\\
\hfuzz=500pt\includegraphics[width=1em]{connector.pdf}\includegraphics[width=1em]{element-mustset.pdf}~input & \hfuzz=500pt string & \hfuzz=500pt Brief description of the input used to generate this solution\\
\hfuzz=500pt\includegraphics[width=1em]{connector.pdf}\includegraphics[width=1em]{element-mustset.pdf}~versionNumber & \hfuzz=500pt string & \hfuzz=500pt Unique identifier of the product, same as in file name, e.g. 000\\
\hfuzz=500pt\includegraphics[width=1em]{connector.pdf}\includegraphics[width=1em]{element.pdf}~inputfileComment & \hfuzz=500pt filename & \hfuzz=500pt comments in the comment block from a file (truncated at 80 characters per line)\\
\hfuzz=500pt\includegraphics[width=1em]{connector.pdf}\includegraphics[width=1em]{element-unbounded.pdf}~comment & \hfuzz=500pt string & \hfuzz=500pt comments in the comment block\\
\hline
\end{tabularx}

\clearpage
%==================================
\subsection{GoceXml2Gradiometer}\label{GoceXml2Gradiometer}
Read ESA XML GOCE Data.
The \config{outputfileGradiometer} is written as \file{instrument file (GRADIOMETER)}{instrument}.


\keepXColumns
\begin{tabularx}{\textwidth}{N T A}
\hline
Name & Type & Annotation\\
\hline
\hfuzz=500pt\includegraphics[width=1em]{element-mustset.pdf}~outputfileGradiometer & \hfuzz=500pt filename & \hfuzz=500pt \\
\hfuzz=500pt\includegraphics[width=1em]{element-mustset-unbounded.pdf}~inputfile & \hfuzz=500pt filename & \hfuzz=500pt \\
\hline
\end{tabularx}

\clearpage
%==================================
\subsection{GoceXml2Orbit}\label{GoceXml2Orbit}
Read ESA XML GOCE Data.


\keepXColumns
\begin{tabularx}{\textwidth}{N T A}
\hline
Name & Type & Annotation\\
\hline
\hfuzz=500pt\includegraphics[width=1em]{element-mustset.pdf}~outputfileOrbit & \hfuzz=500pt filename & \hfuzz=500pt \\
\hfuzz=500pt\includegraphics[width=1em]{element.pdf}~earthRotation & \hfuzz=500pt \hyperref[earthRotationType]{earthRotation} & \hfuzz=500pt rotation from TRF to CRF\\
\hfuzz=500pt\includegraphics[width=1em]{element-mustset-unbounded.pdf}~inputfile & \hfuzz=500pt filename & \hfuzz=500pt \\
\hline
\end{tabularx}

\clearpage
%==================================
\subsection{GoceXml2StarCamera}\label{GoceXml2StarCamera}
Read ESA XML GOCE Data.


\keepXColumns
\begin{tabularx}{\textwidth}{N T A}
\hline
Name & Type & Annotation\\
\hline
\hfuzz=500pt\includegraphics[width=1em]{element-mustset.pdf}~outputfileStarCamera & \hfuzz=500pt filename & \hfuzz=500pt \\
\hfuzz=500pt\includegraphics[width=1em]{element-mustset-unbounded.pdf}~inputfile & \hfuzz=500pt filename & \hfuzz=500pt \\
\hline
\end{tabularx}

\clearpage
%==================================
\subsection{GoceXmlEggNom1b}\label{GoceXmlEggNom1b}
Read ESA XML GOCE Data.


\keepXColumns
\begin{tabularx}{\textwidth}{N T A}
\hline
Name & Type & Annotation\\
\hline
\hfuzz=500pt\includegraphics[width=1em]{element.pdf}~outputfileGradiometer & \hfuzz=500pt filename & \hfuzz=500pt \\
\hfuzz=500pt\includegraphics[width=1em]{element.pdf}~outputfileAccelerometer & \hfuzz=500pt filename & \hfuzz=500pt \\
\hfuzz=500pt\includegraphics[width=1em]{element.pdf}~outputfileStarCamera & \hfuzz=500pt filename & \hfuzz=500pt \\
\hfuzz=500pt\includegraphics[width=1em]{element.pdf}~outputfileAngularRate & \hfuzz=500pt filename & \hfuzz=500pt \\
\hfuzz=500pt\includegraphics[width=1em]{element.pdf}~outputfileAngularAcc & \hfuzz=500pt filename & \hfuzz=500pt \\
\hfuzz=500pt\includegraphics[width=1em]{element-mustset-unbounded.pdf}~inputfile & \hfuzz=500pt filename & \hfuzz=500pt \\
\hline
\end{tabularx}

\clearpage
%==================================
\subsection{Grace2PotentialCoefficients}\label{Grace2PotentialCoefficients}
This program converts potential coefficients from the GRACE SDS format
into \file{potential coefficients file}{potentialCoefficients}.
The program supports file formats for RL04 to RL06.

Within the program, the variables \verb|epochStart|, \verb|epochEnd| and \verb|epochMid|
are populated with the corresponding time-stamps in the file.
These can be used in to \configFile{outputfilePotentialCoefficients}{potentialCoefficients}
to auto-generate the file name.


\keepXColumns
\begin{tabularx}{\textwidth}{N T A}
\hline
Name & Type & Annotation\\
\hline
\hfuzz=500pt\includegraphics[width=1em]{element-mustset.pdf}~outputfilePotentialCoefficients & \hfuzz=500pt filename & \hfuzz=500pt variables: epochStart, epochEnd, epochMid\\
\hfuzz=500pt\includegraphics[width=1em]{element-mustset.pdf}~inputfile & \hfuzz=500pt filename & \hfuzz=500pt \\
\hline
\end{tabularx}

\clearpage
%==================================
\subsection{GraceAccelerometer2L1bAscii}\label{GraceAccelerometer2L1bAscii}
Convert GROOPS accelerometer files to the GRACE SDS L1B ASCII format.


\keepXColumns
\begin{tabularx}{\textwidth}{N T A}
\hline
Name & Type & Annotation\\
\hline
\hfuzz=500pt\includegraphics[width=1em]{element-mustset.pdf}~outputfileAscii & \hfuzz=500pt filename & \hfuzz=500pt ASCII outputfile\\
\hfuzz=500pt\includegraphics[width=1em]{element-mustset.pdf}~inputfileAccelerometer & \hfuzz=500pt filename & \hfuzz=500pt GROOPS acceleromter file\\
\hfuzz=500pt\includegraphics[width=1em]{element-mustset.pdf}~satelliteIdentifier & \hfuzz=500pt string & \hfuzz=500pt satellite identifier (A or B for GRACE, C or D for GRACE-FO)\\
\hfuzz=500pt\includegraphics[width=1em]{element-unbounded.pdf}~globalAttributes & \hfuzz=500pt string & \hfuzz=500pt additional attributes as 'key: value' pairs\\
\hline
\end{tabularx}

\clearpage
%==================================
\subsection{GraceAod2DoodsonHarmonics}\label{GraceAod2DoodsonHarmonics}
This program converts the atmospheric and ocean tidal products (AOD1B)
from the GRACE SDS format into \configFile{outputfileDoodsonHarmonics}{doodsonHarmonic}.


\keepXColumns
\begin{tabularx}{\textwidth}{N T A}
\hline
Name & Type & Annotation\\
\hline
\hfuzz=500pt\includegraphics[width=1em]{element-mustset.pdf}~outputfileDoodsonHarmonics & \hfuzz=500pt filename & \hfuzz=500pt \\
\hfuzz=500pt\includegraphics[width=1em]{element-mustset.pdf}~inputfileTideGeneratingPotential & \hfuzz=500pt filename & \hfuzz=500pt to compute Xi phase correction\\
\hfuzz=500pt\includegraphics[width=1em]{element-mustset-unbounded.pdf}~inputfile & \hfuzz=500pt filename & \hfuzz=500pt \\
\hline
\end{tabularx}

\clearpage
%==================================
\subsection{GraceAod2TimeSplines}\label{GraceAod2TimeSplines}
This program converts the atmospheric and ocean de-aliasing product (AOD1B)
from the GRACE SDS format into \file{time spline files}{timeSplinesGravityField}.
Multiple \config{inputfile}s must be given in the correct time order.
A linear method is assumed for the interpolation between the given points in time.

The GRACE SDS format is described in "AOD1B Product Description Document"
given at \url{http://podaac.jpl.nasa.gov/grace/documentation.html}.


\keepXColumns
\begin{tabularx}{\textwidth}{N T A}
\hline
Name & Type & Annotation\\
\hline
\hfuzz=500pt\includegraphics[width=1em]{element.pdf}~outputfileDealiasing & \hfuzz=500pt filename & \hfuzz=500pt \\
\hfuzz=500pt\includegraphics[width=1em]{element.pdf}~outputfileAtmosphere & \hfuzz=500pt filename & \hfuzz=500pt \\
\hfuzz=500pt\includegraphics[width=1em]{element.pdf}~outputfileOcean & \hfuzz=500pt filename & \hfuzz=500pt \\
\hfuzz=500pt\includegraphics[width=1em]{element.pdf}~outputfileBottomPressure & \hfuzz=500pt filename & \hfuzz=500pt \\
\hfuzz=500pt\includegraphics[width=1em]{element.pdf}~outputfileMisc & \hfuzz=500pt filename & \hfuzz=500pt \\
\hfuzz=500pt\includegraphics[width=1em]{element-mustset-unbounded.pdf}~inputfile & \hfuzz=500pt filename & \hfuzz=500pt \\
\hline
\end{tabularx}

\clearpage
%==================================
\subsection{GraceCoefficients2BlockMeanTimeSplines}\label{GraceCoefficients2BlockMeanTimeSplines}
This program converts potential coefficients from the GRACE SDS RL06 format
into \configFile{outputfileTimeSplines}{timeSplinesGravityField}.

The \configFile{outputfileTimeSeries}{instrument} contains the mid points
of non-empty intervals and \configFile{outputfileTimeIntervals}{instrument}
contains the monthly interval boundaries from first to last solution.

The output will always be monthly block means. If the SDS solutions do vary or overlap,
the nearest solution in terms of reference epoch is used.


\keepXColumns
\begin{tabularx}{\textwidth}{N T A}
\hline
Name & Type & Annotation\\
\hline
\hfuzz=500pt\includegraphics[width=1em]{element-mustset.pdf}~outputfileTimeSplines & \hfuzz=500pt filename & \hfuzz=500pt \\
\hfuzz=500pt\includegraphics[width=1em]{element.pdf}~outputfileTimeSplinesCovariance & \hfuzz=500pt filename & \hfuzz=500pt only the variances are saved\\
\hfuzz=500pt\includegraphics[width=1em]{element.pdf}~outputfileTimeSeries & \hfuzz=500pt filename & \hfuzz=500pt mid points of non-empty intervals\\
\hfuzz=500pt\includegraphics[width=1em]{element.pdf}~outputfileTimeIntervals & \hfuzz=500pt filename & \hfuzz=500pt monthly interval boundaries from first to last solution\\
\hfuzz=500pt\includegraphics[width=1em]{element-mustset-unbounded.pdf}~inputfile & \hfuzz=500pt filename & \hfuzz=500pt \\
\hline
\end{tabularx}

\clearpage
%==================================
\subsection{GraceL1a2Accelerometer}\label{GraceL1a2Accelerometer}
This program converts Level-1A accelerometer data (ACC1A) to the GROOPS instrument file format.
The GRACE Level-1A format is described in \verb|GRACEiolib.h| given at
\url{http://podaac-tools.jpl.nasa.gov/drive/files/allData/grace/sw/GraceReadSW_L1_2010-03-31.tar.gz}.
The output is one arc of satellite data which can include data gaps.
To split the arc in multiple gap free arcs use \program{InstrumentSynchronize}.


\keepXColumns
\begin{tabularx}{\textwidth}{N T A}
\hline
Name & Type & Annotation\\
\hline
\hfuzz=500pt\includegraphics[width=1em]{element.pdf}~outputfileAccelerometer & \hfuzz=500pt filename & \hfuzz=500pt ACCELEROMETER in SRF\\
\hfuzz=500pt\includegraphics[width=1em]{element.pdf}~outputfileAngularAccelerometer & \hfuzz=500pt filename & \hfuzz=500pt ACCELEROMETER in SRF\\
\hfuzz=500pt\includegraphics[width=1em]{element-mustset-unbounded.pdf}~inputfile & \hfuzz=500pt filename & \hfuzz=500pt ACC1A\\
\hline
\end{tabularx}

\clearpage
%==================================
\subsection{GraceL1a2SatelliteTracking}\label{GraceL1a2SatelliteTracking}
This program converts Level-1A satellite tracking data (KBR1A) to the GROOPS instrument file format.
The GRACE Level-1A format is described in \verb|GRACEiolib.h| given at
\url{http://podaac-tools.jpl.nasa.gov/drive/files/allData/grace/sw/GraceReadSW_L1_2010-03-31.tar.gz}.
The output is one arc of satellite data which can include data gaps.
To split the arc in multiple gap free arcs use \program{InstrumentSynchronize}.


\keepXColumns
\begin{tabularx}{\textwidth}{N T A}
\hline
Name & Type & Annotation\\
\hline
\hfuzz=500pt\includegraphics[width=1em]{element.pdf}~outputfileSatelliteTracking & \hfuzz=500pt filename & \hfuzz=500pt MISCVALUES(ant\_id, K\_phase, Ka\_phase, K\_SNR, Ka\_SNR)\\
\hfuzz=500pt\includegraphics[width=1em]{element-mustset-unbounded.pdf}~inputfile & \hfuzz=500pt filename & \hfuzz=500pt KBR1A\\
\hline
\end{tabularx}

\clearpage
%==================================
\subsection{GraceL1a2StarCamera}\label{GraceL1a2StarCamera}
This program converts orientation data measured by the star cameras
from the GRACE Level-1A format (SCA1A) to the GROOPS instrument file format.
For further information see \program{GraceL1a2Accelerometer}.


\keepXColumns
\begin{tabularx}{\textwidth}{N T A}
\hline
Name & Type & Annotation\\
\hline
\hfuzz=500pt\includegraphics[width=1em]{element-mustset.pdf}~outputfileStarCamera1 & \hfuzz=500pt filename & \hfuzz=500pt STARCAMERA1A, head 1\\
\hfuzz=500pt\includegraphics[width=1em]{element-mustset.pdf}~outputfileStarCamera2 & \hfuzz=500pt filename & \hfuzz=500pt STARCAMERA1A, head 2\\
\hfuzz=500pt\includegraphics[width=1em]{element-mustset-unbounded.pdf}~inputfile & \hfuzz=500pt filename & \hfuzz=500pt SCA1A, !GRACE-FO is not working!\\
\hline
\end{tabularx}

\clearpage
%==================================
\subsection{GraceL1a2Temperature}\label{GraceL1a2Temperature}
This program converts Level-1A temperature measurments (HRT1B or HRT1A) to the GROOPS instrument file format.
The GRACE Level-1A format is described in GRACE given at \url{http://podaac-tools.jpl.nasa.gov/drive/files/allData/grace/sw/GraceReadSW_L1_2010-03-31.tar.gz}.
Multiple \config{inputfile}s must be given in the correct time order.
The output is one arc of satellite data which can include data gaps.
To split the arc in multiple gap free arcs use \program{InstrumentSynchronize}.


\keepXColumns
\begin{tabularx}{\textwidth}{N T A}
\hline
Name & Type & Annotation\\
\hline
\hfuzz=500pt\includegraphics[width=1em]{element-mustset.pdf}~outputfileTemperature & \hfuzz=500pt filename & \hfuzz=500pt MISCVALUES\\
\hfuzz=500pt\includegraphics[width=1em]{element-mustset-unbounded.pdf}~inputfile & \hfuzz=500pt filename & \hfuzz=500pt HRT1B or HRT1A\\
\hline
\end{tabularx}

\clearpage
%==================================
\subsection{GraceL1b2AccHousekeeping}\label{GraceL1b2AccHousekeeping}
This program converts ACC housekeeping data (AHK1B or AHK1A) from the GRACE SDS format into \file{instrument file (ACCHOUSEKEEPING)}{instrument}.
For further information see \program{GraceL1b2Accelerometer}.


\keepXColumns
\begin{tabularx}{\textwidth}{N T A}
\hline
Name & Type & Annotation\\
\hline
\hfuzz=500pt\includegraphics[width=1em]{element-mustset.pdf}~outputfileAccHousekeeping & \hfuzz=500pt filename & \hfuzz=500pt ACCHOUSEKEEPING\\
\hfuzz=500pt\includegraphics[width=1em]{element-mustset-unbounded.pdf}~inputfile & \hfuzz=500pt filename & \hfuzz=500pt AHK1B or AHK1A\\
\hline
\end{tabularx}

\clearpage
%==================================
\subsection{GraceL1b2Accelerometer}\label{GraceL1b2Accelerometer}
This program converts accelerometer data (ACC1B or ACT1B) from the GRACE SDS format into \file{instrument file (ACCELEROMETER)}{instrument}.

Multiple \config{inputfile}s must be given in the correct time order.
The output is one arc of satellite data which can include data gaps.
To split the arc in multiple gap free arcs use \program{InstrumentSynchronize}.

The GRACE SDS format is described in "GRACE Level 1B Data Product User Handbook JPL D-22027"
given at \url{https://podaac-tools.jpl.nasa.gov/drive/files/allData/grace/docs/Handbook_1B_v1.3.pdf}.


\keepXColumns
\begin{tabularx}{\textwidth}{N T A}
\hline
Name & Type & Annotation\\
\hline
\hfuzz=500pt\includegraphics[width=1em]{element.pdf}~outputfileAccelerometer & \hfuzz=500pt filename & \hfuzz=500pt ACCELEROMETER\\
\hfuzz=500pt\includegraphics[width=1em]{element.pdf}~outputfileAngularAccelerometer & \hfuzz=500pt filename & \hfuzz=500pt ACCELEROMETER\\
\hfuzz=500pt\includegraphics[width=1em]{element.pdf}~outputfileFlags & \hfuzz=500pt filename & \hfuzz=500pt MISCVALUES(qualflg, acl\_res.x, acl\_res.y, acl\_res.z)\\
\hfuzz=500pt\includegraphics[width=1em]{element-mustset-unbounded.pdf}~inputfile & \hfuzz=500pt filename & \hfuzz=500pt ACC1B or ACT1B\\
\hline
\end{tabularx}

\clearpage
%==================================
\subsection{GraceL1b2ClockOffset}\label{GraceL1b2ClockOffset}
This program converts clock data (CLK1B or LLK1B) from the GRACE SDS format into \file{instrument file (MISCVALUE)}{instrument}.
For further information see \program{GraceL1b2Accelerometer}.


\keepXColumns
\begin{tabularx}{\textwidth}{N T A}
\hline
Name & Type & Annotation\\
\hline
\hfuzz=500pt\includegraphics[width=1em]{element-mustset.pdf}~outputfileClock & \hfuzz=500pt filename & \hfuzz=500pt MISCVALUE\\
\hfuzz=500pt\includegraphics[width=1em]{element-mustset-unbounded.pdf}~inputfile & \hfuzz=500pt filename & \hfuzz=500pt CLK1B or LLK1B\\
\hline
\end{tabularx}

\clearpage
%==================================
\subsection{GraceL1b2GnssReceiver}\label{GraceL1b2GnssReceiver}
This program converts GPS receiver data (phase and pseudo range) data
from the GRACE SDS format (GPS1B or GPS1A) into \file{instrument file (GNSSRECEIVER)}{instrument}.
For further information see \program{GraceL1b2Accelerometer}.


\keepXColumns
\begin{tabularx}{\textwidth}{N T A}
\hline
Name & Type & Annotation\\
\hline
\hfuzz=500pt\includegraphics[width=1em]{element-mustset.pdf}~outputfileGnssReceiver & \hfuzz=500pt filename & \hfuzz=500pt GNSSRECEIVER\\
\hfuzz=500pt\includegraphics[width=1em]{element-mustset-unbounded.pdf}~inputfile & \hfuzz=500pt filename & \hfuzz=500pt GPS1B or GPS1A\\
\hline
\end{tabularx}

\clearpage
%==================================
\subsection{GraceL1b2Magnetometer}\label{GraceL1b2Magnetometer}
This program converts magnetometer data (MAG1B or MAG1A) from the GRACE SDS format into \file{instrument file (MAGNETOMETER)}{instrument}.
For further information see \program{GraceL1b2Accelerometer}.


\keepXColumns
\begin{tabularx}{\textwidth}{N T A}
\hline
Name & Type & Annotation\\
\hline
\hfuzz=500pt\includegraphics[width=1em]{element-mustset.pdf}~outputfileMagnetometer & \hfuzz=500pt filename & \hfuzz=500pt MAGNETOMETER\\
\hfuzz=500pt\includegraphics[width=1em]{element-mustset-unbounded.pdf}~inputfile & \hfuzz=500pt filename & \hfuzz=500pt MAG1B or MAG1A\\
\hline
\end{tabularx}

\clearpage
%==================================
\subsection{GraceL1b2Mass}\label{GraceL1b2Mass}
This program converts mass data (MAS1B or MAS1A) from the GRACE SDS format into \file{instrument file (MASS)}{instrument}.
For further information see \program{GraceL1b2Accelerometer}.


\keepXColumns
\begin{tabularx}{\textwidth}{N T A}
\hline
Name & Type & Annotation\\
\hline
\hfuzz=500pt\includegraphics[width=1em]{element-mustset.pdf}~outputfileMass & \hfuzz=500pt filename & \hfuzz=500pt MASS\\
\hfuzz=500pt\includegraphics[width=1em]{element-mustset-unbounded.pdf}~inputfile & \hfuzz=500pt filename & \hfuzz=500pt MAS1B or MAS1A\\
\hline
\end{tabularx}

\clearpage
%==================================
\subsection{GraceL1b2Orbit}\label{GraceL1b2Orbit}
This program converts the reduced dynamical orbit
from the GRACE/GRACE-FO SDS format (GNV1B, GNI1B) into \file{instrument file (ORBIT)}{instrument}.

When GNV1B is used, the orbit can be rotated from the terrestrial reference frame (TRF) transformed into the celestial reference frame (CRF) by
specifying \configClass{earthRotation}{earthRotationType}.

For further information see \program{GraceL1b2Accelerometer}.


\keepXColumns
\begin{tabularx}{\textwidth}{N T A}
\hline
Name & Type & Annotation\\
\hline
\hfuzz=500pt\includegraphics[width=1em]{element-mustset.pdf}~outputfileOrbit & \hfuzz=500pt filename & \hfuzz=500pt \\
\hfuzz=500pt\includegraphics[width=1em]{element.pdf}~earthRotation & \hfuzz=500pt \hyperref[earthRotationType]{earthRotation} & \hfuzz=500pt to rotate GNV1B into CRF\\
\hfuzz=500pt\includegraphics[width=1em]{element-mustset-unbounded.pdf}~inputfile & \hfuzz=500pt filename & \hfuzz=500pt GNV1B/GNI1B\\
\hline
\end{tabularx}

\clearpage
%==================================
\subsection{GraceL1b2SatelliteTracking}\label{GraceL1b2SatelliteTracking}
This program converts low-low satellite data measured by the K-band ranging system
from the GRACE SDS format (KBR1B or LRI1B) into \file{instrument file (SATELLITETRACKING)}{instrument}.
The \config{inputfile}s contain also corrections to antenna offsets
and the so called light time correction. The corrections can be stored in additional files
in the same format as the observations.
If a phase break is found an artificial gap is created.
For further information see \program{GraceL1b2Accelerometer}.


\keepXColumns
\begin{tabularx}{\textwidth}{N T A}
\hline
Name & Type & Annotation\\
\hline
\hfuzz=500pt\includegraphics[width=1em]{element.pdf}~outputfileSatelliteTracking & \hfuzz=500pt filename & \hfuzz=500pt SATELLITETRACKING\\
\hfuzz=500pt\includegraphics[width=1em]{element.pdf}~outputfileAntCentr & \hfuzz=500pt filename & \hfuzz=500pt SATELLITETRACKING\\
\hfuzz=500pt\includegraphics[width=1em]{element.pdf}~outputfileLighttime & \hfuzz=500pt filename & \hfuzz=500pt SATELLITETRACKING\\
\hfuzz=500pt\includegraphics[width=1em]{element.pdf}~outputfileSNR & \hfuzz=500pt filename & \hfuzz=500pt MISCVALUES(K\_A\_SNR, Ka\_A\_SNR, K\_B\_SNR, Ka\_B\_SNR, qualflg)\\
\hfuzz=500pt\includegraphics[width=1em]{element.pdf}~outputfileIonoCorr & \hfuzz=500pt filename & \hfuzz=500pt MISCVALUE\\
\hfuzz=500pt\includegraphics[width=1em]{element-mustset-unbounded.pdf}~inputfile & \hfuzz=500pt filename & \hfuzz=500pt KBR1B or LRI1B\\
\hline
\end{tabularx}

\clearpage
%==================================
\subsection{GraceL1b2StarCamera}\label{GraceL1b2StarCamera}
This program converts orientation data measured by a star camera (SRF to CRF)
from the GRACE SDS format (SCA1B) into \file{instrument file (STARCAMERA)}{instrument}.
For further information see \program{GraceL1b2Accelerometer}.


\keepXColumns
\begin{tabularx}{\textwidth}{N T A}
\hline
Name & Type & Annotation\\
\hline
\hfuzz=500pt\includegraphics[width=1em]{element.pdf}~outputfileStarCamera & \hfuzz=500pt filename & \hfuzz=500pt \\
\hfuzz=500pt\includegraphics[width=1em]{element.pdf}~outputfileStarCameraFlags & \hfuzz=500pt filename & \hfuzz=500pt MISCVALUES(sca\_id, qual\_rss, qualflg)\\
\hfuzz=500pt\includegraphics[width=1em]{element-mustset-unbounded.pdf}~inputfile & \hfuzz=500pt filename & \hfuzz=500pt SCA1B\\
\hline
\end{tabularx}

\clearpage
%==================================
\subsection{GraceL1b2StarCameraCovariance}\label{GraceL1b2StarCameraCovariance}
This program computes star camera covariance matrices (\file{instrument file, COVARIANE3D}{instrument})
for a GRACE satellite under consideration of the active camera heads and an a priori variance factor.


\keepXColumns
\begin{tabularx}{\textwidth}{N T A}
\hline
Name & Type & Annotation\\
\hline
\hfuzz=500pt\includegraphics[width=1em]{element-mustset.pdf}~outputfileStarCameraCovariance & \hfuzz=500pt filename & \hfuzz=500pt \\
\hfuzz=500pt\includegraphics[width=1em]{element-mustset.pdf}~inputfileStarCameraFlags & \hfuzz=500pt filename & \hfuzz=500pt \\
\hfuzz=500pt\includegraphics[width=1em]{element.pdf}~inputfileSequenceOfEventsQSA & \hfuzz=500pt filename & \hfuzz=500pt \\
\hfuzz=500pt\includegraphics[width=1em]{element.pdf}~sigma0 & \hfuzz=500pt double & \hfuzz=500pt [seconds of arc]\\
\hline
\end{tabularx}

This program is \reference{parallelized}{general.parallelization}.
\clearpage
%==================================
\subsection{GraceL1b2SteeringMirror}\label{GraceL1b2SteeringMirror}
This program converts GRACE-FO Steering Mirror output (LSM1B) to an \file{instrument file (STARCAMERA)}{instrument}.


\keepXColumns
\begin{tabularx}{\textwidth}{N T A}
\hline
Name & Type & Annotation\\
\hline
\hfuzz=500pt\includegraphics[width=1em]{element.pdf}~outputfileStarCamera & \hfuzz=500pt filename & \hfuzz=500pt \\
\hfuzz=500pt\includegraphics[width=1em]{element-mustset-unbounded.pdf}~inputfile & \hfuzz=500pt filename & \hfuzz=500pt LSM1B\\
\hline
\end{tabularx}

\clearpage
%==================================
\subsection{GraceL1b2Thruster}\label{GraceL1b2Thruster}
This program converts thruster data (THR1B or THR1A) from the GRACE SDS format into \file{instrument file (THRUSTER)}{instrument}.
For further information see \program{GraceL1b2Accelerometer}.


\keepXColumns
\begin{tabularx}{\textwidth}{N T A}
\hline
Name & Type & Annotation\\
\hline
\hfuzz=500pt\includegraphics[width=1em]{element-mustset.pdf}~outputfileThruster & \hfuzz=500pt filename & \hfuzz=500pt THRUSTER\\
\hfuzz=500pt\includegraphics[width=1em]{element-mustset-unbounded.pdf}~inputfile & \hfuzz=500pt filename & \hfuzz=500pt THR1B or THR1A\\
\hline
\end{tabularx}

\clearpage
%==================================
\subsection{GraceL1b2TimeOffset}\label{GraceL1b2TimeOffset}
This program converts time data (TIM1A or TIM1B) from the GRACE SDS format into \file{instrument file (MISCVALUE)}{instrument}.
For further information see \program{GraceL1b2Accelerometer}.


\keepXColumns
\begin{tabularx}{\textwidth}{N T A}
\hline
Name & Type & Annotation\\
\hline
\hfuzz=500pt\includegraphics[width=1em]{element-mustset.pdf}~outputfileTime & \hfuzz=500pt filename & \hfuzz=500pt MISCVALUE\\
\hfuzz=500pt\includegraphics[width=1em]{element-mustset.pdf}~fractionalScale & \hfuzz=500pt double & \hfuzz=500pt 1e-6 for GRACE, 1e-9 for GRACE-FO\\
\hfuzz=500pt\includegraphics[width=1em]{element-mustset-unbounded.pdf}~inputfile & \hfuzz=500pt filename & \hfuzz=500pt TIM1A or TIM1B\\
\hline
\end{tabularx}

\clearpage
%==================================
\subsection{GraceL1b2Uso}\label{GraceL1b2Uso}
This program converts clock data (USO1B) from the GRACE SDS format into \file{instrument file (MISCVALUES)}{instrument}.
For further information see \program{GraceL1b2Accelerometer}.


\keepXColumns
\begin{tabularx}{\textwidth}{N T A}
\hline
Name & Type & Annotation\\
\hline
\hfuzz=500pt\includegraphics[width=1em]{element-mustset.pdf}~outputfileUso & \hfuzz=500pt filename & \hfuzz=500pt MISCVALUES(uso\_freq, K\_freq, Ka\_freq)\\
\hfuzz=500pt\includegraphics[width=1em]{element-mustset-unbounded.pdf}~inputfile & \hfuzz=500pt filename & \hfuzz=500pt USO1B\\
\hline
\end{tabularx}

\clearpage
%==================================
\subsection{GraceL1b2Vector}\label{GraceL1b2Vector}
This program reads vector orientation data (positions of instruments in the satellite frame) from the GRACE SDS format
(VGB1B, VGN1B, VGO1B, VKB1B, or VCM1B).
The \configFile{outputfileVector}{matrix} is a $(3n\times1)$ matrix containing $(x,y,z)$ for each record.
The GRACE SDS format is described in "GRACE Level 1B Data Product User Handbook JPL D-22027"
given at \url{http://podaac.jpl.nasa.gov/grace/documentation.html}.


\keepXColumns
\begin{tabularx}{\textwidth}{N T A}
\hline
Name & Type & Annotation\\
\hline
\hfuzz=500pt\includegraphics[width=1em]{element-mustset.pdf}~outputfileVector & \hfuzz=500pt filename & \hfuzz=500pt \\
\hfuzz=500pt\includegraphics[width=1em]{element-mustset.pdf}~inputfile & \hfuzz=500pt filename & \hfuzz=500pt VGB1B, VGN1B, VGO1B, VKB1B, or VCM1B\\
\hline
\end{tabularx}

\clearpage
%==================================
\subsection{GraceSequenceOfEvents}\label{GraceSequenceOfEvents}
This program converts the GRACE SOE (sequence of events) file/format into \file{instrument file (MISCVALUES)}{instrument}.
The GRACE SOE format is described in "GRACE Level 1B Data Product User Handbook JPL D-22027" and "TN-03\_SOE\_format.txt"
given at \url{http://podaac.jpl.nasa.gov/grace/documentation.html}.
The output is one arc of satellite data which can include data gaps.


\keepXColumns
\begin{tabularx}{\textwidth}{N T A}
\hline
Name & Type & Annotation\\
\hline
\hfuzz=500pt\includegraphics[width=1em]{element.pdf}~outputfileGraceA & \hfuzz=500pt filename & \hfuzz=500pt \\
\hfuzz=500pt\includegraphics[width=1em]{element.pdf}~outputfileGraceB & \hfuzz=500pt filename & \hfuzz=500pt \\
\hfuzz=500pt\includegraphics[width=1em]{element-mustset.pdf}~inputfile & \hfuzz=500pt filename & \hfuzz=500pt SoE file\\
\hfuzz=500pt\includegraphics[width=1em]{element-mustset.pdf}~events & \hfuzz=500pt choice & \hfuzz=500pt \\
\hfuzz=500pt\includegraphics[width=1em]{connector.pdf}\includegraphics[width=1em]{element-mustset.pdf}~ACCT & \hfuzz=500pt sequence & \hfuzz=500pt DSHL HeaterDisconnect\\
\hfuzz=500pt\quad\includegraphics[width=1em]{connector.pdf}\includegraphics[width=1em]{element.pdf}~mode & \hfuzz=500pt choice & \hfuzz=500pt \\
\hfuzz=500pt\quad\quad\includegraphics[width=1em]{connector.pdf}\includegraphics[width=1em]{element-mustset.pdf}~Heater & \hfuzz=500pt  & \hfuzz=500pt DSHL HeaterDisconnect\\
\hfuzz=500pt\quad\quad\includegraphics[width=1em]{connector.pdf}\includegraphics[width=1em]{element-mustset.pdf}~SetPoint & \hfuzz=500pt  & \hfuzz=500pt temperature set point\\
\hfuzz=500pt\includegraphics[width=1em]{connector.pdf}\includegraphics[width=1em]{element-mustset.pdf}~AOCS & \hfuzz=500pt sequence & \hfuzz=500pt coarse pointing mode or attitude hold mode\\
\hfuzz=500pt\quad\includegraphics[width=1em]{connector.pdf}\includegraphics[width=1em]{element.pdf}~mode & \hfuzz=500pt choice & \hfuzz=500pt \\
\hfuzz=500pt\quad\quad\includegraphics[width=1em]{connector.pdf}\includegraphics[width=1em]{element-mustset.pdf}~CPM & \hfuzz=500pt  & \hfuzz=500pt coarse pointing mode\\
\hfuzz=500pt\quad\quad\includegraphics[width=1em]{connector.pdf}\includegraphics[width=1em]{element-mustset.pdf}~AHM & \hfuzz=500pt  & \hfuzz=500pt attitude hold mode\\
\hfuzz=500pt\quad\quad\includegraphics[width=1em]{connector.pdf}\includegraphics[width=1em]{element-mustset.pdf}~SM & \hfuzz=500pt  & \hfuzz=500pt science mode\\
\hfuzz=500pt\includegraphics[width=1em]{connector.pdf}\includegraphics[width=1em]{element-mustset.pdf}~ACCR & \hfuzz=500pt  & \hfuzz=500pt ACCR\\
\hfuzz=500pt\includegraphics[width=1em]{connector.pdf}\includegraphics[width=1em]{element-mustset.pdf}~CMCAL & \hfuzz=500pt sequence & \hfuzz=500pt CoM calibration maneuver\\
\hfuzz=500pt\quad\includegraphics[width=1em]{connector.pdf}\includegraphics[width=1em]{element.pdf}~sampling & \hfuzz=500pt double & \hfuzz=500pt [seconds] create events between start and end of maneuver\\
\hfuzz=500pt\includegraphics[width=1em]{connector.pdf}\includegraphics[width=1em]{element-mustset.pdf}~KBRCAL & \hfuzz=500pt sequence & \hfuzz=500pt KBR calibration maneuver\\
\hfuzz=500pt\quad\includegraphics[width=1em]{connector.pdf}\includegraphics[width=1em]{element.pdf}~sampling & \hfuzz=500pt double & \hfuzz=500pt [seconds] create events between start and end of maneuver\\
\hfuzz=500pt\includegraphics[width=1em]{connector.pdf}\includegraphics[width=1em]{element-mustset.pdf}~VCM & \hfuzz=500pt  & \hfuzz=500pt CoM coordinates in SRF (m)\\
\hfuzz=500pt\includegraphics[width=1em]{connector.pdf}\includegraphics[width=1em]{element-mustset.pdf}~VKB & \hfuzz=500pt  & \hfuzz=500pt KBR phase center coordinates in SRF (m)\\
\hfuzz=500pt\includegraphics[width=1em]{connector.pdf}\includegraphics[width=1em]{element-mustset.pdf}~ICUVP & \hfuzz=500pt  & \hfuzz=500pt ICUVP\\
\hfuzz=500pt\includegraphics[width=1em]{connector.pdf}\includegraphics[width=1em]{element-mustset.pdf}~IPU & \hfuzz=500pt  & \hfuzz=500pt IPU\\
\hfuzz=500pt\includegraphics[width=1em]{connector.pdf}\includegraphics[width=1em]{element-mustset.pdf}~IPUR & \hfuzz=500pt  & \hfuzz=500pt IPUR\\
\hfuzz=500pt\includegraphics[width=1em]{connector.pdf}\includegraphics[width=1em]{element-mustset.pdf}~KAMI & \hfuzz=500pt  & \hfuzz=500pt KAMI: time tag offset to Ka-phase meas.\\
\hfuzz=500pt\includegraphics[width=1em]{connector.pdf}\includegraphics[width=1em]{element-mustset.pdf}~KMI & \hfuzz=500pt  & \hfuzz=500pt K\_MI: time tag offset to K-phase meas.\\
\hfuzz=500pt\includegraphics[width=1em]{connector.pdf}\includegraphics[width=1em]{element-mustset.pdf}~KTOFF & \hfuzz=500pt  & \hfuzz=500pt KTOFF: time tag offset to KBR meas.\\
\hfuzz=500pt\includegraphics[width=1em]{connector.pdf}\includegraphics[width=1em]{element-mustset.pdf}~MANV & \hfuzz=500pt  & \hfuzz=500pt MANV\\
\hfuzz=500pt\includegraphics[width=1em]{connector.pdf}\includegraphics[width=1em]{element-mustset.pdf}~MTE1 & \hfuzz=500pt  & \hfuzz=500pt MTE1\\
\hfuzz=500pt\includegraphics[width=1em]{connector.pdf}\includegraphics[width=1em]{element-mustset.pdf}~MTE2 & \hfuzz=500pt  & \hfuzz=500pt MTE2\\
\hfuzz=500pt\includegraphics[width=1em]{connector.pdf}\includegraphics[width=1em]{element-mustset.pdf}~OCC & \hfuzz=500pt  & \hfuzz=500pt OCC\\
\hfuzz=500pt\includegraphics[width=1em]{connector.pdf}\includegraphics[width=1em]{element-mustset.pdf}~QSA & \hfuzz=500pt  & \hfuzz=500pt SCA to SRF frame rotation\\
\hfuzz=500pt\includegraphics[width=1em]{connector.pdf}\includegraphics[width=1em]{element-mustset.pdf}~QKS & \hfuzz=500pt  & \hfuzz=500pt SCA to KBR frame rotation\\
\hline
\end{tabularx}

\clearpage
%==================================
\subsection{GrailCdr2Orbit}\label{GrailCdr2Orbit}
This program converts the orbit from the GRAIL SDS format into  \file{instrument file (ORBIT)}{instrument}.


\keepXColumns
\begin{tabularx}{\textwidth}{N T A}
\hline
Name & Type & Annotation\\
\hline
\hfuzz=500pt\includegraphics[width=1em]{element-mustset.pdf}~outputfileOrbit & \hfuzz=500pt filename & \hfuzz=500pt \\
\hfuzz=500pt\includegraphics[width=1em]{element-mustset-unbounded.pdf}~inputfile & \hfuzz=500pt filename & \hfuzz=500pt \\
\hline
\end{tabularx}

\clearpage
%==================================
\subsection{GrailCdr2SatelliteTracking}\label{GrailCdr2SatelliteTracking}
This program converts low-low satellite data measured by the K-band ranging system
from the GRAIL format into \file{instrument file (SATELLITETRACKING)}{instrument}.
The \config{inputfile}s contain also corrections for antenna offsets
and the so called light time correction.
The corrections can be stored in additional files in the same format as the observations.
If a phase break is found an artificial gap is created.


\keepXColumns
\begin{tabularx}{\textwidth}{N T A}
\hline
Name & Type & Annotation\\
\hline
\hfuzz=500pt\includegraphics[width=1em]{element.pdf}~outputfileSatelliteTracking & \hfuzz=500pt filename & \hfuzz=500pt \\
\hfuzz=500pt\includegraphics[width=1em]{element.pdf}~outputfileAntCentr & \hfuzz=500pt filename & \hfuzz=500pt \\
\hfuzz=500pt\includegraphics[width=1em]{element.pdf}~outputfileLighttime & \hfuzz=500pt filename & \hfuzz=500pt \\
\hfuzz=500pt\includegraphics[width=1em]{element.pdf}~outputfileTemperature & \hfuzz=500pt filename & \hfuzz=500pt \\
\hfuzz=500pt\includegraphics[width=1em]{element.pdf}~approximateTimeBias & \hfuzz=500pt double & \hfuzz=500pt [seconds]\\
\hfuzz=500pt\includegraphics[width=1em]{element-mustset-unbounded.pdf}~inputfile & \hfuzz=500pt filename & \hfuzz=500pt \\
\hline
\end{tabularx}

\clearpage
%==================================
\subsection{GrailCdr2StarCamera}\label{GrailCdr2StarCamera}
This program converts orientation data measured by a star camera (SRF to CRF)
from the GRAIL SDS format into \file{instrument file (STARCAMERA)}{instrument}.
For further information see \program{GraceL1b2Accelerometer}.


\keepXColumns
\begin{tabularx}{\textwidth}{N T A}
\hline
Name & Type & Annotation\\
\hline
\hfuzz=500pt\includegraphics[width=1em]{element.pdf}~outputfileStarCamera & \hfuzz=500pt filename & \hfuzz=500pt \\
\hfuzz=500pt\includegraphics[width=1em]{element-mustset-unbounded.pdf}~inputfile & \hfuzz=500pt filename & \hfuzz=500pt \\
\hline
\end{tabularx}

\clearpage
%==================================
\subsection{GriddedData2NetCdf}\label{GriddedData2NetCdf}
This program converts a \configFile{inputfileGriddedData}{griddedData}
to a COARDS compliant NetCDF file. The output data can be defined with \config{dataVariable}.
You should add at least the attributes \verb|units|, \verb|long_name|, and maybe \verb|_FillValue|
to the variables. For the \config{dataVariable:value} the standard \reference{dataVariables}{general.parser:dataVariables}
are available to select the data columns of \configFile{inputfileGriddedData}{griddedData}.

See also \program{NetCdfInfo}, \program{GriddedDataTimeSeries2NetCdf}, \program{NetCdf2GriddedData}.


\keepXColumns
\begin{tabularx}{\textwidth}{N T A}
\hline
Name & Type & Annotation\\
\hline
\hfuzz=500pt\includegraphics[width=1em]{element-mustset.pdf}~outputfileNetCdf & \hfuzz=500pt filename & \hfuzz=500pt file name of NetCDF output\\
\hfuzz=500pt\includegraphics[width=1em]{element-mustset.pdf}~inputfileGriddedData & \hfuzz=500pt filename & \hfuzz=500pt input grid sequence\\
\hfuzz=500pt\includegraphics[width=1em]{element-mustset-unbounded.pdf}~dataVariable & \hfuzz=500pt sequence & \hfuzz=500pt metadata for data variables\\
\hfuzz=500pt\includegraphics[width=1em]{connector.pdf}\includegraphics[width=1em]{element-mustset.pdf}~name & \hfuzz=500pt string & \hfuzz=500pt netCDF variable name\\
\hfuzz=500pt\includegraphics[width=1em]{connector.pdf}\includegraphics[width=1em]{element.pdf}~value & \hfuzz=500pt expression & \hfuzz=500pt expression (variables 'height', 'data', 'longitude', 'latitude' and, 'area' are taken from the gridded data\\
\hfuzz=500pt\includegraphics[width=1em]{connector.pdf}\includegraphics[width=1em]{element-mustset.pdf}~dataType & \hfuzz=500pt choice & \hfuzz=500pt \\
\hfuzz=500pt\quad\includegraphics[width=1em]{connector.pdf}\includegraphics[width=1em]{element-mustset.pdf}~double & \hfuzz=500pt  & \hfuzz=500pt \\
\hfuzz=500pt\quad\includegraphics[width=1em]{connector.pdf}\includegraphics[width=1em]{element-mustset.pdf}~float & \hfuzz=500pt  & \hfuzz=500pt \\
\hfuzz=500pt\quad\includegraphics[width=1em]{connector.pdf}\includegraphics[width=1em]{element-mustset.pdf}~int & \hfuzz=500pt  & \hfuzz=500pt \\
\hfuzz=500pt\includegraphics[width=1em]{connector.pdf}\includegraphics[width=1em]{element-unbounded.pdf}~attribute & \hfuzz=500pt choice & \hfuzz=500pt netCDF attributes\\
\hfuzz=500pt\quad\includegraphics[width=1em]{connector.pdf}\includegraphics[width=1em]{element-mustset.pdf}~text & \hfuzz=500pt sequence & \hfuzz=500pt \\
\hfuzz=500pt\quad\quad\includegraphics[width=1em]{connector.pdf}\includegraphics[width=1em]{element-mustset.pdf}~name & \hfuzz=500pt string & \hfuzz=500pt \\
\hfuzz=500pt\quad\quad\includegraphics[width=1em]{connector.pdf}\includegraphics[width=1em]{element-mustset.pdf}~value & \hfuzz=500pt string & \hfuzz=500pt \\
\hfuzz=500pt\quad\includegraphics[width=1em]{connector.pdf}\includegraphics[width=1em]{element-mustset.pdf}~value & \hfuzz=500pt sequence & \hfuzz=500pt \\
\hfuzz=500pt\quad\quad\includegraphics[width=1em]{connector.pdf}\includegraphics[width=1em]{element-mustset.pdf}~name & \hfuzz=500pt string & \hfuzz=500pt \\
\hfuzz=500pt\quad\quad\includegraphics[width=1em]{connector.pdf}\includegraphics[width=1em]{element-mustset-unbounded.pdf}~value & \hfuzz=500pt double & \hfuzz=500pt \\
\hfuzz=500pt\quad\quad\includegraphics[width=1em]{connector.pdf}\includegraphics[width=1em]{element-mustset.pdf}~dataType & \hfuzz=500pt choice & \hfuzz=500pt \\
\hfuzz=500pt\quad\quad\quad\includegraphics[width=1em]{connector.pdf}\includegraphics[width=1em]{element-mustset.pdf}~double & \hfuzz=500pt  & \hfuzz=500pt \\
\hfuzz=500pt\quad\quad\quad\includegraphics[width=1em]{connector.pdf}\includegraphics[width=1em]{element-mustset.pdf}~float & \hfuzz=500pt  & \hfuzz=500pt \\
\hfuzz=500pt\quad\quad\quad\includegraphics[width=1em]{connector.pdf}\includegraphics[width=1em]{element-mustset.pdf}~int & \hfuzz=500pt  & \hfuzz=500pt \\
\hfuzz=500pt\includegraphics[width=1em]{element-unbounded.pdf}~globalAttribute & \hfuzz=500pt choice & \hfuzz=500pt additional meta data\\
\hfuzz=500pt\includegraphics[width=1em]{connector.pdf}\includegraphics[width=1em]{element-mustset.pdf}~text & \hfuzz=500pt sequence & \hfuzz=500pt \\
\hfuzz=500pt\quad\includegraphics[width=1em]{connector.pdf}\includegraphics[width=1em]{element-mustset.pdf}~name & \hfuzz=500pt string & \hfuzz=500pt \\
\hfuzz=500pt\quad\includegraphics[width=1em]{connector.pdf}\includegraphics[width=1em]{element-mustset.pdf}~value & \hfuzz=500pt string & \hfuzz=500pt \\
\hfuzz=500pt\includegraphics[width=1em]{connector.pdf}\includegraphics[width=1em]{element-mustset.pdf}~value & \hfuzz=500pt sequence & \hfuzz=500pt \\
\hfuzz=500pt\quad\includegraphics[width=1em]{connector.pdf}\includegraphics[width=1em]{element-mustset.pdf}~name & \hfuzz=500pt string & \hfuzz=500pt \\
\hfuzz=500pt\quad\includegraphics[width=1em]{connector.pdf}\includegraphics[width=1em]{element-mustset-unbounded.pdf}~value & \hfuzz=500pt double & \hfuzz=500pt \\
\hfuzz=500pt\quad\includegraphics[width=1em]{connector.pdf}\includegraphics[width=1em]{element-mustset.pdf}~dataType & \hfuzz=500pt choice & \hfuzz=500pt \\
\hfuzz=500pt\quad\quad\includegraphics[width=1em]{connector.pdf}\includegraphics[width=1em]{element-mustset.pdf}~double & \hfuzz=500pt  & \hfuzz=500pt \\
\hfuzz=500pt\quad\quad\includegraphics[width=1em]{connector.pdf}\includegraphics[width=1em]{element-mustset.pdf}~float & \hfuzz=500pt  & \hfuzz=500pt \\
\hfuzz=500pt\quad\quad\includegraphics[width=1em]{connector.pdf}\includegraphics[width=1em]{element-mustset.pdf}~int & \hfuzz=500pt  & \hfuzz=500pt \\
\hline
\end{tabularx}

\clearpage
%==================================
\subsection{GriddedDataTimeSeries2NetCdf}\label{GriddedDataTimeSeries2NetCdf}
Read a \configFile{inputfileGriddedDataTimeSeries}{griddedDataTimeSeries}
and converts it to a COARDS compliant NetCDF file.

The output data can be defined with \config{dataVariable}.
You should add at least the attributes \verb|units|, \verb|long_name|, and maybe \verb|_FillValue|
to the variables. The \config{dataVariable:inputColumn} selects the data from the input file.

If \configClass{timeSeries}{timeSeriesType} is not set
the temporal nodal points from the inputfile are used.

See also \program{NetCdfInfo}, \program{GriddedData2NetCdf}, \program{NetCdf2GriddedDataTimeSeries}.


\keepXColumns
\begin{tabularx}{\textwidth}{N T A}
\hline
Name & Type & Annotation\\
\hline
\hfuzz=500pt\includegraphics[width=1em]{element-mustset.pdf}~outputfileNetCdf & \hfuzz=500pt filename & \hfuzz=500pt file name of NetCDF output\\
\hfuzz=500pt\includegraphics[width=1em]{element-mustset.pdf}~inputfileGriddedDataTimeSeries & \hfuzz=500pt filename & \hfuzz=500pt \\
\hfuzz=500pt\includegraphics[width=1em]{element-unbounded.pdf}~timeSeries & \hfuzz=500pt \hyperref[timeSeriesType]{timeSeries} & \hfuzz=500pt otherwise times from inputfile are used\\
\hfuzz=500pt\includegraphics[width=1em]{element-mustset-unbounded.pdf}~dataVariable & \hfuzz=500pt sequence & \hfuzz=500pt metadata for data variables\\
\hfuzz=500pt\includegraphics[width=1em]{connector.pdf}\includegraphics[width=1em]{element-mustset.pdf}~name & \hfuzz=500pt string & \hfuzz=500pt netCDF variable name\\
\hfuzz=500pt\includegraphics[width=1em]{connector.pdf}\includegraphics[width=1em]{element.pdf}~inputColumn & \hfuzz=500pt uint & \hfuzz=500pt input data column\\
\hfuzz=500pt\includegraphics[width=1em]{connector.pdf}\includegraphics[width=1em]{element-mustset.pdf}~dataType & \hfuzz=500pt choice & \hfuzz=500pt \\
\hfuzz=500pt\quad\includegraphics[width=1em]{connector.pdf}\includegraphics[width=1em]{element-mustset.pdf}~double & \hfuzz=500pt  & \hfuzz=500pt \\
\hfuzz=500pt\quad\includegraphics[width=1em]{connector.pdf}\includegraphics[width=1em]{element-mustset.pdf}~float & \hfuzz=500pt  & \hfuzz=500pt \\
\hfuzz=500pt\quad\includegraphics[width=1em]{connector.pdf}\includegraphics[width=1em]{element-mustset.pdf}~int & \hfuzz=500pt  & \hfuzz=500pt \\
\hfuzz=500pt\includegraphics[width=1em]{connector.pdf}\includegraphics[width=1em]{element-unbounded.pdf}~attribute & \hfuzz=500pt choice & \hfuzz=500pt netCDF attributes\\
\hfuzz=500pt\quad\includegraphics[width=1em]{connector.pdf}\includegraphics[width=1em]{element-mustset.pdf}~text & \hfuzz=500pt sequence & \hfuzz=500pt \\
\hfuzz=500pt\quad\quad\includegraphics[width=1em]{connector.pdf}\includegraphics[width=1em]{element-mustset.pdf}~name & \hfuzz=500pt string & \hfuzz=500pt \\
\hfuzz=500pt\quad\quad\includegraphics[width=1em]{connector.pdf}\includegraphics[width=1em]{element-mustset.pdf}~value & \hfuzz=500pt string & \hfuzz=500pt \\
\hfuzz=500pt\quad\includegraphics[width=1em]{connector.pdf}\includegraphics[width=1em]{element-mustset.pdf}~value & \hfuzz=500pt sequence & \hfuzz=500pt \\
\hfuzz=500pt\quad\quad\includegraphics[width=1em]{connector.pdf}\includegraphics[width=1em]{element-mustset.pdf}~name & \hfuzz=500pt string & \hfuzz=500pt \\
\hfuzz=500pt\quad\quad\includegraphics[width=1em]{connector.pdf}\includegraphics[width=1em]{element-mustset-unbounded.pdf}~value & \hfuzz=500pt double & \hfuzz=500pt \\
\hfuzz=500pt\quad\quad\includegraphics[width=1em]{connector.pdf}\includegraphics[width=1em]{element-mustset.pdf}~dataType & \hfuzz=500pt choice & \hfuzz=500pt \\
\hfuzz=500pt\quad\quad\quad\includegraphics[width=1em]{connector.pdf}\includegraphics[width=1em]{element-mustset.pdf}~double & \hfuzz=500pt  & \hfuzz=500pt \\
\hfuzz=500pt\quad\quad\quad\includegraphics[width=1em]{connector.pdf}\includegraphics[width=1em]{element-mustset.pdf}~float & \hfuzz=500pt  & \hfuzz=500pt \\
\hfuzz=500pt\quad\quad\quad\includegraphics[width=1em]{connector.pdf}\includegraphics[width=1em]{element-mustset.pdf}~int & \hfuzz=500pt  & \hfuzz=500pt \\
\hfuzz=500pt\includegraphics[width=1em]{element-unbounded.pdf}~globalAttribute & \hfuzz=500pt choice & \hfuzz=500pt additional meta data\\
\hfuzz=500pt\includegraphics[width=1em]{connector.pdf}\includegraphics[width=1em]{element-mustset.pdf}~text & \hfuzz=500pt sequence & \hfuzz=500pt \\
\hfuzz=500pt\quad\includegraphics[width=1em]{connector.pdf}\includegraphics[width=1em]{element-mustset.pdf}~name & \hfuzz=500pt string & \hfuzz=500pt \\
\hfuzz=500pt\quad\includegraphics[width=1em]{connector.pdf}\includegraphics[width=1em]{element-mustset.pdf}~value & \hfuzz=500pt string & \hfuzz=500pt \\
\hfuzz=500pt\includegraphics[width=1em]{connector.pdf}\includegraphics[width=1em]{element-mustset.pdf}~value & \hfuzz=500pt sequence & \hfuzz=500pt \\
\hfuzz=500pt\quad\includegraphics[width=1em]{connector.pdf}\includegraphics[width=1em]{element-mustset.pdf}~name & \hfuzz=500pt string & \hfuzz=500pt \\
\hfuzz=500pt\quad\includegraphics[width=1em]{connector.pdf}\includegraphics[width=1em]{element-mustset-unbounded.pdf}~value & \hfuzz=500pt double & \hfuzz=500pt \\
\hfuzz=500pt\quad\includegraphics[width=1em]{connector.pdf}\includegraphics[width=1em]{element-mustset.pdf}~dataType & \hfuzz=500pt choice & \hfuzz=500pt \\
\hfuzz=500pt\quad\quad\includegraphics[width=1em]{connector.pdf}\includegraphics[width=1em]{element-mustset.pdf}~double & \hfuzz=500pt  & \hfuzz=500pt \\
\hfuzz=500pt\quad\quad\includegraphics[width=1em]{connector.pdf}\includegraphics[width=1em]{element-mustset.pdf}~float & \hfuzz=500pt  & \hfuzz=500pt \\
\hfuzz=500pt\quad\quad\includegraphics[width=1em]{connector.pdf}\includegraphics[width=1em]{element-mustset.pdf}~int & \hfuzz=500pt  & \hfuzz=500pt \\
\hline
\end{tabularx}

\clearpage
%==================================
\subsection{GroopsAscii2Orbit}\label{GroopsAscii2Orbit}
Read Orbits given in groops kinematic orbit ASCII format with covariance information.

See also \program{Orbit2GroopsAscii}.


\keepXColumns
\begin{tabularx}{\textwidth}{N T A}
\hline
Name & Type & Annotation\\
\hline
\hfuzz=500pt\includegraphics[width=1em]{element.pdf}~outputfileOrbit & \hfuzz=500pt filename & \hfuzz=500pt \\
\hfuzz=500pt\includegraphics[width=1em]{element.pdf}~outputfileCovariance & \hfuzz=500pt filename & \hfuzz=500pt \\
\hfuzz=500pt\includegraphics[width=1em]{element-mustset.pdf}~earthRotation & \hfuzz=500pt \hyperref[earthRotationType]{earthRotation} & \hfuzz=500pt \\
\hfuzz=500pt\includegraphics[width=1em]{element-mustset-unbounded.pdf}~inputfile & \hfuzz=500pt filename & \hfuzz=500pt \\
\hline
\end{tabularx}

\clearpage
%==================================
\subsection{Hw2TideGeneratingPotential}\label{Hw2TideGeneratingPotential}
Write \file{tide generating potential}{tideGeneratingPotential}
from Hartmann and Wenzel 1995 file, \url{https://doi.org/10.1029/95GL03324}.


\keepXColumns
\begin{tabularx}{\textwidth}{N T A}
\hline
Name & Type & Annotation\\
\hline
\hfuzz=500pt\includegraphics[width=1em]{element-mustset.pdf}~outputfileTideGeneratingPotential & \hfuzz=500pt filename & \hfuzz=500pt \\
\hfuzz=500pt\includegraphics[width=1em]{element-mustset.pdf}~inputfile & \hfuzz=500pt filename & \hfuzz=500pt \\
\hfuzz=500pt\includegraphics[width=1em]{element.pdf}~headerLines & \hfuzz=500pt uint & \hfuzz=500pt skip number of header lines\\
\hfuzz=500pt\includegraphics[width=1em]{element.pdf}~referenceTime & \hfuzz=500pt time & \hfuzz=500pt reference time\\
\hline
\end{tabularx}

\clearpage
%==================================
\subsection{Icgem2PotentialCoefficients}\label{Icgem2PotentialCoefficients}
Read spherical harmonics in ICGEM format (\url{http://icgem.gfz-potsdam.de/}).


\keepXColumns
\begin{tabularx}{\textwidth}{N T A}
\hline
Name & Type & Annotation\\
\hline
\hfuzz=500pt\includegraphics[width=1em]{element-mustset.pdf}~outputfileStaticCoefficients & \hfuzz=500pt filename & \hfuzz=500pt static potential coefficients in GROOPS gfc format. Available variables (icgem2.0): epochStart, epochEnd, epochMid; (icgem1.0) epochReference\\
\hfuzz=500pt\includegraphics[width=1em]{element.pdf}~outputfileTrendCoefficients & \hfuzz=500pt filename & \hfuzz=500pt trend potential coefficients in GROOPS gfc format.  Available variables (icgem2.0): epochStart, epochEnd, epochMid; (icgem1.0) epochReference\\
\hfuzz=500pt\includegraphics[width=1em]{element.pdf}~outputfileOscillationCosine & \hfuzz=500pt filename & \hfuzz=500pt oscillation cosine coefficients in GROOPS gfc format. Available variables (icgem2.0): epochStart, epochEnd, epochMid, oscillationPeriod; (icgem1.0) epochReference, oscillationPeriod\\
\hfuzz=500pt\includegraphics[width=1em]{element.pdf}~outputfileOscillationSine & \hfuzz=500pt filename & \hfuzz=500pt oscillation sine coefficients in GROOPS gfc format. Available variables (icgem2.0): epochStart, epochEnd, epochMid, oscillationPeriod; (icgem1.0) epochReference, oscillationPeriod\\
\hfuzz=500pt\includegraphics[width=1em]{element.pdf}~outputfileIntervals & \hfuzz=500pt filename & \hfuzz=500pt two column ASCII file with all intervals found (only sensible for icgem2.0). The base name will be extended with .static, .trend, .annualCos, and .annualSin.\\
\hfuzz=500pt\includegraphics[width=1em]{element-mustset.pdf}~inputfileIcgem & \hfuzz=500pt filename & \hfuzz=500pt ICGEM GFC file\\
\hfuzz=500pt\includegraphics[width=1em]{element.pdf}~useFormalErrors & \hfuzz=500pt boolean & \hfuzz=500pt use formal errors if both formal and calibrated errors are given\\
\hline
\end{tabularx}

\clearpage
%==================================
\subsection{Iers2OceanPoleTide}\label{Iers2OceanPoleTide}
Read ocean pole tide model according to IERS conventions
and convert into \file{oceanPoleTide file}{oceanPoleTide}.


\keepXColumns
\begin{tabularx}{\textwidth}{N T A}
\hline
Name & Type & Annotation\\
\hline
\hfuzz=500pt\includegraphics[width=1em]{element-mustset.pdf}~outputfileOceanPole & \hfuzz=500pt filename & \hfuzz=500pt \\
\hfuzz=500pt\includegraphics[width=1em]{element-mustset.pdf}~inputfile & \hfuzz=500pt filename & \hfuzz=500pt \\
\hfuzz=500pt\includegraphics[width=1em]{element-mustset.pdf}~inputfileLoadingLoveNumber & \hfuzz=500pt filename & \hfuzz=500pt \\
\hfuzz=500pt\includegraphics[width=1em]{element-mustset.pdf}~maxDegree & \hfuzz=500pt uint & \hfuzz=500pt \\
\hfuzz=500pt\includegraphics[width=1em]{element.pdf}~GM & \hfuzz=500pt double & \hfuzz=500pt Geocentric gravitational constant\\
\hfuzz=500pt\includegraphics[width=1em]{element.pdf}~R & \hfuzz=500pt double & \hfuzz=500pt Reference radius\\
\hfuzz=500pt\includegraphics[width=1em]{element.pdf}~Omega & \hfuzz=500pt double & \hfuzz=500pt [rad/s] earth rotation\\
\hfuzz=500pt\includegraphics[width=1em]{element.pdf}~rho & \hfuzz=500pt double & \hfuzz=500pt [kg/m**3] density of sea water\\
\hfuzz=500pt\includegraphics[width=1em]{element.pdf}~G & \hfuzz=500pt double & \hfuzz=500pt [m**3/(kg*s**2)] gravitational constant\\
\hfuzz=500pt\includegraphics[width=1em]{element.pdf}~g & \hfuzz=500pt double & \hfuzz=500pt [m/s**2] gravity\\
\hline
\end{tabularx}

\clearpage
%==================================
\subsection{IersC04IAU2000EarthOrientationParameter}\label{IersC04IAU2000EarthOrientationParameter}
Read a IERS Earth orientation data C04 (IAU2000A) file
and write it as \configFile{outputfileEOP}{earthOrientationParameter}.


\keepXColumns
\begin{tabularx}{\textwidth}{N T A}
\hline
Name & Type & Annotation\\
\hline
\hfuzz=500pt\includegraphics[width=1em]{element-mustset.pdf}~outputfileEOP & \hfuzz=500pt filename & \hfuzz=500pt \\
\hfuzz=500pt\includegraphics[width=1em]{element-mustset.pdf}~inputfile & \hfuzz=500pt filename & \hfuzz=500pt \\
\hfuzz=500pt\includegraphics[width=1em]{element.pdf}~timeStart & \hfuzz=500pt time & \hfuzz=500pt \\
\hfuzz=500pt\includegraphics[width=1em]{element.pdf}~timeEnd & \hfuzz=500pt time & \hfuzz=500pt \\
\hline
\end{tabularx}

\clearpage
%==================================
\subsection{IersHighFrequentEop2DoodsonEop}\label{IersHighFrequentEop2DoodsonEop}
Read Diurnal and Subdiurnal Earth Orientation variations according to updated IERS 2010 conventions
and write them as \configFile{outputfileDoodsonEOP}{doodsonEarthOrientationParameter}.


\keepXColumns
\begin{tabularx}{\textwidth}{N T A}
\hline
Name & Type & Annotation\\
\hline
\hfuzz=500pt\includegraphics[width=1em]{element-mustset.pdf}~outputfileDoodsonEOP & \hfuzz=500pt filename & \hfuzz=500pt \\
\hfuzz=500pt\includegraphics[width=1em]{element-mustset.pdf}~inputfile & \hfuzz=500pt filename & \hfuzz=500pt \\
\hline
\end{tabularx}

\clearpage
%==================================
\subsection{IersPotential2DoodsonHarmonics}\label{IersPotential2DoodsonHarmonics}
Read ocean tide file in IERS format.


\keepXColumns
\begin{tabularx}{\textwidth}{N T A}
\hline
Name & Type & Annotation\\
\hline
\hfuzz=500pt\includegraphics[width=1em]{element-mustset.pdf}~outputfileDoodsonHarmoncis & \hfuzz=500pt filename & \hfuzz=500pt \\
\hfuzz=500pt\includegraphics[width=1em]{element-mustset.pdf}~inputfile & \hfuzz=500pt filename & \hfuzz=500pt \\
\hfuzz=500pt\includegraphics[width=1em]{element-mustset.pdf}~headerLines & \hfuzz=500pt uint & \hfuzz=500pt skip number of header lines\\
\hfuzz=500pt\includegraphics[width=1em]{element.pdf}~minDegree & \hfuzz=500pt uint & \hfuzz=500pt \\
\hfuzz=500pt\includegraphics[width=1em]{element-mustset.pdf}~maxDegree & \hfuzz=500pt uint & \hfuzz=500pt \\
\hfuzz=500pt\includegraphics[width=1em]{element.pdf}~GM & \hfuzz=500pt double & \hfuzz=500pt Geocentric gravitational constant\\
\hfuzz=500pt\includegraphics[width=1em]{element.pdf}~R & \hfuzz=500pt double & \hfuzz=500pt reference radius\\
\hline
\end{tabularx}

\clearpage
%==================================
\subsection{IersRapidIAU2000EarthOrientationParameter}\label{IersRapidIAU2000EarthOrientationParameter}
Read a IERS Earth orientation rapid data and prediction file (IAU2000)
and write it as \configFile{outputfileEOP}{earthOrientationParameter}.


\keepXColumns
\begin{tabularx}{\textwidth}{N T A}
\hline
Name & Type & Annotation\\
\hline
\hfuzz=500pt\includegraphics[width=1em]{element-mustset.pdf}~outputfileEOP & \hfuzz=500pt filename & \hfuzz=500pt \\
\hfuzz=500pt\includegraphics[width=1em]{element-mustset.pdf}~inputfile & \hfuzz=500pt filename & \hfuzz=500pt \\
\hfuzz=500pt\includegraphics[width=1em]{element.pdf}~timeStart & \hfuzz=500pt time & \hfuzz=500pt \\
\hfuzz=500pt\includegraphics[width=1em]{element.pdf}~timeEnd & \hfuzz=500pt time & \hfuzz=500pt \\
\hline
\end{tabularx}

\clearpage
%==================================
\subsection{IersWaterHeight2DoodsonHarmonics}\label{IersWaterHeight2DoodsonHarmonics}
Read ocean tide file in IERS format.


\keepXColumns
\begin{tabularx}{\textwidth}{N T A}
\hline
Name & Type & Annotation\\
\hline
\hfuzz=500pt\includegraphics[width=1em]{element-mustset.pdf}~outputfileDoodsonHarmoncis & \hfuzz=500pt filename & \hfuzz=500pt \\
\hfuzz=500pt\includegraphics[width=1em]{element-mustset.pdf}~inputfile & \hfuzz=500pt filename & \hfuzz=500pt \\
\hfuzz=500pt\includegraphics[width=1em]{element-mustset.pdf}~headerLines & \hfuzz=500pt uint & \hfuzz=500pt skip number of header lines\\
\hfuzz=500pt\includegraphics[width=1em]{element-mustset.pdf}~inputfileTideGeneratingPotential & \hfuzz=500pt filename & \hfuzz=500pt to compute Xi phase correction\\
\hfuzz=500pt\includegraphics[width=1em]{element-mustset.pdf}~kernel & \hfuzz=500pt \hyperref[kernelType]{kernel} & \hfuzz=500pt data type of input values\\
\hfuzz=500pt\includegraphics[width=1em]{element.pdf}~factor & \hfuzz=500pt double & \hfuzz=500pt to convert in SI units\\
\hfuzz=500pt\includegraphics[width=1em]{element.pdf}~minDegree & \hfuzz=500pt uint & \hfuzz=500pt \\
\hfuzz=500pt\includegraphics[width=1em]{element-mustset.pdf}~maxDegree & \hfuzz=500pt uint & \hfuzz=500pt \\
\hfuzz=500pt\includegraphics[width=1em]{element.pdf}~GM & \hfuzz=500pt double & \hfuzz=500pt Geocentric gravitational constant\\
\hfuzz=500pt\includegraphics[width=1em]{element.pdf}~R & \hfuzz=500pt double & \hfuzz=500pt reference radius\\
\hline
\end{tabularx}

\clearpage
%==================================
\subsection{Igs2EarthOrientationParameter}\label{Igs2EarthOrientationParameter}
Read Rapid Earth Orientation Parameter from IGS daily file
and write it as \configFile{outputfileEOP}{earthOrientationParameter}.


\keepXColumns
\begin{tabularx}{\textwidth}{N T A}
\hline
Name & Type & Annotation\\
\hline
\hfuzz=500pt\includegraphics[width=1em]{element-mustset.pdf}~outputfileEOP & \hfuzz=500pt filename & \hfuzz=500pt \\
\hfuzz=500pt\includegraphics[width=1em]{element-mustset-unbounded.pdf}~inputfile & \hfuzz=500pt filename & \hfuzz=500pt \\
\hfuzz=500pt\includegraphics[width=1em]{element.pdf}~timeStart & \hfuzz=500pt time & \hfuzz=500pt \\
\hfuzz=500pt\includegraphics[width=1em]{element.pdf}~timeEnd & \hfuzz=500pt time & \hfuzz=500pt \\
\hline
\end{tabularx}

\clearpage
%==================================
\subsection{Jason2Starcamera}\label{Jason2Starcamera}
This program reads in Jason star camera data given in a special  format.
Files available at: \url{cddis.gsfc.nasa.gov/pub/doris/ancillary/quaternions/ja2/}.
A description of the format can be found under:
\url{ftp://ftp.ids-doris.org/pub/ids/ancillary/quaternions/jason1_2_quaternion_solar_panel.pdf}


\keepXColumns
\begin{tabularx}{\textwidth}{N T A}
\hline
Name & Type & Annotation\\
\hline
\hfuzz=500pt\includegraphics[width=1em]{element-mustset.pdf}~outputfileStarCamera & \hfuzz=500pt filename & \hfuzz=500pt \\
\hfuzz=500pt\includegraphics[width=1em]{element-mustset.pdf}~jasonNumber & \hfuzz=500pt uint & \hfuzz=500pt Jason number (different file format), 1 for Sentinel\\
\hfuzz=500pt\includegraphics[width=1em]{element-mustset-unbounded.pdf}~inputfile & \hfuzz=500pt filename & \hfuzz=500pt \\
\hline
\end{tabularx}

\clearpage
%==================================
\subsection{JplAscii2Ephemerides}\label{JplAscii2Ephemerides}
Read JPL DExxx (ASCII) ephemerides.


\keepXColumns
\begin{tabularx}{\textwidth}{N T A}
\hline
Name & Type & Annotation\\
\hline
\hfuzz=500pt\includegraphics[width=1em]{element-mustset.pdf}~outputfileEphemerides & \hfuzz=500pt filename & \hfuzz=500pt \\
\hfuzz=500pt\includegraphics[width=1em]{element-mustset.pdf}~inputfileHeader & \hfuzz=500pt filename & \hfuzz=500pt \\
\hfuzz=500pt\includegraphics[width=1em]{element-mustset-unbounded.pdf}~inputfileData & \hfuzz=500pt filename & \hfuzz=500pt \\
\hline
\end{tabularx}

\clearpage
%==================================
\subsection{Merit2FullRate}\label{Merit2FullRate}
Converts \href{https://ilrs.gsfc.nasa.gov/data_and_products/data/npt/npt_format.html}{MERIT II file}
and writes an \file{instrument file (METEOROLOGICAL)}{instrument} including meteorological data like
temperature, air pressure and humidity as well as an \file{instrument file (SATELLITELASERRANGING)}{instrument}
including full rate data like range, accuracy, wavelength, azimuth and elevation.


\keepXColumns
\begin{tabularx}{\textwidth}{N T A}
\hline
Name & Type & Annotation\\
\hline
\hfuzz=500pt\includegraphics[width=1em]{element-mustset.pdf}~outputfileNormalPoints & \hfuzz=500pt filename & \hfuzz=500pt variable \{station\} available\\
\hfuzz=500pt\includegraphics[width=1em]{element-mustset.pdf}~outputfileMeteorological & \hfuzz=500pt filename & \hfuzz=500pt variable \{station\} available\\
\hfuzz=500pt\includegraphics[width=1em]{element-mustset.pdf}~inputfileSlrData & \hfuzz=500pt filename & \hfuzz=500pt SLR MERIT II file\\
\hline
\end{tabularx}

\clearpage
%==================================
\subsection{Merit2NormalPoints}\label{Merit2NormalPoints}
Converts \href{https://ilrs.gsfc.nasa.gov/data_and_products/data/npt/npt_format.html}{MERIT II file}
and writes an \file{instrument file (METEOROLOGICAL)}{instrument} including meteorological data like
temperature, air pressure and humidity as well as an \file{instrument file (SATELLITELASERRANGING)}{instrument}
including normal point data like range, accuracy, redundancy, wavelength, window size, azimuth and elevation.


\keepXColumns
\begin{tabularx}{\textwidth}{N T A}
\hline
Name & Type & Annotation\\
\hline
\hfuzz=500pt\includegraphics[width=1em]{element-mustset.pdf}~outputfileNormalPoints & \hfuzz=500pt filename & \hfuzz=500pt variable \{station\} available\\
\hfuzz=500pt\includegraphics[width=1em]{element-mustset.pdf}~outputfileMeteorological & \hfuzz=500pt filename & \hfuzz=500pt variable \{station\} available\\
\hfuzz=500pt\includegraphics[width=1em]{element-mustset.pdf}~inputfileSlrData & \hfuzz=500pt filename & \hfuzz=500pt SLR MERIT II file\\
\hline
\end{tabularx}

\clearpage
%==================================
\subsection{NetCdf2GriddedData}\label{NetCdf2GriddedData}
This program converts a COARDS compliant NetCDF file into an
\configFile{outputfileGriddedData}{griddedData}.
If no specific input \config{variableNameData} are selected all suitable data are used.

If the NETCDF file contains a time axis (\config{variableNameData}) an specific epoch
can be selected with \config{time}. The nearest epoch in file is used.

See also \program{NetCdfInfo}, \program{GriddedData2NetCdf}, \program{NetCdf2GriddedDataTimeSeries}.


\keepXColumns
\begin{tabularx}{\textwidth}{N T A}
\hline
Name & Type & Annotation\\
\hline
\hfuzz=500pt\includegraphics[width=1em]{element-mustset.pdf}~outputfileGriddedData & \hfuzz=500pt filename & \hfuzz=500pt \\
\hfuzz=500pt\includegraphics[width=1em]{element-mustset.pdf}~inputfileNetCdf & \hfuzz=500pt filename & \hfuzz=500pt \\
\hfuzz=500pt\includegraphics[width=1em]{element-mustset.pdf}~variableNameLongitude & \hfuzz=500pt string & \hfuzz=500pt name of NetCDF variable\\
\hfuzz=500pt\includegraphics[width=1em]{element-mustset.pdf}~variableNameLatitude & \hfuzz=500pt string & \hfuzz=500pt name of NetCDF variable\\
\hfuzz=500pt\includegraphics[width=1em]{element.pdf}~variableNameTime & \hfuzz=500pt string & \hfuzz=500pt if with time axis: name of NetCDF variable\\
\hfuzz=500pt\includegraphics[width=1em]{element-unbounded.pdf}~variableNameData & \hfuzz=500pt string & \hfuzz=500pt data variables, otherwise all suitable data are used\\
\hfuzz=500pt\includegraphics[width=1em]{element.pdf}~time & \hfuzz=500pt time & \hfuzz=500pt if with time axis: nearest epoch is used\\
\hfuzz=500pt\includegraphics[width=1em]{element.pdf}~R & \hfuzz=500pt double & \hfuzz=500pt reference radius for ellipsoidal coordinates\\
\hfuzz=500pt\includegraphics[width=1em]{element.pdf}~inverseFlattening & \hfuzz=500pt double & \hfuzz=500pt reference flattening for ellipsoidal coordinates\\
\hline
\end{tabularx}

\clearpage
%==================================
\subsection{NetCdf2GriddedDataTimeSeries}\label{NetCdf2GriddedDataTimeSeries}
This program converts a COARDS compliant NetCDF file into
\configFile{outputfileGriddedDataTimeSeries}{griddedDataTimeSeries}.
If no specific input \config{variableNameData} are selected all suitable data are used.

See also \program{NetCdfInfo}, \program{NetCdf2GriddedData}, \program{GriddedDataTimeSeries2NetCdf}.


\keepXColumns
\begin{tabularx}{\textwidth}{N T A}
\hline
Name & Type & Annotation\\
\hline
\hfuzz=500pt\includegraphics[width=1em]{element-mustset.pdf}~outputfileGriddedDataTimeSeries & \hfuzz=500pt filename & \hfuzz=500pt \\
\hfuzz=500pt\includegraphics[width=1em]{element-mustset.pdf}~inputfileNetCdf & \hfuzz=500pt filename & \hfuzz=500pt \\
\hfuzz=500pt\includegraphics[width=1em]{element-mustset.pdf}~variableNameLongitude & \hfuzz=500pt string & \hfuzz=500pt name of NetCDF variable\\
\hfuzz=500pt\includegraphics[width=1em]{element-mustset.pdf}~variableNameLatitude & \hfuzz=500pt string & \hfuzz=500pt name of NetCDF variable\\
\hfuzz=500pt\includegraphics[width=1em]{element-mustset.pdf}~variableNameTime & \hfuzz=500pt string & \hfuzz=500pt name of NetCDF variable)\\
\hfuzz=500pt\includegraphics[width=1em]{element-unbounded.pdf}~variableNameData & \hfuzz=500pt string & \hfuzz=500pt data variables, otherwise all suitable data are used\\
\hfuzz=500pt\includegraphics[width=1em]{element.pdf}~R & \hfuzz=500pt double & \hfuzz=500pt reference radius for ellipsoidal coordinates\\
\hfuzz=500pt\includegraphics[width=1em]{element.pdf}~inverseFlattening & \hfuzz=500pt double & \hfuzz=500pt reference flattening for ellipsoidal coordinates\\
\hline
\end{tabularx}

\clearpage
%==================================
\subsection{NetCdfInfo}\label{NetCdfInfo}
Print content information of a NetCDF file like
dimensions, variables and attributes.

See also \program{NetCdf2GriddedData}, \program{NetCdf2GriddedDataTimeSeries},
\program{GriddedData2NetCdf}, \program{GriddedDataTimeSeries2NetCdf}.


\keepXColumns
\begin{tabularx}{\textwidth}{N T A}
\hline
Name & Type & Annotation\\
\hline
\hfuzz=500pt\includegraphics[width=1em]{element-mustset.pdf}~inputfileNetCdf & \hfuzz=500pt filename & \hfuzz=500pt \\
\hline
\end{tabularx}

\clearpage
%==================================
\subsection{NormalsSphericalHarmonics2Sinex}\label{NormalsSphericalHarmonics2Sinex}
Write potential coefficients and \file{normal equations}{normalEquation} to
\href{http://www.iers.org/IERS/EN/Organization/AnalysisCoordinator/SinexFormat/sinex.html}{SINEX format}.

See also \program{Sinex2Normals} and \program{GnssNormals2Sinex}.


\keepXColumns
\begin{tabularx}{\textwidth}{N T A}
\hline
Name & Type & Annotation\\
\hline
\hfuzz=500pt\includegraphics[width=1em]{element-mustset.pdf}~outputfileSinex & \hfuzz=500pt filename & \hfuzz=500pt solutions in SINEX format\\
\hfuzz=500pt\includegraphics[width=1em]{element-mustset.pdf}~inputfileNormals & \hfuzz=500pt filename & \hfuzz=500pt normal equation matrix\\
\hfuzz=500pt\includegraphics[width=1em]{element.pdf}~inputfileSolution & \hfuzz=500pt filename & \hfuzz=500pt parameter vector\\
\hfuzz=500pt\includegraphics[width=1em]{element.pdf}~inputfileSigmax & \hfuzz=500pt filename & \hfuzz=500pt standard deviations of the parameters (sqrt of the diagonal of the inverse normal equation)\\
\hfuzz=500pt\includegraphics[width=1em]{element-mustset.pdf}~inputfileApriori & \hfuzz=500pt filename & \hfuzz=500pt apriori parameter vector\\
\hfuzz=500pt\includegraphics[width=1em]{element.pdf}~inputfileAprioriMatrix & \hfuzz=500pt filename & \hfuzz=500pt normal equation matrix of applied constraints\\
\hfuzz=500pt\includegraphics[width=1em]{element-mustset.pdf}~time & \hfuzz=500pt time & \hfuzz=500pt reference time for parameters\\
\hfuzz=500pt\includegraphics[width=1em]{element-mustset.pdf}~sinexHeader & \hfuzz=500pt sequence & \hfuzz=500pt \\
\hfuzz=500pt\includegraphics[width=1em]{connector.pdf}\includegraphics[width=1em]{element.pdf}~agencyCode & \hfuzz=500pt string & \hfuzz=500pt identify the agency providing the data\\
\hfuzz=500pt\includegraphics[width=1em]{connector.pdf}\includegraphics[width=1em]{element.pdf}~timeStart & \hfuzz=500pt time & \hfuzz=500pt start time of the data\\
\hfuzz=500pt\includegraphics[width=1em]{connector.pdf}\includegraphics[width=1em]{element.pdf}~timeEnd & \hfuzz=500pt time & \hfuzz=500pt end time of the data \\
\hfuzz=500pt\includegraphics[width=1em]{connector.pdf}\includegraphics[width=1em]{element.pdf}~observationCode & \hfuzz=500pt string & \hfuzz=500pt technique used to generate the SINEX solution\\
\hfuzz=500pt\includegraphics[width=1em]{connector.pdf}\includegraphics[width=1em]{element.pdf}~constraintCode & \hfuzz=500pt string & \hfuzz=500pt 0: tight constraint, 1: siginficant constraint, 2: unconstrained\\
\hfuzz=500pt\includegraphics[width=1em]{connector.pdf}\includegraphics[width=1em]{element.pdf}~solutionContent & \hfuzz=500pt string & \hfuzz=500pt solution types contained in the SINEX solution (S O E T C A)\\
\hfuzz=500pt\includegraphics[width=1em]{connector.pdf}\includegraphics[width=1em]{element.pdf}~description & \hfuzz=500pt string & \hfuzz=500pt organizitions gathering/alerting the file contents\\
\hfuzz=500pt\includegraphics[width=1em]{connector.pdf}\includegraphics[width=1em]{element.pdf}~contact & \hfuzz=500pt string & \hfuzz=500pt Address of the relevant contact. e-mail\\
\hfuzz=500pt\includegraphics[width=1em]{connector.pdf}\includegraphics[width=1em]{element.pdf}~output & \hfuzz=500pt string & \hfuzz=500pt Description of the file contents\\
\hfuzz=500pt\includegraphics[width=1em]{connector.pdf}\includegraphics[width=1em]{element.pdf}~input & \hfuzz=500pt string & \hfuzz=500pt Brief description of the input used to generate this solution\\
\hfuzz=500pt\includegraphics[width=1em]{connector.pdf}\includegraphics[width=1em]{element.pdf}~software & \hfuzz=500pt string & \hfuzz=500pt Software used to generate the file\\
\hfuzz=500pt\includegraphics[width=1em]{connector.pdf}\includegraphics[width=1em]{element.pdf}~hardware & \hfuzz=500pt string & \hfuzz=500pt Computer hardware on which above software was run\\
\hfuzz=500pt\includegraphics[width=1em]{connector.pdf}\includegraphics[width=1em]{element.pdf}~inputfileComment & \hfuzz=500pt filename & \hfuzz=500pt comments in the comment block from a file (truncated at 80 characters)\\
\hfuzz=500pt\includegraphics[width=1em]{connector.pdf}\includegraphics[width=1em]{element-unbounded.pdf}~comment & \hfuzz=500pt string & \hfuzz=500pt comments in the comment block\\
\hline
\end{tabularx}

\clearpage
%==================================
\subsection{OceanTidesDTU2GriddedData}\label{OceanTidesDTU2GriddedData}
Convert DTU ocean tide grids to griddedData (amplitude, phase).


\keepXColumns
\begin{tabularx}{\textwidth}{N T A}
\hline
Name & Type & Annotation\\
\hline
\hfuzz=500pt\includegraphics[width=1em]{element-mustset.pdf}~outputfileGriddedData & \hfuzz=500pt filename & \hfuzz=500pt data0=amplitude, data1=phase\\
\hfuzz=500pt\includegraphics[width=1em]{element-mustset.pdf}~inputfile & \hfuzz=500pt filename & \hfuzz=500pt \\
\hfuzz=500pt\includegraphics[width=1em]{element.pdf}~R & \hfuzz=500pt double & \hfuzz=500pt reference radius for ellipsoidal coordinates\\
\hfuzz=500pt\includegraphics[width=1em]{element.pdf}~inverseFlattening & \hfuzz=500pt double & \hfuzz=500pt reference flattening for ellipsoidal coordinates\\
\hline
\end{tabularx}

\clearpage
%==================================
\subsection{Orbit2Cpf}\label{Orbit2Cpf}
Writes groops orbits to \href{https://ilrs.gsfc.nasa.gov/data_and_products/formats/cpf.html}{CPF file}.

The coordinate system used in the CPF format is usually presented in ITRF.
The required time format for the input orbit file is GPS.
The time format of the output CPF file is given in UTC.

See also \program{Cpf2Orbit}.


\keepXColumns
\begin{tabularx}{\textwidth}{N T A}
\hline
Name & Type & Annotation\\
\hline
\hfuzz=500pt\includegraphics[width=1em]{element-mustset.pdf}~outputfile & \hfuzz=500pt filename & \hfuzz=500pt \\
\hfuzz=500pt\includegraphics[width=1em]{element-mustset.pdf}~inputfileOrbit & \hfuzz=500pt filename & \hfuzz=500pt \\
\hfuzz=500pt\includegraphics[width=1em]{element-mustset.pdf}~inputfileSatelliteInfo & \hfuzz=500pt filename & \hfuzz=500pt Platform File\\
\hfuzz=500pt\includegraphics[width=1em]{element-mustset.pdf}~earthRotation & \hfuzz=500pt \hyperref[earthRotationType]{earthRotation} & \hfuzz=500pt \\
\hfuzz=500pt\includegraphics[width=1em]{element.pdf}~versionNumber & \hfuzz=500pt uint & \hfuzz=500pt Version number of production day with zero leading fill, e.g. 01\\
\hfuzz=500pt\includegraphics[width=1em]{element.pdf}~targetClass & \hfuzz=500pt uint & \hfuzz=500pt set 1 for passive retroreflector, set 0 for no retroreflector (includes debris)\\
\hline
\end{tabularx}

\clearpage
%==================================
\subsection{Orbit2GroopsAscii}\label{Orbit2GroopsAscii}
Convert groops orbits and corresponding covariance information to ASCII format.
The format is used to publish TUG orbits. It contains a two line header
with a short description of the orbit defined in \config{firstLine}.
The orbit is rotated to the Earth fixed frame (TRF) with \configClass{earthRotation}{earthRotationType} and given as one line per epoch.
The epoch lines contained time [MJD GPS time], position x, y and z [m], and the epoch covariance xx, yy, zz, xy, xz and yz [$m^2$].

See also \program{GroopsAscii2Orbit}.


\keepXColumns
\begin{tabularx}{\textwidth}{N T A}
\hline
Name & Type & Annotation\\
\hline
\hfuzz=500pt\includegraphics[width=1em]{element-mustset.pdf}~outputfile & \hfuzz=500pt filename & \hfuzz=500pt \\
\hfuzz=500pt\includegraphics[width=1em]{element-mustset.pdf}~inputfileOrbit & \hfuzz=500pt filename & \hfuzz=500pt \\
\hfuzz=500pt\includegraphics[width=1em]{element-mustset.pdf}~inputfileCovariance & \hfuzz=500pt filename & \hfuzz=500pt \\
\hfuzz=500pt\includegraphics[width=1em]{element-mustset.pdf}~earthRotation & \hfuzz=500pt \hyperref[earthRotationType]{earthRotation} & \hfuzz=500pt \\
\hfuzz=500pt\includegraphics[width=1em]{element.pdf}~firstLine & \hfuzz=500pt string & \hfuzz=500pt Text for first line\\
\hline
\end{tabularx}

\clearpage
%==================================
\subsection{Orbit2Sp3Format}\label{Orbit2Sp3Format}
Writes orbits to \href{https://files.igs.org/pub/data/format/sp3d.pdf}{SP3 format}.

SP3 orbits are usually given in the terrestrial reference frame (TRF), so providing \configClass{earthRotation}{earthRotationType}
automatically rotates the orbits from the celestial reference frame (CRF) to the TRF.
Since SP3 orbits often use the center of Earth as a reference, a correction from center of mass to center
of Earth can be applied to the orbits by providing \configClass{gravityfield}{gravityfieldType} (e.g. ocean tides).

See also \program{Sp3Format2Orbit}.


\keepXColumns
\begin{tabularx}{\textwidth}{N T A}
\hline
Name & Type & Annotation\\
\hline
\hfuzz=500pt\includegraphics[width=1em]{element-mustset.pdf}~outputfile & \hfuzz=500pt filename & \hfuzz=500pt \\
\hfuzz=500pt\includegraphics[width=1em]{element-mustset-unbounded.pdf}~satellite & \hfuzz=500pt sequence & \hfuzz=500pt \\
\hfuzz=500pt\includegraphics[width=1em]{connector.pdf}\includegraphics[width=1em]{element-mustset.pdf}~inputfileOrbit & \hfuzz=500pt filename & \hfuzz=500pt \\
\hfuzz=500pt\includegraphics[width=1em]{connector.pdf}\includegraphics[width=1em]{element.pdf}~inputfileClock & \hfuzz=500pt filename & \hfuzz=500pt \\
\hfuzz=500pt\includegraphics[width=1em]{connector.pdf}\includegraphics[width=1em]{element.pdf}~inputfileCovariance & \hfuzz=500pt filename & \hfuzz=500pt \\
\hfuzz=500pt\includegraphics[width=1em]{connector.pdf}\includegraphics[width=1em]{element-mustset.pdf}~identifier & \hfuzz=500pt string & \hfuzz=500pt 3 characters (e.g. GNSS PRN: G01)\\
\hfuzz=500pt\includegraphics[width=1em]{connector.pdf}\includegraphics[width=1em]{element.pdf}~orbitAccuracy & \hfuzz=500pt double & \hfuzz=500pt [m] used for accuracy codes in header (0 = unknown)\\
\hfuzz=500pt\includegraphics[width=1em]{element.pdf}~earthRotation & \hfuzz=500pt \hyperref[earthRotationType]{earthRotation} & \hfuzz=500pt rotate data into Earth-fixed frame\\
\hfuzz=500pt\includegraphics[width=1em]{element-unbounded.pdf}~gravityfield & \hfuzz=500pt \hyperref[gravityfieldType]{gravityfield} & \hfuzz=500pt degree 1 fluid mantle for CM2CE correction (SP3 orbits should be in center of Earth)\\
\hfuzz=500pt\includegraphics[width=1em]{element-unbounded.pdf}~comment & \hfuzz=500pt string & \hfuzz=500pt comment lines (77 char max)\\
\hfuzz=500pt\includegraphics[width=1em]{element.pdf}~firstLine & \hfuzz=500pt string & \hfuzz=500pt Text for first line e.g:  u+U  IGb14 KIN ITSG\\
\hfuzz=500pt\includegraphics[width=1em]{element.pdf}~writeVelocity & \hfuzz=500pt boolean & \hfuzz=500pt write velocity in addition to position\\
\hfuzz=500pt\includegraphics[width=1em]{element.pdf}~useSp3kFormat & \hfuzz=500pt boolean & \hfuzz=500pt use the extended sp3k format\\
\hline
\end{tabularx}

\clearpage
%==================================
\subsection{PotentialCoefficients2Icgem}\label{PotentialCoefficients2Icgem}
Write spherical harmonics in ICGEM format.
GROOPS uses this format as default but this program enables
the possibility to include comments and set the modelname.


\keepXColumns
\begin{tabularx}{\textwidth}{N T A}
\hline
Name & Type & Annotation\\
\hline
\hfuzz=500pt\includegraphics[width=1em]{element-mustset.pdf}~outputfile & \hfuzz=500pt filename & \hfuzz=500pt \\
\hfuzz=500pt\includegraphics[width=1em]{element-mustset.pdf}~inputfilePotentialCoefficients & \hfuzz=500pt filename & \hfuzz=500pt \\
\hfuzz=500pt\includegraphics[width=1em]{element.pdf}~inputfileTrend & \hfuzz=500pt filename & \hfuzz=500pt \\
\hfuzz=500pt\includegraphics[width=1em]{element-unbounded.pdf}~oscillation & \hfuzz=500pt sequence & \hfuzz=500pt \\
\hfuzz=500pt\includegraphics[width=1em]{connector.pdf}\includegraphics[width=1em]{element-mustset.pdf}~inputfileCosPotentialCoefficients & \hfuzz=500pt filename & \hfuzz=500pt \\
\hfuzz=500pt\includegraphics[width=1em]{connector.pdf}\includegraphics[width=1em]{element-mustset.pdf}~inputfileSinPotentialCoefficients & \hfuzz=500pt filename & \hfuzz=500pt \\
\hfuzz=500pt\includegraphics[width=1em]{connector.pdf}\includegraphics[width=1em]{element-mustset.pdf}~period & \hfuzz=500pt string & \hfuzz=500pt period of oscillation [year]\\
\hfuzz=500pt\includegraphics[width=1em]{element-unbounded.pdf}~comment & \hfuzz=500pt string & \hfuzz=500pt comment in header\\
\hfuzz=500pt\includegraphics[width=1em]{element.pdf}~inputfileComment & \hfuzz=500pt filename & \hfuzz=500pt file containing comments for header\\
\hfuzz=500pt\includegraphics[width=1em]{element-mustset.pdf}~modelname & \hfuzz=500pt string & \hfuzz=500pt name of the model\\
\hfuzz=500pt\includegraphics[width=1em]{element.pdf}~tideSystem & \hfuzz=500pt choice & \hfuzz=500pt tide system of model\\
\hfuzz=500pt\includegraphics[width=1em]{connector.pdf}\includegraphics[width=1em]{element-mustset.pdf}~zero\_tide & \hfuzz=500pt  & \hfuzz=500pt \\
\hfuzz=500pt\includegraphics[width=1em]{connector.pdf}\includegraphics[width=1em]{element-mustset.pdf}~tide\_free & \hfuzz=500pt  & \hfuzz=500pt \\
\hfuzz=500pt\includegraphics[width=1em]{element.pdf}~minDegree & \hfuzz=500pt uint & \hfuzz=500pt \\
\hfuzz=500pt\includegraphics[width=1em]{element.pdf}~maxDegree & \hfuzz=500pt uint & \hfuzz=500pt \\
\hfuzz=500pt\includegraphics[width=1em]{element.pdf}~GM & \hfuzz=500pt double & \hfuzz=500pt Geocentric gravitational constant\\
\hfuzz=500pt\includegraphics[width=1em]{element.pdf}~R & \hfuzz=500pt double & \hfuzz=500pt reference radius\\
\hfuzz=500pt\includegraphics[width=1em]{element.pdf}~time & \hfuzz=500pt time & \hfuzz=500pt reference time\\
\hline
\end{tabularx}

\clearpage
%==================================
\subsection{PsmslOceanBottomPressure2TimeSeries}\label{PsmslOceanBottomPressure2TimeSeries}
This programs reads ocean bottom pressure time series from the Permanent Service for Mean Sea Level (PSMSL).

In addition to the OBP measurements, the recorder position can be written to a \file{grid file}{griddedData}.


\keepXColumns
\begin{tabularx}{\textwidth}{N T A}
\hline
Name & Type & Annotation\\
\hline
\hfuzz=500pt\includegraphics[width=1em]{element-mustset.pdf}~outputfileTimeSeries & \hfuzz=500pt filename & \hfuzz=500pt \\
\hfuzz=500pt\includegraphics[width=1em]{element.pdf}~outputfilePosition & \hfuzz=500pt filename & \hfuzz=500pt recorder position as gridded data\\
\hfuzz=500pt\includegraphics[width=1em]{element-mustset.pdf}~inputfile & \hfuzz=500pt filename & \hfuzz=500pt \\
\hfuzz=500pt\includegraphics[width=1em]{element-mustset.pdf}~isDaily & \hfuzz=500pt boolean & \hfuzz=500pt \\
\hfuzz=500pt\includegraphics[width=1em]{element-mustset.pdf}~ignoreBadData & \hfuzz=500pt boolean & \hfuzz=500pt \\
\hfuzz=500pt\includegraphics[width=1em]{element.pdf}~R & \hfuzz=500pt double & \hfuzz=500pt \\
\hfuzz=500pt\includegraphics[width=1em]{element.pdf}~inverseFlattening & \hfuzz=500pt double & \hfuzz=500pt \\
\hfuzz=500pt\includegraphics[width=1em]{element-mustset-unbounded.pdf}~timeSeries & \hfuzz=500pt \hyperref[timeSeriesType]{timeSeries} & \hfuzz=500pt \\
\hline
\end{tabularx}

\clearpage
%==================================
\subsection{RinexObservation2GnssReceiver}\label{RinexObservation2GnssReceiver}
Converts \href{https://files.igs.org/pub/data/format/rinex_4.00.pdf}{RINEX} (version 2, 3, and 4) and
\href{https://terras.gsi.go.jp/ja/crx2rnx.html}{Compact RINEX} observation files to
\file{GnssReceiver Instrument file}{instrument}.

In case of \href{https://files.igs.org/pub/data/format/rinex211.txt}{RINEX v2.x} observation files
containing GLONASS satellites, a mapping from PRN
to frequency number must be provided via \config{inputfileMatrixPrn2FrequencyNumber}
in the form of a \file{matrix file}{matrix} with columns: GLONASS PRN, mjdStart, mjdEnd, frequencyNumber.
Source for mapping: \url{http://semisys.gfz-potsdam.de/semisys/api/?symname=2002&format=json&satellite=GLO}.
RINEX v3+ observation files already contain this information.

\configClass{useType}{gnssType} and \configClass{ignoreType}{gnssType} can be used to filter
the observation types that will be exported.

If \configFile{inputfileStationInfo}{platform} is set, RINEX antenna and receiver info
will be cross-checked with the provided file and warnings are raised in case of differences.

A list of semi-codeless GPS receivers (observing C2D instead of C2W) can be provided via
\configFile{inputfileSemiCodelessReceivers}{stringList} in ASCII format with one receiver name per line.
Observation types will be automatically corrected for these receivers.

Some LEO satellites use special RINEX observation types, either from the unofficial RINEX v2.20
or custom ones. These can be provided via \configFile{inputfileSpecialObservationTypes}{stringTable}
in ASCII format. The file must  must contain a table with two columns, the first being the special type,
and the second being the equivalent RINEX v3 type.

%Example for RINEX v2.20:

%\begin{tabular}{ll}
%LA & L1C \\
%L1 & L1W \\
%L2 & L2W \\
%SA & S1C \\
%S1 & S1W \\
%S2 & S2W \\
%\end{tabular}


\keepXColumns
\begin{tabularx}{\textwidth}{N T A}
\hline
Name & Type & Annotation\\
\hline
\hfuzz=500pt\includegraphics[width=1em]{element-mustset.pdf}~outputfileGnssReceiver & \hfuzz=500pt filename & \hfuzz=500pt \\
\hfuzz=500pt\includegraphics[width=1em]{element-mustset-unbounded.pdf}~inputfileRinexObservation & \hfuzz=500pt filename & \hfuzz=500pt RINEX or Compact RINEX observation files\\
\hfuzz=500pt\includegraphics[width=1em]{element.pdf}~inputfileMatrixPrn2FrequencyNumber & \hfuzz=500pt filename & \hfuzz=500pt (required for RINEX v2 files containing GLONASS observations) GROOPS matrix with columns: GLONASS PRN, SVN, mjdStart, mjdEnd, frequencyNumber\\
\hfuzz=500pt\includegraphics[width=1em]{element.pdf}~inputfileStationInfo & \hfuzz=500pt filename & \hfuzz=500pt used to determine semi-codeless receivers and to cross-check antenna and receiver info\\
\hfuzz=500pt\includegraphics[width=1em]{element.pdf}~inputfileSemiCodelessReceivers & \hfuzz=500pt filename & \hfuzz=500pt ASCII list with one receiver name per line\\
\hfuzz=500pt\includegraphics[width=1em]{element.pdf}~inputfileSpecialObservationTypes & \hfuzz=500pt filename & \hfuzz=500pt ASCII table mapping special observation types to RINEX 3 types, e.g.: LA L1C\\
\hfuzz=500pt\includegraphics[width=1em]{element-unbounded.pdf}~useType & \hfuzz=500pt \hyperref[gnssType]{gnssType} & \hfuzz=500pt only use observations that match any of these patterns\\
\hfuzz=500pt\includegraphics[width=1em]{element-unbounded.pdf}~ignoreType & \hfuzz=500pt \hyperref[gnssType]{gnssType} & \hfuzz=500pt ignore observations that match any of these patterns\\
\hline
\end{tabularx}

\clearpage
%==================================
\subsection{Sacc2Orbit}\label{Sacc2Orbit}
This program reads in SACC orbit data.


\keepXColumns
\begin{tabularx}{\textwidth}{N T A}
\hline
Name & Type & Annotation\\
\hline
\hfuzz=500pt\includegraphics[width=1em]{element-mustset.pdf}~outputfileOrbit & \hfuzz=500pt filename & \hfuzz=500pt \\
\hfuzz=500pt\includegraphics[width=1em]{element-mustset-unbounded.pdf}~inputfile & \hfuzz=500pt filename & \hfuzz=500pt \\
\hline
\end{tabularx}

\clearpage
%==================================
\subsection{Sentinel2StarCamera}\label{Sentinel2StarCamera}
This program reads in Sentinel-1/2/3 star camera data given in the special format.


\keepXColumns
\begin{tabularx}{\textwidth}{N T A}
\hline
Name & Type & Annotation\\
\hline
\hfuzz=500pt\includegraphics[width=1em]{element-mustset.pdf}~outputfileStarCamera & \hfuzz=500pt filename & \hfuzz=500pt \\
\hfuzz=500pt\includegraphics[width=1em]{element-mustset-unbounded.pdf}~inputfile & \hfuzz=500pt filename & \hfuzz=500pt \\
\hline
\end{tabularx}

\clearpage
%==================================
\subsection{SentinelXml2Orbit}\label{SentinelXml2Orbit}
Read Sentinel orbits from XML format.


\keepXColumns
\begin{tabularx}{\textwidth}{N T A}
\hline
Name & Type & Annotation\\
\hline
\hfuzz=500pt\includegraphics[width=1em]{element-mustset.pdf}~outputfileOrbit & \hfuzz=500pt filename & \hfuzz=500pt \\
\hfuzz=500pt\includegraphics[width=1em]{element-mustset.pdf}~earthRotation & \hfuzz=500pt \hyperref[earthRotationType]{earthRotation} & \hfuzz=500pt \\
\hfuzz=500pt\includegraphics[width=1em]{element-mustset-unbounded.pdf}~inputfile & \hfuzz=500pt filename & \hfuzz=500pt \\
\hline
\end{tabularx}

\clearpage
%==================================
\subsection{Sinex2Normals}\label{Sinex2Normals}
Convert normal equations from \href{http://www.iers.org/IERS/EN/Organization/AnalysisCoordinator/SinexFormat/sinex.html}{SINEX format}
to \file{normal equations}{normalEquation}.

See also \program{GnssNormals2Sinex} and \program{NormalsSphericalHarmonics2Sinex}.


\keepXColumns
\begin{tabularx}{\textwidth}{N T A}
\hline
Name & Type & Annotation\\
\hline
\hfuzz=500pt\includegraphics[width=1em]{element.pdf}~outputfileNormals & \hfuzz=500pt filename & \hfuzz=500pt N, n: unconstrained normal equations\\
\hfuzz=500pt\includegraphics[width=1em]{element.pdf}~outputfileNormalsConstraint & \hfuzz=500pt filename & \hfuzz=500pt N0, n0: normal equations of applied constraints\\
\hfuzz=500pt\includegraphics[width=1em]{element.pdf}~outputfileSolution & \hfuzz=500pt filename & \hfuzz=500pt x: parameter vector\\
\hfuzz=500pt\includegraphics[width=1em]{element.pdf}~outputfileSolutionApriori & \hfuzz=500pt filename & \hfuzz=500pt x0: a priori parameter vector\\
\hfuzz=500pt\includegraphics[width=1em]{element-mustset.pdf}~inputFileSinex & \hfuzz=500pt filename & \hfuzz=500pt \\
\hline
\end{tabularx}

\clearpage
%==================================
\subsection{Sinex2StationDiscontinuities}\label{Sinex2StationDiscontinuities}
Convert station discontinuities from
\href{http://www.iers.org/IERS/EN/Organization/AnalysisCoordinator/SinexFormat/sinex.html}{SINEX format}
(e.g. ITRF20) to \configFile{outputfileInstrument}{instrument} (MISCVALUE).
A value of 1 means position discontinuity, a value of 2 means velocity discontinuity.
Start and end epochs with value 0 are added in addition to the discontinuities from
SINEX to define continuity interval borders.

See also \program{Sinex2StationPosition} and \program{Sinex2StationPostSeismicDeformation}.


\keepXColumns
\begin{tabularx}{\textwidth}{N T A}
\hline
Name & Type & Annotation\\
\hline
\hfuzz=500pt\includegraphics[width=1em]{element-mustset.pdf}~outputfileInstrument & \hfuzz=500pt filename & \hfuzz=500pt loop variable is replaced with station name (e.g. wtzz)\\
\hfuzz=500pt\includegraphics[width=1em]{element-mustset.pdf}~inputfileDiscontinuities & \hfuzz=500pt filename & \hfuzz=500pt SINEX (e.g. ITRF20) station discontinuities\\
\hfuzz=500pt\includegraphics[width=1em]{element.pdf}~variableLoopStation & \hfuzz=500pt string & \hfuzz=500pt variable name for station loop\\
\hfuzz=500pt\includegraphics[width=1em]{element-unbounded.pdf}~stationName & \hfuzz=500pt string & \hfuzz=500pt only export these stations\\
\hline
\end{tabularx}

\clearpage
%==================================
\subsection{Sinex2StationPositions}\label{Sinex2StationPositions}
Extracts station positions from \config{inputfileSinexSolution}
(\href{http://www.iers.org/IERS/EN/Organization/AnalysisCoordinator/SinexFormat/sinex.html}{SINEX format description})
and writes an \configFile{outputfileInstrument}{instrument} of type VECTOR3D
for each station.  Positions will be computed at \configClass{timeSeries}{timeSeriesType} based on position and velocity
of each provided interval in the SINEX file.
With \config{inputfileSinexDiscontinuities} the bounds of these time spans are adjusted to the exact epochs of discontinuities.
The \config{inputfileSinexPostSeismicDeformations} adds the ITRF post-seismic deformation model to the affected stations.
The \config{inputfileSinexFrequencies} adds annual and semi-annual frequencies.

If \config{extrapolateBackward} or \config{extrapolateForward} are provided, positions will also be computed for epochs
before the first interval/after the last interval, based on the position and velocity of the first/last interval.
Position extrapolation will stop at the first discontinuity before the first interval/after the last interval.

Stations can be limited via \config{stationName}, otherwise all stations in \config{inputfileSinexSolution} will be used.


\keepXColumns
\begin{tabularx}{\textwidth}{N T A}
\hline
Name & Type & Annotation\\
\hline
\hfuzz=500pt\includegraphics[width=1em]{element-mustset.pdf}~outputfileInstrument & \hfuzz=500pt filename & \hfuzz=500pt loop variable is replaced with station name (e.g. wtzz)\\
\hfuzz=500pt\includegraphics[width=1em]{element.pdf}~variableLoopStation & \hfuzz=500pt string & \hfuzz=500pt variable name for station loop\\
\hfuzz=500pt\includegraphics[width=1em]{element-mustset.pdf}~inputfileSinexSolution & \hfuzz=500pt filename & \hfuzz=500pt SINEX file\\
\hfuzz=500pt\includegraphics[width=1em]{element.pdf}~inputfileSinexDiscontinuities & \hfuzz=500pt filename & \hfuzz=500pt SINEX file\\
\hfuzz=500pt\includegraphics[width=1em]{element.pdf}~inputfileSinexPostSeismicDeformations & \hfuzz=500pt filename & \hfuzz=500pt SINEX file\\
\hfuzz=500pt\includegraphics[width=1em]{element.pdf}~inputfileSinexFrequencies & \hfuzz=500pt filename & \hfuzz=500pt SINEX file (XYZ or ENU)\\
\hfuzz=500pt\includegraphics[width=1em]{element-mustset-unbounded.pdf}~timeSeries & \hfuzz=500pt \hyperref[timeSeriesType]{timeSeries} & \hfuzz=500pt compute positions for these epochs based on velocity\\
\hfuzz=500pt\includegraphics[width=1em]{element.pdf}~extrapolateForward & \hfuzz=500pt boolean & \hfuzz=500pt also compute positions for epochs after last interval defined in SINEX file\\
\hfuzz=500pt\includegraphics[width=1em]{element.pdf}~extrapolateBackward & \hfuzz=500pt boolean & \hfuzz=500pt also compute positions for epochs before first interval defined in SINEX file\\
\hfuzz=500pt\includegraphics[width=1em]{element-unbounded.pdf}~stationName & \hfuzz=500pt string & \hfuzz=500pt convert only these stations\\
\hline
\end{tabularx}

\clearpage
%==================================
\subsection{SinexEccentricties2SlrPlatform}\label{SinexEccentricties2SlrPlatform}
Reads metadata like station name, station number, approximate station position and station eccentricities
from \href{https://ilrs.gsfc.nasa.gov/network/site_procedures/eccentricity.html}{Station Eccentricities Sinex File}
(une version) and write them to the \configFile{outputfileStationInfo}{platform} for each station.


\keepXColumns
\begin{tabularx}{\textwidth}{N T A}
\hline
Name & Type & Annotation\\
\hline
\hfuzz=500pt\includegraphics[width=1em]{element-mustset.pdf}~outputfileStationInfo & \hfuzz=500pt filename & \hfuzz=500pt loop variable is replaced with station name\\
\hfuzz=500pt\includegraphics[width=1em]{element.pdf}~variableLoopStation & \hfuzz=500pt string & \hfuzz=500pt variable name for station loop\\
\hfuzz=500pt\includegraphics[width=1em]{element-mustset.pdf}~inputfileSinex & \hfuzz=500pt filename & \hfuzz=500pt SINEX file (.snx or .ssc)\\
\hfuzz=500pt\includegraphics[width=1em]{element-unbounded.pdf}~stationName & \hfuzz=500pt string & \hfuzz=500pt convert only these stations\\
\hline
\end{tabularx}

\clearpage
%==================================
\subsection{SinexMetadata2GlonassFrequencyNumber}\label{SinexMetadata2GlonassFrequencyNumber}
Create \configFile{outputfileMatrixPrn2FrequencyNumber}{matrix} matrix
from \href{https://www.igs.org/mgex/metadata/#metadata}{IGS SINEX metadata format}
with the columns: GLONASS PRN, SVN, mjdStart, mjdEnd, frequencyNumber.

See also \program{RinexObservation2GnssReceiver}.


\keepXColumns
\begin{tabularx}{\textwidth}{N T A}
\hline
Name & Type & Annotation\\
\hline
\hfuzz=500pt\includegraphics[width=1em]{element-mustset.pdf}~outputfileMatrixPrn2FrequencyNumber & \hfuzz=500pt filename & \hfuzz=500pt GROOPS matrix with columns: GLONASS PRN, SVN, mjdStart, mjdEnd, frequencyNumber\\
\hfuzz=500pt\includegraphics[width=1em]{element-mustset.pdf}~inputfileSinexMetadata & \hfuzz=500pt filename & \hfuzz=500pt IGS SINEX metadata file\\
\hline
\end{tabularx}

\clearpage
%==================================
\subsection{SinexMetadata2SatelliteModel}\label{SinexMetadata2SatelliteModel}
Create \configFile{outputfileSatelliteModel}{satelliteModel} from \href{https://www.igs.org/mgex/metadata/#metadata}{IGS SINEX metadata format}.

If \configFile{inputfileSatelliteModel}{satelliteModel} is provided it is used as a basis and values are updated from the metadata file.

See also \program{SatelliteModelCreate}.


\keepXColumns
\begin{tabularx}{\textwidth}{N T A}
\hline
Name & Type & Annotation\\
\hline
\hfuzz=500pt\includegraphics[width=1em]{element-mustset.pdf}~outputfileSatelliteModel & \hfuzz=500pt filename & \hfuzz=500pt \\
\hfuzz=500pt\includegraphics[width=1em]{element-mustset.pdf}~inputfileSinexMetadata & \hfuzz=500pt filename & \hfuzz=500pt IGS SINEX metadata file\\
\hfuzz=500pt\includegraphics[width=1em]{element.pdf}~inputfileSatelliteModel & \hfuzz=500pt filename & \hfuzz=500pt base satellite model\\
\hfuzz=500pt\includegraphics[width=1em]{element-mustset.pdf}~svn & \hfuzz=500pt string & \hfuzz=500pt e.g. G040, R736, E204, C211\\
\hline
\end{tabularx}

\clearpage
%==================================
\subsection{SlrComModel2RangeBiasStationSatellite}\label{SlrComModel2RangeBiasStationSatellite}
Converts the tables of CoM corrections of José Rodríguez
(\url{https://icts-yebes.oan.es/slr/com_models/models/}) into station/satellite
specific \configFile{outputfileRangeBias}{instrument}. Only the deviations to the default value in
\configFile{inputfileSatelliteInfo}{platform} are written. This program must be called for every
provided satellite. The range bias values can be used in
\configClass{parametrization:rangeBiasStationSatelliteApriori}{slrParametrizationType:rangeBiasStationSatelliteApriori}
in \program {SlrProcessing}.

Reference:
Rodriguez J., Otsubo T., Appleby G. Upgraded Modelling for the
Determination of Centre of Mass Corrections of Geodetic SLR
Satellites: Impact on Key Parameters of the Terrestrial Reference
Frame. Journal of Geodesy, 2019. doi: 10.1007/s00190-019-01315-0


\keepXColumns
\begin{tabularx}{\textwidth}{N T A}
\hline
Name & Type & Annotation\\
\hline
\hfuzz=500pt\includegraphics[width=1em]{element-mustset.pdf}~outputfileRangeBias & \hfuzz=500pt filename & \hfuzz=500pt MISCVALUE, variable \{station\} available\\
\hfuzz=500pt\includegraphics[width=1em]{element-mustset.pdf}~inputfileSatelliteInfo & \hfuzz=500pt filename & \hfuzz=500pt \\
\hfuzz=500pt\includegraphics[width=1em]{element-mustset.pdf}~inputfile & \hfuzz=500pt filename & \hfuzz=500pt from Rodriguez model\\
\hfuzz=500pt\includegraphics[width=1em]{element.pdf}~variableLoopStation & \hfuzz=500pt string & \hfuzz=500pt variable name for station loop\\
\hfuzz=500pt\includegraphics[width=1em]{element-unbounded.pdf}~stationName & \hfuzz=500pt string & \hfuzz=500pt convert only these stations\\
\hline
\end{tabularx}

\clearpage
%==================================
\subsection{SlrSinexDataHandling2Files}\label{SlrSinexDataHandling2Files}
Converts SLR range and time bias from \verb|ILRS_Data_Handling_File_xxxx.xx.xx.snx| provided at
\url{https://cddis.nasa.gov/archive/slr/products/resource/}. The range and time bias values can be used in
\configClass{parametrization:rangeBiasXxxApriori}{slrParametrizationType:rangeBiasStationApriori}
in \program {SlrProcessing}.


\keepXColumns
\begin{tabularx}{\textwidth}{N T A}
\hline
Name & Type & Annotation\\
\hline
\hfuzz=500pt\includegraphics[width=1em]{element.pdf}~outputfileRangeBiasStation & \hfuzz=500pt filename & \hfuzz=500pt MISCVALUE [m]\\
\hfuzz=500pt\includegraphics[width=1em]{element.pdf}~outputfileRangeBiasStationSatellite & \hfuzz=500pt filename & \hfuzz=500pt MISCVALUE [m]\\
\hfuzz=500pt\includegraphics[width=1em]{element.pdf}~outputfileTimeBias & \hfuzz=500pt filename & \hfuzz=500pt MISCVALUES(bias [s], drift [s/d])\\
\hfuzz=500pt\includegraphics[width=1em]{element.pdf}~variableLoopStation & \hfuzz=500pt string & \hfuzz=500pt variable name for station loop\\
\hfuzz=500pt\includegraphics[width=1em]{element.pdf}~variableLoopSatellite & \hfuzz=500pt string & \hfuzz=500pt variable name for satellite loop\\
\hfuzz=500pt\includegraphics[width=1em]{element-mustset.pdf}~inputfileSinex & \hfuzz=500pt filename & \hfuzz=500pt SINEX file (.snx)\\
\hfuzz=500pt\includegraphics[width=1em]{element.pdf}~inputfileSatelliteId & \hfuzz=500pt filename & \hfuzz=500pt table SP3 and satellite name\\
\hfuzz=500pt\includegraphics[width=1em]{element-unbounded.pdf}~stationName & \hfuzz=500pt string & \hfuzz=500pt convert only these stations\\
\hline
\end{tabularx}

\clearpage
%==================================
\subsection{Sp3Format2Orbit}\label{Sp3Format2Orbit}
Read orbits from \href{https://files.igs.org/pub/data/format/sp3d.pdf}{SP3 format}
and write an \file{instrument file (ORBIT)}{instrument}.
The additional \config{outputfileClock} is an \file{instrument file (MISCVALUE)}{instrument}
and \config{outputfileCovariance} is an \file{instrument file (COVARIANCE3D)}{instrument}.

With \config{satelliteIdentifier} a single satellite can be selected if the \config{inputfile}s
contain more than one satellites. If \config{satelliteIdentifier} is empty the first satellite is taken.
All satellites can be selected with \config{satelliteIdentifier}=\verb|<all>|.
In this case the identifier is appended to each output file.

If \configClass{earthRotation}{earthRotationType} is provided the data are transformed
from terrestrial (TRF) to celestial reference frame (CRF).
Since SP3 orbits often use the center of Earth as a reference, a correction from center
of Earth to center of mass can be applied to the orbits by providing \configClass{gravityfield}{gravityfieldType} (e.g. ocean tides).

See also \program{Orbit2Sp3Format}.


\keepXColumns
\begin{tabularx}{\textwidth}{N T A}
\hline
Name & Type & Annotation\\
\hline
\hfuzz=500pt\includegraphics[width=1em]{element-mustset.pdf}~outputfileOrbit & \hfuzz=500pt filename & \hfuzz=500pt \\
\hfuzz=500pt\includegraphics[width=1em]{element.pdf}~outputfileClock & \hfuzz=500pt filename & \hfuzz=500pt \\
\hfuzz=500pt\includegraphics[width=1em]{element.pdf}~outputfileCovariance & \hfuzz=500pt filename & \hfuzz=500pt 3x3 epoch covariance\\
\hfuzz=500pt\includegraphics[width=1em]{element.pdf}~satelliteIdentifier & \hfuzz=500pt string & \hfuzz=500pt e.g. L09 for GRACE A, empty: take first satellite, \$<\$all\$>\$: identifier is appended to each file\\
\hfuzz=500pt\includegraphics[width=1em]{element.pdf}~earthRotation & \hfuzz=500pt \hyperref[earthRotationType]{earthRotation} & \hfuzz=500pt rotation from TRF to CRF\\
\hfuzz=500pt\includegraphics[width=1em]{element-unbounded.pdf}~gravityfield & \hfuzz=500pt \hyperref[gravityfieldType]{gravityfield} & \hfuzz=500pt degree 1 fluid mantle for CM2CE correction (SP3 orbits should be in center of Earth)\\
\hfuzz=500pt\includegraphics[width=1em]{element-mustset-unbounded.pdf}~inputfile & \hfuzz=500pt filename & \hfuzz=500pt orbits in SP3 format\\
\hline
\end{tabularx}

\clearpage
%==================================
\subsection{StarCamera2Orbex}\label{StarCamera2Orbex}
Converts satellite attitude from \file{instrument file (STARCAMERA)}{instrument} to
\href{http://acc.igs.org/misc/ORBEX009.pdf}{ORBEX file format} (quaternions).

If \configClass{earthRotation}{earthRotationType} is provided, the output file contains quaternions
for rotation from TRF to satellite body frame (IGS/ORBEX convention),
otherwise the rotation is from CRF to satellite body frame.

See also \program{GnssOrbex2StarCamera}, \program{SimulateStarCameraGnss}.


\keepXColumns
\begin{tabularx}{\textwidth}{N T A}
\hline
Name & Type & Annotation\\
\hline
\hfuzz=500pt\includegraphics[width=1em]{element-mustset.pdf}~outputfileOrbex & \hfuzz=500pt filename & \hfuzz=500pt ORBEX file\\
\hfuzz=500pt\includegraphics[width=1em]{element-mustset-unbounded.pdf}~satellite & \hfuzz=500pt sequence & \hfuzz=500pt \\
\hfuzz=500pt\includegraphics[width=1em]{connector.pdf}\includegraphics[width=1em]{element-mustset.pdf}~inputfileStarCamera & \hfuzz=500pt filename & \hfuzz=500pt \\
\hfuzz=500pt\includegraphics[width=1em]{connector.pdf}\includegraphics[width=1em]{element-mustset.pdf}~identifier & \hfuzz=500pt string & \hfuzz=500pt string identifier (e.g. GNSS PRN: G01)\\
\hfuzz=500pt\includegraphics[width=1em]{connector.pdf}\includegraphics[width=1em]{element.pdf}~description & \hfuzz=500pt string & \hfuzz=500pt e.g. BLOCK IIR-B, GRACE\\
\hfuzz=500pt\includegraphics[width=1em]{element.pdf}~earthRotation & \hfuzz=500pt \hyperref[earthRotationType]{earthRotation} & \hfuzz=500pt rotate data into Earth-fixed frame\\
\hfuzz=500pt\includegraphics[width=1em]{element-unbounded.pdf}~timeSeries & \hfuzz=500pt \hyperref[timeSeriesType]{timeSeries} & \hfuzz=500pt resample to these epochs (otherwise input file epochs are used)\\
\hfuzz=500pt\includegraphics[width=1em]{element-mustset.pdf}~interpolationDegree & \hfuzz=500pt uint & \hfuzz=500pt for attitude and Earth rotation interpolation\\
\hfuzz=500pt\includegraphics[width=1em]{element-mustset.pdf}~description & \hfuzz=500pt string & \hfuzz=500pt description of file contents\\
\hfuzz=500pt\includegraphics[width=1em]{element-mustset.pdf}~createdBy & \hfuzz=500pt string & \hfuzz=500pt name of agency\\
\hfuzz=500pt\includegraphics[width=1em]{element-mustset.pdf}~inputData & \hfuzz=500pt string & \hfuzz=500pt description of input data (see ORBEX description)\\
\hfuzz=500pt\includegraphics[width=1em]{element-mustset.pdf}~contact & \hfuzz=500pt string & \hfuzz=500pt email address\\
\hfuzz=500pt\includegraphics[width=1em]{element-mustset.pdf}~referenceFrame & \hfuzz=500pt string & \hfuzz=500pt reference frame used in file\\
\hfuzz=500pt\includegraphics[width=1em]{element-unbounded.pdf}~comment & \hfuzz=500pt string & \hfuzz=500pt \\
\hline
\end{tabularx}

\clearpage
%==================================
\subsection{Swarm2Starcamera}\label{Swarm2Starcamera}
This program reads SWARM star camera data given in the cdf format
and before converted to an ascii file using the program \verb|cdfexport|
provided by the Goddard Space Flight Center (\url{http://cdf.gsfc.nasa.gov/}).


\keepXColumns
\begin{tabularx}{\textwidth}{N T A}
\hline
Name & Type & Annotation\\
\hline
\hfuzz=500pt\includegraphics[width=1em]{element-mustset.pdf}~outputfileStarCamera & \hfuzz=500pt filename & \hfuzz=500pt \\
\hfuzz=500pt\includegraphics[width=1em]{element-mustset.pdf}~earthRotation & \hfuzz=500pt \hyperref[earthRotationType]{earthRotation} & \hfuzz=500pt \\
\hfuzz=500pt\includegraphics[width=1em]{element-mustset-unbounded.pdf}~inputfile & \hfuzz=500pt filename & \hfuzz=500pt \\
\hline
\end{tabularx}

\clearpage
%==================================
\subsection{TerraSarTandem2Orbit}\label{TerraSarTandem2Orbit}
This program reads in TerraSar-X or Tandem-X orbits in the special CHORB format and takes the appropriate
time frame as stated in the document header.
A description of the format can be found under: \url{http://op.gfz-potsdam.de/champ/docs_CHAMP/CH-GFZ-FD-002.pdf}


\keepXColumns
\begin{tabularx}{\textwidth}{N T A}
\hline
Name & Type & Annotation\\
\hline
\hfuzz=500pt\includegraphics[width=1em]{element-mustset.pdf}~outputfileOrbit & \hfuzz=500pt filename & \hfuzz=500pt \\
\hfuzz=500pt\includegraphics[width=1em]{element-mustset.pdf}~earthRotation & \hfuzz=500pt \hyperref[earthRotationType]{earthRotation} & \hfuzz=500pt \\
\hfuzz=500pt\includegraphics[width=1em]{element-mustset-unbounded.pdf}~inputfile & \hfuzz=500pt filename & \hfuzz=500pt orbits in CHORB format\\
\hline
\end{tabularx}

\clearpage
%==================================
\subsection{TerraSarTandem2StarCamera}\label{TerraSarTandem2StarCamera}
This program reads in TerraSar-X or Tandem-X star camera data given in the special format.


\keepXColumns
\begin{tabularx}{\textwidth}{N T A}
\hline
Name & Type & Annotation\\
\hline
\hfuzz=500pt\includegraphics[width=1em]{element-mustset.pdf}~outputfileStarCamera & \hfuzz=500pt filename & \hfuzz=500pt \\
\hfuzz=500pt\includegraphics[width=1em]{element-mustset-unbounded.pdf}~inputfile & \hfuzz=500pt filename & \hfuzz=500pt \\
\hline
\end{tabularx}

\clearpage
%==================================
\subsection{Tle2Orbit}\label{Tle2Orbit}
This program computes the \configFile{outputfileOrbit}{instrument}
from two-line elements (TLE/3LE)
as can be found at e.g. \url{http://celestrak.org/NORAD/elements/}.
The first satellite in the input file that matches the wildcard of \config{satelliteName} is used.

The program uses the Simplified General Perturbation (SGP) model. More information can
be found in the Revisiting Spacetrack Report 3 by Vallado et al. 2006.


\keepXColumns
\begin{tabularx}{\textwidth}{N T A}
\hline
Name & Type & Annotation\\
\hline
\hfuzz=500pt\includegraphics[width=1em]{element-mustset.pdf}~outputfileOrbit & \hfuzz=500pt filename & \hfuzz=500pt \\
\hfuzz=500pt\includegraphics[width=1em]{element-mustset.pdf}~inputfileTLE & \hfuzz=500pt filename & \hfuzz=500pt two line elements (TLE/3LE)\\
\hfuzz=500pt\includegraphics[width=1em]{element.pdf}~satelliteName & \hfuzz=500pt string & \hfuzz=500pt first name of wildcard match is used\\
\hfuzz=500pt\includegraphics[width=1em]{element-mustset-unbounded.pdf}~timeSeries & \hfuzz=500pt \hyperref[timeSeriesType]{timeSeries} & \hfuzz=500pt output orbit at these times\\
\hfuzz=500pt\includegraphics[width=1em]{element-mustset.pdf}~earthRotation & \hfuzz=500pt \hyperref[earthRotationType]{earthRotation} & \hfuzz=500pt rotation to CRF\\
\hline
\end{tabularx}

\clearpage
%==================================
\subsection{ViennaMappingFunctionGrid2File}\label{ViennaMappingFunctionGrid2File}
This program converts the gridded time series of the Vienna Mapping Functions (VMF) into
the \file{GROOPS file format}{griddedDataTimeSeries}.

Gridded VMF data is available at: \url{https://vmf.geo.tuwien.ac.at/trop_products/GRID/}


\keepXColumns
\begin{tabularx}{\textwidth}{N T A}
\hline
Name & Type & Annotation\\
\hline
\hfuzz=500pt\includegraphics[width=1em]{element.pdf}~outputfileVmfCoefficients & \hfuzz=500pt filename & \hfuzz=500pt \\
\hfuzz=500pt\includegraphics[width=1em]{element-mustset-unbounded.pdf}~inputfile & \hfuzz=500pt filename & \hfuzz=500pt files must be given for each point in time\\
\hfuzz=500pt\includegraphics[width=1em]{element-mustset-unbounded.pdf}~timeSeries & \hfuzz=500pt \hyperref[timeSeriesType]{timeSeries} & \hfuzz=500pt times of input files\\
\hfuzz=500pt\includegraphics[width=1em]{element.pdf}~deltaLambda & \hfuzz=500pt angle & \hfuzz=500pt [deg] sampling in longitude\\
\hfuzz=500pt\includegraphics[width=1em]{element.pdf}~deltaPhi & \hfuzz=500pt angle & \hfuzz=500pt [deg] sampling in latitude\\
\hfuzz=500pt\includegraphics[width=1em]{element.pdf}~isCellRegistered & \hfuzz=500pt boolean & \hfuzz=500pt grid points represent cells (VMF3), not grid corners (VMF1)\\
\hline
\end{tabularx}

\clearpage
%==================================
\subsection{ViennaMappingFunctionStation2File}\label{ViennaMappingFunctionStation2File}
Converts Vienna Mapping Functions (VMF) station time series into \file{GROOPS file format}{griddedDataTimeSeries}.

Station-wise VMF data for GNSS is available at: \url{https://vmf.geo.tuwien.ac.at/trop_products/GNSS/}


\keepXColumns
\begin{tabularx}{\textwidth}{N T A}
\hline
Name & Type & Annotation\\
\hline
\hfuzz=500pt\includegraphics[width=1em]{element.pdf}~outputfileVmfCoefficients & \hfuzz=500pt filename & \hfuzz=500pt \\
\hfuzz=500pt\includegraphics[width=1em]{element-mustset.pdf}~inputfileStationInfo & \hfuzz=500pt filename & \hfuzz=500pt \\
\hfuzz=500pt\includegraphics[width=1em]{element-mustset.pdf}~inputfileStation & \hfuzz=500pt filename & \hfuzz=500pt \\
\hfuzz=500pt\includegraphics[width=1em]{element-mustset.pdf}~inputfileVmf & \hfuzz=500pt filename & \hfuzz=500pt \\
\hline
\end{tabularx}

\clearpage
%==================================
\section{Programs: Deprecated}
\subsection{GnssAttitude2Orbex}\label{GnssAttitude2Orbex}
DEPRECATED. Please use \program{StarCamera2Orbex} instead.


\keepXColumns
\begin{tabularx}{\textwidth}{N T A}
\hline
Name & Type & Annotation\\
\hline
\hfuzz=500pt\includegraphics[width=1em]{element-mustset.pdf}~outputfileOrbex & \hfuzz=500pt filename & \hfuzz=500pt ORBEX file\\
\hfuzz=500pt\includegraphics[width=1em]{element-mustset-unbounded.pdf}~inputfileTransmitterList & \hfuzz=500pt filename & \hfuzz=500pt ASCII list with transmitter PRNs\\
\hfuzz=500pt\includegraphics[width=1em]{element-mustset.pdf}~inputfileAttitude & \hfuzz=500pt filename & \hfuzz=500pt instrument file containing attitude\\
\hfuzz=500pt\includegraphics[width=1em]{element.pdf}~variablePrn & \hfuzz=500pt string & \hfuzz=500pt loop variable for PRNs from transmitter list\\
\hfuzz=500pt\includegraphics[width=1em]{element-unbounded.pdf}~timeSeries & \hfuzz=500pt \hyperref[timeSeriesType]{timeSeries} & \hfuzz=500pt resample to these epochs (otherwise input file epochs are used)\\
\hfuzz=500pt\includegraphics[width=1em]{element.pdf}~earthRotation & \hfuzz=500pt \hyperref[earthRotationType]{earthRotation} & \hfuzz=500pt rotate data into Earth-fixed frame\\
\hfuzz=500pt\includegraphics[width=1em]{element-mustset.pdf}~interpolationDegree & \hfuzz=500pt uint & \hfuzz=500pt for attitude and Earth rotation interpolation\\
\hfuzz=500pt\includegraphics[width=1em]{element-mustset.pdf}~description & \hfuzz=500pt string & \hfuzz=500pt description of file contents\\
\hfuzz=500pt\includegraphics[width=1em]{element-mustset.pdf}~createdBy & \hfuzz=500pt string & \hfuzz=500pt name of agency\\
\hfuzz=500pt\includegraphics[width=1em]{element-mustset.pdf}~inputData & \hfuzz=500pt string & \hfuzz=500pt description of input data (see ORBEX description)\\
\hfuzz=500pt\includegraphics[width=1em]{element-mustset.pdf}~contact & \hfuzz=500pt string & \hfuzz=500pt email address\\
\hfuzz=500pt\includegraphics[width=1em]{element-mustset.pdf}~referenceFrame & \hfuzz=500pt string & \hfuzz=500pt reference frame used in file\\
\hfuzz=500pt\includegraphics[width=1em]{element-unbounded.pdf}~comment & \hfuzz=500pt string & \hfuzz=500pt \\
\hline
\end{tabularx}

\clearpage
%==================================
\subsection{GnssPrn2SvnBlockVariables}\label{GnssPrn2SvnBlockVariables}
DEPRECATED. This program no longer works!

Setup up a \configClass{loop:platformEquipment}{loopType:platformEquipment} instead with
\begin{itemize}
  \item \configFile{inputfilePlatform}{platform}: the old \config{inputfileTransmitterInfo}
  \item \config{equipmentType}         = \verb|gnssAntenna|
  \item \config{variableLoopName}      = \verb|block|
  \item \config{variableLoopSerial}    = \verb|svn|
  \item \config{variableLoopTimeStart} = \verb|svnTimeStart|
  \item \config{variableLoopTimeEnd}   = \verb|svnTimeEnd|
  \item \configClass{condition:expression}{conditionType:expression}
  \begin{itemize}
    \item \config{expression} = \verb|(svnTimeStart <= time) && (time < svnTimeEnd)|
  \end{itemize}
\end{itemize}
Attribute this loop to programs, which uses the variables.

\clearpage
%==================================
\subsection{GridRectangular2NetCdf}\label{GridRectangular2NetCdf}
DEPRECATED. Please use \program{GriddedData2NetCdf} or \program{GriddedDataTimeSeries2NetCdf} instead.


\keepXColumns
\begin{tabularx}{\textwidth}{N T A}
\hline
Name & Type & Annotation\\
\hline
\hfuzz=500pt\includegraphics[width=1em]{element-mustset.pdf}~outputfileNetCdf & \hfuzz=500pt filename & \hfuzz=500pt file name of NetCDF output\\
\hfuzz=500pt\includegraphics[width=1em]{element-mustset-unbounded.pdf}~inputfileGridRectangular & \hfuzz=500pt filename & \hfuzz=500pt input grid sequence\\
\hfuzz=500pt\includegraphics[width=1em]{element-unbounded.pdf}~times & \hfuzz=500pt \hyperref[timeSeriesType]{timeSeries} & \hfuzz=500pt values for time axis (COARDS specification)\\
\hfuzz=500pt\includegraphics[width=1em]{element-mustset-unbounded.pdf}~dataVariable & \hfuzz=500pt sequence & \hfuzz=500pt metadata for data variables\\
\hfuzz=500pt\includegraphics[width=1em]{connector.pdf}\includegraphics[width=1em]{element.pdf}~selectDataField & \hfuzz=500pt uint & \hfuzz=500pt input data column\\
\hfuzz=500pt\includegraphics[width=1em]{connector.pdf}\includegraphics[width=1em]{element-mustset.pdf}~name & \hfuzz=500pt string & \hfuzz=500pt netCDF variable name\\
\hfuzz=500pt\includegraphics[width=1em]{connector.pdf}\includegraphics[width=1em]{element-mustset.pdf}~dataType & \hfuzz=500pt choice & \hfuzz=500pt \\
\hfuzz=500pt\quad\includegraphics[width=1em]{connector.pdf}\includegraphics[width=1em]{element-mustset.pdf}~double & \hfuzz=500pt  & \hfuzz=500pt \\
\hfuzz=500pt\quad\includegraphics[width=1em]{connector.pdf}\includegraphics[width=1em]{element-mustset.pdf}~float & \hfuzz=500pt  & \hfuzz=500pt \\
\hfuzz=500pt\quad\includegraphics[width=1em]{connector.pdf}\includegraphics[width=1em]{element-mustset.pdf}~int & \hfuzz=500pt  & \hfuzz=500pt \\
\hfuzz=500pt\includegraphics[width=1em]{connector.pdf}\includegraphics[width=1em]{element-unbounded.pdf}~attribute & \hfuzz=500pt choice & \hfuzz=500pt netCDF attributes\\
\hfuzz=500pt\quad\includegraphics[width=1em]{connector.pdf}\includegraphics[width=1em]{element-mustset.pdf}~text & \hfuzz=500pt sequence & \hfuzz=500pt \\
\hfuzz=500pt\quad\quad\includegraphics[width=1em]{connector.pdf}\includegraphics[width=1em]{element-mustset.pdf}~name & \hfuzz=500pt string & \hfuzz=500pt \\
\hfuzz=500pt\quad\quad\includegraphics[width=1em]{connector.pdf}\includegraphics[width=1em]{element-mustset.pdf}~value & \hfuzz=500pt string & \hfuzz=500pt \\
\hfuzz=500pt\quad\includegraphics[width=1em]{connector.pdf}\includegraphics[width=1em]{element-mustset.pdf}~value & \hfuzz=500pt sequence & \hfuzz=500pt \\
\hfuzz=500pt\quad\quad\includegraphics[width=1em]{connector.pdf}\includegraphics[width=1em]{element-mustset.pdf}~name & \hfuzz=500pt string & \hfuzz=500pt \\
\hfuzz=500pt\quad\quad\includegraphics[width=1em]{connector.pdf}\includegraphics[width=1em]{element-mustset-unbounded.pdf}~value & \hfuzz=500pt double & \hfuzz=500pt \\
\hfuzz=500pt\quad\quad\includegraphics[width=1em]{connector.pdf}\includegraphics[width=1em]{element-mustset.pdf}~dataType & \hfuzz=500pt choice & \hfuzz=500pt \\
\hfuzz=500pt\quad\quad\quad\includegraphics[width=1em]{connector.pdf}\includegraphics[width=1em]{element-mustset.pdf}~double & \hfuzz=500pt  & \hfuzz=500pt \\
\hfuzz=500pt\quad\quad\quad\includegraphics[width=1em]{connector.pdf}\includegraphics[width=1em]{element-mustset.pdf}~float & \hfuzz=500pt  & \hfuzz=500pt \\
\hfuzz=500pt\quad\quad\quad\includegraphics[width=1em]{connector.pdf}\includegraphics[width=1em]{element-mustset.pdf}~int & \hfuzz=500pt  & \hfuzz=500pt \\
\hfuzz=500pt\includegraphics[width=1em]{element-unbounded.pdf}~globalAttribute & \hfuzz=500pt choice & \hfuzz=500pt additional meta data\\
\hfuzz=500pt\includegraphics[width=1em]{connector.pdf}\includegraphics[width=1em]{element-mustset.pdf}~text & \hfuzz=500pt sequence & \hfuzz=500pt \\
\hfuzz=500pt\quad\includegraphics[width=1em]{connector.pdf}\includegraphics[width=1em]{element-mustset.pdf}~name & \hfuzz=500pt string & \hfuzz=500pt \\
\hfuzz=500pt\quad\includegraphics[width=1em]{connector.pdf}\includegraphics[width=1em]{element-mustset.pdf}~value & \hfuzz=500pt string & \hfuzz=500pt \\
\hfuzz=500pt\includegraphics[width=1em]{connector.pdf}\includegraphics[width=1em]{element-mustset.pdf}~value & \hfuzz=500pt sequence & \hfuzz=500pt \\
\hfuzz=500pt\quad\includegraphics[width=1em]{connector.pdf}\includegraphics[width=1em]{element-mustset.pdf}~name & \hfuzz=500pt string & \hfuzz=500pt \\
\hfuzz=500pt\quad\includegraphics[width=1em]{connector.pdf}\includegraphics[width=1em]{element-mustset-unbounded.pdf}~value & \hfuzz=500pt double & \hfuzz=500pt \\
\hfuzz=500pt\quad\includegraphics[width=1em]{connector.pdf}\includegraphics[width=1em]{element-mustset.pdf}~dataType & \hfuzz=500pt choice & \hfuzz=500pt \\
\hfuzz=500pt\quad\quad\includegraphics[width=1em]{connector.pdf}\includegraphics[width=1em]{element-mustset.pdf}~double & \hfuzz=500pt  & \hfuzz=500pt \\
\hfuzz=500pt\quad\quad\includegraphics[width=1em]{connector.pdf}\includegraphics[width=1em]{element-mustset.pdf}~float & \hfuzz=500pt  & \hfuzz=500pt \\
\hfuzz=500pt\quad\quad\includegraphics[width=1em]{connector.pdf}\includegraphics[width=1em]{element-mustset.pdf}~int & \hfuzz=500pt  & \hfuzz=500pt \\
\hline
\end{tabularx}

\clearpage
%==================================
\subsection{NetCdf2GridRectangular}\label{NetCdf2GridRectangular}
DEPRECATED. Please use \program{NetCdf2GriddedData} or \program{NetCdf2GriddedDataTimeSeries} instead.


\keepXColumns
\begin{tabularx}{\textwidth}{N T A}
\hline
Name & Type & Annotation\\
\hline
\hfuzz=500pt\includegraphics[width=1em]{element-mustset.pdf}~outputfileGridRectangular & \hfuzz=500pt filename & \hfuzz=500pt One grid for each epoch in the NetCDF file is written. Use loopTimeVariable as template.\\
\hfuzz=500pt\includegraphics[width=1em]{element-mustset.pdf}~loopTimeVariable & \hfuzz=500pt string & \hfuzz=500pt \\
\hfuzz=500pt\includegraphics[width=1em]{element-mustset.pdf}~inputfileNetCdf & \hfuzz=500pt filename & \hfuzz=500pt \\
\hfuzz=500pt\includegraphics[width=1em]{element-mustset.pdf}~variableNameLongitude & \hfuzz=500pt string & \hfuzz=500pt name of NetCDF variable\\
\hfuzz=500pt\includegraphics[width=1em]{element-mustset.pdf}~variableNameLatitude & \hfuzz=500pt string & \hfuzz=500pt name of NetCDF variable\\
\hfuzz=500pt\includegraphics[width=1em]{element.pdf}~variableNameTime & \hfuzz=500pt string & \hfuzz=500pt name of NetCDF variable (leave blank for static grids)\\
\hfuzz=500pt\includegraphics[width=1em]{element-unbounded.pdf}~variableNameData & \hfuzz=500pt string & \hfuzz=500pt name of NetCDF variable\\
\hfuzz=500pt\includegraphics[width=1em]{element.pdf}~R & \hfuzz=500pt double & \hfuzz=500pt reference radius for ellipsoidal coordinates\\
\hfuzz=500pt\includegraphics[width=1em]{element.pdf}~inverseFlattening & \hfuzz=500pt double & \hfuzz=500pt reference flattening for ellipsoidal coordinates\\
\hline
\end{tabularx}

\clearpage
%==================================
\subsection{Sinex2StationPosition}\label{Sinex2StationPosition}
DEPRECATED. Please use \program{Sinex2StationPositions} instead.


\keepXColumns
\begin{tabularx}{\textwidth}{N T A}
\hline
Name & Type & Annotation\\
\hline
\hfuzz=500pt\includegraphics[width=1em]{element-mustset.pdf}~outputfileInstrument & \hfuzz=500pt filename & \hfuzz=500pt loop variable is replaced with station name (e.g. wtzz)\\
\hfuzz=500pt\includegraphics[width=1em]{element-mustset.pdf}~inputfileSinex & \hfuzz=500pt filename & \hfuzz=500pt SINEX file (.snx or .ssc)\\
\hfuzz=500pt\includegraphics[width=1em]{element.pdf}~inputfileDiscontinuities & \hfuzz=500pt filename & \hfuzz=500pt discontinuities file per station; loop variable is replaced with station name (e.g. wtzz)\\
\hfuzz=500pt\includegraphics[width=1em]{element.pdf}~variableLoopStation & \hfuzz=500pt string & \hfuzz=500pt variable name for station loop\\
\hfuzz=500pt\includegraphics[width=1em]{element-unbounded.pdf}~stationName & \hfuzz=500pt string & \hfuzz=500pt convert only these stations\\
\hfuzz=500pt\includegraphics[width=1em]{element-unbounded.pdf}~timeSeries & \hfuzz=500pt \hyperref[timeSeriesType]{timeSeries} & \hfuzz=500pt compute positions for these epochs based on velocity\\
\hfuzz=500pt\includegraphics[width=1em]{element.pdf}~extrapolateForward & \hfuzz=500pt boolean & \hfuzz=500pt also compute positions for epochs after last interval defined in SINEX file\\
\hfuzz=500pt\includegraphics[width=1em]{element.pdf}~extrapolateBackward & \hfuzz=500pt boolean & \hfuzz=500pt also compute positions for epochs before first interval defined in SINEX file\\
\hline
\end{tabularx}

\clearpage
%==================================
\subsection{Sinex2StationPostSeismicDeformation}\label{Sinex2StationPostSeismicDeformation}
DEPRECATED. Please use \program{Sinex2StationPositions} instead.


\keepXColumns
\begin{tabularx}{\textwidth}{N T A}
\hline
Name & Type & Annotation\\
\hline
\hfuzz=500pt\includegraphics[width=1em]{element-mustset.pdf}~outputfileInstrument & \hfuzz=500pt filename & \hfuzz=500pt deformation time series\\
\hfuzz=500pt\includegraphics[width=1em]{element-mustset.pdf}~inputfileSinex & \hfuzz=500pt filename & \hfuzz=500pt ITRF post-seismic deformation SINEX file\\
\hfuzz=500pt\includegraphics[width=1em]{element-mustset-unbounded.pdf}~timeSeries & \hfuzz=500pt \hyperref[timeSeriesType]{timeSeries} & \hfuzz=500pt compute deformation for these epochs\\
\hfuzz=500pt\includegraphics[width=1em]{element-mustset.pdf}~stationName & \hfuzz=500pt string & \hfuzz=500pt \\
\hfuzz=500pt\includegraphics[width=1em]{element.pdf}~localLevelFrame & \hfuzz=500pt boolean & \hfuzz=500pt output in North, East, Up local-level frame\\
\hline
\end{tabularx}

\clearpage
%==================================


\chapter{Classes}\label{classes}
This chapter details classes included in GROOPS, describing what they are and what they do. For usage examples see the cookbook in Chapter~\ref{general.cookbook}.
% auto generated by GROOPS

\section{AutoregressiveModelSequence}\label{autoregressiveModelSequenceType}
Represents a sequence of multivariate autoregressive (AR) models with increasing order $p$.
The AR models should be stored as \file{matrix file}{matrix} in the \reference{GROOPS definition of
AR models}{fundamentals.autoregressiveModel}.
The required AR models can be computed with \program{CovarianceMatrix2AutoregressiveModel},
and passed to this class through
\config{inputfileAutoregressiveModel} in increasing order.

The main purpose of AutoregressiveModelSequence is to use AR models of the form
\begin{equation}
  \label{eq:ar-model}
  \mathbf{y}_e(t_i) = \sum_{k=1}^p \mathbf{\Phi}^{(p)}_k\mathbf{y}_e(t_{i-k}) + \mathbf{w}(t_i),
  \hspace{5pt} \mathbf{w}(t_i) \sim \mathcal{N}(0, \mathbf{\Sigma}^{(p)}_\mathbf{w}),
\end{equation}
to create pseudo-observation equations
\begin{equation}
  \label{eq:pseudo-observations-transformed}
  0 = \bar{\mathbf{\Phi}} \Delta\mathbf{y} + \bar{\mathbf{w}}, \hspace{5pt} \bar{\mathbf{w}} \sim
  \mathcal{N}(0, \bar{\mathbf{\Sigma}}_{\bar{\mathbf{w}}}),
\end{equation}
with
\begin{equation}
  \label{eq:pseudo-observations-ar}
  \bar{\mathbf{\Phi}} =
  \begin{bmatrix}
    \mathbf{I} & & & & & \\
    -\mathbf{\Phi}^{(1)}_1 & \mathbf{I} & & & &  \\
    -\mathbf{\Phi}^{(2)}_2 & -\mathbf{\Phi}^{(2)}_1 & \mathbf{I} & & & \\
    -\mathbf{\Phi}^{(3)}_3 & -\mathbf{\Phi}^ {(3)}_2 & -\mathbf{\Phi}^ {(3)}_1 & \mathbf{I} & &  \\
    & -\mathbf{\Phi}^{(3)}_3 & -\mathbf{\Phi}^ {(3)}_2 & -\mathbf{\Phi}^ {(3)}_1 & \mathbf{I} &  \\
    & & \ddots & \ddots & \ddots & \ddots  \\
  \end{bmatrix},
  \hspace{15pt}
  \bar{\mathbf{\Sigma}}_{\bar{\mathbf{w}}} =
  \bar{\mathbf{\Sigma}}_{\bar{\mathbf{w}}} =
  \begin{bmatrix}
    \mathbf{\Sigma}^{(0)}_{\mathbf{w}} & & & & & \\
    & \mathbf{\Sigma}^{(1)}_{\mathbf{w}} & & & & \\
    & & \mathbf{\Sigma}^{(2)}_{\mathbf{w}} & & & \\
    & & & \mathbf{\Sigma}^{(3)}_{\mathbf{w}} & & \\
    & & & & \mathbf{\Sigma}^{(3)}_{\mathbf{w}} &  \\
    & & & & & \ddots \\
  \end{bmatrix}.
\end{equation}
used to constrain high-frequency temporal gravity field variations (see
\program{KalmanSmootherLeastSquares}, \program{NormalsBuildShortTimeStaticLongTime},
\program{PreprocessingSst}).

The corresponding normal equation coefficient matrix is given by
\begin{equation}
  \label{eq:ar-normals}
  \bar{\mathbf{\Phi}}^T\bar{\mathbf{\Sigma}}^{-1}_{\bar{\mathbf{w}}}\bar{\mathbf{\Phi}}
\end{equation}
and if all AR models are estimated from the same sample its inverse is a block-Toeplitz covariance matrix
\begin{equation}
  (\mathbf{\Sigma}_{\mathbf{y}_m})_{ij} =
  \begin{cases}
 \mathbf{\Sigma}(|j-i|) & \text{for } i \leq j \\
 \mathbf{\Sigma}(|j-i|))^T & \text{otherwise}
 \end{cases},
\end{equation}
which can be computed using \program{AutoregressiveModel2CovarianceMatrix}.

A detailed description with applications can be found in:
Kvas, A., Mayer-Gürr, T. GRACE gravity field recovery with background model uncertainties.
J Geod 93, 2543–2552 (2019). \url{https://doi.org/10.1007/s00190-019-01314-1}


\keepXColumns
\begin{tabularx}{\textwidth}{N T A}
\hline
Name & Type & Annotation\\
\hline
\hfuzz=500pt\includegraphics[width=1em]{element-mustset.pdf}~autoregressiveModelSequenceType & \hfuzz=500pt sequence & \hfuzz=500pt \\
\hfuzz=500pt\includegraphics[width=1em]{connector.pdf}\includegraphics[width=1em]{element-mustset-unbounded.pdf}~inputfileAutoregressiveModel & \hfuzz=500pt filename & \hfuzz=500pt matrix file containing an AR model\\
\hfuzz=500pt\includegraphics[width=1em]{connector.pdf}\includegraphics[width=1em]{element.pdf}~sigma0 & \hfuzz=500pt double & \hfuzz=500pt a-priori sigma for white noise covariance\\
\hline
\end{tabularx}

\clearpage
%==================================

\section{Border}\label{borderType}
With this class you can select one or more region on the surface of the Earth.
In every instance of Border you can choose wether the specific region is excluded
from the overall result with the switch \config{exclude}.
To determine wether a specific point will be used furthermore the following algorithm will be applied:
In a first step all points are selected if first border excludes points otherwise all points excluded.
When every point will be tested for each instance of border from top to bottom.
If the point is not in the selected region nothing happens.
Otherwise it will included or excluded depending on the switch \config{exclude}.

First Example: The border excludes all continental areas.
The result are points on the oceans only.

Second Example: First border describes the continent north america. The next borders
excludes the great lakes and the last border describes Washington island.
In this configuration points are selected if they are inside north america
but not in the area of the great lakes. But if the point is on Washington island
it will be included again.


\subsection{Rectangle}
The region is restricted along lines of geographical coordinates.
\config{minPhi} and \config{maxPhi} describe the lower and the upper bound of the region.
\config{minLambda} and \config{maxLambda} define the left and right bound.


\keepXColumns
\begin{tabularx}{\textwidth}{N T A}
\hline
Name & Type & Annotation\\
\hline
\hfuzz=500pt\includegraphics[width=1em]{element-mustset.pdf}~minLambda & \hfuzz=500pt angle & \hfuzz=500pt \\
\hfuzz=500pt\includegraphics[width=1em]{element-mustset.pdf}~maxLambda & \hfuzz=500pt angle & \hfuzz=500pt \\
\hfuzz=500pt\includegraphics[width=1em]{element-mustset.pdf}~minPhi & \hfuzz=500pt angle & \hfuzz=500pt \\
\hfuzz=500pt\includegraphics[width=1em]{element-mustset.pdf}~maxPhi & \hfuzz=500pt angle & \hfuzz=500pt \\
\hfuzz=500pt\includegraphics[width=1em]{element.pdf}~exclude & \hfuzz=500pt boolean & \hfuzz=500pt dismiss points inside\\
\hline
\end{tabularx}


\subsection{Cap}
The region is defined by a spherical cap with the center given in geographical coordinates
longitude (\config{lambdaCenter}) and latitude (\config{phiCenter}).
The radius of the cap is given as aperture angle \config{psi}.

\fig{!hb}{0.4}{borderCap}{fig:borderCap}{spherical cap}


\keepXColumns
\begin{tabularx}{\textwidth}{N T A}
\hline
Name & Type & Annotation\\
\hline
\hfuzz=500pt\includegraphics[width=1em]{element-mustset.pdf}~lambdaCenter & \hfuzz=500pt angle & \hfuzz=500pt longitude of the center of the cap\\
\hfuzz=500pt\includegraphics[width=1em]{element-mustset.pdf}~phiCenter & \hfuzz=500pt angle & \hfuzz=500pt latitude of the center of the cap\\
\hfuzz=500pt\includegraphics[width=1em]{element-mustset.pdf}~psi & \hfuzz=500pt angle & \hfuzz=500pt aperture angle (radius)\\
\hfuzz=500pt\includegraphics[width=1em]{element.pdf}~exclude & \hfuzz=500pt boolean & \hfuzz=500pt dismiss points inside\\
\hline
\end{tabularx}


\subsection{Polygon}\label{borderType:polygon}
The region is defined by \configFile{inputfilePolygon}{polygon}
containing one or more polygons given in longitude and latitude.
An additional \config{buffer} around the polygon can be defined.
Use a negative value to shrink the polygon area.


\keepXColumns
\begin{tabularx}{\textwidth}{N T A}
\hline
Name & Type & Annotation\\
\hline
\hfuzz=500pt\includegraphics[width=1em]{element-mustset.pdf}~inputfilePolygon & \hfuzz=500pt filename & \hfuzz=500pt \\
\hfuzz=500pt\includegraphics[width=1em]{element.pdf}~buffer & \hfuzz=500pt double & \hfuzz=500pt buffer around polygon [km], \$<\$0: inside\\
\hfuzz=500pt\includegraphics[width=1em]{element.pdf}~exclude & \hfuzz=500pt boolean & \hfuzz=500pt dismiss points inside\\
\hline
\end{tabularx}

\clearpage
%==================================

\section{Condition}\label{conditionType}
Test for conditions. See \reference{Loop and conditions}{general.loopsAndConditions} for usage.


\subsection{FileExist}
Check for a file or directory existing.


\keepXColumns
\begin{tabularx}{\textwidth}{N T A}
\hline
Name & Type & Annotation\\
\hline
\hfuzz=500pt\includegraphics[width=1em]{element-mustset.pdf}~file & \hfuzz=500pt filename & \hfuzz=500pt \\
\hline
\end{tabularx}


\subsection{Command}
Execute command and check success.


\keepXColumns
\begin{tabularx}{\textwidth}{N T A}
\hline
Name & Type & Annotation\\
\hline
\hfuzz=500pt\includegraphics[width=1em]{element-mustset.pdf}~command & \hfuzz=500pt filename & \hfuzz=500pt \\
\hfuzz=500pt\includegraphics[width=1em]{element.pdf}~silently & \hfuzz=500pt boolean & \hfuzz=500pt without showing the output.\\
\hline
\end{tabularx}


\subsection{Expression}
Evaluate expression.


\keepXColumns
\begin{tabularx}{\textwidth}{N T A}
\hline
Name & Type & Annotation\\
\hline
\hfuzz=500pt\includegraphics[width=1em]{element-mustset.pdf}~expression & \hfuzz=500pt expression & \hfuzz=500pt \\
\hline
\end{tabularx}


\subsection{StringContainsPattern}
Determines if there is a match between a pattern or a regular expression and some subsequence in a string.


\keepXColumns
\begin{tabularx}{\textwidth}{N T A}
\hline
Name & Type & Annotation\\
\hline
\hfuzz=500pt\includegraphics[width=1em]{element.pdf}~string & \hfuzz=500pt string & \hfuzz=500pt should contain a \{variable\}\\
\hfuzz=500pt\includegraphics[width=1em]{element.pdf}~pattern & \hfuzz=500pt string & \hfuzz=500pt \\
\hfuzz=500pt\includegraphics[width=1em]{element.pdf}~isRegularExpression & \hfuzz=500pt boolean & \hfuzz=500pt pattern is  a regular expression\\
\hfuzz=500pt\includegraphics[width=1em]{element.pdf}~caseSensitive & \hfuzz=500pt boolean & \hfuzz=500pt treat lower and upper case as distinct\\
\hline
\end{tabularx}


\subsection{StringMatchPattern}
Determines if a pattern or a regular expression matches the entire string.


\keepXColumns
\begin{tabularx}{\textwidth}{N T A}
\hline
Name & Type & Annotation\\
\hline
\hfuzz=500pt\includegraphics[width=1em]{element.pdf}~string & \hfuzz=500pt string & \hfuzz=500pt should contain a \{variable\}\\
\hfuzz=500pt\includegraphics[width=1em]{element.pdf}~pattern & \hfuzz=500pt string & \hfuzz=500pt \\
\hfuzz=500pt\includegraphics[width=1em]{element.pdf}~isRegularExpression & \hfuzz=500pt boolean & \hfuzz=500pt pattern is  a regular expression\\
\hfuzz=500pt\includegraphics[width=1em]{element.pdf}~caseSensitive & \hfuzz=500pt boolean & \hfuzz=500pt treat lower and upper case as distinct\\
\hline
\end{tabularx}


\subsection{And}
All conditions must be met (with short-circuit evaluation).


\keepXColumns
\begin{tabularx}{\textwidth}{N T A}
\hline
Name & Type & Annotation\\
\hline
\hfuzz=500pt\includegraphics[width=1em]{element-mustset-unbounded.pdf}~condition & \hfuzz=500pt \hyperref[conditionType]{condition} & \hfuzz=500pt \\
\hline
\end{tabularx}


\subsection{Or}
One of the conditions must be met (with short-circuit evaluation).


\keepXColumns
\begin{tabularx}{\textwidth}{N T A}
\hline
Name & Type & Annotation\\
\hline
\hfuzz=500pt\includegraphics[width=1em]{element-mustset-unbounded.pdf}~condition & \hfuzz=500pt \hyperref[conditionType]{condition} & \hfuzz=500pt \\
\hline
\end{tabularx}


\subsection{Not}
The result of the condition is inverted.


\keepXColumns
\begin{tabularx}{\textwidth}{N T A}
\hline
Name & Type & Annotation\\
\hline
\hfuzz=500pt\includegraphics[width=1em]{element-mustset.pdf}~condition & \hfuzz=500pt \hyperref[conditionType]{condition} & \hfuzz=500pt \\
\hline
\end{tabularx}

\clearpage
%==================================

\section{CovariancePod}\label{covariancePodType}
Provides arc wise covariance matrices for precise orbit data.
Temporal correlations are modeled in the orbit system (along, cross, radial).
The \configFile{inputfileCovarianceFunction}{matrix} provides temporal covariance functions for each axis.
From the diagonal matrix for each time step
\begin{equation}
  Cov_{3\times3}(t) = \text{diag}(cov_x(t), cov_y(t), cov_z(t))
\end{equation}
the Toeplitz covariance matrix for an arc is constructed
\begin{equation}
  \M C = \begin{pmatrix}
    Cov(t_0) & Cov(t_1) & \cdots   &          &        &        \\
    Cov(t_1) & Cov(t_0) & Cov(t_1) & \cdots   &        &        \\
    \cdots   & Cov(t_1) & Cov(t_0) & Cov(t_1) & \cdots &        \\
             & \cdots   & \ddots   & \ddots   & \ddots & \cdots \\
  \end{pmatrix}
\end{equation}

The epoch wise $3\times3$ covariance matrices given by \configFile{inputfileCovariancePodEpoch}{instrument}
are eigen value decomposed
\begin{equation}
  \M C_{3\times3}(t_i) = \M Q \M\Lambda \M Q^T,
\end{equation}
where $\M Q$ is an orthgonal matrix and $\M\Lambda$ diagonal.
This used to split the covariances matrices
\begin{equation}
  \M C_{3\times3}(t_i) = \M D(t_i) \M D(t_i)^T = (\M Q \M\Lambda^{1/2} \M Q^T)(\M Q \M\Lambda^{1/2} \M Q^T)^T,
\end{equation}
and to compose a block diagonal matrix for an arc
\begin{equation}
  \M D = \text{diag}(\M D(t_1), \M D(t_2), \ldots, \M D(t_2)).
\end{equation}

The complete covariance matrix of an arc is given by
\begin{equation}
  \M C_{arc} = \sigma_0^2 \sigma_{arc}^2 \M D \M C \M D^T +
  \text{diag}(\sigma_1^2\M I_{3\times3}, \sigma_2^2\M I_{3\times3}, \ldots, \sigma_n^2\M I_{3\times3})
\end{equation}
where \config{sigma}~$\sigma_0$ is an overall factor
and the arc specific factors $\sigma_{arc}$ can be provided with \configFile{inputfileSigmasPerArc}{matrix}.
The last matrix can be used to downweight outliers in single epochs and will be added if
\configFile{inputfileSigmasPerEpoch}{instrument} is provided.


\keepXColumns
\begin{tabularx}{\textwidth}{N T A}
\hline
Name & Type & Annotation\\
\hline
\hfuzz=500pt\includegraphics[width=1em]{element-mustset.pdf}~covariancePodType & \hfuzz=500pt sequence & \hfuzz=500pt \\
\hfuzz=500pt\includegraphics[width=1em]{connector.pdf}\includegraphics[width=1em]{element.pdf}~sigma & \hfuzz=500pt double & \hfuzz=500pt general variance factor\\
\hfuzz=500pt\includegraphics[width=1em]{connector.pdf}\includegraphics[width=1em]{element.pdf}~inputfileSigmasPerArc & \hfuzz=500pt filename & \hfuzz=500pt different accuaries for each arc (multplicated with sigma)\\
\hfuzz=500pt\includegraphics[width=1em]{connector.pdf}\includegraphics[width=1em]{element.pdf}~inputfileSigmasPerEpoch & \hfuzz=500pt filename & \hfuzz=500pt different accuaries for each epoch (added)\\
\hfuzz=500pt\includegraphics[width=1em]{connector.pdf}\includegraphics[width=1em]{element.pdf}~inputfileCovarianceFunction & \hfuzz=500pt filename & \hfuzz=500pt covariances in time for along, cross, and radial direction\\
\hfuzz=500pt\includegraphics[width=1em]{connector.pdf}\includegraphics[width=1em]{element.pdf}~inputfileCovariancePodEpoch & \hfuzz=500pt filename & \hfuzz=500pt 3x3 epoch wise covariances\\
\hline
\end{tabularx}

\clearpage
%==================================

\section{CovarianceSst}\label{covarianceSstType}
Provides arc wise covariance matrices for satellite-to-satellite observations SST).
The \configFile{inputfileCovarianceFunction}{matrix} provides a temporal covariance function.
From it the Toeplitz covariance matrix is constructed
\begin{equation}
  \M C = \begin{pmatrix}
    cov(t_0) & cov(t_1) & \cdots   &          &        &        \\
    cov(t_1) & cov(t_0) & cov(t_1) & \cdots   &        &        \\
    \cdots   & cov(t_1) & cov(t_0) & cov(t_1) & \cdots &        \\
             & \cdots   & \ddots   & \ddots   & \ddots & \cdots \\
  \end{pmatrix} \\
\end{equation}

The complete covariance matrix of an arc is given by
\begin{equation}
  \M C_{arc} = \sigma_0^2 \sigma_{arc}^2 \M C + \sigma_{S,arc}^2 \M S_{arc}+ \text{diag}(\sigma_1^2, \sigma_2^2, \ldots, \sigma_n^2)
\end{equation}
where \config{sigma}~$\sigma_0$ is an overall factor and the arc specific factors $\sigma_{arc}$
can be provided with \configFile{inputfileSigmasPerArc}{matrix}.
The second term describes general covariance matrices for each arc
\configFile{inputfileCovarianceMatrixArc}{matrix} together with the factors $\sigma_{S,arc}$ from \config{sigmasCovarianceMatrixArc}.
The last matrix can be used to downweight outliers in single epochs and will be added if
\configFile{inputfileSigmasPerEpoch}{instrument} is provided.



\keepXColumns
\begin{tabularx}{\textwidth}{N T A}
\hline
Name & Type & Annotation\\
\hline
\hfuzz=500pt\includegraphics[width=1em]{element-mustset.pdf}~covarianceSstType & \hfuzz=500pt sequence & \hfuzz=500pt \\
\hfuzz=500pt\includegraphics[width=1em]{connector.pdf}\includegraphics[width=1em]{element.pdf}~sigma & \hfuzz=500pt double & \hfuzz=500pt general variance factor\\
\hfuzz=500pt\includegraphics[width=1em]{connector.pdf}\includegraphics[width=1em]{element.pdf}~inputfileSigmasPerArc & \hfuzz=500pt filename & \hfuzz=500pt different accuaries for each arc (multplicated with sigma)\\
\hfuzz=500pt\includegraphics[width=1em]{connector.pdf}\includegraphics[width=1em]{element.pdf}~inputfileSigmasPerEpoch & \hfuzz=500pt filename & \hfuzz=500pt different accuaries for each epoch (added)\\
\hfuzz=500pt\includegraphics[width=1em]{connector.pdf}\includegraphics[width=1em]{element.pdf}~inputfileCovarianceFunction & \hfuzz=500pt filename & \hfuzz=500pt covariance function in time\\
\hfuzz=500pt\includegraphics[width=1em]{connector.pdf}\includegraphics[width=1em]{element-unbounded.pdf}~inputfileCovarianceMatrixArc & \hfuzz=500pt filename & \hfuzz=500pt one matrix file per arc. Use \{arcNo\} as template\\
\hfuzz=500pt\includegraphics[width=1em]{connector.pdf}\includegraphics[width=1em]{element.pdf}~sigmasCovarianceMatrixArc & \hfuzz=500pt filename & \hfuzz=500pt vector with one sigma for each covarianceMatrixArc\\
\hline
\end{tabularx}

\clearpage
%==================================

\section{DigitalFilter}\label{digitalFilterType}
Digital filter implementation for the filtering of equally spaced time series. This class implements the filter equations as
\begin{equation}\label{digitalFilterType:arma}
  \sum_{l=0}^Q a_l y_{n-l} = \sum_{k=-p_0}^{P-p_0-1} b_k x_{n-k}, \hspace{25pt} a_0 = 1,
\end{equation}
where $Q$ is the autoregressive (AR) order and $P$ is the moving average (MA) order. Note that the MA part can also be non-causal.
The characteristics of a filter cascade can be computed by the programs \program{DigitalFilter2FrequencyResponse} and \program{DigitalFilter2ImpulseResponse}.
To apply a filter cascade to a time series (or an instrument file ) use \program{InstrumentFilter}.
Each filter can be applyed in forward and backward direction by setting \config{backwardDirection}.
If the same filter is applied in both directions, the combined filter has zero phase and the squared magnitude response.
Setting \config{inFrequencyDomain} to true applies the transfer function of the filter to the DFT of the input and synthesizes the result, i.e.:
\begin{equation}
  y_n = \mathcal{F}^{-1}\{H\cdot\mathcal{F}\{x_n\}\}.
\end{equation}
This is equivalent to setting \config{padType} to \config{periodic}.

To reduce warmup effects, the input time series can be padded by chosing a \config{padType}:
\begin{itemize}
\item \config{none}: no padding is applied
\item \config{zero}: zeros are appended at the beginning and end of the input time series
\item \config{constant}: the beginning of the input time series is padded with the first value, the end is padded with the last value
\item \config{periodic}: periodic continuation of the input time series (i.,e. the beginning is padded with the last epochs and the end is padded with the first epochs)
\item \config{symmetric}: beginning and end are reflected around the first and last epoch respectively
\end{itemize}


\subsection{MovingAverage}
Moving average (boxcar) filter. For odd lengths, this filter is symmetric and has therefore no phase shift. For even lengths, a phase shift of half a cycle is introduced.

\[
  y_n = \sum_{k=-\lfloor\frac{P}{2}\rfloor}^{\lfloor\frac{P}{2}\rfloor} \frac{1}{P}x_{n-k}
\]



\keepXColumns
\begin{tabularx}{\textwidth}{N T A}
\hline
Name & Type & Annotation\\
\hline
\hfuzz=500pt\includegraphics[width=1em]{element-mustset.pdf}~length & \hfuzz=500pt uint & \hfuzz=500pt number of epochs in averaging operator\\
\hfuzz=500pt\includegraphics[width=1em]{element.pdf}~inFrequencyDomain & \hfuzz=500pt boolean & \hfuzz=500pt apply filter in frequency domain\\
\hfuzz=500pt\includegraphics[width=1em]{element-mustset.pdf}~padType & \hfuzz=500pt choice & \hfuzz=500pt \\
\hfuzz=500pt\includegraphics[width=1em]{connector.pdf}\includegraphics[width=1em]{element-mustset.pdf}~none & \hfuzz=500pt  & \hfuzz=500pt no padding is applied\\
\hfuzz=500pt\includegraphics[width=1em]{connector.pdf}\includegraphics[width=1em]{element-mustset.pdf}~zero & \hfuzz=500pt  & \hfuzz=500pt zero padding\\
\hfuzz=500pt\includegraphics[width=1em]{connector.pdf}\includegraphics[width=1em]{element-mustset.pdf}~constant & \hfuzz=500pt  & \hfuzz=500pt pad using first and last value\\
\hfuzz=500pt\includegraphics[width=1em]{connector.pdf}\includegraphics[width=1em]{element-mustset.pdf}~periodic & \hfuzz=500pt  & \hfuzz=500pt periodic continuation of matrix\\
\hfuzz=500pt\includegraphics[width=1em]{connector.pdf}\includegraphics[width=1em]{element-mustset.pdf}~symmetric & \hfuzz=500pt  & \hfuzz=500pt symmetric continuation around the matrix edges\\
\hline
\end{tabularx}


\subsection{Median}
Moving median filter of length $n$. The filter output at epoch $k$ is the median of the set start at $k-n/2$ to $k+n/2$.
The filter length $n$ should be uneven to avoid a phase shift.


\keepXColumns
\begin{tabularx}{\textwidth}{N T A}
\hline
Name & Type & Annotation\\
\hline
\hfuzz=500pt\includegraphics[width=1em]{element-mustset.pdf}~length & \hfuzz=500pt uint & \hfuzz=500pt length of the moving window [epochs]\\
\hfuzz=500pt\includegraphics[width=1em]{element-mustset.pdf}~padType & \hfuzz=500pt choice & \hfuzz=500pt \\
\hfuzz=500pt\includegraphics[width=1em]{connector.pdf}\includegraphics[width=1em]{element-mustset.pdf}~none & \hfuzz=500pt  & \hfuzz=500pt no padding is applied\\
\hfuzz=500pt\includegraphics[width=1em]{connector.pdf}\includegraphics[width=1em]{element-mustset.pdf}~zero & \hfuzz=500pt  & \hfuzz=500pt zero padding\\
\hfuzz=500pt\includegraphics[width=1em]{connector.pdf}\includegraphics[width=1em]{element-mustset.pdf}~constant & \hfuzz=500pt  & \hfuzz=500pt pad using first and last value\\
\hfuzz=500pt\includegraphics[width=1em]{connector.pdf}\includegraphics[width=1em]{element-mustset.pdf}~periodic & \hfuzz=500pt  & \hfuzz=500pt periodic continuation of matrix\\
\hfuzz=500pt\includegraphics[width=1em]{connector.pdf}\includegraphics[width=1em]{element-mustset.pdf}~symmetric & \hfuzz=500pt  & \hfuzz=500pt symmetric continuation around the matrix edges\\
\hline
\end{tabularx}


\subsection{Derivative}
Symmetric MA filter for numerical differentiation using polynomial approximation. The input time series is approximated by a moving polynomial of degree \config{polynomialDegree}, by solving
\begin{equation}
  \begin{bmatrix} x(t_k+\tau_0) \\ \vdots \\ x(t_k+\tau_M) \end{bmatrix}
  =
  \begin{bmatrix}
  1      & \tau_0 & \tau_0^2 & \cdots & \tau_0^M \\
  \vdots & \vdots & \vdots   &        & \vdots   \\
  1      & \tau_M & \tau_M^2 & \cdots & \tau_M^M \\
  \end{bmatrix}%^{-1}
  \begin{bmatrix}
  a_0 \\ \vdots \\ a_M
  \end{bmatrix}
  \qquad\text{with}\quad
  \tau_j =  (j-M/2)\cdot \Delta t,
\end{equation}
for each time step $t_k$ ($\Delta t$ is the \config{sampling} of the time series).
The filter coefficients for the $k$-th derivative are obtained by taking the appropriate row of the inverse coefficient matrix $\mathbf{W}$:
\begin{equation}
  b_n = \prod_{i=0}^{k-1} (k-i) \mathbf{w}_{2,:}.
\end{equation}
The \config{polynomialDegree} shoud be even if no phase shift should be introduced.


\keepXColumns
\begin{tabularx}{\textwidth}{N T A}
\hline
Name & Type & Annotation\\
\hline
\hfuzz=500pt\includegraphics[width=1em]{element-mustset.pdf}~polynomialDegree & \hfuzz=500pt uint & \hfuzz=500pt degree of approximation polynomial\\
\hfuzz=500pt\includegraphics[width=1em]{element.pdf}~derivative & \hfuzz=500pt uint & \hfuzz=500pt take kth derivative\\
\hfuzz=500pt\includegraphics[width=1em]{element.pdf}~sampling & \hfuzz=500pt double & \hfuzz=500pt assumed time step between points\\
\hfuzz=500pt\includegraphics[width=1em]{element-mustset.pdf}~padType & \hfuzz=500pt choice & \hfuzz=500pt \\
\hfuzz=500pt\includegraphics[width=1em]{connector.pdf}\includegraphics[width=1em]{element-mustset.pdf}~none & \hfuzz=500pt  & \hfuzz=500pt no padding is applied\\
\hfuzz=500pt\includegraphics[width=1em]{connector.pdf}\includegraphics[width=1em]{element-mustset.pdf}~zero & \hfuzz=500pt  & \hfuzz=500pt zero padding\\
\hfuzz=500pt\includegraphics[width=1em]{connector.pdf}\includegraphics[width=1em]{element-mustset.pdf}~constant & \hfuzz=500pt  & \hfuzz=500pt pad using first and last value\\
\hfuzz=500pt\includegraphics[width=1em]{connector.pdf}\includegraphics[width=1em]{element-mustset.pdf}~periodic & \hfuzz=500pt  & \hfuzz=500pt periodic continuation of matrix\\
\hfuzz=500pt\includegraphics[width=1em]{connector.pdf}\includegraphics[width=1em]{element-mustset.pdf}~symmetric & \hfuzz=500pt  & \hfuzz=500pt symmetric continuation around the matrix edges\\
\hline
\end{tabularx}


\subsection{Integral}
Numerical integration using polynomial approximation.
The input time series is approximated by a moving polynomial of degree \config{polynomialDegree}
by solving
\begin{equation}
  \begin{bmatrix} x(t_k+\tau_0) \\ \vdots \\ x(t_k+\tau_M) \end{bmatrix}
  =
  \begin{bmatrix}
  1      & \tau_0 & \tau_0^2 & \cdots & \tau_0^M \\
  \vdots & \vdots & \vdots   &        & \vdots   \\
  1      & \tau_M & \tau_M^2 & \cdots & \tau_M^M \\
  \end{bmatrix}%^{-1}
  \begin{bmatrix}
  a_0 \\ \vdots \\ a_M
  \end{bmatrix}
  \qquad\text{with}\quad
  \tau_j =  (j-M/2)\cdot \Delta t,
\end{equation}
for each time step $t_k$ ($\Delta t$ is the \config{sampling} of the time series).
The numerical integral for each time step $t_k$ is approximated by the center interval of the estimated polynomial.

\fig{!hb}{0.7}{DigitalFilter_integral}{fig:DigitalFilterIntegral}{Numerical integration by polynomial approximation.}

\config{polynomialDegree} should be even to avoid a phase shift.


\keepXColumns
\begin{tabularx}{\textwidth}{N T A}
\hline
Name & Type & Annotation\\
\hline
\hfuzz=500pt\includegraphics[width=1em]{element-mustset.pdf}~polynomialDegree & \hfuzz=500pt uint & \hfuzz=500pt degree of approximation polynomial\\
\hfuzz=500pt\includegraphics[width=1em]{element.pdf}~sampling & \hfuzz=500pt double & \hfuzz=500pt assumed time step between points\\
\hfuzz=500pt\includegraphics[width=1em]{element-mustset.pdf}~padType & \hfuzz=500pt choice & \hfuzz=500pt \\
\hfuzz=500pt\includegraphics[width=1em]{connector.pdf}\includegraphics[width=1em]{element-mustset.pdf}~none & \hfuzz=500pt  & \hfuzz=500pt no padding is applied\\
\hfuzz=500pt\includegraphics[width=1em]{connector.pdf}\includegraphics[width=1em]{element-mustset.pdf}~zero & \hfuzz=500pt  & \hfuzz=500pt zero padding\\
\hfuzz=500pt\includegraphics[width=1em]{connector.pdf}\includegraphics[width=1em]{element-mustset.pdf}~constant & \hfuzz=500pt  & \hfuzz=500pt pad using first and last value\\
\hfuzz=500pt\includegraphics[width=1em]{connector.pdf}\includegraphics[width=1em]{element-mustset.pdf}~periodic & \hfuzz=500pt  & \hfuzz=500pt periodic continuation of matrix\\
\hfuzz=500pt\includegraphics[width=1em]{connector.pdf}\includegraphics[width=1em]{element-mustset.pdf}~symmetric & \hfuzz=500pt  & \hfuzz=500pt symmetric continuation around the matrix edges\\
\hline
\end{tabularx}


\subsection{Correlation}
Correlation ($\rho$) of \config{corr} is introduced into the time series:
\begin{equation}
  y_n = \rho\cdot y_{n-1} + \sqrt{1-\rho^2}x_n.
\end{equation}


\keepXColumns
\begin{tabularx}{\textwidth}{N T A}
\hline
Name & Type & Annotation\\
\hline
\hfuzz=500pt\includegraphics[width=1em]{element-mustset.pdf}~correlation & \hfuzz=500pt double & \hfuzz=500pt correlation\\
\hfuzz=500pt\includegraphics[width=1em]{element.pdf}~backwardDirection & \hfuzz=500pt boolean & \hfuzz=500pt apply filter in backward direction\\
\hfuzz=500pt\includegraphics[width=1em]{element.pdf}~inFrequencyDomain & \hfuzz=500pt boolean & \hfuzz=500pt apply filter in frequency domain\\
\hfuzz=500pt\includegraphics[width=1em]{element-mustset.pdf}~padType & \hfuzz=500pt choice & \hfuzz=500pt \\
\hfuzz=500pt\includegraphics[width=1em]{connector.pdf}\includegraphics[width=1em]{element-mustset.pdf}~none & \hfuzz=500pt  & \hfuzz=500pt no padding is applied\\
\hfuzz=500pt\includegraphics[width=1em]{connector.pdf}\includegraphics[width=1em]{element-mustset.pdf}~zero & \hfuzz=500pt  & \hfuzz=500pt zero padding\\
\hfuzz=500pt\includegraphics[width=1em]{connector.pdf}\includegraphics[width=1em]{element-mustset.pdf}~constant & \hfuzz=500pt  & \hfuzz=500pt pad using first and last value\\
\hfuzz=500pt\includegraphics[width=1em]{connector.pdf}\includegraphics[width=1em]{element-mustset.pdf}~periodic & \hfuzz=500pt  & \hfuzz=500pt periodic continuation of matrix\\
\hfuzz=500pt\includegraphics[width=1em]{connector.pdf}\includegraphics[width=1em]{element-mustset.pdf}~symmetric & \hfuzz=500pt  & \hfuzz=500pt symmetric continuation around the matrix edges\\
\hline
\end{tabularx}


\subsection{GraceLowpass}
Low pass and differentation filter as used for GRACE KBR and ACC data in the Level1A processing.

\fig{!hb}{0.8}{DigitalFilter_graceLowpass}{fig:DigitalFilterGraceLowpass}{Amplitude response of the low pass filter used in the L1A processing.}


\keepXColumns
\begin{tabularx}{\textwidth}{N T A}
\hline
Name & Type & Annotation\\
\hline
\hfuzz=500pt\includegraphics[width=1em]{element-mustset.pdf}~rawDataRate & \hfuzz=500pt double & \hfuzz=500pt sampling frequency in Hz (fs).\\
\hfuzz=500pt\includegraphics[width=1em]{element.pdf}~convolutionNumber & \hfuzz=500pt uint & \hfuzz=500pt number of self convolutions of the filter kernel\\
\hfuzz=500pt\includegraphics[width=1em]{element.pdf}~fitInterval & \hfuzz=500pt double & \hfuzz=500pt length of the filter kernel [seconds]\\
\hfuzz=500pt\includegraphics[width=1em]{element.pdf}~lowPassBandwith & \hfuzz=500pt double & \hfuzz=500pt target low pass bandwidth\\
\hfuzz=500pt\includegraphics[width=1em]{element.pdf}~normFrequency & \hfuzz=500pt double & \hfuzz=500pt norm filter at this frequency [Hz] (default: GRACE dominant (J2) signal frequency)\\
\hfuzz=500pt\includegraphics[width=1em]{element.pdf}~reduceQuadraticFit & \hfuzz=500pt boolean & \hfuzz=500pt remove-\$>\$filter-\$>\$restore quadratic fit\\
\hfuzz=500pt\includegraphics[width=1em]{element.pdf}~derivative & \hfuzz=500pt choice & \hfuzz=500pt \\
\hfuzz=500pt\includegraphics[width=1em]{connector.pdf}\includegraphics[width=1em]{element-mustset.pdf}~derivative1st & \hfuzz=500pt  & \hfuzz=500pt range rate\\
\hfuzz=500pt\includegraphics[width=1em]{connector.pdf}\includegraphics[width=1em]{element-mustset.pdf}~derivative2nd & \hfuzz=500pt  & \hfuzz=500pt range acceleration\\
\hfuzz=500pt\includegraphics[width=1em]{element.pdf}~inFrequencyDomain & \hfuzz=500pt boolean & \hfuzz=500pt apply filter in frequency domain\\
\hfuzz=500pt\includegraphics[width=1em]{element-mustset.pdf}~padType & \hfuzz=500pt choice & \hfuzz=500pt \\
\hfuzz=500pt\includegraphics[width=1em]{connector.pdf}\includegraphics[width=1em]{element-mustset.pdf}~none & \hfuzz=500pt  & \hfuzz=500pt no padding is applied\\
\hfuzz=500pt\includegraphics[width=1em]{connector.pdf}\includegraphics[width=1em]{element-mustset.pdf}~zero & \hfuzz=500pt  & \hfuzz=500pt zero padding\\
\hfuzz=500pt\includegraphics[width=1em]{connector.pdf}\includegraphics[width=1em]{element-mustset.pdf}~constant & \hfuzz=500pt  & \hfuzz=500pt pad using first and last value\\
\hfuzz=500pt\includegraphics[width=1em]{connector.pdf}\includegraphics[width=1em]{element-mustset.pdf}~periodic & \hfuzz=500pt  & \hfuzz=500pt periodic continuation of matrix\\
\hfuzz=500pt\includegraphics[width=1em]{connector.pdf}\includegraphics[width=1em]{element-mustset.pdf}~symmetric & \hfuzz=500pt  & \hfuzz=500pt symmetric continuation around the matrix edges\\
\hline
\end{tabularx}


\subsection{Butterworth}
Digital implementation of the Butterworth filter. The design of the filter is done by modifying the analog (continuous time) transfer function, which is
then transformed into the digital domain by using the bilinear transform. The filter coefficients are then determined by a least squares adjustment in time domain.

The \config{filterType} can be \config{lowpass}, \config{highpass}, where one cutoff frequency has to be specified, and \config{bandpass} and \config{bandstop} where to cutoff frequencies have to be specified.
Cutoff frequencies must be given as normalized frequency $w_n = f/f_{\text{nyq}}$. For a cutoff frequency of 30~mHz for a time series sampled with 5~seconds gives a normalized frequency of $0.03/0.1 = 0.3$.


\keepXColumns
\begin{tabularx}{\textwidth}{N T A}
\hline
Name & Type & Annotation\\
\hline
\hfuzz=500pt\includegraphics[width=1em]{element-mustset.pdf}~order & \hfuzz=500pt uint & \hfuzz=500pt filter order\\
\hfuzz=500pt\includegraphics[width=1em]{element-mustset.pdf}~type & \hfuzz=500pt choice & \hfuzz=500pt filter type\\
\hfuzz=500pt\includegraphics[width=1em]{connector.pdf}\includegraphics[width=1em]{element-mustset.pdf}~lowpass & \hfuzz=500pt sequence & \hfuzz=500pt \\
\hfuzz=500pt\quad\includegraphics[width=1em]{connector.pdf}\includegraphics[width=1em]{element-mustset.pdf}~Wn & \hfuzz=500pt double & \hfuzz=500pt normalized cutoff frequency (f\_c / f\_nyq)\\
\hfuzz=500pt\includegraphics[width=1em]{connector.pdf}\includegraphics[width=1em]{element-mustset.pdf}~highpass & \hfuzz=500pt sequence & \hfuzz=500pt \\
\hfuzz=500pt\quad\includegraphics[width=1em]{connector.pdf}\includegraphics[width=1em]{element-mustset.pdf}~Wn & \hfuzz=500pt double & \hfuzz=500pt normalized cutoff frequency (f\_c / f\_nyq)\\
\hfuzz=500pt\includegraphics[width=1em]{connector.pdf}\includegraphics[width=1em]{element-mustset.pdf}~bandpass & \hfuzz=500pt sequence & \hfuzz=500pt \\
\hfuzz=500pt\quad\includegraphics[width=1em]{connector.pdf}\includegraphics[width=1em]{element-mustset.pdf}~Wn1 & \hfuzz=500pt double & \hfuzz=500pt lower normalized cutoff frequency (f\_c / f\_nyq)\\
\hfuzz=500pt\quad\includegraphics[width=1em]{connector.pdf}\includegraphics[width=1em]{element-mustset.pdf}~Wn2 & \hfuzz=500pt double & \hfuzz=500pt upper normalized cutoff frequency (f\_c / f\_nyq)\\
\hfuzz=500pt\includegraphics[width=1em]{connector.pdf}\includegraphics[width=1em]{element-mustset.pdf}~bandstop & \hfuzz=500pt sequence & \hfuzz=500pt \\
\hfuzz=500pt\quad\includegraphics[width=1em]{connector.pdf}\includegraphics[width=1em]{element-mustset.pdf}~Wn1 & \hfuzz=500pt double & \hfuzz=500pt lower normalized cutoff frequency (f\_c / f\_nyq)\\
\hfuzz=500pt\quad\includegraphics[width=1em]{connector.pdf}\includegraphics[width=1em]{element-mustset.pdf}~Wn2 & \hfuzz=500pt double & \hfuzz=500pt upper normalized cutoff frequency (f\_c / f\_nyq)\\
\hfuzz=500pt\includegraphics[width=1em]{element.pdf}~backwardDirection & \hfuzz=500pt boolean & \hfuzz=500pt apply filter in backward direction\\
\hfuzz=500pt\includegraphics[width=1em]{element.pdf}~inFrequencyDomain & \hfuzz=500pt boolean & \hfuzz=500pt apply filter in frequency domain\\
\hfuzz=500pt\includegraphics[width=1em]{element-mustset.pdf}~padType & \hfuzz=500pt choice & \hfuzz=500pt \\
\hfuzz=500pt\includegraphics[width=1em]{connector.pdf}\includegraphics[width=1em]{element-mustset.pdf}~none & \hfuzz=500pt  & \hfuzz=500pt no padding is applied\\
\hfuzz=500pt\includegraphics[width=1em]{connector.pdf}\includegraphics[width=1em]{element-mustset.pdf}~zero & \hfuzz=500pt  & \hfuzz=500pt zero padding\\
\hfuzz=500pt\includegraphics[width=1em]{connector.pdf}\includegraphics[width=1em]{element-mustset.pdf}~constant & \hfuzz=500pt  & \hfuzz=500pt pad using first and last value\\
\hfuzz=500pt\includegraphics[width=1em]{connector.pdf}\includegraphics[width=1em]{element-mustset.pdf}~periodic & \hfuzz=500pt  & \hfuzz=500pt periodic continuation of matrix\\
\hfuzz=500pt\includegraphics[width=1em]{connector.pdf}\includegraphics[width=1em]{element-mustset.pdf}~symmetric & \hfuzz=500pt  & \hfuzz=500pt symmetric continuation around the matrix edges\\
\hline
\end{tabularx}


\subsection{File}
Read filter coefficients of \eqref{digitalFilterType:arma} from a coefficient file.
One column might define the index $n$
of the coefficients $a_n$ and $b_n$ in the other columns.


\keepXColumns
\begin{tabularx}{\textwidth}{N T A}
\hline
Name & Type & Annotation\\
\hline
\hfuzz=500pt\includegraphics[width=1em]{element-mustset.pdf}~inputfileMatrix & \hfuzz=500pt filename & \hfuzz=500pt matrix with filter coefficients\\
\hfuzz=500pt\includegraphics[width=1em]{element.pdf}~index & \hfuzz=500pt expression & \hfuzz=500pt index of coefficients (input columns are named data0, data1, ...)\\
\hfuzz=500pt\includegraphics[width=1em]{element.pdf}~bn & \hfuzz=500pt expression & \hfuzz=500pt MA coefficients (moving average) (input columns are named data0, data1, ...)\\
\hfuzz=500pt\includegraphics[width=1em]{element.pdf}~an & \hfuzz=500pt expression & \hfuzz=500pt AR coefficients (autoregressive) (input columns are named data0, data1, ...)\\
\hfuzz=500pt\includegraphics[width=1em]{element.pdf}~backwardDirection & \hfuzz=500pt boolean & \hfuzz=500pt apply filter in backward direction\\
\hfuzz=500pt\includegraphics[width=1em]{element.pdf}~inFrequencyDomain & \hfuzz=500pt boolean & \hfuzz=500pt apply filter in frequency domain\\
\hfuzz=500pt\includegraphics[width=1em]{element-mustset.pdf}~padType & \hfuzz=500pt choice & \hfuzz=500pt \\
\hfuzz=500pt\includegraphics[width=1em]{connector.pdf}\includegraphics[width=1em]{element-mustset.pdf}~none & \hfuzz=500pt  & \hfuzz=500pt no padding is applied\\
\hfuzz=500pt\includegraphics[width=1em]{connector.pdf}\includegraphics[width=1em]{element-mustset.pdf}~zero & \hfuzz=500pt  & \hfuzz=500pt zero padding\\
\hfuzz=500pt\includegraphics[width=1em]{connector.pdf}\includegraphics[width=1em]{element-mustset.pdf}~constant & \hfuzz=500pt  & \hfuzz=500pt pad using first and last value\\
\hfuzz=500pt\includegraphics[width=1em]{connector.pdf}\includegraphics[width=1em]{element-mustset.pdf}~periodic & \hfuzz=500pt  & \hfuzz=500pt periodic continuation of matrix\\
\hfuzz=500pt\includegraphics[width=1em]{connector.pdf}\includegraphics[width=1em]{element-mustset.pdf}~symmetric & \hfuzz=500pt  & \hfuzz=500pt symmetric continuation around the matrix edges\\
\hline
\end{tabularx}


\subsection{Wavelet}
Filter representation of a wavelet.


\keepXColumns
\begin{tabularx}{\textwidth}{N T A}
\hline
Name & Type & Annotation\\
\hline
\hfuzz=500pt\includegraphics[width=1em]{element-mustset.pdf}~inputfileWavelet & \hfuzz=500pt filename & \hfuzz=500pt wavelet coefficients\\
\hfuzz=500pt\includegraphics[width=1em]{element-mustset.pdf}~type & \hfuzz=500pt choice & \hfuzz=500pt filter type\\
\hfuzz=500pt\includegraphics[width=1em]{connector.pdf}\includegraphics[width=1em]{element-mustset.pdf}~lowpass & \hfuzz=500pt  & \hfuzz=500pt \\
\hfuzz=500pt\includegraphics[width=1em]{connector.pdf}\includegraphics[width=1em]{element-mustset.pdf}~highpass & \hfuzz=500pt  & \hfuzz=500pt \\
\hfuzz=500pt\includegraphics[width=1em]{element.pdf}~level & \hfuzz=500pt uint & \hfuzz=500pt compute filter for specific decomposition level\\
\hfuzz=500pt\includegraphics[width=1em]{element.pdf}~backwardDirection & \hfuzz=500pt boolean & \hfuzz=500pt apply filter in backward direction\\
\hfuzz=500pt\includegraphics[width=1em]{element.pdf}~inFrequencyDomain & \hfuzz=500pt boolean & \hfuzz=500pt apply filter in frequency domain\\
\hfuzz=500pt\includegraphics[width=1em]{element-mustset.pdf}~padType & \hfuzz=500pt choice & \hfuzz=500pt \\
\hfuzz=500pt\includegraphics[width=1em]{connector.pdf}\includegraphics[width=1em]{element-mustset.pdf}~none & \hfuzz=500pt  & \hfuzz=500pt no padding is applied\\
\hfuzz=500pt\includegraphics[width=1em]{connector.pdf}\includegraphics[width=1em]{element-mustset.pdf}~zero & \hfuzz=500pt  & \hfuzz=500pt zero padding\\
\hfuzz=500pt\includegraphics[width=1em]{connector.pdf}\includegraphics[width=1em]{element-mustset.pdf}~constant & \hfuzz=500pt  & \hfuzz=500pt pad using first and last value\\
\hfuzz=500pt\includegraphics[width=1em]{connector.pdf}\includegraphics[width=1em]{element-mustset.pdf}~periodic & \hfuzz=500pt  & \hfuzz=500pt periodic continuation of matrix\\
\hfuzz=500pt\includegraphics[width=1em]{connector.pdf}\includegraphics[width=1em]{element-mustset.pdf}~symmetric & \hfuzz=500pt  & \hfuzz=500pt symmetric continuation around the matrix edges\\
\hline
\end{tabularx}


\subsection{Notch}
Implemented after Christian Siemes' dissertation, page 106.

\fig{!hb}{0.6}{DigitalFilter_notch}{fig:DigitalFilterNotch}{Amplitude response of a notch filter of order three with default settings.}


\keepXColumns
\begin{tabularx}{\textwidth}{N T A}
\hline
Name & Type & Annotation\\
\hline
\hfuzz=500pt\includegraphics[width=1em]{element-mustset.pdf}~notchFrequency & \hfuzz=500pt double & \hfuzz=500pt normalized notch frequency w\_n = (f\_n/f\_nyq)\\
\hfuzz=500pt\includegraphics[width=1em]{element-mustset.pdf}~bandWidth & \hfuzz=500pt double & \hfuzz=500pt bandwidth at -3db. Quality factor of filter Q = w\_n/bw\\
\hfuzz=500pt\includegraphics[width=1em]{element.pdf}~backwardDirection & \hfuzz=500pt boolean & \hfuzz=500pt apply filter in backward direction\\
\hfuzz=500pt\includegraphics[width=1em]{element.pdf}~inFrequencyDomain & \hfuzz=500pt boolean & \hfuzz=500pt apply filter in frequency domain\\
\hfuzz=500pt\includegraphics[width=1em]{element-mustset.pdf}~padType & \hfuzz=500pt choice & \hfuzz=500pt \\
\hfuzz=500pt\includegraphics[width=1em]{connector.pdf}\includegraphics[width=1em]{element-mustset.pdf}~none & \hfuzz=500pt  & \hfuzz=500pt no padding is applied\\
\hfuzz=500pt\includegraphics[width=1em]{connector.pdf}\includegraphics[width=1em]{element-mustset.pdf}~zero & \hfuzz=500pt  & \hfuzz=500pt zero padding\\
\hfuzz=500pt\includegraphics[width=1em]{connector.pdf}\includegraphics[width=1em]{element-mustset.pdf}~constant & \hfuzz=500pt  & \hfuzz=500pt pad using first and last value\\
\hfuzz=500pt\includegraphics[width=1em]{connector.pdf}\includegraphics[width=1em]{element-mustset.pdf}~periodic & \hfuzz=500pt  & \hfuzz=500pt periodic continuation of matrix\\
\hfuzz=500pt\includegraphics[width=1em]{connector.pdf}\includegraphics[width=1em]{element-mustset.pdf}~symmetric & \hfuzz=500pt  & \hfuzz=500pt symmetric continuation around the matrix edges\\
\hline
\end{tabularx}


\subsection{Decorrelation}
Moving average decorrelation filter based on eigendecomposition of a Toeplitz covariance matrix.


\keepXColumns
\begin{tabularx}{\textwidth}{N T A}
\hline
Name & Type & Annotation\\
\hline
\hfuzz=500pt\includegraphics[width=1em]{element-mustset.pdf}~inputfileCovarianceFunction & \hfuzz=500pt filename & \hfuzz=500pt covariance function of time series\\
\hfuzz=500pt\includegraphics[width=1em]{element.pdf}~inFrequencyDomain & \hfuzz=500pt boolean & \hfuzz=500pt apply filter in frequency domain\\
\hfuzz=500pt\includegraphics[width=1em]{element-mustset.pdf}~padType & \hfuzz=500pt choice & \hfuzz=500pt \\
\hfuzz=500pt\includegraphics[width=1em]{connector.pdf}\includegraphics[width=1em]{element-mustset.pdf}~none & \hfuzz=500pt  & \hfuzz=500pt no padding is applied\\
\hfuzz=500pt\includegraphics[width=1em]{connector.pdf}\includegraphics[width=1em]{element-mustset.pdf}~zero & \hfuzz=500pt  & \hfuzz=500pt zero padding\\
\hfuzz=500pt\includegraphics[width=1em]{connector.pdf}\includegraphics[width=1em]{element-mustset.pdf}~constant & \hfuzz=500pt  & \hfuzz=500pt pad using first and last value\\
\hfuzz=500pt\includegraphics[width=1em]{connector.pdf}\includegraphics[width=1em]{element-mustset.pdf}~periodic & \hfuzz=500pt  & \hfuzz=500pt periodic continuation of matrix\\
\hfuzz=500pt\includegraphics[width=1em]{connector.pdf}\includegraphics[width=1em]{element-mustset.pdf}~symmetric & \hfuzz=500pt  & \hfuzz=500pt symmetric continuation around the matrix edges\\
\hline
\end{tabularx}


\subsection{TimeLag}
Lag operator in digital filter representation.


\keepXColumns
\begin{tabularx}{\textwidth}{N T A}
\hline
Name & Type & Annotation\\
\hline
\hfuzz=500pt\includegraphics[width=1em]{element-mustset.pdf}~lag & \hfuzz=500pt int & \hfuzz=500pt lag epochs: 1 (lag); -1 (lead)\\
\hfuzz=500pt\includegraphics[width=1em]{element.pdf}~inFrequencyDomain & \hfuzz=500pt boolean & \hfuzz=500pt apply filter in frequency domain\\
\hfuzz=500pt\includegraphics[width=1em]{element-mustset.pdf}~padType & \hfuzz=500pt choice & \hfuzz=500pt \\
\hfuzz=500pt\includegraphics[width=1em]{connector.pdf}\includegraphics[width=1em]{element-mustset.pdf}~none & \hfuzz=500pt  & \hfuzz=500pt no padding is applied\\
\hfuzz=500pt\includegraphics[width=1em]{connector.pdf}\includegraphics[width=1em]{element-mustset.pdf}~zero & \hfuzz=500pt  & \hfuzz=500pt zero padding\\
\hfuzz=500pt\includegraphics[width=1em]{connector.pdf}\includegraphics[width=1em]{element-mustset.pdf}~constant & \hfuzz=500pt  & \hfuzz=500pt pad using first and last value\\
\hfuzz=500pt\includegraphics[width=1em]{connector.pdf}\includegraphics[width=1em]{element-mustset.pdf}~periodic & \hfuzz=500pt  & \hfuzz=500pt periodic continuation of matrix\\
\hfuzz=500pt\includegraphics[width=1em]{connector.pdf}\includegraphics[width=1em]{element-mustset.pdf}~symmetric & \hfuzz=500pt  & \hfuzz=500pt symmetric continuation around the matrix edges\\
\hline
\end{tabularx}


\subsection{ReduceFilterOutput}
Removes the filtered signal from the input, i.e. the input is passed
through a \configClass{digitalFilter}{digitalFilterType} with a frequency response of $1-H(f)$.


\keepXColumns
\begin{tabularx}{\textwidth}{N T A}
\hline
Name & Type & Annotation\\
\hline
\hfuzz=500pt\includegraphics[width=1em]{element-mustset-unbounded.pdf}~filter & \hfuzz=500pt \hyperref[digitalFilterType]{digitalFilter} & \hfuzz=500pt remove filter output from input signal\\
\hline
\end{tabularx}

\clearpage
%==================================

\section{Doodson}\label{doodson}
This is a string which describes a tidal frequency either coded as Doodson number
or using Darwin´s name, e.g. \verb|255.555| or \verb|M2|.

The following names are defined:
\begin{itemize}
\item \verb|055.565|: \verb|om1|  \item \verb|055.575|: \verb|om2|     \item \verb|056.554|: \verb|sa|
\item \verb|056.555|: \verb|sa|   \item \verb|057.555|: \verb|ssa|     \item \verb|058.554|: \verb|sta|
\item \verb|063.655|: \verb|msm|  \item \verb|065.455|: \verb|mm|      \item \verb|073.555|: \verb|msf|
\item \verb|075.555|: \verb|mf|   \item \verb|083.655|: \verb|mstm|    \item \verb|085.455|: \verb|mtm|
\item \verb|093.555|: \verb|msq|  \item \verb|093.555|: \verb|msqm|    \item \verb|125.755|: \verb|2q1|
\item \verb|127.555|: \verb|sig1| \item \verb|127.555|: \verb|sigma1|  \item \verb|135.655|: \verb|q1|
\item \verb|137.455|: \verb|ro1|  \item \verb|137.455|: \verb|rho1|    \item \verb|145.555|: \verb|o1|
\item \verb|147.555|: \verb|tau1| \item \verb|155.655|: \verb|m1|      \item \verb|157.455|: \verb|chi1|
\item \verb|162.556|: \verb|pi1|  \item \verb|163.555|: \verb|p1|      \item \verb|164.555|: \verb|s1|
\item \verb|165.555|: \verb|k1|   \item \verb|166.554|: \verb|psi1|    \item \verb|167.555|: \verb|fi1|
\item \verb|167.555|: \verb|phi1| \item \verb|173.655|: \verb|the1|    \item \verb|173.655|: \verb|theta1|
\item \verb|175.455|: \verb|j1|   \item \verb|183.555|: \verb|so1|     \item \verb|185.555|: \verb|oo1|
\item \verb|195.455|: \verb|v1|   \item \verb|225.855|: \verb|3n2|     \item \verb|227.655|: \verb|eps2|
\item \verb|235.755|: \verb|2n2|  \item \verb|237.555|: \verb|mu2|     \item \verb|237.555|: \verb|mi2|
\item \verb|245.655|: \verb|n2|   \item \verb|247.455|: \verb|nu2|     \item \verb|247.455|: \verb|ni2|
\item \verb|253.755|: \verb|gam2| \item \verb|254.556|: \verb|alf2|    \item \verb|255.555|: \verb|m2|
\item \verb|256.554|: \verb|bet2| \item \verb|257.555|: \verb|dlt2|    \item \verb|263.655|: \verb|la2|
\item \verb|263.655|: \verb|lmb2| \item \verb|263.655|: \verb|lambda2| \item \verb|265.455|: \verb|l2|
\item \verb|271.557|: \verb|2t2|  \item \verb|272.556|: \verb|t2|      \item \verb|273.555|: \verb|s2|
\item \verb|274.554|: \verb|r2|   \item \verb|275.555|: \verb|k2|      \item \verb|283.655|: \verb|ksi2|
\item \verb|285.455|: \verb|eta2| \item \verb|355.555|: \verb|m3|      \item \verb|381.555|: \verb|t3|
\item \verb|382.555|: \verb|s3|   \item \verb|383.555|: \verb|r3|      \item \verb|435.755|: \verb|n4|
\item \verb|445.655|: \verb|mn4|  \item \verb|455.555|: \verb|m4|      \item \verb|473.555|: \verb|ms4|
\item \verb|491.555|: \verb|s4|   \item \verb|655.555|: \verb|m6|      \item \verb|855.555|: \verb|m8|
\end{itemize}

\clearpage
%==================================

\section{EarthRotation}\label{earthRotationType}
This class realize the transformation between a terestrial
reference frame (TRF) and a celestial reference frame (CRF).


\subsection{File}\label{earthRotationType:file}
This class realize the transformation by interpolation from file.
This file can be created with \program{EarthOrientationParameterTimeSeries}.


\keepXColumns
\begin{tabularx}{\textwidth}{N T A}
\hline
Name & Type & Annotation\\
\hline
\hfuzz=500pt\includegraphics[width=1em]{element-mustset.pdf}~inputfileEOP & \hfuzz=500pt filename & \hfuzz=500pt \\
\hfuzz=500pt\includegraphics[width=1em]{element.pdf}~interpolationDegree & \hfuzz=500pt uint & \hfuzz=500pt for polynomial interpolation\\
\hline
\end{tabularx}


\subsection{Iers2010}
This class realize the transformation according to the IERS2010 conventions
given by the \emph{International Earth Rotation and Reference Systems Service} (IERS).
A file with the earth orientation parameter is needed (\configFile{inputfileEOP}{earthOrientationParameter}).


\keepXColumns
\begin{tabularx}{\textwidth}{N T A}
\hline
Name & Type & Annotation\\
\hline
\hfuzz=500pt\includegraphics[width=1em]{element.pdf}~inputfileEOP & \hfuzz=500pt filename & \hfuzz=500pt \\
\hfuzz=500pt\includegraphics[width=1em]{element.pdf}~truncatedNutation & \hfuzz=500pt boolean & \hfuzz=500pt use truncated nutation model (IAU2006B)\\
\hline
\end{tabularx}


\subsection{Iers2010b}\label{earthRotationType:iers2010b}
This class realize the transformation according to the IERS2010 conventions
given by the \emph{International Earth Rotation and Reference Systems Service} (IERS).
A file with the earth orientation parameter is needed (\configFile{inputfileEOP}{earthOrientationParameter}).
Includes additional high-frequency EOP models (\configFile{inputfileDoodsonEOP}{doodsonEarthOrientationParameter}).


\keepXColumns
\begin{tabularx}{\textwidth}{N T A}
\hline
Name & Type & Annotation\\
\hline
\hfuzz=500pt\includegraphics[width=1em]{element.pdf}~inputfileEOP & \hfuzz=500pt filename & \hfuzz=500pt \\
\hfuzz=500pt\includegraphics[width=1em]{element.pdf}~inputfileDoodsonEOP & \hfuzz=500pt filename & \hfuzz=500pt \\
\hline
\end{tabularx}


\subsection{Iers2003}
This class realize the transformation according to IERS2003 conventions
given by the \emph{International Earth Rotation and Reference Systems Service} (IERS).
A file with the earth orientation parameter is needed (\configFile{inputfileEOP}{earthOrientationParameter}).

The following subroutines are used:
\begin{itemize}
\item BPN2000.f,
\item ERA2000.f,
\item pmsdnut.f,
\item POM2000.f,
\item SP2000.f,
\item T2C2000.f,
\item XYS2000A.f
\end{itemize}
from \url{ftp://maia.usno.navy.mil/conv2000/chapter5/} and
\begin{itemize}
\item orthoeop.f
\end{itemize}
from \url{ftp://maia.usno.navy.mil/conv2000/chapter8/}


\keepXColumns
\begin{tabularx}{\textwidth}{N T A}
\hline
Name & Type & Annotation\\
\hline
\hfuzz=500pt\includegraphics[width=1em]{element-mustset.pdf}~inputfileEOP & \hfuzz=500pt filename & \hfuzz=500pt \\
\hline
\end{tabularx}


\subsection{Iers1996}
Very old.


\keepXColumns
\begin{tabularx}{\textwidth}{N T A}
\hline
Name & Type & Annotation\\
\hline
\hfuzz=500pt\includegraphics[width=1em]{element.pdf}~inputfileEOP & \hfuzz=500pt filename & \hfuzz=500pt \\
\hfuzz=500pt\includegraphics[width=1em]{element-mustset.pdf}~inputfileNutation & \hfuzz=500pt filename & \hfuzz=500pt \\
\hline
\end{tabularx}


\subsection{Gmst}
The transformation is realized as rotation about the z-axis.
The angle ist given by the Greenwich Mean Siderial Time (GMST).
\begin{verbatim}
  Double Tu0 = (timeUTC.mjdInt()-51544.5)/36525.0;

  Double GMST0 = (6.0/24 + 41.0/(24*60) + 50.54841/(24*60*60))
               + (8640184.812866/(24*60*60))*Tu0
               + (0.093104/(24*60*60))*Tu0*Tu0
               + (-6.2e-6/(24*60*60))*Tu0*Tu0*Tu0;
  Double r = 1.002737909350795 + 5.9006e-11*Tu0 - 5.9e-15*Tu0*Tu0;
  GMST = fmod(2*PI*(GMST0 + r * timeUTC.mjdMod()), 2*PI);
\end{verbatim}


\subsection{Earth Rotation Angle (ERA)}
The transformation is realized as rotation about the z-axis.
The angle ist given by the Earth Rotation Angle (ERA).


\subsection{Z-Axis}
The transformation is realized as rotation about the z-axis.
You must specify the angle (\config{initialAngle}) at \config{time0} and
the angular velocity (\config{angularVelocity}).


\keepXColumns
\begin{tabularx}{\textwidth}{N T A}
\hline
Name & Type & Annotation\\
\hline
\hfuzz=500pt\includegraphics[width=1em]{element-mustset.pdf}~initialAngle & \hfuzz=500pt double & \hfuzz=500pt Angle at time0 [rad]\\
\hfuzz=500pt\includegraphics[width=1em]{element-mustset.pdf}~angularVelocity & \hfuzz=500pt double & \hfuzz=500pt [rad/s]\\
\hfuzz=500pt\includegraphics[width=1em]{element-mustset.pdf}~time0 & \hfuzz=500pt time & \hfuzz=500pt \\
\hline
\end{tabularx}


\subsection{StarCamera}
This class reads quaternions from an instrument file and interpolates to the given time stamp.


\keepXColumns
\begin{tabularx}{\textwidth}{N T A}
\hline
Name & Type & Annotation\\
\hline
\hfuzz=500pt\includegraphics[width=1em]{element-mustset.pdf}~inputfileStarCamera & \hfuzz=500pt filename & \hfuzz=500pt \\
\hfuzz=500pt\includegraphics[width=1em]{element.pdf}~interpolationDegree & \hfuzz=500pt uint & \hfuzz=500pt degree of interpolation polynomial\\
\hline
\end{tabularx}


\subsection{MoonRotation}
This class realizes the transformation between the moon-fixed system
(Principal Axis System (PA) or Mean Earth System (ME))
and the ICRS according to the JPL ephemeris file.


\keepXColumns
\begin{tabularx}{\textwidth}{N T A}
\hline
Name & Type & Annotation\\
\hline
\hfuzz=500pt\includegraphics[width=1em]{element.pdf}~inputfileEphemerides & \hfuzz=500pt filename & \hfuzz=500pt librations\\
\hfuzz=500pt\includegraphics[width=1em]{element-mustset.pdf}~moonfixedSystem & \hfuzz=500pt choice & \hfuzz=500pt \\
\hfuzz=500pt\includegraphics[width=1em]{connector.pdf}\includegraphics[width=1em]{element-mustset.pdf}~PA & \hfuzz=500pt  & \hfuzz=500pt Principal Axis System\\
\hfuzz=500pt\includegraphics[width=1em]{connector.pdf}\includegraphics[width=1em]{element-mustset.pdf}~ME & \hfuzz=500pt  & \hfuzz=500pt Mean Earth System\\
\hline
\end{tabularx}

\clearpage
%==================================

\section{Eclipse}\label{eclipseType}
Shadowing of satellites by moon and Earth provided as factor
between $[0,1]$ with 0: full shadow and 1: full sun light.


\subsection{Conical}
\fig{!hb}{0.8}{eclipseConical}{fig:eclipseConical}{Modelling umbra and penumbra.}


\subsection{SOLAARS}
Earth’s penumbra modeling with Solar radiation pressure with
Oblateness and Lower Atmospheric Absorption, Refraction, and Scattering (SOLAARS).
See Robertson, Robbie. (2015),
Highly Physical Solar Radiation Pressure Modeling During Penumbra Transitions (pp. 67-75).

\clearpage
%==================================

\section{Ephemerides}\label{ephemeridesType}
Ephemerides of Sun, Moon and planets.
The coordinate system is defined as center of \configClass{origin}{planetType}.


\section{JPL}\label{ephemeridesType:jpl}
Using \verb|DExxx| ephemerides from NASA Jet Propulsion Laboratory (JPL).


\keepXColumns
\begin{tabularx}{\textwidth}{N T A}
\hline
Name & Type & Annotation\\
\hline
\hfuzz=500pt\includegraphics[width=1em]{element-mustset.pdf}~inputfileEphemerides & \hfuzz=500pt filename & \hfuzz=500pt \\
\hfuzz=500pt\includegraphics[width=1em]{element.pdf}~origin & \hfuzz=500pt \hyperref[planetType]{planet} & \hfuzz=500pt center of coordinate system\\
\hline
\end{tabularx}

\clearpage
%==================================

\section{Forces}\label{forcesType}
This class provides the forces acting on a satellite.
This encompasses \configClass{gravityfield}{gravityfieldType}, \configClass{tides}{tidesType}
and \configClass{miscAccelerations}{miscAccelerationsType}.


\keepXColumns
\begin{tabularx}{\textwidth}{N T A}
\hline
Name & Type & Annotation\\
\hline
\hfuzz=500pt\includegraphics[width=1em]{element-mustset.pdf}~forcesType & \hfuzz=500pt sequence & \hfuzz=500pt \\
\hfuzz=500pt\includegraphics[width=1em]{connector.pdf}\includegraphics[width=1em]{element-unbounded.pdf}~gravityfield & \hfuzz=500pt \hyperref[gravityfieldType]{gravityfield} & \hfuzz=500pt \\
\hfuzz=500pt\includegraphics[width=1em]{connector.pdf}\includegraphics[width=1em]{element-unbounded.pdf}~tides & \hfuzz=500pt \hyperref[tidesType]{tides} & \hfuzz=500pt \\
\hfuzz=500pt\includegraphics[width=1em]{connector.pdf}\includegraphics[width=1em]{element-unbounded.pdf}~miscAccelerations & \hfuzz=500pt \hyperref[miscAccelerationsType]{miscAccelerations} & \hfuzz=500pt \\
\hline
\end{tabularx}

\clearpage
%==================================

\section{GnssAntennaDefintionList}\label{gnssAntennaDefintionListType}
Provides a list of GnssAntennaDefinitions as used in \program{GnssAntennaDefinitionCreate}.


\subsection{New}
Creates a new antenna.


\keepXColumns
\begin{tabularx}{\textwidth}{N T A}
\hline
Name & Type & Annotation\\
\hline
\hfuzz=500pt\includegraphics[width=1em]{element.pdf}~name & \hfuzz=500pt string & \hfuzz=500pt \\
\hfuzz=500pt\includegraphics[width=1em]{element.pdf}~serial & \hfuzz=500pt string & \hfuzz=500pt \\
\hfuzz=500pt\includegraphics[width=1em]{element.pdf}~radome & \hfuzz=500pt string & \hfuzz=500pt \\
\hfuzz=500pt\includegraphics[width=1em]{element.pdf}~comment & \hfuzz=500pt string & \hfuzz=500pt \\
\hfuzz=500pt\includegraphics[width=1em]{element-mustset-unbounded.pdf}~pattern & \hfuzz=500pt sequence & \hfuzz=500pt \\
\hfuzz=500pt\includegraphics[width=1em]{connector.pdf}\includegraphics[width=1em]{element-mustset.pdf}~type & \hfuzz=500pt \hyperref[gnssType]{gnssType} & \hfuzz=500pt pattern matching of observation types\\
\hfuzz=500pt\includegraphics[width=1em]{connector.pdf}\includegraphics[width=1em]{element.pdf}~offsetX & \hfuzz=500pt double & \hfuzz=500pt [m] antenna center offset\\
\hfuzz=500pt\includegraphics[width=1em]{connector.pdf}\includegraphics[width=1em]{element.pdf}~offsetY & \hfuzz=500pt double & \hfuzz=500pt [m] antenna center offset\\
\hfuzz=500pt\includegraphics[width=1em]{connector.pdf}\includegraphics[width=1em]{element.pdf}~offsetZ & \hfuzz=500pt double & \hfuzz=500pt [m] antenna center offset\\
\hfuzz=500pt\includegraphics[width=1em]{connector.pdf}\includegraphics[width=1em]{element-mustset.pdf}~deltaAzimuth & \hfuzz=500pt angle & \hfuzz=500pt [degree] step size\\
\hfuzz=500pt\includegraphics[width=1em]{connector.pdf}\includegraphics[width=1em]{element-mustset.pdf}~deltaZenith & \hfuzz=500pt angle & \hfuzz=500pt [degree] step size\\
\hfuzz=500pt\includegraphics[width=1em]{connector.pdf}\includegraphics[width=1em]{element.pdf}~maxZenith & \hfuzz=500pt angle & \hfuzz=500pt [degree]\\
\hfuzz=500pt\includegraphics[width=1em]{connector.pdf}\includegraphics[width=1em]{element.pdf}~values & \hfuzz=500pt expression & \hfuzz=500pt [m] expression (zenith, azimuth: variables)\\
\hline
\end{tabularx}


\subsection{FromFile}
Select all or the first antenna from an \file{antenna definition file}{gnssAntennaDefinition}
which matches the wildcards.


\keepXColumns
\begin{tabularx}{\textwidth}{N T A}
\hline
Name & Type & Annotation\\
\hline
\hfuzz=500pt\includegraphics[width=1em]{element-mustset.pdf}~inputfileAntennaDefinition & \hfuzz=500pt filename & \hfuzz=500pt \\
\hfuzz=500pt\includegraphics[width=1em]{element.pdf}~name & \hfuzz=500pt string & \hfuzz=500pt \\
\hfuzz=500pt\includegraphics[width=1em]{element.pdf}~serial & \hfuzz=500pt string & \hfuzz=500pt \\
\hfuzz=500pt\includegraphics[width=1em]{element.pdf}~radome & \hfuzz=500pt string & \hfuzz=500pt \\
\hfuzz=500pt\includegraphics[width=1em]{element.pdf}~onlyFirstMatch & \hfuzz=500pt boolean & \hfuzz=500pt otherwise all machting antennas included\\
\hline
\end{tabularx}


\subsection{FromStationInfo}
Select all antennas from an \file{antenna definition file}{gnssAntennaDefinition}
which are used by a station within a defined time interval.
With \config{specializeAntenna} an individual antenna is created for each different serial number
using the general type specific values from file.


\keepXColumns
\begin{tabularx}{\textwidth}{N T A}
\hline
Name & Type & Annotation\\
\hline
\hfuzz=500pt\includegraphics[width=1em]{element-mustset.pdf}~inputfileStationInfo & \hfuzz=500pt filename & \hfuzz=500pt \\
\hfuzz=500pt\includegraphics[width=1em]{element-mustset.pdf}~inputfileAntennaDefinition & \hfuzz=500pt filename & \hfuzz=500pt \\
\hfuzz=500pt\includegraphics[width=1em]{element.pdf}~timeStart & \hfuzz=500pt time & \hfuzz=500pt only antennas used in this time interval\\
\hfuzz=500pt\includegraphics[width=1em]{element.pdf}~timeEnd & \hfuzz=500pt time & \hfuzz=500pt only antennas used in this time interval\\
\hfuzz=500pt\includegraphics[width=1em]{element.pdf}~specializeAntenna & \hfuzz=500pt boolean & \hfuzz=500pt e.g. separate different serial numbers from stationInfo\\
\hline
\end{tabularx}


\subsection{Resample}
The azimuth and elevation dependend antenna center variations (patterns) of all \config{antenna}s
are resampled to a new resolution.


\keepXColumns
\begin{tabularx}{\textwidth}{N T A}
\hline
Name & Type & Annotation\\
\hline
\hfuzz=500pt\includegraphics[width=1em]{element-mustset-unbounded.pdf}~antenna & \hfuzz=500pt \hyperref[gnssAntennaDefintionListType]{gnssAntennaDefintionList} & \hfuzz=500pt \\
\hfuzz=500pt\includegraphics[width=1em]{element.pdf}~deltaAzimuth & \hfuzz=500pt angle & \hfuzz=500pt [degree] step size, empty: no change\\
\hfuzz=500pt\includegraphics[width=1em]{element.pdf}~deltaZenith & \hfuzz=500pt angle & \hfuzz=500pt [degree] step size, empty: no change\\
\hfuzz=500pt\includegraphics[width=1em]{element.pdf}~maxZenith & \hfuzz=500pt angle & \hfuzz=500pt [degree], empty: no change\\
\hline
\end{tabularx}


\subsection{Transform}
This class can be used to separate general antenna patterns for different \configClass{gnssType}{gnssType}s.
If the \config{antenna}s contain only one pattern for all GPS observations on the L1 frequency (\verb|*1*G**|),
the \config{patternTypes}=\verb|C1*G**| and \verb|L1*G**| create two patterns with the \verb|*1*G**| patterm as template.
The first matching pattern in the \config{antenna} is used as template.
Also new \config{additionalPattern} can be added (e.g. for \verb|*5*G**|).
With \config{addExistingPatterns} all already existing patterns that don't match completely to any of the above are added.


\keepXColumns
\begin{tabularx}{\textwidth}{N T A}
\hline
Name & Type & Annotation\\
\hline
\hfuzz=500pt\includegraphics[width=1em]{element-mustset-unbounded.pdf}~antenna & \hfuzz=500pt \hyperref[gnssAntennaDefintionListType]{gnssAntennaDefintionList} & \hfuzz=500pt \\
\hfuzz=500pt\includegraphics[width=1em]{element-unbounded.pdf}~patternTypes & \hfuzz=500pt \hyperref[gnssType]{gnssType} & \hfuzz=500pt gnssType for each pattern (first match is used)\\
\hfuzz=500pt\includegraphics[width=1em]{element-unbounded.pdf}~additionalPattern & \hfuzz=500pt sequence & \hfuzz=500pt additional new patterns\\
\hfuzz=500pt\includegraphics[width=1em]{connector.pdf}\includegraphics[width=1em]{element-mustset.pdf}~type & \hfuzz=500pt \hyperref[gnssType]{gnssType} & \hfuzz=500pt pattern matching of observation types\\
\hfuzz=500pt\includegraphics[width=1em]{connector.pdf}\includegraphics[width=1em]{element.pdf}~offsetX & \hfuzz=500pt double & \hfuzz=500pt [m] antenna center offset\\
\hfuzz=500pt\includegraphics[width=1em]{connector.pdf}\includegraphics[width=1em]{element.pdf}~offsetY & \hfuzz=500pt double & \hfuzz=500pt [m] antenna center offset\\
\hfuzz=500pt\includegraphics[width=1em]{connector.pdf}\includegraphics[width=1em]{element.pdf}~offsetZ & \hfuzz=500pt double & \hfuzz=500pt [m] antenna center offset\\
\hfuzz=500pt\includegraphics[width=1em]{connector.pdf}\includegraphics[width=1em]{element-mustset.pdf}~deltaAzimuth & \hfuzz=500pt angle & \hfuzz=500pt [degree] step size\\
\hfuzz=500pt\includegraphics[width=1em]{connector.pdf}\includegraphics[width=1em]{element-mustset.pdf}~deltaZenith & \hfuzz=500pt angle & \hfuzz=500pt [degree] step size\\
\hfuzz=500pt\includegraphics[width=1em]{connector.pdf}\includegraphics[width=1em]{element.pdf}~maxZenith & \hfuzz=500pt angle & \hfuzz=500pt [degree]\\
\hfuzz=500pt\includegraphics[width=1em]{connector.pdf}\includegraphics[width=1em]{element.pdf}~values & \hfuzz=500pt expression & \hfuzz=500pt [m] expression (zenith, azimuth: variables)\\
\hfuzz=500pt\includegraphics[width=1em]{element.pdf}~addExistingPatterns & \hfuzz=500pt boolean & \hfuzz=500pt add existing patterns that don't match completely any of the above\\
\hline
\end{tabularx}

\clearpage
%==================================

\section{GnssParametrizationAmbiguities}\label{gnssParametrizationAmbiguitiesType}

This class provides functionality for GNSS phase ambiguities and integer ambiguity resolution.

Integer ambiguity resolution is performed based on the least squares ambiguity decorrelation adjustment
(LAMBDA) method (Teunissen 1995, DOI \href{https://doi.org/10.1007/BF00863419}{10.1007/BF00863419}), specifically
the modified algorithm (MLAMBDA) by Chang et al. (2005, DOI \href{https://doi.org/10.1007/s00190-005-0004-x}{10.1007/s00190-005-0004-x}).
First the covariance matrix of the integer ambiguity parameters is computed by eliminating all but those parameters
from the full normal equation matrix and inverting it. Then, a Z-transformation is performed as described by
Chang et al. (2005) to decorrelate the ambiguity parameters without losing their integer nature.

The search process follows MLAMBDA and uses integer minimization of the weighted sum of squared residuals.
It is computationally infeasible to search a hyper-ellipsoid with a dimension of ten thousand or more.
Instead, a blocked search algorithm is performed by moving a window with a length of, for example,
\config{searchBlockSize}=\verb|200| parameters over the decorrelated ambiguities, starting from the most accurate.
In each step, the window is moved by half of its length and the overlapping parts are compared to each other.
If all fixed ambiguities in the overlap agree, the algorithm continues.
Otherwise, both windows are combined and the search is repeated using the combined window, again comparing with the overlapping
part of the preceding window. If not all solutions could be checked for a block after \config{maxSearchSteps},
the selected \config{incompleteAction} is performed.
If the algorithm reaches ambiguities with a standard deviation higher than \config{sigmaMaxResolve},
ambiguity resolution stops and the remaining ambiguities are left as float values.
Otherwise, all ambiguity parameters are fixed to integer values.

In contrast to an integer least squares solution over the full ambiguity vector, it is not guaranteed that the resulting solution
is optimal in the sense of minimal variance with given covariance.
This trade-off is necessary to cope with large numbers of ambiguities.

See also \program{GnssProcessing}.


\keepXColumns
\begin{tabularx}{\textwidth}{N T A}
\hline
Name & Type & Annotation\\
\hline
\hfuzz=500pt\includegraphics[width=1em]{element-mustset.pdf}~gnssParametrizationAmbiguitiesType & \hfuzz=500pt sequence & \hfuzz=500pt \\
\hfuzz=500pt\includegraphics[width=1em]{connector.pdf}\includegraphics[width=1em]{element.pdf}~outputfileAmbiguities & \hfuzz=500pt filename & \hfuzz=500pt resolved ambiguities, station name is appended to file name for single receivers\\
\hfuzz=500pt\includegraphics[width=1em]{connector.pdf}\includegraphics[width=1em]{element.pdf}~sigmaMaxResolve & \hfuzz=500pt double & \hfuzz=500pt max. allowed std. dev. of ambiguity to resolve [cycles]\\
\hfuzz=500pt\includegraphics[width=1em]{connector.pdf}\includegraphics[width=1em]{element.pdf}~searchBlockSize & \hfuzz=500pt uint & \hfuzz=500pt block size for blocked integer search\\
\hfuzz=500pt\includegraphics[width=1em]{connector.pdf}\includegraphics[width=1em]{element.pdf}~maxSearchSteps & \hfuzz=500pt uint & \hfuzz=500pt max. steps of integer search for each block\\
\hfuzz=500pt\includegraphics[width=1em]{connector.pdf}\includegraphics[width=1em]{element-mustset.pdf}~incompleteAction & \hfuzz=500pt choice & \hfuzz=500pt if not all solutions tested after maxSearchSteps\\
\hfuzz=500pt\quad\includegraphics[width=1em]{connector.pdf}\includegraphics[width=1em]{element-mustset.pdf}~stop & \hfuzz=500pt  & \hfuzz=500pt stop searching, ambiguities remain float in this block\\
\hfuzz=500pt\quad\includegraphics[width=1em]{connector.pdf}\includegraphics[width=1em]{element-mustset.pdf}~resolve & \hfuzz=500pt  & \hfuzz=500pt use best integer solution found so far\\
\hfuzz=500pt\quad\includegraphics[width=1em]{connector.pdf}\includegraphics[width=1em]{element-mustset.pdf}~shrinkBlockSize & \hfuzz=500pt  & \hfuzz=500pt try again with half block size\\
\hfuzz=500pt\quad\includegraphics[width=1em]{connector.pdf}\includegraphics[width=1em]{element-mustset.pdf}~throwException & \hfuzz=500pt  & \hfuzz=500pt stop and throw an exception\\
\hline
\end{tabularx}

\clearpage
%==================================

\section{GnssParametrizationAntenna}\label{gnssParametrizationAntennaType}

Antenna center offsets or variations for satellite or ground station antennas can be estimated by defining
the parametrization via \configClass{antennaCenterVariations}{parametrizationGnssAntennaType}.
The defined parametrization is set up for each type listed in \configClass{patternTypes}{gnssType}, so,
for example, GPS phase center offsets can be estimated separately for \verb|L1*G|, \verb|L2*G|, and \verb|L5*G|.

If \config{groupAntennas}=\verb|yes| parameters are set up per antenna type instead of per individual antenna.
This enables the estimation of, for example, one set of antenna center variations per GPS satellite block.

See also \program{GnssProcessing}.


\keepXColumns
\begin{tabularx}{\textwidth}{N T A}
\hline
Name & Type & Annotation\\
\hline
\hfuzz=500pt\includegraphics[width=1em]{element-mustset.pdf}~gnssParametrizationAntennaType & \hfuzz=500pt sequence & \hfuzz=500pt \\
\hfuzz=500pt\includegraphics[width=1em]{connector.pdf}\includegraphics[width=1em]{element-mustset-unbounded.pdf}~antennaCenterVariations & \hfuzz=500pt \hyperref[parametrizationGnssAntennaType]{parametrizationGnssAntenna} & \hfuzz=500pt estimate antenna center variations\\
\hfuzz=500pt\includegraphics[width=1em]{connector.pdf}\includegraphics[width=1em]{element-unbounded.pdf}~patternTypes & \hfuzz=500pt \hyperref[gnssType]{gnssType} & \hfuzz=500pt gnssType for each pattern (first match is used)\\
\hfuzz=500pt\includegraphics[width=1em]{connector.pdf}\includegraphics[width=1em]{element.pdf}~addNonMatchingTypes & \hfuzz=500pt boolean & \hfuzz=500pt add patterns for additional observed gnssTypes that don't match any of the above\\
\hfuzz=500pt\includegraphics[width=1em]{connector.pdf}\includegraphics[width=1em]{element.pdf}~groupAntennas & \hfuzz=500pt boolean & \hfuzz=500pt common ACVs for same antenna build types (ignores antenna serial number)\\
\hline
\end{tabularx}

\clearpage
%==================================

\section{GnssParametrizationConstraints}\label{gnssParametrizationConstraintsType}

Constraints for selected \configClass{parameters}{parameterSelectorType} can be added to the normal equation system
in case there are no direct constraint options in the respective config element (e.g. \config{noNetRotationSigma} in
\configClass{gnssParametrizationReceiver:stationNetwork}{gnssParametrizationReceiverType:stationNetwork}).

For details on possible constraints see \configClass{gnssParametrizationReceiver:stationNetwork}{gnssParametrizationReceiverType:stationNetwork}
(tropospheric parameters), \configClass{gnssParametrizationTransmitter}{gnssParametrizationTransmitterType} (antenna center offsets),
or \configClass{gnssParametrizationEarthRotation}{gnssParametrizationEarthRotationType} (dUT1).

See also \program{GnssProcessing}.


\keepXColumns
\begin{tabularx}{\textwidth}{N T A}
\hline
Name & Type & Annotation\\
\hline
\hfuzz=500pt\includegraphics[width=1em]{element-mustset.pdf}~gnssParametrizationConstraintsType & \hfuzz=500pt sequence & \hfuzz=500pt \\
\hfuzz=500pt\includegraphics[width=1em]{connector.pdf}\includegraphics[width=1em]{element.pdf}~comment & \hfuzz=500pt string & \hfuzz=500pt \\
\hfuzz=500pt\includegraphics[width=1em]{connector.pdf}\includegraphics[width=1em]{element-mustset-unbounded.pdf}~parameters & \hfuzz=500pt \hyperref[parameterSelectorType]{parameterSelector} & \hfuzz=500pt parameter to constrain\\
\hfuzz=500pt\includegraphics[width=1em]{connector.pdf}\includegraphics[width=1em]{element-mustset.pdf}~sigma & \hfuzz=500pt double & \hfuzz=500pt sigma of the constraint (same unit as parameter)\\
\hfuzz=500pt\includegraphics[width=1em]{connector.pdf}\includegraphics[width=1em]{element.pdf}~bias & \hfuzz=500pt double & \hfuzz=500pt constrain all selected parameters towards this value\\
\hfuzz=500pt\includegraphics[width=1em]{connector.pdf}\includegraphics[width=1em]{element.pdf}~relativeToApriori & \hfuzz=500pt boolean & \hfuzz=500pt constrain only dx and not full x=dx+x0\\
\hline
\end{tabularx}

\clearpage
%==================================

\section{GnssParametrizationEarthRotation}\label{gnssParametrizationEarthRotationType}

Earth rotation parameters (ERPs) can be estimated by defining
\config{estimatePole} ($x_p$, $y_p$) and \config{estimateUT1} (dUT1, LOD).

Estimating length of day (LOD) with the sign according to IGS conventions requires a negative
value in \configClass{parametrizationTemporal:trend:timeStep}{parametrizationTemporalType:trend}.

Constraints on the defined parameters can be added via
\configClass{gnssParametrizationConstraints}{gnssParametrizationConstraintsType}.
An example would be to set up \configClass{estimateUT1:constant}{parametrizationTemporalType:constant}
so the \emph{dUT1} parameter is included in the normal equation system . Since \emph{dUT1} cannot be
determined by GNSS, a hard constraint to its a priori value can then be added.

See also \program{GnssProcessing} and \program{GnssSimulateReceiver}.


\keepXColumns
\begin{tabularx}{\textwidth}{N T A}
\hline
Name & Type & Annotation\\
\hline
\hfuzz=500pt\includegraphics[width=1em]{element-mustset.pdf}~gnssParametrizationEarthRotationType & \hfuzz=500pt sequence & \hfuzz=500pt \\
\hfuzz=500pt\includegraphics[width=1em]{connector.pdf}\includegraphics[width=1em]{element.pdf}~outputfileEOP & \hfuzz=500pt filename & \hfuzz=500pt EOP time series (mjd, xp, yp, sp, dUT1, LOD, X, Y, S)\\
\hfuzz=500pt\includegraphics[width=1em]{connector.pdf}\includegraphics[width=1em]{element-mustset.pdf}~earthRotation & \hfuzz=500pt \hyperref[earthRotationType]{earthRotation} & \hfuzz=500pt apriori earth rotation\\
\hfuzz=500pt\includegraphics[width=1em]{connector.pdf}\includegraphics[width=1em]{element.pdf}~earthRotationModels & \hfuzz=500pt \hyperref[earthRotationType]{earthRotation} & \hfuzz=500pt will be reduced from reported apriori parameters\\
\hfuzz=500pt\includegraphics[width=1em]{connector.pdf}\includegraphics[width=1em]{element-unbounded.pdf}~estimatePole & \hfuzz=500pt \hyperref[parametrizationTemporalType]{parametrizationTemporal} & \hfuzz=500pt xp, yp [mas]\\
\hfuzz=500pt\includegraphics[width=1em]{connector.pdf}\includegraphics[width=1em]{element-unbounded.pdf}~estimateUT1 & \hfuzz=500pt \hyperref[parametrizationTemporalType]{parametrizationTemporal} & \hfuzz=500pt rotation angle [ms]\\
\hfuzz=500pt\includegraphics[width=1em]{connector.pdf}\includegraphics[width=1em]{element-unbounded.pdf}~estimateNutation & \hfuzz=500pt \hyperref[parametrizationTemporalType]{parametrizationTemporal} & \hfuzz=500pt dX, dY [mas]\\
\hfuzz=500pt\includegraphics[width=1em]{connector.pdf}\includegraphics[width=1em]{element.pdf}~interpolationDegree & \hfuzz=500pt uint & \hfuzz=500pt for quaternion interpolation\\
\hline
\end{tabularx}

\clearpage
%==================================

\section{GnssParametrizationGravityField}\label{gnssParametrizationGravityFieldType}

Estimation of low-degree gravity field parameters via GNSS.

See also \program{GnssProcessing}.


\keepXColumns
\begin{tabularx}{\textwidth}{N T A}
\hline
Name & Type & Annotation\\
\hline
\hfuzz=500pt\includegraphics[width=1em]{element-mustset.pdf}~gnssParametrizationGravityFieldType & \hfuzz=500pt sequence & \hfuzz=500pt \\
\hfuzz=500pt\includegraphics[width=1em]{connector.pdf}\includegraphics[width=1em]{element-mustset-unbounded.pdf}~parametrization & \hfuzz=500pt \hyperref[parametrizationGravityType]{parametrizationGravity} & \hfuzz=500pt gravity field parametrization\\
\hline
\end{tabularx}

\clearpage
%==================================

\section{GnssParametrizationIonosphere}\label{gnssParametrizationIonosphereType}

Definition of settings and constants for ionospheric corrections in GNSS observation equations.

See also \program{GnssProcessing} and \program{GnssSimulateReceiver}.


\keepXColumns
\begin{tabularx}{\textwidth}{N T A}
\hline
Name & Type & Annotation\\
\hline
\hfuzz=500pt\includegraphics[width=1em]{element-mustset.pdf}~gnssParametrizationIonosphereType & \hfuzz=500pt sequence & \hfuzz=500pt \\
\hfuzz=500pt\includegraphics[width=1em]{connector.pdf}\includegraphics[width=1em]{element.pdf}~outputfileVTEC & \hfuzz=500pt filename & \hfuzz=500pt \\
\hfuzz=500pt\includegraphics[width=1em]{connector.pdf}\includegraphics[width=1em]{element-mustset.pdf}~magnetosphere & \hfuzz=500pt \hyperref[magnetosphereType]{magnetosphere} & \hfuzz=500pt \\
\hfuzz=500pt\includegraphics[width=1em]{connector.pdf}\includegraphics[width=1em]{element.pdf}~apply2ndOrderCorrection & \hfuzz=500pt boolean & \hfuzz=500pt apply ionospheric correction\\
\hfuzz=500pt\includegraphics[width=1em]{connector.pdf}\includegraphics[width=1em]{element.pdf}~apply3rdOrderCorrection & \hfuzz=500pt boolean & \hfuzz=500pt apply ionospheric correction\\
\hfuzz=500pt\includegraphics[width=1em]{connector.pdf}\includegraphics[width=1em]{element.pdf}~applyBendingCorrection & \hfuzz=500pt boolean & \hfuzz=500pt apply ionospheric correction\\
\hfuzz=500pt\includegraphics[width=1em]{connector.pdf}\includegraphics[width=1em]{element.pdf}~singleLayerHeight & \hfuzz=500pt double & \hfuzz=500pt [m] height of the ionospheric pierce point\\
\hfuzz=500pt\includegraphics[width=1em]{connector.pdf}\includegraphics[width=1em]{element.pdf}~mapR & \hfuzz=500pt double & \hfuzz=500pt constant of MSLM mapping function\\
\hfuzz=500pt\includegraphics[width=1em]{connector.pdf}\includegraphics[width=1em]{element.pdf}~mapH & \hfuzz=500pt double & \hfuzz=500pt constant of MSLM mapping function\\
\hfuzz=500pt\includegraphics[width=1em]{connector.pdf}\includegraphics[width=1em]{element.pdf}~mapAlpha & \hfuzz=500pt double & \hfuzz=500pt constant of MSLM mapping function\\
\hline
\end{tabularx}

\clearpage
%==================================

\section{GnssParametrizationReceiver}\label{gnssParametrizationReceiverType}

Definition and parametrization of GNSS receivers.

Input files used in \config{stationNetwork} and \config{lowEarthOrbiter}:
\begin{itemize}
  \item \configFile{inputfileStationInfo}{gnssStationInfo}:
        Created via \program{GnssStationLog2StationInfo} or \program{GnssStationInfoCreate}.
  \item \configFile{inputfileAntennaDefinition}{gnssAntennaDefinition}:
        Created via \program{GnssAntex2AntennaDefinition} or \program{GnssAntennaDefinitionCreate}.
  \item \configFile{inputfileReceiverDefinition}{gnssReceiverDefinition}:
        Created via \program{GnssReceiverDefinitionCreate} in case you want to define which signal
        types a receiver model can observe.
  \item \configFile{inputfileAccuracyDefinition}{gnssAntennaDefinition}:
        Created via \program{GnssAntennaDefinitionCreate}.
  \item \configFile{inputfileSignalBias}{gnssSignalBias}:
        Estimated via \program{GnssProcessing}.
  \item \configFile{inputfileObservation}{instrument}:
        Converted from RINEX observation files via \program{RinexObservation2GnssReceiver}.
\end{itemize}

See also \program{GnssProcessing} and \program{GnssSimulateReceiver}.


\section{StationNetwork}\label{gnssParametrizationReceiverType:stationNetwork}

A network of GNSS ground stations is defined via \configFile{inputfileStationList}{stringList}.
All input and output files except \configFile{outputfileGriddedPosition}{griddedData},
\configFile{outputfileUsedStationList}{stringList}, \configFile{inputfileAntennaDefinition}{gnssAntennaDefinition},
\configFile{inputfileReceiverDefinition}{gnssReceiverDefinition}, and
\configFile{inputfileAccuracyDefinition}{gnssAntennaDefinition} are read/written for each station.
The file name is interpreted as a template with the variable \verb|{station}| being replaced by the station name.

If \config{kinematicPositionEstimation}=\verb|yes| station positions are estimated every epoch,
otherwise they are parameterized according to \configClass{positionEstimation}{parametrizationTemporalType}.

In case you estimate the station network together with a GNSS constellations, \config{noNetTranslationSigma} and \config{noNetRotationSigma}
can be set to remove the rank deficiency. No-net constraints are only applied to stations defined in \configFile{noNetStationList}{stringList} that have
precise station coordinates provided via \configFile{inputfileStationPosition}{instrument}. Those will be used instead of the approximate
coordinates from \configFile{inputfileStationInfo}{gnssStationInfo}. In case you want to align to an ITRF/IGS reference frame, precise coordinates can be
generated via combining \program{Sinex2StationPosition}, \program{Sinex2StationDiscontinuities}, and \program{Sinex2StationPostSeismicDeformation}.

A priori tropospheric correction is handled by a \configClass{troposphere}{troposphereType} model (e.g. Vienna Mapping Functions 3).
Additional parameters for zenith wet delay and gradients can be set up via \configClass{troposphereWetEstimation}{parametrizationTemporalType}
and \configClass{troposphereGradientEstimation}{parametrizationTemporalType}. These parameters can be soft-constrained using
\configClass{constraints}{gnssParametrizationConstraintsType} in \program{GnssProcessing}.

Antenna center variations can be estimated by defining \configClass{antennaCenter}{gnssParametrizationAntennaType}.
It is possible to estimate patterns for each individual station or grouped by antenna model.
Antenna center parameters can be constrained by adding \configClass{constraints}{gnssParametrizationConstraintsType} in \program{GnssProcessing}.

The effects of loading and tidal deformation on station positions can be corrected for via \configClass{deformation}{gravityfieldType} and
\configClass{tides}{tidesType}, respectively. These typically include:
\begin{itemize}
  \item \configClass{tides:earthTide}{tidesType:earthTide}: Earth tidal deformations (IERS conventions)
  \item \configClass{tides:doodsonHarmonicTide}{tidesType:doodsonHarmonicTide}: ocean tidal deformations
        (e.g. FES 2014b, \config{minDegree}=\verb|1|)
  \item \configClass{tides:doodsonHarmonicTide}{tidesType:doodsonHarmonicTide}: atmospheric tidal deformation
        (e.g. AOD1B RL06, \config{minDegree}=\verb|1|)
  \item \configClass{tides:poleTide}{tidesType:poleTide}: pole tidal deformations (IERS conventions)
  \item \configClass{tides:poleOceanTide}{tidesType:oceanPoleTide}: ocean pole tidal deformations (IERS conventions)
\end{itemize}

Each station goes through a \config{preprocessing} step individually, where observation outliers are removed or downweighted,
continuous tracks of phase observations are defined for ambiguity parametrization, cycle slips are detected, and stations are
disabled if they do not fulfill certain requirements.


\keepXColumns
\begin{tabularx}{\textwidth}{N T A}
\hline
Name & Type & Annotation\\
\hline
\hfuzz=500pt\includegraphics[width=1em]{element.pdf}~outputfileGriddedPosition & \hfuzz=500pt filename & \hfuzz=500pt delta north east up for all stations\\
\hfuzz=500pt\includegraphics[width=1em]{element.pdf}~outputfilePosition & \hfuzz=500pt filename & \hfuzz=500pt full estimated coordinates (in TRF)\\
\hfuzz=500pt\includegraphics[width=1em]{element.pdf}~outputfilePositionSeries & \hfuzz=500pt filename & \hfuzz=500pt time series of position deltas etc.\\
\hfuzz=500pt\includegraphics[width=1em]{element.pdf}~outputfileResiduals & \hfuzz=500pt filename & \hfuzz=500pt \\
\hfuzz=500pt\includegraphics[width=1em]{element.pdf}~outputfileSignalBias & \hfuzz=500pt filename & \hfuzz=500pt estimated signal biases\\
\hfuzz=500pt\includegraphics[width=1em]{element.pdf}~outputfileClock & \hfuzz=500pt filename & \hfuzz=500pt estimated clock errors\\
\hfuzz=500pt\includegraphics[width=1em]{element.pdf}~outputfileUsedStationList & \hfuzz=500pt filename & \hfuzz=500pt ascii file with names of used stations\\
\hfuzz=500pt\includegraphics[width=1em]{element.pdf}~outputfileTroposphere & \hfuzz=500pt filename & \hfuzz=500pt columns: MJD, ZHD, ZWD, dry north gradient, wet north gradient, dry east gradient, wet east gradient\\
\hfuzz=500pt\includegraphics[width=1em]{element-mustset.pdf}~inputfileStationList & \hfuzz=500pt filename & \hfuzz=500pt ascii file with station names\\
\hfuzz=500pt\includegraphics[width=1em]{element-mustset.pdf}~inputfileStationInfo & \hfuzz=500pt filename & \hfuzz=500pt station metadata (antennas, receivers, ...)\\
\hfuzz=500pt\includegraphics[width=1em]{element-mustset.pdf}~inputfileAntennaDefinition & \hfuzz=500pt filename & \hfuzz=500pt antenna center offsets and variations\\
\hfuzz=500pt\includegraphics[width=1em]{element.pdf}~inputfileReceiverDefinition & \hfuzz=500pt filename & \hfuzz=500pt observed signal types\\
\hfuzz=500pt\includegraphics[width=1em]{element-mustset.pdf}~inputfileAccuracyDefinition & \hfuzz=500pt filename & \hfuzz=500pt elevation and azimuth dependent accuracy\\
\hfuzz=500pt\includegraphics[width=1em]{element.pdf}~inputfileStationPosition & \hfuzz=500pt filename & \hfuzz=500pt precise coordinates used for no-net constraints (in TRF)\\
\hfuzz=500pt\includegraphics[width=1em]{element.pdf}~inputfileSignalBias & \hfuzz=500pt filename & \hfuzz=500pt a priori signal biases\\
\hfuzz=500pt\includegraphics[width=1em]{element.pdf}~inputfileObservations & \hfuzz=500pt filename & \hfuzz=500pt \\
\hfuzz=500pt\includegraphics[width=1em]{element.pdf}~kinematicPositionEstimation & \hfuzz=500pt boolean & \hfuzz=500pt positions are estimated every epoch (in TRF)\\
\hfuzz=500pt\includegraphics[width=1em]{element-unbounded.pdf}~positionEstimation & \hfuzz=500pt \hyperref[parametrizationTemporalType]{parametrizationTemporal} & \hfuzz=500pt parametrization of estimated position (in TRF)\\
\hfuzz=500pt\includegraphics[width=1em]{element.pdf}~noNetTranslationSigma & \hfuzz=500pt double & \hfuzz=500pt (0 = unconstrained) sigma [m] for no-net translation constraint on station coordinates\\
\hfuzz=500pt\includegraphics[width=1em]{element.pdf}~noNetRotationSigma & \hfuzz=500pt double & \hfuzz=500pt (0 = unconstrained) sigma [m] at Earth's surface for no-net rotation constraint on station coordinates\\
\hfuzz=500pt\includegraphics[width=1em]{element.pdf}~noNetStationList & \hfuzz=500pt filename & \hfuzz=500pt ascii file with names of stations used for no-net constraints (empty = all stations)\\
\hfuzz=500pt\includegraphics[width=1em]{element.pdf}~estimateCodeBias & \hfuzz=500pt boolean & \hfuzz=500pt \\
\hfuzz=500pt\includegraphics[width=1em]{element.pdf}~estimatePhaseBias & \hfuzz=500pt boolean & \hfuzz=500pt \\
\hfuzz=500pt\includegraphics[width=1em]{element.pdf}~supportsIntegerAmbiguities & \hfuzz=500pt boolean & \hfuzz=500pt receiver tracks full cycle integer ambiguities\\
\hfuzz=500pt\includegraphics[width=1em]{element.pdf}~troposphere & \hfuzz=500pt \hyperref[troposphereType]{troposphere} & \hfuzz=500pt a priori troposphere model\\
\hfuzz=500pt\includegraphics[width=1em]{element-unbounded.pdf}~troposphereWetEstimation & \hfuzz=500pt \hyperref[parametrizationTemporalType]{parametrizationTemporal} & \hfuzz=500pt [m] parametrization of zenith wet delays\\
\hfuzz=500pt\includegraphics[width=1em]{element-unbounded.pdf}~troposphereGradientEstimation & \hfuzz=500pt \hyperref[parametrizationTemporalType]{parametrizationTemporal} & \hfuzz=500pt [degree] parametrization of north and east gradients\\
\hfuzz=500pt\includegraphics[width=1em]{element-unbounded.pdf}~antennaCenter & \hfuzz=500pt \hyperref[gnssParametrizationAntennaType]{gnssParametrizationAntenna} & \hfuzz=500pt estimate antenna center variations\\
\hfuzz=500pt\includegraphics[width=1em]{element-unbounded.pdf}~deformation & \hfuzz=500pt \hyperref[gravityfieldType]{gravityfield} & \hfuzz=500pt loading deformation\\
\hfuzz=500pt\includegraphics[width=1em]{element-unbounded.pdf}~tides & \hfuzz=500pt \hyperref[tidesType]{tides} & \hfuzz=500pt tidal deformation\\
\hfuzz=500pt\includegraphics[width=1em]{element.pdf}~ephemerides & \hfuzz=500pt \hyperref[ephemeridesType]{ephemerides} & \hfuzz=500pt for tidal deformation\\
\hfuzz=500pt\includegraphics[width=1em]{element-mustset.pdf}~inputfileDeformationLoadLoveNumber & \hfuzz=500pt filename & \hfuzz=500pt \\
\hfuzz=500pt\includegraphics[width=1em]{element.pdf}~inputfilePotentialLoadLoveNumber & \hfuzz=500pt filename & \hfuzz=500pt if full potential is given and not only loading potential\\
\hfuzz=500pt\includegraphics[width=1em]{element-mustset.pdf}~preprocessing & \hfuzz=500pt sequence & \hfuzz=500pt settings for preprocessing of observations/stations\\
\hfuzz=500pt\includegraphics[width=1em]{connector.pdf}\includegraphics[width=1em]{element.pdf}~outputfileTrackBefore & \hfuzz=500pt filename & \hfuzz=500pt TEC and MW-like combinations in cycles for track before preprocessing (variables: loopTime, station, prn, trackTimeStart, trackTimeEnd)\\
\hfuzz=500pt\includegraphics[width=1em]{connector.pdf}\includegraphics[width=1em]{element.pdf}~outputfileTrackAfter & \hfuzz=500pt filename & \hfuzz=500pt TEC and MW-like combinations in cycles for track after preprocessing (variables: loopTime, station, prn, trackTimeStart, trackTimeEnd)\\
\hfuzz=500pt\includegraphics[width=1em]{connector.pdf}\includegraphics[width=1em]{element.pdf}~maxStationCount & \hfuzz=500pt uint & \hfuzz=500pt maximum number of stations to be used\\
\hfuzz=500pt\includegraphics[width=1em]{connector.pdf}\includegraphics[width=1em]{element.pdf}~elevationCutOff & \hfuzz=500pt angle & \hfuzz=500pt [degree] ignore observations below cutoff\\
\hfuzz=500pt\includegraphics[width=1em]{connector.pdf}\includegraphics[width=1em]{element.pdf}~elevationTrackMinimum & \hfuzz=500pt angle & \hfuzz=500pt [degree] ignore tracks that never exceed minimum elevation\\
\hfuzz=500pt\includegraphics[width=1em]{connector.pdf}\includegraphics[width=1em]{element-unbounded.pdf}~useType & \hfuzz=500pt \hyperref[gnssType]{gnssType} & \hfuzz=500pt only use observations that match any of these patterns\\
\hfuzz=500pt\includegraphics[width=1em]{connector.pdf}\includegraphics[width=1em]{element-unbounded.pdf}~ignoreType & \hfuzz=500pt \hyperref[gnssType]{gnssType} & \hfuzz=500pt ignore observations that match any of these patterns\\
\hfuzz=500pt\includegraphics[width=1em]{connector.pdf}\includegraphics[width=1em]{element-mustset.pdf}~noAntennaPatternFound & \hfuzz=500pt choice & \hfuzz=500pt what should happen if no antenna pattern is found for an observation\\
\hfuzz=500pt\quad\includegraphics[width=1em]{connector.pdf}\includegraphics[width=1em]{element-mustset.pdf}~ignoreObservation & \hfuzz=500pt  & \hfuzz=500pt ignore observation if no matching pattern is found\\
\hfuzz=500pt\quad\includegraphics[width=1em]{connector.pdf}\includegraphics[width=1em]{element-mustset.pdf}~useNearestFrequency & \hfuzz=500pt  & \hfuzz=500pt use pattern of nearest frequency if no matching pattern is found\\
\hfuzz=500pt\quad\includegraphics[width=1em]{connector.pdf}\includegraphics[width=1em]{element-mustset.pdf}~throwException & \hfuzz=500pt  & \hfuzz=500pt throw exception if no matching pattern is found\\
\hfuzz=500pt\includegraphics[width=1em]{connector.pdf}\includegraphics[width=1em]{element.pdf}~huber & \hfuzz=500pt double & \hfuzz=500pt residuals \$>\$ huber*sigma0 are downweighted\\
\hfuzz=500pt\includegraphics[width=1em]{connector.pdf}\includegraphics[width=1em]{element.pdf}~huberPower & \hfuzz=500pt double & \hfuzz=500pt residuals \$>\$ huber: sigma=(e/huber)\textasciicircum{}huberPower*sigma0\\
\hfuzz=500pt\includegraphics[width=1em]{connector.pdf}\includegraphics[width=1em]{element.pdf}~codeMaxPositionDiff & \hfuzz=500pt double & \hfuzz=500pt [m] max. allowed position error by PPP code only clock error estimation\\
\hfuzz=500pt\includegraphics[width=1em]{connector.pdf}\includegraphics[width=1em]{element.pdf}~minObsCountPerTrack & \hfuzz=500pt uint & \hfuzz=500pt tracks with less number of epochs with observations are dropped\\
\hfuzz=500pt\includegraphics[width=1em]{connector.pdf}\includegraphics[width=1em]{element.pdf}~minEstimableEpochsRatio & \hfuzz=500pt double & \hfuzz=500pt [0,1] drop stations with lower ratio of estimable epochs to total epochs\\
\hfuzz=500pt\includegraphics[width=1em]{connector.pdf}\includegraphics[width=1em]{element.pdf}~denoisingLambda & \hfuzz=500pt double & \hfuzz=500pt regularization parameter for total variation denoising used in cylce slip detection\\
\hfuzz=500pt\includegraphics[width=1em]{connector.pdf}\includegraphics[width=1em]{element.pdf}~tecWindowSize & \hfuzz=500pt uint & \hfuzz=500pt (0 = disabled) window size for TEC smoothness evaluation used in cycle slip detection\\
\hfuzz=500pt\includegraphics[width=1em]{connector.pdf}\includegraphics[width=1em]{element.pdf}~tecSigmaFactor & \hfuzz=500pt double & \hfuzz=500pt factor applied to moving standard deviation used as threshold in TEC smoothness evaluation during cycle slip detection\\
\hline
\end{tabularx}


\section{LowEarthOrbiter}\label{gnssParametrizationReceiverType:lowEarthOrbiter}

A single low-Earth orbiting (LEO) satellite with an onboard GNSS receiver. Orbit positions and receiver clock errors are estimated per epoch.

Attitude data must be provided via \configFile{inputfileStarCamera}{instrument}. If no attitude data is available from the satellite operator,
the star camera data can be simulated by using \program{SimulateStarCamera} or \program{SimulateStarCameraSentinel1}.

Antenna center variations can be estimated by defining \configClass{antennaCenter}{gnssParametrizationAntennaType}.

In a \config{preprocessing} step observation outliers are removed or downweighted, continuous tracks of phase observations are
defined for ambiguity parametrization, and cycle slips are detected.


\keepXColumns
\begin{tabularx}{\textwidth}{N T A}
\hline
Name & Type & Annotation\\
\hline
\hfuzz=500pt\includegraphics[width=1em]{element.pdf}~outputfileOrbit & \hfuzz=500pt filename & \hfuzz=500pt estimated kinematic orbit\\
\hfuzz=500pt\includegraphics[width=1em]{element.pdf}~outputfileCovariancePodEpoch & \hfuzz=500pt filename & \hfuzz=500pt 3x3 epoch covariances\\
\hfuzz=500pt\includegraphics[width=1em]{element.pdf}~outputfileClock & \hfuzz=500pt filename & \hfuzz=500pt estimtated clock errors\\
\hfuzz=500pt\includegraphics[width=1em]{element.pdf}~outputfileResiduals & \hfuzz=500pt filename & \hfuzz=500pt \\
\hfuzz=500pt\includegraphics[width=1em]{element.pdf}~outputfileSignalBias & \hfuzz=500pt filename & \hfuzz=500pt estimated signal biases\\
\hfuzz=500pt\includegraphics[width=1em]{element-mustset.pdf}~inputfileStationInfo & \hfuzz=500pt filename & \hfuzz=500pt satellite metadata (antenna, receiver, ...)\\
\hfuzz=500pt\includegraphics[width=1em]{element-mustset.pdf}~inputfileAntennaDefinition & \hfuzz=500pt filename & \hfuzz=500pt antenna center offsets and variations\\
\hfuzz=500pt\includegraphics[width=1em]{element.pdf}~inputfileReceiverDefinition & \hfuzz=500pt filename & \hfuzz=500pt observed signal types\\
\hfuzz=500pt\includegraphics[width=1em]{element-mustset.pdf}~inputfileAccuracyDefinition & \hfuzz=500pt filename & \hfuzz=500pt elevation and azimut dependent accuracy\\
\hfuzz=500pt\includegraphics[width=1em]{element.pdf}~inputfileSignalBias & \hfuzz=500pt filename & \hfuzz=500pt a priori signal biases\\
\hfuzz=500pt\includegraphics[width=1em]{element-mustset.pdf}~inputfileObservations & \hfuzz=500pt filename & \hfuzz=500pt \\
\hfuzz=500pt\includegraphics[width=1em]{element-mustset.pdf}~inputfileOrbit & \hfuzz=500pt filename & \hfuzz=500pt approximate positions\\
\hfuzz=500pt\includegraphics[width=1em]{element-mustset.pdf}~inputfileStarCamera & \hfuzz=500pt filename & \hfuzz=500pt satellite attitude\\
\hfuzz=500pt\includegraphics[width=1em]{element.pdf}~estimateCodeBias & \hfuzz=500pt boolean & \hfuzz=500pt \\
\hfuzz=500pt\includegraphics[width=1em]{element.pdf}~estimatePhaseBias & \hfuzz=500pt boolean & \hfuzz=500pt \\
\hfuzz=500pt\includegraphics[width=1em]{element.pdf}~supportsIntegerAmbiguities & \hfuzz=500pt boolean & \hfuzz=500pt receiver tracks full cycle integer ambiguities\\
\hfuzz=500pt\includegraphics[width=1em]{element.pdf}~wavelengthFactor & \hfuzz=500pt double & \hfuzz=500pt factor to account for half-wavelength observations (collected by codeless squaring techniques)\\
\hfuzz=500pt\includegraphics[width=1em]{element.pdf}~sigmaFactorPhase & \hfuzz=500pt expression & \hfuzz=500pt PHASE: factor = f(FREQ, ELE, SNR, ROTI, dTEc, IONOINDEX)\\
\hfuzz=500pt\includegraphics[width=1em]{element.pdf}~sigmaFactorCode & \hfuzz=500pt expression & \hfuzz=500pt CODE: factor = f(FREQ, ELE, SNR, ROTI, dTEc, IONOINDEX)\\
\hfuzz=500pt\includegraphics[width=1em]{element-unbounded.pdf}~antennaCenter & \hfuzz=500pt \hyperref[gnssParametrizationAntennaType]{gnssParametrizationAntenna} & \hfuzz=500pt estimate antenna center variations\\
\hfuzz=500pt\includegraphics[width=1em]{element-mustset.pdf}~preprocessing & \hfuzz=500pt sequence & \hfuzz=500pt \\
\hfuzz=500pt\includegraphics[width=1em]{connector.pdf}\includegraphics[width=1em]{element.pdf}~elevationCutOff & \hfuzz=500pt angle & \hfuzz=500pt [degree] ignore observations below cutoff\\
\hfuzz=500pt\includegraphics[width=1em]{connector.pdf}\includegraphics[width=1em]{element.pdf}~elevationTrackMinimum & \hfuzz=500pt angle & \hfuzz=500pt [degree] ignore tracks that never exceed minimum elevation\\
\hfuzz=500pt\includegraphics[width=1em]{connector.pdf}\includegraphics[width=1em]{element-unbounded.pdf}~useType & \hfuzz=500pt \hyperref[gnssType]{gnssType} & \hfuzz=500pt only use observations that match any of these patterns\\
\hfuzz=500pt\includegraphics[width=1em]{connector.pdf}\includegraphics[width=1em]{element-unbounded.pdf}~ignoreType & \hfuzz=500pt \hyperref[gnssType]{gnssType} & \hfuzz=500pt ignore observations that match any of these patterns\\
\hfuzz=500pt\includegraphics[width=1em]{connector.pdf}\includegraphics[width=1em]{element-mustset.pdf}~noAntennaPatternFound & \hfuzz=500pt choice & \hfuzz=500pt what should happen if no antenna pattern is found for an observation\\
\hfuzz=500pt\quad\includegraphics[width=1em]{connector.pdf}\includegraphics[width=1em]{element-mustset.pdf}~ignoreObservation & \hfuzz=500pt  & \hfuzz=500pt ignore observation if no matching pattern is found\\
\hfuzz=500pt\quad\includegraphics[width=1em]{connector.pdf}\includegraphics[width=1em]{element-mustset.pdf}~useNearestFrequency & \hfuzz=500pt  & \hfuzz=500pt use pattern of nearest frequency if not matching pattern is found\\
\hfuzz=500pt\quad\includegraphics[width=1em]{connector.pdf}\includegraphics[width=1em]{element-mustset.pdf}~throwException & \hfuzz=500pt  & \hfuzz=500pt throw exception if no matching pattern is found\\
\hfuzz=500pt\includegraphics[width=1em]{connector.pdf}\includegraphics[width=1em]{element.pdf}~huber & \hfuzz=500pt double & \hfuzz=500pt residuals \$>\$ huber*sigma0 are downweighted\\
\hfuzz=500pt\includegraphics[width=1em]{connector.pdf}\includegraphics[width=1em]{element.pdf}~huberPower & \hfuzz=500pt double & \hfuzz=500pt residuals \$>\$ huber: sigma=(e/huber)\textasciicircum{}huberPower*sigma0\\
\hfuzz=500pt\includegraphics[width=1em]{connector.pdf}\includegraphics[width=1em]{element.pdf}~codeMaxPositionDiff & \hfuzz=500pt double & \hfuzz=500pt [m] max. allowed position error by PPP code only clock error estimation\\
\hfuzz=500pt\includegraphics[width=1em]{connector.pdf}\includegraphics[width=1em]{element.pdf}~minObsCountPerTrack & \hfuzz=500pt uint & \hfuzz=500pt tracks with less number of epochs with observations are dropped\\
\hfuzz=500pt\includegraphics[width=1em]{connector.pdf}\includegraphics[width=1em]{element.pdf}~windowROTI & \hfuzz=500pt double & \hfuzz=500pt [seconds] moving time window to compute Rate Of Tec Index (ROTI) for each observation\\
\hline
\end{tabularx}

\clearpage
%==================================

\section{GnssParametrizationTransmitter}\label{gnssParametrizationTransmitterType}

Definition and parametrization of GNSS satelllite constellations. For precise point positioning (PPP) the
constellations are held fixed and no additional parameters are set up. In case of
\reference{GNSS satellite orbit determination}{cookbook.gnssNetwork} orbits, clock errors, and signal biases
are usually estimated.

A list of satellite PRNs (i.e for GPS: G01, G02, G03, ...) must be provided via
\configFile{inputfileTransmitterList}{stringList}. Satellite system codes follow the
\href{https://files.igs.org/pub/data/format/rinex305.pdf}{RINEX 3 definition}, see \reference{GnssType}{gnssType}.
All input and output files except \configFile{inputfileAntennaDefinition}{gnssAntennaDefinition},
\configFile{inputfileReceiverDefinition}{gnssReceiverDefinition}, and
\configFile{outputfileUsedTransmitterList}{stringList} are read/written for each satellite.
The file name is interpreted as a template with the variable \verb|{prn}| being replaced by the satellite PRN.

In case of PPP precise orbits must be provided via \configFile{inputfileVariational}{variationalEquation}.
Otherwise, dynamic orbits can be estimated by setting \config{estimateOrbit:dynamic}, with
\configFile{inputfileVariational}{variationalEquation} coming from
\reference{orbit integration}{cookbook.gnssNetwork:orbitIntegration}.
The \configClass{parametrizationAcceleration}{parametrizationAccelerationType} must include at least those
parameters that were estimated in \program{PreprocessingVariationalEquationOrbitFit}.
Additional \config{stochasticPulse} parameters can be set up to reduce orbit mismodeling.

Precise satellite clock errors have to be provided via \configFile{inputfileClock}{instrument} in case of PPP.
Otherwise, they can be estimated via \config{estimateClockError}, either \config{epochWise} or based
on a \config{noiseModel}. In this case it is enough to provide approximate clock errors (e.g. from broadcast ephemeris)
via \configFile{inputfileClock}{instrument}. The clock datum can be defined via a zero-mean constraint over all satellites
of a constellation using \config{sigmaZeroMeanConstraint}. This constraint is applied per epoch.

Signal biases for code and phase observations have to be provided via \configFile{inputfileSignalBias}{gnssSignalBias}
in case of PPP. Otherwise they can be estimated by setting \config{estimateCodeBias}=\verb|yes| and
\config{estimatePhaseBias}=\verb|yes|, in which case no a priori signal biases are needed.

Antenna center offsets and variations can be estimated by defining \configClass{antennaCenter}{gnssParametrizationAntennaType}.
It is possible to estimate patterns for each individual satellite or grouped by satellite block.
Since it is not possible to estimate, for example, antenna center offets based on a single day, they are usually hard-constrained
to their a priori value. This can be done by adding \configClass{constraints}{gnssParametrizationConstraintsType} in \program{GnssProcessing}.

Input files used in all systems:
\begin{itemize}
  \item \configFile{inputfileTransmitterInfo}{gnssStationInfo}:
        Created via \program{GnssAntex2AntennaDefinition} or \program{GnssStationInfoCreate}.
  \item \configFile{inputfileAntennaDefinition}{gnssAntennaDefinition}:
        Created via \program{GnssAntex2AntennaDefinition} or \program{GnssAntennaDefinitionCreate}.
  \item \configFile{inputfileReceiverDefinition}{gnssReceiverDefinition}:
        Created via \program{GnssReceiverDefinitionCreate} in case you want to define which signal
        types a satellite transmits.
  \item \configFile{inputfileSignalBias}{gnssSignalBias}: Converted via (TODO: SinexBias2GnssSignalBias) or estimated via \program{GnssProcessing}.
  \item \configFile{inputfileVariational}{variationalEquation}:
        Integrated via \program{PreprocessingVariationalEquation} and \program{PreprocessingVariationalEquationOrbitFit},
        created via (TODO: GnssVariationalCreate), or output of \program{GnssProcessing}.
  \item \configFile{inputfileClock}{instrument}:
        Converted via \program{GnssClockRinex2InstrumentClock} or \program{GnssRinexNavigation2OrbitClock} or
        output of \program{GnssProcessing}.
\end{itemize}

See also \program{GnssProcessing} and \program{GnssSimulateReceiver}.


\subsection{GPS}\label{gnssParametrizationTransmitterType:gps}

Since GPS block IIF satellites are affected by time-variable signal biases
(see Montenbruck et al. 2011, DOI: \href{https://doi.org/10.1007/s10291-011-0232-x}{10.1007/s10291-011-0232-x}),
a \config{biasModel} can be set up for phase observations on the L5 frequency
by defining \configClass{type}{gnssType}=\verb|L5*G| and an appropriate
\configClass{parametrizationAcceleration}{parametrizationAccelerationType}.
The resulting \configFile{outputfileBias}{instrument} contains the time-variable part of the
respective bias evaluated at each epoch, while the constant part is included in
\configFile{outputfileSignalBias}{gnssSignalBias}. In case of PPP the time-variable
part can then be provided via \config{timeVariableBias}.


\keepXColumns
\begin{tabularx}{\textwidth}{N T A}
\hline
Name & Type & Annotation\\
\hline
\hfuzz=500pt\includegraphics[width=1em]{element-mustset.pdf}~inputfileTransmitterList & \hfuzz=500pt filename & \hfuzz=500pt ascii file with transmitter PRNs, used to loop variable \{prn\}\\
\hfuzz=500pt\includegraphics[width=1em]{element-mustset.pdf}~inputfileTransmitterInfo & \hfuzz=500pt filename & \hfuzz=500pt variable \{prn\} available\\
\hfuzz=500pt\includegraphics[width=1em]{element-mustset.pdf}~inputfileAntennaDefintion & \hfuzz=500pt filename & \hfuzz=500pt phase centers and variations (ANTEX)\\
\hfuzz=500pt\includegraphics[width=1em]{element.pdf}~inputfileReceiverDefintion & \hfuzz=500pt filename & \hfuzz=500pt transmitted signal types\\
\hfuzz=500pt\includegraphics[width=1em]{element.pdf}~inputfileSignalBias & \hfuzz=500pt filename & \hfuzz=500pt variable \{prn\} available\\
\hfuzz=500pt\includegraphics[width=1em]{element-mustset.pdf}~inputfileVariational & \hfuzz=500pt filename & \hfuzz=500pt variable \{prn\} available\\
\hfuzz=500pt\includegraphics[width=1em]{element-mustset.pdf}~inputfileClock & \hfuzz=500pt filename & \hfuzz=500pt variable \{prn\} available\\
\hfuzz=500pt\includegraphics[width=1em]{element.pdf}~outputfileSignalBias & \hfuzz=500pt filename & \hfuzz=500pt variable \{prn\} available\\
\hfuzz=500pt\includegraphics[width=1em]{element.pdf}~outputfileUsedTransmitterList & \hfuzz=500pt filename & \hfuzz=500pt ascii file with PRNs\\
\hfuzz=500pt\includegraphics[width=1em]{element.pdf}~estimateOrbit & \hfuzz=500pt choice & \hfuzz=500pt \\
\hfuzz=500pt\includegraphics[width=1em]{connector.pdf}\includegraphics[width=1em]{element-mustset.pdf}~dynamic & \hfuzz=500pt sequence & \hfuzz=500pt initial state and model parameters are estimated (in CRF)\\
\hfuzz=500pt\quad\includegraphics[width=1em]{connector.pdf}\includegraphics[width=1em]{element.pdf}~outputfileVariational & \hfuzz=500pt filename & \hfuzz=500pt variable \{prn\} available\\
\hfuzz=500pt\quad\includegraphics[width=1em]{connector.pdf}\includegraphics[width=1em]{element.pdf}~outputfileOrbit & \hfuzz=500pt filename & \hfuzz=500pt variable \{prn\} available\\
\hfuzz=500pt\quad\includegraphics[width=1em]{connector.pdf}\includegraphics[width=1em]{element.pdf}~outputfileParameters & \hfuzz=500pt filename & \hfuzz=500pt variable \{prn\} available\\
\hfuzz=500pt\quad\includegraphics[width=1em]{connector.pdf}\includegraphics[width=1em]{element-unbounded.pdf}~parametrizationAcceleration & \hfuzz=500pt \hyperref[parametrizationAccelerationType]{parametrizationAcceleration} & \hfuzz=500pt orbit force parameters\\
\hfuzz=500pt\quad\includegraphics[width=1em]{connector.pdf}\includegraphics[width=1em]{element.pdf}~stochasticPulse & \hfuzz=500pt sequence & \hfuzz=500pt [mum/s] parametrization of stochastic pulses\\
\hfuzz=500pt\quad\quad\includegraphics[width=1em]{connector.pdf}\includegraphics[width=1em]{element-mustset-unbounded.pdf}~timeSeries & \hfuzz=500pt \hyperref[timeSeriesType]{timeSeries} & \hfuzz=500pt \\
\hfuzz=500pt\quad\quad\includegraphics[width=1em]{connector.pdf}\includegraphics[width=1em]{element.pdf}~onlyEclipsing & \hfuzz=500pt boolean & \hfuzz=500pt only apply stochastic pulses to eclipsing satellites\\
\hfuzz=500pt\quad\includegraphics[width=1em]{connector.pdf}\includegraphics[width=1em]{element-mustset.pdf}~minEstimableEpochsRatio & \hfuzz=500pt double & \hfuzz=500pt drop satellites with lower ratio of estimable epochs to total epochs\\
\hfuzz=500pt\includegraphics[width=1em]{element.pdf}~estimateCodeBias & \hfuzz=500pt boolean & \hfuzz=500pt \\
\hfuzz=500pt\includegraphics[width=1em]{element.pdf}~estimatePhaseBias & \hfuzz=500pt boolean & \hfuzz=500pt \\
\hfuzz=500pt\includegraphics[width=1em]{element.pdf}~supportsIntegerAmbiguities & \hfuzz=500pt boolean & \hfuzz=500pt receiver tracks full cycle integer ambiguities\\
\hfuzz=500pt\includegraphics[width=1em]{element-unbounded.pdf}~biasModel & \hfuzz=500pt sequence & \hfuzz=500pt estimate time-variable signal bias\\
\hfuzz=500pt\includegraphics[width=1em]{connector.pdf}\includegraphics[width=1em]{element.pdf}~outputfileBias & \hfuzz=500pt filename & \hfuzz=500pt variable \{prn\} available\\
\hfuzz=500pt\includegraphics[width=1em]{connector.pdf}\includegraphics[width=1em]{element-mustset.pdf}~type & \hfuzz=500pt \hyperref[gnssType]{gnssType} & \hfuzz=500pt \\
\hfuzz=500pt\includegraphics[width=1em]{connector.pdf}\includegraphics[width=1em]{element-mustset-unbounded.pdf}~parametrizationTemporal & \hfuzz=500pt \hyperref[parametrizationTemporalType]{parametrizationTemporal} & \hfuzz=500pt \\
\hfuzz=500pt\includegraphics[width=1em]{element-unbounded.pdf}~timeVariableBias & \hfuzz=500pt sequence & \hfuzz=500pt a priori time-variable signal bias\\
\hfuzz=500pt\includegraphics[width=1em]{connector.pdf}\includegraphics[width=1em]{element-mustset.pdf}~inputfileSignalBias & \hfuzz=500pt filename & \hfuzz=500pt variable \{prn\} available; instrument file with bias [m]\\
\hfuzz=500pt\includegraphics[width=1em]{connector.pdf}\includegraphics[width=1em]{element-mustset.pdf}~type & \hfuzz=500pt \hyperref[gnssType]{gnssType} & \hfuzz=500pt bias type\\
\hfuzz=500pt\includegraphics[width=1em]{element.pdf}~estimateClockError & \hfuzz=500pt choice & \hfuzz=500pt \\
\hfuzz=500pt\includegraphics[width=1em]{connector.pdf}\includegraphics[width=1em]{element-mustset.pdf}~epochWise & \hfuzz=500pt sequence & \hfuzz=500pt clock error is estimated for each epoch\\
\hfuzz=500pt\quad\includegraphics[width=1em]{connector.pdf}\includegraphics[width=1em]{element.pdf}~outputfileClock & \hfuzz=500pt filename & \hfuzz=500pt variable \{prn\} available\\
\hfuzz=500pt\quad\includegraphics[width=1em]{connector.pdf}\includegraphics[width=1em]{element.pdf}~sigmaZeroMeanConstraint & \hfuzz=500pt double & \hfuzz=500pt (0 = unconstrained) sigma [m] for zero-mean constraint over all satellite clocks\\
\hfuzz=500pt\includegraphics[width=1em]{connector.pdf}\includegraphics[width=1em]{element-mustset.pdf}~noiseModel & \hfuzz=500pt sequence & \hfuzz=500pt clock errors are estimated based on a noise model\\
\hfuzz=500pt\quad\includegraphics[width=1em]{connector.pdf}\includegraphics[width=1em]{element.pdf}~outputfileClock & \hfuzz=500pt filename & \hfuzz=500pt variable \{prn\} available\\
\hfuzz=500pt\quad\includegraphics[width=1em]{connector.pdf}\includegraphics[width=1em]{element-mustset.pdf}~inputfileClockNoiseModel & \hfuzz=500pt filename & \hfuzz=500pt variable \{prn\} available\\
\hfuzz=500pt\quad\includegraphics[width=1em]{connector.pdf}\includegraphics[width=1em]{element.pdf}~estimateVarianceFactor & \hfuzz=500pt boolean & \hfuzz=500pt via Variance Component Estimation\\
\hfuzz=500pt\quad\includegraphics[width=1em]{connector.pdf}\includegraphics[width=1em]{element-unbounded.pdf}~clockModel & \hfuzz=500pt \hyperref[parametrizationTemporalType]{parametrizationTemporal} & \hfuzz=500pt \\
\hfuzz=500pt\quad\includegraphics[width=1em]{connector.pdf}\includegraphics[width=1em]{element.pdf}~sigmaZeroMeanConstraint & \hfuzz=500pt double & \hfuzz=500pt (0 = unconstrained) sigma [m] for zero-mean constraint over all satellite clocks\\
\hfuzz=500pt\includegraphics[width=1em]{element-unbounded.pdf}~antennaCenter & \hfuzz=500pt \hyperref[gnssParametrizationAntennaType]{gnssParametrizationAntenna} & \hfuzz=500pt estimate antenna center variations\\
\hfuzz=500pt\includegraphics[width=1em]{element-mustset.pdf}~integrationDegree & \hfuzz=500pt uint & \hfuzz=500pt integration of forces by polynomial approximation of degree n\\
\hfuzz=500pt\includegraphics[width=1em]{element-mustset.pdf}~interpolationDegree & \hfuzz=500pt uint & \hfuzz=500pt for orbit interpolation and velocity calculation\\
\hfuzz=500pt\includegraphics[width=1em]{element.pdf}~disableShadowEpochs & \hfuzz=500pt boolean & \hfuzz=500pt disable epochs if satellite is in Earth's/Moon's shadow\\
\hfuzz=500pt\includegraphics[width=1em]{element.pdf}~disablePostShadowRecoveryEpochs & \hfuzz=500pt boolean & \hfuzz=500pt disable epochs if satellite is in post-shadow recovery maneuver (e.g. GPS block IIA)\\
\hfuzz=500pt\includegraphics[width=1em]{element-mustset.pdf}~ephemerides & \hfuzz=500pt \hyperref[ephemeridesType]{ephemerides} & \hfuzz=500pt \\
\hfuzz=500pt\includegraphics[width=1em]{element-mustset.pdf}~eclipse & \hfuzz=500pt \hyperref[eclipseType]{eclipse} & \hfuzz=500pt eclipse model used to determine if a satellite is in Earth's shadow\\
\hfuzz=500pt\includegraphics[width=1em]{element-mustset.pdf}~noAntennaPatternFound & \hfuzz=500pt choice & \hfuzz=500pt what should happen is no antenna pattern is found for an observation\\
\hfuzz=500pt\includegraphics[width=1em]{connector.pdf}\includegraphics[width=1em]{element-mustset.pdf}~ignoreObservation & \hfuzz=500pt  & \hfuzz=500pt ignore observation if no matching pattern is found\\
\hfuzz=500pt\includegraphics[width=1em]{connector.pdf}\includegraphics[width=1em]{element-mustset.pdf}~useNearestFrequency & \hfuzz=500pt  & \hfuzz=500pt use pattern of nearest frequency if no matching pattern is found\\
\hfuzz=500pt\includegraphics[width=1em]{connector.pdf}\includegraphics[width=1em]{element-mustset.pdf}~throwException & \hfuzz=500pt  & \hfuzz=500pt throw exception if no matching pattern is found\\
\hline
\end{tabularx}


\subsection{GLONASS}\label{gnssParametrizationTransmitterType:glonass}


\keepXColumns
\begin{tabularx}{\textwidth}{N T A}
\hline
Name & Type & Annotation\\
\hline
\hfuzz=500pt\includegraphics[width=1em]{element-mustset.pdf}~inputfileTransmitterList & \hfuzz=500pt filename & \hfuzz=500pt ascii file with transmitter PRNs, used to loop variable \{prn\}\\
\hfuzz=500pt\includegraphics[width=1em]{element-mustset.pdf}~inputfileTransmitterInfo & \hfuzz=500pt filename & \hfuzz=500pt variable \{prn\} available\\
\hfuzz=500pt\includegraphics[width=1em]{element-mustset.pdf}~inputfileAntennaDefintion & \hfuzz=500pt filename & \hfuzz=500pt phase centers and variations (ANTEX)\\
\hfuzz=500pt\includegraphics[width=1em]{element.pdf}~inputfileReceiverDefintion & \hfuzz=500pt filename & \hfuzz=500pt transmitted signal types\\
\hfuzz=500pt\includegraphics[width=1em]{element.pdf}~inputfileSignalBias & \hfuzz=500pt filename & \hfuzz=500pt variable \{prn\} available\\
\hfuzz=500pt\includegraphics[width=1em]{element-mustset.pdf}~inputfileVariational & \hfuzz=500pt filename & \hfuzz=500pt variable \{prn\} available\\
\hfuzz=500pt\includegraphics[width=1em]{element-mustset.pdf}~inputfileClock & \hfuzz=500pt filename & \hfuzz=500pt variable \{prn\} available\\
\hfuzz=500pt\includegraphics[width=1em]{element.pdf}~outputfileSignalBias & \hfuzz=500pt filename & \hfuzz=500pt variable \{prn\} available\\
\hfuzz=500pt\includegraphics[width=1em]{element.pdf}~outputfileUsedTransmitterList & \hfuzz=500pt filename & \hfuzz=500pt ascii file with PRNs\\
\hfuzz=500pt\includegraphics[width=1em]{element.pdf}~estimateOrbit & \hfuzz=500pt choice & \hfuzz=500pt \\
\hfuzz=500pt\includegraphics[width=1em]{connector.pdf}\includegraphics[width=1em]{element-mustset.pdf}~dynamic & \hfuzz=500pt sequence & \hfuzz=500pt initial state and model parameters are estimated (in CRF)\\
\hfuzz=500pt\quad\includegraphics[width=1em]{connector.pdf}\includegraphics[width=1em]{element.pdf}~outputfileVariational & \hfuzz=500pt filename & \hfuzz=500pt variable \{prn\} available\\
\hfuzz=500pt\quad\includegraphics[width=1em]{connector.pdf}\includegraphics[width=1em]{element.pdf}~outputfileOrbit & \hfuzz=500pt filename & \hfuzz=500pt variable \{prn\} available\\
\hfuzz=500pt\quad\includegraphics[width=1em]{connector.pdf}\includegraphics[width=1em]{element.pdf}~outputfileParameters & \hfuzz=500pt filename & \hfuzz=500pt variable \{prn\} available\\
\hfuzz=500pt\quad\includegraphics[width=1em]{connector.pdf}\includegraphics[width=1em]{element-unbounded.pdf}~parametrizationAcceleration & \hfuzz=500pt \hyperref[parametrizationAccelerationType]{parametrizationAcceleration} & \hfuzz=500pt orbit force parameters\\
\hfuzz=500pt\quad\includegraphics[width=1em]{connector.pdf}\includegraphics[width=1em]{element.pdf}~stochasticPulse & \hfuzz=500pt sequence & \hfuzz=500pt [mum/s] parametrization of stochastic pulses\\
\hfuzz=500pt\quad\quad\includegraphics[width=1em]{connector.pdf}\includegraphics[width=1em]{element-mustset-unbounded.pdf}~timeSeries & \hfuzz=500pt \hyperref[timeSeriesType]{timeSeries} & \hfuzz=500pt \\
\hfuzz=500pt\quad\quad\includegraphics[width=1em]{connector.pdf}\includegraphics[width=1em]{element.pdf}~onlyEclipsing & \hfuzz=500pt boolean & \hfuzz=500pt only apply stochastic pulses to eclipsing satellites\\
\hfuzz=500pt\quad\includegraphics[width=1em]{connector.pdf}\includegraphics[width=1em]{element-mustset.pdf}~minEstimableEpochsRatio & \hfuzz=500pt double & \hfuzz=500pt drop satellites with lower ratio of estimable epochs to total epochs\\
\hfuzz=500pt\includegraphics[width=1em]{element.pdf}~estimateCodeBias & \hfuzz=500pt boolean & \hfuzz=500pt \\
\hfuzz=500pt\includegraphics[width=1em]{element.pdf}~estimatePhaseBias & \hfuzz=500pt boolean & \hfuzz=500pt \\
\hfuzz=500pt\includegraphics[width=1em]{element.pdf}~supportsIntegerAmbiguities & \hfuzz=500pt boolean & \hfuzz=500pt receiver tracks full cycle integer ambiguities\\
\hfuzz=500pt\includegraphics[width=1em]{element-unbounded.pdf}~biasModel & \hfuzz=500pt sequence & \hfuzz=500pt estimate time-variable signal bias\\
\hfuzz=500pt\includegraphics[width=1em]{connector.pdf}\includegraphics[width=1em]{element.pdf}~outputfileBias & \hfuzz=500pt filename & \hfuzz=500pt variable \{prn\} available\\
\hfuzz=500pt\includegraphics[width=1em]{connector.pdf}\includegraphics[width=1em]{element-mustset.pdf}~type & \hfuzz=500pt \hyperref[gnssType]{gnssType} & \hfuzz=500pt \\
\hfuzz=500pt\includegraphics[width=1em]{connector.pdf}\includegraphics[width=1em]{element-mustset-unbounded.pdf}~parametrizationTemporal & \hfuzz=500pt \hyperref[parametrizationTemporalType]{parametrizationTemporal} & \hfuzz=500pt \\
\hfuzz=500pt\includegraphics[width=1em]{element-unbounded.pdf}~timeVariableBias & \hfuzz=500pt sequence & \hfuzz=500pt a priori time-variable signal bias\\
\hfuzz=500pt\includegraphics[width=1em]{connector.pdf}\includegraphics[width=1em]{element-mustset.pdf}~inputfileSignalBias & \hfuzz=500pt filename & \hfuzz=500pt variable \{prn\} available; instrument file with bias [m]\\
\hfuzz=500pt\includegraphics[width=1em]{connector.pdf}\includegraphics[width=1em]{element-mustset.pdf}~type & \hfuzz=500pt \hyperref[gnssType]{gnssType} & \hfuzz=500pt bias type\\
\hfuzz=500pt\includegraphics[width=1em]{element.pdf}~estimateClockError & \hfuzz=500pt choice & \hfuzz=500pt \\
\hfuzz=500pt\includegraphics[width=1em]{connector.pdf}\includegraphics[width=1em]{element-mustset.pdf}~epochWise & \hfuzz=500pt sequence & \hfuzz=500pt clock error is estimated for each epoch\\
\hfuzz=500pt\quad\includegraphics[width=1em]{connector.pdf}\includegraphics[width=1em]{element.pdf}~outputfileClock & \hfuzz=500pt filename & \hfuzz=500pt variable \{prn\} available\\
\hfuzz=500pt\quad\includegraphics[width=1em]{connector.pdf}\includegraphics[width=1em]{element.pdf}~sigmaZeroMeanConstraint & \hfuzz=500pt double & \hfuzz=500pt (0 = unconstrained) sigma [m] for zero-mean constraint over all satellite clocks\\
\hfuzz=500pt\includegraphics[width=1em]{connector.pdf}\includegraphics[width=1em]{element-mustset.pdf}~noiseModel & \hfuzz=500pt sequence & \hfuzz=500pt clock errors are estimated based on a noise model\\
\hfuzz=500pt\quad\includegraphics[width=1em]{connector.pdf}\includegraphics[width=1em]{element.pdf}~outputfileClock & \hfuzz=500pt filename & \hfuzz=500pt variable \{prn\} available\\
\hfuzz=500pt\quad\includegraphics[width=1em]{connector.pdf}\includegraphics[width=1em]{element-mustset.pdf}~inputfileClockNoiseModel & \hfuzz=500pt filename & \hfuzz=500pt variable \{prn\} available\\
\hfuzz=500pt\quad\includegraphics[width=1em]{connector.pdf}\includegraphics[width=1em]{element.pdf}~estimateVarianceFactor & \hfuzz=500pt boolean & \hfuzz=500pt via Variance Component Estimation\\
\hfuzz=500pt\quad\includegraphics[width=1em]{connector.pdf}\includegraphics[width=1em]{element-unbounded.pdf}~clockModel & \hfuzz=500pt \hyperref[parametrizationTemporalType]{parametrizationTemporal} & \hfuzz=500pt \\
\hfuzz=500pt\quad\includegraphics[width=1em]{connector.pdf}\includegraphics[width=1em]{element.pdf}~sigmaZeroMeanConstraint & \hfuzz=500pt double & \hfuzz=500pt (0 = unconstrained) sigma [m] for zero-mean constraint over all satellite clocks\\
\hfuzz=500pt\includegraphics[width=1em]{element-unbounded.pdf}~antennaCenter & \hfuzz=500pt \hyperref[gnssParametrizationAntennaType]{gnssParametrizationAntenna} & \hfuzz=500pt estimate antenna center variations\\
\hfuzz=500pt\includegraphics[width=1em]{element.pdf}~integrationDegree & \hfuzz=500pt uint & \hfuzz=500pt integration of forces by polynomial approximation of degree n\\
\hfuzz=500pt\includegraphics[width=1em]{element.pdf}~interpolationDegree & \hfuzz=500pt uint & \hfuzz=500pt for orbit interpolation and velocity calculation\\
\hfuzz=500pt\includegraphics[width=1em]{element.pdf}~disableShadowEpochs & \hfuzz=500pt boolean & \hfuzz=500pt disable epochs if satellite is in Earth's/Moon's shadow\\
\hfuzz=500pt\includegraphics[width=1em]{element-mustset.pdf}~ephemerides & \hfuzz=500pt \hyperref[ephemeridesType]{ephemerides} & \hfuzz=500pt \\
\hfuzz=500pt\includegraphics[width=1em]{element-mustset.pdf}~eclipse & \hfuzz=500pt \hyperref[eclipseType]{eclipse} & \hfuzz=500pt eclipse model used to determine if a satellite is in Earth's shadow\\
\hfuzz=500pt\includegraphics[width=1em]{element-mustset.pdf}~noAntennaPatternFound & \hfuzz=500pt choice & \hfuzz=500pt what should happen is no antenna pattern is found for an observation\\
\hfuzz=500pt\includegraphics[width=1em]{connector.pdf}\includegraphics[width=1em]{element-mustset.pdf}~ignoreObservation & \hfuzz=500pt  & \hfuzz=500pt ignore observation if no matching pattern is found\\
\hfuzz=500pt\includegraphics[width=1em]{connector.pdf}\includegraphics[width=1em]{element-mustset.pdf}~useNearestFrequency & \hfuzz=500pt  & \hfuzz=500pt use pattern of nearest frequency if no matching pattern is found\\
\hfuzz=500pt\includegraphics[width=1em]{connector.pdf}\includegraphics[width=1em]{element-mustset.pdf}~throwException & \hfuzz=500pt  & \hfuzz=500pt throw exception if no matching pattern is found\\
\hline
\end{tabularx}


\subsection{Galileo}\label{gnssParametrizationTransmitterType:galileo}


\keepXColumns
\begin{tabularx}{\textwidth}{N T A}
\hline
Name & Type & Annotation\\
\hline
\hfuzz=500pt\includegraphics[width=1em]{element-mustset.pdf}~inputfileTransmitterList & \hfuzz=500pt filename & \hfuzz=500pt ascii file with transmitter PRNs, used to loop variable \{prn\}\\
\hfuzz=500pt\includegraphics[width=1em]{element-mustset.pdf}~inputfileTransmitterInfo & \hfuzz=500pt filename & \hfuzz=500pt variable \{prn\} available\\
\hfuzz=500pt\includegraphics[width=1em]{element-mustset.pdf}~inputfileAntennaDefintion & \hfuzz=500pt filename & \hfuzz=500pt phase centers and variations (ANTEX)\\
\hfuzz=500pt\includegraphics[width=1em]{element.pdf}~inputfileReceiverDefintion & \hfuzz=500pt filename & \hfuzz=500pt transmitted signal types\\
\hfuzz=500pt\includegraphics[width=1em]{element.pdf}~inputfileSignalBias & \hfuzz=500pt filename & \hfuzz=500pt variable \{prn\} available\\
\hfuzz=500pt\includegraphics[width=1em]{element-mustset.pdf}~inputfileVariational & \hfuzz=500pt filename & \hfuzz=500pt variable \{prn\} available\\
\hfuzz=500pt\includegraphics[width=1em]{element-mustset.pdf}~inputfileClock & \hfuzz=500pt filename & \hfuzz=500pt variable \{prn\} available\\
\hfuzz=500pt\includegraphics[width=1em]{element.pdf}~outputfileSignalBias & \hfuzz=500pt filename & \hfuzz=500pt variable \{prn\} available\\
\hfuzz=500pt\includegraphics[width=1em]{element.pdf}~outputfileUsedTransmitterList & \hfuzz=500pt filename & \hfuzz=500pt ascii file with PRNs\\
\hfuzz=500pt\includegraphics[width=1em]{element.pdf}~estimateOrbit & \hfuzz=500pt choice & \hfuzz=500pt \\
\hfuzz=500pt\includegraphics[width=1em]{connector.pdf}\includegraphics[width=1em]{element-mustset.pdf}~dynamic & \hfuzz=500pt sequence & \hfuzz=500pt initial state and model parameters are estimated (in CRF)\\
\hfuzz=500pt\quad\includegraphics[width=1em]{connector.pdf}\includegraphics[width=1em]{element.pdf}~outputfileVariational & \hfuzz=500pt filename & \hfuzz=500pt variable \{prn\} available\\
\hfuzz=500pt\quad\includegraphics[width=1em]{connector.pdf}\includegraphics[width=1em]{element.pdf}~outputfileOrbit & \hfuzz=500pt filename & \hfuzz=500pt variable \{prn\} available\\
\hfuzz=500pt\quad\includegraphics[width=1em]{connector.pdf}\includegraphics[width=1em]{element.pdf}~outputfileParameters & \hfuzz=500pt filename & \hfuzz=500pt variable \{prn\} available\\
\hfuzz=500pt\quad\includegraphics[width=1em]{connector.pdf}\includegraphics[width=1em]{element-unbounded.pdf}~parametrizationAcceleration & \hfuzz=500pt \hyperref[parametrizationAccelerationType]{parametrizationAcceleration} & \hfuzz=500pt orbit force parameters\\
\hfuzz=500pt\quad\includegraphics[width=1em]{connector.pdf}\includegraphics[width=1em]{element.pdf}~stochasticPulse & \hfuzz=500pt sequence & \hfuzz=500pt [mum/s] parametrization of stochastic pulses\\
\hfuzz=500pt\quad\quad\includegraphics[width=1em]{connector.pdf}\includegraphics[width=1em]{element-mustset-unbounded.pdf}~timeSeries & \hfuzz=500pt \hyperref[timeSeriesType]{timeSeries} & \hfuzz=500pt \\
\hfuzz=500pt\quad\quad\includegraphics[width=1em]{connector.pdf}\includegraphics[width=1em]{element.pdf}~onlyEclipsing & \hfuzz=500pt boolean & \hfuzz=500pt only apply stochastic pulses to eclipsing satellites\\
\hfuzz=500pt\quad\includegraphics[width=1em]{connector.pdf}\includegraphics[width=1em]{element-mustset.pdf}~minEstimableEpochsRatio & \hfuzz=500pt double & \hfuzz=500pt drop satellites with lower ratio of estimable epochs to total epochs\\
\hfuzz=500pt\includegraphics[width=1em]{element.pdf}~estimateCodeBias & \hfuzz=500pt boolean & \hfuzz=500pt \\
\hfuzz=500pt\includegraphics[width=1em]{element.pdf}~estimatePhaseBias & \hfuzz=500pt boolean & \hfuzz=500pt \\
\hfuzz=500pt\includegraphics[width=1em]{element.pdf}~supportsIntegerAmbiguities & \hfuzz=500pt boolean & \hfuzz=500pt receiver tracks full cycle integer ambiguities\\
\hfuzz=500pt\includegraphics[width=1em]{element-unbounded.pdf}~biasModel & \hfuzz=500pt sequence & \hfuzz=500pt estimate time-variable signal bias\\
\hfuzz=500pt\includegraphics[width=1em]{connector.pdf}\includegraphics[width=1em]{element.pdf}~outputfileBias & \hfuzz=500pt filename & \hfuzz=500pt variable \{prn\} available\\
\hfuzz=500pt\includegraphics[width=1em]{connector.pdf}\includegraphics[width=1em]{element-mustset.pdf}~type & \hfuzz=500pt \hyperref[gnssType]{gnssType} & \hfuzz=500pt \\
\hfuzz=500pt\includegraphics[width=1em]{connector.pdf}\includegraphics[width=1em]{element-mustset-unbounded.pdf}~parametrizationTemporal & \hfuzz=500pt \hyperref[parametrizationTemporalType]{parametrizationTemporal} & \hfuzz=500pt \\
\hfuzz=500pt\includegraphics[width=1em]{element-unbounded.pdf}~timeVariableBias & \hfuzz=500pt sequence & \hfuzz=500pt a priori time-variable signal bias\\
\hfuzz=500pt\includegraphics[width=1em]{connector.pdf}\includegraphics[width=1em]{element-mustset.pdf}~inputfileSignalBias & \hfuzz=500pt filename & \hfuzz=500pt variable \{prn\} available; instrument file with bias [m]\\
\hfuzz=500pt\includegraphics[width=1em]{connector.pdf}\includegraphics[width=1em]{element-mustset.pdf}~type & \hfuzz=500pt \hyperref[gnssType]{gnssType} & \hfuzz=500pt bias type\\
\hfuzz=500pt\includegraphics[width=1em]{element.pdf}~estimateClockError & \hfuzz=500pt choice & \hfuzz=500pt \\
\hfuzz=500pt\includegraphics[width=1em]{connector.pdf}\includegraphics[width=1em]{element-mustset.pdf}~epochWise & \hfuzz=500pt sequence & \hfuzz=500pt clock error is estimated for each epoch\\
\hfuzz=500pt\quad\includegraphics[width=1em]{connector.pdf}\includegraphics[width=1em]{element.pdf}~outputfileClock & \hfuzz=500pt filename & \hfuzz=500pt variable \{prn\} available\\
\hfuzz=500pt\quad\includegraphics[width=1em]{connector.pdf}\includegraphics[width=1em]{element.pdf}~sigmaZeroMeanConstraint & \hfuzz=500pt double & \hfuzz=500pt (0 = unconstrained) sigma [m] for zero-mean constraint over all satellite clocks\\
\hfuzz=500pt\includegraphics[width=1em]{connector.pdf}\includegraphics[width=1em]{element-mustset.pdf}~noiseModel & \hfuzz=500pt sequence & \hfuzz=500pt clock errors are estimated based on a noise model\\
\hfuzz=500pt\quad\includegraphics[width=1em]{connector.pdf}\includegraphics[width=1em]{element.pdf}~outputfileClock & \hfuzz=500pt filename & \hfuzz=500pt variable \{prn\} available\\
\hfuzz=500pt\quad\includegraphics[width=1em]{connector.pdf}\includegraphics[width=1em]{element-mustset.pdf}~inputfileClockNoiseModel & \hfuzz=500pt filename & \hfuzz=500pt variable \{prn\} available\\
\hfuzz=500pt\quad\includegraphics[width=1em]{connector.pdf}\includegraphics[width=1em]{element.pdf}~estimateVarianceFactor & \hfuzz=500pt boolean & \hfuzz=500pt via Variance Component Estimation\\
\hfuzz=500pt\quad\includegraphics[width=1em]{connector.pdf}\includegraphics[width=1em]{element-unbounded.pdf}~clockModel & \hfuzz=500pt \hyperref[parametrizationTemporalType]{parametrizationTemporal} & \hfuzz=500pt \\
\hfuzz=500pt\quad\includegraphics[width=1em]{connector.pdf}\includegraphics[width=1em]{element.pdf}~sigmaZeroMeanConstraint & \hfuzz=500pt double & \hfuzz=500pt (0 = unconstrained) sigma [m] for zero-mean constraint over all satellite clocks\\
\hfuzz=500pt\includegraphics[width=1em]{element-unbounded.pdf}~antennaCenter & \hfuzz=500pt \hyperref[gnssParametrizationAntennaType]{gnssParametrizationAntenna} & \hfuzz=500pt estimate antenna center variations\\
\hfuzz=500pt\includegraphics[width=1em]{element-mustset.pdf}~integrationDegree & \hfuzz=500pt uint & \hfuzz=500pt integration of forces by polynomial approximation of degree n\\
\hfuzz=500pt\includegraphics[width=1em]{element-mustset.pdf}~interpolationDegree & \hfuzz=500pt uint & \hfuzz=500pt for orbit interpolation and velocity calculation\\
\hfuzz=500pt\includegraphics[width=1em]{element.pdf}~disableShadowEpochs & \hfuzz=500pt boolean & \hfuzz=500pt disable epochs if satellite is in Earth's/Moon's shadow\\
\hfuzz=500pt\includegraphics[width=1em]{element-mustset.pdf}~ephemerides & \hfuzz=500pt \hyperref[ephemeridesType]{ephemerides} & \hfuzz=500pt \\
\hfuzz=500pt\includegraphics[width=1em]{element-mustset.pdf}~eclipse & \hfuzz=500pt \hyperref[eclipseType]{eclipse} & \hfuzz=500pt eclipse model used to determine if a satellite is in Earth's shadow\\
\hfuzz=500pt\includegraphics[width=1em]{element-mustset.pdf}~noAntennaPatternFound & \hfuzz=500pt choice & \hfuzz=500pt what should happen is no antenna pattern is found for an observation\\
\hfuzz=500pt\includegraphics[width=1em]{connector.pdf}\includegraphics[width=1em]{element-mustset.pdf}~ignoreObservation & \hfuzz=500pt  & \hfuzz=500pt ignore observation if no matching pattern is found\\
\hfuzz=500pt\includegraphics[width=1em]{connector.pdf}\includegraphics[width=1em]{element-mustset.pdf}~useNearestFrequency & \hfuzz=500pt  & \hfuzz=500pt use pattern of nearest frequency if no matching pattern is found\\
\hfuzz=500pt\includegraphics[width=1em]{connector.pdf}\includegraphics[width=1em]{element-mustset.pdf}~throwException & \hfuzz=500pt  & \hfuzz=500pt throw exception if no matching pattern is found\\
\hline
\end{tabularx}


\subsection{BeiDou}\label{gnssParametrizationTransmitterType:beidou}


\keepXColumns
\begin{tabularx}{\textwidth}{N T A}
\hline
Name & Type & Annotation\\
\hline
\hfuzz=500pt\includegraphics[width=1em]{element-mustset.pdf}~inputfileTransmitterList & \hfuzz=500pt filename & \hfuzz=500pt ascii file with transmitter PRNs, used to loop variable \{prn\}\\
\hfuzz=500pt\includegraphics[width=1em]{element-mustset.pdf}~inputfileTransmitterInfo & \hfuzz=500pt filename & \hfuzz=500pt variable \{prn\} available\\
\hfuzz=500pt\includegraphics[width=1em]{element-mustset.pdf}~inputfileAntennaDefintion & \hfuzz=500pt filename & \hfuzz=500pt phase centers and variations (ANTEX)\\
\hfuzz=500pt\includegraphics[width=1em]{element.pdf}~inputfileReceiverDefintion & \hfuzz=500pt filename & \hfuzz=500pt transmitted signal types\\
\hfuzz=500pt\includegraphics[width=1em]{element.pdf}~inputfileSignalBias & \hfuzz=500pt filename & \hfuzz=500pt variable \{prn\} available\\
\hfuzz=500pt\includegraphics[width=1em]{element-mustset.pdf}~inputfileVariational & \hfuzz=500pt filename & \hfuzz=500pt variable \{prn\} available\\
\hfuzz=500pt\includegraphics[width=1em]{element-mustset.pdf}~inputfileClock & \hfuzz=500pt filename & \hfuzz=500pt variable \{prn\} available\\
\hfuzz=500pt\includegraphics[width=1em]{element.pdf}~outputfileSignalBias & \hfuzz=500pt filename & \hfuzz=500pt variable \{prn\} available\\
\hfuzz=500pt\includegraphics[width=1em]{element.pdf}~outputfileUsedTransmitterList & \hfuzz=500pt filename & \hfuzz=500pt ascii file with PRNs\\
\hfuzz=500pt\includegraphics[width=1em]{element.pdf}~estimateOrbit & \hfuzz=500pt choice & \hfuzz=500pt \\
\hfuzz=500pt\includegraphics[width=1em]{connector.pdf}\includegraphics[width=1em]{element-mustset.pdf}~dynamic & \hfuzz=500pt sequence & \hfuzz=500pt initial state and model parameters are estimated (in CRF)\\
\hfuzz=500pt\quad\includegraphics[width=1em]{connector.pdf}\includegraphics[width=1em]{element.pdf}~outputfileVariational & \hfuzz=500pt filename & \hfuzz=500pt variable \{prn\} available\\
\hfuzz=500pt\quad\includegraphics[width=1em]{connector.pdf}\includegraphics[width=1em]{element.pdf}~outputfileOrbit & \hfuzz=500pt filename & \hfuzz=500pt variable \{prn\} available\\
\hfuzz=500pt\quad\includegraphics[width=1em]{connector.pdf}\includegraphics[width=1em]{element.pdf}~outputfileParameters & \hfuzz=500pt filename & \hfuzz=500pt variable \{prn\} available\\
\hfuzz=500pt\quad\includegraphics[width=1em]{connector.pdf}\includegraphics[width=1em]{element-unbounded.pdf}~parametrizationAcceleration & \hfuzz=500pt \hyperref[parametrizationAccelerationType]{parametrizationAcceleration} & \hfuzz=500pt orbit force parameters\\
\hfuzz=500pt\quad\includegraphics[width=1em]{connector.pdf}\includegraphics[width=1em]{element.pdf}~stochasticPulse & \hfuzz=500pt sequence & \hfuzz=500pt [mum/s] parametrization of stochastic pulses\\
\hfuzz=500pt\quad\quad\includegraphics[width=1em]{connector.pdf}\includegraphics[width=1em]{element-mustset-unbounded.pdf}~timeSeries & \hfuzz=500pt \hyperref[timeSeriesType]{timeSeries} & \hfuzz=500pt \\
\hfuzz=500pt\quad\quad\includegraphics[width=1em]{connector.pdf}\includegraphics[width=1em]{element.pdf}~onlyEclipsing & \hfuzz=500pt boolean & \hfuzz=500pt only apply stochastic pulses to eclipsing satellites\\
\hfuzz=500pt\quad\includegraphics[width=1em]{connector.pdf}\includegraphics[width=1em]{element-mustset.pdf}~minEstimableEpochsRatio & \hfuzz=500pt double & \hfuzz=500pt drop satellites with lower ratio of estimable epochs to total epochs\\
\hfuzz=500pt\includegraphics[width=1em]{element.pdf}~estimateCodeBias & \hfuzz=500pt boolean & \hfuzz=500pt \\
\hfuzz=500pt\includegraphics[width=1em]{element.pdf}~estimatePhaseBias & \hfuzz=500pt boolean & \hfuzz=500pt \\
\hfuzz=500pt\includegraphics[width=1em]{element.pdf}~supportsIntegerAmbiguities & \hfuzz=500pt boolean & \hfuzz=500pt receiver tracks full cycle integer ambiguities\\
\hfuzz=500pt\includegraphics[width=1em]{element-unbounded.pdf}~biasModel & \hfuzz=500pt sequence & \hfuzz=500pt estimate time-variable signal bias\\
\hfuzz=500pt\includegraphics[width=1em]{connector.pdf}\includegraphics[width=1em]{element.pdf}~outputfileBias & \hfuzz=500pt filename & \hfuzz=500pt variable \{prn\} available\\
\hfuzz=500pt\includegraphics[width=1em]{connector.pdf}\includegraphics[width=1em]{element-mustset.pdf}~type & \hfuzz=500pt \hyperref[gnssType]{gnssType} & \hfuzz=500pt \\
\hfuzz=500pt\includegraphics[width=1em]{connector.pdf}\includegraphics[width=1em]{element-mustset-unbounded.pdf}~parametrizationTemporal & \hfuzz=500pt \hyperref[parametrizationTemporalType]{parametrizationTemporal} & \hfuzz=500pt \\
\hfuzz=500pt\includegraphics[width=1em]{element-unbounded.pdf}~timeVariableBias & \hfuzz=500pt sequence & \hfuzz=500pt a priori time-variable signal bias\\
\hfuzz=500pt\includegraphics[width=1em]{connector.pdf}\includegraphics[width=1em]{element-mustset.pdf}~inputfileSignalBias & \hfuzz=500pt filename & \hfuzz=500pt variable \{prn\} available; instrument file with bias [m]\\
\hfuzz=500pt\includegraphics[width=1em]{connector.pdf}\includegraphics[width=1em]{element-mustset.pdf}~type & \hfuzz=500pt \hyperref[gnssType]{gnssType} & \hfuzz=500pt bias type\\
\hfuzz=500pt\includegraphics[width=1em]{element.pdf}~estimateClockError & \hfuzz=500pt choice & \hfuzz=500pt \\
\hfuzz=500pt\includegraphics[width=1em]{connector.pdf}\includegraphics[width=1em]{element-mustset.pdf}~epochWise & \hfuzz=500pt sequence & \hfuzz=500pt clock error is estimated for each epoch\\
\hfuzz=500pt\quad\includegraphics[width=1em]{connector.pdf}\includegraphics[width=1em]{element.pdf}~outputfileClock & \hfuzz=500pt filename & \hfuzz=500pt variable \{prn\} available\\
\hfuzz=500pt\quad\includegraphics[width=1em]{connector.pdf}\includegraphics[width=1em]{element.pdf}~sigmaZeroMeanConstraint & \hfuzz=500pt double & \hfuzz=500pt (0 = unconstrained) sigma [m] for zero-mean constraint over all satellite clocks\\
\hfuzz=500pt\includegraphics[width=1em]{connector.pdf}\includegraphics[width=1em]{element-mustset.pdf}~noiseModel & \hfuzz=500pt sequence & \hfuzz=500pt clock errors are estimated based on a noise model\\
\hfuzz=500pt\quad\includegraphics[width=1em]{connector.pdf}\includegraphics[width=1em]{element.pdf}~outputfileClock & \hfuzz=500pt filename & \hfuzz=500pt variable \{prn\} available\\
\hfuzz=500pt\quad\includegraphics[width=1em]{connector.pdf}\includegraphics[width=1em]{element-mustset.pdf}~inputfileClockNoiseModel & \hfuzz=500pt filename & \hfuzz=500pt variable \{prn\} available\\
\hfuzz=500pt\quad\includegraphics[width=1em]{connector.pdf}\includegraphics[width=1em]{element.pdf}~estimateVarianceFactor & \hfuzz=500pt boolean & \hfuzz=500pt via Variance Component Estimation\\
\hfuzz=500pt\quad\includegraphics[width=1em]{connector.pdf}\includegraphics[width=1em]{element-unbounded.pdf}~clockModel & \hfuzz=500pt \hyperref[parametrizationTemporalType]{parametrizationTemporal} & \hfuzz=500pt \\
\hfuzz=500pt\quad\includegraphics[width=1em]{connector.pdf}\includegraphics[width=1em]{element.pdf}~sigmaZeroMeanConstraint & \hfuzz=500pt double & \hfuzz=500pt (0 = unconstrained) sigma [m] for zero-mean constraint over all satellite clocks\\
\hfuzz=500pt\includegraphics[width=1em]{element-unbounded.pdf}~antennaCenter & \hfuzz=500pt \hyperref[gnssParametrizationAntennaType]{gnssParametrizationAntenna} & \hfuzz=500pt estimate antenna center variations\\
\hfuzz=500pt\includegraphics[width=1em]{element-mustset.pdf}~integrationDegree & \hfuzz=500pt uint & \hfuzz=500pt integration of forces by polynomial approximation of degree n\\
\hfuzz=500pt\includegraphics[width=1em]{element-mustset.pdf}~interpolationDegree & \hfuzz=500pt uint & \hfuzz=500pt for orbit interpolation and velocity calculation\\
\hfuzz=500pt\includegraphics[width=1em]{element.pdf}~disableShadowEpochs & \hfuzz=500pt boolean & \hfuzz=500pt disable epochs if satellite is in Earth's/Moon's shadow\\
\hfuzz=500pt\includegraphics[width=1em]{element-mustset.pdf}~ephemerides & \hfuzz=500pt \hyperref[ephemeridesType]{ephemerides} & \hfuzz=500pt \\
\hfuzz=500pt\includegraphics[width=1em]{element-mustset.pdf}~eclipse & \hfuzz=500pt \hyperref[eclipseType]{eclipse} & \hfuzz=500pt eclipse model used to determine if a satellite is in Earth's shadow\\
\hfuzz=500pt\includegraphics[width=1em]{element-mustset.pdf}~noAntennaPatternFound & \hfuzz=500pt choice & \hfuzz=500pt what should happen is no antenna pattern is found for an observation\\
\hfuzz=500pt\includegraphics[width=1em]{connector.pdf}\includegraphics[width=1em]{element-mustset.pdf}~ignoreObservation & \hfuzz=500pt  & \hfuzz=500pt ignore observation if no matching pattern is found\\
\hfuzz=500pt\includegraphics[width=1em]{connector.pdf}\includegraphics[width=1em]{element-mustset.pdf}~useNearestFrequency & \hfuzz=500pt  & \hfuzz=500pt use pattern of nearest frequency if no matching pattern is found\\
\hfuzz=500pt\includegraphics[width=1em]{connector.pdf}\includegraphics[width=1em]{element-mustset.pdf}~throwException & \hfuzz=500pt  & \hfuzz=500pt throw exception if no matching pattern is found\\
\hline
\end{tabularx}

\clearpage
%==================================

\section{GnssProcessingStep}\label{gnssProcessingStepType}
Processing step in \program{GnssProcessing}.

Processing steps enable a dynamic definition of the consecutive steps performed during any kind of GNSS processing.
The most common steps are \configClass{estimate}{gnssProcessingStepType:estimate}, which performs an iterative least
squares adjustment, and \configClass{writeResults}{gnssProcessingStepType:writeResults}, which writes all output files
defined in \program{GnssProcessing} and is usually the last step.
Some steps such as \configClass{selectParameters}{gnssProcessingStepType:selectParameters},
\configClass{selectEpochs}{gnssProcessingStepType:selectEpochs},
\configClass{selectNormalsBlockStructure}{gnssProcessingStepType:selectNormalsBlockStructure}, and
\configClass{selectReceivers}{gnssProcessingStepType:selectReceivers} affect all subsequent steps.
In case these steps are used within a \configClass{group}{gnssProcessingStepType:group} or
\configClass{forEachReceiverSeparately}{gnssProcessingStepType:forEachReceiverSeparately} step,
they only affect the steps within this level.

For usage examples see cookbooks on \reference{GNSS satellite orbit determination and network analysis}{cookbook.gnssNetwork:processing}
or \reference{Kinematic orbit determination of LEO satellites}{cookbook.kinematicOrbit}.


\subsection{Estimate}\label{gnssProcessingStepType:estimate}
Iterative least squares adjustment.

Iterates until either \config{convergenceThreshold} or \config{maxIterationCount} are reached.
If \config{computeResiduals}=\verb|yes| and \config{computeWeights} or \config{adjustSigma0} are set,
observation weights and/or $\sigma_0$ are adjusted after each iteration.


\keepXColumns
\begin{tabularx}{\textwidth}{N T A}
\hline
Name & Type & Annotation\\
\hline
\hfuzz=500pt\includegraphics[width=1em]{element.pdf}~computeResiduals & \hfuzz=500pt boolean & \hfuzz=500pt \\
\hfuzz=500pt\includegraphics[width=1em]{element.pdf}~adjustSigma0 & \hfuzz=500pt choice & \hfuzz=500pt adjust sigma0 by scale factor (per receiver)\\
\hfuzz=500pt\includegraphics[width=1em]{connector.pdf}\includegraphics[width=1em]{element-mustset.pdf}~individual & \hfuzz=500pt  & \hfuzz=500pt individual observations (e.g. L1P, L2P, C1P, C2P)\\
\hfuzz=500pt\includegraphics[width=1em]{connector.pdf}\includegraphics[width=1em]{element-mustset.pdf}~grouped & \hfuzz=500pt  & \hfuzz=500pt grouped by phase and code observation (e.g. L**, C**)\\
\hfuzz=500pt\includegraphics[width=1em]{connector.pdf}\includegraphics[width=1em]{element-mustset.pdf}~groupedPhases & \hfuzz=500pt  & \hfuzz=500pt grouped phases and individual code (e.g. L**, C1P, C2P)\\
\hfuzz=500pt\includegraphics[width=1em]{element.pdf}~computeWeights & \hfuzz=500pt choice & \hfuzz=500pt downweight outliers\\
\hfuzz=500pt\includegraphics[width=1em]{connector.pdf}\includegraphics[width=1em]{element-mustset.pdf}~individual & \hfuzz=500pt  & \hfuzz=500pt individual observations (e.g. L1P, L2P, C1P, C2P)\\
\hfuzz=500pt\includegraphics[width=1em]{connector.pdf}\includegraphics[width=1em]{element-mustset.pdf}~grouped & \hfuzz=500pt  & \hfuzz=500pt grouped by phase and code observation (e.g. L**, C**)\\
\hfuzz=500pt\includegraphics[width=1em]{connector.pdf}\includegraphics[width=1em]{element-mustset.pdf}~groupedPhases & \hfuzz=500pt  & \hfuzz=500pt grouped phases and individual code (e.g. L**, C1P, C2P)\\
\hfuzz=500pt\includegraphics[width=1em]{element.pdf}~convergenceThreshold & \hfuzz=500pt double & \hfuzz=500pt [m] stop iteration once full convergence is reached\\
\hfuzz=500pt\includegraphics[width=1em]{element.pdf}~maxIterationCount & \hfuzz=500pt uint & \hfuzz=500pt maximum number of iterations\\
\hline
\end{tabularx}


\subsection{ResolveAmbiguities}\label{gnssProcessingStepType:resolveAmbiguities}
Least squares adjustment with ambiguity resolution.

After this step all resolved ambiguities are removed from the normal equation system, except if \config{dryRun}=\verb|yes|.
If \config{computeResiduals}=\verb|yes| and \config{computeWeights} or \config{adjustSigma0} are set,
observation weights and/or $\sigma_0$ are adjusted after resolving the ambiguities.

For a description of the ambiguity resolution algorithm see \configClass{GnssParametrizationAmbiguities}{gnssParametrizationAmbiguitiesType}.


\keepXColumns
\begin{tabularx}{\textwidth}{N T A}
\hline
Name & Type & Annotation\\
\hline
\hfuzz=500pt\includegraphics[width=1em]{element.pdf}~computeResiduals & \hfuzz=500pt boolean & \hfuzz=500pt \\
\hfuzz=500pt\includegraphics[width=1em]{element.pdf}~adjustSigma0 & \hfuzz=500pt choice & \hfuzz=500pt adjust sigma0 by scale factor (per receiver)\\
\hfuzz=500pt\includegraphics[width=1em]{connector.pdf}\includegraphics[width=1em]{element-mustset.pdf}~individual & \hfuzz=500pt  & \hfuzz=500pt individual observations (e.g. L1P, L2P, C1P, C2P)\\
\hfuzz=500pt\includegraphics[width=1em]{connector.pdf}\includegraphics[width=1em]{element-mustset.pdf}~grouped & \hfuzz=500pt  & \hfuzz=500pt grouped by phase and code observation (e.g. L**, C**)\\
\hfuzz=500pt\includegraphics[width=1em]{connector.pdf}\includegraphics[width=1em]{element-mustset.pdf}~groupedPhases & \hfuzz=500pt  & \hfuzz=500pt grouped phases and individual code (e.g. L**, C1P, C2P)\\
\hfuzz=500pt\includegraphics[width=1em]{element.pdf}~computeWeights & \hfuzz=500pt choice & \hfuzz=500pt downweight outliers\\
\hfuzz=500pt\includegraphics[width=1em]{connector.pdf}\includegraphics[width=1em]{element-mustset.pdf}~individual & \hfuzz=500pt  & \hfuzz=500pt individual observations (e.g. L1P, L2P, C1P, C2P)\\
\hfuzz=500pt\includegraphics[width=1em]{connector.pdf}\includegraphics[width=1em]{element-mustset.pdf}~grouped & \hfuzz=500pt  & \hfuzz=500pt grouped by phase and code observation (e.g. L**, C**)\\
\hfuzz=500pt\includegraphics[width=1em]{connector.pdf}\includegraphics[width=1em]{element-mustset.pdf}~groupedPhases & \hfuzz=500pt  & \hfuzz=500pt grouped phases and individual code (e.g. L**, C1P, C2P)\\
\hfuzz=500pt\includegraphics[width=1em]{element.pdf}~dryRun & \hfuzz=500pt boolean & \hfuzz=500pt compute integer ambiguities but don't mark as solved\\
\hline
\end{tabularx}


\subsection{ComputeCovarianceMatrix}\label{gnssProcessingStepType:computeCovarianceMatrix}
Accumulate normal equations and compute the sparse inverse.


\subsection{WriteNormalEquations}\label{gnssProcessingStepType:writeNormalEquations}
Write unconstrained normal equations and constraint equations.

Epoch parameters (e.g. clock errors) are eliminated before writing the normal equations
if \config{eliminateEpochParameters}=\verb|yes|. If \configClass{remainingParameters}{parameterSelectorType}
is set only the selected parameters are written to the normal equations and all other parameters are eliminated beforehand.


\keepXColumns
\begin{tabularx}{\textwidth}{N T A}
\hline
Name & Type & Annotation\\
\hline
\hfuzz=500pt\includegraphics[width=1em]{element.pdf}~outputfileNormalEquations & \hfuzz=500pt filename & \hfuzz=500pt unconstrained normals\\
\hfuzz=500pt\includegraphics[width=1em]{element.pdf}~outputfileConstraintEquations & \hfuzz=500pt filename & \hfuzz=500pt applied constraints\\
\hfuzz=500pt\includegraphics[width=1em]{element.pdf}~outputfileAprioriSolution & \hfuzz=500pt filename & \hfuzz=500pt a priori parameters\\
\hfuzz=500pt\includegraphics[width=1em]{element.pdf}~outputfileParameterName & \hfuzz=500pt filename & \hfuzz=500pt parameter names\\
\hfuzz=500pt\includegraphics[width=1em]{element.pdf}~variableReceiver & \hfuzz=500pt string & \hfuzz=500pt variable in single receiver loop\\
\hfuzz=500pt\includegraphics[width=1em]{element.pdf}~eliminateEpochParameters & \hfuzz=500pt boolean & \hfuzz=500pt \\
\hfuzz=500pt\includegraphics[width=1em]{element-unbounded.pdf}~remainingParameters & \hfuzz=500pt \hyperref[parameterSelectorType]{parameterSelector} & \hfuzz=500pt parameter order/selection of output normal equations\\
\hfuzz=500pt\includegraphics[width=1em]{element.pdf}~defaultNormalsBlockSize & \hfuzz=500pt uint & \hfuzz=500pt block size for distributing the normal equations, 0: one block, empty: original block size\\
\hline
\end{tabularx}


\subsection{WriteResults}\label{gnssProcessingStepType:writeResults}
Write estimated solution to output files.

All output files defined in \program{GnssProcessing} are written in this step.
It is usually the last processing step, but can also be used at other points in the
processing in combination with \config{suffix} to write intermediate results, for example
before \configClass{resolveAmbiguities}{gnssProcessingStepType:resolveAmbiguities} to
output the float solution.


\keepXColumns
\begin{tabularx}{\textwidth}{N T A}
\hline
Name & Type & Annotation\\
\hline
\hfuzz=500pt\includegraphics[width=1em]{element.pdf}~suffix & \hfuzz=500pt string & \hfuzz=500pt appended to every output file name (e.g. orbit.G01.\$<\$suffix\$>\$.dat)\\
\hline
\end{tabularx}


\subsection{PrintResidualStatistics}\label{gnssProcessingStepType:printResidualStatistics}
Print residual statistics.


\subsection{SelectParameters}\label{gnssProcessingStepType:selectParameters}
Select parameters for subsequent steps.

This step can be used to control which parameters are estimated and which constraints are applied in subsequent
\configClass{estimate}{gnssProcessingStepType:estimate} steps. An example would be to estimate only the TEC bias
in an intermediate step as seen in \reference{GNSS satellite orbit determination and network analysis}{cookbook.gnssNetwork:processing}.
Another example is to process at a 5-minute sampling using \configClass{selectEpochs}{gnssProcessingStepType:selectEpochs} and then at
the end to densify the clock parameters to the full 30-second observation sampling while keeping all other parameters fixed.


\keepXColumns
\begin{tabularx}{\textwidth}{N T A}
\hline
Name & Type & Annotation\\
\hline
\hfuzz=500pt\includegraphics[width=1em]{element.pdf}~estimateReceiverClock & \hfuzz=500pt boolean & \hfuzz=500pt \\
\hfuzz=500pt\includegraphics[width=1em]{element.pdf}~estimateReceiverKinematicPosition & \hfuzz=500pt boolean & \hfuzz=500pt \\
\hfuzz=500pt\includegraphics[width=1em]{element.pdf}~estimateReceiverPosition & \hfuzz=500pt boolean & \hfuzz=500pt \\
\hfuzz=500pt\includegraphics[width=1em]{element.pdf}~constraintReceiverNoNetTranslation & \hfuzz=500pt boolean & \hfuzz=500pt \\
\hfuzz=500pt\includegraphics[width=1em]{element.pdf}~constraintReceiverNoNetRotation & \hfuzz=500pt boolean & \hfuzz=500pt \\
\hfuzz=500pt\includegraphics[width=1em]{element.pdf}~estimateReceiverTroposphereWet & \hfuzz=500pt boolean & \hfuzz=500pt \\
\hfuzz=500pt\includegraphics[width=1em]{element.pdf}~estimateReceiverTroposphereGradient & \hfuzz=500pt boolean & \hfuzz=500pt \\
\hfuzz=500pt\includegraphics[width=1em]{element.pdf}~estimateReceiverSignalBias & \hfuzz=500pt boolean & \hfuzz=500pt \\
\hfuzz=500pt\includegraphics[width=1em]{element.pdf}~constraintReceiverSignalBiasDiffClock & \hfuzz=500pt boolean & \hfuzz=500pt separate clock errors and biases of different satellite constellations\\
\hfuzz=500pt\includegraphics[width=1em]{element.pdf}~estimateReceiverTecBias & \hfuzz=500pt boolean & \hfuzz=500pt \\
\hfuzz=500pt\includegraphics[width=1em]{element.pdf}~estimateReceiverAntennaCenterVariations & \hfuzz=500pt boolean & \hfuzz=500pt \\
\hfuzz=500pt\includegraphics[width=1em]{element.pdf}~estimateTransmitterClock & \hfuzz=500pt boolean & \hfuzz=500pt \\
\hfuzz=500pt\includegraphics[width=1em]{element.pdf}~constraintTransmitterClockZeroMean & \hfuzz=500pt boolean & \hfuzz=500pt defines the time system\\
\hfuzz=500pt\includegraphics[width=1em]{element.pdf}~constraintTransmitterClockModel & \hfuzz=500pt boolean & \hfuzz=500pt \\
\hfuzz=500pt\includegraphics[width=1em]{element.pdf}~estimateTransmitterKinematicPosition & \hfuzz=500pt boolean & \hfuzz=500pt \\
\hfuzz=500pt\includegraphics[width=1em]{element.pdf}~estimateTransmitterOrbit & \hfuzz=500pt boolean & \hfuzz=500pt incl. SRP parameters and pulses\\
\hfuzz=500pt\includegraphics[width=1em]{element.pdf}~estimateTransmitterClockModel & \hfuzz=500pt boolean & \hfuzz=500pt \\
\hfuzz=500pt\includegraphics[width=1em]{element.pdf}~estimateTransmitterSignalBias & \hfuzz=500pt boolean & \hfuzz=500pt \\
\hfuzz=500pt\includegraphics[width=1em]{element.pdf}~constraintTransmitterSignalBiasZeroMean & \hfuzz=500pt boolean & \hfuzz=500pt separate transmitter and receiver biases\\
\hfuzz=500pt\includegraphics[width=1em]{element.pdf}~estimateTransmitterTecBias & \hfuzz=500pt boolean & \hfuzz=500pt \\
\hfuzz=500pt\includegraphics[width=1em]{element.pdf}~estimateTransmitterBiasModel & \hfuzz=500pt boolean & \hfuzz=500pt \\
\hfuzz=500pt\includegraphics[width=1em]{element.pdf}~constraintTransmitterSignalBiasModelZeroMean & \hfuzz=500pt boolean & \hfuzz=500pt separate temporal biases from constant biases\\
\hfuzz=500pt\includegraphics[width=1em]{element.pdf}~estimateTransmitterAntennaCenterVariations & \hfuzz=500pt boolean & \hfuzz=500pt \\
\hfuzz=500pt\includegraphics[width=1em]{element.pdf}~estimateEarthRotationPole & \hfuzz=500pt boolean & \hfuzz=500pt \\
\hfuzz=500pt\includegraphics[width=1em]{element.pdf}~estimateEarthRotationUT1 & \hfuzz=500pt boolean & \hfuzz=500pt \\
\hfuzz=500pt\includegraphics[width=1em]{element.pdf}~estimateEarthRotationNutation & \hfuzz=500pt boolean & \hfuzz=500pt \\
\hfuzz=500pt\includegraphics[width=1em]{element.pdf}~estimateGravity & \hfuzz=500pt boolean & \hfuzz=500pt \\
\hfuzz=500pt\includegraphics[width=1em]{element.pdf}~estimateAmbiguities & \hfuzz=500pt boolean & \hfuzz=500pt \\
\hfuzz=500pt\includegraphics[width=1em]{element.pdf}~estimateIonosphereSTEC & \hfuzz=500pt boolean & \hfuzz=500pt \\
\hfuzz=500pt\includegraphics[width=1em]{element.pdf}~constraintIonosphereSTEC & \hfuzz=500pt boolean & \hfuzz=500pt \\
\hfuzz=500pt\includegraphics[width=1em]{element.pdf}~estimateIonosphereVTec & \hfuzz=500pt boolean & \hfuzz=500pt \\
\hfuzz=500pt\includegraphics[width=1em]{element.pdf}~estimateIonosphereTecMap & \hfuzz=500pt boolean & \hfuzz=500pt \\
\hfuzz=500pt\includegraphics[width=1em]{element.pdf}~constraintOther & \hfuzz=500pt boolean & \hfuzz=500pt \\
\hline
\end{tabularx}


\subsection{SelectEpochs}\label{gnssProcessingStepType:selectEpochs}
Select epochs for subsequent steps.

This step can be used to reduce the processing sampling while keeping the original observation sampling
for all preprocessing steps (e.g. outlier and cycle slip detection).
Another example is to process at a 5-minute sampling by setting \config{nthEpoch}=\verb|10| and then
at the end to densify only the clock parameters (see \configClass{selectParameters}{gnssProcessingStepType:selectParameters})
to the full 30-second observation sampling by setting \config{nthEpoch}=\verb|1| while keeping all other parameters fixed.


\keepXColumns
\begin{tabularx}{\textwidth}{N T A}
\hline
Name & Type & Annotation\\
\hline
\hfuzz=500pt\includegraphics[width=1em]{element-mustset.pdf}~nthEpoch & \hfuzz=500pt uint & \hfuzz=500pt use only every nth epoch in all subsequent processing steps\\
\hline
\end{tabularx}


\subsection{SelectNormalsBlockStructure}\label{gnssProcessingStepType:selectNormalsBlockStructure}
Select block structure of sparse normal equations for subsequent steps.

This step can be used to define the structure of the different parts of the normal equation system,
which can have a major impact on computing performance and memory consumption depending on the processing setup.

\fig{!hb}{0.4}{gnss_normals_structure}{fig:gnss_normals_structure}{Structure of normal equations in GNSS processing}

The normal equation system is divided into three parts for epoch, interval, and ambiguity parameters.
The epoch part is subdivided further into one subpart per epoch. Each part is divided into blocks and only non-zero
blocks are stored in memory to reduce memory consumption and to prevent unnecessary matrix computations.
\config{defaultBlockSizeEpoch}, \config{defaultBlockSizeInterval}, and \config{defaultBlockSizeAmbiguity} control
the size of the blocks within each part of the normal equations. \config{defaultBlockReceiverCount} can be set to group
a number of receivers into one block within the epoch and interval parts.

If \config{keepEpochNormalsInMemory}=\verb|no| epoch blocks are eliminated after they are set up to reduce the number
of parameters in the normal equation system. \config{defaultBlockCountReduction} controls after how many epoch blocks
an elimination step is performed. For larger processing setups or high sampling rates epoch block elimination is recommended
as the large number of clock parameters require a lot of memory.


\keepXColumns
\begin{tabularx}{\textwidth}{N T A}
\hline
Name & Type & Annotation\\
\hline
\hfuzz=500pt\includegraphics[width=1em]{element.pdf}~defaultBlockSizeEpoch & \hfuzz=500pt uint & \hfuzz=500pt block size of epoch parameters, 0: one block\\
\hfuzz=500pt\includegraphics[width=1em]{element.pdf}~defaultBlockSizeInterval & \hfuzz=500pt uint & \hfuzz=500pt block size of interval parameters, 0: one block\\
\hfuzz=500pt\includegraphics[width=1em]{element.pdf}~defaultBlockSizeAmbiguity & \hfuzz=500pt uint & \hfuzz=500pt block size of ambiguity parameters, 0: one block\\
\hfuzz=500pt\includegraphics[width=1em]{element.pdf}~defaultBlockReceiverCount & \hfuzz=500pt uint & \hfuzz=500pt number of receivers to group into one block for epoch and interval\\
\hfuzz=500pt\includegraphics[width=1em]{element.pdf}~defaultBlockCountReduction & \hfuzz=500pt uint & \hfuzz=500pt minimum number of blocks for epoch reduction\\
\hfuzz=500pt\includegraphics[width=1em]{element.pdf}~keepEpochNormalsInMemory & \hfuzz=500pt boolean & \hfuzz=500pt speeds up processing but uses much more memory\\
\hfuzz=500pt\includegraphics[width=1em]{element.pdf}~accumulateEpochObservations & \hfuzz=500pt boolean & \hfuzz=500pt set up all observations per epoch and receiver at once\\
\hline
\end{tabularx}


\subsection{SelectReceivers}\label{gnssProcessingStepType:selectReceivers}
Select receivers of a \configClass{stationNetwork}{gnssParametrizationReceiverType:stationNetwork} to process in subsequent steps.

This step can be used to process only a subset of stations in subsequent processing steps.
The most common use is to start the processing with a well-distributed network of core stations as seen in
\reference{GNSS satellite orbit determination and network analysis}{cookbook.gnssNetwork:processing}.
In this case \configFile{inputfileStationList}{stringList} should contain one line per station with alternative
stations on the same line separated by whitespace. To later process all other stations individually, use
the processing step \configClass{forEachReceiverSeparately}{gnssProcessingStepType:forEachReceiverSeparately}
and provide the same station list as \configFile{inputfileExcludeStationList}{stringList} in that step.


\keepXColumns
\begin{tabularx}{\textwidth}{N T A}
\hline
Name & Type & Annotation\\
\hline
\hfuzz=500pt\includegraphics[width=1em]{element.pdf}~inputfileStationList & \hfuzz=500pt filename & \hfuzz=500pt empty = all stations\\
\hfuzz=500pt\includegraphics[width=1em]{element.pdf}~inputfileExcludeStationList & \hfuzz=500pt filename & \hfuzz=500pt exclude these stations\\
\hline
\end{tabularx}


\subsection{ForEachReceiverSeparately}\label{gnssProcessingStepType:forEachReceiverSeparately}
Perform these processing steps for each receiver of a \configClass{stationNetwork}{gnssParametrizationReceiverType:stationNetwork} separately.
All \configClass{transmitter}{gnssParametrizationTransmitterType} parameters are disabled in these processing steps.

This step can be used for individual precise point positioning (PPP) of all stations.
During \reference{GNSS satellite orbit determination and network analysis}{cookbook.gnssNetwork:processing} this step is used after the
initial processing of the core network to process all other stations individually. In that case provide the same station list as
\configFile{inputfileExcludeStationList}{stringList} in this step that was used as \configFile{inputfileStationList}{stringList} in the
\configClass{selectReceivers}{gnssProcessingStepType:selectReceivers} step where the core network was selected.


\keepXColumns
\begin{tabularx}{\textwidth}{N T A}
\hline
Name & Type & Annotation\\
\hline
\hfuzz=500pt\includegraphics[width=1em]{element.pdf}~inputfileStationList & \hfuzz=500pt filename & \hfuzz=500pt empty = all stations\\
\hfuzz=500pt\includegraphics[width=1em]{element.pdf}~inputfileExcludeStationList & \hfuzz=500pt filename & \hfuzz=500pt exclude these stations\\
\hfuzz=500pt\includegraphics[width=1em]{element-mustset-unbounded.pdf}~processingStep & \hfuzz=500pt \hyperref[gnssProcessingStepType]{gnssProcessingStep} & \hfuzz=500pt steps are processed consecutively\\
\hline
\end{tabularx}


\subsection{Group}\label{gnssProcessingStepType:group}
Perform these processing steps. This step can be used to structure complex processing flows and has no further effect on the processing itself.


\keepXColumns
\begin{tabularx}{\textwidth}{N T A}
\hline
Name & Type & Annotation\\
\hline
\hfuzz=500pt\includegraphics[width=1em]{element-mustset-unbounded.pdf}~processingStep & \hfuzz=500pt \hyperref[gnssProcessingStepType]{gnssProcessingStep} & \hfuzz=500pt steps are processed consecutively\\
\hline
\end{tabularx}

\clearpage
%==================================

\section{GnssType}\label{gnssType}
A GnssType string consists of six parts (type, frequency, attribute, system, PRN, frequency number)
represented by seven characters.
\begin{itemize}
\item The first three characters (representing type, frequency, and attribute) correspond to the observation codes of the
      \href{https://files.igs.org/pub/data/format/rinex305.pdf}{RINEX 3 definition}.
\item The satellite system character also follows the RINEX 3 definition:
      \begin{itemize}
        \item \verb|G| = GPS
        \item \verb|R| = GLONASS
        \item \verb|E| = Galileo
        \item \verb|C| = BeiDou
        \item \verb|S| = SBAS
        \item \verb|J| = QZSS
        \item \verb|I| = IRNSS
      \end{itemize}
\item PRN is a two-digit number identifying a satellite.
\item Frequency number is only used for GLONASS, where the range -7 to 14 is represented by letters starting with A.
\end{itemize}

Each part of a GnssType string can be replaced by a wildcard '\verb|*|', enabling the use of these strings as patterns,
for example to select a subset of observations (e.g. \verb|C**G**| matches all GPS code/range observations).
Trailing wildcards are optional, meaning \verb|L1*R| is automatically expanded to \verb|L1*R***|.
For some RINEX 2 types (e.g. Galileo L5) the RINEX 3 attribute is unknown/undefined and can be replaced by \verb|?|,
for example \verb|L5?E01|.

Examples:
\begin{itemize}
\item \verb|C1CG23| = code/range observation, L1 frequency, derived from C/A code, GPS, PRN 23
\item \verb|L2PR05B| = phase observation, G2 frequency, derived from P code, GLONASS, PRN 05, frequency number -6
\item \verb|*5*E**| = all observation types, E5a frequency, all attributes, Galileo, all PRNs
\end{itemize}

\clearpage
%==================================

\section{Gravityfield}\label{gravityfieldType}
This class computes functionals of the time depending gravity field,
e.g potential, gravity anomalies or gravity gradients.

If several instances of the class are given the results are summed up.
Before summation every single result is multiplicated by a \config{factor}.
To substract a normal field like GRS80 from a potential
to get the disturbance potential you must choose one factor by 1
and the other by -1. To get the mean of two fields just set each factor to 0.5.

Some of the instances gives also information about the accuracy.
The variance of the result (sum) is computed by means of variance propagation.


\subsection{PotentialCoefficients}\label{gravityfieldType:potentialCoefficients}
Reads coefficients of a spherical harmonics expansion from file.
The potential is given by
\begin{equation}
V(\lambda,\vartheta,r) = \frac{GM}{R}\sum_{n=0}^\infty \sum_{m=0}^n \left(\frac{R}{r}\right)^{n+1}
  \left(c_{nm} C_{nm}(\lambda,\vartheta) + s_{nm} S_{nm}(\lambda,\vartheta)\right).
\end{equation}
If set the expansion is limited in the range between \config{minDegree}
and \config{maxDegree} inclusivly. The computed result
is multiplied with \config{factor}. If \config{setSigmasToZero} is true
the variances are set to zero. This option is only important for variance propagation
and does not change the result of the gravity field functionals.


\keepXColumns
\begin{tabularx}{\textwidth}{N T A}
\hline
Name & Type & Annotation\\
\hline
\hfuzz=500pt\includegraphics[width=1em]{element-mustset.pdf}~inputfilePotentialCoefficients & \hfuzz=500pt filename & \hfuzz=500pt \\
\hfuzz=500pt\includegraphics[width=1em]{element.pdf}~minDegree & \hfuzz=500pt uint & \hfuzz=500pt \\
\hfuzz=500pt\includegraphics[width=1em]{element.pdf}~maxDegree & \hfuzz=500pt uint & \hfuzz=500pt \\
\hfuzz=500pt\includegraphics[width=1em]{element.pdf}~factor & \hfuzz=500pt double & \hfuzz=500pt the result is multplied by this factor, set -1 to substract the field\\
\hfuzz=500pt\includegraphics[width=1em]{element.pdf}~setSigmasToZero & \hfuzz=500pt boolean & \hfuzz=500pt set variances to zero, should be used by adding back reference fields\\
\hline
\end{tabularx}


\subsection{PotentialCoefficientsInterior}
Reads coefficients of a spherical harmonics expansion (for inner space) from file.
If set the expansion is limited in the range between \config{minDegree}
and \config{maxDegree} inclusivly. The computed result is multiplied with \config{factor}.
If \config{setSigmasToZero} is true the variances are set to zero.
This option is only important for error propagation
and does not change the result of the gravity field functionals.


\keepXColumns
\begin{tabularx}{\textwidth}{N T A}
\hline
Name & Type & Annotation\\
\hline
\hfuzz=500pt\includegraphics[width=1em]{element-mustset.pdf}~inputfilePotentialCoefficients & \hfuzz=500pt filename & \hfuzz=500pt \\
\hfuzz=500pt\includegraphics[width=1em]{element.pdf}~minDegree & \hfuzz=500pt uint & \hfuzz=500pt \\
\hfuzz=500pt\includegraphics[width=1em]{element.pdf}~maxDegree & \hfuzz=500pt uint & \hfuzz=500pt \\
\hfuzz=500pt\includegraphics[width=1em]{element.pdf}~factor & \hfuzz=500pt double & \hfuzz=500pt the result is multplied by this factor, set -1 to substract the field\\
\hfuzz=500pt\includegraphics[width=1em]{element.pdf}~setSigmasToZero & \hfuzz=500pt boolean & \hfuzz=500pt set variances to zero, should be used by adding back reference fields\\
\hline
\end{tabularx}


\subsection{FromParametrization}\label{gravityfieldType:fromParametrization}
Reads a solution vector from file \configFile{inputfileSolution}{matrix}
which may be computed by a least squares adjustment (e.g. by \program{NormalsSolverVCE}).
The coefficients of the vector are interpreted from position \config{indexStart}
(counting from zero) with help of \configClass{parametrizationGravity}{parametrizationGravityType}.
If the solution file contains solution of several right hand sides you can choose
one with number \config{rightSide} (counting from zero).
You can also read a vector from file \configFile{inputfileSigmax}{matrix}
containing the accuracies of the coefficients.

The computed result is multiplied with \config{factor}.


\keepXColumns
\begin{tabularx}{\textwidth}{N T A}
\hline
Name & Type & Annotation\\
\hline
\hfuzz=500pt\includegraphics[width=1em]{element-mustset-unbounded.pdf}~parametrization & \hfuzz=500pt \hyperref[parametrizationGravityType]{parametrizationGravity} & \hfuzz=500pt \\
\hfuzz=500pt\includegraphics[width=1em]{element-mustset.pdf}~inputfileSolution & \hfuzz=500pt filename & \hfuzz=500pt solution vector\\
\hfuzz=500pt\includegraphics[width=1em]{element.pdf}~inputfileSigmax & \hfuzz=500pt filename & \hfuzz=500pt standards deviations or covariance matrix of the solution\\
\hfuzz=500pt\includegraphics[width=1em]{element.pdf}~indexStart & \hfuzz=500pt uint & \hfuzz=500pt position in the solution vector\\
\hfuzz=500pt\includegraphics[width=1em]{element.pdf}~rightSide & \hfuzz=500pt uint & \hfuzz=500pt if solution contains several right hand sides, select one\\
\hfuzz=500pt\includegraphics[width=1em]{element.pdf}~factor & \hfuzz=500pt double & \hfuzz=500pt the result is multplied by this factor, set -1 to substract the field\\
\hline
\end{tabularx}


\subsection{TimeSplines}\label{gravityfieldType:timeSplines}
Read a time variable gravity field from file
\configFile{inputfileTimeSplinesGravityfield}{timeSplinesGravityField}
represented by a spherical harmonics expansion in the spatial domain and spline functions
in the time domain. If set the expansion is limited in the range between
\config{minDegree} and \config{maxDegree} inclusivly.

This file can be created for example by \program{Gravityfield2TimeSplines} or
\program{PotentialCoefficients2BlockMeanTimeSplines}.

The computed result is multiplied with \config{factor}.


\keepXColumns
\begin{tabularx}{\textwidth}{N T A}
\hline
Name & Type & Annotation\\
\hline
\hfuzz=500pt\includegraphics[width=1em]{element-mustset.pdf}~inputfileTimeSplinesGravityfield & \hfuzz=500pt filename & \hfuzz=500pt \\
\hfuzz=500pt\includegraphics[width=1em]{element.pdf}~inputfileTimeSplinesCovariance & \hfuzz=500pt filename & \hfuzz=500pt \\
\hfuzz=500pt\includegraphics[width=1em]{element.pdf}~minDegree & \hfuzz=500pt uint & \hfuzz=500pt \\
\hfuzz=500pt\includegraphics[width=1em]{element.pdf}~maxDegree & \hfuzz=500pt uint & \hfuzz=500pt \\
\hfuzz=500pt\includegraphics[width=1em]{element.pdf}~factor & \hfuzz=500pt double & \hfuzz=500pt the result is multplied by this factor, set -1 to substract the field\\
\hline
\end{tabularx}


\subsection{Trend}\label{gravityfieldType:trend}
The given \configClass{gravityfield}{gravityfieldType} is interpreted
as trend function and the result is computed at time $t$ as follows
\begin{equation}
V(\M x,t) = \frac{t-t_0}{\Delta t}V(\M x),
\end{equation}
with $t_0$ is \config{timeStart} and $\Delta t$ is \config{timeStep}.


\keepXColumns
\begin{tabularx}{\textwidth}{N T A}
\hline
Name & Type & Annotation\\
\hline
\hfuzz=500pt\includegraphics[width=1em]{element-mustset-unbounded.pdf}~gravityfield & \hfuzz=500pt \hyperref[gravityfieldType]{gravityfield} & \hfuzz=500pt this field is multiplicated by (time-time0)/timeStep\\
\hfuzz=500pt\includegraphics[width=1em]{element-mustset.pdf}~timeStart & \hfuzz=500pt time & \hfuzz=500pt reference time\\
\hfuzz=500pt\includegraphics[width=1em]{element-mustset.pdf}~timeStep & \hfuzz=500pt time & \hfuzz=500pt \\
\hline
\end{tabularx}


\subsection{Oscillation}\label{gravityfieldType:oscillation}
The given \configClass{gravityfield}{gravityfieldType} is interpreted
as oscillation function and the result is computed at time $t$ as follows
\begin{equation}
V(\M x,t) = \cos(\omega)V_{cos}(\M x)+\sin(\omega)V_{sin}(\M x),
\end{equation}
with $\omega=\frac{2\pi}{T}(t-t_0)$.


\keepXColumns
\begin{tabularx}{\textwidth}{N T A}
\hline
Name & Type & Annotation\\
\hline
\hfuzz=500pt\includegraphics[width=1em]{element-mustset-unbounded.pdf}~gravityfieldCos & \hfuzz=500pt \hyperref[gravityfieldType]{gravityfield} & \hfuzz=500pt multiplicated by cos(2pi/T(time-time0))\\
\hfuzz=500pt\includegraphics[width=1em]{element-mustset-unbounded.pdf}~gravityfieldSin & \hfuzz=500pt \hyperref[gravityfieldType]{gravityfield} & \hfuzz=500pt multiplicated by sin(2pi/T(time-time0))\\
\hfuzz=500pt\includegraphics[width=1em]{element-mustset.pdf}~time0 & \hfuzz=500pt time & \hfuzz=500pt reference time\\
\hfuzz=500pt\includegraphics[width=1em]{element-mustset.pdf}~period & \hfuzz=500pt time & \hfuzz=500pt [day]\\
\hline
\end{tabularx}


\subsection{InInterval}
A \configClass{gravityfield}{gravityfieldType} is only evaluated in the interval between
\config{timeStart} inclusively and \config{timeEnd} exclusively.
Outside the interval the result is zero.

This class is useful to get a time series of monthly mean GRACE gravity field solutions.
In each month another file of potentialCoefficients is valid.
This can easily be created with \configClass{loop}{loopType}.


\keepXColumns
\begin{tabularx}{\textwidth}{N T A}
\hline
Name & Type & Annotation\\
\hline
\hfuzz=500pt\includegraphics[width=1em]{element-mustset-unbounded.pdf}~gravityfield & \hfuzz=500pt \hyperref[gravityfieldType]{gravityfield} & \hfuzz=500pt \\
\hfuzz=500pt\includegraphics[width=1em]{element-mustset.pdf}~timeStart & \hfuzz=500pt time & \hfuzz=500pt first point in time\\
\hfuzz=500pt\includegraphics[width=1em]{element-mustset.pdf}~timeEnd & \hfuzz=500pt time & \hfuzz=500pt last point in time will be less or equal timeEnd\\
\hline
\end{tabularx}


\subsection{Tides}\label{gravityfieldType:tides}
Treat \configClass{tides}{tidesType} as gravitational forces.
The tides need a realization of \configClass{earthRotation}{earthRotationType}
to transform between the CRF and TRF and to compute rotational deformation
from polar motion.
It also needs \configClass{ephemerides}{ephemeridesType} from Sun, moon, and planets.


\keepXColumns
\begin{tabularx}{\textwidth}{N T A}
\hline
Name & Type & Annotation\\
\hline
\hfuzz=500pt\includegraphics[width=1em]{element-mustset-unbounded.pdf}~tides & \hfuzz=500pt \hyperref[tidesType]{tides} & \hfuzz=500pt \\
\hfuzz=500pt\includegraphics[width=1em]{element-mustset.pdf}~earthRotation & \hfuzz=500pt \hyperref[earthRotationType]{earthRotation} & \hfuzz=500pt \\
\hfuzz=500pt\includegraphics[width=1em]{element.pdf}~ephemerides & \hfuzz=500pt \hyperref[ephemeridesType]{ephemerides} & \hfuzz=500pt \\
\hline
\end{tabularx}


\subsection{Topography}\label{gravityfieldType:topography}
The gravity is integrated from a topographic mass distribution.
For each grid point in \configFile{inputfileGridRectangular}{griddedData} a prisma with
\config{density} is assumed. The horizontal extension is computed from the grid spacing
and the vertical extension is given by \config{radialLowerBound}
and \config{radialUpperBound} above ellipsoid. All values are expressions and computed
for each point with given data in the grid file. The standard variables for grids
are available, see~\reference{dataVariables}{general.parser:dataVariables}.

Example: The grid file contains the orthometric height of the topography in the first
column, the geoid height in the second and the mean density of each prism in the third
column. In this case the following settings should be used:
\begin{itemize}
\item \config{radialUpperBound} = \verb|data0+data1|,
\item \config{radialLowerBound} = \verb|data1|,
\item \config{density} = \verb|data2|.
\end{itemize}

As the prim computation is time consuming a maximum distance around the evaluation point
can defined with \config{distancePrism}. Afterwards a simplified radial line
(the prism mass is concentrated to a line in the center) is used up to
a distance of \config{distanceLine}. At last the prim is approximated by a point mass
in the center up to a distance \config{distanceMax} (if set). Prisms nearby the evaluation
point can be excluded with \config{distanceMin}.


\keepXColumns
\begin{tabularx}{\textwidth}{N T A}
\hline
Name & Type & Annotation\\
\hline
\hfuzz=500pt\includegraphics[width=1em]{element-mustset.pdf}~inputfileGridRectangular & \hfuzz=500pt filename & \hfuzz=500pt Digital Terrain Model\\
\hfuzz=500pt\includegraphics[width=1em]{element.pdf}~density & \hfuzz=500pt expression & \hfuzz=500pt expression [kg/m**3]\\
\hfuzz=500pt\includegraphics[width=1em]{element.pdf}~radialUpperBound & \hfuzz=500pt expression & \hfuzz=500pt expression (variables 'height', 'data', 'L', 'B' and, 'area' are taken from the gridded data\\
\hfuzz=500pt\includegraphics[width=1em]{element.pdf}~radialLowerBound & \hfuzz=500pt expression & \hfuzz=500pt expression (variables 'height', 'data', 'L', 'B' and, 'area' are taken from the gridded data\\
\hfuzz=500pt\includegraphics[width=1em]{element.pdf}~distanceMin & \hfuzz=500pt double & \hfuzz=500pt [km] min. influence distance (ignore near zone)\\
\hfuzz=500pt\includegraphics[width=1em]{element.pdf}~distancePrism & \hfuzz=500pt double & \hfuzz=500pt [km] max. distance for prism formular\\
\hfuzz=500pt\includegraphics[width=1em]{element.pdf}~distanceLine & \hfuzz=500pt double & \hfuzz=500pt [km] max. distance for radial integration\\
\hfuzz=500pt\includegraphics[width=1em]{element.pdf}~distanceMax & \hfuzz=500pt double & \hfuzz=500pt [km] max. influence distance (ignore far zone)\\
\hfuzz=500pt\includegraphics[width=1em]{element.pdf}~factor & \hfuzz=500pt double & \hfuzz=500pt the result is multplied by this factor, set -1 to substract the field\\
\hline
\end{tabularx}


\subsection{EarthquakeOscillation}
The given \configClass{gravityfield}{gravityfieldType} is interpreted as an oscillation function
in the gravitational potential field, caused by large earthquakes.
The result is computed at time $t$ as follows:
\begin{equation}
C_{lm}(\M t) = \sum_{n=0}^NC_{nlm}(1-\cos(\omega)\exp(\frac{-\omega}{2Q_{nlm}})),
\end{equation}
with $\omega=\frac{2\pi}{T_{nlm}}(t-t_0)$. In this equation, $Q_{nlm}$ is the attenuation factor,
$n$ is the overtone factor, $m$ is degree, $l$ is order, and $t$ is time in second.
$T_{nlm}$ and $Q_{nlm}$ are computed with the elastic Earth model or observed from the long
period record of superconducting gravimeter measurements after the earthquakes.


\keepXColumns
\begin{tabularx}{\textwidth}{N T A}
\hline
Name & Type & Annotation\\
\hline
\hfuzz=500pt\includegraphics[width=1em]{element-mustset.pdf}~inputCoefficientMatrix & \hfuzz=500pt filename & \hfuzz=500pt oscillation model parameters\\
\hfuzz=500pt\includegraphics[width=1em]{element-mustset.pdf}~time0 & \hfuzz=500pt time & \hfuzz=500pt the time earthquake happened\\
\hfuzz=500pt\includegraphics[width=1em]{element.pdf}~minDegree & \hfuzz=500pt uint & \hfuzz=500pt \\
\hfuzz=500pt\includegraphics[width=1em]{element.pdf}~maxDegree & \hfuzz=500pt uint & \hfuzz=500pt \\
\hfuzz=500pt\includegraphics[width=1em]{element.pdf}~GM & \hfuzz=500pt double & \hfuzz=500pt Geocentric gravitational constant\\
\hfuzz=500pt\includegraphics[width=1em]{element.pdf}~R & \hfuzz=500pt double & \hfuzz=500pt reference radius\\
\hline
\end{tabularx}


\subsection{Filter}
Convert \configClass{gravityfield}{gravityfieldType} to spherical harmonics
and \configClass{filter}{sphericalHarmonicsFilterType} the coefficients.


\keepXColumns
\begin{tabularx}{\textwidth}{N T A}
\hline
Name & Type & Annotation\\
\hline
\hfuzz=500pt\includegraphics[width=1em]{element-mustset-unbounded.pdf}~gravityfield & \hfuzz=500pt \hyperref[gravityfieldType]{gravityfield} & \hfuzz=500pt \\
\hfuzz=500pt\includegraphics[width=1em]{element-mustset-unbounded.pdf}~filter & \hfuzz=500pt \hyperref[sphericalHarmonicsFilterType]{sphericalHarmonicsFilter} & \hfuzz=500pt \\
\hline
\end{tabularx}

\clearpage
%==================================

\section{Grid}\label{gridType}
This class generates a set of grid points. In a first step, the grid
is always generated globally, with \configClass{border}{borderType} a regional
subset of points can be extracted from the global grid. The parameters
\config{R} and \config{inverseFlattening} define the shape of the ellipsoid
on which the grid is generated. In case \config{inverseFlattening} is
chosen as zero, a sphere is used. With \config{height} the distance of
the points above the ellipsoid can be defined. In addition to the location
of the points, weights are assigned to each of the points. These weights
can be regarded as the surface element associated with each grid point.


\subsection{Geograph}
The geographical grid is an equal-angular point distribution with points
located along meridians and along circles of latitude. \config{deltaLambda}
denotes the angular difference between adjacent points along meridians and
\config{deltaPhi} describes the angular difference between adjacent points
along circles of latitude. The point setting results as follows:
\begin{equation}
\lambda_i=\frac{\Delta\lambda}{2}+i\cdot\Delta\lambda\qquad\mbox{with}\qquad 0\leq i< \frac{360^\circ}{\Delta\lambda},
\end{equation}
\begin{equation}
\varphi_j=-90^\circ+\frac{\Delta\varphi}{2}+j\cdot\Delta\varphi\qquad\mbox{with}\qquad 0\leq j<\frac{180^\circ}{\Delta\varphi}.
\end{equation}
The number of grid points can be determined by
\begin{equation}
I=\frac{360^\circ}{\Delta\lambda}\cdot\frac{180^\circ}{\Delta\varphi}.
\end{equation}
The weights are calculated according to
\begin{equation}
w_i=\int\limits_{\lambda_i-\frac{\Delta\lambda}{2}}^{\lambda_i+\frac{\Delta\lambda}{2}}\int\limits_{\vartheta_i-\frac{\Delta\vartheta}{2}}^{\vartheta_i+\frac{\Delta\vartheta}{2}}=2\cdot\Delta\lambda\sin(\Delta\vartheta)\sin(\vartheta_i).
\end{equation}


\keepXColumns
\begin{tabularx}{\textwidth}{N T A}
\hline
Name & Type & Annotation\\
\hline
\hfuzz=500pt\includegraphics[width=1em]{element-mustset.pdf}~deltaLambda & \hfuzz=500pt angle & \hfuzz=500pt \\
\hfuzz=500pt\includegraphics[width=1em]{element-mustset.pdf}~deltaPhi & \hfuzz=500pt angle & \hfuzz=500pt \\
\hfuzz=500pt\includegraphics[width=1em]{element.pdf}~height & \hfuzz=500pt double & \hfuzz=500pt ellipsoidal height expression (variables 'height', 'L', 'B')\\
\hfuzz=500pt\includegraphics[width=1em]{element.pdf}~R & \hfuzz=500pt double & \hfuzz=500pt major axsis of the ellipsoid/sphere\\
\hfuzz=500pt\includegraphics[width=1em]{element.pdf}~inverseFlattening & \hfuzz=500pt double & \hfuzz=500pt flattening of the ellipsoid, 0: sphere\\
\hfuzz=500pt\includegraphics[width=1em]{element-unbounded.pdf}~border & \hfuzz=500pt \hyperref[borderType]{border} & \hfuzz=500pt \\
\hline
\end{tabularx}


\subsection{TriangleVertex}
The zeroth level of densification
coincides with the 12 icosahedron vertices, as displayed in the upper left part
of Fig.~\ref{fig:triangle_grid}. Then, depending on the envisaged densification,
each triangle edge is divided into $n$ parts, illustrated in the upper right
part of Fig.~\ref{fig:triangle_grid}. The new nodes on the edges are then connected
by arcs of great circles parallel to the triangle edges. The intersections of
each three corresponding parallel lines become nodes of the densified grid as well.
As in case of a spherical triangle those three connecting lines do not exactly
intersect in one point, the center of the resulting triangle is used as location
for the new node (lower left part of Fig.~\ref{fig:triangle_grid}). The lower right
side of Fig.~\ref{fig:triangle_grid} finally shows the densified triangle vertex
grid for a level of $n=3$. The number of grid points in dependence of the chosen
level of densification can be calculated by
\begin{equation}\label{eq:numberVertex}
I=10\cdot(n+1)^2+2.
\end{equation}

\fig{!hb}{0.6}{icogrid}{fig:triangle_grid}{TriangleVertex grid.}


\keepXColumns
\begin{tabularx}{\textwidth}{N T A}
\hline
Name & Type & Annotation\\
\hline
\hfuzz=500pt\includegraphics[width=1em]{element-mustset.pdf}~level & \hfuzz=500pt uint & \hfuzz=500pt division of icosahedron, point count = 10*(n+1)**2+2\\
\hfuzz=500pt\includegraphics[width=1em]{element.pdf}~R & \hfuzz=500pt double & \hfuzz=500pt major axsis of the ellipsoid/sphere\\
\hfuzz=500pt\includegraphics[width=1em]{element.pdf}~inverseFlattening & \hfuzz=500pt double & \hfuzz=500pt flattening of the ellipsoid, 0: sphere\\
\hfuzz=500pt\includegraphics[width=1em]{element-unbounded.pdf}~border & \hfuzz=500pt \hyperref[borderType]{border} & \hfuzz=500pt \\
\hline
\end{tabularx}


\subsection{TriangleCenter}
The points of the zeroth level are located at the centers of the icosahedron triangles.
To achieve a finer grid, each of the triangles is divided into four smaller triangles by
connecting the midpoints of the triangle edges. The refined grid points are again located
at the center of the triangles. Subsequently, the triangles can be further densified up to
the desired level of densification $n$, which is defined by \config{level}.

The number of global grid points for a certain level can be determined by
\begin{equation}\label{eq:numberCenter}
I=20\cdot 4^n.
\end{equation}
Thus the quantity of grid points depends exponentially on the level $n$, as with
every additional level the number of grid points quadruplicates.


\keepXColumns
\begin{tabularx}{\textwidth}{N T A}
\hline
Name & Type & Annotation\\
\hline
\hfuzz=500pt\includegraphics[width=1em]{element-mustset.pdf}~level & \hfuzz=500pt uint & \hfuzz=500pt division of icosahedron, point count = 5*4**(n+1)\\
\hfuzz=500pt\includegraphics[width=1em]{element.pdf}~R & \hfuzz=500pt double & \hfuzz=500pt major axsis of the ellipsoid/sphere\\
\hfuzz=500pt\includegraphics[width=1em]{element.pdf}~inverseFlattening & \hfuzz=500pt double & \hfuzz=500pt flattening of the ellipsoid, 0: sphere\\
\hfuzz=500pt\includegraphics[width=1em]{element-unbounded.pdf}~border & \hfuzz=500pt \hyperref[borderType]{border} & \hfuzz=500pt \\
\hline
\end{tabularx}


\subsection{Gauss}
 The grid features equiangular spacing along circles of latitude with
 \config{parallelsCount} defining the number $L$ of the parallels.
\begin{equation}
\Delta\lambda=\frac{\pi}{L}\qquad\Rightarrow\qquad\lambda_i=\frac{\Delta\lambda}{2}+i\cdot\Delta\lambda\qquad\mbox{with}\qquad 0\leq i< 2L.
\end{equation}
Along the meridians the points are located at $L$ parallels at
the $L$ zeros $\vartheta_j$ of the Legendre polynomial of degree~$L$,
\begin{equation}
P_L(\cos\vartheta_j)=0.
\end{equation}
Consequently, the number of grid points sums up to
\begin{equation}
I=2\cdot L^2.
\end{equation}
The weights can be calculated according to
\begin{equation}
w_i(L)=\Delta\lambda\frac{2}{(1-t_i^2)(P'_{L}(\cos(\vartheta _i)))^2},\label{weights}
\end{equation}


\keepXColumns
\begin{tabularx}{\textwidth}{N T A}
\hline
Name & Type & Annotation\\
\hline
\hfuzz=500pt\includegraphics[width=1em]{element-mustset.pdf}~parallelsCount & \hfuzz=500pt uint & \hfuzz=500pt \\
\hfuzz=500pt\includegraphics[width=1em]{element.pdf}~R & \hfuzz=500pt double & \hfuzz=500pt major axsis of the ellipsoid/sphere\\
\hfuzz=500pt\includegraphics[width=1em]{element.pdf}~inverseFlattening & \hfuzz=500pt double & \hfuzz=500pt flattening of the ellipsoid, 0: sphere\\
\hfuzz=500pt\includegraphics[width=1em]{element-unbounded.pdf}~border & \hfuzz=500pt \hyperref[borderType]{border} & \hfuzz=500pt \\
\hline
\end{tabularx}


\subsection{Reuter}
The Reuter grid features equi-distant spacing along the meridians determined
by the control parameter~$\gamma$ according to
\begin{equation}
\Delta\vartheta=\frac{\pi}{\gamma}\qquad\Rightarrow\vartheta_j=j\Delta\vartheta,\qquad\mbox{with}\qquad 1\leq j\leq \gamma-1.
\end{equation}
Thus $\gamma+1$ denotes the number of points per meridian, as the two poles
are included in the point distribution as well. Along the circles of latitude,
the number of grid points decreases with increasing latitude in order to achieve
an evenly distributed point pattern. This number is chosen, so that the points
along each circle of latitude have the same spherical distance as two adjacent
latitudes. The resulting relationship is given by
\begin{equation}\label{eq:sphericalDistance}
\Delta\vartheta=\arccos\left( \cos^2\vartheta_j+\sin^2\vartheta_j\cos\Delta\lambda_j\right).
\end{equation}
The left hand side of this equation is the spherical distance between adjacent
latitudes, the right hand side stands for the spherical distance between two points
with the same polar distance $\vartheta_j$ and a longitudinal difference of
$\Delta\lambda_i$. This longitudinal distance can be adjusted depending on
$\vartheta_j$ to fulfill Eq.~\eqref{eq:sphericalDistance}. The resulting
formula for $\Delta\lambda_i$ is
\begin{equation}\label{eq:deltaLambdai}
\Delta\lambda_j=\arccos\left( \frac{\sin\Delta\vartheta -\cos^2\vartheta_j}{\sin^2\vartheta_j}\right).
\end{equation}
The number of points~$\gamma_j$ for each circle of latitude can then be determined by
\begin{equation}\label{eq:gammai}
\gamma_j=\left[ \frac{2\pi}{\Delta\lambda_j}\right] .
\end{equation}
Here the Gauss bracket $[x]$ specifies the largest integer equal to or less than $x$.
The longitudes are subsequently determined by
\begin{equation}
\lambda_{ij}=\frac{\Delta\lambda_j}{2}+i\cdot(2\pi/\gamma_j),\qquad \mbox{with\qquad}0\leq i< \gamma_j.
\end{equation}
The number of grid points can be estimated by
\begin{equation}\label{eq:numberReuter}
I=\leq 2+\frac{4}{\pi}\gamma^2,
\end{equation}
The $\leq$ results from the fact that the $\gamma_j$ are restricted to integer values.


\keepXColumns
\begin{tabularx}{\textwidth}{N T A}
\hline
Name & Type & Annotation\\
\hline
\hfuzz=500pt\includegraphics[width=1em]{element-mustset.pdf}~gamma & \hfuzz=500pt uint & \hfuzz=500pt number of parallels\\
\hfuzz=500pt\includegraphics[width=1em]{element.pdf}~height & \hfuzz=500pt double & \hfuzz=500pt ellipsoidal height\\
\hfuzz=500pt\includegraphics[width=1em]{element.pdf}~R & \hfuzz=500pt double & \hfuzz=500pt major axsis of the ellipsoid/sphere\\
\hfuzz=500pt\includegraphics[width=1em]{element.pdf}~inverseFlattening & \hfuzz=500pt double & \hfuzz=500pt flattening of the ellipsoid, 0: sphere\\
\hfuzz=500pt\includegraphics[width=1em]{element-unbounded.pdf}~border & \hfuzz=500pt \hyperref[borderType]{border} & \hfuzz=500pt \\
\hline
\end{tabularx}


\subsection{Corput}
This kind of grid distributes an arbitrarily chosen number of $I$ points
(defined by \config{globalPointsCount}) following a recursive, quasi random sequence.
In longitudinal direction the pattern follows
\begin{equation}
\Delta\lambda=\frac{2\pi}{I}\qquad\Rightarrow\qquad\frac{\Delta\lambda}{2}+\lambda_i=i\cdot\Delta\lambda\qquad\mbox{with}\qquad 1\leq i\leq I.
\end{equation}
This implies that every grid point features a unique longitude, with equi-angular
longitudinal differences.

The polar distance in the form $t_i=\cos\vartheta_i$ for each point is determined
by the following recursive sequence:
\begin{itemize}
\item Starting from an interval $t\in[-1,1]$.
\item If $I=1$, then the midpoint of the interval is returned as result of
the sequence, and the sequence is terminated.
\item If the number of points is uneven, the  midpoint is included into the list of $t_i$.
\item Subsequently, the interval is bisected into an upper and lower half,
       and the sequence is called for both halves.
\item $t$ from upper and lower half are alternately sorted into the list of $t_i$.
\item The polar distances are calculated by
\begin{equation}
\vartheta_i=\arccos\, t_i.
\end{equation}
\end{itemize}


\keepXColumns
\begin{tabularx}{\textwidth}{N T A}
\hline
Name & Type & Annotation\\
\hline
\hfuzz=500pt\includegraphics[width=1em]{element-mustset.pdf}~globalPointsCount & \hfuzz=500pt uint & \hfuzz=500pt \\
\hfuzz=500pt\includegraphics[width=1em]{element.pdf}~height & \hfuzz=500pt double & \hfuzz=500pt ellipsoidal height\\
\hfuzz=500pt\includegraphics[width=1em]{element.pdf}~R & \hfuzz=500pt double & \hfuzz=500pt major axsis of the ellipsoid/sphere\\
\hfuzz=500pt\includegraphics[width=1em]{element.pdf}~inverseFlattening & \hfuzz=500pt double & \hfuzz=500pt flattening of the ellipsoid, 0: sphere\\
\hfuzz=500pt\includegraphics[width=1em]{element-unbounded.pdf}~border & \hfuzz=500pt \hyperref[borderType]{border} & \hfuzz=500pt \\
\hline
\end{tabularx}


\subsection{Driscoll}
The Driscoll-Healy grid, has equiangular spacing along the meridians as well
as along the circles of latitude. In longitudinal direction (along the parallels),
these angular differences for a given \config{dimension} $L$ coincide with those
described for the corresponding geographical grid and Gauss grid. Along the meridians,
the size of the latitudinal differences is half the size compared to the geographical
grid. This results in the following point pattern,
\begin{equation}
\begin{split}
\Delta\lambda=\frac{\pi}{L}\qquad&\Rightarrow\qquad\lambda_i=\frac{\Delta\lambda}{2}+i\cdot\Delta\lambda\qquad&\mbox{with}\qquad 0\leq i< 2L, \\
\Delta\vartheta=\frac{\pi}{2L}\qquad&\Rightarrow\qquad\vartheta_j=j\cdot\Delta\vartheta\qquad&\mbox{with}\qquad 1\leq j\leq 2L.
\end{split}
\end{equation}
Consequently, the number of grid points is
\begin{equation}
I=4\cdot L^2.
\end{equation}
The weights are given by
\begin{equation}
w_i=\Delta\lambda\frac{4}{2L}\sin(\vartheta_i)\sum_{l=0}^{L-1}\frac{\sin\left[ (2l+1)\;\vartheta_i\right] }{2l+1}.
\end{equation}


\keepXColumns
\begin{tabularx}{\textwidth}{N T A}
\hline
Name & Type & Annotation\\
\hline
\hfuzz=500pt\includegraphics[width=1em]{element-mustset.pdf}~dimension & \hfuzz=500pt uint & \hfuzz=500pt number of parallels = 2*dimension\\
\hfuzz=500pt\includegraphics[width=1em]{element.pdf}~height & \hfuzz=500pt double & \hfuzz=500pt ellipsoidal height\\
\hfuzz=500pt\includegraphics[width=1em]{element.pdf}~R & \hfuzz=500pt double & \hfuzz=500pt major axsis of the ellipsoid/sphere\\
\hfuzz=500pt\includegraphics[width=1em]{element.pdf}~inverseFlattening & \hfuzz=500pt double & \hfuzz=500pt flattening of the ellipsoid, 0: sphere\\
\hfuzz=500pt\includegraphics[width=1em]{element-unbounded.pdf}~border & \hfuzz=500pt \hyperref[borderType]{border} & \hfuzz=500pt \\
\hline
\end{tabularx}


\subsection{SinglePoint}
Creates one single point.


\keepXColumns
\begin{tabularx}{\textwidth}{N T A}
\hline
Name & Type & Annotation\\
\hline
\hfuzz=500pt\includegraphics[width=1em]{element-mustset.pdf}~L & \hfuzz=500pt angle & \hfuzz=500pt longitude\\
\hfuzz=500pt\includegraphics[width=1em]{element-mustset.pdf}~B & \hfuzz=500pt angle & \hfuzz=500pt latitude\\
\hfuzz=500pt\includegraphics[width=1em]{element.pdf}~height & \hfuzz=500pt double & \hfuzz=500pt ellipsoidal height\\
\hfuzz=500pt\includegraphics[width=1em]{element.pdf}~area & \hfuzz=500pt double & \hfuzz=500pt associated area element on unit sphere\\
\hfuzz=500pt\includegraphics[width=1em]{element.pdf}~R & \hfuzz=500pt double & \hfuzz=500pt major axsis of the ellipsoid/sphere\\
\hfuzz=500pt\includegraphics[width=1em]{element.pdf}~inverseFlattening & \hfuzz=500pt double & \hfuzz=500pt flattening of the ellipsoid, 0: sphere\\
\hline
\end{tabularx}


\subsection{SinglePointCartesian}
Creates one single point.


\keepXColumns
\begin{tabularx}{\textwidth}{N T A}
\hline
Name & Type & Annotation\\
\hline
\hfuzz=500pt\includegraphics[width=1em]{element-mustset.pdf}~x & \hfuzz=500pt double & \hfuzz=500pt [m]\\
\hfuzz=500pt\includegraphics[width=1em]{element-mustset.pdf}~y & \hfuzz=500pt double & \hfuzz=500pt [m]\\
\hfuzz=500pt\includegraphics[width=1em]{element-mustset.pdf}~z & \hfuzz=500pt double & \hfuzz=500pt [m]\\
\hfuzz=500pt\includegraphics[width=1em]{element.pdf}~area & \hfuzz=500pt double & \hfuzz=500pt associated area element on unit sphere\\
\hline
\end{tabularx}


\subsection{File}\label{gridType:file}
In this class grid is read from a file, which is given by \configFile{inputfileGrid}{griddedData}.
A corresponding file can be generated with \program{GriddedDataCreate} or with \program{Matrix2GriddedData}.


\keepXColumns
\begin{tabularx}{\textwidth}{N T A}
\hline
Name & Type & Annotation\\
\hline
\hfuzz=500pt\includegraphics[width=1em]{element-mustset.pdf}~inputfileGrid & \hfuzz=500pt filename & \hfuzz=500pt \\
\hfuzz=500pt\includegraphics[width=1em]{element-unbounded.pdf}~border & \hfuzz=500pt \hyperref[borderType]{border} & \hfuzz=500pt \\
\hline
\end{tabularx}

\clearpage
%==================================

\section{InstrumentType}\label{instrumentTypeType}
Defines the type of an \file{instrument file}{instrument}.


\keepXColumns
\begin{tabularx}{\textwidth}{N T A}
\hline
Name & Type & Annotation\\
\hline
\hfuzz=500pt\includegraphics[width=1em]{element-mustset.pdf}~instrumentTypeType & \hfuzz=500pt choice & \hfuzz=500pt instrument type\\
\hfuzz=500pt\includegraphics[width=1em]{connector.pdf}\includegraphics[width=1em]{element-mustset.pdf}~INSTRUMENTTIME & \hfuzz=500pt  & \hfuzz=500pt time without data\\
\hfuzz=500pt\includegraphics[width=1em]{connector.pdf}\includegraphics[width=1em]{element-mustset.pdf}~MISCVALUE & \hfuzz=500pt  & \hfuzz=500pt single value\\
\hfuzz=500pt\includegraphics[width=1em]{connector.pdf}\includegraphics[width=1em]{element-mustset.pdf}~MISCVALUES & \hfuzz=500pt  & \hfuzz=500pt multiple values\\
\hfuzz=500pt\includegraphics[width=1em]{connector.pdf}\includegraphics[width=1em]{element-mustset.pdf}~VECTOR3D & \hfuzz=500pt  & \hfuzz=500pt x, y, z\\
\hfuzz=500pt\includegraphics[width=1em]{connector.pdf}\includegraphics[width=1em]{element-mustset.pdf}~COVARIANCE3D & \hfuzz=500pt  & \hfuzz=500pt xx, yy, zz, xy, xz, yz\\
\hfuzz=500pt\includegraphics[width=1em]{connector.pdf}\includegraphics[width=1em]{element-mustset.pdf}~ORBIT & \hfuzz=500pt  & \hfuzz=500pt position [m], velocity [m/s], acceleration [m/s\textasciicircum{}2] (each x, y, z)\\
\hfuzz=500pt\includegraphics[width=1em]{connector.pdf}\includegraphics[width=1em]{element-mustset.pdf}~STARCAMERA & \hfuzz=500pt  & \hfuzz=500pt quaternions (q0, qx, qy, qz)\\
\hfuzz=500pt\includegraphics[width=1em]{connector.pdf}\includegraphics[width=1em]{element-mustset.pdf}~ACCELEROMETER & \hfuzz=500pt  & \hfuzz=500pt x, y, z [m/s\textasciicircum{}2]\\
\hfuzz=500pt\includegraphics[width=1em]{connector.pdf}\includegraphics[width=1em]{element-mustset.pdf}~SATELLITETRACKING & \hfuzz=500pt  & \hfuzz=500pt range [m], range rate [m/s], range acceleration [m/s\textasciicircum{}2]\\
\hfuzz=500pt\includegraphics[width=1em]{connector.pdf}\includegraphics[width=1em]{element-mustset.pdf}~GRADIOMETER & \hfuzz=500pt  & \hfuzz=500pt xx, yy, zz, xy, xz, yz [1/s\textasciicircum{}2]\\
\hfuzz=500pt\includegraphics[width=1em]{connector.pdf}\includegraphics[width=1em]{element-mustset.pdf}~GNSSRECEIVER & \hfuzz=500pt  & \hfuzz=500pt GNSS phase/code observations [m]\\
\hfuzz=500pt\includegraphics[width=1em]{connector.pdf}\includegraphics[width=1em]{element-mustset.pdf}~OBSERVATIONSIGMA & \hfuzz=500pt  & \hfuzz=500pt accuracy\\
\hfuzz=500pt\includegraphics[width=1em]{connector.pdf}\includegraphics[width=1em]{element-mustset.pdf}~MASS & \hfuzz=500pt  & \hfuzz=500pt \\
\hfuzz=500pt\includegraphics[width=1em]{connector.pdf}\includegraphics[width=1em]{element-mustset.pdf}~THRUSTER & \hfuzz=500pt  & \hfuzz=500pt \\
\hfuzz=500pt\includegraphics[width=1em]{connector.pdf}\includegraphics[width=1em]{element-mustset.pdf}~MAGNETOMETER & \hfuzz=500pt  & \hfuzz=500pt \\
\hfuzz=500pt\includegraphics[width=1em]{connector.pdf}\includegraphics[width=1em]{element-mustset.pdf}~ACCHOUSEKEEPING & \hfuzz=500pt  & \hfuzz=500pt \\
\hline
\end{tabularx}

\clearpage
%==================================

\section{Kernel}\label{kernelType}
Kernel defines harmonic isotropic integral kernels $K$.
\begin{equation}
T(P) = \frac{1}{4\pi}\int_\Omega K(P,Q)\cdot f(Q)\,d\Omega(Q),
\end{equation}
where $T$ is the (disturbance)potential and $f$ is a functional on the spherical surface~$\Omega$.
The Kernel can be exapanded into a series of (fully normalized) legendre polynomials
\begin{equation}\label{eq.kernel}
K(\cos\psi,r,R) = \sum_n \left(\frac{R}{r}\right)^{n+1}
k_n\sqrt{2n+1}\bar{P}_n(\cos\psi).
\end{equation}
On the one hand the kernel defines the type of the functionals~$f$ that are measured
or have to be computed, e.g. gravity anomalies given by the Stokes-kernel.
On the other hand the kernel functions can be used as basis functions to represent
the gravity field, e.g. as spline functions or wavelets.


\subsection{GeoidHeight}\label{kernelType:geoidHeight}
The geoid height is defined by Bruns formula
\begin{equation}
N = \frac{1}{\gamma}T
\end{equation}
with $T$ the disturbance potential and the normal gravity
\begin{equation}\label{normalgravity}
\gamma  = \gamma_0 - 0.30877\cdot 10^{-5}/s^2(1-0.00142\sin^2(B))h
\end{equation}
and
\begin{equation}
\gamma_0 = 9.780327\,m/s^2(1+0.0053024\sin^2(B)-0.0000058\sin^2(2B))
\end{equation}
where $h$ is the ellipsoidal height in meter and $B$ the longitude.

The kernel is given by
\begin{equation}
K(\cos\psi,r,R) = \gamma\frac{R(r^2-R^2)}{l^3},
\end{equation}
and the coefficients in \eqref{eq.kernel} are
\begin{equation}
k_n = \gamma.
\end{equation}


\subsection{Anomalies}
Gravity anomalies in linearized form are defined by
\begin{equation}
\Delta g = -\frac{\partial T}{\partial r}-\frac{2}{r}T.
\end{equation}
The Stokes kernel is given by
\begin{equation}
K(\cos\psi,r,R) = \frac{2R^2}{l}-3\frac{Rl}{r^2}-\frac{R^2}{r^2}\cos\psi
\left(5+3\ln\frac{l+r-R\cos\psi}{2r}\right),
\end{equation}
and the coefficients in \eqref{eq.kernel} are
\begin{equation}
k_n = \frac{R}{n-1}.
\end{equation}


\subsection{Disturbance}\label{kernelType:disturbance}
Gravity disturbances in linearized form are defined by
\begin{equation}
\delta g = -\frac{dT}{dr}.
\end{equation}
The Hotine kernel is given by
\begin{equation}
K(\cos\psi,r,R) = \frac{2R^2}{l}-R\ln\frac{l+R-r\cos\psi}{r(1-\cos\psi)},
\end{equation}
and the coefficients in \eqref{eq.kernel} are
\begin{equation}
k_n = \frac{R}{n+1}.
\end{equation}


\subsection{Potential}
The Abel-Poisson kernel is given by
\begin{equation}
K(\cos\psi,r,R) = \frac{R(r^2-R^2)}{l^3},
\end{equation}
and the coefficients in \eqref{eq.kernel} are
\begin{equation}
k_n = 1.
\end{equation}


\subsection{Density}
This kernel defines a point mass or mass on a single layer ($1/l$-kernel)
taking the effect of the loading into account.

The coefficients of the kernel defined in \eqref{eq.kernel} are
\begin{equation}
k_n = 4\pi G R\frac{1+k_n'}{2n+1},
\end{equation}
where $G$ is the gravitational constant and $k_n'$ are the load Love numbers.


\keepXColumns
\begin{tabularx}{\textwidth}{N T A}
\hline
Name & Type & Annotation\\
\hline
\hfuzz=500pt\includegraphics[width=1em]{element.pdf}~inputfileLoadingLoveNumber & \hfuzz=500pt filename & \hfuzz=500pt \\
\hline
\end{tabularx}


\subsection{WaterHeight}\label{kernelType:waterHeight}
Height of equivalent water columns taking the effect of the loading into account.

The coefficients of the kernel defined in \eqref{eq.kernel} are
\begin{equation}
k_n = 4\pi G \rho R\frac{1+k_n'}{2n+1},
\end{equation}
where $G$ is the gravitational constant, $\rho$ is the \config{density} of water
and $k_n'$ are the load Love numbers.


\keepXColumns
\begin{tabularx}{\textwidth}{N T A}
\hline
Name & Type & Annotation\\
\hline
\hfuzz=500pt\includegraphics[width=1em]{element-mustset.pdf}~density & \hfuzz=500pt double & \hfuzz=500pt [kg/m**3]\\
\hfuzz=500pt\includegraphics[width=1em]{element.pdf}~inputfileLoadingLoveNumber & \hfuzz=500pt filename & \hfuzz=500pt \\
\hline
\end{tabularx}


\subsection{BottomPressure}
Ocean bottom pressure caused by water and atmosphere masses columns taking the effect of the loading into account.

The coefficients of the kernel defined in \eqref{eq.kernel} are
\begin{equation}
k_n = \frac{4\pi G R }{\gamma}\frac{1+k_n'}{2n+1},
\end{equation}
where $G$ is the gravitational constant, $\gamma$ is the normal gravity and $k_n'$ are the load Love numbers.


\keepXColumns
\begin{tabularx}{\textwidth}{N T A}
\hline
Name & Type & Annotation\\
\hline
\hfuzz=500pt\includegraphics[width=1em]{element.pdf}~inputfileLoadingLoveNumber & \hfuzz=500pt filename & \hfuzz=500pt \\
\hline
\end{tabularx}


\subsection{Deformation}
Computes the radial deformation caused by loading.

The coefficients of the kernel defined in \eqref{eq.kernel} are
\begin{equation}
k_n = \gamma\frac{1+k_n'}{h_n'},
\end{equation}
where $\gamma$ is the normal gravity defined in \eqref{normalgravity},
$h_n'$ and $k_n'$ are the load Love numbers and the load deformation Love numbers.


\keepXColumns
\begin{tabularx}{\textwidth}{N T A}
\hline
Name & Type & Annotation\\
\hline
\hfuzz=500pt\includegraphics[width=1em]{element-mustset.pdf}~inputfileDeformationLoadLoveNumber & \hfuzz=500pt filename & \hfuzz=500pt \\
\hfuzz=500pt\includegraphics[width=1em]{element.pdf}~inputfilePotentialLoadLoveNumber & \hfuzz=500pt filename & \hfuzz=500pt if full potential is given and not only loading potential\\
\hline
\end{tabularx}


\subsection{RadialGradient}
This kernel defines the second radial derivative of the (disturbance) potential.
\begin{equation}
T_{rr} = \frac{\partial^2 T}{\partial r^2}.
\end{equation}
The coefficients of the kernel defined in \eqref{eq.kernel} are
\begin{equation}
k_n = \frac{r^2}{(n+1)(n+2)}.
\end{equation}


\subsection{Coefficients}\label{kernelType:coefficients}
The kernel is defined by the coefficients $k_n$ given by file.


\keepXColumns
\begin{tabularx}{\textwidth}{N T A}
\hline
Name & Type & Annotation\\
\hline
\hfuzz=500pt\includegraphics[width=1em]{element-mustset.pdf}~inputfileCoefficients & \hfuzz=500pt filename & \hfuzz=500pt \\
\hline
\end{tabularx}


\subsection{FilterGauss}
Another \configClass{kernel}{kernelType} is smoothed by a gauss filter
which is defined by
\begin{equation}
F(\cos\psi) = \frac{b\cdot e^{-b(1-\cos\psi)}}{1-e^{-2b}}
\end{equation}
with $b = \frac{ln(2)}{1-\cos(r/R)}$ where $r$ is the given
smoothing \config{radius} in km and $R=6378.1366$~km is the
Earth radius.
The coefficients~$k_n$ of the \config{kernel} are multiplicated by
\begin{equation}
f_n = \frac{1}{2n+1} \int_{-1}^1 F(t)\cdot \bar{P}_n(t)\,dt.
\end{equation}


\keepXColumns
\begin{tabularx}{\textwidth}{N T A}
\hline
Name & Type & Annotation\\
\hline
\hfuzz=500pt\includegraphics[width=1em]{element-mustset.pdf}~kernel & \hfuzz=500pt \hyperref[kernelType]{kernel} & \hfuzz=500pt \\
\hfuzz=500pt\includegraphics[width=1em]{element-mustset.pdf}~radius & \hfuzz=500pt double & \hfuzz=500pt filter radius [km]\\
\hline
\end{tabularx}


\subsection{BlackmanLowpass}
Another \configClass{kernel}{kernelType} is smoothed by a Blackman low-pass filter. The filter is
defined through the beginning and end of the transition from pass-band to stop-band. This
transition band is specified by \config{startDegreeTransition} ($n_1$) and \config{stopDegreeTransition} ($n_2$).

The coefficients of this kernel are defined as
\begin{equation}
\begin{cases}
1 & \text{for } n < n_1 \\
A_n^2 & \text{for } n_1\leq n \leq n_2 \\
0 & \text{for } n > n_2 \\
\end{cases}
\end{equation}
with
\begin{equation}
A_n = 0.42 + 0.5\cos(\pi \frac{n-n_1}{n_2-n_1}) + 0.08 \cos(2\pi\frac{n-n_1}{n_2-n_1}).
\end{equation}


\keepXColumns
\begin{tabularx}{\textwidth}{N T A}
\hline
Name & Type & Annotation\\
\hline
\hfuzz=500pt\includegraphics[width=1em]{element-mustset.pdf}~kernel & \hfuzz=500pt \hyperref[kernelType]{kernel} & \hfuzz=500pt \\
\hfuzz=500pt\includegraphics[width=1em]{element-mustset.pdf}~startDegreeTransition & \hfuzz=500pt uint & \hfuzz=500pt minimum degree in transition band\\
\hfuzz=500pt\includegraphics[width=1em]{element-mustset.pdf}~stopDegreeTransition & \hfuzz=500pt uint & \hfuzz=500pt maximum degree in transition band\\
\hline
\end{tabularx}


\subsection{Truncation}
Another \configClass{kernel}{kernelType} is truncated before \config{minDegree} and after \config{maxDegree}.
The coefficients of this kernel are defined as
\begin{equation}
  k_n =
  \begin{cases}
  1 & \text{for } n_{\text{minDegree}} \leq n \leq n_{\text{maxDegree}}\\
  0 & \text{else.} \\
  \end{cases}
\end{equation}


\keepXColumns
\begin{tabularx}{\textwidth}{N T A}
\hline
Name & Type & Annotation\\
\hline
\hfuzz=500pt\includegraphics[width=1em]{element-mustset.pdf}~kernel & \hfuzz=500pt \hyperref[kernelType]{kernel} & \hfuzz=500pt \\
\hfuzz=500pt\includegraphics[width=1em]{element.pdf}~mminDegree & \hfuzz=500pt uint & \hfuzz=500pt truncate before minDegree\\
\hfuzz=500pt\includegraphics[width=1em]{element-mustset.pdf}~maxDegree & \hfuzz=500pt uint & \hfuzz=500pt truncate after maxDegree\\
\hline
\end{tabularx}


\subsection{SelenoidHeight}
The selenoid height is defined by Bruns formula
\begin{equation}
N = \frac{1}{\gamma}T
\end{equation}
with $T$ the disturbance potential and the normal gravity $\gamma=\frac{GM}{R^2}$ of the moon.

The kernel is given by
\begin{equation}
K(\cos\psi,r,R) = \gamma\frac{R(r^2-R^2)}{l^3},
\end{equation}
and the coefficients in \eqref{eq.kernel} are
\begin{equation}
k_n = \gamma.
\end{equation}

\clearpage
%==================================

\section{Loop}\label{loopType}
Generates a sequence with variables to loop over.
The variable names can be set with \config{variableLoop...} and
the current values are assigned to the variables for each loop step.
See \reference{Loop and conditions}{general.loopsAndConditions} for usage.


\subsection{TimeSeries}
Loop over points in time.


\keepXColumns
\begin{tabularx}{\textwidth}{N T A}
\hline
Name & Type & Annotation\\
\hline
\hfuzz=500pt\includegraphics[width=1em]{element-mustset-unbounded.pdf}~timeSeries & \hfuzz=500pt \hyperref[timeSeriesType]{timeSeries} & \hfuzz=500pt loop is called for every point in time\\
\hfuzz=500pt\includegraphics[width=1em]{element.pdf}~variableLoopTime & \hfuzz=500pt string & \hfuzz=500pt variable with time of each loop\\
\hfuzz=500pt\includegraphics[width=1em]{element.pdf}~variableLoopIndex & \hfuzz=500pt string & \hfuzz=500pt variable with index of current iteration (starts with zero)\\
\hfuzz=500pt\includegraphics[width=1em]{element.pdf}~variableLoopCount & \hfuzz=500pt string & \hfuzz=500pt variable with total number of iterations\\
\hline
\end{tabularx}


\subsection{TimeIntervals}
Loop over the intervals between points in time.


\keepXColumns
\begin{tabularx}{\textwidth}{N T A}
\hline
Name & Type & Annotation\\
\hline
\hfuzz=500pt\includegraphics[width=1em]{element-mustset-unbounded.pdf}~timeIntervals & \hfuzz=500pt \hyperref[timeSeriesType]{timeSeries} & \hfuzz=500pt loop is called for every interval\\
\hfuzz=500pt\includegraphics[width=1em]{element.pdf}~variableLoopTimeStart & \hfuzz=500pt string & \hfuzz=500pt variable with starting time of each interval\\
\hfuzz=500pt\includegraphics[width=1em]{element.pdf}~variableLoopTimeEnd & \hfuzz=500pt string & \hfuzz=500pt variable with ending time of each interval\\
\hfuzz=500pt\includegraphics[width=1em]{element.pdf}~variableLoopIndex & \hfuzz=500pt string & \hfuzz=500pt variable with index of current iteration (starts with zero)\\
\hfuzz=500pt\includegraphics[width=1em]{element.pdf}~variableLoopCount & \hfuzz=500pt string & \hfuzz=500pt variable with total number of iterations\\
\hline
\end{tabularx}


\subsection{ManualList}\label{loopType:manualList}
Loop over list of strings.


\keepXColumns
\begin{tabularx}{\textwidth}{N T A}
\hline
Name & Type & Annotation\\
\hline
\hfuzz=500pt\includegraphics[width=1em]{element-mustset-unbounded.pdf}~string & \hfuzz=500pt string & \hfuzz=500pt explicit list of strings\\
\hfuzz=500pt\includegraphics[width=1em]{element.pdf}~variableLoopString & \hfuzz=500pt string & \hfuzz=500pt name of the variable to be replaced\\
\hfuzz=500pt\includegraphics[width=1em]{element.pdf}~variableLoopIndex & \hfuzz=500pt string & \hfuzz=500pt variable with index of current iteration (starts with zero)\\
\hfuzz=500pt\includegraphics[width=1em]{element.pdf}~variableLoopCount & \hfuzz=500pt string & \hfuzz=500pt variable with total number of iterations\\
\hline
\end{tabularx}


\subsection{ManualTable}
Loop over rows of a table containing strings.
The table can be defined \config{rowWise} or \config{columnWise}.
Each row/column must have the same number of cells.


\keepXColumns
\begin{tabularx}{\textwidth}{N T A}
\hline
Name & Type & Annotation\\
\hline
\hfuzz=500pt\includegraphics[width=1em]{element-mustset.pdf}~table & \hfuzz=500pt choice & \hfuzz=500pt define table by rows/columns\\
\hfuzz=500pt\includegraphics[width=1em]{connector.pdf}\includegraphics[width=1em]{element-mustset.pdf}~rowWise & \hfuzz=500pt sequence & \hfuzz=500pt define table by rows\\
\hfuzz=500pt\quad\includegraphics[width=1em]{connector.pdf}\includegraphics[width=1em]{element-mustset-unbounded.pdf}~row & \hfuzz=500pt sequence & \hfuzz=500pt define table by rows\\
\hfuzz=500pt\quad\quad\includegraphics[width=1em]{connector.pdf}\includegraphics[width=1em]{element-mustset-unbounded.pdf}~cell & \hfuzz=500pt string & \hfuzz=500pt explicit list of cells in row/column\\
\hfuzz=500pt\includegraphics[width=1em]{connector.pdf}\includegraphics[width=1em]{element-mustset.pdf}~columnWise & \hfuzz=500pt sequence & \hfuzz=500pt define table by columns\\
\hfuzz=500pt\quad\includegraphics[width=1em]{connector.pdf}\includegraphics[width=1em]{element-mustset-unbounded.pdf}~column & \hfuzz=500pt sequence & \hfuzz=500pt define table by columns\\
\hfuzz=500pt\quad\quad\includegraphics[width=1em]{connector.pdf}\includegraphics[width=1em]{element-mustset-unbounded.pdf}~cell & \hfuzz=500pt string & \hfuzz=500pt explicit list of cells in row/column\\
\hfuzz=500pt\includegraphics[width=1em]{element-mustset-unbounded.pdf}~variableLoopString & \hfuzz=500pt string & \hfuzz=500pt 1. variable name for the 1. column, next variable name for the 2. column, ... \\
\hfuzz=500pt\includegraphics[width=1em]{element.pdf}~variableLoopIndex & \hfuzz=500pt string & \hfuzz=500pt variable with index of current iteration (starts with zero)\\
\hfuzz=500pt\includegraphics[width=1em]{element.pdf}~variableLoopCount & \hfuzz=500pt string & \hfuzz=500pt variable with total number of iterations\\
\hline
\end{tabularx}


\subsection{FileAscii}
Loop over list of strings from files.


\keepXColumns
\begin{tabularx}{\textwidth}{N T A}
\hline
Name & Type & Annotation\\
\hline
\hfuzz=500pt\includegraphics[width=1em]{element-mustset-unbounded.pdf}~inputfile & \hfuzz=500pt filename & \hfuzz=500pt simple ASCII file with strings (separated by whitespace)\\
\hfuzz=500pt\includegraphics[width=1em]{element.pdf}~sort & \hfuzz=500pt boolean & \hfuzz=500pt sort entries alphabetically (ascending)\\
\hfuzz=500pt\includegraphics[width=1em]{element.pdf}~removeDuplicates & \hfuzz=500pt boolean & \hfuzz=500pt remove duplicate entries (order is preserved)\\
\hfuzz=500pt\includegraphics[width=1em]{element.pdf}~startIndex & \hfuzz=500pt uint & \hfuzz=500pt start at element startIndex (counting from 0)\\
\hfuzz=500pt\includegraphics[width=1em]{element.pdf}~count & \hfuzz=500pt uint & \hfuzz=500pt use count elements (default: use all)\\
\hfuzz=500pt\includegraphics[width=1em]{element.pdf}~variableLoopString & \hfuzz=500pt string & \hfuzz=500pt name of the variable to be replaced\\
\hfuzz=500pt\includegraphics[width=1em]{element.pdf}~variableLoopIndex & \hfuzz=500pt string & \hfuzz=500pt variable with index of current iteration (starts with zero)\\
\hfuzz=500pt\includegraphics[width=1em]{element.pdf}~variableLoopCount & \hfuzz=500pt string & \hfuzz=500pt variable with total number of iterations\\
\hline
\end{tabularx}


\subsection{FileAsciiTable}
Loop over rows of a table containing strings.
Each row must have the same number of columns.


\keepXColumns
\begin{tabularx}{\textwidth}{N T A}
\hline
Name & Type & Annotation\\
\hline
\hfuzz=500pt\includegraphics[width=1em]{element-mustset.pdf}~inputfile & \hfuzz=500pt filename & \hfuzz=500pt simple ASCII file with mutiple columns (separated by whitespace)\\
\hfuzz=500pt\includegraphics[width=1em]{element.pdf}~startLine & \hfuzz=500pt uint & \hfuzz=500pt start at line startLine (counting from 0)\\
\hfuzz=500pt\includegraphics[width=1em]{element.pdf}~countLines & \hfuzz=500pt uint & \hfuzz=500pt read count lines (default: all)\\
\hfuzz=500pt\includegraphics[width=1em]{element-mustset-unbounded.pdf}~variableLoopString & \hfuzz=500pt string & \hfuzz=500pt 1. variable name for the 1. column, next variable name for the 2. column, ... \\
\hfuzz=500pt\includegraphics[width=1em]{element.pdf}~variableLoopIndex & \hfuzz=500pt string & \hfuzz=500pt variable with index of current iteration (starts with zero)\\
\hfuzz=500pt\includegraphics[width=1em]{element.pdf}~variableLoopCount & \hfuzz=500pt string & \hfuzz=500pt variable with total number of iterations\\
\hline
\end{tabularx}


\subsection{Matrix}
Loop over rows of a matrix. To define the loop variables the standard
data variables of the matrix are available, see~\reference{dataVariables}{general.parser:dataVariables}.


\keepXColumns
\begin{tabularx}{\textwidth}{N T A}
\hline
Name & Type & Annotation\\
\hline
\hfuzz=500pt\includegraphics[width=1em]{element-mustset.pdf}~inputfile & \hfuzz=500pt filename & \hfuzz=500pt \\
\hfuzz=500pt\includegraphics[width=1em]{element.pdf}~transpose & \hfuzz=500pt boolean & \hfuzz=500pt effectively loop over columns\\
\hfuzz=500pt\includegraphics[width=1em]{element.pdf}~startRow & \hfuzz=500pt expression & \hfuzz=500pt start at this row (variable: rows)\\
\hfuzz=500pt\includegraphics[width=1em]{element.pdf}~countRows & \hfuzz=500pt expression & \hfuzz=500pt use this many rows (variable: rows)\\
\hfuzz=500pt\includegraphics[width=1em]{element-mustset-unbounded.pdf}~variableLoop & \hfuzz=500pt expression & \hfuzz=500pt define a variable by name = expression (input columns are named data0, data1, ...)\\
\hfuzz=500pt\includegraphics[width=1em]{element.pdf}~variableLoopIndex & \hfuzz=500pt string & \hfuzz=500pt variable with index of current iteration (starts with zero)\\
\hfuzz=500pt\includegraphics[width=1em]{element.pdf}~variableLoopCount & \hfuzz=500pt string & \hfuzz=500pt variable with total number of iterations\\
\hline
\end{tabularx}


\subsection{UniformSampling}
Loop over sequence of numbers.


\keepXColumns
\begin{tabularx}{\textwidth}{N T A}
\hline
Name & Type & Annotation\\
\hline
\hfuzz=500pt\includegraphics[width=1em]{element-mustset.pdf}~rangeStart & \hfuzz=500pt double & \hfuzz=500pt start of range\\
\hfuzz=500pt\includegraphics[width=1em]{element-mustset.pdf}~rangeEnd & \hfuzz=500pt double & \hfuzz=500pt end of range (inclusive)\\
\hfuzz=500pt\includegraphics[width=1em]{element-mustset.pdf}~sampling & \hfuzz=500pt double & \hfuzz=500pt sampling\\
\hfuzz=500pt\includegraphics[width=1em]{element.pdf}~variableLoopNumber & \hfuzz=500pt string & \hfuzz=500pt name of the variable to be replaced\\
\hfuzz=500pt\includegraphics[width=1em]{element.pdf}~variableLoopIndex & \hfuzz=500pt string & \hfuzz=500pt variable with index of current iteration (starts with zero)\\
\hfuzz=500pt\includegraphics[width=1em]{element.pdf}~variableLoopCount & \hfuzz=500pt string & \hfuzz=500pt variable with total number of iterations\\
\hline
\end{tabularx}


\subsection{CommandOutput}\label{loopType:commandOutput}
Loop over lines of command output.


\keepXColumns
\begin{tabularx}{\textwidth}{N T A}
\hline
Name & Type & Annotation\\
\hline
\hfuzz=500pt\includegraphics[width=1em]{element-mustset-unbounded.pdf}~command & \hfuzz=500pt filename & \hfuzz=500pt each output line becomes a loop iteration\\
\hfuzz=500pt\includegraphics[width=1em]{element.pdf}~silently & \hfuzz=500pt boolean & \hfuzz=500pt without showing the output.\\
\hfuzz=500pt\includegraphics[width=1em]{element.pdf}~variableLoopString & \hfuzz=500pt string & \hfuzz=500pt name of the variable to be replaced\\
\hfuzz=500pt\includegraphics[width=1em]{element.pdf}~variableLoopIndex & \hfuzz=500pt string & \hfuzz=500pt variable with index of current iteration (starts with zero)\\
\hfuzz=500pt\includegraphics[width=1em]{element.pdf}~variableLoopCount & \hfuzz=500pt string & \hfuzz=500pt variable with total number of iterations\\
\hline
\end{tabularx}


\subsection{Loop}
Loop over nested loops. First \config{loop} is outermost loop, every subsequent \config{loop} is one level below the previous \config{loop}.


\keepXColumns
\begin{tabularx}{\textwidth}{N T A}
\hline
Name & Type & Annotation\\
\hline
\hfuzz=500pt\includegraphics[width=1em]{element-mustset-unbounded.pdf}~loop & \hfuzz=500pt \hyperref[loopType]{loop} & \hfuzz=500pt subloop\\
\hfuzz=500pt\includegraphics[width=1em]{element.pdf}~variableLoopIndex & \hfuzz=500pt string & \hfuzz=500pt variable with index of current iteration (starts with zero)\\
\hline
\end{tabularx}


\subsection{FileGnssStationInfo}
Loop over receivers or antennas in a GNSS station/transmitter info \file{file}{gnssStationInfo}.


\keepXColumns
\begin{tabularx}{\textwidth}{N T A}
\hline
Name & Type & Annotation\\
\hline
\hfuzz=500pt\includegraphics[width=1em]{element-mustset.pdf}~inputfileGnssStationInfo & \hfuzz=500pt filename & \hfuzz=500pt station/transmitter info file\\
\hfuzz=500pt\includegraphics[width=1em]{element-mustset.pdf}~infoType & \hfuzz=500pt choice & \hfuzz=500pt info to loop over\\
\hfuzz=500pt\includegraphics[width=1em]{connector.pdf}\includegraphics[width=1em]{element-mustset.pdf}~antenna & \hfuzz=500pt  & \hfuzz=500pt loop over antennas\\
\hfuzz=500pt\includegraphics[width=1em]{connector.pdf}\includegraphics[width=1em]{element-mustset.pdf}~receiver & \hfuzz=500pt  & \hfuzz=500pt loop over receivers\\
\hfuzz=500pt\includegraphics[width=1em]{element.pdf}~variableLoopName & \hfuzz=500pt string & \hfuzz=500pt variable with antenna/receiver name\\
\hfuzz=500pt\includegraphics[width=1em]{element.pdf}~variableLoopSerial & \hfuzz=500pt string & \hfuzz=500pt variable with antenna/receiver serial\\
\hfuzz=500pt\includegraphics[width=1em]{element.pdf}~variableLoopInfo & \hfuzz=500pt string & \hfuzz=500pt variable with radome (antenna) or version (receiver)\\
\hfuzz=500pt\includegraphics[width=1em]{element.pdf}~variableLoopTimeStart & \hfuzz=500pt string & \hfuzz=500pt variable with antenna/receiver start time\\
\hfuzz=500pt\includegraphics[width=1em]{element.pdf}~variableLoopTimeEnd & \hfuzz=500pt string & \hfuzz=500pt variable with antenna/receiver end time\\
\hfuzz=500pt\includegraphics[width=1em]{element.pdf}~variableLoopIndex & \hfuzz=500pt string & \hfuzz=500pt variable with index of current iteration (starts with zero)\\
\hfuzz=500pt\includegraphics[width=1em]{element.pdf}~variableLoopCount & \hfuzz=500pt string & \hfuzz=500pt variable with total number of iterations\\
\hline
\end{tabularx}

\clearpage
%==================================

\section{Magnetosphere}\label{magnetosphereType}
This class provides functions of the magnetic field of the Earth.


\subsection{IGRF}
International Geomagnetic Reference Field.

\clearpage
%==================================

\section{MatrixGenerator}\label{matrixGeneratorType}
This class provides a matrix used e.g. by \program{MatrixCalculate}.
If multiple matrices are given the resulting matrix is the sum all
and the size is exandeded to fit all matrices. Before the computation of each submatrix
the variables \verb|rowsBefore| and \verb|columnsBefore| with current size of the overall matrix
are set. As all matrices can be manipulated before, complex matrix operations are possible.


\subsection{File}
Matrix from \file{file}{matrix}.


\keepXColumns
\begin{tabularx}{\textwidth}{N T A}
\hline
Name & Type & Annotation\\
\hline
\hfuzz=500pt\includegraphics[width=1em]{element-mustset.pdf}~inputfileMatrix & \hfuzz=500pt filename & \hfuzz=500pt \\
\hfuzz=500pt\includegraphics[width=1em]{element.pdf}~factor & \hfuzz=500pt double & \hfuzz=500pt \\
\hline
\end{tabularx}


\subsection{Normals file}
Matrix from a \file{normal equation file}{normalEquation}. The symmetric normal matrix,
the right hand side vector or the lPl vector can be selected.


\keepXColumns
\begin{tabularx}{\textwidth}{N T A}
\hline
Name & Type & Annotation\\
\hline
\hfuzz=500pt\includegraphics[width=1em]{element-mustset.pdf}~inputfileNormalEquation & \hfuzz=500pt filename & \hfuzz=500pt \\
\hfuzz=500pt\includegraphics[width=1em]{element-mustset.pdf}~type & \hfuzz=500pt choice & \hfuzz=500pt \\
\hfuzz=500pt\includegraphics[width=1em]{connector.pdf}\includegraphics[width=1em]{element-mustset.pdf}~normalMatrix & \hfuzz=500pt  & \hfuzz=500pt \\
\hfuzz=500pt\includegraphics[width=1em]{connector.pdf}\includegraphics[width=1em]{element-mustset.pdf}~rightHandSide & \hfuzz=500pt  & \hfuzz=500pt \\
\hfuzz=500pt\includegraphics[width=1em]{connector.pdf}\includegraphics[width=1em]{element-mustset.pdf}~lPl & \hfuzz=500pt  & \hfuzz=500pt \\
\hfuzz=500pt\includegraphics[width=1em]{element.pdf}~factor & \hfuzz=500pt double & \hfuzz=500pt \\
\hline
\end{tabularx}


\subsection{Expression}
Matrix filled by an expression. For each element of the new matrix the variables
\verb|row| and \verb|column| are set and the expression \config{element} is evaluated.

Excample: The \config{element}=\verb|if(row==column,1,0)| generates an identity matrix.


\keepXColumns
\begin{tabularx}{\textwidth}{N T A}
\hline
Name & Type & Annotation\\
\hline
\hfuzz=500pt\includegraphics[width=1em]{element-mustset.pdf}~rows & \hfuzz=500pt expression & \hfuzz=500pt (variables: rowsBefore, columnsBefore)\\
\hfuzz=500pt\includegraphics[width=1em]{element-mustset.pdf}~columns & \hfuzz=500pt expression & \hfuzz=500pt (variables: rowsBefore, columnsBefore)\\
\hfuzz=500pt\includegraphics[width=1em]{element-mustset.pdf}~element & \hfuzz=500pt expression & \hfuzz=500pt for each element of matrix (variables: row, column, rows, columns, rowsBefore, columnsBefore)\\
\hline
\end{tabularx}


\subsection{Element manipulation}
The elements of a matrix are replaced an expression.
For each element of the matrix the variables \verb|data|, \verb|row|, \verb|column|
are set and the expression \config{element} is evaluated and replaces the element.
Additionally the standard data variables are available (assigned each row),
see~\reference{dataVariables}{general.parser:dataVariables}.


\keepXColumns
\begin{tabularx}{\textwidth}{N T A}
\hline
Name & Type & Annotation\\
\hline
\hfuzz=500pt\includegraphics[width=1em]{element-mustset-unbounded.pdf}~matrix & \hfuzz=500pt \hyperref[matrixGeneratorType]{matrixGenerator} & \hfuzz=500pt \\
\hfuzz=500pt\includegraphics[width=1em]{element-mustset.pdf}~element & \hfuzz=500pt expression & \hfuzz=500pt for each element of matrix (variables: data, row, column, rows, columns, rowsBefore, columnsBefore)\\
\hline
\end{tabularx}


\subsection{ElementWiseOperation}
Given two matrices $\mathbf{A}$ and $\mathbf{B}$ this class computes $c_{ij} = f(a_{ij}, b_{ij})$,
where $f$ is an expression (for example \verb|data0*data1|).
For each element of the matrix the variables \verb|data0|, \verb|data1|, \verb|row|, \verb|column|
are set and the expression \config{element} is evaluated.


\keepXColumns
\begin{tabularx}{\textwidth}{N T A}
\hline
Name & Type & Annotation\\
\hline
\hfuzz=500pt\includegraphics[width=1em]{element-mustset-unbounded.pdf}~matrix1 & \hfuzz=500pt \hyperref[matrixGeneratorType]{matrixGenerator} & \hfuzz=500pt \\
\hfuzz=500pt\includegraphics[width=1em]{element-mustset-unbounded.pdf}~matrix2 & \hfuzz=500pt \hyperref[matrixGeneratorType]{matrixGenerator} & \hfuzz=500pt \\
\hfuzz=500pt\includegraphics[width=1em]{element.pdf}~expression & \hfuzz=500pt expression & \hfuzz=500pt for each element of matrix (variables: data0, data1, row, column, rows, columns, rowsBefore, columnsBefore)\\
\hline
\end{tabularx}


\subsection{Append}
Append matrix to the right (first row) or bottom (first column).


\keepXColumns
\begin{tabularx}{\textwidth}{N T A}
\hline
Name & Type & Annotation\\
\hline
\hfuzz=500pt\includegraphics[width=1em]{element-mustset-unbounded.pdf}~matrix & \hfuzz=500pt \hyperref[matrixGeneratorType]{matrixGenerator} & \hfuzz=500pt \\
\hfuzz=500pt\includegraphics[width=1em]{element-mustset.pdf}~side & \hfuzz=500pt choice & \hfuzz=500pt \\
\hfuzz=500pt\includegraphics[width=1em]{connector.pdf}\includegraphics[width=1em]{element-mustset.pdf}~right & \hfuzz=500pt  & \hfuzz=500pt \\
\hfuzz=500pt\includegraphics[width=1em]{connector.pdf}\includegraphics[width=1em]{element-mustset.pdf}~bottom & \hfuzz=500pt  & \hfuzz=500pt \\
\hfuzz=500pt\includegraphics[width=1em]{connector.pdf}\includegraphics[width=1em]{element-mustset.pdf}~diagonal & \hfuzz=500pt  & \hfuzz=500pt \\
\hline
\end{tabularx}


\subsection{Shift}
Shift start row and start column of a matrix.
In other words: zero lines and columns are inserted at the beginning of the matrix.


\keepXColumns
\begin{tabularx}{\textwidth}{N T A}
\hline
Name & Type & Annotation\\
\hline
\hfuzz=500pt\includegraphics[width=1em]{element-mustset-unbounded.pdf}~matrix & \hfuzz=500pt \hyperref[matrixGeneratorType]{matrixGenerator} & \hfuzz=500pt \\
\hfuzz=500pt\includegraphics[width=1em]{element.pdf}~startRow & \hfuzz=500pt expression & \hfuzz=500pt start row (variables: rowsBefore, columnsBefore, rows, columns)\\
\hfuzz=500pt\includegraphics[width=1em]{element.pdf}~startColumn & \hfuzz=500pt expression & \hfuzz=500pt start column (variables: rowsBefore, columnsBefore, rows, columns)\\
\hline
\end{tabularx}


\subsection{Slice}
Slice of a matrix.


\keepXColumns
\begin{tabularx}{\textwidth}{N T A}
\hline
Name & Type & Annotation\\
\hline
\hfuzz=500pt\includegraphics[width=1em]{element-mustset-unbounded.pdf}~matrix & \hfuzz=500pt \hyperref[matrixGeneratorType]{matrixGenerator} & \hfuzz=500pt \\
\hfuzz=500pt\includegraphics[width=1em]{element.pdf}~startRow & \hfuzz=500pt expression & \hfuzz=500pt start row of matrix (variables: rowsBefore, columnsBefore, rows, columns)\\
\hfuzz=500pt\includegraphics[width=1em]{element.pdf}~startColumn & \hfuzz=500pt expression & \hfuzz=500pt start column of matrix (variables: rowsBefore, columnsBefore, rows, columns)\\
\hfuzz=500pt\includegraphics[width=1em]{element.pdf}~rows & \hfuzz=500pt expression & \hfuzz=500pt 0: until end (variables: rowsBefore, columnsBefore, rows, columns)\\
\hfuzz=500pt\includegraphics[width=1em]{element.pdf}~columns & \hfuzz=500pt expression & \hfuzz=500pt 0: until end (variables: rowsBefore, columnsBefore, rows, columns)\\
\hline
\end{tabularx}


\subsection{Reshape}
Matrix reshaped columnwise to new row and columns.


\keepXColumns
\begin{tabularx}{\textwidth}{N T A}
\hline
Name & Type & Annotation\\
\hline
\hfuzz=500pt\includegraphics[width=1em]{element-mustset-unbounded.pdf}~matrix & \hfuzz=500pt \hyperref[matrixGeneratorType]{matrixGenerator} & \hfuzz=500pt \\
\hfuzz=500pt\includegraphics[width=1em]{element-mustset.pdf}~rows & \hfuzz=500pt expression & \hfuzz=500pt 0: auto-determine rows, (variables: rowsBefore, columnsBefore)\\
\hfuzz=500pt\includegraphics[width=1em]{element-mustset.pdf}~columns & \hfuzz=500pt expression & \hfuzz=500pt 0: auto-determine columns (variables: rowsBefore, columnsBefore)\\
\hline
\end{tabularx}


\subsection{Reorder}\label{matrixGeneratorType:reorder}
Reorder rows or columns of a matrix by an index vectors.
The index vector can be created with \program{ParameterSelection2IndexVector}.


\keepXColumns
\begin{tabularx}{\textwidth}{N T A}
\hline
Name & Type & Annotation\\
\hline
\hfuzz=500pt\includegraphics[width=1em]{element-mustset-unbounded.pdf}~matrix & \hfuzz=500pt \hyperref[matrixGeneratorType]{matrixGenerator} & \hfuzz=500pt \\
\hfuzz=500pt\includegraphics[width=1em]{element.pdf}~inputfileIndexVectorRow & \hfuzz=500pt filename & \hfuzz=500pt index in input matrix or -1 for new parameter.\\
\hfuzz=500pt\includegraphics[width=1em]{element.pdf}~inputfileIndexVectorColumn & \hfuzz=500pt filename & \hfuzz=500pt index in input matrix or -1 for new parameter.\\
\hline
\end{tabularx}


\subsection{Sort}
Sort matrix by \config{column} in ascending order by default or in \config{descending} order.


\keepXColumns
\begin{tabularx}{\textwidth}{N T A}
\hline
Name & Type & Annotation\\
\hline
\hfuzz=500pt\includegraphics[width=1em]{element-mustset-unbounded.pdf}~matrix & \hfuzz=500pt \hyperref[matrixGeneratorType]{matrixGenerator} & \hfuzz=500pt \\
\hfuzz=500pt\includegraphics[width=1em]{element-mustset-unbounded.pdf}~column & \hfuzz=500pt uint & \hfuzz=500pt sort by column, top = highest priority\\
\hfuzz=500pt\includegraphics[width=1em]{element.pdf}~descending & \hfuzz=500pt boolean & \hfuzz=500pt \\
\hline
\end{tabularx}


\subsection{Transpose}
Transposed of a matrix $\M A^T$.


\keepXColumns
\begin{tabularx}{\textwidth}{N T A}
\hline
Name & Type & Annotation\\
\hline
\hfuzz=500pt\includegraphics[width=1em]{element-mustset-unbounded.pdf}~matrix & \hfuzz=500pt \hyperref[matrixGeneratorType]{matrixGenerator} & \hfuzz=500pt \\
\hline
\end{tabularx}


\subsection{Multiplication}
Multiplication of matrices.


\keepXColumns
\begin{tabularx}{\textwidth}{N T A}
\hline
Name & Type & Annotation\\
\hline
\hfuzz=500pt\includegraphics[width=1em]{element-mustset-unbounded.pdf}~matrix1 & \hfuzz=500pt \hyperref[matrixGeneratorType]{matrixGenerator} & \hfuzz=500pt \\
\hfuzz=500pt\includegraphics[width=1em]{element-mustset-unbounded.pdf}~matrix2 & \hfuzz=500pt \hyperref[matrixGeneratorType]{matrixGenerator} & \hfuzz=500pt \\
\hfuzz=500pt\includegraphics[width=1em]{element.pdf}~factor & \hfuzz=500pt double & \hfuzz=500pt \\
\hline
\end{tabularx}


\subsection{Inverse}
Inverse of a matrix $\M A^{-1}$.


\keepXColumns
\begin{tabularx}{\textwidth}{N T A}
\hline
Name & Type & Annotation\\
\hline
\hfuzz=500pt\includegraphics[width=1em]{element-mustset-unbounded.pdf}~matrix & \hfuzz=500pt \hyperref[matrixGeneratorType]{matrixGenerator} & \hfuzz=500pt \\
\hfuzz=500pt\includegraphics[width=1em]{element.pdf}~pseudoInverse & \hfuzz=500pt boolean & \hfuzz=500pt compute pseudo inverse instead of regular one\\
\hline
\end{tabularx}


\subsection{Cholesky}
Upper triangular natrix of the cholesky decomposition of a symmetric matrix $\M A=\M W^T\M W$.


\keepXColumns
\begin{tabularx}{\textwidth}{N T A}
\hline
Name & Type & Annotation\\
\hline
\hfuzz=500pt\includegraphics[width=1em]{element-mustset-unbounded.pdf}~matrix & \hfuzz=500pt \hyperref[matrixGeneratorType]{matrixGenerator} & \hfuzz=500pt \\
\hline
\end{tabularx}


\subsection{RankKUpdate}
Symmetric matrix from rank k update: $\M A^T\M A$.


\keepXColumns
\begin{tabularx}{\textwidth}{N T A}
\hline
Name & Type & Annotation\\
\hline
\hfuzz=500pt\includegraphics[width=1em]{element-mustset-unbounded.pdf}~matrix & \hfuzz=500pt \hyperref[matrixGeneratorType]{matrixGenerator} & \hfuzz=500pt \\
\hfuzz=500pt\includegraphics[width=1em]{element.pdf}~factor & \hfuzz=500pt double & \hfuzz=500pt \\
\hline
\end{tabularx}


\subsection{EigenValues}
Computes the eigenvalues of a square matrix and gives a vector of eigenvalues for symmetric matrices
or a matrix with 2 columns with real and imaginary parts in general case.


\keepXColumns
\begin{tabularx}{\textwidth}{N T A}
\hline
Name & Type & Annotation\\
\hline
\hfuzz=500pt\includegraphics[width=1em]{element-mustset-unbounded.pdf}~matrix & \hfuzz=500pt \hyperref[matrixGeneratorType]{matrixGenerator} & \hfuzz=500pt \\
\hfuzz=500pt\includegraphics[width=1em]{element.pdf}~eigenVectors & \hfuzz=500pt boolean & \hfuzz=500pt return eigen vectors instead of eigen values\\
\hline
\end{tabularx}


\subsection{Diagonal}
Extract the diagonal or subdiagnoal ($n\times 1$ vector) of a matrix.
The zero \config{diagonal} means the main diagonal, a positive value the superdiagonal,
and a negative the subdiagonal.


\keepXColumns
\begin{tabularx}{\textwidth}{N T A}
\hline
Name & Type & Annotation\\
\hline
\hfuzz=500pt\includegraphics[width=1em]{element-mustset-unbounded.pdf}~matrix & \hfuzz=500pt \hyperref[matrixGeneratorType]{matrixGenerator} & \hfuzz=500pt \\
\hfuzz=500pt\includegraphics[width=1em]{element.pdf}~diagonal & \hfuzz=500pt int & \hfuzz=500pt zero: main diagonal, positive: superdiagonal, negative: subdiagonal\\
\hline
\end{tabularx}


\subsection{FromDiagonal}
Generate a matrix from a diagonal vector.


\keepXColumns
\begin{tabularx}{\textwidth}{N T A}
\hline
Name & Type & Annotation\\
\hline
\hfuzz=500pt\includegraphics[width=1em]{element-mustset-unbounded.pdf}~matrix & \hfuzz=500pt \hyperref[matrixGeneratorType]{matrixGenerator} & \hfuzz=500pt (nx1) or (1xn) diagonal vector\\
\hfuzz=500pt\includegraphics[width=1em]{element.pdf}~diagonal & \hfuzz=500pt int & \hfuzz=500pt zero: main diagonal, positive: superdiagonal, negative: subdiagonal\\
\hline
\end{tabularx}


\subsection{Set type}
Set type (matrix, matrixSymmetricUpper, matrixSymmetricLower, matrixTriangularUpper, matrixTriangularLower)
of a matrix. If the type is not matrix, the matrix must be quadratic. Symmetric matrices are filled symmetric
and for triangular matrix the other triangle is set to zero.


\keepXColumns
\begin{tabularx}{\textwidth}{N T A}
\hline
Name & Type & Annotation\\
\hline
\hfuzz=500pt\includegraphics[width=1em]{element-mustset-unbounded.pdf}~matrix & \hfuzz=500pt \hyperref[matrixGeneratorType]{matrixGenerator} & \hfuzz=500pt \\
\hfuzz=500pt\includegraphics[width=1em]{element-mustset.pdf}~type & \hfuzz=500pt choice & \hfuzz=500pt \\
\hfuzz=500pt\includegraphics[width=1em]{connector.pdf}\includegraphics[width=1em]{element-mustset.pdf}~matrix & \hfuzz=500pt  & \hfuzz=500pt \\
\hfuzz=500pt\includegraphics[width=1em]{connector.pdf}\includegraphics[width=1em]{element-mustset.pdf}~matrixSymmetricUpper & \hfuzz=500pt  & \hfuzz=500pt \\
\hfuzz=500pt\includegraphics[width=1em]{connector.pdf}\includegraphics[width=1em]{element-mustset.pdf}~matrixSymmetricLower & \hfuzz=500pt  & \hfuzz=500pt \\
\hfuzz=500pt\includegraphics[width=1em]{connector.pdf}\includegraphics[width=1em]{element-mustset.pdf}~matrixTriangularUpper & \hfuzz=500pt  & \hfuzz=500pt \\
\hfuzz=500pt\includegraphics[width=1em]{connector.pdf}\includegraphics[width=1em]{element-mustset.pdf}~matrixTriangularLower & \hfuzz=500pt  & \hfuzz=500pt \\
\hline
\end{tabularx}

\clearpage
%==================================

\section{MiscAccelerations}\label{miscAccelerationsType}
This class gives the non conservative forces acting on satellites.


\subsection{Relativistic effect}\label{miscAccelerationsType:relativisticEffect}
The relativistic effect to the acceleration of an artificial Earth satellite
according to IERS2010 conventions.


\keepXColumns
\begin{tabularx}{\textwidth}{N T A}
\hline
Name & Type & Annotation\\
\hline
\hfuzz=500pt\includegraphics[width=1em]{element.pdf}~beta & \hfuzz=500pt double & \hfuzz=500pt PPN (parameterized post-Newtonian) parameter\\
\hfuzz=500pt\includegraphics[width=1em]{element.pdf}~gamma & \hfuzz=500pt double & \hfuzz=500pt PPN (parameterized post-Newtonian) parameter\\
\hfuzz=500pt\includegraphics[width=1em]{element.pdf}~J & \hfuzz=500pt double & \hfuzz=500pt Earth’s angular momentum per unit mass [m**2/s]\\
\hfuzz=500pt\includegraphics[width=1em]{element.pdf}~GM & \hfuzz=500pt double & \hfuzz=500pt Geocentric gravitational constant\\
\hfuzz=500pt\includegraphics[width=1em]{element.pdf}~factor & \hfuzz=500pt double & \hfuzz=500pt the result is multplied by this factor\\
\hline
\end{tabularx}


\subsection{SolarRadiationPressure}\label{miscAccelerationsType:solarRadiationPressure}
Solar radiation pressure model.
Algorithmm for the solar radiation pressure is described in:

Lemoine, F. G., et al. ( 2013), High-degree gravity models from GRAIL primary mission data,
J. Geophys. Res. Planets, 118, 1676 1698 doi:10.1002/jgre.20118.


\keepXColumns
\begin{tabularx}{\textwidth}{N T A}
\hline
Name & Type & Annotation\\
\hline
\hfuzz=500pt\includegraphics[width=1em]{element.pdf}~solarflux & \hfuzz=500pt double & \hfuzz=500pt solar flux constant in 1 AU [W/m**2]\\
\hfuzz=500pt\includegraphics[width=1em]{element-mustset.pdf}~eclipse & \hfuzz=500pt \hyperref[eclipseType]{eclipse} & \hfuzz=500pt \\
\hfuzz=500pt\includegraphics[width=1em]{element.pdf}~factor & \hfuzz=500pt double & \hfuzz=500pt the result is multplied by this factor, set -1 to substract the field\\
\hline
\end{tabularx}


\subsection{Albedo}\label{miscAccelerationsType:albedo}
Acceleration caused by Earth's albedo.

The acceleration on the satellite caused by reflection and reradiation is computed using the
algorithm in:

Knocke, P. C., Ries, J. C., and Tapley, B. D. (1988). Earth radiation pressure effects on satellites.
Proceedings of the AIAA/AAS Astrodynamics Conference, 88-4292-CP, 577-87. DOI: 10.2514/6.
1988-4292.


\keepXColumns
\begin{tabularx}{\textwidth}{N T A}
\hline
Name & Type & Annotation\\
\hline
\hfuzz=500pt\includegraphics[width=1em]{element.pdf}~inputfileReflectivity & \hfuzz=500pt filename & \hfuzz=500pt \\
\hfuzz=500pt\includegraphics[width=1em]{element.pdf}~inputfileEmissivity & \hfuzz=500pt filename & \hfuzz=500pt \\
\hfuzz=500pt\includegraphics[width=1em]{element.pdf}~solarflux & \hfuzz=500pt double & \hfuzz=500pt solar flux constant in 1 AU [W/m**2]\\
\hfuzz=500pt\includegraphics[width=1em]{element.pdf}~factor & \hfuzz=500pt double & \hfuzz=500pt the result is multplied by this factor, set -1 to substract the field\\
\hline
\end{tabularx}


\subsection{AtmosphericDrag}\label{miscAccelerationsType:atmosphericDrag}
Atmospheric drag model.
Algorithm for the atmospheric drag modelling is based on the free molecule flow
theory by Sentman 1961. An analytical expression of this treatise is given in
Moe and Moe 2005.

Sentman L. (1961), Free molecule flow theory and its application to the determination
of aerodynamic forces, Technical report.

Moe K., Moe M. M. (2005), Gas-surface interactions and satellite drag coefficients,
Planetary and Space Science 53(8), 793-801, doi:10.1016/j.pss.2005.03.005.

Optional determination steps:
Turn temperature on or off.
In the first case, the model mentioned above is applied, which estimates variable drag
and lift coefficients - in the latter case a constant drag coefficient can be specified.

Turn wind on/off:
It enables the usage of the Horizontal Wind Model 2014 to add additional thermospheric
winds in the calculation process.


\keepXColumns
\begin{tabularx}{\textwidth}{N T A}
\hline
Name & Type & Annotation\\
\hline
\hfuzz=500pt\includegraphics[width=1em]{element-mustset.pdf}~thermosphere & \hfuzz=500pt \hyperref[thermosphereType]{thermosphere} & \hfuzz=500pt \\
\hfuzz=500pt\includegraphics[width=1em]{element.pdf}~earthRotation & \hfuzz=500pt double & \hfuzz=500pt [rad/s]\\
\hfuzz=500pt\includegraphics[width=1em]{element.pdf}~considerTemperature & \hfuzz=500pt boolean & \hfuzz=500pt compute drag and lift, otherwise simple drag coefficient is used\\
\hfuzz=500pt\includegraphics[width=1em]{element.pdf}~considerWind & \hfuzz=500pt boolean & \hfuzz=500pt \\
\hfuzz=500pt\includegraphics[width=1em]{element.pdf}~factor & \hfuzz=500pt double & \hfuzz=500pt the result is multplied by this factor\\
\hline
\end{tabularx}


\subsection{Antenna thrust}\label{miscAccelerationsType:antennaThrust}
The thrust (acceleration) in the opposite direction the antenna is facing
which is generated by satellite antenna broadcasts.
The thrust is defined in the satellite macro model.


\keepXColumns
\begin{tabularx}{\textwidth}{N T A}
\hline
Name & Type & Annotation\\
\hline
\hfuzz=500pt\includegraphics[width=1em]{element.pdf}~factor & \hfuzz=500pt double & \hfuzz=500pt the result is multplied by this factor\\
\hline
\end{tabularx}

\clearpage
%==================================

\section{NoiseGenerator}\label{noiseGeneratorType}
This class implements the generation of different types of noise.
It provides a generic interface that can be implemented by different
types of generators. The characteristics of the generated noise
is determined by the generators. See the appropriate documentation
for more information.


\subsection{White}
The noise is Gaussian with a standard deviation \config{sigma}.
The noise is computed via a pseudo random sequence with a start value given
by \config{initRandom}. The same value always yields the same sequence.
Be careful in \reference{parallel}{general.parallelization} mode
as all nodes generates the same pseudo random sequence.
If this value is set to zero a real random value is used as starting value.


\keepXColumns
\begin{tabularx}{\textwidth}{N T A}
\hline
Name & Type & Annotation\\
\hline
\hfuzz=500pt\includegraphics[width=1em]{element-mustset.pdf}~sigma & \hfuzz=500pt double & \hfuzz=500pt standard deviation\\
\hfuzz=500pt\includegraphics[width=1em]{element.pdf}~initRandom & \hfuzz=500pt uint & \hfuzz=500pt start value for pseudo random sequence, 0: real random\\
\hline
\end{tabularx}


\subsection{ExpressionPSD}
This generator creates noise defined by a one sided PSD.
The \config{psd} is an expression controlled by the variable 'freq'.
To determine the frequency \config{sampling} must be given.


\keepXColumns
\begin{tabularx}{\textwidth}{N T A}
\hline
Name & Type & Annotation\\
\hline
\hfuzz=500pt\includegraphics[width=1em]{element-mustset-unbounded.pdf}~noise & \hfuzz=500pt \hyperref[noiseGeneratorType]{noiseGenerator} & \hfuzz=500pt Basis noise\\
\hfuzz=500pt\includegraphics[width=1em]{element-mustset.pdf}~psd & \hfuzz=500pt expression & \hfuzz=500pt one sided PSD (variable: freq [Hz]) [unit\textasciicircum{}2/Hz]\\
\hfuzz=500pt\includegraphics[width=1em]{element-mustset.pdf}~sampling & \hfuzz=500pt double & \hfuzz=500pt to determine frequency [seconds]\\
\hline
\end{tabularx}


\subsection{Filter}
Generated noise \configClass{noise}{noiseGeneratorType} is
filtered by a \configClass{filter}{digitalFilterType}.


\keepXColumns
\begin{tabularx}{\textwidth}{N T A}
\hline
Name & Type & Annotation\\
\hline
\hfuzz=500pt\includegraphics[width=1em]{element-mustset-unbounded.pdf}~filter & \hfuzz=500pt \hyperref[digitalFilterType]{digitalFilter} & \hfuzz=500pt digital filter\\
\hfuzz=500pt\includegraphics[width=1em]{element-mustset-unbounded.pdf}~noise & \hfuzz=500pt \hyperref[noiseGeneratorType]{noiseGenerator} & \hfuzz=500pt Basis noise\\
\hfuzz=500pt\includegraphics[width=1em]{element.pdf}~warmupEpochCount & \hfuzz=500pt uint & \hfuzz=500pt number of additional epochs at before start and after end\\
\hfuzz=500pt\includegraphics[width=1em]{element.pdf}~overSamplingFactor & \hfuzz=500pt uint & \hfuzz=500pt noise with multiple higher sampling -\$>\$ filter -\$>\$ decimate\\
\hline
\end{tabularx}


\subsection{PowerLaw}
This generator creates noise that conforms to a power law relationship, where the power
of the noise at a frequency is proportional to $1/f^\alpha$, with a typically between -2 and 2.


\keepXColumns
\begin{tabularx}{\textwidth}{N T A}
\hline
Name & Type & Annotation\\
\hline
\hfuzz=500pt\includegraphics[width=1em]{element-mustset-unbounded.pdf}~noise & \hfuzz=500pt \hyperref[noiseGeneratorType]{noiseGenerator} & \hfuzz=500pt Basis noise\\
\hfuzz=500pt\includegraphics[width=1em]{element-mustset.pdf}~alpha & \hfuzz=500pt double & \hfuzz=500pt Exponent of the power law relationship 1/f\textasciicircum{}alpha\\
\hline
\end{tabularx}

\clearpage
%==================================

\section{NormalEquation}\label{normalEquationType}
This class provides a system of normal equations.
This total system is the weighted sum of individual normals.
\begin{equation}
 \M N_{total} =  \sum_{k=1} \frac{1}{\sigma_k^2}\M N_k
 \qquad\text{and}\qquad
\M n_{total} = \sum_{k=1} \frac{1}{\sigma_k^2} \M n_k.
\end{equation}
The normals do not need to have the same dimension. The dimension
of the total combined system is chosen to cover all individual systems.
For each normal a \config{startIndex} is required which indicates
the position of the first unknown of the individual normal within the
combined parameter vector.

The $\sigma_k$ of the relative weights are defined by \config{aprioriSigma}
in a first step. If an apriori solution \configFile{inputfileApproxSolution}{matrix} is
given or the normals are solved iteratively the weights are determined by means
of variance compoment estimation (VCE), see \program{NormalsSolverVCE}:
\begin{equation}
\sigma_k^2 =
\frac{\M e_k^T\M P\M e_k}
{n_k-\frac{1}{\sigma_k^2}\text{trace}\left(\M N_k\M N_{total}^{-1}\right)},
\end{equation}
where $n_k$ is the number of observations. The square sum of the residuals
is calculated by
\begin{equation}
\M e_k^T\M P\M e_k = \M x^T\M N_k\M x - 2\M n_k^T\M x + \M l_k^T\M P_k\M l_k.
\end{equation}
The system of normal equations can be solved with several right hand sides at once. But
only one right hand side, which can be selected with the index \config{rightHandSide}
(counting from zero), can be used to compute the variance factors.
The combined normal $\M N_{total}$ and the solution $\M x$ are taken from the previous
iteration step. In case of \configClass{DesignVCE}{normalEquationType:designVCE} the algorithm
is a little bit different as described below.


\subsection{Design}\label{normalEquationType:design}
This class acculumates normal equations from observation equations.
The class \configClass{observation}{observationType} computes
the linearized and decorrelated equation system for each arc $i$:
\begin{equation}
\M l_i  = \M A_i \M x + \M B_i \M y_i + \M e_i.
\end{equation}
The arc depending parameters~$\M y_i$ are eliminated and the system of normal
equations is acculumated according to
\begin{equation}
 \M N = \sum_{i=1}^m \M A_i^T  \M A_i
 \qquad\text{and}\qquad
\M n = \sum_{i=1}^m \M A_i^T \M l_i.
\end{equation}


\keepXColumns
\begin{tabularx}{\textwidth}{N T A}
\hline
Name & Type & Annotation\\
\hline
\hfuzz=500pt\includegraphics[width=1em]{element-mustset.pdf}~observation & \hfuzz=500pt \hyperref[observationType]{observation} & \hfuzz=500pt \\
\hfuzz=500pt\includegraphics[width=1em]{element.pdf}~aprioriSigma & \hfuzz=500pt double & \hfuzz=500pt \\
\hfuzz=500pt\includegraphics[width=1em]{element.pdf}~startIndex & \hfuzz=500pt uint & \hfuzz=500pt add this normals at index of total matrix (counting from 0)\\
\hfuzz=500pt\includegraphics[width=1em]{element.pdf}~inputfileArcList & \hfuzz=500pt filename & \hfuzz=500pt to accelerate computation\\
\hline
\end{tabularx}


\subsection{DesignVCE}\label{normalEquationType:designVCE}
This class acculumates normal equations from observation equations.
The class \configClass{observation}{observationType} computes
the linearized and decorrelated equation system for each arc $i$:
\begin{equation}
\M l_i  = \M A_i \M x + \M B_i \M y_i + \M e_i.
\end{equation}
The arc depending parameters~$\M y_i$ are eliminated and the system of normal
equations is acculumated according to
\begin{equation}
 \M N =  \sum_{i=1} \frac{1}{\sigma_i^2}\M A_i^T  \M A_i
 \qquad\text{and}\qquad
\M n = \sum_{i=1} \frac{1}{\sigma_i^2} \M A_i^T \M l_i.
\end{equation}
The variance $\sigma_i^2$ of each individual arc is determined by
\begin{equation}
\sigma_i^2 =
\frac{(\M l_i-\M A_i\M x)^T(\M l_i-\M A_i\M x)}
{n_i-\frac{1}{\sigma_i^2}\text{trace}\left(\M A_i^T  \M A_i\M N_{total}^{-1}\right)},
\end{equation}
where $n_i$ is the number of observations. If an apriori solution is not given at the first
iteration step a zero vector is assumed.


\keepXColumns
\begin{tabularx}{\textwidth}{N T A}
\hline
Name & Type & Annotation\\
\hline
\hfuzz=500pt\includegraphics[width=1em]{element-mustset.pdf}~observation & \hfuzz=500pt \hyperref[observationType]{observation} & \hfuzz=500pt \\
\hfuzz=500pt\includegraphics[width=1em]{element.pdf}~startIndex & \hfuzz=500pt uint & \hfuzz=500pt add this normals at index of total matrix (counting from 0)\\
\hfuzz=500pt\includegraphics[width=1em]{element.pdf}~inputfileArcList & \hfuzz=500pt filename & \hfuzz=500pt to accelerate computation\\
\hline
\end{tabularx}


\subsection{File}\label{normalEquationType:file}
Reads a system of normal equations from file \configFile{inputfileNormalEquation}{normalEquation}
as generated by e.g. \program{NormalsBuild}.


\keepXColumns
\begin{tabularx}{\textwidth}{N T A}
\hline
Name & Type & Annotation\\
\hline
\hfuzz=500pt\includegraphics[width=1em]{element-mustset.pdf}~inputfileNormalEquation & \hfuzz=500pt filename & \hfuzz=500pt \\
\hfuzz=500pt\includegraphics[width=1em]{element.pdf}~aprioriSigma & \hfuzz=500pt double & \hfuzz=500pt \\
\hfuzz=500pt\includegraphics[width=1em]{element.pdf}~startIndex & \hfuzz=500pt uint & \hfuzz=500pt add this normals at index of total matrix (counting from 0)\\
\hline
\end{tabularx}


\subsection{Regularization}\label{normalEquationType:regularization}
Set up a system of normal equations
\begin{equation}
\M N = \M R
\qquad\text{and}\qquad
\M n = \M R \M b,
\end{equation}
where $\M R$ is a diagonal matrix whose elements are given as a vector by
\configFile{inputfileDiagonalMatrix}{matrix} and $\M b$ is the right hand side towards which will
be regularized. It can be given by \configFile{inputfileBiasVector}{matrix}.
The diagonal matrix can be generated with \program{NormalsRegularizationBorders},
\program{NormalsRegularizationSphericalHarmonics}, or \program{MatrixCalculate}.
If $\M R$ is not given a unit matrix is assumed.
The right hand side $\M b$ may be generated with \program{Gravityfield2SphericalHarmonicsVector}.
If $\M b$ is not given a zero vector is assumed.


\keepXColumns
\begin{tabularx}{\textwidth}{N T A}
\hline
Name & Type & Annotation\\
\hline
\hfuzz=500pt\includegraphics[width=1em]{element.pdf}~inputfileDiagonalMatrix & \hfuzz=500pt filename & \hfuzz=500pt Vector with the diagonal elements of the weight matrix\\
\hfuzz=500pt\includegraphics[width=1em]{element.pdf}~inputfileBias & \hfuzz=500pt filename & \hfuzz=500pt Matrix with right hand sides\\
\hfuzz=500pt\includegraphics[width=1em]{element.pdf}~aprioriSigma & \hfuzz=500pt double & \hfuzz=500pt \\
\hfuzz=500pt\includegraphics[width=1em]{element.pdf}~startIndex & \hfuzz=500pt uint & \hfuzz=500pt regularization of parameters starts at this index (counting from 0)\\
\hline
\end{tabularx}


\subsection{RegularizationGeneralized}

Generalized regularization which is represented by the observation equation
\begin{equation}
\mathbf{x}_0 = \mathbf{I} \mathbf{x} + \mathbf{v}, \mathbf{v} \sim \mathcal{N}(0, \sum_k \sigma^2_k \mathbf{V}_k).
\end{equation}

There are no requirements for partial covariance matrices $\mathbf{V}_k$ except for them being symmetric.
The accumulated covariance matrix $\sum_k \sigma^2_k \mathbf{V}_k$ must be positive definite however.
The variance components $\sigma^2_k$ are estimated during the adjustment process and are assumed to be positive.
All \configFile{inputfilePartialCovarianceMatrix}{matrix} must be of same size
and must match the dimension of \configFile{inputfileBiasMatrix}{matrix}
(if provided, otherwise a zero vector of appropriate dimensions is created).

The parameter \config{aprioriSigma} determines the initial variance factor for the partial covariance matrices. Either one $\sigma_0$ can be
supplied or one for each $\mathbf{V}_k$.

The regularization matrix can be applied to a subset of parameters by adjusting \config{startIndex}.


\keepXColumns
\begin{tabularx}{\textwidth}{N T A}
\hline
Name & Type & Annotation\\
\hline
\hfuzz=500pt\includegraphics[width=1em]{element-mustset-unbounded.pdf}~inputfilePartialCovarianceMatrix & \hfuzz=500pt filename & \hfuzz=500pt symmetric matrix (sum of all matrices must be positive definite)\\
\hfuzz=500pt\includegraphics[width=1em]{element.pdf}~inputfileBiasMatrix & \hfuzz=500pt filename & \hfuzz=500pt bias vector (default: zero vector)\\
\hfuzz=500pt\includegraphics[width=1em]{element-mustset-unbounded.pdf}~aprioriSigma & \hfuzz=500pt double & \hfuzz=500pt apriori sigmas for initial iteration (default: 1.0)\\
\hfuzz=500pt\includegraphics[width=1em]{element.pdf}~startIndex & \hfuzz=500pt uint & \hfuzz=500pt regularization of parameters starts at this index (counting from 0)\\
\hline
\end{tabularx}

\clearpage
%==================================

\section{Observation}\label{observationType}
This class set up the oberservation equations in linearized Gauss-Makoff model
\begin{equation}\label{gmm}
\M l  = \M A \M x + \M e\qquad\text{and}\qquad\mathcal{C}(\M e) = \sigma^2\M P^{-1}.
\end{equation}
The observations are divided into short data blocks which can computed independently
and so easily can be parallized. Usually this data blocks are short arcs of a
satellites orbit. In most cases the unknown parameter vector contains coefficients
of a gravity field parametrization given by \configClass{parametrizationGravity}{parametrizationGravityType}.
Additional parameters like instrument calibrations parameters are appended at the
end of the vector~$\M x$.
It is possible to give several observation vectors in one model.

The observations within each arc are decorrelated in the following way:
In a first step a cholesky decomposition of the covariance matrix is performed
\begin{equation}
\M P^{-1} = \M W^T\M W,
\end{equation}
where $\M W$ is an upper regular triangular matrix.
In a second step the transformation
\begin{equation}\label{dekorrelierung}
\bar{\M A} = \M W^{-T}\M A\qquad\text{and}\qquad \bar{\M l} = \M W^{-T}\M l
\end{equation}
gives an estimation from decorrelated observations with equal variance
\begin{equation}\label{normal.GMM}
\bar{\M l} = \bar{\M A} \M x + \bar{\M e}
\qquad\text{and}\qquad
\mathcal{C}(\bar{\M e})= \sigma^2 \M I.
\end{equation}
Usually the arc depending parameters are eliminated in the next step.


\subsection{PodVariational}\label{observationType:podVariational}
The observation equations for precise orbit data (POD) are formulated as variational equations.
It is based on \file{inputfileVariational}{variationalEquation} calculated with \program{PreprocessingVariationalEquation}.
Necessary integrations are performed by integrating a moving interpolation polynomial of degree \config{integrationDegree}.

The kinematic positions as pseudo observations are taken from
\config{rightHandSide} and should not given equally spaced in time. The observation
equations are interpolated to these times by a moving polynomial of degree \config{interpolationDegree}.

The accuracy or the full covariance matrix of the precise orbit data is provided in
\configClass{covariancePod}{covariancePodType} and can be estimated with \program{PreprocessingPod}.

\config{accelerateComputation}: In the event that the sampling of the kinematic orbit is much higher than the sampling
of the variational equations (e.g. 1 second vs. 5 seconds) the accumulation of the observation equations
can be accelerated by transforming the observation equations
\begin{equation}
  \M l = \M J \M A \M x + \M e,
\end{equation}
where $\M J$ describes the interpolation of the sampling of the variational design matrix~$\M A$
to the sampling of the observations $\M l$ with more rows than columns. The QR decomposition
\begin{equation}
  \M J = \begin{pmatrix} \M Q_1 & \M Q_2 \end{pmatrix}
         \begin{pmatrix} \M R \\ \M 0 \end{pmatrix}.
\end{equation}
can be used to transform the observation equations
\begin{equation}
  \begin{pmatrix} \M Q_1^T \M l \\ \M Q_2^T \M l \end{pmatrix} =
  \begin{pmatrix} \M Q_1^T \M R \\ \M 0 \end{pmatrix} \M A \M x +
  \begin{pmatrix} \M Q_1^T \M e \\ \M Q_2^T \M e \end{pmatrix}.
\end{equation}
As the zero lines should not be considered the computational time for the accumulation is reduced.
This option is not meaningful for evaluating the residuals such in \program{PreprocessingPod}.


\keepXColumns
\begin{tabularx}{\textwidth}{N T A}
\hline
Name & Type & Annotation\\
\hline
\hfuzz=500pt\includegraphics[width=1em]{element-mustset.pdf}~rightHandSide & \hfuzz=500pt sequence & \hfuzz=500pt input for observation vectors\\
\hfuzz=500pt\includegraphics[width=1em]{connector.pdf}\includegraphics[width=1em]{element-mustset.pdf}~inputfileOrbit & \hfuzz=500pt filename & \hfuzz=500pt kinematic positions as observations\\
\hfuzz=500pt\includegraphics[width=1em]{element-mustset.pdf}~inputfileVariational & \hfuzz=500pt filename & \hfuzz=500pt approximate position and integrated state matrix\\
\hfuzz=500pt\includegraphics[width=1em]{element.pdf}~ephemerides & \hfuzz=500pt \hyperref[ephemeridesType]{ephemerides} & \hfuzz=500pt \\
\hfuzz=500pt\includegraphics[width=1em]{element-unbounded.pdf}~parametrizationGravity & \hfuzz=500pt \hyperref[parametrizationGravityType]{parametrizationGravity} & \hfuzz=500pt gravity field parametrization\\
\hfuzz=500pt\includegraphics[width=1em]{element-unbounded.pdf}~parametrizationAcceleration & \hfuzz=500pt \hyperref[parametrizationAccelerationType]{parametrizationAcceleration} & \hfuzz=500pt orbit force parameters\\
\hfuzz=500pt\includegraphics[width=1em]{element.pdf}~integrationDegree & \hfuzz=500pt uint & \hfuzz=500pt integration of forces by polynomial approximation of degree n\\
\hfuzz=500pt\includegraphics[width=1em]{element.pdf}~interpolationDegree & \hfuzz=500pt uint & \hfuzz=500pt orbit interpolation by polynomial approximation of degree n\\
\hfuzz=500pt\includegraphics[width=1em]{element.pdf}~accelerateComputation & \hfuzz=500pt boolean & \hfuzz=500pt acceleration of computation by transforming the observations\\
\hfuzz=500pt\includegraphics[width=1em]{element.pdf}~covariancePod & \hfuzz=500pt \hyperref[covariancePodType]{covariancePod} & \hfuzz=500pt covariance matrix of kinematic orbits\\
\hline
\end{tabularx}


\subsection{PodIntegral}\label{observationType:podIntegral}
The observation equations for precise orbit data (POD) of short arcs are given by
\begin{equation}
  {\M r}_\epsilon(\tau) = {\M r}_A(1-\tau) + {\M r}_B\tau - T^2\int_0^1 K(\tau,\tau')
  \left(\M f_0(\tau')+\nabla V(\tau')\right)\,d\tau'
\end{equation}
with the integral kernel
\begin{equation}
  K(\tau,\tau') = \begin{cases} \tau'(1-\tau) & \text{for }\tau'\le\tau \\
  \tau(1-\tau') & \text{for }\tau'>\tau \end{cases},
\end{equation}
and the normalized time
\begin{equation}
  \tau = \frac{t-t_A}{T}\qquad\text{with}\qquad T=t_B-t_A.
\end{equation}
The kinematic positions~${\M r}_\epsilon(\tau)$ as pseudo observations are taken from
\configClass{rightHandSide}{podRightSideType}. From these positions the influence of the reference forces $\M f_0(\tau)$
is subtracted which are computed with the background models in \configClass{rightHandSide}{podRightSideType}.
The integral is solved by the integration of a moving interpolation polynomial of degree \config{integrationDegree}.
The boundary values ${\M r}_A$ and ${\M r}_B$ (satellite's state vector) are estimated per arc
and are usally directly eliminated if \config{keepSatelliteStates} is not set.

The unknown gravity field $\nabla V(\M r, t)$ parametrized by \configClass{parametrizationGravity}{parametrizationGravityType}
is not evaluated at the observed positions but at the orbit given by \configFile{inputfileOrbit}{instrument}.
The same is true for the reference forces. The linearized effect of the gravity field change by the position
adjustment is taken into account by \config{gradientfield}. This may be a low order field up to a
spherical harmonics degree of $n=2$ or $n=3$.

The \configFile{inputfileOrbit}{instrument}, \configFile{inputfileStarCamera}{instrument}, and \configFile{inputfileAccelerometer}{instrument}
must be synchronous and must be given with a constant sampling and without any gaps in each short arc
(see \program{InstrumentSynchronize}).
The kinematic positions~${\M r}_\epsilon(\tau)$ should not given equally spaced in time
but must be divided into the same arcs as the other instrument data.
The observation equations are interpolated to this time by a polynomial interpolation
with degree \config{interpolationDegree}.

The accuracy or the full covariance matrix of the precise orbit data is provided in
\configClass{covariancePod}{covariancePodType} and can be estimated with \program{PreprocessingPod}.

For \config{accelerateComputation} see \configClass{observation:podVariational}{observationType:podVariational}.


\keepXColumns
\begin{tabularx}{\textwidth}{N T A}
\hline
Name & Type & Annotation\\
\hline
\hfuzz=500pt\includegraphics[width=1em]{element.pdf}~inputfileSatelliteModel & \hfuzz=500pt filename & \hfuzz=500pt satellite macro model\\
\hfuzz=500pt\includegraphics[width=1em]{element-mustset-unbounded.pdf}~rightHandSide & \hfuzz=500pt \hyperref[podRightSideType]{podRightSide} & \hfuzz=500pt input for the reduced observation vector\\
\hfuzz=500pt\includegraphics[width=1em]{element-mustset.pdf}~inputfileOrbit & \hfuzz=500pt filename & \hfuzz=500pt used to evaluate the observation equations, not used as observations\\
\hfuzz=500pt\includegraphics[width=1em]{element-mustset.pdf}~inputfileStarCamera & \hfuzz=500pt filename & \hfuzz=500pt \\
\hfuzz=500pt\includegraphics[width=1em]{element-mustset.pdf}~earthRotation & \hfuzz=500pt \hyperref[earthRotationType]{earthRotation} & \hfuzz=500pt \\
\hfuzz=500pt\includegraphics[width=1em]{element.pdf}~ephemerides & \hfuzz=500pt \hyperref[ephemeridesType]{ephemerides} & \hfuzz=500pt \\
\hfuzz=500pt\includegraphics[width=1em]{element-unbounded.pdf}~gradientfield & \hfuzz=500pt \hyperref[gravityfieldType]{gravityfield} & \hfuzz=500pt low order field to estimate the change of the gravity by position adjustement\\
\hfuzz=500pt\includegraphics[width=1em]{element-unbounded.pdf}~parametrizationGravity & \hfuzz=500pt \hyperref[parametrizationGravityType]{parametrizationGravity} & \hfuzz=500pt gravity field parametrization\\
\hfuzz=500pt\includegraphics[width=1em]{element-unbounded.pdf}~parametrizationAcceleration & \hfuzz=500pt \hyperref[parametrizationAccelerationType]{parametrizationAcceleration} & \hfuzz=500pt orbit force parameters\\
\hfuzz=500pt\includegraphics[width=1em]{element.pdf}~keepSatelliteStates & \hfuzz=500pt boolean & \hfuzz=500pt set boundary values of each arc global\\
\hfuzz=500pt\includegraphics[width=1em]{element.pdf}~integrationDegree & \hfuzz=500pt uint & \hfuzz=500pt integration of forces by polynomial approximation of degree n\\
\hfuzz=500pt\includegraphics[width=1em]{element.pdf}~interpolationDegree & \hfuzz=500pt uint & \hfuzz=500pt orbit interpolation by polynomial approximation of degree n\\
\hfuzz=500pt\includegraphics[width=1em]{element.pdf}~accelerateComputation & \hfuzz=500pt boolean & \hfuzz=500pt acceleration of computation by transforming the observations\\
\hfuzz=500pt\includegraphics[width=1em]{element.pdf}~covariancePod & \hfuzz=500pt \hyperref[covariancePodType]{covariancePod} & \hfuzz=500pt covariance matrix of kinematic orbits\\
\hline
\end{tabularx}


\subsection{PodAcceleration}\label{observationType:podAcceleration}
The observation equations for precise orbit data (POD) are given by
\begin{equation}
\ddot{\M r}(t) - \M g_0(t) = \nabla V(\M r, t),
\end{equation}
where the accelerations of the satellite $\ddot{\M r}(t)$ are derived from the kinematic positions
in \configClass{rightHandSide}{podRightSideType}. The orbit differentation is performed by a moving
polynomial interpolation or approximation with degree \config{interpolationDegree}
and number of used epochs \config{numberOfEpochs}. The reference forces $\M g_0(t)$ are computed
with the background models in \configClass{rightHandSide}{podRightSideType}.

All instrument data \configFile{inputfileOrbit}{instrument}, \configFile{inputfileStarCamera}{instrument},
and \configFile{inputfileAccelerometer}{instrument} must be synchronous and be given
with a constant sampling without any gaps in each short arc (see \program{InstrumentSynchronize}).

The unknown gravity field $\nabla V(\M r, t)$ parametrized by \configClass{parametrizationGravity}{parametrizationGravityType}
is not evaluated at the observed positions but at the orbit given by \configFile{inputfileOrbit}{instrument}.
The same is true for the reference forces. This orbit may be a more accurate dynamical orbit but
in most cases the kinematic orbit provides good results.

The accuracy or the full covariance matrix of the precise orbit data is provided in
\configClass{covariancePod}{covariancePodType} and can be estimated with \program{PreprocessingPod}.


\keepXColumns
\begin{tabularx}{\textwidth}{N T A}
\hline
Name & Type & Annotation\\
\hline
\hfuzz=500pt\includegraphics[width=1em]{element.pdf}~inputfileSatelliteModel & \hfuzz=500pt filename & \hfuzz=500pt satellite macro model\\
\hfuzz=500pt\includegraphics[width=1em]{element-mustset-unbounded.pdf}~rightHandSide & \hfuzz=500pt \hyperref[podRightSideType]{podRightSide} & \hfuzz=500pt input for the reduced observation vector\\
\hfuzz=500pt\includegraphics[width=1em]{element-mustset.pdf}~inputfileOrbit & \hfuzz=500pt filename & \hfuzz=500pt used to evaluate the observation equations, not used as observations\\
\hfuzz=500pt\includegraphics[width=1em]{element-mustset.pdf}~inputfileStarCamera & \hfuzz=500pt filename & \hfuzz=500pt \\
\hfuzz=500pt\includegraphics[width=1em]{element-mustset.pdf}~earthRotation & \hfuzz=500pt \hyperref[earthRotationType]{earthRotation} & \hfuzz=500pt \\
\hfuzz=500pt\includegraphics[width=1em]{element.pdf}~ephemerides & \hfuzz=500pt \hyperref[ephemeridesType]{ephemerides} & \hfuzz=500pt \\
\hfuzz=500pt\includegraphics[width=1em]{element-mustset-unbounded.pdf}~parametrizationGravity & \hfuzz=500pt \hyperref[parametrizationGravityType]{parametrizationGravity} & \hfuzz=500pt gravity field parametrization\\
\hfuzz=500pt\includegraphics[width=1em]{element-unbounded.pdf}~parametrizationAcceleration & \hfuzz=500pt \hyperref[parametrizationAccelerationType]{parametrizationAcceleration} & \hfuzz=500pt orbit force parameters\\
\hfuzz=500pt\includegraphics[width=1em]{element.pdf}~interpolationDegree & \hfuzz=500pt uint & \hfuzz=500pt orbit differentation  by polynomial approximation of degree n\\
\hfuzz=500pt\includegraphics[width=1em]{element.pdf}~numberOfEpochs & \hfuzz=500pt uint & \hfuzz=500pt number of used Epochs for polynom computation\\
\hfuzz=500pt\includegraphics[width=1em]{element.pdf}~covariancePod & \hfuzz=500pt \hyperref[covariancePodType]{covariancePod} & \hfuzz=500pt covariance matrix of kinematic orbits\\
\hline
\end{tabularx}


\subsection{PodEnergy}\label{observationType:podEnergy}
The observation equations for precise orbit data (POD) are given by
\begin{equation}
  \frac{1}{2}\dot{\M r}^2
  -\dot{\M r} \cdot (\M\Omega\times\M r)
  +\int_{t_0}^t(\dot{\M\Omega}\times\M r)\cdot \dot{\M r}\,dt
  - \int_{t_0}^t \M g_0 \cdot\dot{\M r}'\,dt
  = V + E.
\end{equation}
where the velocities of the satellite $\ddot{\M r}(t)$ are derived from
the kinematic positions in \configClass{rightHandSide}{podRightSideType} and the Earth's rotation vector~$\M\Omega(t)$ is modeled
within \configClass{earthRotation}{earthRotationType}. The orbit differentation is
performed by a polynomial interpolation with degree \config{interpolationDegree}.
The integrals are solved a polynomial interpolation with degree \config{integrationDegree}.
The reference forces $\M g_0(t)$ are computed with the background models in \configClass{rightHandSide}{podRightSideType}.

All instrument data \configFile{inputfileOrbit}{instrument}, \configFile{inputfileStarCamera}{instrument}, and \configFile{inputfileAccelerometer}{instrument}
must be synchronous and be given with a constant sampling without any gaps in each short arc
(see \program{InstrumentSynchronize}).

The unknown gravity potential $V(\M r)$ parametrized by \configClass{parametrizationGravity}{parametrizationGravityType}
is not evaluated at the observed positions but at the orbit given by \configFile{inputfileOrbit}{instrument}.
The same is true for the reference forces. This orbit may be a more accurate dynamical orbit but
in most cases the kinematic orbit provides good results.

An unknown energy bias~$E$ per arc is parametrized by \configClass{parametrizationBias}{parametrizationTemporalType}
and should be a constant in theory but temporal changes might help to absorb other unmodelled effects.

The accuracy or the full covariance matrix of the precise orbit data is provided in
\configClass{covariancePod}{covariancePodType} and can be estimated with \program{PreprocessingPod}.


\keepXColumns
\begin{tabularx}{\textwidth}{N T A}
\hline
Name & Type & Annotation\\
\hline
\hfuzz=500pt\includegraphics[width=1em]{element.pdf}~inputfileSatelliteModel & \hfuzz=500pt filename & \hfuzz=500pt satellite macro model\\
\hfuzz=500pt\includegraphics[width=1em]{element-mustset-unbounded.pdf}~rightHandSide & \hfuzz=500pt \hyperref[podRightSideType]{podRightSide} & \hfuzz=500pt input for the reduced observation vector\\
\hfuzz=500pt\includegraphics[width=1em]{element-mustset.pdf}~inputfileOrbit & \hfuzz=500pt filename & \hfuzz=500pt used to evaluate the observation equations, not used as observations\\
\hfuzz=500pt\includegraphics[width=1em]{element-mustset.pdf}~inputfileStarCamera & \hfuzz=500pt filename & \hfuzz=500pt \\
\hfuzz=500pt\includegraphics[width=1em]{element-mustset.pdf}~earthRotation & \hfuzz=500pt \hyperref[earthRotationType]{earthRotation} & \hfuzz=500pt \\
\hfuzz=500pt\includegraphics[width=1em]{element.pdf}~ephemerides & \hfuzz=500pt \hyperref[ephemeridesType]{ephemerides} & \hfuzz=500pt \\
\hfuzz=500pt\includegraphics[width=1em]{element-mustset-unbounded.pdf}~parametrizationGravity & \hfuzz=500pt \hyperref[parametrizationGravityType]{parametrizationGravity} & \hfuzz=500pt gravity field parametrization (potential)\\
\hfuzz=500pt\includegraphics[width=1em]{element-mustset-unbounded.pdf}~parametrizationBias & \hfuzz=500pt \hyperref[parametrizationTemporalType]{parametrizationTemporal} & \hfuzz=500pt unknown total energy per arc\\
\hfuzz=500pt\includegraphics[width=1em]{element.pdf}~interpolationDegree & \hfuzz=500pt uint & \hfuzz=500pt orbit differentation  by polynomial approximation of degree n\\
\hfuzz=500pt\includegraphics[width=1em]{element.pdf}~integrationDegree & \hfuzz=500pt uint & \hfuzz=500pt integration of forces by polynomial approximation of degree n\\
\hfuzz=500pt\includegraphics[width=1em]{element.pdf}~covariancePod & \hfuzz=500pt \hyperref[covariancePodType]{covariancePod} & \hfuzz=500pt covariance matrix of kinematic orbits\\
\hline
\end{tabularx}


\subsection{SstVariational}\label{observationType:sstVariational}
Like \configClass{observation:podVariational}{observationType:podVariational} (see there for details)
but with two satellites and additional satellite-to-satellite (SST) observations.

If multiple \configFile{inputfileSatelliteTracking}{instrument} are given
all data are add together. So corrections in extra files like the light time correction
can easily be added. Empirical parameters for the SST observations can be setup with
\configClass{parametrizationSst}{parametrizationSatelliteTrackingType}.
The accuracy or the full covariance matrix of SST is provided in
\configClass{covarianceSst}{covarianceSstType}.


\keepXColumns
\begin{tabularx}{\textwidth}{N T A}
\hline
Name & Type & Annotation\\
\hline
\hfuzz=500pt\includegraphics[width=1em]{element-mustset.pdf}~rightHandSide & \hfuzz=500pt sequence & \hfuzz=500pt input for observation vectors\\
\hfuzz=500pt\includegraphics[width=1em]{connector.pdf}\includegraphics[width=1em]{element-mustset-unbounded.pdf}~inputfileSatelliteTracking & \hfuzz=500pt filename & \hfuzz=500pt ranging observations and corrections\\
\hfuzz=500pt\includegraphics[width=1em]{connector.pdf}\includegraphics[width=1em]{element.pdf}~inputfileOrbit1 & \hfuzz=500pt filename & \hfuzz=500pt kinematic positions of satellite A as observations\\
\hfuzz=500pt\includegraphics[width=1em]{connector.pdf}\includegraphics[width=1em]{element.pdf}~inputfileOrbit2 & \hfuzz=500pt filename & \hfuzz=500pt kinematic positions of satellite B as observations\\
\hfuzz=500pt\includegraphics[width=1em]{element-mustset.pdf}~sstType & \hfuzz=500pt choice & \hfuzz=500pt \\
\hfuzz=500pt\includegraphics[width=1em]{connector.pdf}\includegraphics[width=1em]{element-mustset.pdf}~range & \hfuzz=500pt  & \hfuzz=500pt \\
\hfuzz=500pt\includegraphics[width=1em]{connector.pdf}\includegraphics[width=1em]{element-mustset.pdf}~rangeRate & \hfuzz=500pt  & \hfuzz=500pt \\
\hfuzz=500pt\includegraphics[width=1em]{connector.pdf}\includegraphics[width=1em]{element-mustset.pdf}~none & \hfuzz=500pt  & \hfuzz=500pt \\
\hfuzz=500pt\includegraphics[width=1em]{element-mustset.pdf}~inputfileVariational1 & \hfuzz=500pt filename & \hfuzz=500pt approximate position and integrated state matrix\\
\hfuzz=500pt\includegraphics[width=1em]{element-mustset.pdf}~inputfileVariational2 & \hfuzz=500pt filename & \hfuzz=500pt approximate position and integrated state matrix\\
\hfuzz=500pt\includegraphics[width=1em]{element.pdf}~ephemerides & \hfuzz=500pt \hyperref[ephemeridesType]{ephemerides} & \hfuzz=500pt \\
\hfuzz=500pt\includegraphics[width=1em]{element-unbounded.pdf}~parametrizationGravity & \hfuzz=500pt \hyperref[parametrizationGravityType]{parametrizationGravity} & \hfuzz=500pt gravity field parametrization\\
\hfuzz=500pt\includegraphics[width=1em]{element-unbounded.pdf}~parametrizationAcceleration1 & \hfuzz=500pt \hyperref[parametrizationAccelerationType]{parametrizationAcceleration} & \hfuzz=500pt orbit1 force parameters\\
\hfuzz=500pt\includegraphics[width=1em]{element-unbounded.pdf}~parametrizationAcceleration2 & \hfuzz=500pt \hyperref[parametrizationAccelerationType]{parametrizationAcceleration} & \hfuzz=500pt orbit2 force parameters\\
\hfuzz=500pt\includegraphics[width=1em]{element-unbounded.pdf}~parametrizationSst & \hfuzz=500pt \hyperref[parametrizationSatelliteTrackingType]{parametrizationSatelliteTracking} & \hfuzz=500pt satellite tracking parameter\\
\hfuzz=500pt\includegraphics[width=1em]{element.pdf}~integrationDegree & \hfuzz=500pt uint & \hfuzz=500pt integration of forces by polynomial approximation of degree n\\
\hfuzz=500pt\includegraphics[width=1em]{element.pdf}~interpolationDegree & \hfuzz=500pt uint & \hfuzz=500pt orbit interpolation by polynomial approximation of degree n\\
\hfuzz=500pt\includegraphics[width=1em]{element-mustset.pdf}~covarianceSst & \hfuzz=500pt \hyperref[covarianceSstType]{covarianceSst} & \hfuzz=500pt covariance matrix of satellite to satellite tracking observations\\
\hfuzz=500pt\includegraphics[width=1em]{element-mustset.pdf}~covariancePod1 & \hfuzz=500pt \hyperref[covariancePodType]{covariancePod} & \hfuzz=500pt covariance matrix of kinematic orbits (satellite 1)\\
\hfuzz=500pt\includegraphics[width=1em]{element-mustset.pdf}~covariancePod2 & \hfuzz=500pt \hyperref[covariancePodType]{covariancePod} & \hfuzz=500pt covariance matrix of kinematic orbits (satellite 2)\\
\hline
\end{tabularx}


\subsection{SstIntegral}\label{observationType:sstIntegral}
Like \configClass{observation:podIntegral}{observationType:podIntegral} (see there for details)
but with two satellites and additional satellite-to-satellite (SST) observations.

If multiple \configFile{inputfileSatelliteTracking}{instrument} are given
all data are add together. So corrections in extra files like the light time correction
can easily be added. Empirical parameters for the SST observations can be setup with
\configClass{parametrizationSst}{parametrizationSatelliteTrackingType}.
The accuracy or the full covariance matrix of SST is provided in
\configClass{covarianceSst}{covarianceSstType}.


\keepXColumns
\begin{tabularx}{\textwidth}{N T A}
\hline
Name & Type & Annotation\\
\hline
\hfuzz=500pt\includegraphics[width=1em]{element.pdf}~inputfileSatelliteModel1 & \hfuzz=500pt filename & \hfuzz=500pt satellite macro model\\
\hfuzz=500pt\includegraphics[width=1em]{element.pdf}~inputfileSatelliteModel2 & \hfuzz=500pt filename & \hfuzz=500pt satellite macro model\\
\hfuzz=500pt\includegraphics[width=1em]{element-mustset-unbounded.pdf}~rightHandSide & \hfuzz=500pt \hyperref[sstRightSideType]{sstRightSide} & \hfuzz=500pt input for the reduced observation vector\\
\hfuzz=500pt\includegraphics[width=1em]{element-mustset.pdf}~sstType & \hfuzz=500pt choice & \hfuzz=500pt \\
\hfuzz=500pt\includegraphics[width=1em]{connector.pdf}\includegraphics[width=1em]{element-mustset.pdf}~range & \hfuzz=500pt  & \hfuzz=500pt \\
\hfuzz=500pt\includegraphics[width=1em]{connector.pdf}\includegraphics[width=1em]{element-mustset.pdf}~rangeRate & \hfuzz=500pt  & \hfuzz=500pt \\
\hfuzz=500pt\includegraphics[width=1em]{connector.pdf}\includegraphics[width=1em]{element-mustset.pdf}~rangeAcceleration & \hfuzz=500pt  & \hfuzz=500pt \\
\hfuzz=500pt\includegraphics[width=1em]{connector.pdf}\includegraphics[width=1em]{element-mustset.pdf}~none & \hfuzz=500pt  & \hfuzz=500pt \\
\hfuzz=500pt\includegraphics[width=1em]{element-mustset.pdf}~inputfileOrbit1 & \hfuzz=500pt filename & \hfuzz=500pt used to evaluate the observation equations, not used as observations\\
\hfuzz=500pt\includegraphics[width=1em]{element-mustset.pdf}~inputfileOrbit2 & \hfuzz=500pt filename & \hfuzz=500pt used to evaluate the observation equations, not used as observations\\
\hfuzz=500pt\includegraphics[width=1em]{element-mustset.pdf}~inputfileStarCamera1 & \hfuzz=500pt filename & \hfuzz=500pt \\
\hfuzz=500pt\includegraphics[width=1em]{element-mustset.pdf}~inputfileStarCamera2 & \hfuzz=500pt filename & \hfuzz=500pt \\
\hfuzz=500pt\includegraphics[width=1em]{element-mustset.pdf}~earthRotation & \hfuzz=500pt \hyperref[earthRotationType]{earthRotation} & \hfuzz=500pt \\
\hfuzz=500pt\includegraphics[width=1em]{element.pdf}~ephemerides & \hfuzz=500pt \hyperref[ephemeridesType]{ephemerides} & \hfuzz=500pt \\
\hfuzz=500pt\includegraphics[width=1em]{element-unbounded.pdf}~gradientfield & \hfuzz=500pt \hyperref[gravityfieldType]{gravityfield} & \hfuzz=500pt low order field to estimate the change of the gravity by position adjustement\\
\hfuzz=500pt\includegraphics[width=1em]{element-unbounded.pdf}~parametrizationGravity & \hfuzz=500pt \hyperref[parametrizationGravityType]{parametrizationGravity} & \hfuzz=500pt gravity field parametrization\\
\hfuzz=500pt\includegraphics[width=1em]{element-unbounded.pdf}~parametrizationAcceleration1 & \hfuzz=500pt \hyperref[parametrizationAccelerationType]{parametrizationAcceleration} & \hfuzz=500pt orbit1 force parameters\\
\hfuzz=500pt\includegraphics[width=1em]{element-unbounded.pdf}~parametrizationAcceleration2 & \hfuzz=500pt \hyperref[parametrizationAccelerationType]{parametrizationAcceleration} & \hfuzz=500pt orbit2 force parameters\\
\hfuzz=500pt\includegraphics[width=1em]{element-unbounded.pdf}~parametrizationSst & \hfuzz=500pt \hyperref[parametrizationSatelliteTrackingType]{parametrizationSatelliteTracking} & \hfuzz=500pt satellite tracking parameter\\
\hfuzz=500pt\includegraphics[width=1em]{element.pdf}~keepSatelliteStates & \hfuzz=500pt boolean & \hfuzz=500pt set boundary values of each arc global\\
\hfuzz=500pt\includegraphics[width=1em]{element.pdf}~integrationDegree & \hfuzz=500pt uint & \hfuzz=500pt integration of forces by polynomial approximation of degree n\\
\hfuzz=500pt\includegraphics[width=1em]{element.pdf}~interpolationDegree & \hfuzz=500pt uint & \hfuzz=500pt orbit interpolation by polynomial approximation of degree n\\
\hfuzz=500pt\includegraphics[width=1em]{element-mustset.pdf}~covarianceSst & \hfuzz=500pt \hyperref[covarianceSstType]{covarianceSst} & \hfuzz=500pt covariance matrix of satellite to satellite tracking observations\\
\hfuzz=500pt\includegraphics[width=1em]{element-mustset.pdf}~covariancePod1 & \hfuzz=500pt \hyperref[covariancePodType]{covariancePod} & \hfuzz=500pt covariance matrix of kinematic orbits (satellite 1)\\
\hfuzz=500pt\includegraphics[width=1em]{element-mustset.pdf}~covariancePod2 & \hfuzz=500pt \hyperref[covariancePodType]{covariancePod} & \hfuzz=500pt covariance matrix of kinematic orbits (satellite 2)\\
\hline
\end{tabularx}


\subsection{DualSstVariational}\label{observationType:dualSstVariational}
Like \configClass{observation:sstVariational}{observationType:sstVariational} (see there for details)
but with two simultaneous satellite-to-satellite (SST) observations.

This class reads two SST observation files (\configFile{inputfileSatelliteTracking1}{instrument} and
\configFile{inputfileSatelliteTracking2}{instrument}).
Empirical parameters for the SST observations can be setup independently for both SST observation
types with \configClass{parametrizationSst1}{parametrizationSatelliteTrackingType} and
\configClass{parametrizationSst2}{parametrizationSatelliteTrackingType}.

Both SST observation types are reduced by the same background models and the same impact
of accelerometer measurements. The covariance matrix of the reduced observations should not consider
the the instrument noise only (\configClass{covarianceSst1/2}{covarianceSstType}) but must
take the cross correlations \configClass{covarianceAcc}{covarianceSstType} into account.
The covariance matrix of the reduced observations is given by
\begin{equation}
  \M\Sigma(\begin{bmatrix} \Delta l_{SST1} \\ \Delta l_{SST2} \end{bmatrix})
  = \begin{bmatrix} \M\Sigma_{SST1} + \M\Sigma_{ACC} & \M\Sigma_{ACC} \\
                   \M\Sigma_{ACC} & \M\Sigma_{SST2} + \M\Sigma_{ACC}
    \end{bmatrix}.
\end{equation}


\keepXColumns
\begin{tabularx}{\textwidth}{N T A}
\hline
Name & Type & Annotation\\
\hline
\hfuzz=500pt\includegraphics[width=1em]{element-mustset.pdf}~rightHandSide & \hfuzz=500pt sequence & \hfuzz=500pt input for observation vectors\\
\hfuzz=500pt\includegraphics[width=1em]{connector.pdf}\includegraphics[width=1em]{element-mustset-unbounded.pdf}~inputfileSatelliteTracking1 & \hfuzz=500pt filename & \hfuzz=500pt ranging observations and corrections\\
\hfuzz=500pt\includegraphics[width=1em]{connector.pdf}\includegraphics[width=1em]{element-mustset-unbounded.pdf}~inputfileSatelliteTracking2 & \hfuzz=500pt filename & \hfuzz=500pt ranging observations and corrections\\
\hfuzz=500pt\includegraphics[width=1em]{connector.pdf}\includegraphics[width=1em]{element.pdf}~inputfileOrbit1 & \hfuzz=500pt filename & \hfuzz=500pt kinematic positions of satellite A as observations\\
\hfuzz=500pt\includegraphics[width=1em]{connector.pdf}\includegraphics[width=1em]{element.pdf}~inputfileOrbit2 & \hfuzz=500pt filename & \hfuzz=500pt kinematic positions of satellite B as observations\\
\hfuzz=500pt\includegraphics[width=1em]{element-mustset.pdf}~sstType & \hfuzz=500pt choice & \hfuzz=500pt \\
\hfuzz=500pt\includegraphics[width=1em]{connector.pdf}\includegraphics[width=1em]{element-mustset.pdf}~range & \hfuzz=500pt  & \hfuzz=500pt \\
\hfuzz=500pt\includegraphics[width=1em]{connector.pdf}\includegraphics[width=1em]{element-mustset.pdf}~rangeRate & \hfuzz=500pt  & \hfuzz=500pt \\
\hfuzz=500pt\includegraphics[width=1em]{connector.pdf}\includegraphics[width=1em]{element-mustset.pdf}~none & \hfuzz=500pt  & \hfuzz=500pt \\
\hfuzz=500pt\includegraphics[width=1em]{element-mustset.pdf}~inputfileVariational1 & \hfuzz=500pt filename & \hfuzz=500pt approximate position and integrated state matrix\\
\hfuzz=500pt\includegraphics[width=1em]{element-mustset.pdf}~inputfileVariational2 & \hfuzz=500pt filename & \hfuzz=500pt approximate position and integrated state matrix\\
\hfuzz=500pt\includegraphics[width=1em]{element.pdf}~ephemerides & \hfuzz=500pt \hyperref[ephemeridesType]{ephemerides} & \hfuzz=500pt \\
\hfuzz=500pt\includegraphics[width=1em]{element-unbounded.pdf}~parametrizationGravity & \hfuzz=500pt \hyperref[parametrizationGravityType]{parametrizationGravity} & \hfuzz=500pt gravity field parametrization\\
\hfuzz=500pt\includegraphics[width=1em]{element-unbounded.pdf}~parametrizationAcceleration1 & \hfuzz=500pt \hyperref[parametrizationAccelerationType]{parametrizationAcceleration} & \hfuzz=500pt orbit1 force parameters\\
\hfuzz=500pt\includegraphics[width=1em]{element-unbounded.pdf}~parametrizationAcceleration2 & \hfuzz=500pt \hyperref[parametrizationAccelerationType]{parametrizationAcceleration} & \hfuzz=500pt orbit2 force parameters\\
\hfuzz=500pt\includegraphics[width=1em]{element-unbounded.pdf}~parametrizationSst1 & \hfuzz=500pt \hyperref[parametrizationSatelliteTrackingType]{parametrizationSatelliteTracking} & \hfuzz=500pt satellite tracking parameter for first ranging observations\\
\hfuzz=500pt\includegraphics[width=1em]{element-unbounded.pdf}~parametrizationSst2 & \hfuzz=500pt \hyperref[parametrizationSatelliteTrackingType]{parametrizationSatelliteTracking} & \hfuzz=500pt satellite tracking parameter for second ranging observations\\
\hfuzz=500pt\includegraphics[width=1em]{element.pdf}~integrationDegree & \hfuzz=500pt uint & \hfuzz=500pt integration of forces by polynomial approximation of degree n\\
\hfuzz=500pt\includegraphics[width=1em]{element.pdf}~interpolationDegree & \hfuzz=500pt uint & \hfuzz=500pt orbit interpolation by polynomial approximation of degree n\\
\hfuzz=500pt\includegraphics[width=1em]{element-mustset.pdf}~covarianceSst1 & \hfuzz=500pt \hyperref[covarianceSstType]{covarianceSst} & \hfuzz=500pt covariance matrix of first satellite to satellite tracking observations\\
\hfuzz=500pt\includegraphics[width=1em]{element-mustset.pdf}~covarianceSst2 & \hfuzz=500pt \hyperref[covarianceSstType]{covarianceSst} & \hfuzz=500pt covariance matrix of second satellite to satellite tracking observations\\
\hfuzz=500pt\includegraphics[width=1em]{element.pdf}~covarianceAcc & \hfuzz=500pt \hyperref[covarianceSstType]{covarianceSst} & \hfuzz=500pt common covariance matrix of reduced satellite to satellite tracking observations\\
\hfuzz=500pt\includegraphics[width=1em]{element-mustset.pdf}~covariancePod1 & \hfuzz=500pt \hyperref[covariancePodType]{covariancePod} & \hfuzz=500pt covariance matrix of kinematic orbits (satellite 1)\\
\hfuzz=500pt\includegraphics[width=1em]{element-mustset.pdf}~covariancePod2 & \hfuzz=500pt \hyperref[covariancePodType]{covariancePod} & \hfuzz=500pt covariance matrix of kinematic orbits (satellite 2)\\
\hline
\end{tabularx}


\subsection{Gradiometer}\label{observationType:gradiometer}
Observation equations for satellite gravity gradiometry (SGG)
\begin{equation}
  \nabla\nabla V(\M r) =
  \begin{pmatrix}
    \frac{\partial^2 V}{\partial x^2}         & \frac{\partial^2 V}{\partial x\partial y} & \frac{\partial^2 V}{\partial x\partial z} \\
    \frac{\partial^2 V}{\partial y\partial x} & \frac{\partial^2 V}{\partial y^2}         & \frac{\partial^2 V}{\partial y\partial z} \\
    \frac{\partial^2 V}{\partial z\partial x} & \frac{\partial^2 V}{\partial z\partial y} & \frac{\partial^2 V}{\partial z^2}
  \end{pmatrix}.
\end{equation}
From the \configFile{inputfileGradiometer}{instrument} observations precomputed \configFile{inputfileReferenceGradiometer}{instrument}
together with other background models are reduced, all given in \configClass{rightHandSide}{sggRightSideType}.

All instrument data \configFile{inputfileGradiometer}{instrument}, \configFile{inputfileOrbit}{instrument},
and \configFile{inputfileStarCamera}{instrument} must be synchronous and be diveded
into each short arcs (see \program{InstrumentSynchronize}).

Additional to the \configClass{parametrizationGravity}{parametrizationGravityType}
an (temporal changing) bias for each gradiometer component and arc can be estimated with
\configClass{parametrizationBias}{parametrizationTemporalType}.

The accuracy or the full covariance matrix of the gradiometer is provided in
\config{covarianceSgg} and can be estimated with \program{PreprocessingGradiometer}.



\keepXColumns
\begin{tabularx}{\textwidth}{N T A}
\hline
Name & Type & Annotation\\
\hline
\hfuzz=500pt\includegraphics[width=1em]{element-mustset-unbounded.pdf}~rightHandSide & \hfuzz=500pt \hyperref[sggRightSideType]{sggRightSide} & \hfuzz=500pt input for the observation vector\\
\hfuzz=500pt\includegraphics[width=1em]{element-mustset.pdf}~inputfileOrbit & \hfuzz=500pt filename & \hfuzz=500pt \\
\hfuzz=500pt\includegraphics[width=1em]{element-mustset.pdf}~inputfileStarCamera & \hfuzz=500pt filename & \hfuzz=500pt \\
\hfuzz=500pt\includegraphics[width=1em]{element-mustset.pdf}~earthRotation & \hfuzz=500pt \hyperref[earthRotationType]{earthRotation} & \hfuzz=500pt \\
\hfuzz=500pt\includegraphics[width=1em]{element.pdf}~ephemerides & \hfuzz=500pt \hyperref[ephemeridesType]{ephemerides} & \hfuzz=500pt \\
\hfuzz=500pt\includegraphics[width=1em]{element-mustset-unbounded.pdf}~parametrizationGravity & \hfuzz=500pt \hyperref[parametrizationGravityType]{parametrizationGravity} & \hfuzz=500pt \\
\hfuzz=500pt\includegraphics[width=1em]{element-unbounded.pdf}~parametrizationBias & \hfuzz=500pt \hyperref[parametrizationTemporalType]{parametrizationTemporal} & \hfuzz=500pt per arc\\
\hfuzz=500pt\includegraphics[width=1em]{element.pdf}~useXX & \hfuzz=500pt boolean & \hfuzz=500pt \\
\hfuzz=500pt\includegraphics[width=1em]{element.pdf}~useYY & \hfuzz=500pt boolean & \hfuzz=500pt \\
\hfuzz=500pt\includegraphics[width=1em]{element.pdf}~useZZ & \hfuzz=500pt boolean & \hfuzz=500pt \\
\hfuzz=500pt\includegraphics[width=1em]{element.pdf}~useXY & \hfuzz=500pt boolean & \hfuzz=500pt \\
\hfuzz=500pt\includegraphics[width=1em]{element.pdf}~useXZ & \hfuzz=500pt boolean & \hfuzz=500pt \\
\hfuzz=500pt\includegraphics[width=1em]{element.pdf}~useYZ & \hfuzz=500pt boolean & \hfuzz=500pt \\
\hfuzz=500pt\includegraphics[width=1em]{element-mustset.pdf}~covarianceSgg & \hfuzz=500pt sequence & \hfuzz=500pt \\
\hfuzz=500pt\includegraphics[width=1em]{connector.pdf}\includegraphics[width=1em]{element.pdf}~sigma & \hfuzz=500pt double & \hfuzz=500pt general variance factor\\
\hfuzz=500pt\includegraphics[width=1em]{connector.pdf}\includegraphics[width=1em]{element.pdf}~inputfileSigmasPerArc & \hfuzz=500pt filename & \hfuzz=500pt different accuaries for each arc (multplicated with sigma)\\
\hfuzz=500pt\includegraphics[width=1em]{connector.pdf}\includegraphics[width=1em]{element.pdf}~inputfileCovarianceFunction & \hfuzz=500pt filename & \hfuzz=500pt covariance function in time\\
\hline
\end{tabularx}


\subsection{Terrestrial}\label{observationType:terrestrial}
The gravity field is estimated from point wise measurements.
The gravity field parametrization is given by \configClass{parametrizationGravity}{parametrizationGravityType}.
There is no need to have the data regular distributed or given on a sphere or ellipsoid.
The type of the gridded data (e.g gravity anomalies or geoid heights)
must be set with \configClass{kernel}{kernelType}.
A \configClass{referencefield}{gravityfieldType} can be reduced beforehand.

The observations at given positions are calculated from
\configFile{inputfileGriddedData}{griddedData}.
The input columns are enumerated by \verb|data0|,~\verb|data1|,~\ldots,
see~\reference{dataVariables}{general.parser:dataVariables}.

The observations can be divided into small blocks for parallelization.
With \config{blockingSize} set the maximum count of observations in each block.


\keepXColumns
\begin{tabularx}{\textwidth}{N T A}
\hline
Name & Type & Annotation\\
\hline
\hfuzz=500pt\includegraphics[width=1em]{element-mustset.pdf}~rightHandSide & \hfuzz=500pt sequence & \hfuzz=500pt input for observation vectors\\
\hfuzz=500pt\includegraphics[width=1em]{connector.pdf}\includegraphics[width=1em]{element-mustset.pdf}~inputfileGriddedData & \hfuzz=500pt filename & \hfuzz=500pt \\
\hfuzz=500pt\includegraphics[width=1em]{connector.pdf}\includegraphics[width=1em]{element-mustset.pdf}~observation & \hfuzz=500pt expression & \hfuzz=500pt [SI units]\\
\hfuzz=500pt\includegraphics[width=1em]{connector.pdf}\includegraphics[width=1em]{element.pdf}~sigma & \hfuzz=500pt expression & \hfuzz=500pt accuracy, 1/sigma used as weighting\\
\hfuzz=500pt\includegraphics[width=1em]{connector.pdf}\includegraphics[width=1em]{element-unbounded.pdf}~referencefield & \hfuzz=500pt \hyperref[gravityfieldType]{gravityfield} & \hfuzz=500pt \\
\hfuzz=500pt\includegraphics[width=1em]{element-mustset.pdf}~kernel & \hfuzz=500pt \hyperref[kernelType]{kernel} & \hfuzz=500pt type of observations\\
\hfuzz=500pt\includegraphics[width=1em]{element-mustset-unbounded.pdf}~parametrizationGravity & \hfuzz=500pt \hyperref[parametrizationGravityType]{parametrizationGravity} & \hfuzz=500pt \\
\hfuzz=500pt\includegraphics[width=1em]{element.pdf}~time & \hfuzz=500pt time & \hfuzz=500pt for reference field and parametrization\\
\hfuzz=500pt\includegraphics[width=1em]{element.pdf}~blockingSize & \hfuzz=500pt uint & \hfuzz=500pt segementation of the obervations if designmatrix can't be build at once\\
\hline
\end{tabularx}


\subsection{Deflections}\label{observationType:deflections}
The gravity field parametrized by \configClass{parametrizationGravity}{parametrizationGravityType}
is estimated from deflections of the vertical measurements.
A \configClass{referencefield}{gravityfieldType} can be reduced beforehand.

The observations $\xi$ in north direction and $\eta$ in east direction
at given positions are calculated from
\configFile{inputfileGriddedData}{griddedData}.
The input columns are enumerated by \verb|data0|,~\verb|data1|,~\ldots,
see~\reference{dataVariables}{general.parser:dataVariables}.

The ellipsoid parameters \config{R} and \config{inverseFlattening} are used
to define the local normal direction.

The observations can be divided into small blocks for parallelization.
With \config{blockingSize} set the maximum count of observations in each block.


\keepXColumns
\begin{tabularx}{\textwidth}{N T A}
\hline
Name & Type & Annotation\\
\hline
\hfuzz=500pt\includegraphics[width=1em]{element-mustset.pdf}~rightHandSide & \hfuzz=500pt sequence & \hfuzz=500pt input for observation vectors\\
\hfuzz=500pt\includegraphics[width=1em]{connector.pdf}\includegraphics[width=1em]{element-mustset.pdf}~inputfileGriddedData & \hfuzz=500pt filename & \hfuzz=500pt \\
\hfuzz=500pt\includegraphics[width=1em]{connector.pdf}\includegraphics[width=1em]{element-mustset.pdf}~observationXi & \hfuzz=500pt expression & \hfuzz=500pt North-South Deflections of the Vertical [rad]\\
\hfuzz=500pt\includegraphics[width=1em]{connector.pdf}\includegraphics[width=1em]{element-mustset.pdf}~observationEta & \hfuzz=500pt expression & \hfuzz=500pt East-West Deflections of the Vertical  [rad]\\
\hfuzz=500pt\includegraphics[width=1em]{connector.pdf}\includegraphics[width=1em]{element.pdf}~sigmaXi & \hfuzz=500pt expression & \hfuzz=500pt accuracy, 1/sigma used as weighting\\
\hfuzz=500pt\includegraphics[width=1em]{connector.pdf}\includegraphics[width=1em]{element.pdf}~sigmaEta & \hfuzz=500pt expression & \hfuzz=500pt accuracy, 1/sigma used as weighting\\
\hfuzz=500pt\includegraphics[width=1em]{connector.pdf}\includegraphics[width=1em]{element-unbounded.pdf}~referencefield & \hfuzz=500pt \hyperref[gravityfieldType]{gravityfield} & \hfuzz=500pt \\
\hfuzz=500pt\includegraphics[width=1em]{element-mustset-unbounded.pdf}~parametrizationGravity & \hfuzz=500pt \hyperref[parametrizationGravityType]{parametrizationGravity} & \hfuzz=500pt \\
\hfuzz=500pt\includegraphics[width=1em]{element.pdf}~time & \hfuzz=500pt time & \hfuzz=500pt for reference field and parametrization\\
\hfuzz=500pt\includegraphics[width=1em]{element.pdf}~R & \hfuzz=500pt double & \hfuzz=500pt reference radius for ellipsoid\\
\hfuzz=500pt\includegraphics[width=1em]{element.pdf}~inverseFlattening & \hfuzz=500pt double & \hfuzz=500pt reference flattening for ellipsoid, 0: sphere\\
\hfuzz=500pt\includegraphics[width=1em]{element.pdf}~blockingSize & \hfuzz=500pt uint & \hfuzz=500pt segementation of the obervations if designmatrix can't be build at once\\
\hline
\end{tabularx}


\subsection{StationLoading}\label{observationType:stationLoading}
Observation equations for displacements of a list of stations
due to the effect of time variable loading masses. The displacement~$\M u$ of a station is calculated according to
\begin{equation}
\M u(\M r) = \frac{1}{\gamma}\sum_{n=0}^\infty \left[\frac{h_n}{1+k_n}V_n(\M r)\,\M e_{up}
+ R\frac{l_n}{1+k_n}\left(
 \frac{\partial V_n(\M r)}{\partial \M e_{north}}\M e_{north}
+\frac{\partial V_n(\M r)}{\partial \M e_{east}} \M e_{east}\right)\right],
\end{equation}
where $\gamma$ is the normal gravity, the load Love and Shida numbers $h_n,l_n$ are given by
\configFile{inputfileDeformationLoadLoveNumber}{matrix} and the load Love numbers $k_n$ are given by
\configFile{inputfilePotentialLoadLoveNumber}{matrix}.
The $V_n$ are the spherical harmonics expansion of degree $n$ of the full time variable
gravitational potential (potential of the loading mass + deformation potential)
parametrized by \configClass{parametrizationGravity}{parametrizationGravityType}.
Additional parameters can be setup to estimate the realization of the reference frame
of the station coordinates (\config{estimateTranslation},
\config{estimateRotation}, and \config{estimateScale}).

The observations at stations coordinates are calculated from
\configFile{inputfileGriddedData}{griddedData}.
The input columns are enumerated by \verb|data0|,~\verb|data1|,~\ldots,
see~\reference{dataVariables}{general.parser:dataVariables}.

The ellipsoid parameters \config{R} and \config{inverseFlattening} are used
to define the local frame (north, east, up).

See also \program{Gravityfield2DisplacementTimeSeries}.

Reference:
Rietbroek (2014): Retrieval of Sea Level and Surface Loading Variations from Geodetic Observations
and Model Simulations: an Integrated Approach, Bonn, 2014. - Dissertation,
\url{https://nbn-resolving.org/urn:nbn:de:hbz:5n-35460}


\keepXColumns
\begin{tabularx}{\textwidth}{N T A}
\hline
Name & Type & Annotation\\
\hline
\hfuzz=500pt\includegraphics[width=1em]{element-mustset.pdf}~rightHandSide & \hfuzz=500pt sequence & \hfuzz=500pt input for observation vectors\\
\hfuzz=500pt\includegraphics[width=1em]{connector.pdf}\includegraphics[width=1em]{element-mustset.pdf}~inputfileGriddedData & \hfuzz=500pt filename & \hfuzz=500pt station positions with displacement data\\
\hfuzz=500pt\includegraphics[width=1em]{connector.pdf}\includegraphics[width=1em]{element-mustset.pdf}~observationNorth & \hfuzz=500pt expression & \hfuzz=500pt displacement [m]\\
\hfuzz=500pt\includegraphics[width=1em]{connector.pdf}\includegraphics[width=1em]{element-mustset.pdf}~observationEast & \hfuzz=500pt expression & \hfuzz=500pt displacement [m]\\
\hfuzz=500pt\includegraphics[width=1em]{connector.pdf}\includegraphics[width=1em]{element-mustset.pdf}~observationUp & \hfuzz=500pt expression & \hfuzz=500pt displacement [m]\\
\hfuzz=500pt\includegraphics[width=1em]{connector.pdf}\includegraphics[width=1em]{element.pdf}~sigmaNorth & \hfuzz=500pt expression & \hfuzz=500pt accuracy, 1/sigma used as weighting\\
\hfuzz=500pt\includegraphics[width=1em]{connector.pdf}\includegraphics[width=1em]{element.pdf}~sigmaEast & \hfuzz=500pt expression & \hfuzz=500pt accuracy, 1/sigma used as weighting\\
\hfuzz=500pt\includegraphics[width=1em]{connector.pdf}\includegraphics[width=1em]{element.pdf}~sigmaUp & \hfuzz=500pt expression & \hfuzz=500pt accuracy, 1/sigma used as weighting\\
\hfuzz=500pt\includegraphics[width=1em]{connector.pdf}\includegraphics[width=1em]{element.pdf}~inGlobalFrame & \hfuzz=500pt boolean & \hfuzz=500pt obs/sigmas given in global x,y,z frame instead of north,east,up\\
\hfuzz=500pt\includegraphics[width=1em]{connector.pdf}\includegraphics[width=1em]{element-unbounded.pdf}~referencefield & \hfuzz=500pt \hyperref[gravityfieldType]{gravityfield} & \hfuzz=500pt \\
\hfuzz=500pt\includegraphics[width=1em]{element.pdf}~time & \hfuzz=500pt time & \hfuzz=500pt for reference field and parametrization\\
\hfuzz=500pt\includegraphics[width=1em]{element-mustset-unbounded.pdf}~parametrizationGravity & \hfuzz=500pt \hyperref[parametrizationGravityType]{parametrizationGravity} & \hfuzz=500pt of loading (+defo) potential\\
\hfuzz=500pt\includegraphics[width=1em]{element.pdf}~estimateTranslation & \hfuzz=500pt boolean & \hfuzz=500pt coordinate center\\
\hfuzz=500pt\includegraphics[width=1em]{element.pdf}~estimateScale & \hfuzz=500pt boolean & \hfuzz=500pt scale factor of position\\
\hfuzz=500pt\includegraphics[width=1em]{element.pdf}~estimateRotation & \hfuzz=500pt boolean & \hfuzz=500pt rotation\\
\hfuzz=500pt\includegraphics[width=1em]{element-mustset.pdf}~inputfileDeformationLoadLoveNumber & \hfuzz=500pt filename & \hfuzz=500pt \\
\hfuzz=500pt\includegraphics[width=1em]{element.pdf}~inputfilePotentialLoadLoveNumber & \hfuzz=500pt filename & \hfuzz=500pt if full potential is given and not only loading potential\\
\hfuzz=500pt\includegraphics[width=1em]{element.pdf}~R & \hfuzz=500pt double & \hfuzz=500pt reference radius for ellipsoid\\
\hfuzz=500pt\includegraphics[width=1em]{element.pdf}~inverseFlattening & \hfuzz=500pt double & \hfuzz=500pt reference flattening for ellipsoid, 0: sphere\\
\hline
\end{tabularx}

\clearpage
%==================================

\section{OrbitPropagator}\label{orbitPropagatorType}
Implements the propagation of a satellite orbit under
the influence of \configClass{forces}{forcesType} as
used in \program{SimulateOrbit}
(dynamic orbits from numerical orbit integration).


\subsection{Euler}
This class implements Euler's method to propagate a satellite orbit under the influence of \configClass{Forces}{forcesType}.
Satellite is assumed to be oriented along-track.


\subsection{RungeKutta4}
This class implements the classical Runge-Kutta 4 method of orbit propagation
for satellite orbit under the influence of \configClass{Forces}{forcesType}.
No step-width control or other advanced features are implemented.
Satellite is assumed to be oriented along-track.
See: Montenbruck, Oliver, and Eberhard Gill. 2000. Satellite Orbits


\subsection{AdamsBashforthMoulton}
This class implements the Adams-Moulton class of predictor-corrector orbit propagators
for a satellite orbit under the influence of \configClass{Forces}{forcesType} using an implicit
Adams-Bashforth corrector. The coefficients for the propagator are derived using the equations
given in section 4.2.3 of [1]. Satellite is assumed to be oriented along-track.
[1] Montenbruck, Oliver, and Eberhard Gill. 2000. Satellite Orbits


\keepXColumns
\begin{tabularx}{\textwidth}{N T A}
\hline
Name & Type & Annotation\\
\hline
\hfuzz=500pt\includegraphics[width=1em]{element-mustset.pdf}~order & \hfuzz=500pt uint & \hfuzz=500pt Order of the Adams-Bashforth type propagator.\\
\hfuzz=500pt\includegraphics[width=1em]{element.pdf}~applyMoultonCorrector & \hfuzz=500pt boolean & \hfuzz=500pt Corrector step after Adams-Bashforth predcition.\\
\hfuzz=500pt\includegraphics[width=1em]{element-mustset.pdf}~warmup & \hfuzz=500pt \hyperref[orbitPropagatorType]{orbitPropagator} & \hfuzz=500pt \\
\hline
\end{tabularx}


\subsection{StoermerCowell}
This class implements the Stoermer-Cowell class of predictor-corrector orbit propagators for a satellite orbit
under the influence of \configClass{Forces}{forcesType}. The coefficients for the Stoermer predictor and Cowell corrector
are derived using the equations given in section 4.2.6 of [1]. Stoermer-Cowell is a double integration algorithm,
yielding positions directly from accelertions. It does not produce velocities. The velocities are derived using
Adams-type propagators as suggested in [2]. Satellite is assumed to be oriented along-track.
[1] Montenbruck, Oliver, and Eberhard Gill. 2000. Satellite Orbits
[2] Berry, Matthew M., and Liam M. Healy. 2004. “Implementation of Gauss-Jackson Integration for Orbit Propagation.”


\keepXColumns
\begin{tabularx}{\textwidth}{N T A}
\hline
Name & Type & Annotation\\
\hline
\hfuzz=500pt\includegraphics[width=1em]{element-mustset.pdf}~order & \hfuzz=500pt uint & \hfuzz=500pt Order of the Stoermer-Cowell type propagator.\\
\hfuzz=500pt\includegraphics[width=1em]{element-mustset.pdf}~warmup & \hfuzz=500pt \hyperref[orbitPropagatorType]{orbitPropagator} & \hfuzz=500pt \\
\hline
\end{tabularx}


\subsection{GaussJackson}
This class implements the Gauss-Jackson multi-step predictor-corrector method to
propagate a satellite orbit under the influence of \configClass{Forces}{forcesType}.
Satellite is assumed to be oriented along-track. Implementation is based on [1].
[1] Berry, Matthew M., and Liam M. Healy. 2004. “Implementation of Gauss-Jackson Integration for Orbit Propagation.”


\keepXColumns
\begin{tabularx}{\textwidth}{N T A}
\hline
Name & Type & Annotation\\
\hline
\hfuzz=500pt\includegraphics[width=1em]{element-mustset.pdf}~order & \hfuzz=500pt uint & \hfuzz=500pt of Gauss-Jackson method.\\
\hfuzz=500pt\includegraphics[width=1em]{element-mustset.pdf}~warmup & \hfuzz=500pt \hyperref[orbitPropagatorType]{orbitPropagator} & \hfuzz=500pt \\
\hfuzz=500pt\includegraphics[width=1em]{element.pdf}~correctorIterations & \hfuzz=500pt uint & \hfuzz=500pt Maximum number of iterations to run the corrector step for.\\
\hfuzz=500pt\includegraphics[width=1em]{element.pdf}~epsilon & \hfuzz=500pt double & \hfuzz=500pt Convergence criteria for position, velocity, and acceleration tests.\\
\hline
\end{tabularx}


\subsection{Polynomial}\label{orbitPropagatorType:Polynomial}
This class implements an integration Polynomial method to propagate a satellite orbit under
the influence of \configClass{Forces}{forcesType}. Satellite is assumed to be oriented along-track.
Implementation is based on code by Torsten Mayer-Gürr.


\keepXColumns
\begin{tabularx}{\textwidth}{N T A}
\hline
Name & Type & Annotation\\
\hline
\hfuzz=500pt\includegraphics[width=1em]{element-mustset.pdf}~degree & \hfuzz=500pt uint & \hfuzz=500pt polynomial degree to integrate accelerations\\
\hfuzz=500pt\includegraphics[width=1em]{element-mustset.pdf}~shift & \hfuzz=500pt int & \hfuzz=500pt shift polynomial in future (predicted accelerations)\\
\hfuzz=500pt\includegraphics[width=1em]{element-mustset.pdf}~epsilon & \hfuzz=500pt double & \hfuzz=500pt [m] max. position change to recompute forces\\
\hfuzz=500pt\includegraphics[width=1em]{element-mustset.pdf}~warmup & \hfuzz=500pt \hyperref[orbitPropagatorType]{orbitPropagator} & \hfuzz=500pt to compute epochs before start epoch\\
\hfuzz=500pt\includegraphics[width=1em]{element.pdf}~corrector & \hfuzz=500pt boolean & \hfuzz=500pt apply corrector iteration if position change is larger than epsilon\\
\hline
\end{tabularx}


\subsection{File}
Reads an orbit from file. If the needed epochs are not given an exception is thrown.


\keepXColumns
\begin{tabularx}{\textwidth}{N T A}
\hline
Name & Type & Annotation\\
\hline
\hfuzz=500pt\includegraphics[width=1em]{element-mustset.pdf}~inputfileOrbit & \hfuzz=500pt filename & \hfuzz=500pt epoch at timeStart is not used\\
\hfuzz=500pt\includegraphics[width=1em]{element-mustset.pdf}~margin & \hfuzz=500pt double & \hfuzz=500pt [seconds] to find identical times\\
\hfuzz=500pt\includegraphics[width=1em]{element.pdf}~recomputeForces & \hfuzz=500pt boolean & \hfuzz=500pt \\
\hline
\end{tabularx}

\clearpage
%==================================

\section{ParameterNames}\label{parameterNamesType}
Generate a list of parameter names. All parameters are appended.


\subsection{Name}
The parameter is given by explicitly by four parts:
\begin{enumerate}
\item object: Object this parameter refers to, e.g. \verb|graceA|, \verb|G023|, \verb|earth|, \ldots
\item type: Type of this parameter, e.g. \verb|accBias|, \verb|position.x|, \ldots
\item temporal: Temporal representation of this parameter, e.g. \verb|trend|, \verb|polynomial.degree1|, \ldots
\item interval: Interval/epoch this parameter represents, e.g. \verb|2017-01-01_00-00-00_2017-01-02_00-00-00|, \verb|2018-01-01_00-00-00|.
\end{enumerate}


\keepXColumns
\begin{tabularx}{\textwidth}{N T A}
\hline
Name & Type & Annotation\\
\hline
\hfuzz=500pt\includegraphics[width=1em]{element.pdf}~object & \hfuzz=500pt string & \hfuzz=500pt object this parameter refers to, e.g. graceA, G023, earth\\
\hfuzz=500pt\includegraphics[width=1em]{element.pdf}~type & \hfuzz=500pt string & \hfuzz=500pt type of this parameter, e.g. accBias, position.x\\
\hfuzz=500pt\includegraphics[width=1em]{element.pdf}~temporal & \hfuzz=500pt string & \hfuzz=500pt temporal representation of this parameter, e.g. trend, polynomial.degree1\\
\hfuzz=500pt\includegraphics[width=1em]{element.pdf}~interval & \hfuzz=500pt string & \hfuzz=500pt interval/epoch this parameter refers to, e.g. 2017-01-01\_00-00-00\_2017-01-02\_00-00-00, 2008-01-01\_00-00-00\\
\hline
\end{tabularx}


\subsection{File}
Read parameter names from \file{file}{parameterName}.


\keepXColumns
\begin{tabularx}{\textwidth}{N T A}
\hline
Name & Type & Annotation\\
\hline
\hfuzz=500pt\includegraphics[width=1em]{element-mustset.pdf}~inputfileParameterNames & \hfuzz=500pt filename & \hfuzz=500pt file with parameter names\\
\hline
\end{tabularx}


\subsection{Gravity}
Parameter names of gravity \configClass{parametrization}{parametrizationGravityType}.
An additional \config{object} name can be included in the parameter names.


\keepXColumns
\begin{tabularx}{\textwidth}{N T A}
\hline
Name & Type & Annotation\\
\hline
\hfuzz=500pt\includegraphics[width=1em]{element.pdf}~object & \hfuzz=500pt string & \hfuzz=500pt object these parameters refers to, e.g. earth\\
\hfuzz=500pt\includegraphics[width=1em]{element-mustset-unbounded.pdf}~parametrization & \hfuzz=500pt \hyperref[parametrizationGravityType]{parametrizationGravity} & \hfuzz=500pt \\
\hline
\end{tabularx}


\subsection{Acceleration}
Parameter names of satellite acceleration \configClass{parametrization}{parametrizationAccelerationType}.
Arc related parameters are appended if an \configFile{inputfileInstrument}{instrument} is provided which
defines the arc structure.
An additional \config{object} name can be included in the parameter names.


\keepXColumns
\begin{tabularx}{\textwidth}{N T A}
\hline
Name & Type & Annotation\\
\hline
\hfuzz=500pt\includegraphics[width=1em]{element.pdf}~object & \hfuzz=500pt string & \hfuzz=500pt object these parameters refers to, e.g. graceA, G023\\
\hfuzz=500pt\includegraphics[width=1em]{element-mustset-unbounded.pdf}~parameterization & \hfuzz=500pt \hyperref[parametrizationAccelerationType]{parametrizationAcceleration} & \hfuzz=500pt \\
\hfuzz=500pt\includegraphics[width=1em]{element.pdf}~inputfileInstrument & \hfuzz=500pt filename & \hfuzz=500pt defines the arc structure for arc related parameters\\
\hline
\end{tabularx}


\subsection{SatelliteTracking}
Parameter names of satellite tracking \configClass{parametrization}{parametrizationSatelliteTrackingType}.
An additional \config{object} name can be included in the parameter names.


\keepXColumns
\begin{tabularx}{\textwidth}{N T A}
\hline
Name & Type & Annotation\\
\hline
\hfuzz=500pt\includegraphics[width=1em]{element.pdf}~object & \hfuzz=500pt string & \hfuzz=500pt object these parameters refers to, e.g. grace1.grace2\\
\hfuzz=500pt\includegraphics[width=1em]{element-mustset-unbounded.pdf}~parameterization & \hfuzz=500pt \hyperref[parametrizationSatelliteTrackingType]{parametrizationSatelliteTracking} & \hfuzz=500pt \\
\hline
\end{tabularx}


\subsection{Temporal}
Parameter names from temporal parametrization.
It is possible to setup the temporal parameters for each \configClass{parameterNameBase}{parameterNamesType}.


\keepXColumns
\begin{tabularx}{\textwidth}{N T A}
\hline
Name & Type & Annotation\\
\hline
\hfuzz=500pt\includegraphics[width=1em]{element-unbounded.pdf}~parameterNameBase & \hfuzz=500pt \hyperref[parameterNamesType]{parameterNames} & \hfuzz=500pt \\
\hfuzz=500pt\includegraphics[width=1em]{element-mustset-unbounded.pdf}~parametrizationTemporal & \hfuzz=500pt \hyperref[parametrizationTemporalType]{parametrizationTemporal} & \hfuzz=500pt \\
\hline
\end{tabularx}


\subsection{GnssAntenna}
Parameter names of GNSS antenna center variation \configClass{parametrization}{parametrizationGnssAntennaType}.
An additional \config{object} name (antenna name) can be included in the parameter names.
It is possible to setup the parameters for each \configClass{gnssType}{gnssType}.


\keepXColumns
\begin{tabularx}{\textwidth}{N T A}
\hline
Name & Type & Annotation\\
\hline
\hfuzz=500pt\includegraphics[width=1em]{element.pdf}~object & \hfuzz=500pt string & \hfuzz=500pt antenna name\\
\hfuzz=500pt\includegraphics[width=1em]{element-mustset-unbounded.pdf}~parametrization & \hfuzz=500pt \hyperref[parametrizationGnssAntennaType]{parametrizationGnssAntenna} & \hfuzz=500pt \\
\hfuzz=500pt\includegraphics[width=1em]{element-unbounded.pdf}~gnssType & \hfuzz=500pt \hyperref[gnssType]{gnssType} & \hfuzz=500pt e.g. C1CG**\\
\hline
\end{tabularx}


\subsection{Observation}
Parameter names used in \configClass{observation equations}{observationType}.


\keepXColumns
\begin{tabularx}{\textwidth}{N T A}
\hline
Name & Type & Annotation\\
\hline
\hfuzz=500pt\includegraphics[width=1em]{element-mustset.pdf}~observation & \hfuzz=500pt \hyperref[observationType]{observation} & \hfuzz=500pt \\
\hline
\end{tabularx}


\subsection{Rename}
Replaces parts of \configClass{parameterName}{parameterNamesType}s.
The star "\verb|*|" left this part untouched.


\keepXColumns
\begin{tabularx}{\textwidth}{N T A}
\hline
Name & Type & Annotation\\
\hline
\hfuzz=500pt\includegraphics[width=1em]{element-mustset-unbounded.pdf}~parameterName & \hfuzz=500pt \hyperref[parameterNamesType]{parameterNames} & \hfuzz=500pt \\
\hfuzz=500pt\includegraphics[width=1em]{element.pdf}~object & \hfuzz=500pt string & \hfuzz=500pt *: left this part untouched, object\\
\hfuzz=500pt\includegraphics[width=1em]{element.pdf}~type & \hfuzz=500pt string & \hfuzz=500pt *: left this part untouched, type\\
\hfuzz=500pt\includegraphics[width=1em]{element.pdf}~temporal & \hfuzz=500pt string & \hfuzz=500pt *: left this part untouched, temporal representation\\
\hfuzz=500pt\includegraphics[width=1em]{element.pdf}~interval & \hfuzz=500pt string & \hfuzz=500pt *: left this part untouched, interval/epoch\\
\hline
\end{tabularx}


\subsection{Selection}
Select a subset of \configClass{parameterName}{parameterNamesType}s
using \configClass{parameterSelection}{parameterSelectorType}.


\keepXColumns
\begin{tabularx}{\textwidth}{N T A}
\hline
Name & Type & Annotation\\
\hline
\hfuzz=500pt\includegraphics[width=1em]{element-mustset-unbounded.pdf}~parameterName & \hfuzz=500pt \hyperref[parameterNamesType]{parameterNames} & \hfuzz=500pt \\
\hfuzz=500pt\includegraphics[width=1em]{element-mustset-unbounded.pdf}~parameterSelection & \hfuzz=500pt \hyperref[parameterSelectorType]{parameterSelector} & \hfuzz=500pt parameter order/selection\\
\hline
\end{tabularx}


\subsection{WithoutDuplicates}
Removes all duplicate names (keep first) from \configClass{parameterName}{parameterNamesType}.


\keepXColumns
\begin{tabularx}{\textwidth}{N T A}
\hline
Name & Type & Annotation\\
\hline
\hfuzz=500pt\includegraphics[width=1em]{element-mustset-unbounded.pdf}~parameterName & \hfuzz=500pt \hyperref[parameterNamesType]{parameterNames} & \hfuzz=500pt \\
\hline
\end{tabularx}

\clearpage
%==================================

\section{ParameterSelector}\label{parameterSelectorType}
This class provides an index vector from selected parameters,
which can be used e.g. to reorder a normal equation matrix.
The size of the index vector determines the size of the new matrix.
Entries are the indices of the selected parameters in the provided
parameter list or NULLINDEX for zero/new parameters.


\subsection{Wildcard}\label{parameterSelectorType:wildcard}
Parameter index vector from name. Name matching supports wildcards * for any number of characters and ? for exactly one character.
Does not add zero/empty parameters if there are no matches.


\keepXColumns
\begin{tabularx}{\textwidth}{N T A}
\hline
Name & Type & Annotation\\
\hline
\hfuzz=500pt\includegraphics[width=1em]{element.pdf}~object & \hfuzz=500pt string & \hfuzz=500pt object this parameter refers to, e.g. graceA, G023, earth (wildcards: * and ?)\\
\hfuzz=500pt\includegraphics[width=1em]{element.pdf}~type & \hfuzz=500pt string & \hfuzz=500pt type of this parameter, e.g. accBias, position.x (wildcards: * and ?)\\
\hfuzz=500pt\includegraphics[width=1em]{element.pdf}~temporal & \hfuzz=500pt string & \hfuzz=500pt temporal representation of this parameter, e.g. trend, polynomial.degree1 (wildcards: * and ?)\\
\hfuzz=500pt\includegraphics[width=1em]{element.pdf}~interval & \hfuzz=500pt string & \hfuzz=500pt interval/epoch this parameter refers to, e.g. 2017-01-01\_00-00-00\_2017-01-02\_00-00-00, 2008-01-01\_00-00-00 (wildcards: * and ?)\\
\hline
\end{tabularx}


\subsection{Names}
Parameter index vector from list of parameter names.


\keepXColumns
\begin{tabularx}{\textwidth}{N T A}
\hline
Name & Type & Annotation\\
\hline
\hfuzz=500pt\includegraphics[width=1em]{element-mustset-unbounded.pdf}~parameterName & \hfuzz=500pt \hyperref[parameterNamesType]{parameterNames} & \hfuzz=500pt \\
\hline
\end{tabularx}


\subsection{Range}
Parameter index vector from range.


\keepXColumns
\begin{tabularx}{\textwidth}{N T A}
\hline
Name & Type & Annotation\\
\hline
\hfuzz=500pt\includegraphics[width=1em]{element-mustset.pdf}~start & \hfuzz=500pt expression & \hfuzz=500pt start at this index (variables: length)\\
\hfuzz=500pt\includegraphics[width=1em]{element.pdf}~count & \hfuzz=500pt expression & \hfuzz=500pt count of parameters, default: all parameters to the end (variables: length)\\
\hline
\end{tabularx}


\subsection{Matrix}
Parameter index vector from matrix.


\keepXColumns
\begin{tabularx}{\textwidth}{N T A}
\hline
Name & Type & Annotation\\
\hline
\hfuzz=500pt\includegraphics[width=1em]{element-mustset.pdf}~inputfileMatrix & \hfuzz=500pt filename & \hfuzz=500pt index in old parameter list or -1 for new parameter\\
\hfuzz=500pt\includegraphics[width=1em]{element.pdf}~column & \hfuzz=500pt expression & \hfuzz=500pt use this column (counting from 0, variables: columns)\\
\hfuzz=500pt\includegraphics[width=1em]{element.pdf}~startRow & \hfuzz=500pt expression & \hfuzz=500pt start at this row (counting from 0, variables: rows)\\
\hfuzz=500pt\includegraphics[width=1em]{element.pdf}~countRows & \hfuzz=500pt expression & \hfuzz=500pt use these many rows (default: use all, variables: rows)\\
\hline
\end{tabularx}


\subsection{Zeros}
Expand parameter index vector by adding zero parameters.


\keepXColumns
\begin{tabularx}{\textwidth}{N T A}
\hline
Name & Type & Annotation\\
\hline
\hfuzz=500pt\includegraphics[width=1em]{element-mustset.pdf}~count & \hfuzz=500pt expression & \hfuzz=500pt count of zero parameters (variables: length)\\
\hline
\end{tabularx}


\subsection{Complement}\label{parameterSelectorType:complement}
Parameter index vector from a complement of other parameter selector(s).


\keepXColumns
\begin{tabularx}{\textwidth}{N T A}
\hline
Name & Type & Annotation\\
\hline
\hfuzz=500pt\includegraphics[width=1em]{element-mustset-unbounded.pdf}~parameterSelection & \hfuzz=500pt \hyperref[parameterSelectorType]{parameterSelector} & \hfuzz=500pt parameter order/selection\\
\hline
\end{tabularx}

\clearpage
%==================================

\section{ParametrizationAcceleration}\label{parametrizationAccelerationType}
This class defines parameters of satellite accelerations.
It will be used to set up the design matrix in a least squares adjustment.
If multiple parametrizations are given the coefficents in the parameter vector
are sequently appended.


\subsection{PerRevolution}\label{parametrizationAccelerationType:perRevolution}
Oscillation per revolution.


\keepXColumns
\begin{tabularx}{\textwidth}{N T A}
\hline
Name & Type & Annotation\\
\hline
\hfuzz=500pt\includegraphics[width=1em]{element-mustset.pdf}~order & \hfuzz=500pt uint & \hfuzz=500pt once, twice, ...\\
\hfuzz=500pt\includegraphics[width=1em]{element.pdf}~estimateX & \hfuzz=500pt boolean & \hfuzz=500pt along\\
\hfuzz=500pt\includegraphics[width=1em]{element.pdf}~estimateY & \hfuzz=500pt boolean & \hfuzz=500pt cross\\
\hfuzz=500pt\includegraphics[width=1em]{element.pdf}~estimateZ & \hfuzz=500pt boolean & \hfuzz=500pt radial\\
\hfuzz=500pt\includegraphics[width=1em]{element-unbounded.pdf}~interval & \hfuzz=500pt \hyperref[timeSeriesType]{timeSeries} & \hfuzz=500pt setup new parameters each interval\\
\hfuzz=500pt\includegraphics[width=1em]{element.pdf}~perArc & \hfuzz=500pt boolean & \hfuzz=500pt \\
\hline
\end{tabularx}


\subsection{AccBias}\label{parametrizationAccelerationType:accBias}
Temporal changing accelerometer bias per axis.


\keepXColumns
\begin{tabularx}{\textwidth}{N T A}
\hline
Name & Type & Annotation\\
\hline
\hfuzz=500pt\includegraphics[width=1em]{element.pdf}~estimateX & \hfuzz=500pt boolean & \hfuzz=500pt along\\
\hfuzz=500pt\includegraphics[width=1em]{element.pdf}~estimateY & \hfuzz=500pt boolean & \hfuzz=500pt cross\\
\hfuzz=500pt\includegraphics[width=1em]{element.pdf}~estimateZ & \hfuzz=500pt boolean & \hfuzz=500pt radial\\
\hfuzz=500pt\includegraphics[width=1em]{element-mustset-unbounded.pdf}~temporal & \hfuzz=500pt \hyperref[parametrizationTemporalType]{parametrizationTemporal} & \hfuzz=500pt \\
\hfuzz=500pt\includegraphics[width=1em]{element.pdf}~perArc & \hfuzz=500pt boolean & \hfuzz=500pt \\
\hline
\end{tabularx}


\subsection{AccelerometerScaleFactors}\label{parametrizationAccelerationType:accelerometerScaleFactors}
Accelerometer scale factor per axis.


\keepXColumns
\begin{tabularx}{\textwidth}{N T A}
\hline
Name & Type & Annotation\\
\hline
\hfuzz=500pt\includegraphics[width=1em]{element-mustset.pdf}~inputfileAccelerometer & \hfuzz=500pt filename & \hfuzz=500pt \\
\hfuzz=500pt\includegraphics[width=1em]{element.pdf}~estimateX & \hfuzz=500pt boolean & \hfuzz=500pt along\\
\hfuzz=500pt\includegraphics[width=1em]{element.pdf}~estimateY & \hfuzz=500pt boolean & \hfuzz=500pt cross\\
\hfuzz=500pt\includegraphics[width=1em]{element.pdf}~estimateZ & \hfuzz=500pt boolean & \hfuzz=500pt radial\\
\hfuzz=500pt\includegraphics[width=1em]{element.pdf}~estimateCrossTalk & \hfuzz=500pt boolean & \hfuzz=500pt non-orthognality of axes\\
\hfuzz=500pt\includegraphics[width=1em]{element.pdf}~estimateRotation & \hfuzz=500pt boolean & \hfuzz=500pt misalignment\\
\hfuzz=500pt\includegraphics[width=1em]{element-mustset-unbounded.pdf}~temporal & \hfuzz=500pt \hyperref[parametrizationTemporalType]{parametrizationTemporal} & \hfuzz=500pt \\
\hfuzz=500pt\includegraphics[width=1em]{element.pdf}~perArc & \hfuzz=500pt boolean & \hfuzz=500pt \\
\hline
\end{tabularx}


\subsection{GnssSolarRadiation}\label{parametrizationAccelerationType:gnssSolarRadiation}
GNSS solar radiation pressure model.


\keepXColumns
\begin{tabularx}{\textwidth}{N T A}
\hline
Name & Type & Annotation\\
\hline
\hfuzz=500pt\includegraphics[width=1em]{element.pdf}~estimateD0 & \hfuzz=500pt boolean & \hfuzz=500pt constant term along D-axis (sat-sun vector)\\
\hfuzz=500pt\includegraphics[width=1em]{element.pdf}~estimateD2 & \hfuzz=500pt boolean & \hfuzz=500pt 2-per-rev terms along D-axis\\
\hfuzz=500pt\includegraphics[width=1em]{element.pdf}~estimateD4 & \hfuzz=500pt boolean & \hfuzz=500pt 4-per-rev terms along D-axis\\
\hfuzz=500pt\includegraphics[width=1em]{element.pdf}~estimateY0 & \hfuzz=500pt boolean & \hfuzz=500pt constant term along Y-axis (solar panel axis)\\
\hfuzz=500pt\includegraphics[width=1em]{element.pdf}~estimateB0 & \hfuzz=500pt boolean & \hfuzz=500pt constant term along B-axis (cross product D x Y)\\
\hfuzz=500pt\includegraphics[width=1em]{element.pdf}~estimateB1 & \hfuzz=500pt boolean & \hfuzz=500pt 1-per-rev terms along B-axis\\
\hfuzz=500pt\includegraphics[width=1em]{element.pdf}~estimateB3 & \hfuzz=500pt boolean & \hfuzz=500pt 3-per-rev terms along B-axis\\
\hfuzz=500pt\includegraphics[width=1em]{element.pdf}~perArc & \hfuzz=500pt boolean & \hfuzz=500pt \\
\hfuzz=500pt\includegraphics[width=1em]{element-mustset.pdf}~eclipse & \hfuzz=500pt \hyperref[eclipseType]{eclipse} & \hfuzz=500pt \\
\hline
\end{tabularx}


\subsection{ThermosphericDensity}\label{parametrizationAccelerationType:thermosphericDensity}
Estimate the thermospheric density along the orbit using a satllite macro model.
An optional thermospheric model can be used to compute temperature and wind.
The temperature is used to estimate variable drag and lift coefficients, otherwise a constant drag coefficient is used.


\keepXColumns
\begin{tabularx}{\textwidth}{N T A}
\hline
Name & Type & Annotation\\
\hline
\hfuzz=500pt\includegraphics[width=1em]{element.pdf}~thermosphere & \hfuzz=500pt \hyperref[thermosphereType]{thermosphere} & \hfuzz=500pt for wind and temperature\\
\hfuzz=500pt\includegraphics[width=1em]{element.pdf}~earthRotation & \hfuzz=500pt double & \hfuzz=500pt [rad/s]\\
\hfuzz=500pt\includegraphics[width=1em]{element.pdf}~considerTemperature & \hfuzz=500pt boolean & \hfuzz=500pt compute drag and lift, otherwise simple drag coefficient is used\\
\hfuzz=500pt\includegraphics[width=1em]{element.pdf}~considerWind & \hfuzz=500pt boolean & \hfuzz=500pt \\
\hfuzz=500pt\includegraphics[width=1em]{element-mustset-unbounded.pdf}~temporalDensity & \hfuzz=500pt \hyperref[parametrizationTemporalType]{parametrizationTemporal} & \hfuzz=500pt parameters along orbit\\
\hfuzz=500pt\includegraphics[width=1em]{element.pdf}~perArc & \hfuzz=500pt boolean & \hfuzz=500pt \\
\hline
\end{tabularx}


\subsection{ModelScale}\label{parametrizationAccelerationType:modelScale}
Estimate a scale factor for a given model.


\keepXColumns
\begin{tabularx}{\textwidth}{N T A}
\hline
Name & Type & Annotation\\
\hline
\hfuzz=500pt\includegraphics[width=1em]{element-mustset-unbounded.pdf}~miscAccelerations & \hfuzz=500pt \hyperref[miscAccelerationsType]{miscAccelerations} & \hfuzz=500pt \\
\hfuzz=500pt\includegraphics[width=1em]{element-mustset-unbounded.pdf}~temporal & \hfuzz=500pt \hyperref[parametrizationTemporalType]{parametrizationTemporal} & \hfuzz=500pt \\
\hfuzz=500pt\includegraphics[width=1em]{element.pdf}~perArc & \hfuzz=500pt boolean & \hfuzz=500pt \\
\hline
\end{tabularx}

\clearpage
%==================================

\section{ParametrizationGnssAntenna}\label{parametrizationGnssAntennaType}
Parametrization of antenna center variations. It will be used to set up the design matrix in a least squares adjustment.
Usually the paramtrization is setup separately for different \configClass{gnssType}{gnssType}.

If multiple parametrizations are given the parameters are sequently appended in the design matrix and parameter vector.


\subsection{Center}\label{parametrizationGnssAntennaType:center}
Antenna center or, if setup for a specific \configClass{gnssType}{gnssType},
phase/code center offset (e.g. \verb|*1*G| for GPS L1 phase center offset).


\keepXColumns
\begin{tabularx}{\textwidth}{N T A}
\hline
Name & Type & Annotation\\
\hline
\hfuzz=500pt\includegraphics[width=1em]{element.pdf}~estimateX & \hfuzz=500pt boolean & \hfuzz=500pt \\
\hfuzz=500pt\includegraphics[width=1em]{element.pdf}~estimateY & \hfuzz=500pt boolean & \hfuzz=500pt \\
\hfuzz=500pt\includegraphics[width=1em]{element.pdf}~estimateZ & \hfuzz=500pt boolean & \hfuzz=500pt \\
\hline
\end{tabularx}


\subsection{SphericalHarmonics}
Parametrization of antenna center variations in terms of spherical harmonics.
As usually only data above the horizon are observed only the even spherical harmonics
(degree/order $m+n$ even), which are symmetric to the equator, are setup.


\keepXColumns
\begin{tabularx}{\textwidth}{N T A}
\hline
Name & Type & Annotation\\
\hline
\hfuzz=500pt\includegraphics[width=1em]{element-mustset.pdf}~minDegree & \hfuzz=500pt uint & \hfuzz=500pt min degree\\
\hfuzz=500pt\includegraphics[width=1em]{element-mustset.pdf}~maxDegree & \hfuzz=500pt uint & \hfuzz=500pt max degree\\
\hline
\end{tabularx}


\subsection{RadialBasis}
Parametrization of antenna center variations with radial basis functions
\begin{equation}
  ACV(\M x(A, E)) = \sum_i a_i \Phi(\M x\cdot\M x_i)
\end{equation}
where $a_i$ the coefficients which has to be estimated and $\Phi$ are the basis
functions
\begin{equation}
  \Phi(\cos\psi) = \sum_n \sqrt{2n+1}P_n(\cos\psi).
\end{equation}

\fig{!hb}{0.4}{parametrizationGnssAntennaRadialBasis}{fig:parametrizationGnssAntennaRadialBasis}{Nodal points of the basis functions
using a Reuter grid for transmitting satellites (view angle of 18 deg). The red line indicates the view angle of 14 deg of ground stations.}


\keepXColumns
\begin{tabularx}{\textwidth}{N T A}
\hline
Name & Type & Annotation\\
\hline
\hfuzz=500pt\includegraphics[width=1em]{element-mustset-unbounded.pdf}~grid & \hfuzz=500pt \hyperref[gridType]{grid} & \hfuzz=500pt nodal points of shannon kernels\\
\hfuzz=500pt\includegraphics[width=1em]{element-mustset.pdf}~minDegree & \hfuzz=500pt uint & \hfuzz=500pt min degree of shannon kernel\\
\hfuzz=500pt\includegraphics[width=1em]{element-mustset.pdf}~maxDegree & \hfuzz=500pt uint & \hfuzz=500pt max degree of shannon kernel\\
\hline
\end{tabularx}

\clearpage
%==================================

\section{ParametrizationGravity}\label{parametrizationGravityType}
This class gives a parametrization of the time depending gravity field.
Together with the class \configClass{oberservation}{observationType} it will be used
to set up the design matrix in a least squares adjustment.
If multiple parametrizations are given the coefficents in the parameter vector
are sequently appended.


\subsection{SphericalHarmonics}\label{parametrizationGravityType:sphericalHarmonics}
The potential~$V$ is parametrized by a expansion of (fully normalized) spherical harmonics
\begin{equation}
V(\lambda,\vartheta,r) = \frac{GM}{R}\sum_{n=0}^\infty \sum_{m=0}^n \left(\frac{R}{r}\right)^{n+1}
  \left(c_{nm} C_{nm}(\lambda,\vartheta) + s_{nm} S_{nm}(\lambda,\vartheta)\right).
\end{equation}
You can set the range of degree~$n$ with \config{minDegree} and \config{maxDegree}.
The sorting sequence of the potential coefficients in the parameter vector can be defined by
\configClass{numbering}{sphericalHarmonicsNumberingType}.

The total count of parameters is $(n_{max}+1)^2-n_{min}^2$.


\keepXColumns
\begin{tabularx}{\textwidth}{N T A}
\hline
Name & Type & Annotation\\
\hline
\hfuzz=500pt\includegraphics[width=1em]{element-mustset.pdf}~minDegree & \hfuzz=500pt uint & \hfuzz=500pt \\
\hfuzz=500pt\includegraphics[width=1em]{element-mustset.pdf}~maxDegree & \hfuzz=500pt uint & \hfuzz=500pt \\
\hfuzz=500pt\includegraphics[width=1em]{element.pdf}~GM & \hfuzz=500pt double & \hfuzz=500pt Geocentric gravitational constant\\
\hfuzz=500pt\includegraphics[width=1em]{element.pdf}~R & \hfuzz=500pt double & \hfuzz=500pt reference radius\\
\hfuzz=500pt\includegraphics[width=1em]{element-mustset.pdf}~numbering & \hfuzz=500pt \hyperref[sphericalHarmonicsNumberingType]{sphericalHarmonicsNumbering} & \hfuzz=500pt numbering scheme\\
\hline
\end{tabularx}


\subsection{RadialBasis}\label{parametrizationGravityType:radialBasis}
The potential~$V$ is represented by a sum of space localizing basis functions
\begin{equation}
  V(\M x) = \sum_i a_i \Phi(\M x, \M x_i)
\end{equation}
where $a_i$ the coefficients which has to be estimated and $\Phi$ are the basis
functions given by isotropic radial \configClass{kernel}{kernelType} functions
\begin{equation}
  \Phi(\cos\psi,r,R) = \sum_n \left(\frac{R}{r}\right)^{n+1} k_n\sqrt{2n+1}\bar{P}_n(\cos\psi).
\end{equation}
The basis functions are located on a grid~$\M x_i$ given by \configClass{grid}{gridType}.
This class can also be used to estimate point masses if \configClass{kernel}{kernelType} is set to density.


\keepXColumns
\begin{tabularx}{\textwidth}{N T A}
\hline
Name & Type & Annotation\\
\hline
\hfuzz=500pt\includegraphics[width=1em]{element-mustset.pdf}~kernel & \hfuzz=500pt \hyperref[kernelType]{kernel} & \hfuzz=500pt shape of the radial basis function\\
\hfuzz=500pt\includegraphics[width=1em]{element-mustset-unbounded.pdf}~grid & \hfuzz=500pt \hyperref[gridType]{grid} & \hfuzz=500pt nodal point distribution\\
\hline
\end{tabularx}


\subsection{Temporal}
The time variable potential is given by
\begin{equation}
  V(\M x,t) = \sum_i V_i(\M x)\Psi_i(t),
\end{equation}
wehre $V_i(\M x)$ is the spatial parametrization of the gravity field
and can be choosen with \configClass{parametrizationGravity}{parametrizationGravityType}.
The parametrization in time domain $\Psi_i(t)$ is selected by
\configClass{parametrizationTemporal}{parametrizationTemporalType}.
The total parameter count is the parameter count of \configClass{parametrizationTemporal}{parametrizationTemporalType}
times the parameter count of \configClass{parametrizationGravity}{parametrizationGravityType}.


\keepXColumns
\begin{tabularx}{\textwidth}{N T A}
\hline
Name & Type & Annotation\\
\hline
\hfuzz=500pt\includegraphics[width=1em]{element-mustset-unbounded.pdf}~parametrizationGravity & \hfuzz=500pt \hyperref[parametrizationGravityType]{parametrizationGravity} & \hfuzz=500pt \\
\hfuzz=500pt\includegraphics[width=1em]{element-mustset-unbounded.pdf}~parametrizationTemporal & \hfuzz=500pt \hyperref[parametrizationTemporalType]{parametrizationTemporal} & \hfuzz=500pt \\
\hline
\end{tabularx}


\subsection{LinearTransformation}
Parametrization of the gravity field on the basis of a linear transformation of a source parametrization.
The linear transformation changes the original solution space represented by
\configClass{pararametrizationGravitySource}{parametrizationGravityType} from
\begin{equation}
  \mathbf{l} = \mathbf{A}\mathbf{x} + \mathbf{e}
\end{equation}
to
\begin{equation}
  \mathbf{l} = \mathbf{A}\mathbf{F}\mathbf{y} + \mathbf{e}
\end{equation}
through the linear transformation $\mathbf{x}=\mathbf{F}\mathbf{y}$.
It follows that the rows of the matrix $\mathbf{F}$ in \file{inputfileTransformationMatrix}{matrix} coincides with
the number of parameters in \configClass{pararametrizationGravitySource}{parametrizationGravityType}.
The new parameter count is given by the number of columns in $\mathbf{F}$ and may be smaller, equal or larger
than the original parameter count.


\keepXColumns
\begin{tabularx}{\textwidth}{N T A}
\hline
Name & Type & Annotation\\
\hline
\hfuzz=500pt\includegraphics[width=1em]{element-mustset-unbounded.pdf}~parametrizationGravitySource & \hfuzz=500pt \hyperref[parametrizationGravityType]{parametrizationGravity} & \hfuzz=500pt \\
\hfuzz=500pt\includegraphics[width=1em]{element-mustset.pdf}~inputfileTransformationMatrix & \hfuzz=500pt filename & \hfuzz=500pt transformation matrix from target to source parametrization (rows of this matrix must coincide with the parameter count of the source parametrization)\\
\hline
\end{tabularx}


\subsection{EarthquakeOscillation}
This class is used to estimate the earthquake oscillation function parameters,
i.e. $C_{nlm}$, $\omega_{nlm}$, and $P_{nlm}$.
The results describes the variation in the gravitational potential field caused by large earthquakes.
\begin{equation}
C_{lm}(\M t) = \sum_{n=0}^NC_{nlm}(1-\cos(\omega_{nlm}d\M t)\exp(P_{nlm}\omega_{nlm}d\M t)),
\end{equation}
with $\omega_{nlm}=\frac{2\pi}{T_{nlm}}$ and $P_{nlm}=\frac{-1}{2Q_{nlm}}$ . In this equation, $Q_{nlm}$ is the attenuation factor,
$n$ is the overtone factor, $m$ is degree, $l$ is order, and $t$ is time after earthquake in second.


\keepXColumns
\begin{tabularx}{\textwidth}{N T A}
\hline
Name & Type & Annotation\\
\hline
\hfuzz=500pt\includegraphics[width=1em]{element-mustset.pdf}~inputInitialCoefficient & \hfuzz=500pt filename & \hfuzz=500pt initial values for oscillation parameters\\
\hfuzz=500pt\includegraphics[width=1em]{element-mustset.pdf}~time0 & \hfuzz=500pt time & \hfuzz=500pt the time earthquake happened\\
\hfuzz=500pt\includegraphics[width=1em]{element-mustset.pdf}~minDegree & \hfuzz=500pt uint & \hfuzz=500pt \\
\hfuzz=500pt\includegraphics[width=1em]{element-mustset.pdf}~maxDegree & \hfuzz=500pt uint & \hfuzz=500pt \\
\hfuzz=500pt\includegraphics[width=1em]{element.pdf}~GM & \hfuzz=500pt double & \hfuzz=500pt Geocentric gravitational constant\\
\hfuzz=500pt\includegraphics[width=1em]{element.pdf}~R & \hfuzz=500pt double & \hfuzz=500pt reference radius\\
\hfuzz=500pt\includegraphics[width=1em]{element-mustset.pdf}~numbering & \hfuzz=500pt \hyperref[sphericalHarmonicsNumberingType]{sphericalHarmonicsNumbering} & \hfuzz=500pt numbering scheme\\
\hline
\end{tabularx}

\clearpage
%==================================

\section{ParametrizationSatelliteTracking}\label{parametrizationSatelliteTrackingType}
This class defines parameters of Satellite-to-Satellite tracking observations.
It will be used to set up the design matrix in a least squares adjustment.
If multiple parametrizations are given the coefficents in the parameter vector
are sequently appended.


\subsection{AntennaCenter}\label{parametrizationSatelliteTrackingType:antennaCenter}
Estimate the KBR antenna phase centre (APC) coordinates for each spacecraft in satellite reference frame (SRF)
as constant per axis, once per month. The observation equations are computed by taking the derivative
of the antenna offset correction equation w.r.t. the KBR APC coordinates.


\keepXColumns
\begin{tabularx}{\textwidth}{N T A}
\hline
Name & Type & Annotation\\
\hline
\hfuzz=500pt\includegraphics[width=1em]{element.pdf}~estimate1X & \hfuzz=500pt boolean & \hfuzz=500pt along (satellite 1)\\
\hfuzz=500pt\includegraphics[width=1em]{element.pdf}~estimate1Y & \hfuzz=500pt boolean & \hfuzz=500pt cross (satellite 1)\\
\hfuzz=500pt\includegraphics[width=1em]{element.pdf}~estimate1Z & \hfuzz=500pt boolean & \hfuzz=500pt nadir (satellite 1)\\
\hfuzz=500pt\includegraphics[width=1em]{element.pdf}~estimate2X & \hfuzz=500pt boolean & \hfuzz=500pt along (satellite 2)\\
\hfuzz=500pt\includegraphics[width=1em]{element.pdf}~estimate2Y & \hfuzz=500pt boolean & \hfuzz=500pt cross (satellite 2)\\
\hfuzz=500pt\includegraphics[width=1em]{element.pdf}~estimate2Z & \hfuzz=500pt boolean & \hfuzz=500pt nadir (satellite 2)\\
\hfuzz=500pt\includegraphics[width=1em]{element.pdf}~interpolationDegree & \hfuzz=500pt uint & \hfuzz=500pt differentiation by polynomial approximation of degree n\\
\hline
\end{tabularx}


\subsection{Bias}\label{parametrizationSatelliteTrackingType:bias}
Estimate bias for SST observations. The temporal variation is defined by \configClass{parametrizationTemporal}{parametrizationTemporalType}.


\keepXColumns
\begin{tabularx}{\textwidth}{N T A}
\hline
Name & Type & Annotation\\
\hline
\hfuzz=500pt\includegraphics[width=1em]{element-mustset-unbounded.pdf}~temporal & \hfuzz=500pt \hyperref[parametrizationTemporalType]{parametrizationTemporal} & \hfuzz=500pt \\
\hfuzz=500pt\includegraphics[width=1em]{element.pdf}~perArc & \hfuzz=500pt boolean & \hfuzz=500pt \\
\hline
\end{tabularx}


\subsection{Scale}\label{parametrizationSatelliteTrackingType:scale}
Estimate scale factor for SST observations with respect to reference SST data \configFile{inputfileSatelliteTracking}{instrument}. The temporal variation is defined by \configClass{parametrizationTemporal}{parametrizationTemporalType}.


\keepXColumns
\begin{tabularx}{\textwidth}{N T A}
\hline
Name & Type & Annotation\\
\hline
\hfuzz=500pt\includegraphics[width=1em]{element-mustset.pdf}~inputfileSatelliteTracking & \hfuzz=500pt filename & \hfuzz=500pt \\
\hfuzz=500pt\includegraphics[width=1em]{element-mustset-unbounded.pdf}~temporal & \hfuzz=500pt \hyperref[parametrizationTemporalType]{parametrizationTemporal} & \hfuzz=500pt \\
\hfuzz=500pt\includegraphics[width=1em]{element.pdf}~perArc & \hfuzz=500pt boolean & \hfuzz=500pt \\
\hline
\end{tabularx}


\subsection{TimeBias}\label{parametrizationSatelliteTrackingType:timeBias}
Estimate time shift in seconds in SST observations, with defined temporal variation by \configClass{parametrizationTemporal}{parametrizationTemporalType}. The design matrix is computed by taking the derivative of the ranging data w.r.t. time.


\keepXColumns
\begin{tabularx}{\textwidth}{N T A}
\hline
Name & Type & Annotation\\
\hline
\hfuzz=500pt\includegraphics[width=1em]{element-mustset.pdf}~polynomialDegree & \hfuzz=500pt uint & \hfuzz=500pt polynomial degree\\
\hfuzz=500pt\includegraphics[width=1em]{element-mustset-unbounded.pdf}~temporal & \hfuzz=500pt \hyperref[parametrizationTemporalType]{parametrizationTemporal} & \hfuzz=500pt \\
\hfuzz=500pt\includegraphics[width=1em]{element.pdf}~perArc & \hfuzz=500pt boolean & \hfuzz=500pt \\
\hline
\end{tabularx}


\subsection{ScaleModel}\label{parametrizationSatelliteTrackingType:scaleModel}
Estimate scale factors for deterministic signal models from satellite tracking instrument file \configFile{inputfileSatelliteTracking}{instrument}, see \program{EnsembleAveragingScaleModel}.
Amplitude variation of model waveforms is defined by \configClass{parametrizationTemporal}{parametrizationTemporalType}.


\keepXColumns
\begin{tabularx}{\textwidth}{N T A}
\hline
Name & Type & Annotation\\
\hline
\hfuzz=500pt\includegraphics[width=1em]{element-mustset.pdf}~inputfileSatelliteTracking & \hfuzz=500pt filename & \hfuzz=500pt \\
\hfuzz=500pt\includegraphics[width=1em]{element-mustset-unbounded.pdf}~temporal & \hfuzz=500pt \hyperref[parametrizationTemporalType]{parametrizationTemporal} & \hfuzz=500pt \\
\hfuzz=500pt\includegraphics[width=1em]{element.pdf}~perArc & \hfuzz=500pt boolean & \hfuzz=500pt \\
\hline
\end{tabularx}


\subsection{SpecialEffect}\label{parametrizationSatelliteTrackingType:specialEffect}
Estimate deterministic signals in the GRACE K-Band measurements caused by Sun intrusions
into the star camera baffles of GRACE-A and eclipse transits of the satellites.
These events can be time-indexed beforehand using satellite position and orientation,
see \program{GraceSstSpecialEvents}. The shape of this short-period waveform is nearly
constant within one month and can be approximated by a polynomial.
The amplitude variation of the waveform can also be taken into account
by \configClass{parametrizationTemporal}{parametrizationTemporalType}.


\keepXColumns
\begin{tabularx}{\textwidth}{N T A}
\hline
Name & Type & Annotation\\
\hline
\hfuzz=500pt\includegraphics[width=1em]{element-mustset.pdf}~inputfileEvents & \hfuzz=500pt filename & \hfuzz=500pt instrument with GRACE events\\
\hfuzz=500pt\includegraphics[width=1em]{element-mustset.pdf}~type & \hfuzz=500pt choice & \hfuzz=500pt \\
\hfuzz=500pt\includegraphics[width=1em]{connector.pdf}\includegraphics[width=1em]{element-mustset.pdf}~eclipse1 & \hfuzz=500pt  & \hfuzz=500pt \\
\hfuzz=500pt\includegraphics[width=1em]{connector.pdf}\includegraphics[width=1em]{element-mustset.pdf}~eclipse2 & \hfuzz=500pt  & \hfuzz=500pt \\
\hfuzz=500pt\includegraphics[width=1em]{connector.pdf}\includegraphics[width=1em]{element-mustset.pdf}~starCameraBox1 & \hfuzz=500pt  & \hfuzz=500pt \\
\hfuzz=500pt\includegraphics[width=1em]{connector.pdf}\includegraphics[width=1em]{element-mustset.pdf}~starCameraBox2 & \hfuzz=500pt  & \hfuzz=500pt \\
\hfuzz=500pt\includegraphics[width=1em]{connector.pdf}\includegraphics[width=1em]{element-mustset.pdf}~starCameraBox3 & \hfuzz=500pt  & \hfuzz=500pt \\
\hfuzz=500pt\includegraphics[width=1em]{connector.pdf}\includegraphics[width=1em]{element-mustset.pdf}~starCameraBox4 & \hfuzz=500pt  & \hfuzz=500pt \\
\hfuzz=500pt\includegraphics[width=1em]{connector.pdf}\includegraphics[width=1em]{element-mustset.pdf}~starCameraBox5 & \hfuzz=500pt  & \hfuzz=500pt \\
\hfuzz=500pt\includegraphics[width=1em]{connector.pdf}\includegraphics[width=1em]{element-mustset.pdf}~starCameraBox6 & \hfuzz=500pt  & \hfuzz=500pt \\
\hfuzz=500pt\includegraphics[width=1em]{element-mustset.pdf}~marginLeft & \hfuzz=500pt double & \hfuzz=500pt margin size (on both sides) [seconds]\\
\hfuzz=500pt\includegraphics[width=1em]{element-mustset.pdf}~marginRight & \hfuzz=500pt double & \hfuzz=500pt margin size (on both sides) [seconds]\\
\hfuzz=500pt\includegraphics[width=1em]{element.pdf}~minNumberOfEvents & \hfuzz=500pt uint & \hfuzz=500pt min. number of events to setup parameters\\
\hfuzz=500pt\includegraphics[width=1em]{element-mustset.pdf}~polynomialDegree & \hfuzz=500pt uint & \hfuzz=500pt polynomial degree\\
\hfuzz=500pt\includegraphics[width=1em]{element-mustset-unbounded.pdf}~temporal & \hfuzz=500pt \hyperref[parametrizationTemporalType]{parametrizationTemporal} & \hfuzz=500pt \\
\hline
\end{tabularx}

\clearpage
%==================================

\section{ParametrizationTemporal}\label{parametrizationTemporalType}
This class gives a parametrization of time depending parameters (gravity field, positions, ...).
It will be used to set up the design matrix in a least squares adjustment.
If multiple parametrizations are given the coefficents in the parameter vector
are sequently appended.


\subsection{Constant}\label{parametrizationTemporalType:constant}
Represents a parameter being constant in time in each \config{interval}.


\keepXColumns
\begin{tabularx}{\textwidth}{N T A}
\hline
Name & Type & Annotation\\
\hline
\hfuzz=500pt\includegraphics[width=1em]{element-unbounded.pdf}~interval & \hfuzz=500pt \hyperref[timeSeriesType]{timeSeries} & \hfuzz=500pt \\
\hline
\end{tabularx}


\subsection{Trend}\label{parametrizationTemporalType:trend}
A time variable function is given by a linear trend
\begin{equation}
f(x,t) = \frac{1}{T}(t-t_0) \cdot f_t(x),
\end{equation}
with $t_0$ is \config{timeStart} and $T$ is \config{timeStep} in days.
A constant term is not included and must added separately.


\keepXColumns
\begin{tabularx}{\textwidth}{N T A}
\hline
Name & Type & Annotation\\
\hline
\hfuzz=500pt\includegraphics[width=1em]{element-mustset.pdf}~timeStart & \hfuzz=500pt time & \hfuzz=500pt reference time\\
\hfuzz=500pt\includegraphics[width=1em]{element-mustset.pdf}~timeStep & \hfuzz=500pt time & \hfuzz=500pt \\
\hline
\end{tabularx}


\subsection{Splines}\label{parametrizationTemporalType:splines}
A time variable function is given by
\begin{equation}
f(x,t) =  \sum_i f_i(x)\Psi_i(t),
\end{equation}
with the (spatial) coefficients $f_i(x)$ as parameters and the temporal basis functions~$\Psi_i(t)$.
Basis splines are defined as polynomials of degree~$n$ in intervals between nodal points in time $t_i$,
for details see~\reference{basis splines}{fundamentals.basisSplines}.

The parameters are ordered timewise. First all parameters of $f_{i=1}(x)$ then
$f_{i=2}(x)$ and so on. The total parameter count in each \config{interval} is $N=N_t+d-1$,
where $N_t$ is the count of time points from \configClass{timeSeries}{timeSeriesType} in each interval and $d$
is the \config{degree}.


\keepXColumns
\begin{tabularx}{\textwidth}{N T A}
\hline
Name & Type & Annotation\\
\hline
\hfuzz=500pt\includegraphics[width=1em]{element-mustset.pdf}~degree & \hfuzz=500pt uint & \hfuzz=500pt degree of splines\\
\hfuzz=500pt\includegraphics[width=1em]{element-mustset-unbounded.pdf}~timeSeries & \hfuzz=500pt \hyperref[timeSeriesType]{timeSeries} & \hfuzz=500pt nodal points in time domain\\
\hfuzz=500pt\includegraphics[width=1em]{element-unbounded.pdf}~intervals & \hfuzz=500pt \hyperref[timeSeriesType]{timeSeries} & \hfuzz=500pt \\
\hline
\end{tabularx}


\subsection{Polynomial}
A time variable function is represented by Legendre polynomials in each \config{interval}.
The time is normed to $[-1,1)$ in each interval.

The total parameter count is $(N+1)M$,
where $N$ is the polynmial degree and $M$ the number of intervals.


\keepXColumns
\begin{tabularx}{\textwidth}{N T A}
\hline
Name & Type & Annotation\\
\hline
\hfuzz=500pt\includegraphics[width=1em]{element-mustset.pdf}~polynomialDegree & \hfuzz=500pt uint & \hfuzz=500pt polynomial degree\\
\hfuzz=500pt\includegraphics[width=1em]{element-unbounded.pdf}~interval & \hfuzz=500pt \hyperref[timeSeriesType]{timeSeries} & \hfuzz=500pt intervals of polynomials\\
\hline
\end{tabularx}


\subsection{Oscillation}
A time variable function is given by a oscillation
\begin{equation}
f(x,t) = f^c(\M x)\cos(\omega_i(t)) + f^s(\M x)\sin(\omega_i(t))
\end{equation}
with $\omega_i=\frac{2\pi}{T_i}(t-t_0)$,
$t_0$ is \config{timeStart} and $T$ is \config{timePeriod} in days.


\keepXColumns
\begin{tabularx}{\textwidth}{N T A}
\hline
Name & Type & Annotation\\
\hline
\hfuzz=500pt\includegraphics[width=1em]{element-mustset-unbounded.pdf}~period & \hfuzz=500pt time & \hfuzz=500pt [day]\\
\hfuzz=500pt\includegraphics[width=1em]{element-mustset.pdf}~time0 & \hfuzz=500pt time & \hfuzz=500pt reference time\\
\hline
\end{tabularx}


\subsection{Fourier}
A time variable function is given by a fourier expansion
\begin{equation}
f(x,t) = \sum_{m=1}^M f_m^c(\M x)\cos(2\pi m \tau) + f_m^s(\M x)\sin(2\pi m \tau)
\end{equation}
with the normalized time
\begin{equation}
\tau = \frac{t-t_A}{t_B-t_A},
\end{equation}
and $t_A$ is \config{timeStart}, $t_B$ is \config{timeEnd} in each \config{interval}
and $M$ is the \config{fourierDegree}.

The total parameter count is $2MN$, where $N$ is the number of intervals.
The parameters are sorted in following order: $f_1^c, f_1^s, f_2^c, \ldots$.


\keepXColumns
\begin{tabularx}{\textwidth}{N T A}
\hline
Name & Type & Annotation\\
\hline
\hfuzz=500pt\includegraphics[width=1em]{element-mustset.pdf}~fourierDegree & \hfuzz=500pt uint & \hfuzz=500pt \\
\hfuzz=500pt\includegraphics[width=1em]{element-unbounded.pdf}~interval & \hfuzz=500pt \hyperref[timeSeriesType]{timeSeries} & \hfuzz=500pt intervals of fourier series\\
\hline
\end{tabularx}


\subsection{DoodsonHarmonic}
The time variable function is given by a fourier expansion
\begin{equation}
  f(x,t) = \sum_{i} f_i^c(x)\cos(\Theta_i(t)) + f_i^s(x)\sin(\Theta_i(t)),
\end{equation}
where $\Theta_i(t)$ are the arguments of the tide constituents $i$
\begin{equation}
  \Theta_i(t) = \sum_{k=1}^6 n_i^k\beta_k(t),
\end{equation}
where $\beta_k(t)$ are the Doodson's fundamental arguments ($\tau,s,h,p,N',p_s$) and $n_i^k$
are the Doodson multipliers for the term at frequency~$i$.
The multipliers must be given by \configClass{doodson}{doodson} coded as Doodson number
(e.g. 255.555) or as names intoduced by Darwin (e.g. M2).

The major constituents given by \configClass{doodson}{doodson} can be used to interpolate minor tidal constituents
using the file \configFile{inputfileAdmittance}{admittance}. This file can be created with
\program{DoodsonHarmonicsCalculateAdmittance}.

The total parameter count is $2N$ with $N$ the number of doodson frequencies.
The parameters are sorted in following order: $f_1^c, f_1^s, f_2^c, \ldots$.


\keepXColumns
\begin{tabularx}{\textwidth}{N T A}
\hline
Name & Type & Annotation\\
\hline
\hfuzz=500pt\includegraphics[width=1em]{element-mustset-unbounded.pdf}~doodson & \hfuzz=500pt \hyperref[doodson]{doodson} & \hfuzz=500pt code number (e.g. 255.555) or darwin name (e.g. M2)\\
\hfuzz=500pt\includegraphics[width=1em]{element.pdf}~inputfileAdmittance & \hfuzz=500pt filename & \hfuzz=500pt interpolation of minor constituents\\
\hline
\end{tabularx}

\clearpage
%==================================

\section{Planet}\label{planetType}
Defines the planet to compute the \configClass{ephemeris}{ephemeridesType}.


\keepXColumns
\begin{tabularx}{\textwidth}{N T A}
\hline
Name & Type & Annotation\\
\hline
\hfuzz=500pt\includegraphics[width=1em]{element-mustset.pdf}~planetType & \hfuzz=500pt choice & \hfuzz=500pt planet\\
\hfuzz=500pt\includegraphics[width=1em]{connector.pdf}\includegraphics[width=1em]{element-mustset.pdf}~earth & \hfuzz=500pt  & \hfuzz=500pt \\
\hfuzz=500pt\includegraphics[width=1em]{connector.pdf}\includegraphics[width=1em]{element-mustset.pdf}~sun & \hfuzz=500pt  & \hfuzz=500pt \\
\hfuzz=500pt\includegraphics[width=1em]{connector.pdf}\includegraphics[width=1em]{element-mustset.pdf}~moon & \hfuzz=500pt  & \hfuzz=500pt \\
\hfuzz=500pt\includegraphics[width=1em]{connector.pdf}\includegraphics[width=1em]{element-mustset.pdf}~mercury & \hfuzz=500pt  & \hfuzz=500pt \\
\hfuzz=500pt\includegraphics[width=1em]{connector.pdf}\includegraphics[width=1em]{element-mustset.pdf}~venus & \hfuzz=500pt  & \hfuzz=500pt \\
\hfuzz=500pt\includegraphics[width=1em]{connector.pdf}\includegraphics[width=1em]{element-mustset.pdf}~mars & \hfuzz=500pt  & \hfuzz=500pt \\
\hfuzz=500pt\includegraphics[width=1em]{connector.pdf}\includegraphics[width=1em]{element-mustset.pdf}~jupiter & \hfuzz=500pt  & \hfuzz=500pt \\
\hfuzz=500pt\includegraphics[width=1em]{connector.pdf}\includegraphics[width=1em]{element-mustset.pdf}~saturn & \hfuzz=500pt  & \hfuzz=500pt \\
\hfuzz=500pt\includegraphics[width=1em]{connector.pdf}\includegraphics[width=1em]{element-mustset.pdf}~uranus & \hfuzz=500pt  & \hfuzz=500pt \\
\hfuzz=500pt\includegraphics[width=1em]{connector.pdf}\includegraphics[width=1em]{element-mustset.pdf}~neptune & \hfuzz=500pt  & \hfuzz=500pt \\
\hfuzz=500pt\includegraphics[width=1em]{connector.pdf}\includegraphics[width=1em]{element-mustset.pdf}~pluto & \hfuzz=500pt  & \hfuzz=500pt \\
\hfuzz=500pt\includegraphics[width=1em]{connector.pdf}\includegraphics[width=1em]{element-mustset.pdf}~solarBaryCenter & \hfuzz=500pt  & \hfuzz=500pt \\
\hfuzz=500pt\includegraphics[width=1em]{connector.pdf}\includegraphics[width=1em]{element-mustset.pdf}~earthMoonBaryCenter & \hfuzz=500pt  & \hfuzz=500pt \\
\hline
\end{tabularx}

\clearpage
%==================================

\section{PlotAxis}\label{plotAxisType}
Defines the style of the axes of \program{PlotGraph}.


\subsection{Standard}
General axis for arbitrary input data.


\keepXColumns
\begin{tabularx}{\textwidth}{N T A}
\hline
Name & Type & Annotation\\
\hline
\hfuzz=500pt\includegraphics[width=1em]{element.pdf}~min & \hfuzz=500pt double & \hfuzz=500pt The minimum value of the axis. If no value is given, the minimum scale value is set automatically.\\
\hfuzz=500pt\includegraphics[width=1em]{element.pdf}~max & \hfuzz=500pt double & \hfuzz=500pt The maximum value of the axis. If no value is given, the maximum scale value is set automatically.\\
\hfuzz=500pt\includegraphics[width=1em]{element.pdf}~majorTickSpacing & \hfuzz=500pt double & \hfuzz=500pt The boundary annotation.\\
\hfuzz=500pt\includegraphics[width=1em]{element.pdf}~minorTickSpacing & \hfuzz=500pt double & \hfuzz=500pt The spacing of the frame tick intervals.\\
\hfuzz=500pt\includegraphics[width=1em]{element.pdf}~gridLineSpacing & \hfuzz=500pt double & \hfuzz=500pt The spacing of the grid line intervals\\
\hfuzz=500pt\includegraphics[width=1em]{element.pdf}~gridLine & \hfuzz=500pt \hyperref[plotLineType]{plotLine} & \hfuzz=500pt The style of the grid lines.\\
\hfuzz=500pt\includegraphics[width=1em]{element.pdf}~unit & \hfuzz=500pt string & \hfuzz=500pt Naming unit to append to the axis values.\\
\hfuzz=500pt\includegraphics[width=1em]{element.pdf}~label & \hfuzz=500pt string & \hfuzz=500pt The description of the axis.\\
\hfuzz=500pt\includegraphics[width=1em]{element.pdf}~logarithmic & \hfuzz=500pt boolean & \hfuzz=500pt If set to 'yes', a logarithmic scale is used for the axis.\\
\hfuzz=500pt\includegraphics[width=1em]{element-mustset.pdf}~color & \hfuzz=500pt \hyperref[plotColorType]{plotColor} & \hfuzz=500pt Setting the color of the axis bars and labels.\\
\hfuzz=500pt\includegraphics[width=1em]{element.pdf}~changeDirection & \hfuzz=500pt boolean & \hfuzz=500pt If set to 'yes', the directions right/up are changed to left/down.\\
\hline
\end{tabularx}


\subsection{Time}
The input data are interpreted as MJD (modified Julian date).
The unit of the tick spacings should be appenend to the number and can be any of
\begin{itemize}
\item Y (year, plot with 4 digits)
\item y (year, plot with 2 digits)
\item O (month, plot using \verb|FORMAT_DATE_MAP|)
\item o (month, plot with 2 digits)
\item U (ISO week, plot using \verb|FORMAT_DATE_MAP|)
\item u (ISO week, plot using 2 digits)
\item r (Gregorian week, 7-day stride from start of week \verb|TIME_WEEK_START|)
\item K (ISO weekday, plot name of day)
\item D (date, plot using \verb|FORMAT_DATE_MAP|)
\item d (day, plot day of month 0-31 or year 1-366, via \verb|FORMAT_DATE_MAP|)
\item R (day, same as d, aligned with \verb|TIME_WEEK_START|)
\item H (hour, plot using \verb|FORMAT_CLOCK_MAP|)
\item h (hour, plot with 2 digits)
\item M (minute, plot using \verb|FORMAT_CLOCK_MAP|)
\item m (minute, plot with 2 digits)
\item S (second, plot using \verb|FORMAT_CLOCK_MAP|)
\item s (second, plot with 2 digits).
\end{itemize}

A secondary time axis can be added to specify larger intervals (e.g dates of hourly data).

Examples: Settings for Fig.~\ref{plotAxisType:plotAxisTime1}: \config{majorTickSpacing}=\verb|6H|, secondary: \config{majorTickSpacing}=\verb|1D|.
\fig{!hb}{1.0}{plotAxisTime1}{plotAxisType:plotAxisTime1}{Time axis for daily data.}

Settings for Fig.~\ref{plotAxisType:plotAxisTime2}: \config{majorTickSpacing}=\verb|2d|, secondary: \config{majorTickSpacing}=\verb|1O|, \config{options}=\verb|FORMAT_DATE_MAP="o yyyy"|.
\fig{!hb}{1.0}{plotAxisTime2}{plotAxisType:plotAxisTime2}{Time axis for monthly data.}

Settings for Fig.~\ref{plotAxisType:plotAxisTime3}: \config{majorTickSpacing}=\verb|1o|, secondary: \config{majorTickSpacing}=\verb|1Y|, \config{options}=\verb|FORMAT_DATE_MAP="mm"|.
\fig{!hb}{1.0}{plotAxisTime3}{plotAxisType:plotAxisTime3}{Time axis for yearly data.}



\keepXColumns
\begin{tabularx}{\textwidth}{N T A}
\hline
Name & Type & Annotation\\
\hline
\hfuzz=500pt\includegraphics[width=1em]{element.pdf}~min & \hfuzz=500pt time & \hfuzz=500pt The minimum value of the time axis. If no value is given, the minimum scale value is set automatically.\\
\hfuzz=500pt\includegraphics[width=1em]{element.pdf}~max & \hfuzz=500pt time & \hfuzz=500pt The maximum value of the time axis. If no value is given, the maximum scale value is set automatically.\\
\hfuzz=500pt\includegraphics[width=1em]{element.pdf}~majorTickSpacing & \hfuzz=500pt string & \hfuzz=500pt Y: year, o: month\\
\hfuzz=500pt\includegraphics[width=1em]{element.pdf}~minorTickSpacing & \hfuzz=500pt string & \hfuzz=500pt D: date, d: day\\
\hfuzz=500pt\includegraphics[width=1em]{element.pdf}~gridLineSpacing & \hfuzz=500pt string & \hfuzz=500pt H: clock, h: hour, m: minute, s: second\\
\hfuzz=500pt\includegraphics[width=1em]{element.pdf}~secondary & \hfuzz=500pt sequence & \hfuzz=500pt secondary time axis\\
\hfuzz=500pt\includegraphics[width=1em]{connector.pdf}\includegraphics[width=1em]{element.pdf}~majorTickSpacing & \hfuzz=500pt string & \hfuzz=500pt Y: year, o: month\\
\hfuzz=500pt\includegraphics[width=1em]{connector.pdf}\includegraphics[width=1em]{element.pdf}~minorTickSpacing & \hfuzz=500pt string & \hfuzz=500pt D: date, d: day\\
\hfuzz=500pt\includegraphics[width=1em]{connector.pdf}\includegraphics[width=1em]{element.pdf}~gridLineSpacing & \hfuzz=500pt string & \hfuzz=500pt H: clock, h: hour, m: minute, s: second\\
\hfuzz=500pt\includegraphics[width=1em]{element-mustset.pdf}~color & \hfuzz=500pt \hyperref[plotColorType]{plotColor} & \hfuzz=500pt color of axis bars and labels\\
\hfuzz=500pt\includegraphics[width=1em]{element.pdf}~gridLine & \hfuzz=500pt \hyperref[plotLineType]{plotLine} & \hfuzz=500pt The style of the grid lines.\\
\hfuzz=500pt\includegraphics[width=1em]{element.pdf}~changeDirection & \hfuzz=500pt boolean & \hfuzz=500pt right-\$>\$left / up-\$>\$down\\
\hfuzz=500pt\includegraphics[width=1em]{element-unbounded.pdf}~options & \hfuzz=500pt string & \hfuzz=500pt adjust date format\\
\hline
\end{tabularx}


\subsection{Labeled}
Axis with string labels. The coordinate system is based on the label indices (e.g. 0, 1, 2).


\keepXColumns
\begin{tabularx}{\textwidth}{N T A}
\hline
Name & Type & Annotation\\
\hline
\hfuzz=500pt\includegraphics[width=1em]{element-mustset-unbounded.pdf}~labels & \hfuzz=500pt string & \hfuzz=500pt tick labels (ticks are placed at their index. e.g. 0, 1, ..., 5)\\
\hfuzz=500pt\includegraphics[width=1em]{element.pdf}~min & \hfuzz=500pt expression & \hfuzz=500pt minimum value of the axis\\
\hfuzz=500pt\includegraphics[width=1em]{element.pdf}~max & \hfuzz=500pt expression & \hfuzz=500pt maximum values of the axis\\
\hfuzz=500pt\includegraphics[width=1em]{element.pdf}~orthogonalLabels & \hfuzz=500pt boolean & \hfuzz=500pt labels are oriented orthogonal to axis\\
\hfuzz=500pt\includegraphics[width=1em]{element.pdf}~gridLine & \hfuzz=500pt \hyperref[plotLineType]{plotLine} & \hfuzz=500pt The style of the grid lines.\\
\hfuzz=500pt\includegraphics[width=1em]{element-mustset.pdf}~color & \hfuzz=500pt \hyperref[plotColorType]{plotColor} & \hfuzz=500pt set the color of the axis and labels\\
\hfuzz=500pt\includegraphics[width=1em]{element.pdf}~changeDirection & \hfuzz=500pt boolean & \hfuzz=500pt If set to 'yes', the directions right/up are changed to left/down.\\
\hline
\end{tabularx}

\clearpage
%==================================

\section{PlotColor}\label{plotColorType}
Selects a color.
Used in \program{PlotDegreeAmplitudes}, \program{PlotGraph}, \program{PlotMap},
\program{PlotMatrix}, \program{PlotSphericalHarmonicsTriangle}.



\keepXColumns
\begin{tabularx}{\textwidth}{N T A}
\hline
Name & Type & Annotation\\
\hline
\hfuzz=500pt\includegraphics[width=1em]{element-mustset.pdf}~plotColorType & \hfuzz=500pt choice & \hfuzz=500pt color\\
\hfuzz=500pt\includegraphics[width=1em]{connector.pdf}\includegraphics[width=1em]{element-mustset.pdf}~black & \hfuzz=500pt  & \hfuzz=500pt \\
\hfuzz=500pt\includegraphics[width=1em]{connector.pdf}\includegraphics[width=1em]{element-mustset.pdf}~red & \hfuzz=500pt  & \hfuzz=500pt \\
\hfuzz=500pt\includegraphics[width=1em]{connector.pdf}\includegraphics[width=1em]{element-mustset.pdf}~blue & \hfuzz=500pt  & \hfuzz=500pt \\
\hfuzz=500pt\includegraphics[width=1em]{connector.pdf}\includegraphics[width=1em]{element-mustset.pdf}~green & \hfuzz=500pt  & \hfuzz=500pt \\
\hfuzz=500pt\includegraphics[width=1em]{connector.pdf}\includegraphics[width=1em]{element-mustset.pdf}~orange & \hfuzz=500pt  & \hfuzz=500pt \\
\hfuzz=500pt\includegraphics[width=1em]{connector.pdf}\includegraphics[width=1em]{element-mustset.pdf}~darkred & \hfuzz=500pt  & \hfuzz=500pt \\
\hfuzz=500pt\includegraphics[width=1em]{connector.pdf}\includegraphics[width=1em]{element-mustset.pdf}~yellow & \hfuzz=500pt  & \hfuzz=500pt \\
\hfuzz=500pt\includegraphics[width=1em]{connector.pdf}\includegraphics[width=1em]{element-mustset.pdf}~lightgreen & \hfuzz=500pt  & \hfuzz=500pt \\
\hfuzz=500pt\includegraphics[width=1em]{connector.pdf}\includegraphics[width=1em]{element-mustset.pdf}~gray & \hfuzz=500pt  & \hfuzz=500pt \\
\hfuzz=500pt\includegraphics[width=1em]{connector.pdf}\includegraphics[width=1em]{element-mustset.pdf}~rgb & \hfuzz=500pt sequence & \hfuzz=500pt \\
\hfuzz=500pt\quad\includegraphics[width=1em]{connector.pdf}\includegraphics[width=1em]{element-mustset.pdf}~red & \hfuzz=500pt uint & \hfuzz=500pt 0..255\\
\hfuzz=500pt\quad\includegraphics[width=1em]{connector.pdf}\includegraphics[width=1em]{element-mustset.pdf}~green & \hfuzz=500pt uint & \hfuzz=500pt 0..255\\
\hfuzz=500pt\quad\includegraphics[width=1em]{connector.pdf}\includegraphics[width=1em]{element-mustset.pdf}~blue & \hfuzz=500pt uint & \hfuzz=500pt 0..255\\
\hfuzz=500pt\includegraphics[width=1em]{connector.pdf}\includegraphics[width=1em]{element-mustset.pdf}~grayscale & \hfuzz=500pt sequence & \hfuzz=500pt \\
\hfuzz=500pt\quad\includegraphics[width=1em]{connector.pdf}\includegraphics[width=1em]{element-mustset.pdf}~value & \hfuzz=500pt uint & \hfuzz=500pt 0..255\\
\hfuzz=500pt\includegraphics[width=1em]{connector.pdf}\includegraphics[width=1em]{element-mustset.pdf}~namedColor & \hfuzz=500pt sequence & \hfuzz=500pt \\
\hfuzz=500pt\quad\includegraphics[width=1em]{connector.pdf}\includegraphics[width=1em]{element-mustset.pdf}~colorName & \hfuzz=500pt string & \hfuzz=500pt name after GMT definition\\
\hfuzz=500pt\includegraphics[width=1em]{connector.pdf}\includegraphics[width=1em]{element-mustset.pdf}~cycler & \hfuzz=500pt sequence & \hfuzz=500pt \\
\hfuzz=500pt\quad\includegraphics[width=1em]{connector.pdf}\includegraphics[width=1em]{element-mustset.pdf}~index & \hfuzz=500pt uint & \hfuzz=500pt pick color based on index expression\\
\hfuzz=500pt\quad\includegraphics[width=1em]{connector.pdf}\includegraphics[width=1em]{element.pdf}~inputfileColorList & \hfuzz=500pt filename & \hfuzz=500pt list of colors as defined by GMT\\
\hline
\end{tabularx}

\clearpage
%==================================

\section{PlotColorbar}\label{plotColorbarType}
A colorbar as used in \program{PlotMap}, \program{PlotMatrix}, \program{PlotSphericalHarmonicsTriangle}.


\keepXColumns
\begin{tabularx}{\textwidth}{N T A}
\hline
Name & Type & Annotation\\
\hline
\hfuzz=500pt\includegraphics[width=1em]{element-mustset.pdf}~plotColorbarType & \hfuzz=500pt sequence & \hfuzz=500pt \\
\hfuzz=500pt\includegraphics[width=1em]{connector.pdf}\includegraphics[width=1em]{element.pdf}~min & \hfuzz=500pt double & \hfuzz=500pt \\
\hfuzz=500pt\includegraphics[width=1em]{connector.pdf}\includegraphics[width=1em]{element.pdf}~max & \hfuzz=500pt double & \hfuzz=500pt \\
\hfuzz=500pt\includegraphics[width=1em]{connector.pdf}\includegraphics[width=1em]{element.pdf}~annotation & \hfuzz=500pt double & \hfuzz=500pt boundary annotation\\
\hfuzz=500pt\includegraphics[width=1em]{connector.pdf}\includegraphics[width=1em]{element.pdf}~unit & \hfuzz=500pt string & \hfuzz=500pt appended to axis values\\
\hfuzz=500pt\includegraphics[width=1em]{connector.pdf}\includegraphics[width=1em]{element.pdf}~label & \hfuzz=500pt string & \hfuzz=500pt description of the axis\\
\hfuzz=500pt\includegraphics[width=1em]{connector.pdf}\includegraphics[width=1em]{element.pdf}~logarithmic & \hfuzz=500pt boolean & \hfuzz=500pt use logarithmic scale\\
\hfuzz=500pt\includegraphics[width=1em]{connector.pdf}\includegraphics[width=1em]{element.pdf}~triangleLeft & \hfuzz=500pt boolean & \hfuzz=500pt \\
\hfuzz=500pt\includegraphics[width=1em]{connector.pdf}\includegraphics[width=1em]{element.pdf}~triangleRight & \hfuzz=500pt boolean & \hfuzz=500pt \\
\hfuzz=500pt\includegraphics[width=1em]{connector.pdf}\includegraphics[width=1em]{element.pdf}~illuminate & \hfuzz=500pt boolean & \hfuzz=500pt illuminate\\
\hfuzz=500pt\includegraphics[width=1em]{connector.pdf}\includegraphics[width=1em]{element.pdf}~vertical & \hfuzz=500pt boolean & \hfuzz=500pt plot vertical color bar on the right\\
\hfuzz=500pt\includegraphics[width=1em]{connector.pdf}\includegraphics[width=1em]{element.pdf}~length & \hfuzz=500pt double & \hfuzz=500pt length of colorbar in percent\\
\hfuzz=500pt\includegraphics[width=1em]{connector.pdf}\includegraphics[width=1em]{element.pdf}~margin & \hfuzz=500pt double & \hfuzz=500pt between colorbar and figure [cm]\\
\hfuzz=500pt\includegraphics[width=1em]{connector.pdf}\includegraphics[width=1em]{element.pdf}~colorTable & \hfuzz=500pt string & \hfuzz=500pt name of the color bar\\
\hfuzz=500pt\includegraphics[width=1em]{connector.pdf}\includegraphics[width=1em]{element.pdf}~reverse & \hfuzz=500pt boolean & \hfuzz=500pt reverse direction\\
\hfuzz=500pt\includegraphics[width=1em]{connector.pdf}\includegraphics[width=1em]{element.pdf}~showColorbar & \hfuzz=500pt boolean & \hfuzz=500pt \\
\hline
\end{tabularx}

\clearpage
%==================================

\section{PlotGraphLayer}\label{plotGraphLayerType}
Defines the content of an xy-plot of \program{PlotGraph}.
Multiple layers are are plotted sequentially. With \config{plotOnSecondAxis}
the alternative y-axis on the right hand side can be selected if provided.


\subsection{LinesAndPoints}\label{plotGraphLayerType:linesAndPoints}
Draws a \configClass{line}{plotLineType} and/or points (\configClass{symbol}{plotSymbolType})
of xy data. The standard \reference{dataVariables}{general.parser:dataVariables}
are available to select the data columns of \configFile{inputfileMatrix}{matrix}.
If no \configClass{color}{plotColorType} of the \configClass{symbol}{plotSymbolType}
is given a \configClass{colorbar}{plotColorbarType}
is required and the color is determined by \config{valueZ}.
Additionally a vertical error bar can be plotted at each data point with
size \config{valueErrorBar}.

See \program{Gravityfield2AreaMeanTimeSeries} for an example plot.


\keepXColumns
\begin{tabularx}{\textwidth}{N T A}
\hline
Name & Type & Annotation\\
\hline
\hfuzz=500pt\includegraphics[width=1em]{element-mustset.pdf}~inputfileMatrix & \hfuzz=500pt filename & \hfuzz=500pt each line contains x,y\\
\hfuzz=500pt\includegraphics[width=1em]{element.pdf}~valueX & \hfuzz=500pt expression & \hfuzz=500pt expression for x-values (input columns are named data0, data1, ...)\\
\hfuzz=500pt\includegraphics[width=1em]{element-mustset.pdf}~valueY & \hfuzz=500pt expression & \hfuzz=500pt expression for y-values (input columns are named data0, data1, ...)\\
\hfuzz=500pt\includegraphics[width=1em]{element.pdf}~valueZ & \hfuzz=500pt expression & \hfuzz=500pt expression for the colorbar\\
\hfuzz=500pt\includegraphics[width=1em]{element.pdf}~valueErrorBar & \hfuzz=500pt expression & \hfuzz=500pt expression for error bars (input columns are named data0, data1, ...)\\
\hfuzz=500pt\includegraphics[width=1em]{element.pdf}~description & \hfuzz=500pt string & \hfuzz=500pt text of the legend\\
\hfuzz=500pt\includegraphics[width=1em]{element.pdf}~line & \hfuzz=500pt \hyperref[plotLineType]{plotLine} & \hfuzz=500pt \\
\hfuzz=500pt\includegraphics[width=1em]{element.pdf}~symbol & \hfuzz=500pt \hyperref[plotSymbolType]{plotSymbol} & \hfuzz=500pt \\
\hfuzz=500pt\includegraphics[width=1em]{element.pdf}~plotOnSecondAxis & \hfuzz=500pt boolean & \hfuzz=500pt draw dataset on a second Y-axis (if available).\\
\hline
\end{tabularx}


\subsection{ErrorEnvelope}
Draws a symmetrical envelope around \config{valueY} as function of \config{valueX}
using deviations \config{valueErrors}.
The standard \reference{dataVariables}{general.parser:dataVariables}
are available to select the data columns of \configFile{inputfileMatrix}{matrix}.
The data line itself is not plotted but must be added as extra
\configClass{layer:linesAndPoints}{plotGraphLayerType:linesAndPoints}.


\keepXColumns
\begin{tabularx}{\textwidth}{N T A}
\hline
Name & Type & Annotation\\
\hline
\hfuzz=500pt\includegraphics[width=1em]{element-mustset.pdf}~inputfileMatrix & \hfuzz=500pt filename & \hfuzz=500pt each line contains x,y\\
\hfuzz=500pt\includegraphics[width=1em]{element.pdf}~valueX & \hfuzz=500pt expression & \hfuzz=500pt expression for x-values (input columns are named data0, data1, ...)\\
\hfuzz=500pt\includegraphics[width=1em]{element-mustset.pdf}~valueY & \hfuzz=500pt expression & \hfuzz=500pt expression for y-values (input columns are named data0, data1, ...)\\
\hfuzz=500pt\includegraphics[width=1em]{element-mustset.pdf}~valueErrors & \hfuzz=500pt expression & \hfuzz=500pt expression for error values\\
\hfuzz=500pt\includegraphics[width=1em]{element.pdf}~description & \hfuzz=500pt string & \hfuzz=500pt text of the legend\\
\hfuzz=500pt\includegraphics[width=1em]{element.pdf}~fillColor & \hfuzz=500pt \hyperref[plotColorType]{plotColor} & \hfuzz=500pt fill color of the envelope\\
\hfuzz=500pt\includegraphics[width=1em]{element.pdf}~edgeLine & \hfuzz=500pt \hyperref[plotLineType]{plotLine} & \hfuzz=500pt edge line style of the envelope\\
\hfuzz=500pt\includegraphics[width=1em]{element.pdf}~plotOnSecondAxis & \hfuzz=500pt boolean & \hfuzz=500pt draw dataset on a second Y-axis (if available).\\
\hline
\end{tabularx}


\subsection{Bars}
Creates a bar plot with vertical or \config{horizontal} bars out of the given
x- and y-values. The standard \reference{dataVariables}{general.parser:dataVariables}
are available to select the data columns of \configFile{inputfileMatrix}{matrix}.
The bars ranges from \config{valueBase} (can be also an expression) to the \config{valueY}.
If no \configClass{color}{plotColorType} is given a \configClass{colorbar}{plotColorbarType}
is required and the color is determined by \config{valueZ}.

See \program{Instrument2Histogram} for an example plot.


\keepXColumns
\begin{tabularx}{\textwidth}{N T A}
\hline
Name & Type & Annotation\\
\hline
\hfuzz=500pt\includegraphics[width=1em]{element-mustset.pdf}~inputfileMatrix & \hfuzz=500pt filename & \hfuzz=500pt each line contains x,y\\
\hfuzz=500pt\includegraphics[width=1em]{element.pdf}~valueX & \hfuzz=500pt expression & \hfuzz=500pt expression for x-values (input columns are named data0, data1, ...)\\
\hfuzz=500pt\includegraphics[width=1em]{element-mustset.pdf}~valueY & \hfuzz=500pt expression & \hfuzz=500pt expression for y-values (input columns are named data0, data1, ...)\\
\hfuzz=500pt\includegraphics[width=1em]{element.pdf}~valueZ & \hfuzz=500pt expression & \hfuzz=500pt expression for the colorbar\\
\hfuzz=500pt\includegraphics[width=1em]{element.pdf}~valueBase & \hfuzz=500pt expression & \hfuzz=500pt base value of bars (default: minimum y-value)\\
\hfuzz=500pt\includegraphics[width=1em]{element.pdf}~width & \hfuzz=500pt expression & \hfuzz=500pt width of bars (default: minimum x-gap)\\
\hfuzz=500pt\includegraphics[width=1em]{element.pdf}~horizontal & \hfuzz=500pt boolean & \hfuzz=500pt draw horizontal bars instead of vertical\\
\hfuzz=500pt\includegraphics[width=1em]{element.pdf}~description & \hfuzz=500pt string & \hfuzz=500pt text of the legend\\
\hfuzz=500pt\includegraphics[width=1em]{element.pdf}~color & \hfuzz=500pt \hyperref[plotColorType]{plotColor} & \hfuzz=500pt \\
\hfuzz=500pt\includegraphics[width=1em]{element.pdf}~edgeLine & \hfuzz=500pt \hyperref[plotLineType]{plotLine} & \hfuzz=500pt line\\
\hfuzz=500pt\includegraphics[width=1em]{element.pdf}~plotOnSecondAxis & \hfuzz=500pt boolean & \hfuzz=500pt draw dataset on a second Y-axis (if available).\\
\hline
\end{tabularx}


\subsection{Gridded}
Creates a regular grid of yxz values. The standard \reference{dataVariables}{general.parser:dataVariables}
are available to select the data columns of \configFile{inputfileMatrix}{matrix}.
Empty grid cells are not plotted. Cells with more than one value will be set to the mean value.
The grid spacing is determined by the median spacing of the input data or set by \config{incrementX/Y}.

See \program{Orbit2ArgumentOfLatitude} for an example plot.


\keepXColumns
\begin{tabularx}{\textwidth}{N T A}
\hline
Name & Type & Annotation\\
\hline
\hfuzz=500pt\includegraphics[width=1em]{element-mustset.pdf}~inputfileMatrix & \hfuzz=500pt filename & \hfuzz=500pt each line contains x,y,z\\
\hfuzz=500pt\includegraphics[width=1em]{element.pdf}~valueX & \hfuzz=500pt expression & \hfuzz=500pt expression for x-values (input columns are named data0, data1, ...)\\
\hfuzz=500pt\includegraphics[width=1em]{element-mustset.pdf}~valueY & \hfuzz=500pt expression & \hfuzz=500pt expression for y-values (input columns are named data0, data1, ...)\\
\hfuzz=500pt\includegraphics[width=1em]{element-mustset.pdf}~valueZ & \hfuzz=500pt expression & \hfuzz=500pt expression for the colorbar\\
\hfuzz=500pt\includegraphics[width=1em]{element.pdf}~incrementX & \hfuzz=500pt double & \hfuzz=500pt the grid spacing\\
\hfuzz=500pt\includegraphics[width=1em]{element.pdf}~incrementY & \hfuzz=500pt double & \hfuzz=500pt the grid spacing\\
\hfuzz=500pt\includegraphics[width=1em]{element.pdf}~plotOnSecondAxis & \hfuzz=500pt boolean & \hfuzz=500pt draw dataset on a second Y-axis (if available).\\
\hline
\end{tabularx}


\subsection{Rectangle}
Plots a rectangle to highlight an area.


\keepXColumns
\begin{tabularx}{\textwidth}{N T A}
\hline
Name & Type & Annotation\\
\hline
\hfuzz=500pt\includegraphics[width=1em]{element.pdf}~minX & \hfuzz=500pt double & \hfuzz=500pt empty: left\\
\hfuzz=500pt\includegraphics[width=1em]{element.pdf}~maxX & \hfuzz=500pt double & \hfuzz=500pt empty: right\\
\hfuzz=500pt\includegraphics[width=1em]{element.pdf}~minY & \hfuzz=500pt double & \hfuzz=500pt empty: bottom\\
\hfuzz=500pt\includegraphics[width=1em]{element.pdf}~maxY & \hfuzz=500pt double & \hfuzz=500pt empty: top\\
\hfuzz=500pt\includegraphics[width=1em]{element.pdf}~description & \hfuzz=500pt string & \hfuzz=500pt text of the legend\\
\hfuzz=500pt\includegraphics[width=1em]{element.pdf}~edgeLine & \hfuzz=500pt \hyperref[plotLineType]{plotLine} & \hfuzz=500pt \\
\hfuzz=500pt\includegraphics[width=1em]{element.pdf}~fillColor & \hfuzz=500pt \hyperref[plotColorType]{plotColor} & \hfuzz=500pt \\
\hfuzz=500pt\includegraphics[width=1em]{element.pdf}~plotOnSecondAxis & \hfuzz=500pt boolean & \hfuzz=500pt draw dataset on a second Y-axis (if available).\\
\hline
\end{tabularx}


\subsection{Text}
Writes a \config{text} at \config{originX} and \config{originY} position in the graph.
With \config{clip} the text is cutted at the boundaries of the plotting area.


\keepXColumns
\begin{tabularx}{\textwidth}{N T A}
\hline
Name & Type & Annotation\\
\hline
\hfuzz=500pt\includegraphics[width=1em]{element-mustset.pdf}~text & \hfuzz=500pt string & \hfuzz=500pt \\
\hfuzz=500pt\includegraphics[width=1em]{element-mustset.pdf}~originX & \hfuzz=500pt double & \hfuzz=500pt \\
\hfuzz=500pt\includegraphics[width=1em]{element-mustset.pdf}~originY & \hfuzz=500pt double & \hfuzz=500pt \\
\hfuzz=500pt\includegraphics[width=1em]{element.pdf}~offsetX & \hfuzz=500pt double & \hfuzz=500pt [cm] x-offset from origin\\
\hfuzz=500pt\includegraphics[width=1em]{element.pdf}~offsetY & \hfuzz=500pt double & \hfuzz=500pt [cm] y-offset from origin\\
\hfuzz=500pt\includegraphics[width=1em]{element.pdf}~alignment & \hfuzz=500pt string & \hfuzz=500pt L, C, R (left, center, right) and T, M, B (top, middle, bottom)\\
\hfuzz=500pt\includegraphics[width=1em]{element.pdf}~fontSize & \hfuzz=500pt double & \hfuzz=500pt [pt]\\
\hfuzz=500pt\includegraphics[width=1em]{element-mustset.pdf}~fontColor & \hfuzz=500pt \hyperref[plotColorType]{plotColor} & \hfuzz=500pt \\
\hfuzz=500pt\includegraphics[width=1em]{element.pdf}~clip & \hfuzz=500pt boolean & \hfuzz=500pt clip at boundaries\\
\hfuzz=500pt\includegraphics[width=1em]{element.pdf}~plotOnSecondAxis & \hfuzz=500pt boolean & \hfuzz=500pt draw dataset on a second Y-axis (if available).\\
\hline
\end{tabularx}


\subsection{DegreeAmplitudes}
Plot degree amplitudes of potential coefficients computed by \program{Gravityfield2DegreeAmplitudes}
or \program{PotentialCoefficients2DegreeAmplitudes}.
The standard \reference{dataVariables}{general.parser:dataVariables} are available to select
the data columns of \configFile{inputfileMatrix}{matrix}. It plots a solid line for the
\config{valueSignal} and a dotted line for the \config{valueError} per default.


\keepXColumns
\begin{tabularx}{\textwidth}{N T A}
\hline
Name & Type & Annotation\\
\hline
\hfuzz=500pt\includegraphics[width=1em]{element-mustset.pdf}~inputfileMatrix & \hfuzz=500pt filename & \hfuzz=500pt degree amplitudes\\
\hfuzz=500pt\includegraphics[width=1em]{element.pdf}~valueDegree & \hfuzz=500pt expression & \hfuzz=500pt expression for x-values (degrees) (input columns are named data0, data1, ...)\\
\hfuzz=500pt\includegraphics[width=1em]{element.pdf}~valueSignal & \hfuzz=500pt expression & \hfuzz=500pt expression for y-values (signal) (input columns are named data0, data1, ...)\\
\hfuzz=500pt\includegraphics[width=1em]{element.pdf}~valueErrors & \hfuzz=500pt expression & \hfuzz=500pt expression for y-values (formal errors)\\
\hfuzz=500pt\includegraphics[width=1em]{element.pdf}~description & \hfuzz=500pt string & \hfuzz=500pt text of the legend\\
\hfuzz=500pt\includegraphics[width=1em]{element.pdf}~lineSignal & \hfuzz=500pt \hyperref[plotLineType]{plotLine} & \hfuzz=500pt \\
\hfuzz=500pt\includegraphics[width=1em]{element.pdf}~lineErrors & \hfuzz=500pt \hyperref[plotLineType]{plotLine} & \hfuzz=500pt \\
\hfuzz=500pt\includegraphics[width=1em]{element.pdf}~plotOnSecondAxis & \hfuzz=500pt boolean & \hfuzz=500pt draw dataset on a second Y-axis (if available).\\
\hline
\end{tabularx}

\clearpage
%==================================

\section{PlotLegend}\label{plotLegendType}
Plot a legend of the descriptions provided in
\configClass{plotGraphLayer}{plotGraphLayerType} in \program{PlotGraph}.


\keepXColumns
\begin{tabularx}{\textwidth}{N T A}
\hline
Name & Type & Annotation\\
\hline
\hfuzz=500pt\includegraphics[width=1em]{element-mustset.pdf}~plotLegendType & \hfuzz=500pt sequence & \hfuzz=500pt \\
\hfuzz=500pt\includegraphics[width=1em]{connector.pdf}\includegraphics[width=1em]{element.pdf}~width & \hfuzz=500pt double & \hfuzz=500pt legend width [cm]\\
\hfuzz=500pt\includegraphics[width=1em]{connector.pdf}\includegraphics[width=1em]{element.pdf}~height & \hfuzz=500pt double & \hfuzz=500pt legend height [cm] (default: estimated)\\
\hfuzz=500pt\includegraphics[width=1em]{connector.pdf}\includegraphics[width=1em]{element.pdf}~positionX & \hfuzz=500pt double & \hfuzz=500pt legend x-position in normalized (0-1) coordinates.\\
\hfuzz=500pt\includegraphics[width=1em]{connector.pdf}\includegraphics[width=1em]{element.pdf}~positionY & \hfuzz=500pt double & \hfuzz=500pt legend y-position in normalized (0-1) coordinates.\\
\hfuzz=500pt\includegraphics[width=1em]{connector.pdf}\includegraphics[width=1em]{element.pdf}~anchorPoint & \hfuzz=500pt string & \hfuzz=500pt Two character combination of L, C, R (for left, center, or right) and T, M, B for top, middle, or bottom. e.g., TL for top left\\
\hfuzz=500pt\includegraphics[width=1em]{connector.pdf}\includegraphics[width=1em]{element.pdf}~columns & \hfuzz=500pt uint & \hfuzz=500pt number of columns in legend\\
\hfuzz=500pt\includegraphics[width=1em]{connector.pdf}\includegraphics[width=1em]{element.pdf}~fillColor & \hfuzz=500pt \hyperref[plotColorType]{plotColor} & \hfuzz=500pt fill color of the legend box\\
\hfuzz=500pt\includegraphics[width=1em]{connector.pdf}\includegraphics[width=1em]{element.pdf}~edgeLine & \hfuzz=500pt \hyperref[plotLineType]{plotLine} & \hfuzz=500pt style of the legend box edge\\
\hline
\end{tabularx}

\clearpage
%==================================

\section{PlotLine}\label{plotLineType}
Defines the line style to be plotted.


\subsection{Solid}
Draws a solid line.


\keepXColumns
\begin{tabularx}{\textwidth}{N T A}
\hline
Name & Type & Annotation\\
\hline
\hfuzz=500pt\includegraphics[width=1em]{element.pdf}~width & \hfuzz=500pt double & \hfuzz=500pt line width [p]\\
\hfuzz=500pt\includegraphics[width=1em]{element-mustset.pdf}~color & \hfuzz=500pt \hyperref[plotColorType]{plotColor} & \hfuzz=500pt \\
\hline
\end{tabularx}


\subsection{Dashed}
Draws a dashed line.


\keepXColumns
\begin{tabularx}{\textwidth}{N T A}
\hline
Name & Type & Annotation\\
\hline
\hfuzz=500pt\includegraphics[width=1em]{element.pdf}~width & \hfuzz=500pt double & \hfuzz=500pt line width [p]\\
\hfuzz=500pt\includegraphics[width=1em]{element-mustset.pdf}~color & \hfuzz=500pt \hyperref[plotColorType]{plotColor} & \hfuzz=500pt \\
\hline
\end{tabularx}


\subsection{Dotted}
Draws a dotted line.


\keepXColumns
\begin{tabularx}{\textwidth}{N T A}
\hline
Name & Type & Annotation\\
\hline
\hfuzz=500pt\includegraphics[width=1em]{element.pdf}~width & \hfuzz=500pt double & \hfuzz=500pt line width [p]\\
\hfuzz=500pt\includegraphics[width=1em]{element-mustset.pdf}~color & \hfuzz=500pt \hyperref[plotColorType]{plotColor} & \hfuzz=500pt   \\
\hline
\end{tabularx}


\subsection{Custom}
Draws a custom line. The line \config{style} code is described in
\url{https://docs.generic-mapping-tools.org/latest/cookbook/features.html#specifying-pen-attributes}.


\keepXColumns
\begin{tabularx}{\textwidth}{N T A}
\hline
Name & Type & Annotation\\
\hline
\hfuzz=500pt\includegraphics[width=1em]{element-mustset.pdf}~style & \hfuzz=500pt string & \hfuzz=500pt line style code\\
\hfuzz=500pt\includegraphics[width=1em]{element.pdf}~width & \hfuzz=500pt double & \hfuzz=500pt line width [p]\\
\hfuzz=500pt\includegraphics[width=1em]{element-mustset.pdf}~color & \hfuzz=500pt \hyperref[plotColorType]{plotColor} & \hfuzz=500pt \\
\hline
\end{tabularx}

\clearpage
%==================================

\section{PlotMapLayer}\label{plotMapLayerType}
Defines the content of a map of \program{PlotMap}. Multiple layers are are plotted sequentially.


\subsection{GriddedData}
Creates a regular grid of yxz values. The standard \reference{dataVariables}{general.parser:dataVariables}
are available to select the data column of \configFile{inputfileGriddedData}{griddedData}.
Empty grid cells are not plotted. Cells with more than one value will be set to the mean value.
The grid spacing can be determined automatically for regular rectangular grids otherwise
it must be set with \config{increment}. To get a better display together with some projections
the grid should be internally \config{resample}d to higher resolution.
It is assumed that the points of \configFile{inputfileGriddedData}{griddedData} represents centers of grid cells.
This assumption can be changed with \config{gridlineRegistered} (e.g if the data starts at the north pole).


\keepXColumns
\begin{tabularx}{\textwidth}{N T A}
\hline
Name & Type & Annotation\\
\hline
\hfuzz=500pt\includegraphics[width=1em]{element-mustset.pdf}~inputfileGriddedData & \hfuzz=500pt filename & \hfuzz=500pt \\
\hfuzz=500pt\includegraphics[width=1em]{element-mustset.pdf}~value & \hfuzz=500pt expression & \hfuzz=500pt expression to compute values (input columns are named data0, data1, ...)\\
\hfuzz=500pt\includegraphics[width=1em]{element.pdf}~increment & \hfuzz=500pt angle & \hfuzz=500pt the grid spacing [degrees]\\
\hfuzz=500pt\includegraphics[width=1em]{element.pdf}~illuminate & \hfuzz=500pt boolean & \hfuzz=500pt illuminate grid\\
\hfuzz=500pt\includegraphics[width=1em]{element.pdf}~resample & \hfuzz=500pt sequence & \hfuzz=500pt \\
\hfuzz=500pt\includegraphics[width=1em]{connector.pdf}\includegraphics[width=1em]{element.pdf}~intermediateDpi & \hfuzz=500pt double & \hfuzz=500pt oversample grid for a smoother visual effect\\
\hfuzz=500pt\includegraphics[width=1em]{connector.pdf}\includegraphics[width=1em]{element-mustset.pdf}~interpolationMethod & \hfuzz=500pt choice & \hfuzz=500pt interpolation method for oversampling\\
\hfuzz=500pt\quad\includegraphics[width=1em]{connector.pdf}\includegraphics[width=1em]{element-mustset.pdf}~bspline & \hfuzz=500pt  & \hfuzz=500pt B-Spline interpolation\\
\hfuzz=500pt\quad\includegraphics[width=1em]{connector.pdf}\includegraphics[width=1em]{element-mustset.pdf}~bicubic & \hfuzz=500pt  & \hfuzz=500pt bicubic interpolation\\
\hfuzz=500pt\quad\includegraphics[width=1em]{connector.pdf}\includegraphics[width=1em]{element-mustset.pdf}~bilinear & \hfuzz=500pt  & \hfuzz=500pt bilinear interpolation\\
\hfuzz=500pt\quad\includegraphics[width=1em]{connector.pdf}\includegraphics[width=1em]{element-mustset.pdf}~nearest & \hfuzz=500pt  & \hfuzz=500pt nearest neighbour interpolation\\
\hfuzz=500pt\includegraphics[width=1em]{connector.pdf}\includegraphics[width=1em]{element.pdf}~threshold & \hfuzz=500pt double & \hfuzz=500pt A threshold of 1.0 requires all (4 or 16) nodes involved in interpolation to be non-NaN. 0.5 will interpolate about half way from a non-NaN value; 0.1 will go about 90\% of the way.\\
\hfuzz=500pt\includegraphics[width=1em]{element.pdf}~gridlineRegistered & \hfuzz=500pt boolean & \hfuzz=500pt treat input as point values instead of cell means\\
\hline
\end{tabularx}


\subsection{Points}\label{plotMapLayerType:points}
Draws points (\configClass{symbol}{plotSymbolType}) and/or \configClass{line}{plotLineType}s
between the points. If no \configClass{color}{plotColorType} of the \configClass{symbol}{plotSymbolType}
is given a \configClass{colorbar}{plotColorbarType} is required and the color is determined
by the \config{value} expression. The standard \reference{dataVariables}{general.parser:dataVariables}
are available to select the data column of \configFile{inputfileGriddedData}{griddedData}.


\keepXColumns
\begin{tabularx}{\textwidth}{N T A}
\hline
Name & Type & Annotation\\
\hline
\hfuzz=500pt\includegraphics[width=1em]{element-mustset.pdf}~inputfileGriddedData & \hfuzz=500pt filename & \hfuzz=500pt \\
\hfuzz=500pt\includegraphics[width=1em]{element.pdf}~value & \hfuzz=500pt expression & \hfuzz=500pt expression to compute color (input columns are named data0, data1, ...)\\
\hfuzz=500pt\includegraphics[width=1em]{element.pdf}~symbol & \hfuzz=500pt \hyperref[plotSymbolType]{plotSymbol} & \hfuzz=500pt \\
\hfuzz=500pt\includegraphics[width=1em]{element.pdf}~line & \hfuzz=500pt \hyperref[plotLineType]{plotLine} & \hfuzz=500pt style of connecting lines\\
\hfuzz=500pt\includegraphics[width=1em]{element.pdf}~drawLineAsGreatCircle & \hfuzz=500pt boolean & \hfuzz=500pt draw connecting lines as great circles (otherwise, a straight line is drawn instead)\\
\hline
\end{tabularx}


\subsection{Arrows}
Draws an arrow for each point in \configFile{inputfileGriddedData}{griddedData}.
The arrows are defined by the expressions \config{valueNorth/East}.
The standard \reference{dataVariables}{general.parser:dataVariables}
are available to select the correspondent data columns of \configFile{inputfileGriddedData}{griddedData}.
The \config{scale} factor converts the input units to cm in the plot.
If no \configClass{color}{plotColorType} is given a \configClass{colorbar}{plotColorbarType} is required
and the color is determined by the \config{value} expression.
With \config{scaleArrow} a reference arrow as legend can be plotted inside or outside the map.


\keepXColumns
\begin{tabularx}{\textwidth}{N T A}
\hline
Name & Type & Annotation\\
\hline
\hfuzz=500pt\includegraphics[width=1em]{element-mustset.pdf}~inputfileGriddedData & \hfuzz=500pt filename & \hfuzz=500pt grid file with north and east values for arrows\\
\hfuzz=500pt\includegraphics[width=1em]{element-mustset.pdf}~valueNorth & \hfuzz=500pt expression & \hfuzz=500pt expression to compute north values (input columns are named data0, data1, ...)\\
\hfuzz=500pt\includegraphics[width=1em]{element-mustset.pdf}~valueEast & \hfuzz=500pt expression & \hfuzz=500pt expression to compute east values (input columns are named data0, data1, ...)\\
\hfuzz=500pt\includegraphics[width=1em]{element.pdf}~value & \hfuzz=500pt expression & \hfuzz=500pt expression to compute arrow color (input columns are named data0, data1, ...)\\
\hfuzz=500pt\includegraphics[width=1em]{element.pdf}~scale & \hfuzz=500pt double & \hfuzz=500pt [cm per input unit] length scale factor\\
\hfuzz=500pt\includegraphics[width=1em]{element-mustset.pdf}~penSize & \hfuzz=500pt double & \hfuzz=500pt [pt] width of arrow shaft\\
\hfuzz=500pt\includegraphics[width=1em]{element-mustset.pdf}~headSize & \hfuzz=500pt double & \hfuzz=500pt [pt] size of arrow head, 0: no head, negative: reverse head\\
\hfuzz=500pt\includegraphics[width=1em]{element.pdf}~color & \hfuzz=500pt \hyperref[plotColorType]{plotColor} & \hfuzz=500pt empty: from value\\
\hfuzz=500pt\includegraphics[width=1em]{element.pdf}~scaleArrow & \hfuzz=500pt sequence & \hfuzz=500pt draw an arrow for scale reference\\
\hfuzz=500pt\includegraphics[width=1em]{connector.pdf}\includegraphics[width=1em]{element-mustset.pdf}~originX & \hfuzz=500pt double & \hfuzz=500pt [0-1] 0: left, 1: right\\
\hfuzz=500pt\includegraphics[width=1em]{connector.pdf}\includegraphics[width=1em]{element-mustset.pdf}~originY & \hfuzz=500pt double & \hfuzz=500pt [0-1] 0: bottom, 1: top\\
\hfuzz=500pt\includegraphics[width=1em]{connector.pdf}\includegraphics[width=1em]{element-mustset.pdf}~length & \hfuzz=500pt double & \hfuzz=500pt in same unit as valueNorth and valueEast\\
\hfuzz=500pt\includegraphics[width=1em]{connector.pdf}\includegraphics[width=1em]{element-mustset.pdf}~unit & \hfuzz=500pt string & \hfuzz=500pt displayed unit text (e.g. 1 cm)\\
\hfuzz=500pt\includegraphics[width=1em]{connector.pdf}\includegraphics[width=1em]{element.pdf}~label & \hfuzz=500pt string & \hfuzz=500pt description of the arrows\\
\hline
\end{tabularx}


\subsection{Polygon}
Draws a \configFile{inputfilePolygon}{polygon}.
If \configClass{fillColor}{plotColorType} is not set and a \config{value}
is given the fill color is taken from a \configClass{colorbar}{plotColorbarType}.


\keepXColumns
\begin{tabularx}{\textwidth}{N T A}
\hline
Name & Type & Annotation\\
\hline
\hfuzz=500pt\includegraphics[width=1em]{element-mustset.pdf}~inputfilePolygon & \hfuzz=500pt filename & \hfuzz=500pt \\
\hfuzz=500pt\includegraphics[width=1em]{element.pdf}~line & \hfuzz=500pt \hyperref[plotLineType]{plotLine} & \hfuzz=500pt style of border lines\\
\hfuzz=500pt\includegraphics[width=1em]{element.pdf}~fillColor & \hfuzz=500pt \hyperref[plotColorType]{plotColor} & \hfuzz=500pt polygon fill color (no fill color: determine from value if given, else: no fill)\\
\hfuzz=500pt\includegraphics[width=1em]{element.pdf}~value & \hfuzz=500pt double & \hfuzz=500pt value to compute fill color from a colorbar (ignored if a fillColor is given)\\
\hfuzz=500pt\includegraphics[width=1em]{element.pdf}~drawLineAsGreatCircle & \hfuzz=500pt boolean & \hfuzz=500pt draw connecting lines as great circles (otherwise, a straight line is drawn instead)\\
\hline
\end{tabularx}


\subsection{Coast}
Plots coastlines. GMT provides them in different \config{resolution}s.
Features with an area smaller than \config{minArea} in $km^2$ will not be plotted.


\keepXColumns
\begin{tabularx}{\textwidth}{N T A}
\hline
Name & Type & Annotation\\
\hline
\hfuzz=500pt\includegraphics[width=1em]{element-mustset.pdf}~resolution & \hfuzz=500pt choice & \hfuzz=500pt \\
\hfuzz=500pt\includegraphics[width=1em]{connector.pdf}\includegraphics[width=1em]{element-mustset.pdf}~crude & \hfuzz=500pt  & \hfuzz=500pt \\
\hfuzz=500pt\includegraphics[width=1em]{connector.pdf}\includegraphics[width=1em]{element-mustset.pdf}~low & \hfuzz=500pt  & \hfuzz=500pt \\
\hfuzz=500pt\includegraphics[width=1em]{connector.pdf}\includegraphics[width=1em]{element-mustset.pdf}~medium & \hfuzz=500pt  & \hfuzz=500pt \\
\hfuzz=500pt\includegraphics[width=1em]{connector.pdf}\includegraphics[width=1em]{element-mustset.pdf}~high & \hfuzz=500pt  & \hfuzz=500pt \\
\hfuzz=500pt\includegraphics[width=1em]{connector.pdf}\includegraphics[width=1em]{element-mustset.pdf}~full & \hfuzz=500pt  & \hfuzz=500pt \\
\hfuzz=500pt\includegraphics[width=1em]{element.pdf}~line & \hfuzz=500pt \hyperref[plotLineType]{plotLine} & \hfuzz=500pt line style for coastlines\\
\hfuzz=500pt\includegraphics[width=1em]{element.pdf}~landColor & \hfuzz=500pt \hyperref[plotColorType]{plotColor} & \hfuzz=500pt fill land area\\
\hfuzz=500pt\includegraphics[width=1em]{element.pdf}~oceanColor & \hfuzz=500pt \hyperref[plotColorType]{plotColor} & \hfuzz=500pt fill ocean area\\
\hfuzz=500pt\includegraphics[width=1em]{element.pdf}~minArea & \hfuzz=500pt uint & \hfuzz=500pt [km\textasciicircum{}2] features with a smaller area than this are dropped\\
\hline
\end{tabularx}


\subsection{Rivers}
Plots rivers and lakes. GMT provides different \config{class}es
(\url{https://docs.generic-mapping-tools.org/latest/coast.html}).


\keepXColumns
\begin{tabularx}{\textwidth}{N T A}
\hline
Name & Type & Annotation\\
\hline
\hfuzz=500pt\includegraphics[width=1em]{element-mustset.pdf}~class & \hfuzz=500pt choice & \hfuzz=500pt \\
\hfuzz=500pt\includegraphics[width=1em]{connector.pdf}\includegraphics[width=1em]{element-mustset.pdf}~riversCanalsLakes & \hfuzz=500pt  & \hfuzz=500pt \\
\hfuzz=500pt\includegraphics[width=1em]{connector.pdf}\includegraphics[width=1em]{element-mustset.pdf}~riversCanals & \hfuzz=500pt  & \hfuzz=500pt \\
\hfuzz=500pt\includegraphics[width=1em]{connector.pdf}\includegraphics[width=1em]{element-mustset.pdf}~permanentRiversLakes & \hfuzz=500pt  & \hfuzz=500pt \\
\hfuzz=500pt\includegraphics[width=1em]{connector.pdf}\includegraphics[width=1em]{element-mustset.pdf}~permanentRivers & \hfuzz=500pt  & \hfuzz=500pt \\
\hfuzz=500pt\includegraphics[width=1em]{connector.pdf}\includegraphics[width=1em]{element-mustset.pdf}~intermittentRivers & \hfuzz=500pt  & \hfuzz=500pt \\
\hfuzz=500pt\includegraphics[width=1em]{connector.pdf}\includegraphics[width=1em]{element-mustset.pdf}~canals & \hfuzz=500pt  & \hfuzz=500pt \\
\hfuzz=500pt\includegraphics[width=1em]{connector.pdf}\includegraphics[width=1em]{element-mustset.pdf}~singleClass & \hfuzz=500pt sequence & \hfuzz=500pt \\
\hfuzz=500pt\quad\includegraphics[width=1em]{connector.pdf}\includegraphics[width=1em]{element-mustset-unbounded.pdf}~class & \hfuzz=500pt uint & \hfuzz=500pt 0-10. See GMT documentation\\
\hfuzz=500pt\includegraphics[width=1em]{element-mustset.pdf}~line & \hfuzz=500pt \hyperref[plotLineType]{plotLine} & \hfuzz=500pt \\
\hline
\end{tabularx}


\subsection{PoliticalBoundary}
Plots national boundaries. GMT provides them in different \config{resolution}s.


\keepXColumns
\begin{tabularx}{\textwidth}{N T A}
\hline
Name & Type & Annotation\\
\hline
\hfuzz=500pt\includegraphics[width=1em]{element-mustset.pdf}~resolution & \hfuzz=500pt choice & \hfuzz=500pt \\
\hfuzz=500pt\includegraphics[width=1em]{connector.pdf}\includegraphics[width=1em]{element-mustset.pdf}~crude & \hfuzz=500pt  & \hfuzz=500pt \\
\hfuzz=500pt\includegraphics[width=1em]{connector.pdf}\includegraphics[width=1em]{element-mustset.pdf}~low & \hfuzz=500pt  & \hfuzz=500pt \\
\hfuzz=500pt\includegraphics[width=1em]{connector.pdf}\includegraphics[width=1em]{element-mustset.pdf}~medium & \hfuzz=500pt  & \hfuzz=500pt \\
\hfuzz=500pt\includegraphics[width=1em]{connector.pdf}\includegraphics[width=1em]{element-mustset.pdf}~high & \hfuzz=500pt  & \hfuzz=500pt \\
\hfuzz=500pt\includegraphics[width=1em]{connector.pdf}\includegraphics[width=1em]{element-mustset.pdf}~full & \hfuzz=500pt  & \hfuzz=500pt \\
\hfuzz=500pt\includegraphics[width=1em]{element-mustset.pdf}~line & \hfuzz=500pt \hyperref[plotLineType]{plotLine} & \hfuzz=500pt \\
\hline
\end{tabularx}


\subsection{BlueMarble}
An image of the Earth's surface as seen from outer space -
the image is known as \emph{blue marble}. The directory of \config{inputfileChannels}
contains several files in different resolutions representing the Earth's surface each
month throughout a year.

\fig{!hb}{0.8}{blueMarble}{fig:blueMarbleMap}{The blue marble.}


\keepXColumns
\begin{tabularx}{\textwidth}{N T A}
\hline
Name & Type & Annotation\\
\hline
\hfuzz=500pt\includegraphics[width=1em]{element-mustset.pdf}~inputfileChannels & \hfuzz=500pt filename & \hfuzz=500pt GMT grid file containing layers for red, green and blue values.\\
\hfuzz=500pt\includegraphics[width=1em]{element.pdf}~brightness & \hfuzz=500pt double & \hfuzz=500pt brightness of bitmap [-1, 1]\\
\hfuzz=500pt\includegraphics[width=1em]{element.pdf}~illuminate & \hfuzz=500pt sequence & \hfuzz=500pt add hillshade based on topography\\
\hfuzz=500pt\includegraphics[width=1em]{connector.pdf}\includegraphics[width=1em]{element-mustset.pdf}~inputfileTopography & \hfuzz=500pt filename & \hfuzz=500pt GMT grid file containing topography.\\
\hfuzz=500pt\includegraphics[width=1em]{connector.pdf}\includegraphics[width=1em]{element.pdf}~azimuth & \hfuzz=500pt angle & \hfuzz=500pt direction of lighting source [deg]\\
\hfuzz=500pt\includegraphics[width=1em]{connector.pdf}\includegraphics[width=1em]{element.pdf}~elevation & \hfuzz=500pt angle & \hfuzz=500pt direction of lighting source [deg]\\
\hfuzz=500pt\includegraphics[width=1em]{connector.pdf}\includegraphics[width=1em]{element.pdf}~ambient & \hfuzz=500pt double & \hfuzz=500pt ambient lighting\\
\hfuzz=500pt\includegraphics[width=1em]{connector.pdf}\includegraphics[width=1em]{element.pdf}~diffuse & \hfuzz=500pt double & \hfuzz=500pt diffuse lighting\\
\hfuzz=500pt\includegraphics[width=1em]{connector.pdf}\includegraphics[width=1em]{element.pdf}~specular & \hfuzz=500pt double & \hfuzz=500pt specular reflection\\
\hfuzz=500pt\includegraphics[width=1em]{connector.pdf}\includegraphics[width=1em]{element.pdf}~shine & \hfuzz=500pt double & \hfuzz=500pt surface shine\\
\hfuzz=500pt\includegraphics[width=1em]{connector.pdf}\includegraphics[width=1em]{element.pdf}~amplitude & \hfuzz=500pt double & \hfuzz=500pt scale gradient by factor\\
\hline
\end{tabularx}


\subsection{Text}
Writes a \config{text} at \config{originLongitude} and \config{originLatitude} position in the map.
With \config{clip} the text is cutted at the boundaries of the plotting area.


\keepXColumns
\begin{tabularx}{\textwidth}{N T A}
\hline
Name & Type & Annotation\\
\hline
\hfuzz=500pt\includegraphics[width=1em]{element-mustset.pdf}~text & \hfuzz=500pt string & \hfuzz=500pt \\
\hfuzz=500pt\includegraphics[width=1em]{element-mustset.pdf}~originLongitude & \hfuzz=500pt angle & \hfuzz=500pt [deg]\\
\hfuzz=500pt\includegraphics[width=1em]{element-mustset.pdf}~originLatitude & \hfuzz=500pt angle & \hfuzz=500pt [deg]\\
\hfuzz=500pt\includegraphics[width=1em]{element.pdf}~offsetX & \hfuzz=500pt double & \hfuzz=500pt [cm] x-offset from origin\\
\hfuzz=500pt\includegraphics[width=1em]{element.pdf}~offsetY & \hfuzz=500pt double & \hfuzz=500pt [cm] y-offset from origin\\
\hfuzz=500pt\includegraphics[width=1em]{element.pdf}~alignment & \hfuzz=500pt string & \hfuzz=500pt L, C, R (left, center, right) and T, M, B (top, middle, bottom)\\
\hfuzz=500pt\includegraphics[width=1em]{element.pdf}~fontSize & \hfuzz=500pt double & \hfuzz=500pt \\
\hfuzz=500pt\includegraphics[width=1em]{element-mustset.pdf}~fontColor & \hfuzz=500pt \hyperref[plotColorType]{plotColor} & \hfuzz=500pt \\
\hfuzz=500pt\includegraphics[width=1em]{element.pdf}~clip & \hfuzz=500pt boolean & \hfuzz=500pt clip at boundaries\\
\hline
\end{tabularx}

\clearpage
%==================================

\section{PlotMapProjection}\label{plotMapProjectionType}
Selects the underlying projection of \program{PlotMap}.


\subsection{Robinson}
The Robinson projection, presented by Arthur H. Robinson in 1963,
is a modified cylindrical projection that is neither conformal nor equal-area.
Central meridian and all parallels are straight lines; other meridians are curved.
It uses lookup tables rather than analytic expressions to make the world map look right.


\keepXColumns
\begin{tabularx}{\textwidth}{N T A}
\hline
Name & Type & Annotation\\
\hline
\hfuzz=500pt\includegraphics[width=1em]{element.pdf}~centralMeridian & \hfuzz=500pt angle & \hfuzz=500pt central meridian [degree]\\
\hline
\end{tabularx}


\subsection{Orthographic}
The orthographic azimuthal projection is a perspective projection from infinite distance.
It is therefore often used to give the appearance of a globe viewed from space.


\keepXColumns
\begin{tabularx}{\textwidth}{N T A}
\hline
Name & Type & Annotation\\
\hline
\hfuzz=500pt\includegraphics[width=1em]{element.pdf}~lambdaCenter & \hfuzz=500pt angle & \hfuzz=500pt central point [degree]\\
\hfuzz=500pt\includegraphics[width=1em]{element.pdf}~phiCenter & \hfuzz=500pt angle & \hfuzz=500pt central point [degree]\\
\hline
\end{tabularx}


\subsection{Perspective sphere}
The orthographic azimuthal projection is a perspective projection from infinite distance.
It is therefore often used to give the appearance of a globe viewed from space.


\keepXColumns
\begin{tabularx}{\textwidth}{N T A}
\hline
Name & Type & Annotation\\
\hline
\hfuzz=500pt\includegraphics[width=1em]{element.pdf}~lambdaCenter & \hfuzz=500pt angle & \hfuzz=500pt longitude of central point in degrees\\
\hfuzz=500pt\includegraphics[width=1em]{element.pdf}~phiCenter & \hfuzz=500pt angle & \hfuzz=500pt latitude of central point in degrees\\
\hfuzz=500pt\includegraphics[width=1em]{element.pdf}~altitude & \hfuzz=500pt double & \hfuzz=500pt [km]\\
\hfuzz=500pt\includegraphics[width=1em]{element.pdf}~azimuth & \hfuzz=500pt angle & \hfuzz=500pt to the east of north of view [degrees]\\
\hfuzz=500pt\includegraphics[width=1em]{element.pdf}~tilt & \hfuzz=500pt angle & \hfuzz=500pt upward tilt of the plane of projection, if negative, then the view is centered on the horizon [degrees]\\
\hfuzz=500pt\includegraphics[width=1em]{element.pdf}~viewpointTwist & \hfuzz=500pt angle & \hfuzz=500pt clockwise twist of the viewpoint [degrees]\\
\hfuzz=500pt\includegraphics[width=1em]{element.pdf}~viewpointWidth & \hfuzz=500pt angle & \hfuzz=500pt width of the viewpoint [degrees]\\
\hfuzz=500pt\includegraphics[width=1em]{element.pdf}~viewpointHeight & \hfuzz=500pt angle & \hfuzz=500pt height of the viewpoint [degrees]\\
\hline
\end{tabularx}


\subsection{Polar}
Stereographic projection around given central point.


\keepXColumns
\begin{tabularx}{\textwidth}{N T A}
\hline
Name & Type & Annotation\\
\hline
\hfuzz=500pt\includegraphics[width=1em]{element.pdf}~lambdaCenter & \hfuzz=500pt angle & \hfuzz=500pt longitude of central point in degrees\\
\hfuzz=500pt\includegraphics[width=1em]{element.pdf}~phiCenter & \hfuzz=500pt angle & \hfuzz=500pt latitude of central point in degrees\\
\hline
\end{tabularx}


\subsection{Skyplot}
Skyplot used to plot azimuth/elevation data as generated by
\program{GnssAntennaDefinition2Skyplot} or \program{GnssResiduals2Skyplot}.


\subsection{UTM}
A particular subset of the transverse Mercator is the Universal Transverse Mercator (UTM)
which was adopted by the US Army for large-scale military maps.
Here, the globe is divided into 60 zones between 84$^{o}$S and 84$^{o}$N, most of which are 6$^{o}$ wide.
Each of these UTM zones have their unique central meridian.


\keepXColumns
\begin{tabularx}{\textwidth}{N T A}
\hline
Name & Type & Annotation\\
\hline
\hfuzz=500pt\includegraphics[width=1em]{element-mustset.pdf}~zone & \hfuzz=500pt string & \hfuzz=500pt UTM zone code (e.g. 33N)\\
\hline
\end{tabularx}


\subsection{Lambert}
This conic projection was designed by Lambert (1772) and has been used extensively for mapping of regions with predominantly east-west orientation.


\keepXColumns
\begin{tabularx}{\textwidth}{N T A}
\hline
Name & Type & Annotation\\
\hline
\hfuzz=500pt\includegraphics[width=1em]{element-mustset.pdf}~lambda0 & \hfuzz=500pt angle & \hfuzz=500pt longitude of projection center [deg]\\
\hfuzz=500pt\includegraphics[width=1em]{element-mustset.pdf}~phi0 & \hfuzz=500pt angle & \hfuzz=500pt latitude of projection centert [deg]\\
\hfuzz=500pt\includegraphics[width=1em]{element-mustset.pdf}~phi1 & \hfuzz=500pt angle & \hfuzz=500pt latitude of first standard parallel [deg]\\
\hfuzz=500pt\includegraphics[width=1em]{element-mustset.pdf}~phi2 & \hfuzz=500pt angle & \hfuzz=500pt latitude of first standard parallel [deg]\\
\hline
\end{tabularx}


\subsection{Linear}
Linear mapping of longitude/latitude to x/y (Plate Caree).

\clearpage
%==================================

\section{PlotSymbol}\label{plotSymbolType}
Plots a symbol as used e.g. in \configClass{plotGraphLayer:linesAndPoints}{plotGraphLayerType:linesAndPoints}
or \configClass{plotMapLayer:points}{plotMapLayerType:points}.


\keepXColumns
\begin{tabularx}{\textwidth}{N T A}
\hline
Name & Type & Annotation\\
\hline
\hfuzz=500pt\includegraphics[width=1em]{element-mustset.pdf}~plotSymbolType & \hfuzz=500pt choice & \hfuzz=500pt symbol\\
\hfuzz=500pt\includegraphics[width=1em]{connector.pdf}\includegraphics[width=1em]{element-mustset.pdf}~circle & \hfuzz=500pt sequence & \hfuzz=500pt \\
\hfuzz=500pt\quad\includegraphics[width=1em]{connector.pdf}\includegraphics[width=1em]{element.pdf}~color & \hfuzz=500pt \hyperref[plotColorType]{plotColor} & \hfuzz=500pt empty: determined from value\\
\hfuzz=500pt\quad\includegraphics[width=1em]{connector.pdf}\includegraphics[width=1em]{element.pdf}~size & \hfuzz=500pt double & \hfuzz=500pt size of symbol [point]\\
\hfuzz=500pt\quad\includegraphics[width=1em]{connector.pdf}\includegraphics[width=1em]{element.pdf}~blackContour & \hfuzz=500pt boolean & \hfuzz=500pt \\
\hfuzz=500pt\includegraphics[width=1em]{connector.pdf}\includegraphics[width=1em]{element-mustset.pdf}~star & \hfuzz=500pt sequence & \hfuzz=500pt \\
\hfuzz=500pt\quad\includegraphics[width=1em]{connector.pdf}\includegraphics[width=1em]{element.pdf}~color & \hfuzz=500pt \hyperref[plotColorType]{plotColor} & \hfuzz=500pt empty: determined from value\\
\hfuzz=500pt\quad\includegraphics[width=1em]{connector.pdf}\includegraphics[width=1em]{element.pdf}~size & \hfuzz=500pt double & \hfuzz=500pt size of symbol [point]\\
\hfuzz=500pt\quad\includegraphics[width=1em]{connector.pdf}\includegraphics[width=1em]{element.pdf}~blackContour & \hfuzz=500pt boolean & \hfuzz=500pt \\
\hfuzz=500pt\includegraphics[width=1em]{connector.pdf}\includegraphics[width=1em]{element-mustset.pdf}~cross & \hfuzz=500pt sequence & \hfuzz=500pt \\
\hfuzz=500pt\quad\includegraphics[width=1em]{connector.pdf}\includegraphics[width=1em]{element.pdf}~color & \hfuzz=500pt \hyperref[plotColorType]{plotColor} & \hfuzz=500pt empty: determined from value\\
\hfuzz=500pt\quad\includegraphics[width=1em]{connector.pdf}\includegraphics[width=1em]{element.pdf}~size & \hfuzz=500pt double & \hfuzz=500pt size of symbol [point]\\
\hfuzz=500pt\quad\includegraphics[width=1em]{connector.pdf}\includegraphics[width=1em]{element.pdf}~blackContour & \hfuzz=500pt boolean & \hfuzz=500pt \\
\hfuzz=500pt\includegraphics[width=1em]{connector.pdf}\includegraphics[width=1em]{element-mustset.pdf}~square & \hfuzz=500pt sequence & \hfuzz=500pt \\
\hfuzz=500pt\quad\includegraphics[width=1em]{connector.pdf}\includegraphics[width=1em]{element.pdf}~color & \hfuzz=500pt \hyperref[plotColorType]{plotColor} & \hfuzz=500pt empty: determined from value\\
\hfuzz=500pt\quad\includegraphics[width=1em]{connector.pdf}\includegraphics[width=1em]{element.pdf}~size & \hfuzz=500pt double & \hfuzz=500pt size of symbol [point]\\
\hfuzz=500pt\quad\includegraphics[width=1em]{connector.pdf}\includegraphics[width=1em]{element.pdf}~blackContour & \hfuzz=500pt boolean & \hfuzz=500pt \\
\hfuzz=500pt\includegraphics[width=1em]{connector.pdf}\includegraphics[width=1em]{element-mustset.pdf}~triangle & \hfuzz=500pt sequence & \hfuzz=500pt \\
\hfuzz=500pt\quad\includegraphics[width=1em]{connector.pdf}\includegraphics[width=1em]{element.pdf}~color & \hfuzz=500pt \hyperref[plotColorType]{plotColor} & \hfuzz=500pt empty: determined from value\\
\hfuzz=500pt\quad\includegraphics[width=1em]{connector.pdf}\includegraphics[width=1em]{element.pdf}~size & \hfuzz=500pt double & \hfuzz=500pt size of symbol [point]\\
\hfuzz=500pt\quad\includegraphics[width=1em]{connector.pdf}\includegraphics[width=1em]{element.pdf}~blackContour & \hfuzz=500pt boolean & \hfuzz=500pt \\
\hfuzz=500pt\includegraphics[width=1em]{connector.pdf}\includegraphics[width=1em]{element-mustset.pdf}~diamond & \hfuzz=500pt sequence & \hfuzz=500pt \\
\hfuzz=500pt\quad\includegraphics[width=1em]{connector.pdf}\includegraphics[width=1em]{element.pdf}~color & \hfuzz=500pt \hyperref[plotColorType]{plotColor} & \hfuzz=500pt empty: determined from value\\
\hfuzz=500pt\quad\includegraphics[width=1em]{connector.pdf}\includegraphics[width=1em]{element.pdf}~size & \hfuzz=500pt double & \hfuzz=500pt size of symbol [point]\\
\hfuzz=500pt\quad\includegraphics[width=1em]{connector.pdf}\includegraphics[width=1em]{element.pdf}~blackContour & \hfuzz=500pt boolean & \hfuzz=500pt \\
\hfuzz=500pt\includegraphics[width=1em]{connector.pdf}\includegraphics[width=1em]{element-mustset.pdf}~dash & \hfuzz=500pt sequence & \hfuzz=500pt \\
\hfuzz=500pt\quad\includegraphics[width=1em]{connector.pdf}\includegraphics[width=1em]{element.pdf}~color & \hfuzz=500pt \hyperref[plotColorType]{plotColor} & \hfuzz=500pt empty: determined from value\\
\hfuzz=500pt\quad\includegraphics[width=1em]{connector.pdf}\includegraphics[width=1em]{element.pdf}~size & \hfuzz=500pt double & \hfuzz=500pt size of symbol [point]\\
\hfuzz=500pt\quad\includegraphics[width=1em]{connector.pdf}\includegraphics[width=1em]{element.pdf}~blackContour & \hfuzz=500pt boolean & \hfuzz=500pt \\
\hline
\end{tabularx}

\clearpage
%==================================

\section{PodRightSide}\label{podRightSideType}
Observation vector for precise orbit data (POD) of \configClass{observation}{observationType}
equations in a least squares adjustment. The observations are reduced by the effect of
\configFile{inputfileAccelerometer}{instrument} and \configClass{forces}{forcesType}
(observed minus computed).


\keepXColumns
\begin{tabularx}{\textwidth}{N T A}
\hline
Name & Type & Annotation\\
\hline
\hfuzz=500pt\includegraphics[width=1em]{element-mustset.pdf}~podRightSideType & \hfuzz=500pt sequence & \hfuzz=500pt \\
\hfuzz=500pt\includegraphics[width=1em]{connector.pdf}\includegraphics[width=1em]{element-mustset.pdf}~inputfileOrbit & \hfuzz=500pt filename & \hfuzz=500pt kinematic positions of satellite as observations\\
\hfuzz=500pt\includegraphics[width=1em]{connector.pdf}\includegraphics[width=1em]{element.pdf}~inputfileAccelerometer & \hfuzz=500pt filename & \hfuzz=500pt non-gravitational forces in satellite reference frame\\
\hfuzz=500pt\includegraphics[width=1em]{connector.pdf}\includegraphics[width=1em]{element-mustset.pdf}~forces & \hfuzz=500pt \hyperref[forcesType]{forces} & \hfuzz=500pt \\
\hline
\end{tabularx}

\clearpage
%==================================

\section{SggRightSide}\label{sggRightSideType}
Observation vector for gradiometer data (satellite gravity gradiometry, SGG)
of \configClass{observation}{observationType} equations in a least squares adjustment.
The observations are reduced by an \configFile{inputfileReferenceGradiometer}{instrument},
the effect of \configClass{referencefield}{gravityfieldType}, and \configClass{tides}{tidesType}
(observed minus computed).

The reference gradiometer data can be precomputed with \program{SimulateGradiometer}.


\keepXColumns
\begin{tabularx}{\textwidth}{N T A}
\hline
Name & Type & Annotation\\
\hline
\hfuzz=500pt\includegraphics[width=1em]{element-mustset.pdf}~sggRightSideType & \hfuzz=500pt sequence & \hfuzz=500pt \\
\hfuzz=500pt\includegraphics[width=1em]{connector.pdf}\includegraphics[width=1em]{element-mustset.pdf}~inputfileGradiometer & \hfuzz=500pt filename & \hfuzz=500pt observed gravity gradients\\
\hfuzz=500pt\includegraphics[width=1em]{connector.pdf}\includegraphics[width=1em]{element-unbounded.pdf}~inputfileReferenceGradiometer & \hfuzz=500pt filename & \hfuzz=500pt precomputed gradients at orbit positions\\
\hfuzz=500pt\includegraphics[width=1em]{connector.pdf}\includegraphics[width=1em]{element-unbounded.pdf}~referencefield & \hfuzz=500pt \hyperref[gravityfieldType]{gravityfield} & \hfuzz=500pt \\
\hfuzz=500pt\includegraphics[width=1em]{connector.pdf}\includegraphics[width=1em]{element-unbounded.pdf}~tides & \hfuzz=500pt \hyperref[tidesType]{tides} & \hfuzz=500pt \\
\hline
\end{tabularx}

\clearpage
%==================================

\section{SphericalHarmonicsFilter}\label{sphericalHarmonicsFilterType}
Filtering of a spherical harmonics expansion.


\subsection{DDK}
Orderwise filtering with the DDK filter by Kusche et al. 2009.


\keepXColumns
\begin{tabularx}{\textwidth}{N T A}
\hline
Name & Type & Annotation\\
\hline
\hfuzz=500pt\includegraphics[width=1em]{element-mustset.pdf}~level & \hfuzz=500pt uint & \hfuzz=500pt DDK filter level (1, 2, 3, ...)\\
\hfuzz=500pt\includegraphics[width=1em]{element-mustset.pdf}~inputfileNormalEquation & \hfuzz=500pt filename & \hfuzz=500pt \\
\hline
\end{tabularx}


\subsection{Gauss}
Filtering the spherical harmonics expansion with a Gaussian filter.
\config{radius} gives the filter radius on the Earth surface in km.


\keepXColumns
\begin{tabularx}{\textwidth}{N T A}
\hline
Name & Type & Annotation\\
\hline
\hfuzz=500pt\includegraphics[width=1em]{element-mustset.pdf}~radius & \hfuzz=500pt double & \hfuzz=500pt filter radius [km]\\
\hline
\end{tabularx}


\subsection{Matrix}
Filtering the spherical harmonics expansion with a matrix filter.


\keepXColumns
\begin{tabularx}{\textwidth}{N T A}
\hline
Name & Type & Annotation\\
\hline
\hfuzz=500pt\includegraphics[width=1em]{element-mustset.pdf}~inputfileMatrix & \hfuzz=500pt filename & \hfuzz=500pt \\
\hfuzz=500pt\includegraphics[width=1em]{element-mustset.pdf}~minDegree & \hfuzz=500pt uint & \hfuzz=500pt of matrix\\
\hfuzz=500pt\includegraphics[width=1em]{element-mustset.pdf}~maxDegree & \hfuzz=500pt uint & \hfuzz=500pt of matrix\\
\hfuzz=500pt\includegraphics[width=1em]{element-mustset.pdf}~numbering & \hfuzz=500pt \hyperref[sphericalHarmonicsNumberingType]{sphericalHarmonicsNumbering} & \hfuzz=500pt numbering scheme of the matrix\\
\hline
\end{tabularx}

\clearpage
%==================================

\section{SphericalHarmonicsNumbering}\label{sphericalHarmonicsNumberingType}
This class organizes the numbering scheme of spherical harmonics coefficients
in a parameter vector (e.g \program{Gravityfield2SphericalHarmonicsVector} and the design matrix of
\configClass{parametrizationGravity:sphericalHarmoncis}{parametrizationGravityType:sphericalHarmonics}.


\subsection{Degree}
Numbering degree by degree:
\[ c20, c21, s21, c22, s22, c30, c31, s31, c32, s32,\ldots \]


\subsection{Order}
Numbering order by order:
\[ c20, c30, c40, \ldots, c21, s21, c31, s31, \ldots, c22, s22 \]


\subsection{OrderNonAlternating}
Numbering order by order with cnm, snm non-alternating:
\[ c20, c30, c40, \ldots, c21, c31, c41, \ldots, s21, s31, s41, \]


\subsection{File}
Numbering as specified in the chosen file.
The \configFile{inputfile}{matrix} is a matrix with the first column indicating cnm/snm with 0 or 1.
The second and third column specify degree and order.


\keepXColumns
\begin{tabularx}{\textwidth}{N T A}
\hline
Name & Type & Annotation\\
\hline
\hfuzz=500pt\includegraphics[width=1em]{element-mustset.pdf}~inputfile & \hfuzz=500pt filename & \hfuzz=500pt \\
\hline
\end{tabularx}

\clearpage
%==================================

\section{SstRightSide}\label{sstRightSideType}
Observation vector for GRACE like data (satellite-tracking and precise orbit data (POD))
of \configClass{observation}{observationType} equations in a least squares adjustment.
The observations are reduced by the effect of \configFile{inputfileAccelerometer}{instrument}
and \configClass{forces}{forcesType} (observed minus computed).


\keepXColumns
\begin{tabularx}{\textwidth}{N T A}
\hline
Name & Type & Annotation\\
\hline
\hfuzz=500pt\includegraphics[width=1em]{element-mustset.pdf}~sstRightSideType & \hfuzz=500pt sequence & \hfuzz=500pt \\
\hfuzz=500pt\includegraphics[width=1em]{connector.pdf}\includegraphics[width=1em]{element-unbounded.pdf}~inputfileSatelliteTracking & \hfuzz=500pt filename & \hfuzz=500pt ranging observations and corrections\\
\hfuzz=500pt\includegraphics[width=1em]{connector.pdf}\includegraphics[width=1em]{element.pdf}~inputfileOrbit1 & \hfuzz=500pt filename & \hfuzz=500pt kinematic positions of satellite A as observations\\
\hfuzz=500pt\includegraphics[width=1em]{connector.pdf}\includegraphics[width=1em]{element.pdf}~inputfileOrbit2 & \hfuzz=500pt filename & \hfuzz=500pt kinematic positions of satellite B as observations\\
\hfuzz=500pt\includegraphics[width=1em]{connector.pdf}\includegraphics[width=1em]{element.pdf}~inputfileAccelerometer1 & \hfuzz=500pt filename & \hfuzz=500pt non-gravitational forces in satellite reference frame A\\
\hfuzz=500pt\includegraphics[width=1em]{connector.pdf}\includegraphics[width=1em]{element.pdf}~inputfileAccelerometer2 & \hfuzz=500pt filename & \hfuzz=500pt non-gravitational forces in satellite reference frame B\\
\hfuzz=500pt\includegraphics[width=1em]{connector.pdf}\includegraphics[width=1em]{element-mustset.pdf}~forces & \hfuzz=500pt \hyperref[forcesType]{forces} & \hfuzz=500pt \\
\hline
\end{tabularx}

\clearpage
%==================================

\section{Thermosphere}\label{thermosphereType}
This class provides functions for calculating the density, temperature and velocity
in the thermosphere.
The wind is computed by HWM14 model if \config{hwm14DataDirectory} is provided.
A quiet thermosphere is assumed if \config{inputfileMagnetic3hAp} is not given.


\subsection{JB2008}
Thermosphere parameters from the JB2008 model:

Bowman, B. R., Tobiska, W. K., Marcos, F. A., Huang, C. Y., Lin, C. S., Burke, W. J. (2008).
A new empirical thermospheric density model JB2008 using new solar and geomagnetic indices.
 In AIAA/AAS Astrodynamics Specialist Conference and Exhibit. \url{https://doi.org/10.2514/6.2008-6438}


\keepXColumns
\begin{tabularx}{\textwidth}{N T A}
\hline
Name & Type & Annotation\\
\hline
\hfuzz=500pt\includegraphics[width=1em]{element-mustset.pdf}~inputfileSolfsmy & \hfuzz=500pt filename & \hfuzz=500pt solar indices\\
\hfuzz=500pt\includegraphics[width=1em]{element-mustset.pdf}~inputfileDtc & \hfuzz=500pt filename & \hfuzz=500pt \\
\hfuzz=500pt\includegraphics[width=1em]{element.pdf}~inputfileMagnetic3hAp & \hfuzz=500pt filename & \hfuzz=500pt indicies for wind model\\
\hfuzz=500pt\includegraphics[width=1em]{element.pdf}~hwm14DataDirectory & \hfuzz=500pt filename & \hfuzz=500pt directory containing dwm07b104i.dat, gd2qd.dat, hwm123114.bin\\
\hline
\end{tabularx}

\clearpage
%==================================

\section{Tides}\label{tidesType}
This class computes functionals of the time depending tide potential,
e.g potential, acceleration or gravity gradients.

If several instances of the class are given the results are summed up.
Before summation every single result is multiplicated by a \config{factor}.
To get the difference between two ocean tide models you must choose one factor by 1
and the other by -1. To get the mean of two models just set each factor to 0.5.


\subsection{AstronomicalTide}\label{tidesType:astronomicalTide}
This class computes the tide generating potential (TGP) of sun, moon
and planets (Mercury, Venus, Mars, Jupiter, Saturn).
It takes into account the flattening of the Earth (At the moment only at the acceleration level).

The computed result is multiplied with \config{factor}.


\keepXColumns
\begin{tabularx}{\textwidth}{N T A}
\hline
Name & Type & Annotation\\
\hline
\hfuzz=500pt\includegraphics[width=1em]{element.pdf}~useMoon & \hfuzz=500pt boolean & \hfuzz=500pt TGP of moon\\
\hfuzz=500pt\includegraphics[width=1em]{element.pdf}~useSun & \hfuzz=500pt boolean & \hfuzz=500pt TGP of sun\\
\hfuzz=500pt\includegraphics[width=1em]{element.pdf}~usePlanets & \hfuzz=500pt boolean & \hfuzz=500pt TGP of planets\\
\hfuzz=500pt\includegraphics[width=1em]{element.pdf}~useEarth & \hfuzz=500pt boolean & \hfuzz=500pt TGP of Earth\\
\hfuzz=500pt\includegraphics[width=1em]{element.pdf}~c20Earth & \hfuzz=500pt double & \hfuzz=500pt J2 flattening of the Earth\\
\hfuzz=500pt\includegraphics[width=1em]{element.pdf}~factor & \hfuzz=500pt double & \hfuzz=500pt the result is multplied by this factor, set -1 to substract the field\\
\hline
\end{tabularx}


\subsection{EarthTide}\label{tidesType:earthTide}
This class computes the earth tide according to the IERS2003 conventions.
The values of solid Earth tide external potential Love numbers and
the frequency dependent corrections of these values are given in the file
\configFile{inputfileEarthtide}{earthTide}. The effect of the permanent tide is removed if
\config{includePermanentTide} is set to false.

The computed result is multiplied with \config{factor}.


\keepXColumns
\begin{tabularx}{\textwidth}{N T A}
\hline
Name & Type & Annotation\\
\hline
\hfuzz=500pt\includegraphics[width=1em]{element-mustset.pdf}~inputfileEarthtide & \hfuzz=500pt filename & \hfuzz=500pt \\
\hfuzz=500pt\includegraphics[width=1em]{element.pdf}~includePermanentTide & \hfuzz=500pt boolean & \hfuzz=500pt results in FALSE: zero tide, TRUE: tide free gravity field\\
\hfuzz=500pt\includegraphics[width=1em]{element.pdf}~factor & \hfuzz=500pt double & \hfuzz=500pt the result is multplied by this factor, set -1 to substract the field\\
\hline
\end{tabularx}


\subsection{PoleTide}\label{tidesType:poleTide}
The potential coefficients of the solid Earth pole tide according to the
IERS2003 conventions are given by
\begin{equation}
\begin{split}
\Delta c_{21} &= s\cdot(m_1 + o\cdot m_2), \\
\Delta s_{21} &= s\cdot(m_2 - o\cdot m_1),
\end{split}
\end{equation}
with $s$ is the \config{scale}, $o$ is the \config{outPhase} and
$(m_1,m_2)$ are the wobble variables in seconds of arc.
They are related to the polar motion variables $(x_p,y_p)$ according to
\begin{equation}
\begin{split}
m_1 &=  (x_p - \bar{x}_p), \\
m_2 &= -(y_p - \bar{y}_p),
\end{split}
\end{equation}
The mean pole $(\bar{x}_p, \bar{y}_p)$ is approximated by a polynomial
read from \configFile{inputfileMeanPole}{meanPolarMotion}.

The displacment is calculated with
\begin{equation}
\begin{split}
S_r          &= -v\sin2\vartheta(m_1\cos\lambda+m_2\sin\lambda),\\
S_\vartheta &= -h\cos2\vartheta(m_1\cos\lambda+m_2\sin\lambda),\\
S_\lambda   &=  h\cos\vartheta(m_1\sin\lambda-m_2\cos\lambda),
\end{split}
\end{equation}
where $h$ is the \config{horizontalDisplacement}
and $v$ is the \config{verticalDisplacement}.

The computed result is multiplied with \config{factor}.


\keepXColumns
\begin{tabularx}{\textwidth}{N T A}
\hline
Name & Type & Annotation\\
\hline
\hfuzz=500pt\includegraphics[width=1em]{element.pdf}~scale & \hfuzz=500pt double & \hfuzz=500pt \\
\hfuzz=500pt\includegraphics[width=1em]{element.pdf}~outPhase & \hfuzz=500pt double & \hfuzz=500pt \\
\hfuzz=500pt\includegraphics[width=1em]{element-mustset.pdf}~inputfileMeanPole & \hfuzz=500pt filename & \hfuzz=500pt \\
\hfuzz=500pt\includegraphics[width=1em]{element.pdf}~horizontalDisplacement & \hfuzz=500pt double & \hfuzz=500pt [m]\\
\hfuzz=500pt\includegraphics[width=1em]{element.pdf}~verticalDisplacement & \hfuzz=500pt double & \hfuzz=500pt [m]\\
\hfuzz=500pt\includegraphics[width=1em]{element.pdf}~factor & \hfuzz=500pt double & \hfuzz=500pt the result is multplied by this factor, set -1 to substract the field\\
\hline
\end{tabularx}


\subsection{OceanPoleTide}\label{tidesType:oceanPoleTide}
The ocean pole tide is generated by the centrifugal effect of polar motion on the oceans.
The potential coefficients of this effect is given by
IERS2003 conventions are given by
\begin{equation}
\begin{Bmatrix}
\Delta c_{nm}  \\
\Delta s_{nm}
\end{Bmatrix}=
\begin{Bmatrix}
c_{nm}^R  \\
s_{nm}^R
\end{Bmatrix}
(m_1\gamma^R+m_2\gamma^I)+
\begin{Bmatrix}
c_{nm}^I  \\
s_{nm}^I
\end{Bmatrix}
(m_2\gamma^R-m_1\gamma^I)
\end{equation}
where the coefficients are read from file \configFile{inputfileOceanPole}{oceanPoleTide},
$\gamma=\gamma^R+i\gamma^I$ is given by \config{gammaReal} and
\config{gammaImaginary} and $(m_1,m_2)$ are the wobble variables in radians.
They are related to the polar motion variables $(x_p,y_p)$ according to
\begin{equation}
\begin{split}
m_1 &=  (x_p - \bar{x}_p), \\
m_2 &= -(y_p - \bar{y}_p),
\end{split}
\end{equation}
The mean pole $(\bar{x}_p, \bar{y}_p)$ is approximated by a polynomial
read from \configFile{inputfileMeanPole}{meanPolarMotion}.

The computed result is multiplied with \config{factor}.


\keepXColumns
\begin{tabularx}{\textwidth}{N T A}
\hline
Name & Type & Annotation\\
\hline
\hfuzz=500pt\includegraphics[width=1em]{element-mustset.pdf}~inputfileOceanPole & \hfuzz=500pt filename & \hfuzz=500pt \\
\hfuzz=500pt\includegraphics[width=1em]{element.pdf}~minDegree & \hfuzz=500pt uint & \hfuzz=500pt \\
\hfuzz=500pt\includegraphics[width=1em]{element.pdf}~maxDegree & \hfuzz=500pt uint & \hfuzz=500pt \\
\hfuzz=500pt\includegraphics[width=1em]{element.pdf}~gammaReal & \hfuzz=500pt double & \hfuzz=500pt \\
\hfuzz=500pt\includegraphics[width=1em]{element.pdf}~gammaImaginary & \hfuzz=500pt double & \hfuzz=500pt \\
\hfuzz=500pt\includegraphics[width=1em]{element-mustset.pdf}~inputfileMeanPole & \hfuzz=500pt filename & \hfuzz=500pt \\
\hfuzz=500pt\includegraphics[width=1em]{element.pdf}~factor & \hfuzz=500pt double & \hfuzz=500pt the result is multplied by this factor, set -1 to substract the field\\
\hline
\end{tabularx}


\subsection{DoodsonHarmonicTide}\label{tidesType:doodsonHarmonicTide}
The time variable potential of ocean tides is given by a fourier expansion
\begin{equation}
V(\M x,t) = \sum_{f} V_f^c(\M x)\cos(\Theta_f(t)) + V_f^s(\M x)\sin(\Theta_f(t)),
\end{equation}
where $V_f^c(\M x)$ and $V_f^s(\M x)$ are spherical harmonics expansions and are
read from the file \configFile{inputfileDoodsonHarmonic}{doodsonHarmonic}.
If set the expansion is limited in the range between \config{minDegree}
and \config{maxDegree} inclusivly.
$\Theta_f(t)$ are the arguments of the tide constituents~$f$:
\begin{equation}
\Theta_f(t) = \sum_{i=1}^6 n_f^i\beta_i(t),
\end{equation}
where $\beta_i(t)$ are the Doodson's fundamental arguments ($\tau,s,h,p,N',p_s$)
and $n_f^i$ are the Doodson multipliers for the term at frequency~$f$.

The major constituents given by \configFile{inputfileDoodsonHarmonic}{doodsonHarmonic} can be used to
interpolate minor tidal constituents using the file \configFile{inputfileAdmittance}{admittance}.
This file can be created with \program{DoodsonHarmonicsCalculateAdmittance}.

After the interpolation step a selection of the computed constituents can be
choosen by \configClass{selectDoodson}{doodson}. Only these constiuents are considered for the results.
If no \configClass{selectDoodson}{doodson} is set all constituents will be used. The constituents can
be coded as Doodson number (e.g. 255.555) or as names intoduced by Darwin (e.g. M2).

The computed result is multiplied with \config{factor}.


\keepXColumns
\begin{tabularx}{\textwidth}{N T A}
\hline
Name & Type & Annotation\\
\hline
\hfuzz=500pt\includegraphics[width=1em]{element-mustset.pdf}~inputfileTides & \hfuzz=500pt filename & \hfuzz=500pt \\
\hfuzz=500pt\includegraphics[width=1em]{element.pdf}~inputfileAdmittance & \hfuzz=500pt filename & \hfuzz=500pt interpolation of minor constituents\\
\hfuzz=500pt\includegraphics[width=1em]{element-unbounded.pdf}~selectDoodson & \hfuzz=500pt \hyperref[doodson]{doodson} & \hfuzz=500pt consider only these constituents, code number (e.g. 255.555) or darwin name (e.g. M2)\\
\hfuzz=500pt\includegraphics[width=1em]{element.pdf}~minDegree & \hfuzz=500pt uint & \hfuzz=500pt \\
\hfuzz=500pt\includegraphics[width=1em]{element.pdf}~maxDegree & \hfuzz=500pt uint & \hfuzz=500pt \\
\hfuzz=500pt\includegraphics[width=1em]{element.pdf}~nodeCorr & \hfuzz=500pt uint & \hfuzz=500pt nodal corrections: 0-no corr, 1-IHO, 2-Schureman\\
\hfuzz=500pt\includegraphics[width=1em]{element.pdf}~factor & \hfuzz=500pt double & \hfuzz=500pt the result is multplied by this factor, set -1 to substract the field\\
\hline
\end{tabularx}


\subsection{Centrifugal}\label{tidesType:centrifugal}
Computes the centrifugal potential in a rotating system
\begin{equation}
V(\M r, t) = \frac{1}{2} (\M\omega(t)\times\M r)^2.
\end{equation}
The current rotation vector $\M\omega(t)$ is computed from the
\configClass{earthRotation}{earthRotationType}
provided by the calling program.
The computed result is multiplied with \config{factor}.

Be careful, the centrifugal potential is not harmonic.
Convolution with a harmonic kernel (e.g. to compute gravity
anomalies) is not meaningful.


\keepXColumns
\begin{tabularx}{\textwidth}{N T A}
\hline
Name & Type & Annotation\\
\hline
\hfuzz=500pt\includegraphics[width=1em]{element.pdf}~factor & \hfuzz=500pt double & \hfuzz=500pt the result is multplied by this factor, set -1 to substract the field\\
\hline
\end{tabularx}


\subsection{SolidMoonTide}
This class computes the solid moon tide according to the IERS2010 conventions.
The values of solid Moon tide external potential Love numbers are given and
there are no frequency dependent corrections of these values.
The computed result is multiplied with \config{factor}.


\keepXColumns
\begin{tabularx}{\textwidth}{N T A}
\hline
Name & Type & Annotation\\
\hline
\hfuzz=500pt\includegraphics[width=1em]{element.pdf}~k20 & \hfuzz=500pt double & \hfuzz=500pt \\
\hfuzz=500pt\includegraphics[width=1em]{element.pdf}~k30 & \hfuzz=500pt double & \hfuzz=500pt \\
\hfuzz=500pt\includegraphics[width=1em]{element.pdf}~factor & \hfuzz=500pt double & \hfuzz=500pt the result is multplied by this factor, set -1 to substract the field\\
\hline
\end{tabularx}

\clearpage
%==================================

\section{TimeSeries}\label{timeSeriesType}
This class generates a series of points in time. The series is always sorted in ascending order.
Depending of the application the series is interpreted as list of points or as intervals between the points.

\fig{!hb}{0.4}{timeSeriesIntervals}{fig:timeSeriesIntervals}{List of points $t_i$ vs. intervals $T_i$.}


\subsection{UniformSampling}\label{timeSeriesType:uniformSampling}
Generates a time series with uniform sampling. The first point in time will be \config{timeStart}.
The last generated point in time will be less or equal \config{timeEnd}.
The time step between generated points in time is given by \config{sampling}.


\keepXColumns
\begin{tabularx}{\textwidth}{N T A}
\hline
Name & Type & Annotation\\
\hline
\hfuzz=500pt\includegraphics[width=1em]{element-mustset.pdf}~timeStart & \hfuzz=500pt time & \hfuzz=500pt first point in time\\
\hfuzz=500pt\includegraphics[width=1em]{element-mustset.pdf}~timeEnd & \hfuzz=500pt time & \hfuzz=500pt last point in time will be less or equal timeEnd\\
\hfuzz=500pt\includegraphics[width=1em]{element-mustset.pdf}~sampling & \hfuzz=500pt time & \hfuzz=500pt time step between points in time\\
\hline
\end{tabularx}


\subsection{UniformInterval}
Generates a time series with uniform sampling between \config{timeStart} and \config{timeEnd}.
\config{intervallCount} gives the count of intervals. This class generates count+1 points in time
inclusive \config{timeStart} and \config{timeEnd}.


\keepXColumns
\begin{tabularx}{\textwidth}{N T A}
\hline
Name & Type & Annotation\\
\hline
\hfuzz=500pt\includegraphics[width=1em]{element-mustset.pdf}~timeStart & \hfuzz=500pt time & \hfuzz=500pt 1st point of the time series\\
\hfuzz=500pt\includegraphics[width=1em]{element-mustset.pdf}~timeEnd & \hfuzz=500pt time & \hfuzz=500pt last point of the time series\\
\hfuzz=500pt\includegraphics[width=1em]{element-mustset.pdf}~intervalCount & \hfuzz=500pt uint & \hfuzz=500pt count of intervals, count+1 points in time will generated\\
\hline
\end{tabularx}


\subsection{Irregular}
The points of the time series are given explicitly with \config{time}.


\keepXColumns
\begin{tabularx}{\textwidth}{N T A}
\hline
Name & Type & Annotation\\
\hline
\hfuzz=500pt\includegraphics[width=1em]{element-mustset-unbounded.pdf}~time & \hfuzz=500pt time & \hfuzz=500pt explicit list of points in time\\
\hline
\end{tabularx}


\subsection{Monthly}
If \config{monthMiddle} is set, time points are generated at mid of each month inclusively
the \config{monthStart} in \config{yearStart} and \config{monthEnd} in \config{yearEnd}.
Otherwise times are given at the first of each month and a time point after the last month.


\keepXColumns
\begin{tabularx}{\textwidth}{N T A}
\hline
Name & Type & Annotation\\
\hline
\hfuzz=500pt\includegraphics[width=1em]{element-mustset.pdf}~monthStart & \hfuzz=500pt uint & \hfuzz=500pt \\
\hfuzz=500pt\includegraphics[width=1em]{element-mustset.pdf}~yearStart & \hfuzz=500pt uint & \hfuzz=500pt \\
\hfuzz=500pt\includegraphics[width=1em]{element-mustset.pdf}~monthEnd & \hfuzz=500pt uint & \hfuzz=500pt \\
\hfuzz=500pt\includegraphics[width=1em]{element-mustset.pdf}~yearEnd & \hfuzz=500pt uint & \hfuzz=500pt \\
\hfuzz=500pt\includegraphics[width=1em]{element.pdf}~useMonthMiddle & \hfuzz=500pt boolean & \hfuzz=500pt time points are mid of months, otherwise the 1st of each month + a time point behind the last month\\
\hline
\end{tabularx}


\subsection{Yearly}
If \config{yearMiddle} is set, time points are generated at mid of each year inclusively \config{yearStart}
and \config{yearEnd}. Otherwise times are given at the first of each year and a time point after the last year.


\keepXColumns
\begin{tabularx}{\textwidth}{N T A}
\hline
Name & Type & Annotation\\
\hline
\hfuzz=500pt\includegraphics[width=1em]{element-mustset.pdf}~yearStart & \hfuzz=500pt uint & \hfuzz=500pt \\
\hfuzz=500pt\includegraphics[width=1em]{element-mustset.pdf}~yearEnd & \hfuzz=500pt uint & \hfuzz=500pt \\
\hfuzz=500pt\includegraphics[width=1em]{element.pdf}~useYearMiddle & \hfuzz=500pt boolean & \hfuzz=500pt time points are mid of years, otherwise the 1st of each year + a time point behind the last year\\
\hline
\end{tabularx}


\subsection{EveryMonth}
Generates a time series with monthly sampling. The first point in time will be \config{timeStart} and the following
points are generated for each month at the same day and time in month.
The last generated point in time will be less or equal \config{timeEnd}.


\keepXColumns
\begin{tabularx}{\textwidth}{N T A}
\hline
Name & Type & Annotation\\
\hline
\hfuzz=500pt\includegraphics[width=1em]{element-mustset.pdf}~timeStart & \hfuzz=500pt time & \hfuzz=500pt first point in time\\
\hfuzz=500pt\includegraphics[width=1em]{element-mustset.pdf}~timeEnd & \hfuzz=500pt time & \hfuzz=500pt last point in time will be less or equal timeEnd\\
\hline
\end{tabularx}


\subsection{EveryYear}
Generates a time series with yearly sampling. The first point in time will be \config{timeStart} and the following
points are generated for each year at the same day and time in year.
The last generated point in time will be less or equal \config{timeEnd}.


\keepXColumns
\begin{tabularx}{\textwidth}{N T A}
\hline
Name & Type & Annotation\\
\hline
\hfuzz=500pt\includegraphics[width=1em]{element-mustset.pdf}~timeStart & \hfuzz=500pt time & \hfuzz=500pt first point in time\\
\hfuzz=500pt\includegraphics[width=1em]{element-mustset.pdf}~timeEnd & \hfuzz=500pt time & \hfuzz=500pt last point in time will be less or equal timeEnd\\
\hline
\end{tabularx}


\subsection{Instrument}
Read a time series (epochs) from an \file{instrument file}{instrument}.
The time series can be restricted to the interval
starting from \config{timeStart} and before \config{timeEnd}.


\keepXColumns
\begin{tabularx}{\textwidth}{N T A}
\hline
Name & Type & Annotation\\
\hline
\hfuzz=500pt\includegraphics[width=1em]{element-mustset.pdf}~inputfileInstrument & \hfuzz=500pt filename & \hfuzz=500pt \\
\hfuzz=500pt\includegraphics[width=1em]{element.pdf}~timeStart & \hfuzz=500pt time & \hfuzz=500pt exclude peochs before this epoch\\
\hfuzz=500pt\includegraphics[width=1em]{element.pdf}~timeEnd & \hfuzz=500pt time & \hfuzz=500pt only epochs before this time are used\\
\hline
\end{tabularx}


\subsection{InstrumentArcIntervals}
Reconstruct a time series from an \file{instrument file}{instrument}.
The time series is the first epoch of each arc plus one time step beyond the last
epoch of the last arc (using median sampling).


\keepXColumns
\begin{tabularx}{\textwidth}{N T A}
\hline
Name & Type & Annotation\\
\hline
\hfuzz=500pt\includegraphics[width=1em]{element-mustset.pdf}~inputfileInstrument & \hfuzz=500pt filename & \hfuzz=500pt Must be regular. Time series is first epoch of each arc plus one time step extrapolated from last epoch of last arc.\\
\hline
\end{tabularx}


\subsection{Revolution}
Reads an \file{orbit file}{instrument} and create a time stamp for each ascending equator crossing.
The time series can be restricted to the interval
starting from \config{timeStart} and before \config{timeEnd}.


\keepXColumns
\begin{tabularx}{\textwidth}{N T A}
\hline
Name & Type & Annotation\\
\hline
\hfuzz=500pt\includegraphics[width=1em]{element-mustset.pdf}~inputfileOrbit & \hfuzz=500pt filename & \hfuzz=500pt \\
\hfuzz=500pt\includegraphics[width=1em]{element.pdf}~timeStart & \hfuzz=500pt time & \hfuzz=500pt exclude peochs before this epoch\\
\hfuzz=500pt\includegraphics[width=1em]{element.pdf}~timeEnd & \hfuzz=500pt time & \hfuzz=500pt only epochs before this time are used\\
\hline
\end{tabularx}


\subsection{Exclude}
In a first step a \configClass{timeSeries}{timeSeriesType} is generated.
In a second step all times are removed which are in range before or after \config{excludeMargin} seconds
of the times given by \configClass{timeSeriesExclude}{timeSeriesType}.


\keepXColumns
\begin{tabularx}{\textwidth}{N T A}
\hline
Name & Type & Annotation\\
\hline
\hfuzz=500pt\includegraphics[width=1em]{element-mustset-unbounded.pdf}~timeSeries & \hfuzz=500pt \hyperref[timeSeriesType]{timeSeries} & \hfuzz=500pt time series to be created\\
\hfuzz=500pt\includegraphics[width=1em]{element-mustset-unbounded.pdf}~timeSeriesExclude & \hfuzz=500pt \hyperref[timeSeriesType]{timeSeries} & \hfuzz=500pt exclude this time points from time series (within margin)\\
\hfuzz=500pt\includegraphics[width=1em]{element.pdf}~excludeMargin & \hfuzz=500pt double & \hfuzz=500pt on both sides [seconds]\\
\hline
\end{tabularx}


\subsection{Conditional}
Only times for which the \configClass{condition}{conditionType} is met are included in the time series.
The \config{variableLoopTime} is set to every time and the \configClass{condition}{conditionType} is evaluated.


\keepXColumns
\begin{tabularx}{\textwidth}{N T A}
\hline
Name & Type & Annotation\\
\hline
\hfuzz=500pt\includegraphics[width=1em]{element-mustset-unbounded.pdf}~timeSeries & \hfuzz=500pt \hyperref[timeSeriesType]{timeSeries} & \hfuzz=500pt only times for which condition is met will be included\\
\hfuzz=500pt\includegraphics[width=1em]{element.pdf}~variableLoopTime & \hfuzz=500pt string & \hfuzz=500pt variable with time of each loop\\
\hfuzz=500pt\includegraphics[width=1em]{element-mustset.pdf}~condition & \hfuzz=500pt \hyperref[conditionType]{condition} & \hfuzz=500pt test for each time\\
\hline
\end{tabularx}


\subsection{Interpolate}
Interpolates \config{nodeInterpolation} count points between
the given \configClass{timeSeries}{timeSeriesType} uniformly.


\keepXColumns
\begin{tabularx}{\textwidth}{N T A}
\hline
Name & Type & Annotation\\
\hline
\hfuzz=500pt\includegraphics[width=1em]{element-mustset-unbounded.pdf}~timeSeries & \hfuzz=500pt \hyperref[timeSeriesType]{timeSeries} & \hfuzz=500pt time series to be created\\
\hfuzz=500pt\includegraphics[width=1em]{element-mustset.pdf}~nodeInterpolation & \hfuzz=500pt uint & \hfuzz=500pt interpolates count points in each time interval given by the time series\\
\hline
\end{tabularx}

\clearpage
%==================================

\section{Troposphere}\label{troposphereType}
This class provides functions for calculating and estimating
the signal delay in the dry and wet atmosphere.


\subsection{ViennaMapping}\label{troposphereType:viennaMapping}

Tropospheric delays based on the Vienna Mapping Functions 3 (VMF3) model
(Landskron and Boehm 2017, DOI: \href{https://doi.org/10.1007/s00190-017-1066-2}{10.1007/s00190-017-1066-2}).

Hydrostatic and wet mapping function coefficients ($a_h$, $a_w$) and zenith delays (ZHD, ZWD) have to be provided
via \configFile{inputfileVmfCoefficients}{griddedDataTimeSeries}. This file can contain either station-specific data
(see \program{ViennaMappingFunctionStation2File}) or data on a regular global grid
(see \program{ViennaMappingFunctionGrid2File}). In the second case mapping coefficients and zenith delays are
interpolated to the requested coordinates. This includes a height correction that requires approximate meteorological
data provided via \configFile{inputfileGpt}{griddedData}.


\keepXColumns
\begin{tabularx}{\textwidth}{N T A}
\hline
Name & Type & Annotation\\
\hline
\hfuzz=500pt\includegraphics[width=1em]{element-mustset-unbounded.pdf}~inputfileVmfCoefficients & \hfuzz=500pt filename & \hfuzz=500pt ah, aw, zhd, zwd coefficients\\
\hfuzz=500pt\includegraphics[width=1em]{element-mustset.pdf}~inputfileGpt & \hfuzz=500pt filename & \hfuzz=500pt gridded GPT data\\
\hfuzz=500pt\includegraphics[width=1em]{element.pdf}~aHeight & \hfuzz=500pt double & \hfuzz=500pt parameter a (height correction)\\
\hfuzz=500pt\includegraphics[width=1em]{element.pdf}~bHeight & \hfuzz=500pt double & \hfuzz=500pt parameter b (height correction)\\
\hfuzz=500pt\includegraphics[width=1em]{element.pdf}~cHeight & \hfuzz=500pt double & \hfuzz=500pt parameter c (height correction)\\
\hline
\end{tabularx}


\subsection{GPT}\label{troposphereType:gpt}

Tropospheric delays based on the Global Pressure and Temperature 3 (GPT3) model
(Landskron and Boehm 2017, DOI: \href{https://doi.org/10.1007/s00190-017-1066-2}{10.1007/s00190-017-1066-2}).

It is an empirical model derived from the Vienna Mapping Functions 3
(VMF3, see \configClass{viennaMapping}{troposphereType:viennaMapping}) and thus does not require
additional mapping coefficients and zenith delay values.


\keepXColumns
\begin{tabularx}{\textwidth}{N T A}
\hline
Name & Type & Annotation\\
\hline
\hfuzz=500pt\includegraphics[width=1em]{element-mustset.pdf}~inputfileGpt & \hfuzz=500pt filename & \hfuzz=500pt gridded GPT data\\
\hfuzz=500pt\includegraphics[width=1em]{element.pdf}~aHeight & \hfuzz=500pt double & \hfuzz=500pt parameter a (height correction)\\
\hfuzz=500pt\includegraphics[width=1em]{element.pdf}~bHeight & \hfuzz=500pt double & \hfuzz=500pt parameter b (height correction)\\
\hfuzz=500pt\includegraphics[width=1em]{element.pdf}~cHeight & \hfuzz=500pt double & \hfuzz=500pt parameter c (height correction)\\
\hline
\end{tabularx}

\clearpage
%==================================


\end{document}
